% !TEX program = xelatex
% 中文自动译本：Mamoru Ishii and Takashi Hibiki,
% Thermo-Fluid Dynamics of Two-Phase Flow, Second Edition (Springer, 2011).
\documentclass[11pt,a4paper,openany,fontset=none]{ctexbook}
\usepackage[margin=2.55cm,headheight=15pt]{geometry}
\usepackage{microtype}
\usepackage{fancyhdr}
\usepackage{hyperref}
\usepackage{xurl}
\usepackage{titlesec}
\usepackage{amsmath,amssymb,amsfonts,bm,mathtools}
\usepackage[export]{adjustbox}
\usepackage{enumitem}
\usepackage{graphicx}
\usepackage{caption}
\usepackage{longtable,booktabs,array}
\usepackage{tikz}
\setmainfont[AutoFakeBold=2,AutoFakeSlant=.2]{Arial Unicode MS}
\setCJKmainfont{Songti SC}
\setCJKsansfont{Songti SC}
\setCJKmonofont{Songti SC}
\setlength{\parindent}{2em}
\setlength{\parskip}{0.65em}
\setlength{\emergencystretch}{3em}
\linespread{1.30}\selectfont
\xeCJKsetup{CJKglue = \hskip 0.075em plus 0.04em minus 0.03em}
\setlist[itemize]{leftmargin=2.8em,itemsep=0.3em,topsep=0.45em}
\titleformat{\chapter}[display]{\bfseries\Huge\centering}{第\thechapter 章}{0.75em}{}
\titleformat{\section}{\large\bfseries}{}{0pt}{}
\titleformat{\subsection}{\normalsize\bfseries}{}{0pt}{}
\titlespacing*{\section}{0pt}{1.45em}{0.7em}
\titlespacing*{\subsection}{0pt}{1.2em}{0.55em}
\pagestyle{fancy}
\fancyhf{}
\fancyhead[C]{\small\nouppercase{\leftmark}}
\fancyfoot[C]{\thepage}
\renewcommand{\headrulewidth}{0.3pt}
\hypersetup{
  colorlinks=true,
  linkcolor=black,
  urlcolor=blue,
  pdftitle={两相流热流体动力学},
  pdfauthor={Mamoru Ishii, Takashi Hibiki}
}
\graphicspath{{./}}
\title{两相流热流体动力学\\[0.5em]\Large 第二版·中文自动译文}
\author{Mamoru Ishii、Takashi Hibiki 著\\
\small 译文由自动化工具生成，供学习与检索使用}
\date{原著第二版：2011\\生成日期：\today}
% 统一译本版式：避免公式、图片和题注跨页或侵入页眉、页脚。
\usepackage{needspace}
\clubpenalty=10000
\widowpenalty=10000
\displaywidowpenalty=10000
\raggedbottom
\setlength{\abovedisplayskip}{0.8\baselineskip plus 0.2\baselineskip minus 0.1\baselineskip}
\setlength{\belowdisplayskip}{0.8\baselineskip plus 0.2\baselineskip minus 0.1\baselineskip}
\setlength{\abovedisplayshortskip}{0.55\baselineskip}
\setlength{\belowdisplayshortskip}{0.55\baselineskip}
\newenvironment{sourceblock}
  {\par\addvspace{0.65\baselineskip}\noindent\begin{minipage}{\linewidth}\centering}
  {\end{minipage}\par\addvspace{0.45\baselineskip}}
\captionsetup{font=small,labelfont=bf,justification=centering,singlelinecheck=false,skip=0.45\baselineskip}
\begin{document}
\frontmatter
\maketitle
\chapter*{说明}
\markboth{说明}{说明}
本文件是 Mamoru Ishii 与 Takashi Hibiki 所著
\emph{Thermo-Fluid Dynamics of Two-Phase Flow}（第二版）的中文自动译本。
正文按章节与小节的正常阅读顺序重新排版，不与原书逐页对齐。

原 PDF 使用 Euclid 专用数学字体，直接抽取时部分积分号、矩阵括号、上下标和希腊字母会错序。
为保证数学内容可靠，本译本把显示公式、插图和复杂表格以原 PDF 的矢量裁片嵌入；
正文中的变量和较短表达式保留在译文中。涉及推导、符号定义、图表数据与数值模型时，
请同时核对原著。原著版权归作者与 Springer Science+Business Media 所有。

\begingroup
\small
\tableofcontents
\endgroup
\clearpage
\makeatletter\@mainmattertrue\makeatother
\pagenumbering{arabic}
\begingroup\let\clearpage\relax
\chapter*{前言}\addcontentsline{toc}{chapter}{前言}
\markboth{前言}{前言}
\endgroup
本书旨在为认真参与该学科的研究生、科学家和实践工程师介绍两相流热流体动力学理论。本书可作为核工程、机械工程、化学工程等专业两相流研究生课程的教材，也可作为解决各技术领域多相流问题的研究人员和工程师的两相流公式的基础参考书。

单相流流体动力学和传热原理相对容易理解，但由于存在移动和可变形界面及其与两相的相互作用，两相流热流体动力学比单相流复杂一个数量级。然而，鉴于两相流在与核能、化学工程过程和先进传热系统相关的各种现代工程技术中的实际重要性，近年来人们在开发精确的通用两相公式、界面输运和界面结构的机械模型以及求解这些预测模型的计算方法方面做出了巨大的努力。

重点强调了推导两相流公式的合理方法，该公式代表了基本物理原理，例如散装流体和界面中各种传递机制的守恒定律和本构模型。提出了几种模型，例如基于界面显式处理的单相流模型的局部即时公式和基于各种平均方法的宏观连续介质公式。

通过界面面积输运方程动态模拟两相流中界面结构的变化。这是一种新方法，可以取代基于流态转变标准的静态且不准确的方法。详细讨论了界面动量传递模型，因为对于大多数两相流，热流体动力学由界面结构和界面动量传递主导。还讨论了其他一些必要的本构关系，例如湍流建模、瞬变力和升力。

石井守博士核工程学院

普渡大学 西拉斐特，印第安纳州，美国
响隆博士核工程学院

普渡大学 西拉斐特，印第安纳州，美国
August 2010

\chapter*{序}\addcontentsline{toc}{chapter}{序}
\markboth{序}{序}
两相流的热流体动力学在我们探索流体通过相、性质和结构的变化而变形的过程中向前迈出了重要一步。就像神话中的罗马神雅努斯（Janus）一样，他有两张脸，流体分为液体和气体，每种状态都可以单独充分理解，这对最敏锐和最有洞察力的科学头脑来说，在破译它们的动态纠缠时提出了重大挑战。

挑战部分源于可能遇到两相现象的巨大规模。在分子动力学的微观纳米尺度和深层建模假设以及宏观尺度的测量之间，存在着一个既宽广又模糊且难以捉摸的中观尺度。在这个尺度上，一切都处于永久的交换状态，赫拉克利特的流动状态，没有什么是一成不变的，只有通过牢牢掌握事物的基本原理才能获得知识。

主题源于作者自己对基础知识的牢固掌握。他们对两相原理的文献贡献反映了一种科学传统，即认为理论和实验的二元性与表象和现实的二元性一样重要。在这一点上，它与流体科学中的其他专题著作不同。例如，贯穿两相流的主要概念是界面速度。这一概念尤其需要在建模和测量方面不断改进。在中观尺度上，这产生了新的界面科学，除了问题的复杂性和结构的模糊性之外，还为创建优雅、简约的公式以及有前途的工程应用提供了充足的空间。

这是一个有趣的主题，正确理解它需要运用严格的高等数学工具。作者们对这个主题充满了极大的关心和智力热爱，并邀请了包括科学家、工程师、技术人员、教授和学生在内的包容性读者。

非常荣幸将两相流的热流体动力学纳入智能能源系统系列：纳瓦到太瓦。这是一部经得起时间考验的作品，其科学价值以及优雅和美学特征。

Lefteri H. Tsoukalas 博士核工程学院
普渡大学 西拉斐特，印第安纳州，美国
September 2005

\chapter*{致谢}\addcontentsline{toc}{chapter}{致谢}
\markboth{致谢}{致谢}
作者谨向为本书的编写做出贡献的人们表示诚挚的谢意。 N. Zuber 和 J. M. Delhaye 教授因其对多相流热流体动力学理论基本方法发展的早期投入和讨论而受到认可。我们要感谢美国核管理委员会的 F. Eltawila 博士长期以来对我们专注于两相流基础物理研究的支持。这项研究产生了本书中包含的一些重要成果。我们许多以前的学生，如 Q. Wu 教授、S. Kim 教授、X. Sun 教授和 Drs. X.Y。 Fu、S. Paranjape、B. Ozar、Y. Liu 和 D. Y. Lee 通过他们的博士学位做出了重大贡献。论文研究。现任博士学生 J. Schlegel、S. W. Chen、S. Rassame、S. Miwa 和​​ C. S. Brooks 非常感谢他们检查手稿。作者感谢 Lefteri Tsoukalas 教授邀请我们在新系列“智能能源系统：纳瓦到太瓦”下撰写这本书。

\chapter{引言}
\section*{1.1 问题的相关性}\phantomsection\addcontentsline{toc}{section}{1.1 问题的相关性}
本书旨在成为两相流热流体动力学理论的基础参考书。两相流或多相流的主题在各种工程系统的最佳设计和安全操作中变得越来越重要。然而，它绝不限于当今的现代工业技术，在许多生物系统和自然现象中都可以观察到多相流现象，这需要更好的理解。下面列出了一些重要的应用。

电力系统沸水和压水核反应堆；液态金属快中子增殖核反应堆；配有锅炉和蒸发器的传统发电厂；朗肯循环液态金属空间发电厂；磁流体发生器；地热能发电厂；内燃机；喷气发动机；液体或固体推进剂火箭；两相推进器等
传热系统热交换器；蒸发器；冷凝器；喷雾冷却塔；烘干机、冰箱和电子冷却系统；低温热交换器；薄膜冷却系统；热管；直接接触式热交换器；通过熔化热等进行蓄热。

工艺系统萃取和蒸馏装置；流化床；化学反应器；海水淡化系统；乳化剂；相分离器；雾化器；洗涤器；吸收器；均化器；搅拌反应器；多孔介质等
运输系统 气升泵；喷射器；天然气和石油混合物、浆料、纤维、小麦和粉状固体颗粒的管道运输；有气穴的泵和水翼；气动输送机；高速公路交通流量和控制等。

信息系统液氦的超流动性；导电或带电的液膜；液晶等
润滑系统 两相流润滑；通过低温技术等进行轴承冷却。环境控制空调；冰箱和冷却器；除尘器；污水处理厂；污染物分离器；空气污染控制；空间应用的生命支持系统等
地质气象现象沉积；土壤侵蚀和风吹移；海浪；雪花飘落；沙丘地层；雨滴的形成和运动；冰层；河流洪水、山体滑坡和雪崩；流冰覆盖的云层、河流或海洋的物理学；余波等
生物系统心血管系统；呼吸系统；胃肠道；血流（量;支气管血流和鼻腔血流；毛细管运输；通过排汗等调节体温。

可以说，上面列出的所有系统和组件本质上都受相同的质量、动量和能量传输物理定律的支配。显然，随着工程技术的快速进步，对感兴趣系统的逐步准确预测的需求不断增加。随着工程系统的规模变得越来越大，运行条件​​也被推向新的极限，准确理解控制这些多相流系统的物理原理对于安全且经济地运行是必不可少的。这意味着设计方法从完全基于静态实验相关性的设计方法转变为基于可以预测系统动态行为（例如瞬态响应和稳定性）的数学模型的设计方法。显然，多相流这一主题具有巨大的意义。

\section*{1.2 多相流特性}\phantomsection\addcontentsline{toc}{section}{1.2 多相流特性}
上面提到了多相流系统的许多例子。乍一看，各种两相流或多相流系统及其物理现象似乎没有什么共同点。因此，趋势是分析特定系统、组件或过程的问题，并开发通用性和适用性有限的系统特定模型和相关性。因此，对两相流热流体动力学的广泛理解只是缓慢发展，因此预测能力尚未达到单相流分析可用的水平。

工程系统的设计和预测其性能的能力取决于实验数据和概念数学模型的可用性，这些模型可用于以所需的准确度描述物理过程。至关重要的是，两相流的各种特性和物理现象应该在合理的基础上进行建模和公式化，并得到详细的科学实验的支持。在连续介质力学中，单相流的概念模型是根据描述质量、动量、能量、电荷等守恒定律的场方程来制定的。然后，这些场方程由热力学状态、应力、能量传递、化学反应等的适当本构方程来补充。这些本构方程指定了特定组成材料的热力学、输运和化学性质。

因此，预计多相流的概念模型也应根据适当的场和本构关系来制定。然而，多相流方程的推导比单相流方程的推导要复杂得多。两相流或多相流的复杂性质源于多个可变形和移动界面的存在以及随之而来的流体性质的显著不连续性和界面附近的复杂流场。通过关注界面结构和传递，我们注意到许多两相系统具有共同的几何结构。单相流根据流动结构可分为层流、过渡流和湍流。相比之下，两相流可以根据界面结构分为以下几种：

对于大多数两相流问题，基于具有显式移动界面的单相流公式的局部即时公式遇到了难以克服的数学和数值困难，因此它不是一种现实或实用的方法。这导致需要基于适当平均的宏观公式，通过有效消除界面不连续性给出两相流连续体公式。该公式的本质是通过级联建模方法考虑各种多尺度物理，将微观和介观尺度物理纳入宏观连续介质公式中。

上述讨论表明了在发展对多相流的广泛理解和分析这种流的通用方法时所遇到的困难的根源。两相流物理学本质上是多尺度的。在两相流公式和闭合关系中，有必要考虑不同尺度下各种物理的级联效应。在多相流中，至少有四种不同的尺度可能很重要。这些是1）系统尺度，2）连续统假设所需的宏观尺度，3）与局部结构相关的介观尺度，以及4）与精细结构和分子输运相关的微观尺度。在最高级别上，规模是系统瞬态和组件交互是主要关注点的系统。例如，核反应堆事故和瞬态分析需要专门的系统分析代码。在下一个层次上，研究了宏观物理，例如界面结构以及质量、动量和能量的传输。然而，描述守恒定律的多相流场方程需要额外的本体传递本构关系。这包括动量和能量的湍流效应以及质量、动量和能量传递的界面交换。这些是中尺度的物理现象，需要集中研究工作。由于界面传输速率可以被视为界面通量和可用界面面积的乘积，因此界面面积浓度的建模至关重要。在两相流分析中，
这可能表明，直接通过输运方程对界面面积浓度变化进行建模是比使用流态转变标准和界面面积浓度的流态相关本构关系的传统方法更好的方法。对于两相流的三维公式尤其如此。多相流中物理的下一个较低层次与局部微观现象有关，例如：壁成核或凝聚；气泡合并和破裂；以及夹带和沉积。

\section*{1.3 两相流的分类}\phantomsection\addcontentsline{toc}{section}{1.3 两相流的分类}
根据两相的组合以及界面结构，存在多种两相流。两相混合物的特征是存在一个或多个界面以及界面处的不连续性。根据两相的组合对两相混合物进行分类很容易，因为在标准条件下，我们只有三种物质状态，最多四种，即固相、液相、气相以及可能的等离子体（Pai，1972）。在这里，我们只考虑前三个阶段，因此我们有：

1、气固混合物； 2、气液混合物； 3、液固混合物； 4. 不相溶的液体混合物。显然，第四组不是两相流，然而，出于所有实际目的，它可以被视为两相混合物。

基于界面结构和各相的形貌分布的第二分类要困难得多，因为这些界面结构的变化是连续发生的。在这里，我们遵循 Wallis (1969)、Hewitt 和 Hall Taylor (1970)、Collier (1972)、Govier 和 Aziz (1972) 审查的标准流态以及 Zuber (1971)、Ishii (1971) 和 Kocamustafaogullari (1971) 的主要分类。根据界面的几何形状，两相流可分为三大类，即分离流、过渡或混合流和分散流，如表1-1所示。

\clearpage
\begin{center}\small 表 1-1。两相流的分类（Ishii，1975）\end{center}
\begin{sourceblock}
\includegraphics[page=22,pagebox=mediabox,trim={46.69bp 130.15bp 47.58bp 96.97bp},clip,width=\textwidth,max height=0.72\textheight]{Ishii.pdf}
\end{sourceblock}
根据界面的类型，分离流的类型可以分为平面流和准轴对称流，每种流又可以分为两种流态。因此，平面流包括薄膜流和分层流，而准轴对称流则包括环形流。

分散流的类别也可以分为几种类型。根据界面的几何形状，可以考虑球形、椭圆形、粒状颗粒等。然而，通过考虑分散体的相来细分分散流的类别更方便。因此，我们可以区分三种状态：气泡、液滴或雾以及颗粒流。在每个状态中，分散体的几何形状可以是球形、椭球形、扭曲形等。相和混合物组分之间的各种构型如表1-1所示。

如上所述，界面结构的变化逐渐发生，因此我们有第三类，其特征是同时存在分离流和分散流。当相变沿着通道进行时，液汽混合物经常发生这种转变。在这里，根据分散相来细分混合流的类别也更方便。因此，我们可以区分五种流态，即帽流、段塞流或搅动湍流、气泡环形流、气泡环形液滴流和夹带薄膜流。各相和混合物组分之间的各种配置如表 1-1 所示。

\begin{center}\small 图 1-1 和 1-2 分别显示了在直径为 25.4 毫米和 50.8 毫米的垂直管道中观察到的典型空气-水流动状态。左起第一、第二、第三、第四和第五图中的流态分别是气泡流、帽泡流、段塞流、搅动湍流和环形流。图1-3还显示了在间隙为10毫米、宽度为200毫米的垂直矩形通道中观察到的典型空气-水流态。左起第一、第二、第三和第四图中的流态分别是气泡流、帽泡流、搅动湍流和环形流。图 1-4 显示了用湍流水射流绝热模拟的倒环流，从大纵横比喷嘴向下喷出，封闭在气体环中（De Jarlais 等，1986）。左侧第一、第二、第三和第四幅图像分别表示对称射流不稳定性、蜿蜒射流不稳定性、大表面波和裙部形成以及高度湍流射流不稳定性。图 1-5 显示了入口液体速度 10.5 cm/s、入口气体速度 43.7 cm/s（氮气）、入口 Freon-113 温度 23 oC、壁温接近 200 oC 时倒环流的典型图像（Ishii 和 De Jarlais，1987）。通过在加热器石英管内同轴居中的薄壁管状喷嘴将测试流体引入测试部分芯，同时将蒸汽或气体引入液体喷嘴和加热石英管之间的环形间隙中，形成倒环形流。每幅图像的绝对垂直尺寸为 12.5 厘米。可视化标高从左图到右图较高。\end{center}
\begin{sourceblock}
\includegraphics[page=24,pagebox=mediabox,trim={50.00bp 371.76bp 44.37bp 55.00bp},clip,width=\textwidth,max height=0.72\textheight]{Ishii.pdf}
\par\vspace{0.45\baselineskip}
\small 图 1-1。在垂直 25.4 毫米直径管道中观察到的典型空气-水流图像
\end{sourceblock}
\begin{sourceblock}
\includegraphics[page=24,pagebox=mediabox,trim={50.00bp 209.40bp 44.37bp 308.38bp},clip,width=\textwidth,max height=0.72\textheight]{Ishii.pdf}
\par\vspace{0.45\baselineskip}
\small 图 1-2。在垂直 50.8 毫米直径管道中观察到的典型空气-水流图像
\end{sourceblock}
\begin{sourceblock}
\includegraphics[page=24,pagebox=mediabox,trim={50.00bp 76.62bp 44.37bp 470.74bp},clip,width=\textwidth,max height=0.72\textheight]{Ishii.pdf}
\par\vspace{0.45\baselineskip}
\small 图 1-3。在矩形通道中观察到的典型空气-水流图像
\end{sourceblock}
200毫米×10毫米

\begin{sourceblock}
\includegraphics[page=25,pagebox=mediabox,trim={50.00bp 322.08bp 44.37bp 55.00bp},clip,width=\textwidth,max height=0.72\textheight]{Ishii.pdf}
\par\vspace{0.45\baselineskip}
\small 图 1-4。模拟空气-水倒环流典型图像（并流向下）
\end{sourceblock}
\begin{sourceblock}
\includegraphics[page=25,pagebox=mediabox,trim={50.00bp 61.98bp 44.37bp 358.06bp},clip,width=\textwidth,max height=0.72\textheight]{Ishii.pdf}
\par\vspace{0.45\baselineskip}
\small 图1-5。逆环流轴向发展（并流向上流）
\end{sourceblock}
\section*{1.4 本书概要}\phantomsection\addcontentsline{toc}{section}{1.4 本书概要}
本书的目的是提出一个详细的两相流公式，该公式是使用机械建模合理导出和开发的。本书是作者（Ishii，1975）早期作品的延伸，特别强调使用界面区域传输方程对界面结构进行建模以及对流体动力学本构关系进行建模。然而，我们付出了特别的努力，使复杂的两相流物理和现象的公式和数学模型对于工程分析来说是现实和实用的。它重点详细讨论基于质量、动量和能量守恒定律的两相流各种数学模型的一般公式。

在第一部分中，两相流公式的基础被给出为基于单相流连续体公式和划分相的界面的显式存在的两相流的局部即时公式。讨论了封闭数学方程组的守恒方程、本构关系、界面处的跳跃条件以及界面处的特殊热力关系。

基于这个局部即时公式，在第二部分中，使用各种平均技术开发了宏观两相连续体公式，这些平均技术本质上是积分变换。时间平均的应用导致了通用的三维公式，有效地消除了界面不连续性并使两相共存连续体。平均微分平衡方程中的界面不连续性被界面传输源项和汇项取代。

三维两相流模型的详细信息将在第三部分中介绍。双流体模型、漂移-通量模型、界面区域传输和界面动量传递是讨论的主要主题。

\enlargethispage{2\baselineskip}
在第四部分中，根据二流体模型和漂移通量模型给出了更实用的两相流一维公式。还提出了考虑控制体积中的结构材料的双流体模型，即多孔介质公式。

\chapter{局部瞬时表述}
两相或两种不混溶混合物的独特特征是存在一个或多个分隔相或组分的界面。这种流动系统的例子可以在大量的工程系统以及各种各样的自然现象中找到。对两相系统的流动和传热过程的理解在核、机械和化学工程以及环境和医学科学中变得越来越重要。

在分析两相流时，显然我们首先遵循连续介质力学的标准方法。因此，两相流被认为是细分为单相区域的场，并且相之间的边界移动。标准微分平衡方程适用于具有适当跳跃和边界条件的每个子区域，以匹配这些微分方程在界面处的解。因此，理论上，可以用局部瞬时变量来表述两相流问题，即 ( ) F F t = x ,。该公式称为局部即时公式，以区别于基于各种平均方法的公式。

这种公式会导致多边界问题，由于场和边界条件的耦合，界面的位置未知。事实上，使用这种局部即时公式所遇到的数学困难可能是相当大的，并且在许多情况下，它们可能是无法克服的。然而，本地即时配方有两个基本重要性。首先重要的是直接应用来研究薄膜流、分层流、环形流和射流流等分离流，见表1-1。该公式可用于研究压降、传热、相变、界面的动态和稳定性以及临界热通量。除了上述应用之外，还可以使用此配方的重要示例
局部即时公式的第二个重要性是作为使用各种平均的宏观两相流模型的基础。当由界面界定的每个子区域可以被视为连续体时，局部瞬时公式在数学上是严格的。因此，两相流模型应通过适当的平均方法从该公式中导出。下面，提出并讨论基于局部瞬时变量的两相流系统的一般公式。这里应该指出的是，单相单组分流的平衡方程已经牢固地建立了一段时间（Truesdell 和 Toupin，1960；Bird 等，1960）。然而，包括状态方程在内的一般本构关系的公理化构造被专家们置于数学严谨之中（Coleman，1964；Bowen，1973；Truesdell，1969）。 Muller (1968) 也将类似的方法用于单相扩散混合物。

在详细推导和讨论局部瞬时公式之前，我们回顾一下与连续介质力学相关的数学物理方法。下图显示了用于获取物理系统数学模型的基本过程。

\begin{sourceblock}
\includegraphics[page=28,pagebox=mediabox,trim={48.21bp 247.21bp 49.22bp 326.71bp},clip,width=\textwidth,max height=0.72\textheight]{Ishii.pdf}
\end{sourceblock}
从图中可以看出，首先通过引入数学概念、一般公理和本构公理，将物理系统替换为数学系统。在连续介质力学中，它们对应于变量、场方程和本构方程，而在奇异表面，数学系统需要界面条件。后者不仅可以应用于两相之间的界面，还可以应用于限制系统的外部边界。从图中可以清楚地看出，连续介质公式由三个基本部分组成，即：场方程的推导、本构方程和界面条件。

现在让我们研究一下用于解决特定问题的基本过程。下图总结了标准方法。使用连续统公式，物理问题由理想化边界几何形状、边界条件、初始条件、场和

\begin{sourceblock}
\includegraphics[page=29,pagebox=mediabox,trim={65.31bp 488.11bp 66.56bp 54.67bp},clip,width=\textwidth,max height=0.72\textheight]{Ishii.pdf}
\end{sourceblock}
本构方程。显然，在两相流系统中，系统内存在可以用一般界面条件表示的界面。可以通过求解这些微分方程组以及一些理想化或简化假设来获得解。对于大多数具有实际意义的问题，实验数据也起着关键作用。首先，可以通过接受模型来获取实验数据，表明测量的可能性。解决方案与实验数据的比较为模型本身和各种假设提供了反馈。这种反馈将改进实验方法和解决方案。该模型的有效性通常通过解决一些简单的物理问题来证明。

单相热流体动力学中的连续介质方法被广泛接受，其有效性得到了充分证明。因此，如果两相系统中由界面界定的每个子区域都可以被视为连续体，则局部瞬时公式的有效性是显而易见的。通过接受这个假设，我们推导并讨论了场方程、本构关系和界面条件。由于界面是连续场的奇异情况，因此界面处有两种不同的条件。与场方程对应的界面处的平衡称为跳跃条件。与空间本构关系相对应的任何附加信息（在界面处也是必需的）称为界面边界条件。

\section*{2.1 单相流守恒方程}\phantomsection\addcontentsline{toc}{section}{2.1 单相流守恒方程}
\subsection*{2.1.1 一般平衡方程}\phantomsection\addcontentsline{toc}{subsection}{2.1.1 一般平衡方程}
微分平衡方程的推导如下图所示。一般积分平衡可以通过引入为单位质量定义的任何量 k ψ 的流体密度 k ρ 、流出量 k J 和体源 k φ 来编写。因此我们有

\begin{sourceblock}
\includegraphics[page=30,pagebox=mediabox,trim={156.03bp 492.97bp 157.10bp 54.67bp},clip,width=\textwidth,max height=0.72\textheight]{Ishii.pdf}
\end{sourceblock}
\begin{sourceblock}
\includegraphics[page=30,pagebox=mediabox,trim={45.00bp 437.26bp 44.37bp 191.96bp},clip,width=0.92\textwidth,max height=0.38\textheight]{Ishii.pdf}
\end{sourceblock}
其中 m V 是具有材料表面 m A 的材料体积。它指出 k k ρ ψ 以 m V 为单位的时间变化率等于通过 m A 的流入量加上体源。下标k表示第k相。如果式（2-1）中出现的函数足够光滑，使得物质坐标和空间坐标之间的雅可比变换存在，则可以得到平衡方程的熟悉的微分形式。这是通过使用雷诺传输定理（Aris，1962）来完成的，表示为

\begin{sourceblock}
\includegraphics[page=30,pagebox=mediabox,trim={45.00bp 288.14bp 44.37bp 336.73bp},clip,width=0.92\textwidth,max height=0.38\textheight]{Ishii.pdf}
\end{sourceblock}
其中 v k 表示流体粒子的速度。格林定理给出了一定体积和表面积分之间的变换，因此

\begin{sourceblock}
\includegraphics[page=30,pagebox=mediabox,trim={45.00bp 212.75bp 44.37bp 421.49bp},clip,width=0.92\textwidth,max height=0.38\textheight]{Ishii.pdf}
\end{sourceblock}
因此，由式（2-2）和（2-3）我们得到

\begin{sourceblock}
\includegraphics[page=30,pagebox=mediabox,trim={45.00bp 138.68bp 44.37bp 484.14bp},clip,width=0.92\textwidth,max height=0.38\textheight]{Ishii.pdf}
\end{sourceblock}
此外，我们注意到雷诺输运定理是莱布尼茨规则的一个特例，由下式给出

\begin{sourceblock}
\includegraphics[page=30,pagebox=mediabox,trim={45.00bp 55.17bp 44.37bp 569.80bp},clip,width=0.92\textwidth,max height=0.38\textheight]{Ishii.pdf}
\end{sourceblock}
\begin{sourceblock}
\includegraphics[page=31,pagebox=mediabox,trim={45.00bp 505.23bp 44.37bp 119.74bp},clip,width=0.92\textwidth,max height=0.38\textheight]{Ishii.pdf}
\end{sourceblock}
上式的第一项是每单位体积的数量的时间变化率，而第二项是每单位体积的对流率。右侧项代表表面通量和体积源。

\subsection*{2.1.2 守恒方程}\phantomsection\addcontentsline{toc}{subsection}{2.1.2 守恒方程}
连续性方程 质量守恒可以通过设置以微分形式表示

\begin{sourceblock}
\includegraphics[page=31,pagebox=mediabox,trim={45.00bp 334.83bp 44.37bp 306.28bp},clip,width=0.92\textwidth,max height=0.38\textheight]{Ishii.pdf}
\end{sourceblock}
因为相对于固定质量体积，不存在质量的表面和体积源。因此从一般平衡方程我们得到

\begin{sourceblock}
\includegraphics[page=31,pagebox=mediabox,trim={45.00bp 250.23bp 44.37bp 374.74bp},clip,width=0.92\textwidth,max height=0.38\textheight]{Ishii.pdf}
\end{sourceblock}
动量方程 引入表面应力张量 k T 和体力 k g 由式（2-6）可得动量守恒，因此有

\begin{sourceblock}
\includegraphics[page=31,pagebox=mediabox,trim={45.00bp 129.51bp 44.37bp 455.68bp},clip,width=0.92\textwidth,max height=0.38\textheight]{Ishii.pdf}
\end{sourceblock}
其中 I 是单位张量。在这里，我们将应力张量分解为压力项和粘性应力 k T。鉴于式(2-6)我们有

\begin{sourceblock}
\includegraphics[page=31,pagebox=mediabox,trim={45.00bp 50.14bp 44.37bp 580.79bp},clip,width=0.92\textwidth,max height=0.38\textheight]{Ishii.pdf}
\end{sourceblock}
角动量守恒如果我们假设没有体扭矩或力偶应力，那么所有扭矩都由表面应力和体力产生。在这种情况下，角动量守恒简化为

\begin{sourceblock}
\includegraphics[page=32,pagebox=mediabox,trim={45.00bp 519.90bp 44.37bp 121.69bp},clip,width=0.92\textwidth,max height=0.38\textheight]{Ishii.pdf}
\end{sourceblock}
+ T 表示转置应力张量。上述结果对于非极性流体是正确的，然而，对于极性流体我们应该引入固有角动量。在这种情况下，我们有一个微分角动量方程（Aris，1962）。

where k
能量守恒 能量平衡可以通过考虑流体的总能量来写出。因此我们设置

\begin{sourceblock}
\includegraphics[page=32,pagebox=mediabox,trim={45.00bp 302.74bp 44.37bp 273.50bp},clip,width=0.92\textwidth,max height=0.38\textheight]{Ishii.pdf}
\end{sourceblock}
其中 k u 、k q 和 k q 分别表示内能、热通量和物体加热。从这里可以看出，通量和体源都由热效应和机械效应组成。将式(2-12)代入式(2-6)可得总能量方程

\begin{sourceblock}
\includegraphics[page=32,pagebox=mediabox,trim={45.00bp 151.47bp 44.37bp 433.87bp},clip,width=0.92\textwidth,max height=0.38\textheight]{Ishii.pdf}
\end{sourceblock}
这四个局部方程，即式(2-8)、(2-10)、(2-11)和(2-13)，表达了质量、动量、角动量和能量守恒定律的四个基本物理定律。为了求解这些方程，必须指定通量和体源以及基本状态方程。这些将在宪法下讨论。除了这些本构关系之外，我们注意到上述方程还有几个重要的变换。很好的评论
材料导数的变换 根据连续性方程我们有

\begin{sourceblock}
\includegraphics[page=33,pagebox=mediabox,trim={45.00bp 489.20bp 44.37bp 132.64bp},clip,width=0.92\textwidth,max height=0.38\textheight]{Ishii.pdf}
\end{sourceblock}
这种特殊的时间导数称为物质导数或实质导数，因为它表示观察者随流体移动时相对于时间的变化率。

运动方程 通过使用上述变换，动量方程变为运动方程

\begin{sourceblock}
\includegraphics[page=33,pagebox=mediabox,trim={45.00bp 341.31bp 44.37bp 287.56bp},clip,width=0.92\textwidth,max height=0.38\textheight]{Ishii.pdf}
\end{sourceblock}
这里注意到 k k D Dt v 是流体加速度，因此运动方程表达了牛顿第二运动定律。

机械能方程 通过将运动方程除以速度，我们得到

\begin{sourceblock}
\includegraphics[page=33,pagebox=mediabox,trim={45.00bp 196.31bp 44.37bp 406.73bp},clip,width=0.92\textwidth,max height=0.38\textheight]{Ishii.pdf}
\end{sourceblock}
对于对称应力张量

\begin{sourceblock}
\includegraphics[page=33,pagebox=mediabox,trim={45.00bp 137.98bp 44.37bp 503.28bp},clip,width=0.92\textwidth,max height=0.38\textheight]{Ishii.pdf}
\end{sourceblock}
因此，式(2-16)可写为

\begin{sourceblock}
\includegraphics[page=34,pagebox=mediabox,trim={45.00bp 547.60bp 44.37bp 55.44bp},clip,width=0.92\textwidth,max height=0.38\textheight]{Ishii.pdf}
\end{sourceblock}
该机械能方程是标量方程，因此它仅代表动量方程控制的流体运动物理定律的一部分。

内能方程 从总能量方程中减去机械能方程，我们得到内能方程

\begin{sourceblock}
\includegraphics[page=34,pagebox=mediabox,trim={45.00bp 396.38bp 44.37bp 228.58bp},clip,width=0.92\textwidth,max height=0.38\textheight]{Ishii.pdf}
\end{sourceblock}
焓方程 通过引入由下式定义的焓

\begin{sourceblock}
\includegraphics[page=34,pagebox=mediabox,trim={45.00bp 302.74bp 44.37bp 324.73bp},clip,width=0.92\textwidth,max height=0.38\textheight]{Ishii.pdf}
\end{sourceblock}
焓能方程可得为

\begin{sourceblock}
\includegraphics[page=34,pagebox=mediabox,trim={45.00bp 230.42bp 44.37bp 394.54bp},clip,width=0.92\textwidth,max height=0.38\textheight]{Ishii.pdf}
\end{sourceblock}
\subsection*{2.1.3 熵不等式与本构律原理}\phantomsection\addcontentsline{toc}{subsection}{2.1.3 熵不等式与本构律原理}
宪法法是建立在三个不同基础上的。熵不等式可以被认为是对本构关系的限制，并且无论材料响应如何，它都应该由适当的本构方程来满足。除了熵不等式之外，还有一组重要的本构公理，它们一般理想化理论中包含的所有材料的响应和行为。决定论和局部作用原理在连续介质力学中经常使用。

上述本构关系的两个基础定义了理论中允许的本构方程的一般形式。本构关系的第三个基础是对材料响应的数学建模
现在我们继续讨论熵不等式。为了阐述热力学第二定律，有必要引入温度 k T 和比熵 k s 的概念。利用这些变量，第二定律可以写成不等式

\begin{sourceblock}
\includegraphics[page=35,pagebox=mediabox,trim={45.00bp 440.56bp 44.37bp 183.85bp},clip,width=0.92\textwidth,max height=0.38\textheight]{Ishii.pdf}
\end{sourceblock}
假设我们获得的变量足够平滑

\begin{sourceblock}
\includegraphics[page=35,pagebox=mediabox,trim={45.00bp 364.54bp 44.37bp 257.06bp},clip,width=0.92\textwidth,max height=0.38\textheight]{Ishii.pdf}
\end{sourceblock}
其中 k Δ 是每单位体积的熵产生率。在这种形式中，方程（2-23）似乎没有给出与守恒方程相关的明确的物理或数学意义，因为 k s 和 k T 与其他因变量的关系没有指定。换句话说，本构方程还没有给出。因此，不平等可以被视为对构成法的限制，而不是对过程本身的限制。

从上一节可以看出，因变量的数量超过了场方程的数量，因此具有适当边界条件的质量、动量、角动量和总能量的平衡方程不足以得出任何具体答案。因此，有必要用定义某种理想材料的各种本构方程来补充它们。因此，本构方程可以被视为特定材料组的数学模型。它们是根据表征材料特定行为的实验数据以及控制它们的假设原理制定的。

从物理意义来看，可以将各种本构方程分为三组：

1.力学本构方程； 2. 能量本构方程； 3.本构状态方程。

正如已经解释的那样，本构律一般形式的推导遵循熵不等式、决定论、框架冷漠和局部作用等假设原理。其中最重要的是决定论原则，它粗略地说明了过去历史对当前状态的可预测性。材料框架无差异原理是材料的响应独立于框架或观察者这一思想的实现。而熵不等式要求本构方程无条件满足不等式（2-23）。进一步的限制，例如变量的等价性，经常被引入到通量的本构方程中，即 k T 和 k q。

\subsection*{2.1.4 本构方程}\phantomsection\addcontentsline{toc}{subsection}{2.1.4 本构方程}
我们将注意力集中在流体力学中最重要且广泛使用的特定类型的材料和本构方程。

基本状态方程 热力学均匀流体的基本状态方程的标准形式由内能与熵和密度相关的函数给出，因此我们有

\begin{sourceblock}
\includegraphics[page=36,pagebox=mediabox,trim={45.00bp 208.00bp 44.37bp 433.29bp},clip,width=0.92\textwidth,max height=0.38\textheight]{Ishii.pdf}
\end{sourceblock}
温度和热力学压力由下式给出

\begin{sourceblock}
\includegraphics[page=36,pagebox=mediabox,trim={45.00bp 128.88bp 44.37bp 490.79bp},clip,width=0.92\textwidth,max height=0.38\textheight]{Ishii.pdf}
\end{sourceblock}
因此，在微分形式中，基本状态方程变为

\begin{sourceblock}
\includegraphics[page=36,pagebox=mediabox,trim={45.00bp 51.64bp 44.37bp 569.97bp},clip,width=0.92\textwidth,max height=0.38\textheight]{Ishii.pdf}
\end{sourceblock}
吉布斯自由能、焓和亥姆霍兹自由能函数定义为

\begin{sourceblock}
\includegraphics[page=37,pagebox=mediabox,trim={45.00bp 533.50bp 44.37bp 93.92bp},clip,width=0.92\textwidth,max height=0.38\textheight]{Ishii.pdf}
\end{sourceblock}
\begin{sourceblock}
\includegraphics[page=37,pagebox=mediabox,trim={45.00bp 476.49bp 44.37bp 150.92bp},clip,width=0.92\textwidth,max height=0.38\textheight]{Ishii.pdf}
\end{sourceblock}
\begin{sourceblock}
\includegraphics[page=37,pagebox=mediabox,trim={45.00bp 433.83bp 44.37bp 207.28bp},clip,width=0.92\textwidth,max height=0.38\textheight]{Ishii.pdf}
\end{sourceblock}
分别。这些可以被视为勒让德变换*（Callen，1960），它将自变量从原始变量更改为一阶导数。因此在我们的例子中我们有

\begin{sourceblock}
\includegraphics[page=37,pagebox=mediabox,trim={45.00bp 350.98bp 44.37bp 292.52bp},clip,width=0.92\textwidth,max height=0.38\textheight]{Ishii.pdf}
\end{sourceblock}
* 如果我们有

\begin{sourceblock}
\includegraphics[page=37,pagebox=mediabox,trim={45.00bp 261.70bp 44.37bp 362.57bp},clip,width=0.92\textwidth,max height=0.38\textheight]{Ishii.pdf}
\end{sourceblock}
那么勒让德变换由下式给出

\begin{sourceblock}
\includegraphics[page=37,pagebox=mediabox,trim={45.00bp 188.43bp 44.37bp 438.28bp},clip,width=0.92\textwidth,max height=0.38\textheight]{Ishii.pdf}
\end{sourceblock}
在这种情况下 Z 变为

\begin{sourceblock}
\includegraphics[page=37,pagebox=mediabox,trim={45.00bp 133.17bp 44.37bp 508.67bp},clip,width=0.92\textwidth,max height=0.38\textheight]{Ishii.pdf}
\end{sourceblock}
因此，我们有

\begin{sourceblock}
\includegraphics[page=37,pagebox=mediabox,trim={45.00bp 58.60bp 44.37bp 568.26bp},clip,width=0.92\textwidth,max height=0.38\textheight]{Ishii.pdf}
\end{sourceblock}
\begin{sourceblock}
\includegraphics[page=38,pagebox=mediabox,trim={45.00bp 584.20bp 44.37bp 59.33bp},clip,width=0.92\textwidth,max height=0.38\textheight]{Ishii.pdf}
\end{sourceblock}
\begin{sourceblock}
\includegraphics[page=38,pagebox=mediabox,trim={45.00bp 538.18bp 44.37bp 103.08bp},clip,width=0.92\textwidth,max height=0.38\textheight]{Ishii.pdf}
\end{sourceblock}
这也是一个基本的状态方程。

由于温度和压力是一阶导数

基本状态方程的 k u 方程（2-24）可以用热和热量状态方程的组合来代替（Bird et al., 1960; Callen, 1960），由下式给出

\begin{sourceblock}
\includegraphics[page=38,pagebox=mediabox,trim={45.00bp 427.18bp 44.37bp 214.08bp},clip,width=0.92\textwidth,max height=0.38\textheight]{Ishii.pdf}
\end{sourceblock}
\begin{sourceblock}
\includegraphics[page=38,pagebox=mediabox,trim={45.00bp 381.10bp 44.37bp 260.16bp},clip,width=0.92\textwidth,max height=0.38\textheight]{Ishii.pdf}
\end{sourceblock}
温度和压力是容易测量的量；因此，从实验中获得这两个状态方程并在公式中使用它们更为实用。这些状态方程的一个简单例子是不可压缩流体

\begin{sourceblock}
\includegraphics[page=38,pagebox=mediabox,trim={45.00bp 264.34bp 44.37bp 354.71bp},clip,width=0.92\textwidth,max height=0.38\textheight]{Ishii.pdf}
\end{sourceblock}
在这种情况下，压力无法用热力学定义，因此我们使用流体动压，它是法向应力的平均值。此外，理想气体具有状态方程

\begin{sourceblock}
\includegraphics[page=38,pagebox=mediabox,trim={45.00bp 159.88bp 44.37bp 459.10bp},clip,width=0.92\textwidth,max height=0.38\textheight]{Ishii.pdf}
\end{sourceblock}
其中 M R 是理想气体常数除以分子量。力学本构方程 最简单的流变本构方程是无粘性流体的方程，表示为

\begin{sourceblock}
\includegraphics[page=38,pagebox=mediabox,trim={45.00bp 50.14bp 44.37bp 591.26bp},clip,width=0.92\textwidth,max height=0.38\textheight]{Ishii.pdf}
\end{sourceblock}
\begin{sourceblock}
\includegraphics[page=39,pagebox=mediabox,trim={45.00bp 532.51bp 44.37bp 93.47bp},clip,width=0.92\textwidth,max height=0.38\textheight]{Ishii.pdf}
\end{sourceblock}
其中 k μ 和 k λ 分别是第 k 相的粘度和体积粘度。

体力源自外力场以及与周围物体或流体粒子的相互作用力。力的来源是牛顿引力、静电力和电磁力。如果相互相互作用的力很重要，则体积力可能不会被视为仅是自变量 x 和 t 的函数。在这种情况下，就不能应用本地操作的原则。然而，对于大多数问题，与引力场力 g 相比，这些相互作用力可以忽略不计。因此我们有

\begin{sourceblock}
\includegraphics[page=39,pagebox=mediabox,trim={45.00bp 359.48bp 44.37bp 283.30bp},clip,width=0.92\textwidth,max height=0.38\textheight]{Ishii.pdf}
\end{sourceblock}
能量本构方程 接触传热由热通量矢量 k q 表示，其本构方程规定了接触能量传递的性质和机制。大多数流体遵循广义傅里叶热传导定律，其形式为

\begin{sourceblock}
\includegraphics[page=39,pagebox=mediabox,trim={45.00bp 240.63bp 44.37bp 400.48bp},clip,width=0.92\textwidth,max height=0.38\textheight]{Ishii.pdf}
\end{sourceblock}
二阶张量 k K 是考虑材料各向异性的电导率张量。对于各向同性流体，本构关系可以用单个系数表示为

\begin{sourceblock}
\includegraphics[page=39,pagebox=mediabox,trim={45.00bp 157.78bp 44.37bp 483.47bp},clip,width=0.92\textwidth,max height=0.38\textheight]{Ishii.pdf}
\end{sourceblock}
这是傅里叶热传导定律和标量的标准形式

k K 称为热导率。

身体热量 k q 来自外部能源和相互的相互作用。能量可以通过核裂变产生，并可以通过辐射、导电和磁感应从远处转移。流体两部分之间的相互辐射最好地说明了能量的相互作用或传递。大多数情况下这些

\begin{sourceblock}
\includegraphics[page=40,pagebox=mediabox,trim={45.00bp 508.83bp 44.37bp 132.38bp},clip,width=0.92\textwidth,max height=0.38\textheight]{Ishii.pdf}
\end{sourceblock}
可以广泛应用于解决实际问题。

最后，我们注意到熵不等式要求传输系数 k μ 、k λ 和 k K 为非负值。因此，粘性应力充当流体运动的阻力，并且不做功。此外，热量仅沿温度较高的方向流动至温度较低的方向。

\section*{2.2 界面平衡和边界条件}\phantomsection\addcontentsline{toc}{section}{2.2 界面平衡和边界条件}
\subsection*{2.2.1 界面平衡（跳跃条件）}\phantomsection\addcontentsline{toc}{subsection}{2.2.1 界面平衡（跳跃条件）}
前面几节中导出的标准微分平衡方程可以应用于直到界面的每个相，但不能应用于整个界面。应在界面处使用特定形式的平衡方程，以考虑奇异特征，即各种变量的急剧变化（或不连续性）。通过将界面视为单个表面，在该表面上流体密度、能量和速度会遭受跳跃不连续性，因此开发了所谓的跳跃条件。这些条件指定了通过界面的质量、动量和能量的交换，并作为两相之间的匹配条件，因此它们在两相流分析中是必不可少的。此外，由于单相流问题中的固体边界也构成界面，因此各种简化形式的跳跃条件被频繁使用而没有引起太多注意。由于其重要性，我们详细讨论跳跃条件的推导和物理意义。

没有任何表面特性的界面跳跃条件首先由 Kotchine (1926) 作为冲击不连续处的动力学相容条件提出一般形式，尽管不同的作者之前已经开发了特殊情况。它可以通过积分平衡方程推导出来，假设它适用于具有不连续表面的材料体积。许多作者（Scriven，1960；Slattery，1964；Standart，1964；Delhaye，1968；Kelly，1964）都试图扩展 Kotchine 定理。这些包括引入界面线通量，例如表面张力、粘性应力和热通量或表面材料特性。该问题有多种解决方法以及上述结果
由于在将时间平均应用于两相流场时考虑有限厚度界面会很方便，因此我们基于控制体积分析推导了通用界面平衡方程。假设界面的位置由数学曲面 ( ) 0 f t = x , 给出。界面对物理变量的影响仅限于表面附近，影响范围由厚度为 δ 的薄层给出，表面每侧有 1 δ 和 2 δ。让我们用 i A 表示表面上的简单连通区域，然后控制体积以垂直于 i A 的表面 i Σ 为界，并且 i A 和 i Σ 的交点是一条闭合曲线 i C。因此 i Σ 形成一个宽度为 δ 的环，而每侧界面区域的边界由下式表示
1A和2A。我们的控制体积 i V 由 i Σ 、 1 A 和 2 A 组成。由于假设 δ 的大小远小于沿表面 i A 的特征尺寸，因此我们将

\begin{sourceblock}
\includegraphics[page=41,pagebox=mediabox,trim={45.00bp 327.78bp 44.37bp 292.85bp},clip,width=0.92\textwidth,max height=0.38\textheight]{Ishii.pdf}
\end{sourceblock}
其中 1 n 和 2 n 分别是第 1 相和第 2 相本体流体的向外单位法向矢量。垂直于 i Σ 的向外单位向量用 N 表示，则控制体积 i V 的扩展一般积分平衡方程为

\begin{sourceblock}
\includegraphics[page=41,pagebox=mediabox,trim={45.00bp 186.63bp 44.37bp 411.92bp},clip,width=0.92\textwidth,max height=0.38\textheight]{Ishii.pdf}
\end{sourceblock}
右侧的前两个积分考虑了来自表面 1 A 、 2 A 和 i Σ 的通量。为了将体积积分平衡减少到 i A 上的表面积分平衡，我们应该引入下面定义的表面属性。

\begin{sourceblock}
\includegraphics[page=41,pagebox=mediabox,trim={45.00bp 67.00bp 44.37bp 540.79bp},clip,width=0.92\textwidth,max height=0.38\textheight]{Ishii.pdf}
\end{sourceblock}
\begin{sourceblock}
\includegraphics[page=42,pagebox=mediabox,trim={50.00bp 397.44bp 44.37bp 55.00bp},clip,width=\textwidth,max height=0.72\textheight]{Ishii.pdf}
\par\vspace{0.45\baselineskip}
\small 图2-1。接口（石井，1975）
\end{sourceblock}
其中平均密度 s ρ 和每单位表面积的平均密度 a ρ 定义为

\begin{sourceblock}
\includegraphics[page=42,pagebox=mediabox,trim={45.00bp 290.02bp 44.37bp 340.58bp},clip,width=0.92\textwidth,max height=0.38\textheight]{Ishii.pdf}
\end{sourceblock}
\begin{sourceblock}
\includegraphics[page=42,pagebox=mediabox,trim={45.00bp 260.91bp 44.37bp 345.59bp},clip,width=0.92\textwidth,max height=0.38\textheight]{Ishii.pdf}
\end{sourceblock}
\begin{sourceblock}
\includegraphics[page=42,pagebox=mediabox,trim={45.00bp 222.04bp 44.37bp 408.64bp},clip,width=0.92\textwidth,max height=0.38\textheight]{Ishii.pdf}
\end{sourceblock}
and

\begin{sourceblock}
\includegraphics[page=42,pagebox=mediabox,trim={45.00bp 154.00bp 44.37bp 476.62bp},clip,width=0.92\textwidth,max height=0.38\textheight]{Ishii.pdf}
\end{sourceblock}
这里的符号是这样的：每单位界面质量和每单位表面积的量分别由下标 s 和 a 表示。

控制表观速度可以分为切向分量和法向分量，因此

\begin{sourceblock}
\includegraphics[page=42,pagebox=mediabox,trim={45.00bp 59.44bp 44.37bp 584.64bp},clip,width=0.92\textwidth,max height=0.38\textheight]{Ishii.pdf}
\end{sourceblock}
\begin{sourceblock}
\includegraphics[page=43,pagebox=mediabox,trim={45.00bp 512.93bp 44.37bp 80.83bp},clip,width=0.92\textwidth,max height=0.38\textheight]{Ishii.pdf}
\end{sourceblock}
因此，法向分量是表面位移速度，切向分量由平均切向粒子速度 ts v 给出。由于单位向量 N 位于切平面并垂直于 i C，我们有

\begin{sourceblock}
\includegraphics[page=43,pagebox=mediabox,trim={45.00bp 416.92bp 44.37bp 224.31bp},clip,width=0.92\textwidth,max height=0.38\textheight]{Ishii.pdf}
\end{sourceblock}
因此，由式（2-45）和（2-51）我们得到

\begin{sourceblock}
\includegraphics[page=43,pagebox=mediabox,trim={45.00bp 349.12bp 44.37bp 280.36bp},clip,width=0.92\textwidth,max height=0.38\textheight]{Ishii.pdf}
\end{sourceblock}
and

\begin{sourceblock}
\includegraphics[page=43,pagebox=mediabox,trim={45.00bp 282.34bp 44.37bp 348.34bp},clip,width=0.92\textwidth,max height=0.38\textheight]{Ishii.pdf}
\end{sourceblock}
鉴于方程（2-44）和（2-53），我们定义沿 i C 的平均线流出量为

\begin{sourceblock}
\includegraphics[page=43,pagebox=mediabox,trim={45.00bp 201.34bp 44.37bp 429.34bp},clip,width=0.92\textwidth,max height=0.38\textheight]{Ishii.pdf}
\end{sourceblock}
使用上述定义，界面区域的积分平衡变为

\begin{sourceblock}
\includegraphics[page=43,pagebox=mediabox,trim={45.00bp 55.53bp 44.37bp 512.30bp},clip,width=0.92\textwidth,max height=0.38\textheight]{Ishii.pdf}
\end{sourceblock}
与场方程的推导一样，这里我们需要两个数学变换，即表面输运定理和表面格林定理（Weatherburn，1927；McConnell，1957；Aris 1962）。表面传输定理由下式给出

\begin{sourceblock}
\includegraphics[page=44,pagebox=mediabox,trim={45.00bp 506.38bp 44.37bp 122.68bp},clip,width=0.92\textwidth,max height=0.38\textheight]{Ishii.pdf}
\end{sourceblock}
其中 s d dt 表示表观速度 i v 的对流导数
由式(2-50)定义，s∇表示表面散度算子。表面格林定理由下式给出

\begin{sourceblock}
\includegraphics[page=44,pagebox=mediabox,trim={45.00bp 379.60bp 44.37bp 253.94bp},clip,width=0.92\textwidth,max height=0.38\textheight]{Ishii.pdf}
\end{sourceblock}
这里，A αβ 、g ln 、nt α 和 ( ) β 分别表示表面度量张量、空间度量张量、混合张量和表面协变导数（Aris，1962）。

空间坐标中的表面通量 a J 用 l a J i 表示，它代表质量和能量平衡的空间矢量以及动量平衡的空间张量。下面给出张量符号的基本概念。首先，笛卡尔空间坐标表示为
( ) 1 2 3 y y , , 和 ( ) 1 2 3 x x x , , 的一般坐标，则空间度量张量定义为

\begin{sourceblock}
\includegraphics[page=44,pagebox=mediabox,trim={45.00bp 201.69bp 44.37bp 423.27bp},clip,width=0.92\textwidth,max height=0.38\textheight]{Ishii.pdf}
\end{sourceblock}
它涉及这两个系统之间的无穷小坐标元素的距离。如图2-2所示，如果笛卡尔坐标 k y 给出表面坐标为 ( ) 1 2 u u 的表面上的一点，则 ( ) 1 2 k k y y u u = ,，则表面度量张量定义为

\begin{sourceblock}
\includegraphics[page=44,pagebox=mediabox,trim={45.00bp 89.67bp 44.37bp 510.37bp},clip,width=0.92\textwidth,max height=0.38\textheight]{Ishii.pdf}
\end{sourceblock}
小距离 ds 由下式给出

\begin{sourceblock}
\includegraphics[page=45,pagebox=mediabox,trim={50.00bp 394.26bp 44.37bp 55.00bp},clip,width=\textwidth,max height=0.72\textheight]{Ishii.pdf}
\par\vspace{0.45\baselineskip}
\small 图2-2。笛卡尔坐标与曲面坐标的关系
\end{sourceblock}
\begin{sourceblock}
\includegraphics[page=45,pagebox=mediabox,trim={45.00bp 320.18bp 44.37bp 319.34bp},clip,width=0.92\textwidth,max height=0.38\textheight]{Ishii.pdf}
\end{sourceblock}
\begin{sourceblock}
\includegraphics[page=45,pagebox=mediabox,trim={45.00bp 232.21bp 44.37bp 368.19bp},clip,width=0.92\textwidth,max height=0.38\textheight]{Ishii.pdf}
\end{sourceblock}
协方差表面导数 ( ) β，与空间导数类似
但它也考虑了弯曲坐标的影响。此外，如果 a ⋅ N J 仅具有切向分量（如表面张力的情况）

\begin{sourceblock}
\includegraphics[page=45,pagebox=mediabox,trim={45.00bp 134.53bp 44.37bp 488.82bp},clip,width=0.92\textwidth,max height=0.38\textheight]{Ishii.pdf}
\end{sourceblock}
矢量符号。值得注意的是，对于动量传递，主要的界面动量通量是各向同性表面张力 σ。然后，
a A αβ αβ σ = J。在这种情况下，表面通量贡献如下

\begin{sourceblock}
\includegraphics[page=45,pagebox=mediabox,trim={45.00bp 59.89bp 44.37bp 580.45bp},clip,width=0.92\textwidth,max height=0.38\textheight]{Ishii.pdf}
\end{sourceblock}
由于我们假设δ足够小，因此表面1 A 和2 A 在几何上与A i 重合。因此，式（2-55）简化为

\begin{sourceblock}
\includegraphics[page=46,pagebox=mediabox,trim={45.00bp 433.84bp 44.37bp 135.71bp},clip,width=0.92\textwidth,max height=0.38\textheight]{Ishii.pdf}
\end{sourceblock}
该平衡方程适用于界面的任意部分

2 i A δ >>，从而得到微分平衡方程

\begin{sourceblock}
\includegraphics[page=46,pagebox=mediabox,trim={45.00bp 294.49bp 44.37bp 280.31bp},clip,width=0.92\textwidth,max height=0.38\textheight]{Ishii.pdf}
\end{sourceblock}
我们注意到这里的结果与 Delhaye (1974) 得出的结果具有完全相同的形式，尽管使用的方法和表观速度 i v 的定义不同。让我们将表面积和每个表面积的源定义为

\begin{sourceblock}
\includegraphics[page=46,pagebox=mediabox,trim={45.00bp 198.88bp 44.37bp 442.32bp},clip,width=0.92\textwidth,max height=0.38\textheight]{Ishii.pdf}
\end{sourceblock}
and

\begin{sourceblock}
\includegraphics[page=46,pagebox=mediabox,trim={45.00bp 143.92bp 44.37bp 497.28bp},clip,width=0.92\textwidth,max height=0.38\textheight]{Ishii.pdf}
\end{sourceblock}
那么表面平衡方程变为

\begin{sourceblock}
\includegraphics[page=46,pagebox=mediabox,trim={45.00bp 50.89bp 44.37bp 556.53bp},clip,width=0.92\textwidth,max height=0.38\textheight]{Ishii.pdf}
\end{sourceblock}
为了获得量的界面跳跃的更简单的表达式，我们做出进一步的假设，这些假设与我们给出的薄层假设一致

\begin{sourceblock}
\includegraphics[page=47,pagebox=mediabox,trim={45.00bp 454.90bp 44.37bp 186.33bp},clip,width=0.92\textwidth,max height=0.38\textheight]{Ishii.pdf}
\end{sourceblock}
首先，界面a ρ 的质量密度小到可以忽略不计，因此其动量和机械能也可以忽略不计。其次，忽略沿线的所有分子扩散通量，即没有表面粘性应力或表面热通量。此外，所有表面源都被忽略，即除了重力之外没有特定的体积力并且没有辐射效应。

热力学张力以及界面能都包含在下面的分析中，因此从决定论原理我们应该假设表面状态方程的存在。在这些假设下我们得到
界面质量平衡

\begin{sourceblock}
\includegraphics[page=47,pagebox=mediabox,trim={45.00bp 239.19bp 44.37bp 390.02bp},clip,width=0.92\textwidth,max height=0.38\textheight]{Ishii.pdf}
\end{sourceblock}
通过将第 k 相的界面质量流出定义为

\begin{sourceblock}
\includegraphics[page=47,pagebox=mediabox,trim={45.00bp 181.90bp 44.37bp 459.35bp},clip,width=0.92\textwidth,max height=0.38\textheight]{Ishii.pdf}
\end{sourceblock}
我们从方程（2-69）得到

\begin{sourceblock}
\includegraphics[page=47,pagebox=mediabox,trim={45.00bp 107.19bp 44.37bp 522.01bp},clip,width=0.92\textwidth,max height=0.38\textheight]{Ishii.pdf}
\end{sourceblock}
该方程简单地表明界面处没有质量容量，因此相变是两相之间的纯粹质量交换。

界面动量平衡

\begin{sourceblock}
\includegraphics[page=48,pagebox=mediabox,trim={45.00bp 543.27bp 44.37bp 85.93bp},clip,width=0.92\textwidth,max height=0.38\textheight]{Ishii.pdf}
\end{sourceblock}
方程（2-72）是来自本体流体的动量通量与界面张力之间的平衡。

界面能量平衡 将单位表面积的界面能量a u 代入a ψ，由式（2-67）得到

\begin{sourceblock}
\includegraphics[page=48,pagebox=mediabox,trim={45.00bp 349.56bp 44.37bp 220.94bp},clip,width=0.92\textwidth,max height=0.38\textheight]{Ishii.pdf}
\end{sourceblock}
左侧代表表面能的变化率，而右侧则说明两侧主体的能量转移以及表面张力所做的功。

\subsection*{2.2.2 界面边界条件}\phantomsection\addcontentsline{toc}{subsection}{2.2.2 界面边界条件}
与三维场方程的情况一样，表面平衡方程应该由各种本构关系来补充。为了建立决定论原理，首先我们引入一个简单的状态方程。由于界面的质量可以忽略不计，我们有

\begin{sourceblock}
\includegraphics[page=48,pagebox=mediabox,trim={45.00bp 173.98bp 44.37bp 469.52bp},clip,width=0.92\textwidth,max height=0.38\textheight]{Ishii.pdf}
\end{sourceblock}
其中 a u 和 a s 分别是单位表面积的比内能和比熵。

热力学张力由下式给出

\begin{sourceblock}
\includegraphics[page=48,pagebox=mediabox,trim={45.00bp 90.88bp 44.37bp 550.36bp},clip,width=0.92\textwidth,max height=0.38\textheight]{Ishii.pdf}
\end{sourceblock}
其中温度 i T 定义为

\begin{sourceblock}
\includegraphics[page=49,pagebox=mediabox,trim={45.00bp 566.80bp 44.37bp 57.56bp},clip,width=0.92\textwidth,max height=0.38\textheight]{Ishii.pdf}
\end{sourceblock}
因此，以微分形式，式(2-74)变为

\begin{sourceblock}
\includegraphics[page=49,pagebox=mediabox,trim={45.00bp 508.24bp 44.37bp 133.07bp},clip,width=0.92\textwidth,max height=0.38\textheight]{Ishii.pdf}
\end{sourceblock}
吉布斯-迪昂关系由下式给出

\begin{sourceblock}
\includegraphics[page=49,pagebox=mediabox,trim={45.00bp 449.14bp 44.37bp 192.17bp},clip,width=0.92\textwidth,max height=0.38\textheight]{Ishii.pdf}
\end{sourceblock}
界面焓定义为

\begin{sourceblock}
\includegraphics[page=49,pagebox=mediabox,trim={45.00bp 390.88bp 44.37bp 250.36bp},clip,width=0.92\textwidth,max height=0.38\textheight]{Ishii.pdf}
\end{sourceblock}
由式(2-78)可得

\begin{sourceblock}
\includegraphics[page=49,pagebox=mediabox,trim={45.00bp 319.54bp 44.37bp 304.78bp},clip,width=0.92\textwidth,max height=0.38\textheight]{Ishii.pdf}
\end{sourceblock}
因此，从式（2-77）和（2-80）我们得到

\begin{sourceblock}
\includegraphics[page=49,pagebox=mediabox,trim={45.00bp 243.51bp 44.37bp 378.03bp},clip,width=0.92\textwidth,max height=0.38\textheight]{Ishii.pdf}
\end{sourceblock}
结合式（2-75）、（2-79）和（2-80）我们得到

\begin{sourceblock}
\includegraphics[page=49,pagebox=mediabox,trim={45.00bp 160.38bp 44.37bp 454.87bp},clip,width=0.92\textwidth,max height=0.38\textheight]{Ishii.pdf}
\end{sourceblock}
因此状态热方程

\begin{sourceblock}
\includegraphics[page=49,pagebox=mediabox,trim={45.00bp 101.92bp 44.37bp 539.35bp},clip,width=0.92\textwidth,max height=0.38\textheight]{Ishii.pdf}
\end{sourceblock}
提供所有必要的信息来相互关联热力学性质。将式（2-81）代入式（2-73），得到表面张力的能量跳跃条件为

\begin{sourceblock}
\includegraphics[page=50,pagebox=mediabox,trim={45.00bp 498.78bp 44.37bp 55.42bp},clip,width=0.92\textwidth,max height=0.38\textheight]{Ishii.pdf}
\end{sourceblock}
界面熵不等式 根据上述讨论，我们假设存在表面温度 i T，这使我们能够在界面处写出熵不等式。因此，在没有表面热通量和源项的情况下，我们有

\begin{sourceblock}
\includegraphics[page=50,pagebox=mediabox,trim={45.00bp 359.74bp 44.37bp 261.85bp},clip,width=0.92\textwidth,max height=0.38\textheight]{Ishii.pdf}
\end{sourceblock}
利用能量平衡方程式(2-73)和状态方程式(2-77)可以消去上述不等式中的熵a s，从而得到

\begin{sourceblock}
\includegraphics[page=50,pagebox=mediabox,trim={45.00bp 207.72bp 44.37bp 364.72bp},clip,width=0.92\textwidth,max height=0.38\textheight]{Ishii.pdf}
\end{sourceblock}
我们注意到，这个表达式与 Delhaye (1974) 获得的表达式具有相同的形式。 Standart (1968) 也得出了类似的结果，没有考虑表面性质和表面张力项，但包括化学反应的影响。

一般来说，界面跳跃条件，方程（2-69），（2-72）和（2-84），不构成唯一定义问题所必需的足够的匹配条件。因此，它们应该辅以限制两相之间运动学、动力学和热关系的各种边界条件。这些关系也可以被视为界面本构关系，满足熵不等式（2-86）施加的限制。它们可以从标准中获得
基于界面本构关系的分析对于两相系统的详细研究可能很重要。然而，它们通常太复杂而无法用作边界条件。此外，势的影响，即温度、化学势、切向速度等的不连续性，作为量传递的驱动力，或由此产生的对热、动量和质量传递的界面阻力在整个系统中相对不显著。

因此，需要一种更简单的理论来提供必要的边界条件。作为限制情况，可以考虑界面的熵产生aΔ变为零的情况。这意味着界面上的量传递不存在阻力。因此，两相之间的交换由每一侧的大量流体的条件决定，而不是由界面本身决定。此外，从经典热力学的角度来看，界面处的传递被认为是可逆的。对于单相流中的激波不连续性来说，情况并非如此。

通过将方程（2-86）的熵产生设置为零，我们得到

\begin{sourceblock}
\includegraphics[page=51,pagebox=mediabox,trim={45.00bp 108.53bp 44.37bp 470.83bp},clip,width=0.92\textwidth,max height=0.38\textheight]{Ishii.pdf}
\end{sourceblock}
此外，我们假设对于质量通量、切向应力和热通量的所有组合，方程（2-87）中的三项独立为零。

热边界条件 因此，从式（2-87）的最后一项，我们得到界面处的热平衡条件

\begin{sourceblock}
\includegraphics[page=52,pagebox=mediabox,trim={45.00bp 521.92bp 44.37bp 122.48bp},clip,width=0.92\textwidth,max height=0.38\textheight]{Ishii.pdf}
\end{sourceblock}
这与界面处状态方程存在的假设（2-74）和（2-83）一致。鉴于方程（2-82）和（2-84），该热边界条件设定了界面的能级。与上述方程相反，能量跳跃条件方程（2-73）指定了传递到界面的能量之间的关系。此外，热平衡条件式(2-88)消除了变量i T，它作为界面各相温度的匹配条件。我们在这里注意到，实际上，界面处的温度不连续性是存在的，并且可以通过动力学理论来估计（Hirschfelder 等，1954）。然而，对于大多数材料来说，它的值与绝对温度相比非常小，除了少数例外，例如液态金属（Brodkey，1971）。因此，在标准条件下对界面转移的影响可以忽略不计。

无滑移条件 根据界面表观速度 i v 的定义（式（2-50）），切向速度 ti v 是一个未知参数，而法向分量与界面位置直接相关。此外，它还出现在熵不等式（2-86）和方程（2-87）的耗散项中。因此，很自然地提供切向应力 tk τ 和切向相对速度 tk ti − v v 之间的本构关系，正如前面所讨论的那样。然而，在本分析中，我们假设界面熵产生完全为零。通过将式(2-87)的第二项独立地取为零，我们得到无滑移条件

\begin{sourceblock}
\includegraphics[page=52,pagebox=mediabox,trim={45.00bp 155.68bp 44.37bp 488.40bp},clip,width=0.92\textwidth,max height=0.38\textheight]{Ishii.pdf}
\end{sourceblock}
移动粘性流体与固体壁接触的无滑移条件已得到充分确立（Goldstein，1938；Serrin，1959）。它被称为经典粘附条件，已通过实验和动力学理论分析得到验证。利用式(2-89)给出的关系可以消除界面切向粒子速度，然后将其用作界面处的速度边界条件。

然而，这里应该注意的是，对于无粘性流体，由于动量跳跃条件的切向分量（式（2-72）），无滑移条件（2-89）不是必需的并且通常不能满足。这与我们的分析完全一致，因为对于无粘性流体，方程（2-87）中的粘性耗散项同样为零，并且不会出现在熵不等式中。因此，无法得到式(2-89)。此外，在无滑移条件下，切线方向和法线方向的动量跳跃条件式（2-72）变为

\begin{sourceblock}
\includegraphics[page=53,pagebox=mediabox,trim={45.00bp 452.20bp 44.37bp 177.02bp},clip,width=0.92\textwidth,max height=0.38\textheight]{Ishii.pdf}
\end{sourceblock}
and

\begin{sourceblock}
\includegraphics[page=53,pagebox=mediabox,trim={45.00bp 376.54bp 44.37bp 245.05bp},clip,width=0.92\textwidth,max height=0.38\textheight]{Ishii.pdf}
\end{sourceblock}
其中法向粘性应力和切向粘性应力由下式给出

\begin{sourceblock}
\includegraphics[page=53,pagebox=mediabox,trim={45.00bp 319.83bp 44.37bp 321.27bp},clip,width=0.92\textwidth,max height=0.38\textheight]{Ishii.pdf}
\end{sourceblock}
平均曲率 21 H 取自第 2 相到第 1 相，即 21 0 H > 如果界面在第 1 相中为凸面。

化学（相变）边界条件与前面的讨论类似，化学（或相变）边界条件可以通过将式（2-87）的第一项对于k m 的所有值独立地设置为零来获得。这意味着相变产生的熵为零，因此相变不被认为是非平衡力引起的转移，而是状态的平衡变换。

将热平衡条件式（2-88）代入式（2-87）第一项并令其为零，可得

\begin{sourceblock}
\includegraphics[page=53,pagebox=mediabox,trim={45.00bp 97.69bp 44.37bp 520.12bp},clip,width=0.92\textwidth,max height=0.38\textheight]{Ishii.pdf}
\end{sourceblock}
上式给出的相变条件表明化学势的差异补偿了机械效应

\begin{sourceblock}
\includegraphics[page=54,pagebox=mediabox,trim={45.00bp 495.88bp 44.37bp 145.37bp},clip,width=0.92\textwidth,max height=0.38\textheight]{Ishii.pdf}
\end{sourceblock}
它取代了化学势的条件。

\subsection*{2.2.3 简化边界条件}\phantomsection\addcontentsline{toc}{subsection}{2.2.3 简化边界条件}
在前面的章节中已经给出了界面跳跃条件和补充边界条件。重要的是要认识到热平衡条件方程（2-88）、动量跳跃条件的法向分量方程（2-91）和相变边界条件方程（2-93）对应于恒温学的标准热、机械和化学平衡条件（Gibbs，1948）。不同之处在于，本分析考虑了质量传递的动态效应以及机械和相变边界条件下的法向应力。动力学分析结果和恒温理论结果之间的这些有趣的特性可以总结在下表中。

从表中可以看出，除了热条件之外，这些界面关系对于许多实际应用来说仍然非常复杂。这主要是由于质量传递和法向应力产生的项。前者由于机械边界条件中的密度变化而产生推力，并且还作为化学（相）边界条件中的冲击动能变化。后者引入了流场与界面热力学性质的复杂耦合效应。然而，在标准条件下，相对于压力项，法向应力可以忽略不计，这大大简化了机械边界条件，方程（2-91）。同样的论点也适用于化学边界条件，因为 k k g ρ 项的数量级为 k p，​​因此式（2-93）中的法向应力项也可以忽略不计。类似地，在大多数实际问题中，传质项很小，可以忽略不计，尽管它们对于传质速率较大或汽膜沸腾的问题可能很重要。

由于在场方程的标准公式中，吉布斯自由能
k g 没有明确出现，需要将化学边界条件式（2-93）中的变量 k g 变换为具有以下性质的其他变量：

\begin{sourceblock}
\includegraphics[page=55,pagebox=mediabox,trim={49.41bp 449.59bp 50.60bp 78.43bp},clip,width=\textwidth,max height=0.72\textheight]{Ishii.pdf}
\end{sourceblock}
已经在场方程中使用。为此，我们在这里回顾一下，表示为温度和压力函数的吉布斯自由能是一个基本状态方程，Eq.(2-30)，因此我们有

\begin{sourceblock}
\includegraphics[page=55,pagebox=mediabox,trim={45.00bp 354.64bp 44.37bp 288.86bp},clip,width=0.92\textwidth,max height=0.38\textheight]{Ishii.pdf}
\end{sourceblock}
and

\begin{sourceblock}
\includegraphics[page=55,pagebox=mediabox,trim={45.00bp 281.26bp 44.37bp 346.06bp},clip,width=0.92\textwidth,max height=0.38\textheight]{Ishii.pdf}
\end{sourceblock}
恒温相平衡条件由下式给出

\begin{sourceblock}
\includegraphics[page=55,pagebox=mediabox,trim={45.00bp 224.71bp 44.37bp 418.27bp},clip,width=0.92\textwidth,max height=0.38\textheight]{Ishii.pdf}
\end{sourceblock}
因此从式（2-95）和（2-97）我们得到

\begin{sourceblock}
\includegraphics[page=55,pagebox=mediabox,trim={45.00bp 167.58bp 44.37bp 475.08bp},clip,width=0.92\textwidth,max height=0.38\textheight]{Ishii.pdf}
\end{sourceblock}
这减少到经典的饱和条件

\begin{sourceblock}
\includegraphics[page=55,pagebox=mediabox,trim={45.00bp 108.60bp 44.37bp 534.06bp},clip,width=0.92\textwidth,max height=0.38\textheight]{Ishii.pdf}
\end{sourceblock}
该关系表明，恒温平衡条件与各相的热力学势唯一相关。此外，

\begin{sourceblock}
\includegraphics[page=56,pagebox=mediabox,trim={45.00bp 509.14bp 44.37bp 96.25bp},clip,width=0.92\textwidth,max height=0.38\textheight]{Ishii.pdf}
\end{sourceblock}
其中右侧的所有值都是在方程（2-99）给出的饱和线上计算的。

如果我们假设各相界面压力与界面温度 i T 对应的饱和压力的偏差与压力水平相比足够小，则吉布斯自由能函数可以在静态饱和点附近展开。因此我们有

\begin{sourceblock}
\includegraphics[page=56,pagebox=mediabox,trim={45.00bp 350.94bp 44.37bp 265.60bp},clip,width=0.92\textwidth,max height=0.38\textheight]{Ishii.pdf}
\end{sourceblock}
其中 k p δ 定义为

\begin{sourceblock}
\includegraphics[page=56,pagebox=mediabox,trim={45.00bp 292.01bp 44.37bp 349.27bp},clip,width=0.92\textwidth,max height=0.38\textheight]{Ishii.pdf}
\end{sourceblock}
既然我们有

\begin{sourceblock}
\includegraphics[page=56,pagebox=mediabox,trim={45.00bp 232.64bp 44.37bp 409.63bp},clip,width=0.92\textwidth,max height=0.38\textheight]{Ishii.pdf}
\end{sourceblock}
方程（2-93）可以简化为

\begin{sourceblock}
\includegraphics[page=56,pagebox=mediabox,trim={45.00bp 155.63bp 44.37bp 466.06bp},clip,width=0.92\textwidth,max height=0.38\textheight]{Ishii.pdf}
\end{sourceblock}
而机械边界条件方程（2-91）根据 k p δ 的定义变为

\begin{sourceblock}
\includegraphics[page=56,pagebox=mediabox,trim={45.00bp 65.63bp 44.37bp 556.06bp},clip,width=0.92\textwidth,max height=0.38\textheight]{Ishii.pdf}
\end{sourceblock}
\begin{sourceblock}
\includegraphics[page=57,pagebox=mediabox,trim={45.00bp 465.97bp 44.37bp 95.26bp},clip,width=0.92\textwidth,max height=0.38\textheight]{Ishii.pdf}
\end{sourceblock}
该结果表明，两相都不处于式（2-99）给出的饱和状态。压力与 sat p 的偏差量取决于平均曲率、表面张力、传质速率和法向应力。如果我们仅考虑表面张力的影响并放弃通常可以忽略不计的其他项，则会出现有趣的结果。在这种情况下，我们可以近似

\begin{sourceblock}
\includegraphics[page=57,pagebox=mediabox,trim={45.00bp 322.27bp 44.37bp 296.68bp},clip,width=0.92\textwidth,max height=0.38\textheight]{Ishii.pdf}
\end{sourceblock}
由于液相的平均曲率 fg H 对于液滴为正，对于气泡为负，因此界面处的相压力对于液滴流均高于饱和压力，对于气泡流均低于饱和压力。

现在我们回顾一下在固定温度下的压力偏差方面存在超出饱和条件的液体加热或蒸汽冷却极限，即恒温学中状态方程的不稳定点。因此，我们写

\begin{sourceblock}
\includegraphics[page=57,pagebox=mediabox,trim={50.00bp 124.63bp 44.37bp 461.49bp},clip,width=\textwidth,max height=0.72\textheight]{Ishii.pdf}
\par\vspace{0.45\baselineskip}
\small 图2-3显示了对应于克劳修斯-克拉佩龙方程或方程(2-99)的饱和线以及亚稳态液相和气相的极限。这两个极限可以从范德华状态方程获得：
\end{sourceblock}
\begin{sourceblock}
\includegraphics[page=58,pagebox=mediabox,trim={50.00bp 446.34bp 44.37bp 55.00bp},clip,width=\textwidth,max height=0.72\textheight]{Ishii.pdf}
\par\vspace{0.45\baselineskip}
\small 图2-3。 p-T 图（石井，1975 年）
\end{sourceblock}
\begin{sourceblock}
\includegraphics[page=58,pagebox=mediabox,trim={45.00bp 336.96bp 44.37bp 265.13bp},clip,width=0.92\textwidth,max height=0.38\textheight]{Ishii.pdf}
\end{sourceblock}
其中 R 和 M 分别是气体常数和分子量。 a 和 b 是经验常数。热力学理论指出，内在的热力学稳定性需要

\begin{sourceblock}
\includegraphics[page=58,pagebox=mediabox,trim={45.00bp 223.37bp 44.37bp 393.10bp},clip,width=0.92\textwidth,max height=0.38\textheight]{Ishii.pdf}
\end{sourceblock}
因此，利用范德瓦尔斯方程， ( ) 1 0 p ρ ∂ ∂ =
可以找到。这些轨迹实际上代表两个极限，即过热液体极限和过冷蒸汽极限。这两个位点如图2-3中的虚线所示。

有趣的是，方程（2-107）在方程（2-108）的限制条件下给出了最小的液滴和气泡尺寸。换句话说，这些大小是统计意义上干扰的最低自然水平。超出这些限制，液相或气相不能在不改变相的情况下保持不变，因为统计波动会产生一个可以生长为气泡或液滴的核心。

在温度i T 下，式(2-107)和式(2-108)给出的关系如图2-4所示。广泛使用的界面条件是蒸气

\begin{sourceblock}
\includegraphics[page=59,pagebox=mediabox,trim={50.00bp 404.07bp 44.37bp 55.00bp},clip,width=\textwidth,max height=0.72\textheight]{Ishii.pdf}
\par\vspace{0.45\baselineskip}
\small 图2-4。 δ p g - δ p f 关系（Ishii，1975）
\end{sourceblock}
界面压力等于温度 i T 下的饱和压力 sat p，可进一步近似式（2-107）推导出。由于在较小的减压下相间密度比很大，即
1 c p p << 其中c p 为临界压力，式(2-107)可近似为

\begin{sourceblock}
\includegraphics[page=59,pagebox=mediabox,trim={45.00bp 242.98bp 44.37bp 376.09bp},clip,width=0.92\textwidth,max height=0.38\textheight]{Ishii.pdf}
\end{sourceblock}
\subsection*{2.2.4 外部边界条件和接触角}\phantomsection\addcontentsline{toc}{subsection}{2.2.4 外部边界条件和接触角}
外部边界条件是上一节中讨论的跳跃和补充界面边界条件的特例。对于标准的单相流问题，这些条件变得特别简单，因为传质速率 k m 、表面张力的影响和固壁界面的速度都设置为零。类似的简化也可以应用于两相流系统，但是，这里应该考虑两个特殊特征。它们是： 1. 壁微观结构对气泡成核的影响； 2. 相界面与外部边界的交点。

\begin{sourceblock}
\includegraphics[page=60,pagebox=mediabox,trim={50.00bp 489.90bp 44.37bp 55.00bp},clip,width=\textwidth,max height=0.72\textheight]{Ishii.pdf}
\par\vspace{0.45\baselineskip}
\small 图2-5。接触角（石井，1975）
\end{sourceblock}
第一个效应的特点是必须考虑表面成核位点的存在，这些位点具有偏离标准理想化壁边界的不规则几何形状。这些微观结构和这些位点的气体含量往往决定气泡的成核条件和热力学非平衡的程度。第二种情况是两个不同界面相遇产生的奇点，见图2-5。当气泡或液滴与外部边界接触时，汽液界面附着在壁上并在相交处形成奇异曲线。当形成这样的接触线时，通过液体测量的接触角θ表征沿曲线的状况。也可以针对该单线进行类似于界面分析的分析。在这种情况下，由于散装流体的传输面积是界面的厚度δ，因此可以忽略传质和流体通量的影响。因此，只有表面通量以及可能与曲线相关的属性（即接触线的能量）才是重要的。通过仅考虑表面通量，我们可以从法向平面中的力平衡到奇异曲线

\begin{sourceblock}
\includegraphics[page=60,pagebox=mediabox,trim={45.00bp 183.76bp 44.37bp 439.30bp},clip,width=0.92\textwidth,max height=0.38\textheight]{Ishii.pdf}
\end{sourceblock}
其中fg σ 、gs σ 和fs σ 分别表示气-液、气-固和液-固之间的表面张力。

我们注意到，如果我们忽略奇异曲线的正切张力以及曲线的热能，则方程（2-112）与跳跃条件一致。如果忽略这些影响，方程（2-112）是与跳跃条件并行获得的唯一条件。因此，正如已经提到的，接触角 θ 表征了该现象，如果 gs σ 和 fs σ 不可用，则应提供适当的本构关系。静态接触角 θ 已被很好地测量并制成表格
此外，移动界面的动态接触角可能与静态值显著不同。然而，由于动态条件下θ缺乏完善的本构关系，因此在实际问题中经常使用静态值。我们在这里只注意到，人们普遍认为静态和动态接触角之间的明显差异是表面张力 fg σ 和奇异曲线法向滑移速度的函数（Schwartz 和 Tejada，1972；Phillips 和 Raddiford，1972）。

在总结本节时，我们列出了实体壁处的标准外部边界条件：

外部边界的位置

\begin{sourceblock}
\includegraphics[page=61,pagebox=mediabox,trim={45.00bp 404.92bp 44.37bp 239.16bp},clip,width=0.92\textwidth,max height=0.38\textheight]{Ishii.pdf}
\end{sourceblock}
无传质条件

\begin{sourceblock}
\includegraphics[page=61,pagebox=mediabox,trim={45.00bp 349.83bp 44.37bp 294.13bp},clip,width=0.92\textwidth,max height=0.38\textheight]{Ishii.pdf}
\end{sourceblock}
粘性流体的无滑移条件

\begin{sourceblock}
\includegraphics[page=61,pagebox=mediabox,trim={45.00bp 294.88bp 44.37bp 349.20bp},clip,width=0.92\textwidth,max height=0.38\textheight]{Ishii.pdf}
\end{sourceblock}
动量跳跃条件下的力平衡

\begin{sourceblock}
\includegraphics[page=61,pagebox=mediabox,trim={45.00bp 239.92bp 44.37bp 401.32bp},clip,width=0.92\textwidth,max height=0.38\textheight]{Ishii.pdf}
\end{sourceblock}
能量跳跃条件下的能量平衡

\begin{sourceblock}
\includegraphics[page=61,pagebox=mediabox,trim={45.00bp 184.83bp 44.37bp 456.27bp},clip,width=0.92\textwidth,max height=0.38\textheight]{Ishii.pdf}
\end{sourceblock}
热平衡条件

\begin{sourceblock}
\includegraphics[page=61,pagebox=mediabox,trim={45.00bp 129.88bp 44.37bp 514.51bp},clip,width=0.92\textwidth,max height=0.38\textheight]{Ishii.pdf}
\end{sourceblock}
上述这些条件可以应用于流体与壁接触的情况。然而，它不能应用于界面与实体边界的相交处。在这样的奇异曲线上，应给出接触角 θ 的本构方程。最后，我们在下图中总结了两相流系统的局部即时公式。

\begin{sourceblock}
\includegraphics[page=62,pagebox=mediabox,trim={68.49bp 389.83bp 69.62bp 54.61bp},clip,width=\textwidth,max height=0.72\textheight]{Ishii.pdf}
\end{sourceblock}
\section*{2.3 局部即时公式在两相流问题中的应用}\phantomsection\addcontentsline{toc}{section}{2.3 局部即时公式在两相流问题中的应用}
\subsection*{2.3.1 在非常慢的流中作用于球形颗粒的曳力}\phantomsection\addcontentsline{toc}{subsection}{2.3.1 在非常慢的流中作用于球形颗粒的曳力}
作为将局部瞬时公式应用于两相流问题的示例，让我们研究半径为 d r 的固体球体在速度为 0 U（蠕动流）的非常慢的流中的阻力本构方程（Stokes，1851；Schlichting，1979）。为了分析地分析这个问题，我们假设（1）具有恒定粘度的牛顿粘性流体，（2）不可压缩流动（流体密度恒定），以及（3）非常小的雷诺数（ ( ) 0 2 1 d c c Re r U ρ μ ≡ << ），其中粘性效应主导流动，并且在动量方程中可以忽略惯性项。那么，连续性方程Eq.(2-8)和动量方程Eq.(2-10)可以线性化为

\begin{sourceblock}
\includegraphics[page=62,pagebox=mediabox,trim={45.00bp 136.53bp 44.37bp 504.57bp},clip,width=0.92\textwidth,max height=0.38\textheight]{Ishii.pdf}
\end{sourceblock}
\begin{sourceblock}
\includegraphics[page=62,pagebox=mediabox,trim={45.00bp 92.56bp 44.37bp 548.68bp},clip,width=0.92\textwidth,max height=0.38\textheight]{Ishii.pdf}
\end{sourceblock}
通过考虑排除流体静力效应的压力场，重力项被丢弃。速度分量和压力

\begin{sourceblock}
\includegraphics[page=63,pagebox=mediabox,trim={45.00bp 529.27bp 44.37bp 93.70bp},clip,width=0.92\textwidth,max height=0.38\textheight]{Ishii.pdf}
\end{sourceblock}
\begin{sourceblock}
\includegraphics[page=63,pagebox=mediabox,trim={45.00bp 467.29bp 44.37bp 155.68bp},clip,width=0.92\textwidth,max height=0.38\textheight]{Ishii.pdf}
\end{sourceblock}
\begin{sourceblock}
\includegraphics[page=63,pagebox=mediabox,trim={45.00bp 409.29bp 44.37bp 216.70bp},clip,width=0.92\textwidth,max height=0.38\textheight]{Ishii.pdf}
\end{sourceblock}
\begin{sourceblock}
\includegraphics[page=63,pagebox=mediabox,trim={45.00bp 321.22bp 44.37bp 274.68bp},clip,width=0.92\textwidth,max height=0.38\textheight]{Ishii.pdf}
\end{sourceblock}
因此，作用在固体球上的总阻力 D F 通过将表面周围的压力和剪切应力积分得出：

\begin{sourceblock}
\includegraphics[page=63,pagebox=mediabox,trim={45.00bp 222.10bp 44.37bp 388.87bp},clip,width=0.92\textwidth,max height=0.38\textheight]{Ishii.pdf}
\end{sourceblock}
其中 A 是表面积。这表明即使在这种粘度主导的流动中，阻力也由压力和剪切力组成。然后，我们定义阻力系数 DC：

\begin{sourceblock}
\includegraphics[page=63,pagebox=mediabox,trim={45.00bp 112.63bp 44.37bp 478.85bp},clip,width=0.92\textwidth,max height=0.38\textheight]{Ishii.pdf}
\end{sourceblock}
其中 p A 是粒子的投影面积。因此，我们有

\begin{sourceblock}
\includegraphics[page=64,pagebox=mediabox,trim={45.00bp 571.53bp 44.37bp 57.53bp},clip,width=0.92\textwidth,max height=0.38\textheight]{Ishii.pdf}
\end{sourceblock}
Rybczynski (1911) 和 Hadamard (1911) 将这一分析扩展到无限纳维-斯托克斯流体中球形流体粒子的蠕动运动（Brodkey，1967；Soo，1967）。因此，作用在流体粒子上的总力由下式给出

\begin{sourceblock}
\includegraphics[page=64,pagebox=mediabox,trim={45.00bp 452.70bp 44.37bp 163.72bp},clip,width=0.92\textwidth,max height=0.38\textheight]{Ishii.pdf}
\end{sourceblock}
\begin{sourceblock}
\includegraphics[page=64,pagebox=mediabox,trim={45.00bp 370.75bp 44.37bp 221.29bp},clip,width=0.92\textwidth,max height=0.38\textheight]{Ishii.pdf}
\end{sourceblock}
粒子雷诺数为

\begin{sourceblock}
\includegraphics[page=64,pagebox=mediabox,trim={45.00bp 296.80bp 44.37bp 326.56bp},clip,width=0.92\textwidth,max height=0.38\textheight]{Ishii.pdf}
\end{sourceblock}
显然，c v ∞ 和 d v 分别是原状流速和粒子速度。因此，我们有

\begin{sourceblock}
\includegraphics[page=64,pagebox=mediabox,trim={45.00bp 202.62bp 44.37bp 413.78bp},clip,width=0.92\textwidth,max height=0.38\textheight]{Ishii.pdf}
\end{sourceblock}
Rybczynski 和 Hadamard 给出的阻力定律在雷诺数约为 1 时效果良好。

\subsection*{2.3.2 开尔文-亥姆霍兹不稳定性}\phantomsection\addcontentsline{toc}{subsection}{2.3.2 开尔文-亥姆霍兹不稳定性}
作为将局部瞬时公式应用于两相流问题的另一个例子，让我们研究 Kelvin-Helmholtz 不稳定性（Helmholtz，1868；Kelvin；1871；Lamb，1945）。开尔文-亥姆霍兹不稳定性出现在水平管道中以平均速度 1 v 和 2 v 流动的不同密度 1 ρ 和 2 ρ 的两个流体层的界面处。为了解析地分析这个问题，我们假设：（1）无粘流

\begin{sourceblock}
\includegraphics[page=65,pagebox=mediabox,trim={45.00bp 479.47bp 44.37bp 120.72bp},clip,width=0.92\textwidth,max height=0.38\textheight]{Ishii.pdf}
\end{sourceblock}
因此，连续性方程 Eq.(2-8) 以速度势的形式给出：

\begin{sourceblock}
\includegraphics[page=65,pagebox=mediabox,trim={45.00bp 393.49bp 44.37bp 231.68bp},clip,width=0.92\textwidth,max height=0.38\textheight]{Ishii.pdf}
\end{sourceblock}
动量方程 Eq.(2-10) 由下式给出

\begin{sourceblock}
\includegraphics[page=65,pagebox=mediabox,trim={45.00bp 320.50bp 44.37bp 303.74bp},clip,width=0.92\textwidth,max height=0.38\textheight]{Ishii.pdf}
\end{sourceblock}
其中 ( ) F t 分别是 t 的函数。两相之间界面的形状用正弦波近似为

\begin{sourceblock}
\includegraphics[page=65,pagebox=mediabox,trim={45.00bp 250.09bp 44.37bp 391.25bp},clip,width=0.92\textwidth,max height=0.38\textheight]{Ishii.pdf}
\end{sourceblock}
其中 0 η 、k 和 C 分别是振幅、波数和波速。那么上层流体的速度势
( ) 1 k = 和下流体 ( ) 2 k = 是在上、下管道表面无流体渗透的边界条件和小摄动假设下推导的。

\begin{sourceblock}
\includegraphics[page=65,pagebox=mediabox,trim={45.00bp 120.47bp 44.37bp 502.87bp},clip,width=0.92\textwidth,max height=0.38\textheight]{Ishii.pdf}
\end{sourceblock}
\begin{sourceblock}
\includegraphics[page=65,pagebox=mediabox,trim={45.00bp 56.51bp 44.37bp 568.19bp},clip,width=0.92\textwidth,max height=0.38\textheight]{Ishii.pdf}
\end{sourceblock}
\begin{sourceblock}
\includegraphics[page=66,pagebox=mediabox,trim={45.00bp 511.00bp 44.37bp 91.67bp},clip,width=0.92\textwidth,max height=0.38\textheight]{Ishii.pdf}
\end{sourceblock}
\begin{sourceblock}
\includegraphics[page=66,pagebox=mediabox,trim={45.00bp 445.00bp 44.37bp 175.63bp},clip,width=0.92\textwidth,max height=0.38\textheight]{Ishii.pdf}
\end{sourceblock}
其中 i p 是光滑界面处的压力。两个流体层之间的界面压力差是由表面张力引起的，可以近似为

\begin{sourceblock}
\includegraphics[page=66,pagebox=mediabox,trim={45.00bp 346.59bp 44.37bp 278.69bp},clip,width=0.92\textwidth,max height=0.38\textheight]{Ishii.pdf}
\end{sourceblock}
则由式(2-135)和式(2-138)-至-(2-140)可得波速为

\begin{sourceblock}
\includegraphics[page=66,pagebox=mediabox,trim={45.00bp 255.53bp 44.37bp 365.56bp},clip,width=0.92\textwidth,max height=0.38\textheight]{Ishii.pdf}
\end{sourceblock}
其中 ( ) k k k kh ρ ρ ′ ≡ coth。在深水假设 ( ) ( ) 1 2 2 2 0 25 h k h k π π > ,。 , k ρ ′ 可以近似为k ρ。在这个
情况下，式(2-141)可简化为

\begin{sourceblock}
\includegraphics[page=66,pagebox=mediabox,trim={45.00bp 135.70bp 44.37bp 485.88bp},clip,width=0.92\textwidth,max height=0.38\textheight]{Ishii.pdf}
\end{sourceblock}
where

\begin{sourceblock}
\includegraphics[page=66,pagebox=mediabox,trim={45.00bp 61.81bp 44.37bp 562.60bp},clip,width=0.92\textwidth,max height=0.38\textheight]{Ishii.pdf}
\end{sourceblock}
\begin{sourceblock}
\includegraphics[page=67,pagebox=mediabox,trim={45.00bp 512.51bp 44.37bp 109.01bp},clip,width=0.92\textwidth,max height=0.38\textheight]{Ishii.pdf}
\end{sourceblock}
在这个稳定性标准中有几个要点需要认识到。首先，忽略流体的粘性效应；因此，雷诺数在这种类型的界面不稳定性中不起任何作用。系统的稳定性由三种效应决定，即重力、表面张力和相对运动。由于伯努利效应的惯性力，相对运动项总是不稳定的。表面张力总是稳定的，因为平坦的界面具有最小的表面积，并且表面张力的作用是抵抗平衡构型的任何变形。仅当上层流体轻于下层流体 ( 2 1 ρ ρ > ) 时，重力项才稳定。在没有流动的情况下（或稳定性准则的左侧），传播速度 C ∞ 是波数 k 的函数。因此，随着波长 2 k λ π = 从零变化到无穷大，波速先减小到最小值，然后增大。这个

\begin{sourceblock}
\includegraphics[page=67,pagebox=mediabox,trim={150.69bp 259.09bp 52.28bp 344.47bp},clip,width=\textwidth,max height=0.72\textheight]{Ishii.pdf}
\end{sourceblock}
发生在 ( ) 2 2 1 c k g ρ ρ σ = −。这对应于 2 c c k λ π = 的临界波长。这称为泰勒波长，是两相流中最重要的内部长度尺度之一。如果相对速度足够小以满足

\begin{sourceblock}
\includegraphics[page=67,pagebox=mediabox,trim={45.00bp 177.67bp 44.37bp 445.77bp},clip,width=0.92\textwidth,max height=0.38\textheight]{Ishii.pdf}
\end{sourceblock}
对于大于此限制的相对速度，系统仅在一定的波长范围内是有条件稳定的。当波长较大时，式(2-143)中2 C ∞ 的值主要由重力项决定。相反，如果 λ 足够小，则毛细管力控制波动。

此外，可以基于一维两相流方程开发类似的稳定性准则（Wallis，1969；Kocamustafaogullari，1971）。值得注意的是（Miles，1957）开尔文-亥姆霍兹不稳定性理论往往会高估临界相对速度

\begin{sourceblock}
\includegraphics[page=68,pagebox=mediabox,trim={45.00bp 515.97bp 44.37bp 126.96bp},clip,width=0.92\textwidth,max height=0.38\textheight]{Ishii.pdf}
\end{sourceblock}
对于具有相对较大波长的重力主导流

c λ λ >>，表面张力的影响可以忽略不计。考虑有限通道流，式（2-141）可以给出不稳定准则：

\begin{sourceblock}
\includegraphics[page=68,pagebox=mediabox,trim={45.00bp 422.51bp 44.37bp 199.01bp},clip,width=0.92\textwidth,max height=0.38\textheight]{Ishii.pdf}
\end{sourceblock}
通过采用泰勒展开式并仅保留双曲函数的一阶项，可以获得以下简化但有用的准则。

\begin{sourceblock}
\includegraphics[page=68,pagebox=mediabox,trim={45.00bp 318.66bp 44.37bp 303.05bp},clip,width=0.92\textwidth,max height=0.38\textheight]{Ishii.pdf}
\end{sourceblock}
当将此标准与通道中段塞形成的实验数据进行比较时，临界相对速度被高估了接近两倍。这种差异可以通过引入有限振幅或波前传播方法的理论分析来解释（Mishima 和 Ishii，1980；Wu 和 Ishii，1996）。

\subsection*{2.3.3 瑞利-泰勒不稳定性}\phantomsection\addcontentsline{toc}{subsection}{2.3.3 瑞利-泰勒不稳定性}
瑞利-泰勒不稳定性是两种不同密度的流体之间的界面不稳定性，这些流体在重力场中分层或垂直于界面加速。通常观察到，如果上层流体密度 1 ρ 大于下层流体密度 2 ρ，则两个静止分层流体层之间的边界不稳定。由于瑞利-泰勒不稳定性会导致单个公共界面的破坏，因此它对于气泡或液滴的形成很重要。特别是，相关稳定性分析预测的临界波长是两相流最重要的长度尺度之一。

瑞利-泰勒不稳定性可以被视为零流量和 1 2 ρ ρ > 的开尔文-亥姆霍兹不稳定性的特例。因此，

\begin{sourceblock}
\includegraphics[page=69,pagebox=mediabox,trim={45.00bp 531.67bp 44.37bp 69.60bp},clip,width=0.92\textwidth,max height=0.38\textheight]{Ishii.pdf}
\end{sourceblock}
如果传播速度的根具有非零虚部，则系统不稳定。因此，式（2-148）表明，对于 1 2 ρ ρ >，重力是不稳定的，而表面张力是稳定的。存在一个临界波长 λ，低于该波长 2 C 始终为正值。这个
由 ( ) 2 1 2 c g λ π σ ρ ρ = − 给出。如果扰动的波长为
大于临界波长（c λ λ > ），则 2 C 变为负值，界面不稳定。对于横向不受限制的流体，扰动的波长可以任意大；因此这样的系统总是不稳定的。然而，如果流体被横向限制，则最大波长被限制为系统尺寸的两倍。这意味着如果横向特征尺寸小于临界波长c λ 的一半，则系统是稳定的。对于空气-水系统，该特征尺寸为 0.86 厘米。通过在稳定性分析中使用极坐标，可以从垂直圆柱体中包含的流体获得类似的尺寸。

对于不稳定系统，任何波长大于
c λ 可以随时间增长。然而，主导波是那些具有最大生长因子的波。由于波幅随着

\begin{sourceblock}
\includegraphics[page=69,pagebox=mediabox,trim={45.00bp 210.67bp 44.37bp 389.71bp},clip,width=0.92\textwidth,max height=0.38\textheight]{Ishii.pdf}
\end{sourceblock}
这些不稳定的波可以被观察为雨天从电线上滴下的水滴，或从水平向下的表面落下的冷凝水滴。在薄膜沸腾中也可以观察到由于瑞利-泰勒不稳定性而产生的非常规则的波形和气泡的产生。请注意，这种不稳定性不仅限于引力场。任何界面以及垂直于界面加速的流体都可能表现出相同的不稳定性。例如，这可能发生在核爆炸和聚变弹丸的惯性约束中。在这种情况下，分析中应使用加速度代替重力场 g。

\chapter{各种平均方法}
\section*{3.1 平均的目的}\phantomsection\addcontentsline{toc}{section}{3.1 平均的目的}
工程系统的设计和预测其性能的能力取决于实验数据和概念模型的可用性，这些数据和概念模型可用于以所需的准确度描述物理过程。从科学和实践的角度来看，这些概念模型和过程的各种特征和性质必须在合理的基础上明确表述并得到实验数据的支持。为此，需要专门设计的实验，这些实验必须与分析研究结合并支持分析研究。在连续介质力学中，气体或液体单相流的概念模型是根据描述质量、动量、能量、电荷等守恒定律的场方程来制定的。然后，这些场方程由适当的本构方程（例如状态、应力、化学反应等的本构方程）来补充，这些方程指定了给定组成材料（即指定的固体、液体或气体）的热力学、输运和化学性质。

因此，预计描述多相或多组分介质的稳态和动态特性的概念模型也应根据适当的场和本构方程来制定。然而，结构介质流的此类方程的推导比单相流的严格连续均质介质的方程的推导要复杂得多。为了理解推导结构平衡方程的困难，即具有界面不连续性的非均匀介质，我们回顾一下，在连续介质力学中，场论是建立在质量、动量和能量的积分平衡基础上的。因此，如果变量在
在多相或多组分流中，界面的存在给问题的数学和物理表述带来了很大的困难。从数学的角度来看，多相流可以被视为细分为单相区域的场，其移动边界将各组成相分开。每个子区域都存在差异平衡。然而，在不违反上述连续性条件的情况下，它不能应用于通常意义上的这些子区域的集合。从物理学的角度来看，推导适用于多相流系统的场和本构方程所遇到的困难源于界面的存在。它还源于这样一个事实：多相流的稳态和动态特性都取决于流动的界面结构。例如，分散两相流系统的稳态和动态特性取决于固体颗粒、气泡或液滴彼此相互作用以及与周围连续相相互作用的集体动力学；然而，在分离流的情况下，这些特性取决于界面的结构和波动力学。为了确定粒子的集体相互作用和界面的动力学，有必要首先描述流动的局部属性，然后通过适当的平均程序获得宏观描述。例如，对于分散流，有必要确定单个液滴（气泡）的成核、蒸发或凝结、运动和分裂的速率以及多个液滴（或气​​泡）的碰撞和聚结过程。

对于分离流，界面的结构和动力学极大地影响质量、热量和动量传递的速率以及系统的稳定性。例如，用于太空应用的冷凝器的性能和流动稳定性取决于蒸汽界面的动力学。类似地，液膜夹带液滴的速率以及膜冷却的有效性取决于气液界面的稳定性。

从这个讨论中可以得出结论，为了推导适合结构化多相流的场和本构方程，有必要描述流动的局部特征。从该流程中，应通过适当的平均程序获得宏观属性。很明显，大量重要技术的设计、性能以及通常的安全操作
第 2 章基于局部即时变量的公式表明，一般来说，它会导致界面位置未知的多边界问题。在这种情况下，获得解决方案所遇到的数学困难非常大，并且在许多实际问题中，它们超出了我们目前的计算能力。为了理解这些困难，我们回顾一下，即使在没有移动界面的单相湍流中，也不可能获得表达局部瞬时波动的精确解。可以说，本土速溶配方遇到的巨大困难源于：

1. 存在多个运动未知的可变形运动界面； 2. 由于湍流和界面运动而存在变量波动； 3. 界面处的性能显著不连续。第一个效应导致各相的场方程和界面条件之间存在复杂的耦合，而第二个效应不可避免地引入了源自纳维-斯托克斯方程和界面波的不稳定性的统计特征。第三个效应在空间和时间上的各种变量中引入了巨大的局部跳跃。由于这些困难几乎存在于所有两相流系统中，因此应用局部即时公式来获得解决方案受到严重限制。然而，对于具有简单界面几何形状的系统，如在单个或多个气泡问题或分离流的情况下，它已被广泛使用并获得了非常有用的信息。由于实际工程系统中观察到的大多数两相流具有极其复杂的界面几何形状和运动，因此不可能求解流体粒子的局部瞬时运动。工程问题很少需要流体运动和其他变量的这种微观细节，但流动的宏观方面更为重要。

通过适当的平均，可以获得流体运动和特性的平均值，有效消除局部瞬时波动。平均过程可以被视为低通滤波，从局部瞬时波动中排除不需要的高频信号。然而，值得注意的是，在基于平均的公式中应考虑这些影响宏观现象的波动的统计特性。

\section*{3.2 平均的分类}\phantomsection\addcontentsline{toc}{section}{3.2 平均的分类}
为了推导结构化两相介质的宏观场和本构方程，平均程序的重要性和必要性已在第 3 章的 1.1 节中讨论。在本节中，我们研究可应用于一般热流体动力学和特别是两相流的各种平均方法。根据用于制定热工水力问题的基本物理概念，平均程序可分为三大类：欧拉平均；拉格朗日平均；和玻尔兹曼统计平均。它们可以根据定义平均数学运算符的变量进一步分为子组。下面给出了分类总结和各种平均的定义。

i) 欧拉平均值 - 欧拉平均值

\begin{sourceblock}
\includegraphics[page=73,pagebox=mediabox,trim={45.00bp 366.05bp 44.37bp 279.54bp},clip,width=0.92\textwidth,max height=0.38\textheight]{Ishii.pdf}
\end{sourceblock}
\begin{sourceblock}
\includegraphics[page=73,pagebox=mediabox,trim={45.00bp 304.24bp 44.37bp 322.29bp},clip,width=0.92\textwidth,max height=0.38\textheight]{Ishii.pdf}
\end{sourceblock}
\begin{sourceblock}
\includegraphics[page=73,pagebox=mediabox,trim={45.00bp 248.20bp 44.37bp 378.33bp},clip,width=0.92\textwidth,max height=0.38\textheight]{Ishii.pdf}
\end{sourceblock}
\begin{sourceblock}
\includegraphics[page=73,pagebox=mediabox,trim={45.00bp 192.22bp 44.37bp 434.31bp},clip,width=0.92\textwidth,max height=0.38\textheight]{Ishii.pdf}
\end{sourceblock}
\begin{sourceblock}
\includegraphics[page=73,pagebox=mediabox,trim={45.00bp 136.24bp 44.37bp 490.29bp},clip,width=0.92\textwidth,max height=0.38\textheight]{Ishii.pdf}
\end{sourceblock}
\begin{sourceblock}
\includegraphics[page=73,pagebox=mediabox,trim={45.00bp 80.20bp 44.37bp 546.33bp},clip,width=0.92\textwidth,max height=0.38\textheight]{Ishii.pdf}
\end{sourceblock}
\begin{sourceblock}
\includegraphics[page=74,pagebox=mediabox,trim={45.00bp 568.47bp 44.37bp 58.60bp},clip,width=0.92\textwidth,max height=0.38\textheight]{Ishii.pdf}
\end{sourceblock}
混合平均值：以上操作的组合
ii) 拉格朗日平均值 - 拉格朗日平均值

\begin{sourceblock}
\includegraphics[page=74,pagebox=mediabox,trim={45.00bp 475.01bp 44.37bp 170.57bp},clip,width=0.92\textwidth,max height=0.38\textheight]{Ishii.pdf}
\end{sourceblock}
\begin{sourceblock}
\includegraphics[page=74,pagebox=mediabox,trim={45.00bp 413.20bp 44.37bp 213.32bp},clip,width=0.92\textwidth,max height=0.38\textheight]{Ishii.pdf}
\end{sourceblock}
\begin{sourceblock}
\includegraphics[page=74,pagebox=mediabox,trim={45.00bp 356.79bp 44.37bp 270.21bp},clip,width=0.92\textwidth,max height=0.38\textheight]{Ishii.pdf}
\end{sourceblock}
iii) 玻尔兹曼统计平均值

\begin{sourceblock}
\includegraphics[page=74,pagebox=mediabox,trim={45.00bp 294.14bp 44.37bp 343.81bp},clip,width=0.92\textwidth,max height=0.38\textheight]{Ishii.pdf}
\end{sourceblock}
\begin{sourceblock}
\includegraphics[page=74,pagebox=mediabox,trim={45.00bp 216.86bp 44.37bp 394.30bp},clip,width=0.92\textwidth,max height=0.38\textheight]{Ishii.pdf}
\end{sourceblock}
这里我们注意到 x 和 X 分别是空间坐标和材料坐标，而 xi 是粒子的相速度或动能。此外，我们指出，真实时间或统计平均是通过取极限 t Δ → ∞ 或 N → ∞ 来定义的，这仅在概念上是可能的。材质坐标可以被视为所有粒子的初始位置，因此如果 X 固定，则意味着跟随粒子的函数值。

连续介质力学中最重要和最广泛使用的平均组是欧拉平均，因为它与人类观察和大多数仪器密切相关。该方法的基本概念是物理现象的时空描述。在所谓的欧拉描述中，时间和空间坐标被视为
拉格朗日平均值与力学的拉格朗日描述直接相关。由于粒子坐标 X 取代了欧拉描述的空间变量 x，因此这种平均自然适合粒子动力学的研究。如果我们的兴趣集中在单个粒子的行为而不是一组粒子的集体力学上，那么拉格朗日平均值对于分析就很重要且有用。拉格朗日时间平均是通过跟踪某个粒子并在一段时间间隔内观察它来获得的。一个简单的例子是特定车辆（例如汽车、火车或飞机）的平均速度。此外，欧拉时间平均可以通过一段时间间隔内经过道路上某一点的所有汽车的平均速度来示例。

与上面解释的平均值相反，欧拉和拉格朗日统计平均值基于统计假设，因为它们涉及由 n F 表示的 N 个相似样本的集合，其中 n N = " 1, , . 当我们问时，出现了一个基本问题：“对于具有波动信号的系统，相似样本是什么？”为了可视化一组相似的样本，将时间平均视为消除不需要的波动的过滤过程是有用的。然后，相似样本可以被视为一组样本，其所有重要变量的时间平均在一定偏差范围内。在这种情况下，平均的时间间隔和偏差范围定义了不需要的波动，因此统计平均取决于它们。对于基于时间平均的稳态流，时域上的随机采样可以构成一组适当的样本，就像在实验测量中经常进行的那样。在这种情况下，时间平均和统计平均是等效的。还有许多其他因素需要考虑，但也可以将其作为一个抽象概念。当结合实验数据研究本构方程时就会出现困难。涉及无限数量的相似样本的真正统计平均仅在概念上是可能的，并且无法实现。因此，如果单独考虑，系综平均面临两个困难，即选择一组相似的样本和将实验数据连接到模型。

当大的集体力学时，具有粒子数密度概念的玻尔兹曼统计平均非常重要
2 1 2 ψ 然后在粒子速度场上进行平均，可以简化为类似于连续介质力学中质量、动量和能量的标准守恒方程的形式。

\section*{3.3 与两相流分析相关的各种平均}\phantomsection\addcontentsline{toc}{section}{3.3 与两相流分析相关的各种平均}
为了研究两相流系统，各种研究人员已经使用了许多上述平均方法。平均的应用可以分为两大类
1. 定义属性，然后关联实验数据。 2. 获得可用于预测宏观过程的可用场和本构方程。最基本的用途是定义平均属性和运动，包括各相或混合物的各种浓度、密度、速度和能量。然后，这些正确定义的平均值可用于各种实验目的和经验相关性的发展。平均和仪器的选择密切相关，因为通常测量量本身代表某种平均值。

欧拉时间平均和空间平均都经常使用，因为实验者倾向于将两相混合物视为准连续介质。此外，它们通常是流体流动系统中最容易测量的平均值。然而，当特定的流体颗粒可以区分和追踪时，如气泡或液滴流的情况，也可以使用拉格朗日平均值。很自然，这些平均值是针对可以被认为具有平均值方面的稳态特征的固定系统获得的。然后通过进一步应用不同数据之间的统计平均来开发各种相关性。这是实验物理学中最小化误差的标准方法。

在我们继续进行平均的第二次应用之前，我们简要讨论宏观场方程的两种根本不同的公式：即二流体模型和漂移通量（混合）模型。双流体模型是通过单独考虑每个相来制定的。因此，它由两组质量、动量和能量守恒方程表示。这六个场方程中的每一个都始终具有通过跳跃条件耦合两个相的相互作用项。混合模型是通过将混合物视为一个整体来制定的。因此，该模型用质量、动量和能量的三混合守恒方程表示，并附加一个考虑了浓度变化的扩散（连续性）方程。在适当的跳跃条件下，将每个相的两个相应的守恒方程相加即可得到混合守恒方程。然而，应该注意的是，应该根据正确定义的混合量来制定适当的混合模型。可以说，漂移-通量模型是混合模型的一个例子，混合模型包括扩散模型、滑流模型和均质流模型。然而，对于大多数实际应用，漂移通量模型是针对正常重力（Ishii，1977）和微重力条件（Hibiki 和 Ishii，2003b；Hibiki 等人，2004）高度开发的最佳混合模型。

现在我们继续讨论这些平均的第二个也是更重要的应用。即得到宏观两相流场方程和均值形式的本构方程。在这里，欧拉空间和时间平均已经被不同的作者广泛使用，尽管欧拉或玻尔兹曼统计平均也被应用。

Zuber (1964a)、Zuber 等人利用欧拉体积平均法，对建立高度分散流的三维模型做出了重要贡献。 (1964)、冯特 (1967)、德尔哈耶 (1968) 和斯莱特里 (1972)。这些分析基于同时包含两个阶段的体积元素。此外，还被认为
人们很早就认识到，管道横截面的欧拉面积平均对于工程应用非常有用，因为场方程简化为一维模型。通过面积平均，垂直于主流方向的变量变化信息基本丢失。因此，壁与流体之间的动量和能量传递应通过经验相关性或通过取代精确界面条件的简化模型来表达。我们注意到，即使在单相流问题中，面积平均法也已被广泛使用，因为它的简单性在许多实际工程应用中是非常理想的。例如，壁面摩擦系数或传热系数的使用与面积平均的概念密切相关。我们在这里还提到它在可压缩流体流动分析中的广泛应用。 Bird 等人对单相流面积平均以及与热力学中的开放系统方程相对应的宏观方程进行了很好的回顾。 （1960）、惠特克（1968）和斯莱特里（1972）。冯卡门边界层积分法也是面积平均法的巧妙应用。此外，在有关润滑膜、明渠流动和力学壳理论的文献中可以找到大量面积平均的例子。

然而，在两相流系统的应用中，许多作者使用现象学方法而不是数学上精确的面积平均，因此Martinelli和Nelson（1948）、Kutateladze（1952）、Brodkey（1967）、Levy（1960）和Wallis（1969）的结果彼此不一致并且没有一个是完整的（Kocamustafaogullari，1971）。获得一维模型的合理方法是对横截面积上的单相微分场方程进行积分。 Meyer (1960) 是使用这种方法获得混合方程的早期用户，但他对各种混合物性质的定义以及缺乏扩散（连续性）方程引起了反对（Zuber，1967）。

Zuber 等人对带有附加扩散（连续性）方程（即漂移通量模型）的一维混合场方程进行了严格推导。 （1964）和祖伯（1967）。结果显示与非均相化学反应单相系统的场方程有显著的相似性。后者已发展为基于同时占据同一点但具有两种不同速度的相互作用连续体的扩散热机械理论。许多作者对该理论做出了贡献，因此我们只记得 Fick (1855)、Stefan (1871)、von Karman (1950)、Prigogine 和 Mazur (1951)、Hirschfelder 等人的著作。 (1954)、Truesdell 和 Toupin (1960) 以及 Truesdell (1969)。从完全不同的情况获得类似的结果
与 Zuber 的分析相反，Delhaye (1968) 以及 Vernier 和 Delhaye (1968) 的分析针对双流体模型，该模型基于每个相的三个场方程以及耦合两个场的三个跳跃条件。 Vernier 和 Delhaye (1968) 采用了一种非常系统的方法，从三个不同的欧拉空间平均以及统计和时间平均导出场方程。这显然是第一篇显示各种平均方法之间重要相似点和差异的出版物。表面张力的影响对于界面稳定性和流动状态的分析非常重要，已被纳入 Kocamustafaogullari (1971) 的研究中。这项研究强调，面积平均模型特别适合研究分离流态和界面波不稳定性。然而，只要可以提供本构方程，它就可以用于任何类型的流态（Bouré 和 Réocreux，1972）。此外，在前一篇参考文献中，漂移通量模型和双流体模型之间的分析方法明确分开。尽管这种区别对于其他类型的混合物来说是众所周知的，例如在 Landau (1941) 的氦 II 的超流动性、Pai (1962) 的等离子体动力学以及 Truesdell (1969) 的扩散理论的研究中，但在两相流分析中它一直是模糊的。传统两相流公式的这一缺点首先由 Zuber 和 Dougherty (1967) 指出。在随后使用时间平均的 Ishii (1971) 和使用面积平均的 Kocamustafaogullari (1971) 的分析中，两个模型之间有了明显的区别。 Bouré 和 Réocreux (1972) 也结合两相声波传播和窒息现象问题讨论了这一点。

欧拉时间平均法广泛应用于分析单相湍流，也可用于分析两相流。在将时间平均方法应用于混合物时，许多作者将其与其他空间平均程序结合起来。俄罗斯研究人员做出了重要贡献（Teletov，1945；Frankl，1953；Teletov，1957；Diunin，1963），他们使用欧拉时间体积平均值并获得了三维场方程。

仅基于欧拉时间平均的分析显然是由 Vernier 和 Delhaye (1968) 发起的，然而，那里没有给出导致数学公式的详细研究。此外，Panton（1968）通过首先在时间间隔中积分然后在体积元素中积分来获得混合模型。他的分析在整合过程中比俄罗斯研究人员的工作更明确，但两个结果非常相似。在 Ishii (1971) 中，双流体模型公式包括
可以说，欧拉时间平均对于湍流两相流或分散两相流特别有用（Ishii，1975；Ishii，1977；Ishii 和 Mishima，1984）。在这些流动中，由于传输过程高度依赖于平均值附近变量的局部波动，因此最好从实验数据中获得时间平均模型的本构方程。标准单相湍流分析支持了这一点。

Vernier 和 Delhaye (1968) 使用欧拉统计平均进行了广泛的研究，他们得出了一个重要的结论。他们得出结论，在稳态流条件下，来自真实时间平均的场方程，即时间平均
t Δ →∞，与统计平均的结果相同。此外，将统计平均与空间平均相结合，然后补充各种本构假设以产生实用的二维模型。玻尔兹曼统计平均也被一些作者（Murray，1954；Buevich，1969；Buyevich，1972；Culick，1964；Kalinin，1970；Pai，1971）用于高度分散的两相流系统。一般来说，考虑粒子密度函数，然后写出该函数的玻尔兹曼输运方程。 Kalinin (1970) 假设粒子密度函数代表特定质量和速度的粒子的预期数量，而 Pai (1971) 则将半径、速度和温度视为函数的参数。然后，通过对除时间和空间变量之外的粒子密度函数的参数进行积分，从麦克斯韦-玻尔兹曼方程获得了每个相的麦克斯韦传递方程的简化版本。由于它涉及对分布以及粒子间和粒子-气体相互作用项的假设，因此结果不是一般性的，而是代表一种特殊的连续体。

然而，有趣的是，值得注意的是，代表了局部意义上的混合物力学的三种不同方法和观点：应用于两相混合物的欧拉时间或统计平均；基于两个连续体的扩散热机械理论；以及应用于气体混合物或高度分散流的玻尔兹曼统计平均。第一种理论认为混合物本质上是一组由界面界定的单相区域，而在第二种理论中，两种组分在同一点和时间共存。与上述建立在连续介质力学基础上的两个理论不同，最后一个理论是基于统计期望和概率的。然而，重要的是，如果正确解释上述模型的每个传递项，所得的场方程具有非常相似的形式。

Arnold 等人使用集合单元平均进行了初步研究。 (1989)，他们由于非扭曲气泡表面的压力变化而导出了湍流应力和界面压力，以实现理想化的无粘性气泡流。他们讨论了空间平均技术固有的缺陷，并建议使用系综平均来制定两相流的双流体模型。张和 Prosperetti (1994a) 通过系综平均法导出了控制无粘液体中等球形可压缩气泡混合物的平均方程。他们的结论是，由于不需要特殊的闭合关系，该方法是系统性和通用性的，并提出了该方法可能应用于各种热流体和固体力学情况的可能性。张和 Prosperetti (1994b) 将这种方法扩展到半径可变的球体的情况。张（1993）总结了热传导和对流、斯托克斯流和热毛细管过程的其他应用。这里，读者还应该参考Prosperetti（1999），其中可以找到一些关于用平均方程模拟分散多相流的考虑。 Kolev (2002) 提出了一种两相流公式，主要用于开发基于多领域方法的安全分析代码。

最后，我们简要讨论拉格朗日平均在两相流系统中的应用。这种方法对于颗粒流很有用，但是，由于扩散和相变，通常会遇到相当大的困难和不切实际。对于无相变的颗粒流，许多实际情况都可以详细得到平均颗粒运动的拉格朗日方程。因此，我们注意到单粒子动力学的拉格朗日描述经常用作高度分散流中颗粒相的动量方程（Carrier，1958；Zuber，1964a）。许多关于气泡上升和终端速度的分析都隐含地使用拉格朗日时间平均，特别是在连续相处于湍流状态的情况下。

\chapter{时间平均的基本关系}
欧拉时间平均在研究单相湍流中的重要性是众所周知的。由于分析标准流体流动系统中最有用的信息是时间平均而不是流体的局部即时响应，因此它在实验和分析目的中的使用在湍流研究中是必不可少的。例如，平均速度、温度和压力或传热系数和摩擦系数是标准问题中通常需要的重要平均值。此外，常用的实验方法和测量非常适合时间平均的应用。因此，利用时间平均场方程和用平均值表示的本构关系对单相湍流进行了深入的研究。尽管这些基于时间平均的模型没有给出湍流的基本起源、结构和传输机制的答案，但它们在工程系统中的应用被广泛认为是解决问题的有效手段。

在讨论欧拉时间平均应用于两相混合物的重要性时，我们首先回顾一下，在两相流中，变量的局部瞬时波动不仅由湍流引起，而且还由快速移动和变形的界面引起。由于这些复杂的流动和波动，无法获得局部瞬时公式的解，因此为了导出适当的场和本构方程，有必要对原始局部瞬时公式应用一些平均过程。鉴于上述关于时间平均在单相湍流分析中的重要性和有用性的讨论，我们将时间平均应用于两相流是自然且符合逻辑的。

预计平均场方程可以清楚地表现出系统从极其复杂的界面和湍流涨落中得到的宏观现象，因为它们仅进入公式
1. 平滑湍流波动，其意义与单相流相同； 2.将交替占据体积元素的两个相带入同时存在于同一点的连续体中，并且每个相具有适当定义的概率。此外，应该认识到，平均场方程中出现的本构关系应该通过时间平均来表达。这些本构关系可以通过两相输运现象的简单建模以及通常用平均值表示的各种实验数据来得出。

在接下来的章节中，我们使用时间平均发展了两相流热流体动力学的详细理论。首先，我们假设任何特定点的占用者在两个阶段之间随机交替，并且时间平均函数在新坐标中足够平滑。即，基于平均时间间隔的具有最小尺度t Δ 的时间坐标，在该时间坐标以下，时间微分算子没有物理意义。

\section*{4.1 时域和函数定义}\phantomsection\addcontentsline{toc}{section}{4.1 时域和函数定义}
首先，我们回顾两相或两种不混溶混合物的独特特征是相或组分之间存在一个或多个界面。此外，单相流可以根据层流、过渡流和湍流中的流动几何形状进行分类，而两相流或不混溶液体的混合物的流动可以根据界面的几何形状分为三类：分离流；过渡流或混合流；和分散流。这些类别的结构化流如表 1-1 所示。

在任何流态中，如果将相界面视为厚度为零且属性具有跳跃不连续性的奇异表面，则各种属性在相界面处会遭受不连续的变化。这可以通过取流体密度ρ来更形象地说明，如图4-1和4-2所示。由于在两相流系统中，各相的质量明显被界面分开，并且不在分子水平上混合，因此局部瞬时流体密度在 1 ρ 和 2 ρ 之间显示出逐步不连续性。图 4-1 显示了 ρ 在空间上的瞬时不连续性，而图 4-2 显示了在某个固定点 0 x 处的时间不连续性。

\begin{sourceblock}
\includegraphics[page=84,pagebox=mediabox,trim={50.00bp 376.08bp 44.37bp 55.00bp},clip,width=\textwidth,max height=0.72\textheight]{Ishii.pdf}
\par\vspace{0.45\baselineskip}
\small 图 4-1。 t = t 0 时空间中的流体密度（Ishii，1975）
\end{sourceblock}
\begin{sourceblock}
\includegraphics[page=84,pagebox=mediabox,trim={50.00bp 155.52bp 44.37bp 304.85bp},clip,width=\textwidth,max height=0.72\textheight]{Ishii.pdf}
\par\vspace{0.45\baselineskip}
\small 图4-2。 x = x 0 时的流体密度（Ishii，1975）
\end{sourceblock}
出于时间平均的目的，从时间坐标进行观察可以更准确地描述问题。很容易看出，在任何固定点都可能发生四种不同的 ρ 过程，我们可以将其分类如下：

我们的目的是对流体特性和场方程进行平均，以便将两相流视为连续体的混合物。首先，我们采用固定的平均时间间隔 t Δ，并假设它足够大以平滑属性的局部变化，但与整体流不稳定性的宏观时间常数相比很小。该假设与分析湍流单相流时所做的假设相同。选择任意特定参考点和时间 ( ) 0 0 t , x 后，我们有确定的时间 1 2 j t t t " " , , ,，指的是经过该点的接口
0 x 从时间 ( ) 0 2 t t Δ − 到 ( ) 0 2 t t Δ +。通过使用任意小的
第2章1.2节的界面厚度δ，与每个界面相关的时间间隔可以定义为

\begin{sourceblock}
\includegraphics[page=85,pagebox=mediabox,trim={45.00bp 246.58bp 44.37bp 377.74bp},clip,width=0.92\textwidth,max height=0.38\textheight]{Ishii.pdf}
\end{sourceblock}
它可能因接口而异，因此我们使用 j ε 作为第 j 个接口。由于我们将把界面视为一个壳，其位置由数学表面表示，因此我们可以将 j ε 作为 1 δ 和 2 δ 的相应时间间隔。那么假设界面是奇异面，或者界面厚度 0 δ →，对应于

\begin{sourceblock}
\includegraphics[page=85,pagebox=mediabox,trim={45.00bp 133.33bp 44.37bp 477.74bp},clip,width=0.92\textwidth,max height=0.38\textheight]{Ishii.pdf}
\end{sourceblock}
在后续的分析中，我们经常使用这种关系来推导宏观场方程。现在我们定义时间间隔的集合，其中界面的特征占主导地位，如下

\begin{sourceblock}
\includegraphics[page=86,pagebox=mediabox,trim={50.00bp 398.64bp 44.37bp 55.00bp},clip,width=\textwidth,max height=0.72\textheight]{Ishii.pdf}
\par\vspace{0.45\baselineskip}
\small 图 4-3。不同的时间间隔（石井，1975）
\end{sourceblock}
\begin{sourceblock}
\includegraphics[page=86,pagebox=mediabox,trim={45.00bp 330.74bp 44.37bp 300.24bp},clip,width=0.92\textwidth,max height=0.38\textheight]{Ishii.pdf}
\end{sourceblock}
时间间隔的剩余部分由 [ ] T t Δ 给出，它可以分为阶段 1 和阶段 2 的间隔。因此

\begin{sourceblock}
\includegraphics[page=86,pagebox=mediabox,trim={45.00bp 260.42bp 44.37bp 378.33bp},clip,width=0.92\textwidth,max height=0.38\textheight]{Ishii.pdf}
\end{sourceblock}
通过介绍

\begin{sourceblock}
\includegraphics[page=86,pagebox=mediabox,trim={45.00bp 186.68bp 44.37bp 436.15bp},clip,width=0.92\textwidth,max height=0.38\textheight]{Ishii.pdf}
\end{sourceblock}
我们有

\begin{sourceblock}
\includegraphics[page=86,pagebox=mediabox,trim={45.00bp 115.23bp 44.37bp 511.27bp},clip,width=0.92\textwidth,max height=0.38\textheight]{Ishii.pdf}
\end{sourceblock}
这些关系如图4-3所示。

由于在分析过程中需要区分三种状态，即相 1、相 2 或界面，我们分配状态密度函数 1 M 、 2 M 和 S M，定义为：

\begin{sourceblock}
\includegraphics[page=87,pagebox=mediabox,trim={45.00bp 522.01bp 44.37bp 59.17bp},clip,width=0.92\textwidth,max height=0.38\textheight]{Ishii.pdf}
\end{sourceblock}
与两相相关的一般函数 F 被认为在除厚度 δ 的界面区域之外的任何地方都是连续可导的。那么平均 0 x 点处的第 k 相 k F 的一般函数定义为

\begin{sourceblock}
\includegraphics[page=87,pagebox=mediabox,trim={45.00bp 392.77bp 44.37bp 215.80bp},clip,width=0.92\textwidth,max height=0.38\textheight]{Ishii.pdf}
\end{sourceblock}
该函数 k F 表示第 2 章给出的局部瞬时公式中每个相的变量。

\section*{4.2 局部时间分数-局部空洞分数}\phantomsection\addcontentsline{toc}{section}{4.2 局部时间分数-局部空洞分数}
通过取极限来定义每个阶段所占用的时间间隔

0 δ → as (见图4-3)

\begin{sourceblock}
\includegraphics[page=87,pagebox=mediabox,trim={45.00bp 195.30bp 44.37bp 382.36bp},clip,width=0.92\textwidth,max height=0.38\textheight]{Ishii.pdf}
\end{sourceblock}
因此，根据假设（4-2）我们有

\begin{sourceblock}
\includegraphics[page=87,pagebox=mediabox,trim={45.00bp 142.21bp 44.37bp 499.11bp},clip,width=0.92\textwidth,max height=0.38\textheight]{Ishii.pdf}
\end{sourceblock}
\begin{sourceblock}
\includegraphics[page=87,pagebox=mediabox,trim={45.00bp 50.14bp 44.37bp 555.48bp},clip,width=0.92\textwidth,max height=0.38\textheight]{Ishii.pdf}
\end{sourceblock}
\begin{sourceblock}
\includegraphics[page=88,pagebox=mediabox,trim={45.00bp 544.23bp 44.37bp 80.74bp},clip,width=0.92\textwidth,max height=0.38\textheight]{Ishii.pdf}
\end{sourceblock}
由式(4-10)和式(4-12)可得如下关系

\begin{sourceblock}
\includegraphics[page=88,pagebox=mediabox,trim={45.00bp 490.92bp 44.37bp 150.28bp},clip,width=0.92\textwidth,max height=0.38\textheight]{Ishii.pdf}
\end{sourceblock}
这是平均界面密度函数 S α 为零的结果。我们注意到，S α 与方程（4-7）和（4-11）并行定义为

\begin{sourceblock}
\includegraphics[page=88,pagebox=mediabox,trim={45.00bp 394.23bp 44.37bp 233.75bp},clip,width=0.92\textwidth,max height=0.38\textheight]{Ishii.pdf}
\end{sourceblock}
函数k α 仅在积分运算后出现，是研究时间平均场方程的基本参数。物理上，k α 表示找到第 k 相的概率，因此它表示该相的几何（静态）重要性。此后我们将 k α 称为第 k 相的局部时间分数或局部空隙率。

\section*{4.3 时间平均和加权平均值}\phantomsection\addcontentsline{toc}{section}{4.3 时间平均和加权平均值}
在本节中，我们定义与两相流场相关的函数的时间平均和加权平均值。

时间平均 通用函数 F 的欧拉时间平均定义为

\begin{sourceblock}
\includegraphics[page=88,pagebox=mediabox,trim={45.00bp 167.19bp 44.37bp 460.72bp},clip,width=0.92\textwidth,max height=0.38\textheight]{Ishii.pdf}
\end{sourceblock}
此后，符号F表示等式(4-15)右侧定义的数学运算。类似地，第 k 相通用函数 k F 的平均值由下式给出

\begin{sourceblock}
\includegraphics[page=88,pagebox=mediabox,trim={45.00bp 71.18bp 44.37bp 532.86bp},clip,width=0.92\textwidth,max height=0.38\textheight]{Ishii.pdf}
\end{sourceblock}
\begin{sourceblock}
\includegraphics[page=89,pagebox=mediabox,trim={45.00bp 558.94bp 44.37bp 65.73bp},clip,width=0.92\textwidth,max height=0.38\textheight]{Ishii.pdf}
\end{sourceblock}
鉴于式(4-8)、(4-15)和(4-16)我们有

\begin{sourceblock}
\includegraphics[page=89,pagebox=mediabox,trim={45.00bp 487.30bp 44.37bp 140.34bp},clip,width=0.92\textwidth,max height=0.38\textheight]{Ishii.pdf}
\end{sourceblock}
因此我们得到一个重要的关系

\begin{sourceblock}
\includegraphics[page=89,pagebox=mediabox,trim={45.00bp 431.05bp 44.37bp 213.44bp},clip,width=0.92\textwidth,max height=0.38\textheight]{Ishii.pdf}
\end{sourceblock}
1F和2F的功能与各相的瞬时、局部物理或流量变量直接相关；然而，1 F 和 2 F 是在总时间间隔 t Δ 上取平均值的。因此，它们可以被视为表面平均值。从这个角度来看，我们引入了各种加权平均值，保留了原始变量的一些重要特征。

我们从一般情况开始，然后再讨论特殊情况。因此，通过采用非零标量权重函数 w，我们将 F 的一般加权平均值定义为

\begin{sourceblock}
\includegraphics[page=89,pagebox=mediabox,trim={45.00bp 254.38bp 44.37bp 374.84bp},clip,width=0.92\textwidth,max height=0.38\textheight]{Ishii.pdf}
\end{sourceblock}
其中函数 w 也属于上一节中定义的通用函数 F 的群。由此可见，各阶段的权重函数也可以通过式（4-8）来定义。

因此，与方程（4-20）并行的一般相位加权平均值应为

\begin{sourceblock}
\includegraphics[page=89,pagebox=mediabox,trim={45.00bp 128.63bp 44.37bp 498.86bp},clip,width=0.92\textwidth,max height=0.38\textheight]{Ishii.pdf}
\end{sourceblock}
由此可见

\begin{sourceblock}
\includegraphics[page=90,pagebox=mediabox,trim={45.00bp 536.79bp 44.37bp 58.35bp},clip,width=0.92\textwidth,max height=0.38\textheight]{Ishii.pdf}
\end{sourceblock}
其将混合物和相平均值联系起来。由于这些公式过于笼统，我们在下面讨论一些重要的特殊情况。

相位平均 k F 与每个相位相关的最自然平均值可以通过将相位密度函数 k M 作为方程（4-21）中的加权函数来定义，因此我们有

\begin{sourceblock}
\includegraphics[page=90,pagebox=mediabox,trim={45.00bp 383.80bp 44.37bp 243.52bp},clip,width=0.92\textwidth,max height=0.38\textheight]{Ishii.pdf}
\end{sourceblock}
\begin{sourceblock}
\includegraphics[page=90,pagebox=mediabox,trim={45.00bp 284.44bp 44.37bp 314.52bp},clip,width=0.92\textwidth,max height=0.38\textheight]{Ishii.pdf}
\end{sourceblock}
质量加权平均值 l ψ 和 m
k ψ 一般情况下，体积、动量、能量和熵等被认为是广延变量（Callen，1960）。如果将函数 F 视为广延特性的每单位体积的量，则它们也可以用每单位质量的变量 ψ 表示为

\begin{sourceblock}
\includegraphics[page=90,pagebox=mediabox,trim={45.00bp 164.71bp 44.37bp 477.11bp},clip,width=0.92\textwidth,max height=0.38\textheight]{Ishii.pdf}
\end{sourceblock}
其中 ρ 是局部瞬时流体密度。因此，每个相位 k F 的属性由下式给出

\begin{sourceblock}
\includegraphics[page=90,pagebox=mediabox,trim={45.00bp 96.03bp 44.37bp 545.08bp},clip,width=0.92\textwidth,max height=0.38\textheight]{Ishii.pdf}
\end{sourceblock}
\begin{sourceblock}
\includegraphics[page=91,pagebox=mediabox,trim={45.00bp 567.52bp 44.37bp 57.49bp},clip,width=0.92\textwidth,max height=0.38\textheight]{Ishii.pdf}
\end{sourceblock}
and

\begin{sourceblock}
\includegraphics[page=91,pagebox=mediabox,trim={45.00bp 490.60bp 44.37bp 133.63bp},clip,width=0.92\textwidth,max height=0.38\textheight]{Ishii.pdf}
\end{sourceblock}
用密度ρ对ψ进行加权的最重要原因是ψ所表示的量是质量的加性集函数。

根据质量加权平均值的定义，我们有

\begin{sourceblock}
\includegraphics[page=91,pagebox=mediabox,trim={45.00bp 391.48bp 44.37bp 235.03bp},clip,width=0.92\textwidth,max height=0.38\textheight]{Ishii.pdf}
\end{sourceblock}
因此，我们得到质量加权混合物性质与两相性质之间最重要的关系：

\begin{sourceblock}
\includegraphics[page=91,pagebox=mediabox,trim={45.00bp 273.76bp 44.37bp 321.39bp},clip,width=0.92\textwidth,max height=0.38\textheight]{Ishii.pdf}
\end{sourceblock}
我们在这里注意到，上述结果类似于扩散热机械理论（Truesdell，1969）和气体动力学理论中麦克斯韦传递方程（Maxwell，1867）中使用的定义。特别是，由于密度是单位体积的属性，因此我们有

\begin{sourceblock}
\includegraphics[page=91,pagebox=mediabox,trim={45.00bp 161.44bp 44.37bp 465.57bp},clip,width=0.92\textwidth,max height=0.38\textheight]{Ishii.pdf}
\end{sourceblock}
关于平均值平滑性的基本假设 我们平均的目的是将两个相（交替占据界面上不连续的点）转换为两个同时连续体。因此，需要对导数的连续性做出假设。因此，让我们在这里引入关于平均值 F 和 k F 的平滑度的基本假设。

通过根据平均值考虑宏观过程，假设如果过程的时间常数足够大于 t Δ，则除了某些孤立的奇点之外，它们足够平滑以具有必要的更高阶导数。换句话说，时间间隔 t Δ 内平均值的任何变化都被认为是无穷小的。这可以通过将平均场中的时间微分算子视为时间增量 t t δ Δ → 而不是 0 t δ → 的有限差分算子来可视化。如果我们对平均值进行平均，我们会得到

\begin{sourceblock}
\includegraphics[page=92,pagebox=mediabox,trim={45.00bp 454.16bp 44.37bp 173.73bp},clip,width=0.92\textwidth,max height=0.38\textheight]{Ishii.pdf}
\end{sourceblock}
然而，由于F是连续的，我们从积分中值定理得到

\begin{sourceblock}
\includegraphics[page=92,pagebox=mediabox,trim={45.00bp 379.64bp 44.37bp 258.31bp},clip,width=0.92\textwidth,max height=0.38\textheight]{Ishii.pdf}
\end{sourceblock}
where

\begin{sourceblock}
\includegraphics[page=92,pagebox=mediabox,trim={45.00bp 306.42bp 44.37bp 319.76bp},clip,width=0.92\textwidth,max height=0.38\textheight]{Ishii.pdf}
\end{sourceblock}
因此，与基本假设类似，假设在宏观领域中我们有

\begin{sourceblock}
\includegraphics[page=92,pagebox=mediabox,trim={45.00bp 231.55bp 44.37bp 406.36bp},clip,width=0.92\textwidth,max height=0.38\textheight]{Ishii.pdf}
\end{sourceblock}
它指出平均不会改变平均值。那么我们就可以简单的得到

\begin{sourceblock}
\includegraphics[page=92,pagebox=mediabox,trim={45.00bp 125.25bp 44.37bp 481.83bp},clip,width=0.92\textwidth,max height=0.38\textheight]{Ishii.pdf}
\end{sourceblock}
对于常数 C，我们有

\begin{sourceblock}
\includegraphics[page=92,pagebox=mediabox,trim={45.00bp 50.14bp 44.37bp 578.88bp},clip,width=0.92\textwidth,max height=0.38\textheight]{Ishii.pdf}
\end{sourceblock}
波动分量 正如在湍流分析中一样，引入变量的波动分量以便在统计上考虑这些影响是最重要的。一般来说，它们被定义为局部即时变量与其加权平均值之间的差，因此我们有

\begin{sourceblock}
\includegraphics[page=93,pagebox=mediabox,trim={45.00bp 506.91bp 44.37bp 133.21bp},clip,width=0.92\textwidth,max height=0.38\textheight]{Ishii.pdf}
\end{sourceblock}
由于一旦指定了变量 k F，则将给出加权平均值的形式，因此可以唯一地定义波动分量。从式（4-35）和（4-37）我们立即得到

\begin{sourceblock}
\includegraphics[page=93,pagebox=mediabox,trim={45.00bp 422.02bp 44.37bp 217.02bp},clip,width=0.92\textwidth,max height=0.38\textheight]{Ishii.pdf}
\end{sourceblock}
此外，波动分量的平均值可以与其他参数相关：

\begin{sourceblock}
\includegraphics[page=93,pagebox=mediabox,trim={45.00bp 342.41bp 44.37bp 291.10bp},clip,width=0.92\textwidth,max height=0.38\textheight]{Ishii.pdf}
\end{sourceblock}
and

\begin{sourceblock}
\includegraphics[page=93,pagebox=mediabox,trim={45.00bp 278.99bp 44.37bp 361.21bp},clip,width=0.92\textwidth,max height=0.38\textheight]{Ishii.pdf}
\end{sourceblock}
这些关系将用于分析平均场方程中的两相湍流通量。

\section*{4.4 导数的时间平均}\phantomsection\addcontentsline{toc}{section}{4.4 导数的时间平均}
在本节中，获得导数的平均值与平均值的导数之间的关系。通过平均值的时间导数，我们的意思是

\begin{sourceblock}
\includegraphics[page=93,pagebox=mediabox,trim={45.00bp 103.69bp 44.37bp 515.72bp},clip,width=0.92\textwidth,max height=0.38\textheight]{Ishii.pdf}
\end{sourceblock}
由于积分域是不连续的，我们对其进行细分，并将莱布尼茨规则应用于方程（4-41）。因此我们有

\begin{sourceblock}
\includegraphics[page=94,pagebox=mediabox,trim={45.00bp 510.51bp 44.37bp 55.51bp},clip,width=0.92\textwidth,max height=0.38\textheight]{Ishii.pdf}
\end{sourceblock}
鉴于方程（4-2）我们定义

\begin{sourceblock}
\includegraphics[page=94,pagebox=mediabox,trim={45.00bp 450.04bp 44.37bp 186.06bp},clip,width=0.92\textwidth,max height=0.38\textheight]{Ishii.pdf}
\end{sourceblock}
并相应地

\begin{sourceblock}
\includegraphics[page=94,pagebox=mediabox,trim={45.00bp 383.01bp 44.37bp 252.07bp},clip,width=0.92\textwidth,max height=0.38\textheight]{Ishii.pdf}
\end{sourceblock}
因此，由式(4-42)可得

\begin{sourceblock}
\includegraphics[page=94,pagebox=mediabox,trim={45.00bp 300.91bp 44.37bp 318.50bp},clip,width=0.92\textwidth,max height=0.38\textheight]{Ishii.pdf}
\end{sourceblock}
\begin{sourceblock}
\includegraphics[page=94,pagebox=mediabox,trim={45.00bp 229.46bp 44.37bp 374.68bp},clip,width=0.92\textwidth,max height=0.38\textheight]{Ishii.pdf}
\end{sourceblock}
因此，通过应用莱布尼茨规则，我们得到

\begin{sourceblock}
\includegraphics[page=94,pagebox=mediabox,trim={45.00bp 90.37bp 44.37bp 466.72bp},clip,width=0.92\textwidth,max height=0.38\textheight]{Ishii.pdf}
\end{sourceblock}
上式右侧最后一项的物理意义尚不清楚，应更详细地研究。为此，我们引入由下式给出的表面方程

\begin{sourceblock}
\includegraphics[page=95,pagebox=mediabox,trim={45.00bp 535.43bp 44.37bp 110.31bp},clip,width=0.92\textwidth,max height=0.38\textheight]{Ishii.pdf}
\end{sourceblock}
\begin{sourceblock}
\includegraphics[page=95,pagebox=mediabox,trim={45.00bp 461.19bp 44.37bp 141.56bp},clip,width=0.92\textwidth,max height=0.38\textheight]{Ishii.pdf}
\end{sourceblock}
因此，根据式（4-43），我们有

\begin{sourceblock}
\includegraphics[page=95,pagebox=mediabox,trim={45.00bp 378.26bp 44.37bp 232.75bp},clip,width=0.92\textwidth,max height=0.38\textheight]{Ishii.pdf}
\end{sourceblock}
然而，法向矢量和位移速度（Truesdell 和 Toupin，1960）由下式给出

\begin{sourceblock}
\includegraphics[page=95,pagebox=mediabox,trim={45.00bp 296.53bp 44.37bp 328.72bp},clip,width=0.92\textwidth,max height=0.38\textheight]{Ishii.pdf}
\end{sourceblock}
and

\begin{sourceblock}
\includegraphics[page=95,pagebox=mediabox,trim={45.00bp 210.63bp 44.37bp 401.22bp},clip,width=0.92\textwidth,max height=0.38\textheight]{Ishii.pdf}
\end{sourceblock}
因此，通过消除式（4-50）中的表面函数，我们得到

\begin{sourceblock}
\includegraphics[page=95,pagebox=mediabox,trim={45.00bp 139.54bp 44.37bp 487.61bp},clip,width=0.92\textwidth,max height=0.38\textheight]{Ishii.pdf}
\end{sourceblock}
则式(4-47)变为

\begin{sourceblock}
\includegraphics[page=96,pagebox=mediabox,trim={45.00bp 543.22bp 44.37bp 58.76bp},clip,width=0.92\textwidth,max height=0.38\textheight]{Ishii.pdf}
\end{sourceblock}
界面的单位法向量定义为

\begin{sourceblock}
\includegraphics[page=96,pagebox=mediabox,trim={45.00bp 487.12bp 44.37bp 154.12bp},clip,width=0.92\textwidth,max height=0.38\textheight]{Ishii.pdf}
\end{sourceblock}
Then

\begin{sourceblock}
\includegraphics[page=96,pagebox=mediabox,trim={45.00bp 431.14bp 44.37bp 209.83bp},clip,width=0.92\textwidth,max height=0.38\textheight]{Ishii.pdf}
\end{sourceblock}
其中 + n 和 − n 对应于界面每侧流体的极限向外法向矢量。

使用简化符号，方程（4-45）和（4-54）变为

\begin{sourceblock}
\includegraphics[page=96,pagebox=mediabox,trim={45.00bp 325.78bp 44.37bp 297.59bp},clip,width=0.92\textwidth,max height=0.38\textheight]{Ishii.pdf}
\end{sourceblock}
\begin{sourceblock}
\includegraphics[page=96,pagebox=mediabox,trim={45.00bp 266.80bp 44.37bp 360.53bp},clip,width=0.92\textwidth,max height=0.38\textheight]{Ishii.pdf}
\end{sourceblock}
我们在这里注意到函数 F 可以是标量、向量或张量，并且
∇ 算子可以是散度算子或梯度算子，在 ± n 和 F ± 之间进行适当的张量运算。上述两种变换和各种平均值的定义是获得平均值形式的宏观场方程的基本工具。与函数 F 中没有不连续性的情况相反，上述变换显示了移动界面在将导数的平均值与平均值的导数相关联方面做出的重要贡献。

作为方程（4-57）和（4-58）的推论，我们有

\begin{sourceblock}
\includegraphics[page=96,pagebox=mediabox,trim={45.00bp 87.76bp 44.37bp 536.57bp},clip,width=0.92\textwidth,max height=0.38\textheight]{Ishii.pdf}
\end{sourceblock}
\begin{sourceblock}
\includegraphics[page=97,pagebox=mediabox,trim={45.00bp 569.80bp 44.37bp 57.53bp},clip,width=0.92\textwidth,max height=0.38\textheight]{Ishii.pdf}
\end{sourceblock}
上述方程的特殊情况是针对时间分数 k α。在这种情况下，直接从原始定义方程（4-11），我们得到

\begin{sourceblock}
\includegraphics[page=97,pagebox=mediabox,trim={45.00bp 483.69bp 44.37bp 140.56bp},clip,width=0.92\textwidth,max height=0.38\textheight]{Ishii.pdf}
\end{sourceblock}
\begin{sourceblock}
\includegraphics[page=97,pagebox=mediabox,trim={45.00bp 424.77bp 44.37bp 202.36bp},clip,width=0.92\textwidth,max height=0.38\textheight]{Ishii.pdf}
\end{sourceblock}
这些方程清楚地证明了微观奇点的存在，这些奇点的解释与第 4 章 1.3 节中的平滑基本假设有关。

\section*{4.5 浓度和混合物性质}\phantomsection\addcontentsline{toc}{section}{4.5 浓度和混合物性质}
本地时间分数 k α 已在第 4 章 1.2 节中定义。参数 k α 表示任意特定点处的物理事件和两相流的结构。因此，预计本地时间分数 k α 出现在所有场方程中。此外，由于两相本构关系还应取决于流动的物理结构，因此预计其在推导这些定律中的重要性。

除了本地时间分数 k α 之外，还可以定义另一个基于质量的浓度。与扩散理论类似，质量分数 k c 由下式给出

\begin{sourceblock}
\includegraphics[page=97,pagebox=mediabox,trim={45.00bp 166.78bp 44.37bp 457.62bp},clip,width=0.92\textwidth,max height=0.38\textheight]{Ishii.pdf}
\end{sourceblock}
它是第 k 相质量相对于混合物质量的相对重要性的度量。由于动量和能量是质量的加性集函数，因此期望这些变量的混合特性可以用质量分数 k c 作为权重函数的各相的混合特性来表示。由式(4-63)可得

\begin{sourceblock}
\includegraphics[page=98,pagebox=mediabox,trim={45.00bp 568.47bp 44.37bp 60.70bp},clip,width=0.92\textwidth,max height=0.38\textheight]{Ishii.pdf}
\end{sourceblock}
and

\begin{sourceblock}
\includegraphics[page=98,pagebox=mediabox,trim={45.00bp 494.80bp 44.37bp 131.01bp},clip,width=0.92\textwidth,max height=0.38\textheight]{Ishii.pdf}
\end{sourceblock}
上述两个参数，即k α 和k c，是静态浓度，因为它们代表两相流中的事件、结构或质量。此外，运动浓度是通过各种平均速度场定义的，因此它们代表了流量或通量的相对重要性。由于运动学变量的这一基本特征，它们通常不能用 3 维公式来定义，因为流量和通量是矢量而不是标量。但是，可以轻松地为一维模型定义它们。例如，质量 x 在文献中经常使用。接下来，我们定义重要的混合物属性。

1) 混合料密度

\begin{sourceblock}
\includegraphics[page=98,pagebox=mediabox,trim={45.00bp 278.43bp 44.37bp 348.04bp},clip,width=0.92\textwidth,max height=0.38\textheight]{Ishii.pdf}
\end{sourceblock}
where

\begin{sourceblock}
\includegraphics[page=98,pagebox=mediabox,trim={45.00bp 206.80bp 44.37bp 417.52bp},clip,width=0.92\textwidth,max height=0.38\textheight]{Ishii.pdf}
\end{sourceblock}
2) 混合物质心速度

\begin{sourceblock}
\includegraphics[page=98,pagebox=mediabox,trim={45.00bp 116.80bp 44.37bp 495.58bp},clip,width=0.92\textwidth,max height=0.38\textheight]{Ishii.pdf}
\end{sourceblock}
where

\begin{sourceblock}
\includegraphics[page=99,pagebox=mediabox,trim={45.00bp 567.88bp 44.37bp 56.51bp},clip,width=0.92\textwidth,max height=0.38\textheight]{Ishii.pdf}
\end{sourceblock}
3）混合能量

\begin{sourceblock}
\includegraphics[page=99,pagebox=mediabox,trim={45.00bp 477.76bp 44.37bp 134.62bp},clip,width=0.92\textwidth,max height=0.38\textheight]{Ishii.pdf}
\end{sourceblock}
where

\begin{sourceblock}
\includegraphics[page=99,pagebox=mediabox,trim={45.00bp 403.78bp 44.37bp 220.61bp},clip,width=0.92\textwidth,max height=0.38\textheight]{Ishii.pdf}
\end{sourceblock}
4）混合压力

\begin{sourceblock}
\includegraphics[page=99,pagebox=mediabox,trim={45.00bp 330.39bp 44.37bp 296.02bp},clip,width=0.92\textwidth,max height=0.38\textheight]{Ishii.pdf}
\end{sourceblock}
where

\begin{sourceblock}
\includegraphics[page=99,pagebox=mediabox,trim={45.00bp 258.76bp 44.37bp 368.71bp},clip,width=0.92\textwidth,max height=0.38\textheight]{Ishii.pdf}
\end{sourceblock}
5) 混合物的焓

\begin{sourceblock}
\includegraphics[page=99,pagebox=mediabox,trim={45.00bp 168.76bp 44.37bp 443.62bp},clip,width=0.92\textwidth,max height=0.38\textheight]{Ishii.pdf}
\end{sourceblock}
where

\begin{sourceblock}
\includegraphics[page=99,pagebox=mediabox,trim={45.00bp 92.74bp 44.37bp 531.58bp},clip,width=0.92\textwidth,max height=0.38\textheight]{Ishii.pdf}
\end{sourceblock}
那么可以证明

\begin{sourceblock}
\includegraphics[page=100,pagebox=mediabox,trim={45.00bp 571.78bp 44.37bp 55.69bp},clip,width=0.92\textwidth,max height=0.38\textheight]{Ishii.pdf}
\end{sourceblock}
and

\begin{sourceblock}
\includegraphics[page=100,pagebox=mediabox,trim={45.00bp 497.80bp 44.37bp 129.69bp},clip,width=0.92\textwidth,max height=0.38\textheight]{Ishii.pdf}
\end{sourceblock}
6) 混合物熵

\begin{sourceblock}
\includegraphics[page=100,pagebox=mediabox,trim={45.00bp 407.80bp 44.37bp 204.58bp},clip,width=0.92\textwidth,max height=0.38\textheight]{Ishii.pdf}
\end{sourceblock}
where

\begin{sourceblock}
\includegraphics[page=100,pagebox=mediabox,trim={45.00bp 333.88bp 44.37bp 290.51bp},clip,width=0.92\textwidth,max height=0.38\textheight]{Ishii.pdf}
\end{sourceblock}
7) 一般混合物通量 J 这里我们回顾一下第 2 章 1.1 节中定义的一般平衡方程 Eq.(2-6)、广义通量和体积源。从平衡方程的形式，很自然地将混合物分子扩散通量 J 定义为

\begin{sourceblock}
\includegraphics[page=100,pagebox=mediabox,trim={45.00bp 208.47bp 44.37bp 418.00bp},clip,width=0.92\textwidth,max height=0.38\textheight]{Ishii.pdf}
\end{sourceblock}
where

\begin{sourceblock}
\includegraphics[page=100,pagebox=mediabox,trim={45.00bp 134.80bp 44.37bp 489.92bp},clip,width=0.92\textwidth,max height=0.38\textheight]{Ishii.pdf}
\end{sourceblock}
8) 混合物一般源项 m φ 由于源项 φ 被定义为单位质量的变量，因此应以密度来衡量。因此，我们有

\begin{sourceblock}
\includegraphics[page=101,pagebox=mediabox,trim={45.00bp 551.80bp 44.37bp 60.58bp},clip,width=0.92\textwidth,max height=0.38\textheight]{Ishii.pdf}
\end{sourceblock}
where

\begin{sourceblock}
\includegraphics[page=101,pagebox=mediabox,trim={45.00bp 474.88bp 44.37bp 149.51bp},clip,width=0.92\textwidth,max height=0.38\textheight]{Ishii.pdf}
\end{sourceblock}
可以看出，基于单位质量的变量由质量浓度加权，而基于单位体积或表面积的变量由时间分数加权。

\section*{4.6 速度场}\phantomsection\addcontentsline{toc}{section}{4.6 速度场}
一般来说，具有质量、动量和能量传输的两相流系统的特征是存在两种不同的密度和速度。因此，有必要在公式中引入两个正确定义的平均速度场，以考虑相之间相对运动的影响，即质量、动量和能量的扩散。然而，有几个速度场可用于分析两相流问题的各个方面。针对特定问题的速度场的选择取决于流动的特征和性质以及可用本构关系的形式。接下来，我们将介绍这些速度场，它们对于研究两相流系统的各个方面非常重要。

正如上一节所解释的，质心速度的定义是基于线性动量的基本特征。首先，我们记得它是质量的加性集函数。换句话说，正如众所周知的质心基本定理一样，物体的总动量由与物体质量相同的质心动量给出。正是上述思想直接扩展到平均过程中，我们获得了质量加权的混合物速度和相速度作为适当的平均速度。

时间平均中质心的概念很简单，其形式为

\begin{sourceblock}
\includegraphics[page=102,pagebox=mediabox,trim={45.00bp 548.52bp 44.37bp 57.46bp},clip,width=0.92\textwidth,max height=0.38\textheight]{Ishii.pdf}
\end{sourceblock}
然而，0 x 在积分过程中保持恒定，因此质心为 0 x。那么混合物密度的定义，方程（4-66）自然就出来了。质心基本定理可以推广为

\begin{sourceblock}
\includegraphics[page=102,pagebox=mediabox,trim={45.00bp 455.94bp 44.37bp 176.23bp},clip,width=0.92\textwidth,max height=0.38\textheight]{Ishii.pdf}
\end{sourceblock}
因此，根据加权平均值的定义，很容易表明混合物和各相的质心速度分别由方程（4-68）和（4-69）正确给出。

我们定义相对速度

\begin{sourceblock}
\includegraphics[page=102,pagebox=mediabox,trim={45.00bp 345.95bp 44.37bp 293.00bp},clip,width=0.92\textwidth,max height=0.38\textheight]{Ishii.pdf}
\end{sourceblock}
各相的体积通量由下式给出

\begin{sourceblock}
\includegraphics[page=102,pagebox=mediabox,trim={45.00bp 293.08bp 44.37bp 348.01bp},clip,width=0.92\textwidth,max height=0.38\textheight]{Ishii.pdf}
\end{sourceblock}
可以认为是在流量总量固定的情况下，其中一相表面占据整个区间tΔ时的速度，因此也称为空塔速度。因此，混合物体积通量，即体积中心的速度，定义为

\begin{sourceblock}
\includegraphics[page=102,pagebox=mediabox,trim={45.00bp 181.47bp 44.37bp 445.52bp},clip,width=0.92\textwidth,max height=0.38\textheight]{Ishii.pdf}
\end{sourceblock}
如果相间存在相对速度，则由于两相密度不同，速度 m v 和 j 不相等。各相的扩散速度，即相对于混合物质量中心的相对速度，定义为

\begin{sourceblock}
\includegraphics[page=102,pagebox=mediabox,trim={45.00bp 87.03bp 44.37bp 554.00bp},clip,width=0.92\textwidth,max height=0.38\textheight]{Ishii.pdf}
\end{sourceblock}
常用于非均相化学反应单相系统的分析。扩散速度也可以表示为

\begin{sourceblock}
\includegraphics[page=103,pagebox=mediabox,trim={45.00bp 472.54bp 44.37bp 95.15bp},clip,width=0.92\textwidth,max height=0.38\textheight]{Ishii.pdf}
\end{sourceblock}
在两相流系统中，各相的漂移速度，即相对于体积中心的相对速度，非常重要，因为混合物配方中这些速度的本构方程相对简单且完善（Zuber 等，1964；Ishii，1977）。根据定义，漂移速度由下式给出

\begin{sourceblock}
\includegraphics[page=103,pagebox=mediabox,trim={45.00bp 364.96bp 44.37bp 276.20bp},clip,width=0.92\textwidth,max height=0.38\textheight]{Ishii.pdf}
\end{sourceblock}
就相对速度而言，则变为

\begin{sourceblock}
\includegraphics[page=103,pagebox=mediabox,trim={45.00bp 267.40bp 44.37bp 332.85bp},clip,width=0.92\textwidth,max height=0.38\textheight]{Ishii.pdf}
\end{sourceblock}
上述速度之间的几个重要关系可以直接从定义中得到。例如，从方程（4-88）、（4-89）和（4-90）我们得到

\begin{sourceblock}
\includegraphics[page=103,pagebox=mediabox,trim={45.00bp 105.40bp 44.37bp 459.33bp},clip,width=0.92\textwidth,max height=0.38\textheight]{Ishii.pdf}
\end{sourceblock}
由式（4-90）和（4-92）我们有

\begin{sourceblock}
\includegraphics[page=104,pagebox=mediabox,trim={45.00bp 568.47bp 44.37bp 60.73bp},clip,width=0.92\textwidth,max height=0.38\textheight]{Ishii.pdf}
\end{sourceblock}
and

\begin{sourceblock}
\includegraphics[page=104,pagebox=mediabox,trim={45.00bp 495.39bp 44.37bp 133.81bp},clip,width=0.92\textwidth,max height=0.38\textheight]{Ishii.pdf}
\end{sourceblock}
最后，我们注意到如果相对速度为零，那么

\begin{sourceblock}
\includegraphics[page=104,pagebox=mediabox,trim={45.00bp 440.14bp 44.37bp 203.93bp},clip,width=0.92\textwidth,max height=0.38\textheight]{Ishii.pdf}
\end{sourceblock}
因此

\begin{sourceblock}
\includegraphics[page=104,pagebox=mediabox,trim={45.00bp 383.08bp 44.37bp 258.01bp},clip,width=0.92\textwidth,max height=0.38\textheight]{Ishii.pdf}
\end{sourceblock}
它表征了均匀速度场。

一般来说，由于质心基本定理，基于质心的速度对于动力学分析很重要。然而，基于体积中心的速度（即体积通量）对于运动学分析非常有用。如果每个相具有恒定的性质，例如恒定的密度、内能或焓等，则尤其如此。

\section*{4.7 基本恒等式}\phantomsection\addcontentsline{toc}{section}{4.7 基本恒等式}
在开发基于混合物特性的漂移通量模型时，需要用各种平均值来表示平均对流通量。在本节中，我们直接从定义中推导出这种关系。由式(4-29)可得混合物的对流通量为

\begin{sourceblock}
\includegraphics[page=104,pagebox=mediabox,trim={45.00bp 139.47bp 44.37bp 487.51bp},clip,width=0.92\textwidth,max height=0.38\textheight]{Ishii.pdf}
\end{sourceblock}
我们这里的目的是将方程的右侧分成由平均值表示的项和代表波动成分的统计影响的项。由于 ψ 和 v 的平均值按质量加权，波动分量由式（4-37）给出：

\begin{sourceblock}
\includegraphics[page=105,pagebox=mediabox,trim={45.00bp 583.24bp 44.37bp 55.56bp},clip,width=0.92\textwidth,max height=0.38\textheight]{Ishii.pdf}
\end{sourceblock}
with

\begin{sourceblock}
\includegraphics[page=105,pagebox=mediabox,trim={45.00bp 518.14bp 44.37bp 122.13bp},clip,width=0.92\textwidth,max height=0.38\textheight]{Ishii.pdf}
\end{sourceblock}
将式(4-99)和式(4-100)代入式(4-98)可得

\begin{sourceblock}
\includegraphics[page=105,pagebox=mediabox,trim={45.00bp 441.39bp 44.37bp 185.53bp},clip,width=0.92\textwidth,max height=0.38\textheight]{Ishii.pdf}
\end{sourceblock}
通过使用混合物性质和扩散速度的定义，上式可简化为

\begin{sourceblock}
\includegraphics[page=105,pagebox=mediabox,trim={45.00bp 355.47bp 44.37bp 271.51bp},clip,width=0.92\textwidth,max height=0.38\textheight]{Ishii.pdf}
\end{sourceblock}
结果表明，根据传输机制的不同，平均对流通量可分为三部分：基于混合物性质的混合物传输； m 的扩散传输
k ψ 由于相速度的差异；以及由于两相和湍流波动引起的传输。为了区分这最后两种传输机制，我们引入了与它们相关的特殊通量。因此，我们将扩散通量 D J 定义为

\begin{sourceblock}
\includegraphics[page=105,pagebox=mediabox,trim={45.00bp 200.80bp 44.37bp 420.41bp},clip,width=0.92\textwidth,max height=0.38\textheight]{Ishii.pdf}
\end{sourceblock}
而协方差或湍流通量 T k J 定义为

\begin{sourceblock}
\includegraphics[page=105,pagebox=mediabox,trim={45.00bp 139.12bp 44.37bp 501.03bp},clip,width=0.92\textwidth,max height=0.38\textheight]{Ishii.pdf}
\end{sourceblock}
因此，混合物的湍流通量应为

\begin{sourceblock}
\includegraphics[page=105,pagebox=mediabox,trim={45.00bp 66.39bp 44.37bp 560.54bp},clip,width=0.92\textwidth,max height=0.38\textheight]{Ishii.pdf}
\end{sourceblock}
\begin{sourceblock}
\includegraphics[page=106,pagebox=mediabox,trim={45.00bp 530.19bp 44.37bp 96.74bp},clip,width=0.92\textwidth,max height=0.38\textheight]{Ishii.pdf}
\end{sourceblock}
因此，根据平均值的定义，我们有

\begin{sourceblock}
\includegraphics[page=106,pagebox=mediabox,trim={45.00bp 438.86bp 44.37bp 167.91bp},clip,width=0.92\textwidth,max height=0.38\textheight]{Ishii.pdf}
\end{sourceblock}
对于每个阶段我们得到

\begin{sourceblock}
\includegraphics[page=106,pagebox=mediabox,trim={45.00bp 344.53bp 44.37bp 260.02bp},clip,width=0.92\textwidth,max height=0.38\textheight]{Ishii.pdf}
\end{sourceblock}
我们注意到方程（4-108）显示了与单相湍流平均的简单类比，因此，最后一项称为雷诺通量。

由式(4-107)可知，式左边不是用导数的平均值来表示，而是用平均值的导数来表示。然而，当我们将欧拉时间平均应用于两相流的局部公式时，我们首先遇到导数的平均值。现在我们回想一下，这两个运算之间的重要变换已经在第 4 章的 1.4 节中导出。因此，通过将式（4-57）和（4-58）代入式（4-107）我们得到

\begin{sourceblock}
\includegraphics[page=106,pagebox=mediabox,trim={45.00bp 110.37bp 44.37bp 473.99bp},clip,width=0.92\textwidth,max height=0.38\textheight]{Ishii.pdf}
\end{sourceblock}
各个相的类似关系如下，因此我们有

\begin{sourceblock}
\includegraphics[page=107,pagebox=mediabox,trim={45.00bp 526.64bp 44.37bp 55.77bp},clip,width=0.92\textwidth,max height=0.38\textheight]{Ishii.pdf}
\end{sourceblock}
上述两个方程显示了除了波动的统计效应之外界面转移所做出的重要贡献。此外，在混合平均值方程（4-107）中，由于相速度的差异而出现扩散项D J。

此外，根据平均值和扩散通量的定义，我们有

\begin{sourceblock}
\includegraphics[page=107,pagebox=mediabox,trim={45.00bp 352.43bp 44.37bp 233.20bp},clip,width=0.92\textwidth,max height=0.38\textheight]{Ishii.pdf}
\end{sourceblock}
这种关系使我们能够将二流体模型的场方程转换为漂移通量模型的场方程。

\chapter{时间平均平衡方程}
\section*{5.1 一般平衡方程}\phantomsection\addcontentsline{toc}{section}{5.1 一般平衡方程}
前面的章节已经给出了重要的定义以及它们之间的基本关系。我们现在将它们应用于两相流介质中平衡定律的时间平均。正如第 4 章 1.1 节中所解释的，由于平均点周围流体性质的不连续变化，有必要引入几组时间间隔。因此，平均域被分为 [ ] S t Δ 和 [ ] T t Δ。在 [ ] T t Δ 期间，标准平衡方程（2-6）成立，因为占据点 0 x 的流体可以被视为连续体。然而，在[ ] S t Δ 中，界面平衡方程，即第2 章1.2 节的跳跃条件，是有效的，因为界面的特性在这个时间间隔内占主导地位。

我们这里的目的是通过正确分配散装流体和界面的适当平衡方程来及时平均平衡定律。现在让我们首先在[ ] T t Δ 中进行分析，此时平均点被其中一相而不是界面占据。对于时间
[ ] [ ] [ ] 1 2 T t t t t Δ Δ Δ ε = +，我们按以下形式考虑数量 ψ 的平衡

\begin{sourceblock}
\includegraphics[page=108,pagebox=mediabox,trim={45.00bp 116.84bp 44.37bp 508.04bp},clip,width=0.92\textwidth,max height=0.38\textheight]{Ishii.pdf}
\end{sourceblock}
这里 J 和 φ 分别表示广义张量流出和 ψ 的来源。由于它乘以密度 ρ，因此数量 ψ 表示为每单位质量的数量。因此，上述方程本身就是单位体积中数量平衡的数学表述。这个
[ ] S t t Δ ε，由于界面的特殊性，应使用不同类型的平衡方程。由于界面平衡方程的详细推导已在第 2 章 1.2 节中给出，我们简单回顾一下这些结果。因此，由式（2-67）可知，物质平衡
界面处的 ψ 变为

\begin{sourceblock}
\includegraphics[page=109,pagebox=mediabox,trim={45.00bp 367.15bp 44.37bp 201.06bp},clip,width=0.92\textwidth,max height=0.38\textheight]{Ishii.pdf}
\end{sourceblock}
为了得到方程（5-2），我们将方程（2-67）除以界面厚度δ。因此，上式就是该区域单位体积内ψ的平衡。

每体积的平均平衡可以通过在时域中积分适当的平衡方程来获得。现在让我们一般性地表达平衡方程

\begin{sourceblock}
\includegraphics[page=109,pagebox=mediabox,trim={45.00bp 246.75bp 44.37bp 399.24bp},clip,width=0.92\textwidth,max height=0.38\textheight]{Ishii.pdf}
\end{sourceblock}
where

\begin{sourceblock}
\includegraphics[page=109,pagebox=mediabox,trim={45.00bp 189.82bp 44.37bp 450.19bp},clip,width=0.92\textwidth,max height=0.38\textheight]{Ishii.pdf}
\end{sourceblock}
\begin{sourceblock}
\includegraphics[page=109,pagebox=mediabox,trim={45.00bp 144.82bp 44.37bp 495.19bp},clip,width=0.92\textwidth,max height=0.38\textheight]{Ishii.pdf}
\end{sourceblock}
通过取 B 的时间平均，我们有

\begin{sourceblock}
\includegraphics[page=109,pagebox=mediabox,trim={45.00bp 73.29bp 44.37bp 554.69bp},clip,width=0.92\textwidth,max height=0.38\textheight]{Ishii.pdf}
\end{sourceblock}
\begin{sourceblock}
\includegraphics[page=110,pagebox=mediabox,trim={45.00bp 519.15bp 44.37bp 108.76bp},clip,width=0.92\textwidth,max height=0.38\textheight]{Ishii.pdf}
\end{sourceblock}
第一部分可以用第 4 章 1.3 节中定义的平均值来表示，因此从式（5-1）和式（4-25）、（4-57）和（4-58）我们得到

\begin{sourceblock}
\includegraphics[page=110,pagebox=mediabox,trim={45.00bp 381.41bp 44.37bp 204.40bp},clip,width=0.92\textwidth,max height=0.38\textheight]{Ishii.pdf}
\end{sourceblock}
或就混合物性质而言

\begin{sourceblock}
\includegraphics[page=110,pagebox=mediabox,trim={45.00bp 248.99bp 44.37bp 313.36bp},clip,width=0.92\textwidth,max height=0.38\textheight]{Ishii.pdf}
\end{sourceblock}
其中使用了第 4 章第 1.7 节的基本恒等式。从方程（4-2）到方程（4-1），源自接口的第二部分变为

\begin{sourceblock}
\includegraphics[page=110,pagebox=mediabox,trim={45.00bp 112.55bp 44.37bp 474.88bp},clip,width=0.92\textwidth,max height=0.38\textheight]{Ishii.pdf}
\end{sourceblock}
显然，上式是一个时间平均界面平衡方程。为了与局部跳跃条件相区别，我们将其称为界面输运条件或宏观跳跃条件。

根据式(5-7)、(5-9)和(5-10)，我们得到混合物的宏观平衡方程

\begin{sourceblock}
\includegraphics[page=111,pagebox=mediabox,trim={45.00bp 464.74bp 44.37bp 92.91bp},clip,width=0.92\textwidth,max height=0.38\textheight]{Ishii.pdf}
\end{sourceblock}
J 、D J 和 T J 给出的项表示由于平均分子扩散、相对于混合物质心的宏观相扩散以及两相和湍流波动的统计效应而产生的流出量，而 m m ρ φ 是混合物体积源。从平衡方程的形式来看，也可以将界面项视为附加源或汇。

人们普遍认为界面的质量和动量可以忽略不计。然而，由于与热力学张力（即表面张力）相关的能量，表面能可能并非微不足道。因此，界面项的第一部分仅在能量平衡方程中重要。尽管忽略了沿界面的分子扩散传递（即表面粘性应力和表面热通量），但动量和能量平衡中出现了表面线通量。这意味着这些线通量仅考虑表面张力的影响。当将式（5-11）应用于质量、动量和能量平衡时，应采用与式（2-69）、（2-72）和（2-73）的简化跳跃条件相对应的适当形式。

每个相的平均平衡方程可以通过考虑仅与特定相相关的函数来获得，方程（4-8）。因此，与式（5-9）类比，我们有

\begin{sourceblock}
\includegraphics[page=111,pagebox=mediabox,trim={45.00bp 91.45bp 44.37bp 492.62bp},clip,width=0.92\textwidth,max height=0.38\textheight]{Ishii.pdf}
\end{sourceblock}
其中我们使用了变换（4-110）和第 k 相的局部瞬时一般平衡方程，方程（2-6）。为了简单起见，让我们定义

\begin{sourceblock}
\includegraphics[page=112,pagebox=mediabox,trim={45.00bp 562.89bp 44.37bp 55.76bp},clip,width=0.92\textwidth,max height=0.38\textheight]{Ishii.pdf}
\end{sourceblock}
\begin{sourceblock}
\includegraphics[page=112,pagebox=mediabox,trim={45.00bp 473.57bp 44.37bp 119.79bp},clip,width=0.92\textwidth,max height=0.38\textheight]{Ishii.pdf}
\end{sourceblock}
混合物总通量由下式给出

\begin{sourceblock}
\includegraphics[page=112,pagebox=mediabox,trim={45.00bp 417.11bp 44.37bp 223.87bp},clip,width=0.92\textwidth,max height=0.38\textheight]{Ishii.pdf}
\end{sourceblock}
这里 k I 和 m I 分别代表第 k 相和混合物的界面源。根据这些定义，混合物一般平衡方程 (5-11) 简化为

\begin{sourceblock}
\includegraphics[page=112,pagebox=mediabox,trim={45.00bp 316.34bp 44.37bp 305.41bp},clip,width=0.92\textwidth,max height=0.38\textheight]{Ishii.pdf}
\end{sourceblock}
而第 k 相的平衡方程变为

\begin{sourceblock}
\includegraphics[page=112,pagebox=mediabox,trim={45.00bp 225.82bp 44.37bp 378.50bp},clip,width=0.92\textwidth,max height=0.38\textheight]{Ishii.pdf}
\end{sourceblock}
此外，界面转移条件（5-10）可以重写为

\begin{sourceblock}
\includegraphics[page=112,pagebox=mediabox,trim={45.00bp 152.43bp 44.37bp 476.77bp},clip,width=0.92\textwidth,max height=0.38\textheight]{Ishii.pdf}
\end{sourceblock}
这三个宏观方程中的每一个都表达了物质的平衡

ψ 分别代表混合物、第 k 相和界面。混合物平衡方程将是制定漂移通量模型的基础。此外，二流体模型公式需要相平衡方程和界面输运条件。

由式(5-16)、式(5-17)、式(5-18)可知，我们求平均的根本目的已经达到。这样，原来交替占据一个点的两相就变成了两个共存的连续体。此外，极其复杂的两相波动和湍流波动已被平滑化，并且协方差（或湍流通量）项已考虑到它们的统计宏观效应。在接下来的两节中，我们分别介绍扩散模型和双流体模型的质量、动量和能量的平衡方程。

\section*{5.2 二流体模型场方程}\phantomsection\addcontentsline{toc}{section}{5.2 二流体模型场方程}
在本节中，将从时间平均导出的宏观平衡方程（5-17）和界面输运条件（5-18）应用于质量、动量和能量守恒定律。这些方程中变量的选择遵循第 2 章的局部即时公式。

质量平衡 为了获得质量平衡方程，我们设置

\begin{sourceblock}
\includegraphics[page=113,pagebox=mediabox,trim={45.00bp 324.94bp 44.37bp 316.41bp},clip,width=0.92\textwidth,max height=0.38\textheight]{Ishii.pdf}
\end{sourceblock}
鉴于方程（2-69）我们定义

\begin{sourceblock}
\includegraphics[page=113,pagebox=mediabox,trim={45.00bp 242.62bp 44.37bp 376.00bp},clip,width=0.92\textwidth,max height=0.38\textheight]{Ishii.pdf}
\end{sourceblock}
\begin{sourceblock}
\includegraphics[page=113,pagebox=mediabox,trim={45.00bp 200.89bp 44.37bp 443.33bp},clip,width=0.92\textwidth,max height=0.38\textheight]{Ishii.pdf}
\end{sourceblock}
将式(2-7)代入式(5-17)和式(5-18)可得

\begin{sourceblock}
\includegraphics[page=113,pagebox=mediabox,trim={45.00bp 127.31bp 44.37bp 497.68bp},clip,width=0.92\textwidth,max height=0.38\textheight]{Ishii.pdf}
\end{sourceblock}
and

\begin{sourceblock}
\includegraphics[page=113,pagebox=mediabox,trim={45.00bp 56.31bp 44.37bp 572.89bp},clip,width=0.92\textwidth,max height=0.38\textheight]{Ishii.pdf}
\end{sourceblock}
动量平衡 宏观动量平衡可由式（5-13）、（5-14）、（5-17）通过设置得到

\begin{sourceblock}
\includegraphics[page=114,pagebox=mediabox,trim={45.00bp 469.96bp 44.37bp 171.30bp},clip,width=0.92\textwidth,max height=0.38\textheight]{Ishii.pdf}
\end{sourceblock}
\begin{sourceblock}
\includegraphics[page=114,pagebox=mediabox,trim={45.00bp 439.63bp 44.37bp 207.53bp},clip,width=0.92\textwidth,max height=0.38\textheight]{Ishii.pdf}
\end{sourceblock}
\begin{sourceblock}
\includegraphics[page=114,pagebox=mediabox,trim={45.00bp 387.71bp 44.37bp 230.93bp},clip,width=0.92\textwidth,max height=0.38\textheight]{Ishii.pdf}
\end{sourceblock}
\begin{sourceblock}
\includegraphics[page=114,pagebox=mediabox,trim={45.00bp 323.61bp 44.37bp 295.01bp},clip,width=0.92\textwidth,max height=0.38\textheight]{Ishii.pdf}
\end{sourceblock}
\begin{sourceblock}
\includegraphics[page=114,pagebox=mediabox,trim={45.00bp 276.88bp 44.37bp 363.27bp},clip,width=0.92\textwidth,max height=0.38\textheight]{Ishii.pdf}
\end{sourceblock}
\begin{sourceblock}
\includegraphics[page=114,pagebox=mediabox,trim={45.00bp 250.63bp 44.37bp 396.53bp},clip,width=0.92\textwidth,max height=0.38\textheight]{Ishii.pdf}
\end{sourceblock}
\begin{sourceblock}
\includegraphics[page=114,pagebox=mediabox,trim={45.00bp 172.75bp 44.37bp 421.40bp},clip,width=0.92\textwidth,max height=0.38\textheight]{Ishii.pdf}
\end{sourceblock}
and

\begin{sourceblock}
\includegraphics[page=114,pagebox=mediabox,trim={45.00bp 101.19bp 44.37bp 527.98bp},clip,width=0.92\textwidth,max height=0.38\textheight]{Ishii.pdf}
\end{sourceblock}
这里，术语 T k T 和 k M 分别表示湍流通量和形成界面输运的第 k 相动量源，而术语 m M 是由于表面张力效应而产生的混合动量源。

能量平衡 首先设置由式（5-17）可以得到宏观场的能量平衡

\begin{sourceblock}
\includegraphics[page=115,pagebox=mediabox,trim={45.00bp 516.58bp 44.37bp 107.87bp},clip,width=0.92\textwidth,max height=0.38\textheight]{Ishii.pdf}
\end{sourceblock}
\begin{sourceblock}
\includegraphics[page=115,pagebox=mediabox,trim={45.00bp 489.67bp 44.37bp 157.49bp},clip,width=0.92\textwidth,max height=0.38\textheight]{Ishii.pdf}
\end{sourceblock}
\begin{sourceblock}
\includegraphics[page=115,pagebox=mediabox,trim={45.00bp 415.15bp 44.37bp 179.86bp},clip,width=0.92\textwidth,max height=0.38\textheight]{Ishii.pdf}
\end{sourceblock}
\begin{sourceblock}
\includegraphics[page=115,pagebox=mediabox,trim={45.00bp 326.22bp 44.37bp 266.68bp},clip,width=0.92\textwidth,max height=0.38\textheight]{Ishii.pdf}
\end{sourceblock}
\begin{sourceblock}
\includegraphics[page=115,pagebox=mediabox,trim={45.00bp 255.35bp 44.37bp 364.56bp},clip,width=0.92\textwidth,max height=0.38\textheight]{Ishii.pdf}
\end{sourceblock}
\begin{sourceblock}
\includegraphics[page=115,pagebox=mediabox,trim={45.00bp 187.22bp 44.37bp 430.08bp},clip,width=0.92\textwidth,max height=0.38\textheight]{Ishii.pdf}
\end{sourceblock}
\begin{sourceblock}
\includegraphics[page=115,pagebox=mediabox,trim={45.00bp 161.65bp 44.37bp 485.51bp},clip,width=0.92\textwidth,max height=0.38\textheight]{Ishii.pdf}
\end{sourceblock}
\begin{sourceblock}
\includegraphics[page=115,pagebox=mediabox,trim={45.00bp 63.16bp 44.37bp 509.21bp},clip,width=0.92\textwidth,max height=0.38\textheight]{Ishii.pdf}
\end{sourceblock}
\begin{sourceblock}
\includegraphics[page=116,pagebox=mediabox,trim={45.00bp 543.27bp 44.37bp 85.93bp},clip,width=0.92\textwidth,max height=0.38\textheight]{Ishii.pdf}
\end{sourceblock}
我们假设

\begin{sourceblock}
\includegraphics[page=116,pagebox=mediabox,trim={45.00bp 487.83bp 44.37bp 153.20bp},clip,width=0.92\textwidth,max height=0.38\textheight]{Ishii.pdf}
\end{sourceblock}
我们还注意到，视内能 l k e 由标准热能和湍流动能组成，见方程（5-31）。术语
k E 代表第 k 相的界面能量供应，而 m E 是混合物的能量来源。这意味着能量可以存储在界面处或从界面释放。从定义中可以看出，湍流热通量 T k q 考虑了湍流能量对流以及湍流​​功。对于大多数两相流问题，内部加热 k q 可以忽略不计。二流体模型是基于上述六个场方程，即两个连续性方程、两个动量方程和两个能量方程。质量、动量和能量的界面输运条件耦合了各相的传输过程。由于这九个方程基本上表达了守恒定律，因此应辅以各种本构方程来补充，这些本构方程指定了分子扩散、湍流传输和界面输运机制以及热力学状态变量之间的关系。

在解决问题时，将总能量方程中的机械效应和热效应分开通常很有用。因此，根据用速度对动量方程进行点分的标准方法，我们得到了机械能方程

\begin{sourceblock}
\includegraphics[page=116,pagebox=mediabox,trim={45.00bp 111.95bp 44.37bp 455.27bp},clip,width=0.92\textwidth,max height=0.38\textheight]{Ishii.pdf}
\end{sourceblock}
然后将式（5-32）减去式（5-35）即可得到内能方程，从而

\begin{sourceblock}
\includegraphics[page=117,pagebox=mediabox,trim={45.00bp 479.57bp 44.37bp 56.78bp},clip,width=0.92\textwidth,max height=0.38\textheight]{Ishii.pdf}
\end{sourceblock}
这里我们回想一下，虚拟内能 L k e 除了标准内能之外还包括湍流动能。

对于两相流分析，焓能方程很重要，并且经常用于解决各种工程问题。因此，与式（5-31）并行，我们引入了虚焓 l
k h 定义为

\begin{sourceblock}
\includegraphics[page=117,pagebox=mediabox,trim={45.00bp 344.80bp 44.37bp 270.87bp},clip,width=0.92\textwidth,max height=0.38\textheight]{Ishii.pdf}
\end{sourceblock}
将式（5-37）代入式： (5-36) 我们得到

\begin{sourceblock}
\includegraphics[page=117,pagebox=mediabox,trim={45.00bp 203.18bp 44.37bp 352.64bp},clip,width=0.92\textwidth,max height=0.38\textheight]{Ishii.pdf}
\end{sourceblock}
\begin{sourceblock}
\includegraphics[page=117,pagebox=mediabox,trim={63.45bp 143.35bp 60.08bp 470.41bp},clip,width=\textwidth,max height=0.72\textheight]{Ishii.pdf}
\end{sourceblock}
其中，基本导数 k D Dt 是通过跟随第 k 相的质心或以速度 l k v 移动而获得的，因此 k D Dt l k t = ∂∂+ ⋅∇ v。由于湍流波动的机械项与热项之间的相互作用，这些热能方程极其复杂。然而，在许多实际的两相流问题中，传热和相变项在能量方程中占主导地位。在这种情况下，上述方程可以简化为简单的形式。

从式（5-38）可以看出，热能方程中的界面输运具有特殊的形式，由质量、动量和能量传递项的组合来表达。因此我们定义

\begin{sourceblock}
\includegraphics[page=118,pagebox=mediabox,trim={45.00bp 568.69bp 44.37bp 57.38bp},clip,width=0.92\textwidth,max height=0.38\textheight]{Ishii.pdf}
\end{sourceblock}
\section*{5.3 扩散（混合）模型场方程}\phantomsection\addcontentsline{toc}{section}{5.3 扩散（混合）模型场方程}
扩散（混合物）模型的基本概念是将混合物视为一个整体，因此应根据混合物性质编写混合物质量、动量和能量平衡的场方程。然后，这三个宏观混合物守恒方程由考虑浓度变化的扩散方程补充。

混合物连续性和扩散方程由混合物一般平衡方程Eq.(5-16)，根据m ρ 和m v 的定义，我们得到混合物连续性方程

\begin{sourceblock}
\includegraphics[page=118,pagebox=mediabox,trim={45.00bp 352.35bp 44.37bp 247.37bp},clip,width=0.92\textwidth,max height=0.38\textheight]{Ishii.pdf}
\end{sourceblock}
上述方程的形式与没有内部不连续性的连续体的形式完全相同。

\begin{sourceblock}
\includegraphics[page=118,pagebox=mediabox,trim={45.00bp 288.85bp 44.37bp 341.37bp},clip,width=0.92\textwidth,max height=0.38\textheight]{Ishii.pdf}
\end{sourceblock}
\begin{sourceblock}
\includegraphics[page=118,pagebox=mediabox,trim={45.00bp 242.49bp 44.37bp 382.49bp},clip,width=0.92\textwidth,max height=0.38\textheight]{Ishii.pdf}
\end{sourceblock}
它有一个质量源项 1 Γ，仅在连续性方程在时间间隔内求平均值后出现，因为它考虑了界面处的质量传递。另外，式(5-41)的右侧有一个扩散项，因为对流通量已用混合物质心速度m v 表示。混合动量方程 将一般平衡方程 (5-16) 应用于动量守恒，我们得到

\begin{sourceblock}
\includegraphics[page=118,pagebox=mediabox,trim={45.00bp 53.80bp 44.37bp 556.24bp},clip,width=0.92\textwidth,max height=0.38\textheight]{Ishii.pdf}
\end{sourceblock}
\begin{sourceblock}
\includegraphics[page=119,pagebox=mediabox,trim={45.00bp 387.65bp 44.37bp 85.22bp},clip,width=0.92\textwidth,max height=0.38\textheight]{Ishii.pdf}
\end{sourceblock}
此外，界面动量源m M 由式(5-24)给出。

三个张量通量 T 、T T 和 D T 分别代表平均粘性应力、湍流应力和扩散应力。显然，如果忽略表面张力项，则混合物动量方程中不存在直接的界面项。

混合物总能量方程 将式（5-16）应用于总能量平衡，可得混合物能量方程，因此

\begin{sourceblock}
\includegraphics[page=119,pagebox=mediabox,trim={45.00bp 137.86bp 44.37bp 419.93bp},clip,width=0.92\textwidth,max height=0.38\textheight]{Ishii.pdf}
\end{sourceblock}
我们在哪里

\begin{sourceblock}
\includegraphics[page=119,pagebox=mediabox,trim={45.00bp 64.23bp 44.37bp 562.28bp},clip,width=0.92\textwidth,max height=0.38\textheight]{Ishii.pdf}
\end{sourceblock}
\begin{sourceblock}
\includegraphics[page=120,pagebox=mediabox,trim={45.00bp 542.10bp 44.37bp 55.12bp},clip,width=0.92\textwidth,max height=0.38\textheight]{Ishii.pdf}
\end{sourceblock}
\begin{sourceblock}
\includegraphics[page=120,pagebox=mediabox,trim={45.00bp 442.49bp 44.37bp 144.22bp},clip,width=0.92\textwidth,max height=0.38\textheight]{Ishii.pdf}
\end{sourceblock}
并且，根据定义，我们具有以下混合属性

\begin{sourceblock}
\includegraphics[page=120,pagebox=mediabox,trim={45.00bp 346.77bp 44.37bp 252.84bp},clip,width=0.92\textwidth,max height=0.38\textheight]{Ishii.pdf}
\end{sourceblock}
\begin{sourceblock}
\includegraphics[page=120,pagebox=mediabox,trim={45.00bp 268.06bp 44.37bp 339.46bp},clip,width=0.92\textwidth,max height=0.38\textheight]{Ishii.pdf}
\end{sourceblock}
界面能量m E 由式(5-29)给出。重要的特殊情况是当身体力场恒定时

\begin{sourceblock}
\includegraphics[page=120,pagebox=mediabox,trim={45.00bp 199.12bp 44.37bp 441.97bp},clip,width=0.92\textwidth,max height=0.38\textheight]{Ishii.pdf}
\end{sourceblock}
那么扩散体功项变为零，因此

\begin{sourceblock}
\includegraphics[page=120,pagebox=mediabox,trim={45.00bp 126.40bp 44.37bp 502.81bp},clip,width=0.92\textwidth,max height=0.38\textheight]{Ishii.pdf}
\end{sourceblock}
我们使用恒等式（4-94）。

因此，在恒定体积力场的标准条件下，总混合能量方程简化为

\begin{sourceblock}
\includegraphics[page=121,pagebox=mediabox,trim={45.00bp 519.70bp 44.37bp 55.75bp},clip,width=0.92\textwidth,max height=0.38\textheight]{Ishii.pdf}
\end{sourceblock}
可见，式（5-52）的形式与单相流能量方程十分相似。这些差异表现为附加热通量，即湍流通量 T q 和扩散通量 D q 以及界面体源 m E。然而，混合物最有趣的特性可以在动能项中找到，方程（5-49）。从方程中我们看到，总混合物动能由质心动能加上两相的扩散动能组成。还应该记住，由于将湍流动能与热效应分开非常困难，因此湍流动能已包含在虚拟内能中。我们还指出，如果忽略表面张力效应，则界面项不会出现在混合物总能量方程中（与混合物动量方程的情况相同）。

混合热能方程 在单相流中，通过从总能量平衡中减去机械能方程，可以很容易地分离机械能和热能。如果我们像方程（5-36）和（5-37）中所做的那样将湍流动能包含在虚拟热能中，则可以在双流体模型公式中使用完全相同的方法。然而，在扩散模型公式中，由于扩散动能传输的存在，情况变得更加复杂。因此，没有明确的方法来获得混合物的相应热能方程。在下文中，我们演示了两种不同的方法，它们给出了截然不同的结果。

第一种方法是从总能量方程(5-44)中减去两相动能方程之和。这样就可以消除扩散动能。因此，由式（5-35）和式（5-44）可得

\begin{sourceblock}
\includegraphics[page=121,pagebox=mediabox,trim={45.00bp 50.14bp 44.37bp 541.32bp},clip,width=0.92\textwidth,max height=0.38\textheight]{Ishii.pdf}
\end{sourceblock}
\begin{sourceblock}
\includegraphics[page=122,pagebox=mediabox,trim={45.00bp 550.71bp 44.37bp 76.27bp},clip,width=0.92\textwidth,max height=0.38\textheight]{Ishii.pdf}
\end{sourceblock}
\begin{sourceblock}
\includegraphics[page=122,pagebox=mediabox,trim={45.00bp 452.80bp 44.37bp 124.56bp},clip,width=0.92\textwidth,max height=0.38\textheight]{Ishii.pdf}
\end{sourceblock}
将各相的焓方程（5-38）相加也可以得到相同的方程（5-53）。除了最后一项外，方程的形式相当简单，但我们应该意识到界面项

\begin{sourceblock}
\includegraphics[page=122,pagebox=mediabox,trim={45.00bp 369.97bp 44.37bp 265.70bp},clip,width=0.92\textwidth,max height=0.38\textheight]{Ishii.pdf}
\end{sourceblock}
energies.

值得注意的是，通过使用以质心速度 m v 表示的混合物动能方程，可以避免界面项中的这一困难。然而，所得的热能方程具有来自扩散动能的附加项。从式(5-52)中减去混合机械能方程，即由m v 点缀的动量方程(5-42)，可得

\begin{sourceblock}
\includegraphics[page=122,pagebox=mediabox,trim={45.00bp 90.08bp 44.37bp 398.39bp},clip,width=0.92\textwidth,max height=0.38\textheight]{Ishii.pdf}
\end{sourceblock}
由式（5-53）和式（5-55）这两个热能方程可以看出，混合物能量传递高度

\section*{5.4 vni=0的奇异情况（准静态接口）}\phantomsection\addcontentsline{toc}{section}{5.4 vni=0的奇异情况（准静态接口）}
在前面的分析中，假设界面位移速度 ni v 不为零，但实际上在孤立奇点处它可以为零。例如，当界面静止或界面的运动与其完全相切时，就会发生这种情况。自从
0 ni v = 是与平衡方程中的所有界面项（例如 k I 和 m I ）相关的重要奇点，我们对其进行了一些详细研究。关于这种奇点，我们首先引入单位体积的表面积浓度。通过仅考虑一个接口，它可以给出为

\begin{sourceblock}
\includegraphics[page=123,pagebox=mediabox,trim={45.00bp 264.04bp 44.37bp 356.02bp},clip,width=0.92\textwidth,max height=0.38\textheight]{Ishii.pdf}
\end{sourceblock}
而总面积浓度由下式给出

\begin{sourceblock}
\includegraphics[page=123,pagebox=mediabox,trim={45.00bp 150.34bp 44.37bp 434.33bp},clip,width=0.92\textwidth,max height=0.38\textheight]{Ishii.pdf}
\end{sourceblock}
S L 给出的 i a 是两相流的内部长度尺度。该变量将在第 10 章中详细讨论。

现在我们回到 0 ni v = 的奇异情况。鉴于式(5-56)和(5-57)，可以说在这种情况下传输长度S L 也变得奇异，从而失去其物理意义。因此，平衡方程中的所有界面项 k I 和 m I 也是奇异的，并且时间分数 S α 或其导数可能会出现不连续性。为了解决这个困难，我们首先回顾一下 S α 是界面的时间分数。我们在这里注意到，从方程（4-14）中，情况
0 ni v ≠ 对应于 0 S α =。此外，如果界面 ni v 的法向速度为零，则它可能会在有限时间内保持在 0 x。因此我们有

\begin{sourceblock}
\includegraphics[page=124,pagebox=mediabox,trim={45.00bp 310.22bp 44.37bp 314.73bp},clip,width=0.92\textwidth,max height=0.38\textheight]{Ishii.pdf}
\end{sourceblock}
式中，S t Δ 为t Δ 时间内接口占用的总时间。因此

\begin{sourceblock}
\includegraphics[page=124,pagebox=mediabox,trim={45.00bp 226.23bp 44.37bp 400.27bp},clip,width=0.92\textwidth,max height=0.38\textheight]{Ishii.pdf}
\end{sourceblock}
或者就时间分数而言

\begin{sourceblock}
\includegraphics[page=124,pagebox=mediabox,trim={45.00bp 170.92bp 44.37bp 470.31bp},clip,width=0.92\textwidth,max height=0.38\textheight]{Ishii.pdf}
\end{sourceblock}
现在很明显，我们有两种不同的奇点情况，即 1. 0 0 和 ; ni S v α == 2. 0 0 和 . ni S v α => 正如 k α 的定义所示，情况 1 并没有给 k α 带来不连续性，而只是在微观意义上导致其导数不连续性。然而，第 4 章 1.3 节已就基本原理进行了解释。

我们注意到，案例 2 所代表的宏观界面最容易通过固定界面（例如实体墙）来例证。因为对于
0 S α > 不可能考虑体积平衡方程的平均值，宏观平衡的正确形式只是t Δ 时间间隔内跳跃条件的时间平均。因此我们有

\begin{sourceblock}
\includegraphics[page=125,pagebox=mediabox,trim={45.00bp 337.16bp 44.37bp 266.76bp},clip,width=0.92\textwidth,max height=0.38\textheight]{Ishii.pdf}
\end{sourceblock}
正常两相流中的情况2对应于k α 的不连续性，因此，它可以被视为浓度冲击，见方程（4-61）和（4-62）。对于大部分流场，假设它是连续的。由于时间平均宏观公式旨在应用于此类两相流，因此在大多数应用中可以忽略这些界面奇点。

\section*{5.5 宏观跳跃条件}\phantomsection\addcontentsline{toc}{section}{5.5 宏观跳跃条件}
我们在上一节中讨论了与准稳态界面相关的奇点。现在我们研究与时间平均场中宏观激波不连续性相关的重要奇点。分析的基本部分可以与第 2 章的 1.2 节并行进行，其中导出了接口处的标准跳转条件。

众所周知，在单相流和两相流中，存在各种属性发生极大变化的区域。它可以通过压缩性效应引起的冲击波或混合物中的浓度冲击来举例说明。这些区域异常高的梯度需要对本构关系进行特殊考虑，以便将它们视为连续介质力学的一部分。然而，正如一直以来

\begin{sourceblock}
\includegraphics[page=126,pagebox=mediabox,trim={50.00bp 409.32bp 44.37bp 55.00bp},clip,width=\textwidth,max height=0.72\textheight]{Ishii.pdf}
\par\vspace{0.45\baselineskip}
\small 图 5-1。宏观不连续性（Ishii，1975）
\end{sourceblock}
在第 2 章 1.2 节中提到，对于大多数实际流动问题，用具有跳跃条件的不连续面替换这些区域会产生足够准确的模型。因此，在本节中，我们推导出宏观跳跃条件，该条件作为平均场中不连续表面的平衡方程。通过将时间平均宏观场视为连续体，可以直接应用第 2 章 1.2 节的分析。表面U的速度的定义类似于式(2-49)和(2-50)中的iv，因此其法向分量是表面位移速度，而其切向分量是厚度δ区域内的平均混合切向速度。通过用 + 和 - 表示区域的每一边，我们可以从方程（2-64）

\begin{sourceblock}
\includegraphics[page=126,pagebox=mediabox,trim={45.00bp 118.92bp 44.37bp 461.00bp},clip,width=0.92\textwidth,max height=0.38\textheight]{Ishii.pdf}
\end{sourceblock}
式中，求和代表区域两侧，下标a和s代表式（2-45）至（2-54）定义的区域内的平均值，表面导数具有标准显著性。方程(5-62)是冲击区域m ψ 的平衡。的左侧
显然，对于两相中的每一相也可以得到类似于式（5-62）的平衡方程，因此

\begin{sourceblock}
\includegraphics[page=127,pagebox=mediabox,trim={45.00bp 424.73bp 44.37bp 148.16bp},clip,width=0.92\textwidth,max height=0.38\textheight]{Ishii.pdf}
\end{sourceblock}
从该区域平均值的定义很容易看出，如果将方程（5-63）中的每一项对两个相求和，则得到的项就成为方程（5-62）中出现的混合项。

由于式（5-62）和（5-63）引入了与不连续性相关的新变量，即表面性质和线通量，因此有必要对这些项做出一些具体假设或给出足够的本构关系。通过考虑 0 δ → 的极限，并且忽略冲击的表面能和相关的热力学张力，可以获得具有实际重要性的简单结果。在这些条件下，我们对混合物有

\begin{sourceblock}
\includegraphics[page=127,pagebox=mediabox,trim={45.00bp 240.83bp 44.37bp 387.97bp},clip,width=0.92\textwidth,max height=0.38\textheight]{Ishii.pdf}
\end{sourceblock}
对于 k = 1 和 2 的每个阶段

\begin{sourceblock}
\includegraphics[page=127,pagebox=mediabox,trim={45.00bp 169.92bp 44.37bp 458.91bp},clip,width=0.92\textwidth,max height=0.38\textheight]{Ishii.pdf}
\end{sourceblock}
在这里，我们注意到术语 ma I 和 ka I 的重要性，它们允许激波层内质量、动量和能量的交换。仅仅因为这些术语显示为体积源而忽略它们是不正确的。根据表达它们的本构关系，它们可能采取通量的形式。此外，从物理角度来看，由于两相激波层的高度非平衡状态，在宏观跳跃条件下自然存在这些界面输运项 ka I。

\section*{5.6 宏观场方程和跳跃条件总结}\phantomsection\addcontentsline{toc}{section}{5.6 宏观场方程和跳跃条件总结}
\begin{sourceblock}
\includegraphics[page=128,pagebox=mediabox,trim={45.00bp 535.27bp 44.37bp 98.93bp},clip,width=0.92\textwidth,max height=0.38\textheight]{Ishii.pdf}
\end{sourceblock}
\begin{sourceblock}
\includegraphics[page=128,pagebox=mediabox,trim={45.00bp 489.75bp 44.37bp 135.15bp},clip,width=0.92\textwidth,max height=0.38\textheight]{Ishii.pdf}
\end{sourceblock}
and

\begin{sourceblock}
\includegraphics[page=128,pagebox=mediabox,trim={45.00bp 399.22bp 44.37bp 205.11bp},clip,width=0.92\textwidth,max height=0.38\textheight]{Ishii.pdf}
\end{sourceblock}
混合物总通量定义为

\begin{sourceblock}
\includegraphics[page=128,pagebox=mediabox,trim={45.00bp 341.50bp 44.37bp 299.48bp},clip,width=0.92\textwidth,max height=0.38\textheight]{Ishii.pdf}
\end{sourceblock}
with

\begin{sourceblock}
\includegraphics[page=128,pagebox=mediabox,trim={45.00bp 268.83bp 44.37bp 357.67bp},clip,width=0.92\textwidth,max height=0.38\textheight]{Ishii.pdf}
\end{sourceblock}
\begin{sourceblock}
\includegraphics[page=128,pagebox=mediabox,trim={45.00bp 208.83bp 44.37bp 417.67bp},clip,width=0.92\textwidth,max height=0.38\textheight]{Ishii.pdf}
\end{sourceblock}
\begin{sourceblock}
\includegraphics[page=128,pagebox=mediabox,trim={45.00bp 148.83bp 44.37bp 478.14bp},clip,width=0.92\textwidth,max height=0.38\textheight]{Ishii.pdf}
\end{sourceblock}
界面转移条件为

\begin{sourceblock}
\includegraphics[page=128,pagebox=mediabox,trim={45.00bp 75.87bp 44.37bp 550.63bp},clip,width=0.92\textwidth,max height=0.38\textheight]{Ishii.pdf}
\end{sourceblock}
\begin{sourceblock}
\includegraphics[page=129,pagebox=mediabox,trim={45.00bp 536.75bp 44.37bp 91.95bp},clip,width=0.92\textwidth,max height=0.38\textheight]{Ishii.pdf}
\end{sourceblock}
and

\begin{sourceblock}
\includegraphics[page=129,pagebox=mediabox,trim={45.00bp 465.89bp 44.37bp 162.93bp},clip,width=0.92\textwidth,max height=0.38\textheight]{Ishii.pdf}
\end{sourceblock}
其中 n（带有 + 和 -）是表面每侧的向外单位法向矢量，U 表示冲击表观速度。

\section*{5.7 湍流热流的替代形式}\phantomsection\addcontentsline{toc}{section}{5.7 湍流热流的替代形式}
第5章1.2和1.3节中的能量方程是根据式(5-30)的湍流热通量定义推导的。这里，方程（5-30）’的湍流热通量的另一种定义可能为

\begin{sourceblock}
\includegraphics[page=129,pagebox=mediabox,trim={45.00bp 264.38bp 44.37bp 341.77bp},clip,width=0.92\textwidth,max height=0.38\textheight]{Ishii.pdf}
\end{sourceblock}
然后，我们可以将停滞内能方程Eq.(5-32)、内能方程Eq.(5-36)和停滞焓方程Eq.(5-38)改写为分别由方程(5-32)’、(5-36)’和(5-38)’给出。

\begin{sourceblock}
\includegraphics[page=129,pagebox=mediabox,trim={45.00bp 104.87bp 44.37bp 456.45bp},clip,width=0.92\textwidth,max height=0.38\textheight]{Ishii.pdf}
\end{sourceblock}
\begin{sourceblock}
\includegraphics[page=130,pagebox=mediabox,trim={45.00bp 458.87bp 44.37bp 57.14bp},clip,width=0.92\textwidth,max height=0.38\textheight]{Ishii.pdf}
\end{sourceblock}
\begin{sourceblock}
\includegraphics[page=130,pagebox=mediabox,trim={45.00bp 332.14bp 44.37bp 223.77bp},clip,width=0.92\textwidth,max height=0.38\textheight]{Ishii.pdf}
\end{sourceblock}
我们还可以重新构建总混合物能量方程 Eq.(5-44) 和混合物停滞焓方程 Eq.(5-53)，分别由方程 (5-44)’ 和 (5-53)’ 给出。

\begin{sourceblock}
\includegraphics[page=130,pagebox=mediabox,trim={45.00bp 139.83bp 44.37bp 392.68bp},clip,width=0.92\textwidth,max height=0.38\textheight]{Ishii.pdf}
\end{sourceblock}
我们在哪里

\begin{sourceblock}
\includegraphics[page=131,pagebox=mediabox,trim={45.00bp 540.19bp 44.37bp 55.18bp},clip,width=0.92\textwidth,max height=0.38\textheight]{Ishii.pdf}
\end{sourceblock}
\begin{sourceblock}
\includegraphics[page=131,pagebox=mediabox,trim={45.00bp 440.38bp 44.37bp 146.31bp},clip,width=0.92\textwidth,max height=0.38\textheight]{Ishii.pdf}
\end{sourceblock}
\begin{sourceblock}
\includegraphics[page=131,pagebox=mediabox,trim={45.00bp 313.65bp 44.37bp 241.54bp},clip,width=0.92\textwidth,max height=0.38\textheight]{Ishii.pdf}
\end{sourceblock}
相应地，方程（5-52）和（5-55）可以分别改写为方程（5-52）’和（5-55）’。

\begin{sourceblock}
\includegraphics[page=131,pagebox=mediabox,trim={45.00bp 172.16bp 44.37bp 396.88bp},clip,width=0.92\textwidth,max height=0.38\textheight]{Ishii.pdf}
\end{sourceblock}
\begin{sourceblock}
\includegraphics[page=131,pagebox=mediabox,trim={45.00bp 77.55bp 44.37bp 513.17bp},clip,width=0.92\textwidth,max height=0.38\textheight]{Ishii.pdf}
\end{sourceblock}
\begin{sourceblock}
\includegraphics[page=132,pagebox=mediabox,trim={45.00bp 480.96bp 44.37bp 73.96bp},clip,width=0.92\textwidth,max height=0.38\textheight]{Ishii.pdf}
\end{sourceblock}
通过湍流热通量的替代定义，例如方程（5-30）’，能量方程关于应力张量项变得对称。式（5-30）定义的能量方程与式（5-30）定义的能量方程相同。然而，具有式（5-30）’定义的能量方程可能因以下原因而引起争议和误导。例如，由于湍流应力而产生的表观能量耗散项如 l : k k k Τ α ∇ v T 出现在式（5-36）’的内能方程中。由于此项是由对流或湍流引起的，因此它可能取负值或正值。因此，由于湍流功而产生的术语可能不被视为能量耗散，而能量耗散应始终为正。为了使湍流功项始终为正，可以对熵不等式应用一些附加约束。因此，具有式（5-30）’定义的能量方程对于解释不可逆热力学是有问题的。

\chapter{与其他统计平均的联系}
\section*{6.1 欧拉统计平均值（系综平均值）}\phantomsection\addcontentsline{toc}{section}{6.1 欧拉统计平均值（系综平均值）}
欧拉统计平均的基本概念已在第3章中解释。通过考虑一组N个相似的样本或系统，统计平均值由它们之间的简单算术平均值定义，方程（3-7）。因此，时间平均的积分数学运算应该被统计平均中的求和数学运算所取代。显然，基于统计平均的场论的整个推导可以与第 4 章和第 5 章并行进行，只需用有限统计平均算子方程（3-7）代替时间平均方程（3-2）即可。

最重要的参数是 k N 和 S N，它们代表当界面厚度为 δ 时第 k 相和界面出现的次数。然后，空隙率可以定义为 k N 与 N 的比率，并取极限 0 δ →。类似于式(4-8)，可以在时空域中定义通用函数( ) n F 和( ) k n F。

由于统计平均方程的推导遵循与时间平均情况完全相同的步骤，因此我们仅列出它们之间最重要和最典型的关系。在表 6-1 中，我们看到了平均的四个基本参数，即： 空隙率 k α；两相通用函数 F 的平均值；第k相通用函数k F 的平均值；以及单位体积内的界面面积浓度。显然，方程（5-16）和（5-17）形式的一般平衡方程可以从应用于一组的统计平均中获得
N 个相似样本。我们在此注意到，时间和统计平均之间的重要差异并不存在于余额的结果形式中
统计平均时间平均变量 统计平均时间平均变量

\begin{sourceblock}
\includegraphics[page=134,pagebox=mediabox,trim={45.00bp 433.23bp 44.37bp 99.42bp},clip,width=0.92\textwidth,max height=0.38\textheight]{Ishii.pdf}
\end{sourceblock}
方程，但在相对于实际流量的变量解释中，正如第 3 章 1.3 节中详细讨论的那样。

\section*{6.2 玻尔兹曼统计平均值}\phantomsection\addcontentsline{toc}{section}{6.2 玻尔兹曼统计平均值}
由于其在各种平均过程中的独特特性，我们详细讨论了应用于两相流系统的玻尔兹曼统计平均。首先我们回想一下 ( ) f t x , , xi 是粒子密度函数，其中 x 、 t 和 xi 分别代表粒子的位置、时间和速度。在气体动力学理论的标准分析中，各组分的颗粒质量被认为是恒定的，因为它代表了分子质量。然而，在两相高度分散流的应用中，可能有必要假设颗粒质量发生变化。因此，通过考虑可变粒子质量的存在，我们定义

\begin{sourceblock}
\includegraphics[page=134,pagebox=mediabox,trim={45.00bp 177.10bp 44.37bp 464.14bp},clip,width=0.92\textwidth,max height=0.38\textheight]{Ishii.pdf}
\end{sourceblock}
其中 kn f 是质量为 kn m 的 k 个相颗粒的颗粒密度函数。可以说kn·m是单个分子质量的倍数。 x 和 xi 处相空间元素 d d x ψ 中第 k 相粒子的总数由下式给出

\begin{sourceblock}
\includegraphics[page=134,pagebox=mediabox,trim={45.00bp 68.44bp 44.37bp 560.84bp},clip,width=0.92\textwidth,max height=0.38\textheight]{Ishii.pdf}
\end{sourceblock}
引入外力场 k g，并取粒子数平衡，可得到 kn m 个粒子的玻尔兹曼方程，即

\begin{sourceblock}
\includegraphics[page=135,pagebox=mediabox,trim={45.00bp 514.81bp 44.37bp 106.67bp},clip,width=0.92\textwidth,max height=0.38\textheight]{Ishii.pdf}
\end{sourceblock}
其中 kn C + 和 kn C − 表示源项和汇项，即由粒子质量的变化以及将粒子抛入和抛出相元素 d ε 的碰撞引起的 kn m 个粒子的增益和损失。如果 kn g 与速度无关，那么我们有

\begin{sourceblock}
\includegraphics[page=135,pagebox=mediabox,trim={45.00bp 396.85bp 44.37bp 224.63bp},clip,width=0.92\textwidth,max height=0.38\textheight]{Ishii.pdf}
\end{sourceblock}
上述带有用麦克斯韦二元碰撞积分的简单模型表示的碰撞项的玻尔兹曼输运方程称为麦克斯韦-玻尔兹曼方程。它是气体动力学理论的基础。类似的方法可用于中子传输。

\begin{sourceblock}
\includegraphics[page=135,pagebox=mediabox,trim={45.00bp 276.82bp 44.37bp 333.88bp},clip,width=0.92\textwidth,max height=0.38\textheight]{Ishii.pdf}
\end{sourceblock}
因此，基于概率 kn f 的期望可以定义为

\begin{sourceblock}
\includegraphics[page=135,pagebox=mediabox,trim={45.00bp 190.29bp 44.37bp 420.92bp},clip,width=0.92\textwidth,max height=0.38\textheight]{Ishii.pdf}
\end{sourceblock}
由于总质量是每个颗粒的质量之和，因此第 k 相的部分密度由下式给出

\begin{sourceblock}
\includegraphics[page=135,pagebox=mediabox,trim={45.00bp 110.38bp 44.37bp 518.95bp},clip,width=0.92\textwidth,max height=0.38\textheight]{Ishii.pdf}
\end{sourceblock}
由此可见，质量加权平均值定义为

\begin{sourceblock}
\includegraphics[page=136,pagebox=mediabox,trim={45.00bp 544.95bp 44.37bp 56.59bp},clip,width=0.92\textwidth,max height=0.38\textheight]{Ishii.pdf}
\end{sourceblock}
用 l k v 表示平均速度，每个粒子的特有速度由下式给出

\begin{sourceblock}
\includegraphics[page=136,pagebox=mediabox,trim={45.00bp 476.08bp 44.37bp 165.01bp},clip,width=0.92\textwidth,max height=0.38\textheight]{Ishii.pdf}
\end{sourceblock}
如果我们将麦克斯韦-玻尔兹曼输运方程，即方程（6-4）乘以
kn kn m ψ，对相速度 ψ 积分，然后对第 k 相的各种粒子求和，得到平均值的麦克斯韦传递方程。因此，

\begin{sourceblock}
\includegraphics[page=136,pagebox=mediabox,trim={45.00bp 262.90bp 44.37bp 259.55bp},clip,width=0.92\textwidth,max height=0.38\textheight]{Ishii.pdf}
\end{sourceblock}
我们注意到，方程（6-10）右侧的每一项代表由于粒子间扩散引起的传递、由于体力场引起的源、当输运性质 kn ψ 是相速度 ψ 的函数时出现的多分子粒子效应，以及由于相变和/或碰撞引起的源。

质量守恒 通过在式（6-10）中设定 1 kn ψ = 即可得到质量场方程，因此有

\begin{sourceblock}
\includegraphics[page=136,pagebox=mediabox,trim={45.00bp 86.43bp 44.37bp 538.54bp},clip,width=0.92\textwidth,max height=0.38\textheight]{Ishii.pdf}
\end{sourceblock}
where

\begin{sourceblock}
\includegraphics[page=137,pagebox=mediabox,trim={45.00bp 573.58bp 44.37bp 55.75bp},clip,width=0.92\textwidth,max height=0.38\textheight]{Ishii.pdf}
\end{sourceblock}
总体质量守恒定律的假设给出

\begin{sourceblock}
\includegraphics[page=137,pagebox=mediabox,trim={45.00bp 500.43bp 44.37bp 128.77bp},clip,width=0.92\textwidth,max height=0.38\textheight]{Ishii.pdf}
\end{sourceblock}
由式(6-11)和式(6-13)可得混合连续性方程

\begin{sourceblock}
\includegraphics[page=137,pagebox=mediabox,trim={45.00bp 425.44bp 44.37bp 196.42bp},clip,width=0.92\textwidth,max height=0.38\textheight]{Ishii.pdf}
\end{sourceblock}
where

\begin{sourceblock}
\includegraphics[page=137,pagebox=mediabox,trim={45.00bp 353.80bp 44.37bp 270.45bp},clip,width=0.92\textwidth,max height=0.38\textheight]{Ishii.pdf}
\end{sourceblock}
回想一下 k ρ 是第 k 相的部分密度，我们发现时间平均的结果（方程（5-21）、（5-22）和（5-40））与玻尔兹曼统计平均的结果（方程（6-11）、（6-13）和（6-14）完全相似。

动量守恒 由式（6-10）可得线性动量方程，设 kn ψ = ψ，则

\begin{sourceblock}
\includegraphics[page=137,pagebox=mediabox,trim={45.00bp 187.34bp 44.37bp 409.38bp},clip,width=0.92\textwidth,max height=0.38\textheight]{Ishii.pdf}
\end{sourceblock}
\begin{sourceblock}
\includegraphics[page=137,pagebox=mediabox,trim={45.00bp 122.44bp 44.37bp 492.47bp},clip,width=0.92\textwidth,max height=0.38\textheight]{Ishii.pdf}
\end{sourceblock}
\begin{sourceblock}
\includegraphics[page=137,pagebox=mediabox,trim={45.00bp 54.58bp 44.37bp 552.33bp},clip,width=0.92\textwidth,max height=0.38\textheight]{Ishii.pdf}
\end{sourceblock}
这里我们注意到，如果分压张量 k − P 的负值被解释为作用在第 k 相上的组合应力，则动量方程 (6-16) 的形式与通过时间平均获得的形式完全相同，即方程 (5-26)。然而，该方程中通量项和动量源 k M 的物理意义可能与时间平均方程中的物理意义显著不同。

作为一个例子，让我们考虑一个分散两相流系统，其中第 1 相作为分散相。如果颗粒尺寸相当小并且连续相是稀释气体，则配方基本上减少到反应气体混合物的配方。在这种情况下，粒子与位于真实空间中的特定体积元素外部的分子和/或粒子的碰撞的影响可以被忽略。因此，总碰撞项
1 k k = Σ M 可以视为零。因此，由式（6-18）可知，
每个相的碰撞项由质量传递产生的动量源和碰撞过程中动量交换产生的阻力组成。

此外，如果与之前的点质量假设相反，粒子具有确定的体积，则粒子与分子以及中心位于体积元素之外的其他粒子的多次碰撞变得重要。在这种情况下，分散相的碰撞项可以分为三个不同的部分：由于碰撞而产生的内部动量传递；相变化的影响；以及外部碰撞的影响。因此，在配方中引入空隙率 k α 是必要的，其中 1 α 是颗粒所占体积与总体积元素的比率，因此 2 1 1 α α = -。然而，为了积分动量碰撞项，有必要引入分子-粒子和粒子-粒子碰撞过程的模型。由于粒子和分子云之间的多次碰撞以及相变的影响，这些模型将极其复杂。因此，我们在这里不详细讨论碰撞积分。如果引入标准流体力学的观点，就可以将动量源k M 解释为相关的物理术语。根据上面的讨论，我们将各相的总碰撞项展开为以下形式

\begin{sourceblock}
\includegraphics[page=138,pagebox=mediabox,trim={45.00bp 114.93bp 44.37bp 526.04bp},clip,width=0.92\textwidth,max height=0.38\textheight]{Ishii.pdf}
\end{sourceblock}
右侧的每一项分别代表内部碰撞力相互作用、由于相变而产生的动量源以及由于有限颗粒尺寸而导致的外部碰撞力相互作用。第一项可以被认为是标准阻力和压力

\begin{sourceblock}
\includegraphics[page=139,pagebox=mediabox,trim={45.00bp 546.81bp 44.37bp 85.31bp},clip,width=0.92\textwidth,max height=0.38\textheight]{Ishii.pdf}
\end{sourceblock}
\begin{sourceblock}
\includegraphics[page=139,pagebox=mediabox,trim={45.00bp 489.70bp 44.37bp 116.73bp},clip,width=0.92\textwidth,max height=0.38\textheight]{Ishii.pdf}
\end{sourceblock}
其中 m M 可以视为粒子间碰撞效应。第二项是连续相对位于体积单元边界上的粒子的影响。连续相 2 的实际压力张量应为

\begin{sourceblock}
\includegraphics[page=139,pagebox=mediabox,trim={45.00bp 375.67bp 44.37bp 249.22bp},clip,width=0.92\textwidth,max height=0.38\textheight]{Ishii.pdf}
\end{sourceblock}
那么混合动量方程变为

\begin{sourceblock}
\includegraphics[page=139,pagebox=mediabox,trim={45.00bp 266.81bp 44.37bp 320.97bp},clip,width=0.92\textwidth,max height=0.38\textheight]{Ishii.pdf}
\end{sourceblock}
因此，混合物总应力由分散相的部分应力、连续相的真实应力和由于两相之间的相对运动而产生的扩散应力组成。这里，平均扩散速度km V 的定义采用式(4-89)的形式。可见，基于玻尔兹曼输运方程的混合动量方程的应力张量项的物理意义与时间平均得到的应力张量项的物理意义存在显著差异，即式(5-42)。出现这种基本差异的原因是，在前一种方法中，应力是根据粒子的运动来定义的。因此，颗粒内部的应力在分析中不占有一席之地。然而，在后一种方法中，应力是在系统中的各处定义的。我们可以在这里补充一点，通过将涡流传输视为粒子传输，可以轻松地将玻尔兹曼统计平均值扩展到包括连续相中的湍流波动。

能量守恒与简单的动力学理论相反，由于粒子能量状态的内部自由度，多分子粒子的能量传递相当复杂。显然，如果每个粒子由大量分子组成，基于扩散动能的动力学理论平移温度就没有用处。

\begin{sourceblock}
\includegraphics[page=140,pagebox=mediabox,trim={45.00bp 455.47bp 44.37bp 149.76bp},clip,width=0.92\textwidth,max height=0.38\textheight]{Ishii.pdf}
\end{sourceblock}
其中第一项是平移动能，kn u 是粒子中分子所包含的内能。由于 kn u 项不是 kn xi 的函数，因此输运方程 (6-10) 不是玻尔兹曼输运方程的速度矩。

我们将第 k 相的平均能量和通量引入为

\begin{sourceblock}
\includegraphics[page=140,pagebox=mediabox,trim={45.00bp 315.63bp 44.37bp 292.55bp},clip,width=0.92\textwidth,max height=0.38\textheight]{Ishii.pdf}
\end{sourceblock}
and

\begin{sourceblock}
\includegraphics[page=140,pagebox=mediabox,trim={45.00bp 240.34bp 44.37bp 381.20bp},clip,width=0.92\textwidth,max height=0.38\textheight]{Ishii.pdf}
\end{sourceblock}
其中右侧平均值的定义遵循式（6-8）。将式（6-23）、（6-24）、（6-25）代入式（6-10），可得第k相的总能量方程

\begin{sourceblock}
\includegraphics[page=140,pagebox=mediabox,trim={45.00bp 82.99bp 44.37bp 482.65bp},clip,width=0.92\textwidth,max height=0.38\textheight]{Ishii.pdf}
\end{sourceblock}
平均内能 l k e 是随机热平动能与粒子真实内能之和。这个定义与时间平均的方程(5-31)完全类似。最重要的特征出现在方程（6-26）的最后两项中。很明显，最后第二项给出的单个粒子内能的变化与碰撞项 kE 耦合，因为在没有长程能量交换的情况下，粒子内能仅通过与周围分子和粒子的相互作用而变化。因此，与方程（6-20）给出的动量交换类似，总能量相互作用可以给出为

\begin{sourceblock}
\includegraphics[page=141,pagebox=mediabox,trim={45.00bp 378.54bp 44.37bp 197.82bp},clip,width=0.92\textwidth,max height=0.38\textheight]{Ishii.pdf}
\end{sourceblock}
这里方程（6-27）右边的第一项考虑了有限体积粒子的粒子-分子碰撞。 m E 表示的项表示粒子间碰撞能量传输等效应。通过将每个相的两个能量方程相加，可以得到混合方程

\begin{sourceblock}
\includegraphics[page=141,pagebox=mediabox,trim={45.00bp 176.63bp 44.37bp 368.93bp},clip,width=0.92\textwidth,max height=0.38\textheight]{Ishii.pdf}
\end{sourceblock}
在这里，我们再次注意到相 1 是分散相。总能量方程每一项的物理意义与动量方程的物理意义相似。分散相的部分通量 1 q 和 1 P 以及连续相的真实通量 2 q 和 P 2 是由于相 1 的颗粒占据的有限体积而出现的。

与混合动量方程的情况一样，玻尔兹曼统计模型和时间平均模型之间的显著差异出现在热和功的总通量项中。由于前一个模型没有考虑颗粒内部的传热和做功，因此部分
将玻尔兹曼统计平均值应用于两相流系统的兴趣主要在于研究碰撞项的简单模型以及随机假设的本构方程。这不太可能成为分散两相流系统的通用模型，因为玻尔兹曼输运方程中包含颗粒形状和变形的影响会给分析带来相当大的困难。即使没有这些效应，由于三种效应，预计每个相的碰撞项也会非常复杂：1）粒子间碰撞、合并和分裂； 2）一个粒子与大量分子之间的多重碰撞； 3）相变的存在。此外，从式（6-27）可以看出，分散相能量传递项需要特殊的颗粒与流体之间的传热本构关系。

总结本节，可以说玻尔兹曼统计平均对于高度分散的流非常有用，其中每个粒子被视为集中实体而不是分布式系统本身。例如，颗粒流的混合应力张量在本模型中具有比时间平均的形式更自然的形式。这是因为在这种两相流中，在颗粒内部引入应力是不切实际的。然而，为了使数量密度的统计处理更加现实，体积单元中的粒子数量应该非常大。此外，如果界面的变形和颗粒内的性质变化很重要，则不能使用玻尔兹曼统计方法。与稀气体的情况相比，有限粒径的分散两相体系的碰撞积分项极其复杂，因此从统计力学中获得本构关系非常困难。在许多情况下，这些碰撞项应从连续介质力学考虑中提供。因此，也可以结合连续介质理论和玻尔兹曼统计方法来构建模型。例如，我们仅对分散相取统计平均值，其中阻力包含在体积力中。然后我们使用另一相的标准体积平均场方程以及相应的相互作用项。这种制剂可用于分散两相系统以及流化床。此外，玻尔兹曼统计方法对于获得一些本构关系非常有用，如第 10 章所示。

\chapter{平均场运动学}
\section*{7.1 对流坐标和对流导数}\phantomsection\addcontentsline{toc}{section}{7.1 对流坐标和对流导数}
时间平均一致地通过空间描述来表达，如定义（4-15）和（4-16）所示，并且由于相变和扩散，平均两相流场的粒子坐标的想法并不清晰也不平凡。相变对应于整个场中每个相的流体粒子的产生或消失。之所以出现困难，是因为每个相本身显然并不遵守连续性公理的推论，即物质的永久性。然而，各相的扩散允许其他流体颗粒渗透混合物颗粒。显然，作为标准连续介质力学基础的材料坐标并不是从时间平均获得的一般两相流场所固有的。然而，可以引入数学上特殊的对流坐标，这对于研究每个相和混合物的运动学很有用。

各相的路径线由系统的积分曲线定义

\begin{sourceblock}
\includegraphics[page=143,pagebox=mediabox,trim={45.00bp 158.15bp 44.37bp 483.07bp},clip,width=0.92\textwidth,max height=0.38\textheight]{Ishii.pdf}
\end{sourceblock}
与初始条件

\begin{sourceblock}
\includegraphics[page=143,pagebox=mediabox,trim={45.00bp 99.05bp 44.37bp 545.14bp},clip,width=0.92\textwidth,max height=0.38\textheight]{Ishii.pdf}
\end{sourceblock}
其中我们将 k X 定义为第 k 相的对流坐标。因此，对式（7-1）积分，我们得到

\begin{sourceblock}
\includegraphics[page=144,pagebox=mediabox,trim={45.00bp 582.95bp 44.37bp 59.33bp},clip,width=0.92\textwidth,max height=0.38\textheight]{Ishii.pdf}
\end{sourceblock}
该方程给出了对流坐标 k X 上的固定点的路径线，该固定点以粒子速度 l k v 移动。根据光滑度的标准假设，或者雅可比行列式的存在，我们可以将方程（7-3）变换为

\begin{sourceblock}
\includegraphics[page=144,pagebox=mediabox,trim={45.00bp 487.12bp 44.37bp 156.96bp},clip,width=0.92\textwidth,max height=0.38\textheight]{Ishii.pdf}
\end{sourceblock}
该方程表示以第 k 相 l k v 的局部平均速度移动的假想粒子的位置。以x和t为自变量的问题的表述称为空间描述，而以k X 和t为自变量的问题的表述称为对流描述。一般来说，由于相变，占据 ( ) k t = x x X 邻域的粒子的含量可能与初始粒子不同。因此，不可能考虑固定粒子的变化。然而，观察具有固定 k X 的过程很简单。例如，第 k 相的速度可以与单相流类似地给出：

\begin{sourceblock}
\includegraphics[page=144,pagebox=mediabox,trim={45.00bp 295.23bp 44.37bp 328.00bp},clip,width=0.92\textwidth,max height=0.38\textheight]{Ishii.pdf}
\end{sourceblock}
上述分析也可以应用于混合物质心，因此我们定义混合物路径线为

\begin{sourceblock}
\includegraphics[page=144,pagebox=mediabox,trim={45.00bp 225.16bp 44.37bp 418.91bp},clip,width=0.92\textwidth,max height=0.38\textheight]{Ishii.pdf}
\end{sourceblock}
其中 m = x X 在 0 t t =。对式(7-6)积分可得路径线

\begin{sourceblock}
\includegraphics[page=144,pagebox=mediabox,trim={45.00bp 161.92bp 44.37bp 480.32bp},clip,width=0.92\textwidth,max height=0.38\textheight]{Ishii.pdf}
\end{sourceblock}
通过逆变换我们得到

\begin{sourceblock}
\includegraphics[page=144,pagebox=mediabox,trim={45.00bp 105.16bp 44.37bp 538.91bp},clip,width=0.92\textwidth,max height=0.38\textheight]{Ishii.pdf}
\end{sourceblock}
因此，如果混合对流坐标固定在方程（7-7）中，则观察者以局部混合速度 m v 移动。然而，由于m X 沿着路径线不断变化。

此外，有趣的是，如果流场是均匀的，或者 l l 1 2 m == v v v，那么无论相位如何变化，混合对流坐标都会变成材料坐标。从方程（7-6）和（7-7）中，混合物速度可以象征性地给出：

\begin{sourceblock}
\includegraphics[page=145,pagebox=mediabox,trim={45.00bp 477.03bp 44.37bp 146.19bp},clip,width=0.92\textwidth,max height=0.38\textheight]{Ishii.pdf}
\end{sourceblock}
从式（7-1）和式（7-6）容易看出，各相和混合物的路径线可以相互交叉。

由于欧拉时间均值是空间描述，所以固定点的时间变化率表示为

\begin{sourceblock}
\includegraphics[page=145,pagebox=mediabox,trim={45.00bp 364.17bp 44.37bp 260.19bp},clip,width=0.92\textwidth,max height=0.38\textheight]{Ishii.pdf}
\end{sourceblock}
然而，从观察者看到的随流体速度移动的变化率称为对流导数或实质导数。它由下式给出

\begin{sourceblock}
\includegraphics[page=145,pagebox=mediabox,trim={45.00bp 277.05bp 44.37bp 346.18bp},clip,width=0.92\textwidth,max height=0.38\textheight]{Ishii.pdf}
\end{sourceblock}
and

\begin{sourceblock}
\includegraphics[page=145,pagebox=mediabox,trim={45.00bp 200.43bp 44.37bp 422.80bp},clip,width=0.92\textwidth,max height=0.38\textheight]{Ishii.pdf}
\end{sourceblock}
式(7-11)和式(7-12)的对流导数分别是沿着第k相和混合物的质心求得的。

如果用相位对流导数来表达场方程(5-17)左边，我们得到

\begin{sourceblock}
\includegraphics[page=145,pagebox=mediabox,trim={45.00bp 70.65bp 44.37bp 548.66bp},clip,width=0.92\textwidth,max height=0.38\textheight]{Ishii.pdf}
\end{sourceblock}
\begin{sourceblock}
\includegraphics[page=146,pagebox=mediabox,trim={45.00bp 531.14bp 44.37bp 93.73bp},clip,width=0.92\textwidth,max height=0.38\textheight]{Ishii.pdf}
\end{sourceblock}
我们注意到，质量源项的贡献出现在方程（7-13）中，因为表观速度为 l k v 的体积内的质量量不是恒定的。通过结合基本恒等式的推论，方程（4-111）和上述两个关系，我们在混合对流导数和相对流导数之间进行了重要的变换，因此

\begin{sourceblock}
\includegraphics[page=146,pagebox=mediabox,trim={45.00bp 405.09bp 44.37bp 215.57bp},clip,width=0.92\textwidth,max height=0.38\textheight]{Ishii.pdf}
\end{sourceblock}
\section*{7.2 精简}\phantomsection\addcontentsline{toc}{section}{7.2 精简}
停滞点定义为所有速度都消失的点，因此

\begin{sourceblock}
\includegraphics[page=146,pagebox=mediabox,trim={45.00bp 320.92bp 44.37bp 320.17bp},clip,width=0.92\textwidth,max height=0.38\textheight]{Ishii.pdf}
\end{sourceblock}
其中一个相的 l k v 为零的点称为第 k 相驻点。如果混合速度 m v 在某一点为零，我们称其为伪驻点。在这一点上，两相的运动是纯粹的扩散。如果每个相速度与时间无关，则流动是完全稳定的

\begin{sourceblock}
\includegraphics[page=146,pagebox=mediabox,trim={45.00bp 213.88bp 44.37bp 427.21bp},clip,width=0.92\textwidth,max height=0.38\textheight]{Ishii.pdf}
\end{sourceblock}
如果 ( ) mm = v v x，则混合运动是稳定的，但它并不对应于完全稳定的运动，因为扩散速度可以是时间的函数。

矢量场的矢量线是一条处处与该矢量相切的曲线。特别地，速度场l k v 的矢量线被称为第k相流线。因此它可以由联立方程的积分曲线给出

\begin{sourceblock}
\includegraphics[page=146,pagebox=mediabox,trim={45.00bp 81.04bp 44.37bp 560.17bp},clip,width=0.92\textwidth,max height=0.38\textheight]{Ishii.pdf}
\end{sourceblock}
其中 l 是沿流线的参数。一般来说，流线型是

\begin{sourceblock}
\includegraphics[page=147,pagebox=mediabox,trim={50.00bp 448.14bp 44.37bp 55.00bp},clip,width=\textwidth,max height=0.72\textheight]{Ishii.pdf}
\par\vspace{0.45\baselineskip}
\small 图7-1。流线（Ishii，1975）是时间的函数并且与路径线不重合。每个相的两条流线也不同，因为两个速度场不平行。混合物的流线可以类似地定义为
\end{sourceblock}
\begin{sourceblock}
\includegraphics[page=147,pagebox=mediabox,trim={45.00bp 352.42bp 44.37bp 288.75bp},clip,width=0.92\textwidth,max height=0.38\textheight]{Ishii.pdf}
\end{sourceblock}
各种流线之间的关系如图7-1所示。这里我们还注意到，各相流线的重合并不意味着流场均匀。

\section*{7.3 质量守恒}\phantomsection\addcontentsline{toc}{section}{7.3 质量守恒}
基于质心速度的公式 各相的连续性方程已在第 5 章 1.2 节中推导，因此我们有

\begin{sourceblock}
\includegraphics[page=147,pagebox=mediabox,trim={45.00bp 172.82bp 44.37bp 452.24bp},clip,width=0.92\textwidth,max height=0.38\textheight]{Ishii.pdf}
\end{sourceblock}
界面传质条件由下式给出

\begin{sourceblock}
\includegraphics[page=147,pagebox=mediabox,trim={45.00bp 101.79bp 44.37bp 527.38bp},clip,width=0.92\textwidth,max height=0.38\textheight]{Ishii.pdf}
\end{sourceblock}
方程（7-20）简单地表明，每单位体积的局部密度 k k α ρ 的局部时间变化率等于第 k 相 l ( ) k k k α ρ −∇⋅ v 的净质量流入加上由于相变而产生的质量源。正如第 5 章 1.2 节所示，通过将这三个方程相加，我们得到

\begin{sourceblock}
\includegraphics[page=148,pagebox=mediabox,trim={45.00bp 515.17bp 44.37bp 106.62bp},clip,width=0.92\textwidth,max height=0.38\textheight]{Ishii.pdf}
\end{sourceblock}
这是混合物的连续性方程。

为了指定两相混合物中的质量守恒，有必要采用两个连续性方程。我们通过方程（7-20）中每个相的质心速度表达了这些关系，然而，现在通过引入不同的观测观点来考虑连续性方程的替代形式是很有趣的。如果观察者随混合质心移动，则扩散项明确出现在相连续性方程中，因此我们有

\begin{sourceblock}
\includegraphics[page=148,pagebox=mediabox,trim={45.00bp 350.26bp 44.37bp 271.61bp},clip,width=0.92\textwidth,max height=0.38\textheight]{Ishii.pdf}
\end{sourceblock}
基于质量分数的公式 除了使用时间分数 k α 之外，我们还可以用方程（4-63）定义的质量分数 k c 将上式表示为

\begin{sourceblock}
\includegraphics[page=148,pagebox=mediabox,trim={45.00bp 250.53bp 44.37bp 348.87bp},clip,width=0.92\textwidth,max height=0.38\textheight]{Ishii.pdf}
\end{sourceblock}
其中我们使用了混合连续性方程。此外，与非均相单相混合物类似，扩散系数 k D 可用于表示扩散通量，因此我们设置

\begin{sourceblock}
\includegraphics[page=148,pagebox=mediabox,trim={45.00bp 168.94bp 44.37bp 472.37bp},clip,width=0.92\textwidth,max height=0.38\textheight]{Ishii.pdf}
\end{sourceblock}
这里我们注意到，只有当扩散是由于浓度梯度引起的并且可以用线性本构律表示时，方程（7-25）才是正确的。然而，预计菲克扩散定律可能不适用于一般的两相流系统，因为界面几何形状、体力场和界面动量传递项是影响相扩散的重要因素。式（7-25）给出的线性本构关系与牛顿粘度定律和傅里叶热传导定律完全类似。这些线性本构关系是
通过考虑方程（7-24）的非常简化的形式，可以表明具有方程（7-25）的线性本构律的扩散方程由于方程中k c 的二阶导数而表现出浓度k c 的扩散特性。这与本节后面部分基于运动波速度的公式形成鲜明对比，后者表现出传播的特性。

基于体积通量的公式 连续性关系也可以通过体积通量 j 和漂移速度 kj V 来表达。因此，从式（7-20）和（4-91）我们有

\begin{sourceblock}
\includegraphics[page=149,pagebox=mediabox,trim={45.00bp 356.74bp 44.37bp 265.21bp},clip,width=0.92\textwidth,max height=0.38\textheight]{Ishii.pdf}
\end{sourceblock}
上式右侧的最后一项表示第 k 相质量相对于混合物体积中心的漂移。对式(7-20)左边部分微分，可得

\begin{sourceblock}
\includegraphics[page=149,pagebox=mediabox,trim={45.00bp 256.90bp 44.37bp 364.22bp},clip,width=0.92\textwidth,max height=0.38\textheight]{Ishii.pdf}
\end{sourceblock}
其中实质导数由式(7-11)定义。上式可以被认为是时间分数或空隙率的连续性方程，因此它代表体积传输。因此，从 k α 的角度来看，连续性方程具有由于质量传递而产生的源项和由于相的真实压缩性而产生的汇项。此外，如果我们使用方程（4-87），我们得到

\begin{sourceblock}
\includegraphics[page=149,pagebox=mediabox,trim={45.00bp 115.90bp 44.37bp 505.23bp},clip,width=0.92\textwidth,max height=0.38\textheight]{Ishii.pdf}
\end{sourceblock}
通过将每个相的这两个方程相加，我们得到

\begin{sourceblock}
\includegraphics[page=150,pagebox=mediabox,trim={45.00bp 562.89bp 44.37bp 55.75bp},clip,width=0.92\textwidth,max height=0.38\textheight]{Ishii.pdf}
\end{sourceblock}
它描述了体积速度中心的散度。右侧第一项是由于相变而产生的体积源，第二项是由于可压缩性而产生的体积沉。

如果每个相都经历由下式定义的不可压缩或等容过程，那么基于体积通量的公式就很重要

\begin{sourceblock}
\includegraphics[page=150,pagebox=mediabox,trim={45.00bp 440.62bp 44.37bp 185.53bp},clip,width=0.92\textwidth,max height=0.38\textheight]{Ishii.pdf}
\end{sourceblock}
在这种情况下，方程（7-29）简化为

\begin{sourceblock}
\includegraphics[page=150,pagebox=mediabox,trim={45.00bp 367.84bp 44.37bp 256.51bp},clip,width=0.92\textwidth,max height=0.38\textheight]{Ishii.pdf}
\end{sourceblock}
它简单地说明体积通量的散度与相变量以及每个相的比体积之间的差异成正比。我们记得，对于不可压缩的单相流，速度发散为零。因此，两相流等价是用体积中心的速度而不是质量的速度来表示的，此外，它还有一个由于相变而产生的来源。可以看出，在没有传质的情况下，上式简化为
0 ∇⋅ = j。由于在许多实际的两相流问题中，不可压缩流体假设是有效的，并且相变速率k Г可以表示为位置和时间的已知函数，因此方程（7-31）可以在解决这些问题中发挥重要作用。

运动波速度和空洞传播方程与单相混合物扩散热机械理论一样，漂移通量（或混合物）模型公式的基本假设之一是两相之间的相对运动可以用本构关系而不是两个动量方程来表达。在这方面，我们已经讨论了菲克扩散定律，它有效地消除了两个动量方程之一。从式（7-25）可以清楚地看出，本构关系给出了 l k v 和 m v 之间的特殊运动学关系，因此在混合动量方程存在的情况下，第 k 相动量方程变得多余。

然而，我们注意到，一般来说，对两相混合物使用菲克定律是不正确的，因此应在漂移通量（或混合物）模型公式中使用不同类型的扩散本构关系。相间相对运动更有用的本构关系之一是用漂移速度 kj V 来表示（Zuber et al., 1964; Ishii, 1977; Kataoka and Ishii, 1987; Hibiki and Ishii, 2003a, 2003b; Hibiki et al., 2003; Goda et al., 2003; Hibiki等，2004）。特别是如果漂移速度仅是时间浓度 k α 的函数，则流体力学中非常重要的理论之一，即运动波理论（Kynch，1952；Lighthill 和 Whitham，1955；Hayes，1970）可以应用于两相流系统。 Zuber (1964b) 表明，许多具有实际意义的流态都是这种情况，并且它对于分散流态特别有用。

在不可压缩流体假设下，式（7-27）可以用 j 和 kj V 表示为以下形式

\begin{sourceblock}
\includegraphics[page=151,pagebox=mediabox,trim={45.00bp 345.67bp 44.37bp 260.71bp},clip,width=0.92\textwidth,max height=0.38\textheight]{Ishii.pdf}
\end{sourceblock}
\begin{sourceblock}
\includegraphics[page=151,pagebox=mediabox,trim={45.00bp 284.98bp 44.37bp 339.58bp},clip,width=0.92\textwidth,max height=0.38\textheight]{Ishii.pdf}
\end{sourceblock}
将式(7-33)代入式(7-32)可得空洞传播方程

\begin{sourceblock}
\includegraphics[page=151,pagebox=mediabox,trim={45.00bp 194.58bp 44.37bp 426.62bp},clip,width=0.92\textwidth,max height=0.38\textheight]{Ishii.pdf}
\end{sourceblock}
\begin{sourceblock}
\includegraphics[page=151,pagebox=mediabox,trim={45.00bp 118.54bp 44.37bp 483.54bp},clip,width=0.92\textwidth,max height=0.38\textheight]{Ishii.pdf}
\end{sourceblock}
因此，将运动波之后的特殊对流导数表示为

\begin{sourceblock}
\includegraphics[page=152,pagebox=mediabox,trim={45.00bp 569.43bp 44.37bp 55.54bp},clip,width=0.92\textwidth,max height=0.38\textheight]{Ishii.pdf}
\end{sourceblock}
空隙传播方程简化为

\begin{sourceblock}
\includegraphics[page=152,pagebox=mediabox,trim={45.00bp 499.87bp 44.37bp 124.54bp},clip,width=0.92\textwidth,max height=0.38\textheight]{Ishii.pdf}
\end{sourceblock}
因此，如果我们通过随运动波速度移动来观察 k α 的时间变化率，则由于相位变化，它与源项成正比。在相间不存在质量传递的情况下，k α 的扰动随运动波速传播。此外，在ρ 1 和ρ 2 恒定的条件下，我们可以将式（7-34）用混合物密度表示如下：

\begin{sourceblock}
\includegraphics[page=152,pagebox=mediabox,trim={45.00bp 295.93bp 44.37bp 265.00bp},clip,width=0.92\textwidth,max height=0.38\textheight]{Ishii.pdf}
\end{sourceblock}
这称为密度传播方程。

运动冲击波已经表明，如果漂移速度仅是k α 的函数，则空隙率方程可以转化为空隙传播方程。相反，如果相的扩散可以用菲克扩散定律形式的本构方程来表示，则场表现出扩散介质的特性，并且由于空间中的二阶导数而无法观察到清晰的空洞传播。

对于几种类型的混合物，前一种空穴传播现象是已知的。例如，它们在明渠、气泡两相流和高速公路交通流中非常重要。在此类系统中，观察到在某些条件下这些运动波传播导致浓度激波的形成。由于其起源以及由于压缩效应而需要将其与冲击波区分开来，我们将其称为运动冲击波。正如 Lighthill 和 Whitham (1955) 以及 Kynch (1952) 所表明的，这种现象可以通过运动学来分析

\begin{sourceblock}
\includegraphics[page=153,pagebox=mediabox,trim={45.00bp 484.79bp 44.37bp 144.01bp},clip,width=0.92\textwidth,max height=0.38\textheight]{Ishii.pdf}
\end{sourceblock}
and

\begin{sourceblock}
\includegraphics[page=153,pagebox=mediabox,trim={45.00bp 417.83bp 44.37bp 210.97bp},clip,width=0.92\textwidth,max height=0.38\textheight]{Ishii.pdf}
\end{sourceblock}
显然，方程（7-39）表达了总质量守恒，而方程（7-40）给出了激波上第k相质量的平衡。这里 ka Г 表示激波层内的相变量。因此，方程（7-40）表明，由于激波中的相变，激波每一侧的第 k 相质量通量与批量生产相平衡。

求解式(7-40)求出激波位移速度，可得

\begin{sourceblock}
\includegraphics[page=153,pagebox=mediabox,trim={45.00bp 256.24bp 44.37bp 359.11bp},clip,width=0.92\textwidth,max height=0.38\textheight]{Ishii.pdf}
\end{sourceblock}
其中+和-表示冲击层的每一侧。应该记住，方程（7-41）给出的条件不仅适用于运动冲击波，而且还适用于由于压缩性效应而产生的动态冲击波。

通过将我们的情况限制为严格的运动学现象，我们在这里假设相密度在整个激波中是连续的，并且层中没有相变化。因此我们有

\begin{sourceblock}
\includegraphics[page=153,pagebox=mediabox,trim={45.00bp 130.83bp 44.37bp 508.01bp},clip,width=0.92\textwidth,max height=0.38\textheight]{Ishii.pdf}
\end{sourceblock}
and

\begin{sourceblock}
\includegraphics[page=153,pagebox=mediabox,trim={45.00bp 75.88bp 44.37bp 565.31bp},clip,width=0.92\textwidth,max height=0.38\textheight]{Ishii.pdf}
\end{sourceblock}
因此式(7-40)简化为

\begin{sourceblock}
\includegraphics[page=154,pagebox=mediabox,trim={45.00bp 575.03bp 44.37bp 53.77bp},clip,width=0.92\textwidth,max height=0.38\textheight]{Ishii.pdf}
\end{sourceblock}
由此可见，运动冲击波速度U应满足

\begin{sourceblock}
\includegraphics[page=154,pagebox=mediabox,trim={45.00bp 497.74bp 44.37bp 122.33bp},clip,width=0.92\textwidth,max height=0.38\textheight]{Ishii.pdf}
\end{sourceblock}
如果我们使用方程（4-87）定义的体积通量 k j 的定义，我们有

\begin{sourceblock}
\includegraphics[page=154,pagebox=mediabox,trim={45.00bp 407.73bp 44.37bp 215.00bp},clip,width=0.92\textwidth,max height=0.38\textheight]{Ishii.pdf}
\end{sourceblock}
由式(7-46)可得

\begin{sourceblock}
\includegraphics[page=154,pagebox=mediabox,trim={45.00bp 351.08bp 44.37bp 290.85bp},clip,width=0.92\textwidth,max height=0.38\textheight]{Ishii.pdf}
\end{sourceblock}
这意味着在简单的运动冲击中，即相密度连续并且冲击中没有相变化，总体积通量 j 是守恒的。鉴于式(7-47)，我们可以将式(7-46)变换为如下形式

\begin{sourceblock}
\includegraphics[page=154,pagebox=mediabox,trim={45.00bp 234.76bp 44.37bp 387.15bp},clip,width=0.92\textwidth,max height=0.38\textheight]{Ishii.pdf}
\end{sourceblock}
这里我们看到方程（7-35）给出的运动波速度与方程（7-48）给出的激波位移速度之间的密切联系。

\section*{7.4 膨胀}\phantomsection\addcontentsline{toc}{section}{7.4 膨胀}
第 7 章 1.1 节的对流和空间坐标的雅可比行列式由下式给出

\begin{sourceblock}
\includegraphics[page=154,pagebox=mediabox,trim={45.00bp 68.10bp 44.37bp 546.56bp},clip,width=0.92\textwidth,max height=0.38\textheight]{Ishii.pdf}
\end{sourceblock}
\begin{sourceblock}
\includegraphics[page=155,pagebox=mediabox,trim={45.00bp 534.94bp 44.37bp 109.54bp},clip,width=0.92\textwidth,max height=0.38\textheight]{Ishii.pdf}
\end{sourceblock}
其中 k V 和 0 k V 分别表示当前体积和初始体积。鉴于式(7-5)我们有

\begin{sourceblock}
\includegraphics[page=155,pagebox=mediabox,trim={45.00bp 446.55bp 44.37bp 180.76bp},clip,width=0.92\textwidth,max height=0.38\textheight]{Ishii.pdf}
\end{sourceblock}
这给出了相速度发散的重要物理解释。回顾第 k 相的连续性方程，它也可以表示为

\begin{sourceblock}
\includegraphics[page=155,pagebox=mediabox,trim={45.00bp 345.57bp 44.37bp 276.42bp},clip,width=0.92\textwidth,max height=0.38\textheight]{Ishii.pdf}
\end{sourceblock}
这两个方程表明 l k v 的散度与体积元素的膨胀直接相关，而不是与体积内的密度变化相关。此外，我们看到膨胀是由三种效应引起的，即相位变化 k Γ 、相位重新分布 k k D Dt α 和
流体的实际压缩率 k k D Dt ρ。对于作为整体的混合物，膨胀和的特征
m ∇⋅ v 基本上简化为单相流。因此我们有

\begin{sourceblock}
\includegraphics[page=155,pagebox=mediabox,trim={45.00bp 186.57bp 44.37bp 434.57bp},clip,width=0.92\textwidth,max height=0.38\textheight]{Ishii.pdf}
\end{sourceblock}
\begin{sourceblock}
\includegraphics[page=155,pagebox=mediabox,trim={45.00bp 110.50bp 44.37bp 488.42bp},clip,width=0.92\textwidth,max height=0.38\textheight]{Ishii.pdf}
\end{sourceblock}
\chapter{界面输运}
质量、动量和能量交换的界面输运项 k I 和 m I 的精确形式已在第 5 章 1.2 节中给出。然而，它们是用局部瞬时变量表示的，因此不可能将它们用作平均场方程中的本构关系。显然，在构建两相流系统的任何特定本构方程之前，我们需要详细理解这些术语的物理含义。考虑到这一点，我们阐明了控制这些术语的不同物理机制，并确定了它们所依赖的重要参数。此外，重要的是要接受并非局部瞬时两相流固有的所有特征都可以带入时间平均模型中。我们认为平均场方程表达了控制宏观两相流的一般物理原理，而本构方程则用简单的数学模型近似了特定系统组的材料响应。在这方面，我们对界面输运项做出了一些假设，以便区分主要的传递机制，并消除一些在宏观领域影响不显著的复杂项。

\section*{8.1 界面传质}\phantomsection\addcontentsline{toc}{section}{8.1 界面传质}
界面传质项 k Γ 由式（5-19）给出，因此根据式（2-70）和（5-57）我们有

\begin{sourceblock}
\includegraphics[page=156,pagebox=mediabox,trim={45.00bp 78.28bp 44.37bp 549.05bp},clip,width=0.92\textwidth,max height=0.38\textheight]{Ishii.pdf}
\end{sourceblock}
其中 km 是单位时间内每个界面面积的质量损失率，从 k th -
我们将表面平均值定义为

\begin{sourceblock}
\includegraphics[page=157,pagebox=mediabox,trim={45.00bp 511.28bp 44.37bp 98.54bp},clip,width=0.92\textwidth,max height=0.38\textheight]{Ishii.pdf}
\end{sourceblock}
此后，我们仅对可能与整体流体性质相混淆的变量使用下标( i )，对于仅出现在界面处的变量，例如σ和21 H 等，我们省略下标( i )。式(8-2)定义的平均值对应于式(4-23)的相位平均值，因此我们使用与下标i相同的符号。

类似地，也可以与方程(4-28)的相质量加权平均值类似地定义表面平均值。然而，在界面处，这些广泛特征的变量总是出现在通量项中，因此定义由传质速率 k m 加权的平均值比由密度加权更方便。因此我们有

\begin{sourceblock}
\includegraphics[page=157,pagebox=mediabox,trim={45.00bp 289.61bp 44.37bp 288.64bp},clip,width=0.92\textwidth,max height=0.38\textheight]{Ishii.pdf}
\end{sourceblock}
利用式（8-2）的定义，单位表面积的平均传质变为

\begin{sourceblock}
\includegraphics[page=157,pagebox=mediabox,trim={45.00bp 190.53bp 44.37bp 420.55bp},clip,width=0.92\textwidth,max height=0.38\textheight]{Ishii.pdf}
\end{sourceblock}
由式（8-1）和式（8-4）可知，界面传质条件可改写为

\begin{sourceblock}
\includegraphics[page=157,pagebox=mediabox,trim={45.00bp 104.19bp 44.37bp 524.98bp},clip,width=0.92\textwidth,max height=0.38\textheight]{Ishii.pdf}
\end{sourceblock}
\section*{8.2 界面动量传递}\phantomsection\addcontentsline{toc}{section}{8.2 界面动量传递}
我们回想一下，宏观界面动量传递项 k M 已在第 5 章 1.2 节中得到。因此，根据方程（2-9）、（2-70）、（5-23）和（5-57），我们有

\begin{sourceblock}
\includegraphics[page=158,pagebox=mediabox,trim={45.00bp 497.86bp 44.37bp 131.41bp},clip,width=0.92\textwidth,max height=0.38\textheight]{Ishii.pdf}
\end{sourceblock}
其中括号内的项是第 k 相单位面积的界面动量损失率。由于 k M 代表净界面动量增益，因此它由表面积浓度 ij a 加权。类似地，来自界面的混合动量源由下式给出

\begin{sourceblock}
\includegraphics[page=158,pagebox=mediabox,trim={45.00bp 388.90bp 44.37bp 240.37bp},clip,width=0.92\textwidth,max height=0.38\textheight]{Ishii.pdf}
\end{sourceblock}
其中上式右侧的两项代表平均曲率和表面张力梯度的影响。

在我们研究界面动量传递方程的矢量形式之前，让我们检查一下动量跳跃条件的法向分量，方程（2-91）。这是因为原始跳转条件方程（2-72）包含两条不同的信息；一个在法线方向，另一个在切线方向。为了在界面输运方程中保留这一特性，我们应该特别注意。通过将法向跳跃条件 Eq.(2-91) 用单位法向向量 1 n 进行点化并取时间平均，我们得到

\begin{sourceblock}
\includegraphics[page=158,pagebox=mediabox,trim={45.00bp 159.07bp 44.37bp 439.67bp},clip,width=0.92\textwidth,max height=0.38\textheight]{Ishii.pdf}
\end{sourceblock}
现在我们用方程（8-2）定义的表面平均值来表达这个方程。为了简化结果，我们假设

\begin{sourceblock}
\includegraphics[page=158,pagebox=mediabox,trim={45.00bp 83.32bp 44.37bp 556.72bp},clip,width=0.92\textwidth,max height=0.38\textheight]{Ishii.pdf}
\end{sourceblock}
因此传质速率 k m 、表面张力 σ 和密度

\begin{sourceblock}
\includegraphics[page=159,pagebox=mediabox,trim={45.00bp 527.56bp 44.37bp 71.75bp},clip,width=0.92\textwidth,max height=0.38\textheight]{Ishii.pdf}
\end{sourceblock}
我们在这里注意到，在类似的假设下，我们应该能够从矢量界面输运方程中恢复方程（8-10）。通过使用表面平均值，第 k 相界面动量增益 k M 变为

\begin{sourceblock}
\includegraphics[page=159,pagebox=mediabox,trim={45.00bp 439.90bp 44.37bp 201.08bp},clip,width=0.92\textwidth,max height=0.38\textheight]{Ishii.pdf}
\end{sourceblock}
where

\begin{sourceblock}
\includegraphics[page=159,pagebox=mediabox,trim={45.00bp 317.80bp 44.37bp 257.66bp},clip,width=0.92\textwidth,max height=0.38\textheight]{Ishii.pdf}
\end{sourceblock}
这里需要注意的是，界面处的剪切可以分解为法向分量和切向分量，因此 k k nk tk ⋅ = + n τ τ T。然而，法向应力小得可以忽略不计，因此可以假设
k k tk ⋅ = n τ T。式(8-11)右边的前三项原本是法向分量，而后两项实质上是切向分量。这里由于方程（4-62）而出现浓度梯度。

\begin{sourceblock}
\includegraphics[page=159,pagebox=mediabox,trim={45.00bp 126.04bp 44.37bp 456.64bp},clip,width=0.92\textwidth,max height=0.38\textheight]{Ishii.pdf}
\end{sourceblock}
第二项考虑了平均曲率变化的影响。然而，在微观尺度上σ的梯度被认为很小并且其矢量方向是相当随机的，因此最后一项可以被忽略。因此我们近似

\begin{sourceblock}
\includegraphics[page=160,pagebox=mediabox,trim={45.00bp 582.16bp 44.37bp 59.08bp},clip,width=0.92\textwidth,max height=0.38\textheight]{Ishii.pdf}
\end{sourceblock}
其中 H m M 是平均曲率变化对混合动量源的影响。

由式（8-11）和式（8-14）容易看出，界面动量传递条件的法向分量式（5-27）可由下式给出：

\begin{sourceblock}
\includegraphics[page=160,pagebox=mediabox,trim={45.00bp 445.73bp 44.37bp 170.84bp},clip,width=0.92\textwidth,max height=0.38\textheight]{Ishii.pdf}
\end{sourceblock}
因此，通过将法向分量方程（8-10）的标量形式与向量形式方程（8-15）进行比较，我们得到

\begin{sourceblock}
\includegraphics[page=160,pagebox=mediabox,trim={45.00bp 361.47bp 44.37bp 267.73bp},clip,width=0.92\textwidth,max height=0.38\textheight]{Ishii.pdf}
\end{sourceblock}
这里n k M 表示由于界面处压力不平衡而产生的形状阻力和升力。 t k M 表示由于剪切力不平衡而产生的表皮阻力。因此，剪切分量应满足

\begin{sourceblock}
\includegraphics[page=160,pagebox=mediabox,trim={45.00bp 262.47bp 44.37bp 366.73bp},clip,width=0.92\textwidth,max height=0.38\textheight]{Ishii.pdf}
\end{sourceblock}
由式（8-17）可知，各相表皮阻力之间以及界面剪切力之间存在作用-反作用关系。

为简单起见，我们将这两个阻力结合起来，并通过以下方式定义总广义阻力 ik M

\begin{sourceblock}
\includegraphics[page=160,pagebox=mediabox,trim={45.00bp 154.13bp 44.37bp 471.51bp},clip,width=0.92\textwidth,max height=0.38\textheight]{Ishii.pdf}
\end{sourceblock}
where

\begin{sourceblock}
\includegraphics[page=160,pagebox=mediabox,trim={45.00bp 54.59bp 44.37bp 541.23bp},clip,width=0.92\textwidth,max height=0.38\textheight]{Ishii.pdf}
\end{sourceblock}
\begin{sourceblock}
\includegraphics[page=161,pagebox=mediabox,trim={45.00bp 543.27bp 44.37bp 85.93bp},clip,width=0.92\textwidth,max height=0.38\textheight]{Ishii.pdf}
\end{sourceblock}
此外，根据式（4-61）和（4-62）给出的关系对传质速率 k m 进行简单分析，我们可以得出

\begin{sourceblock}
\includegraphics[page=161,pagebox=mediabox,trim={45.00bp 458.56bp 44.37bp 165.76bp},clip,width=0.92\textwidth,max height=0.38\textheight]{Ishii.pdf}
\end{sourceblock}
这使我们能够用单个参数 l i v 替换 m 1 i v 和 m 2 i v。作为界面动量传递条件的总结，我们有以下关系式

\begin{sourceblock}
\includegraphics[page=161,pagebox=mediabox,trim={45.00bp 371.93bp 44.37bp 269.06bp},clip,width=0.92\textwidth,max height=0.38\textheight]{Ishii.pdf}
\end{sourceblock}
其中 ik M 包括形状阻力、升力和蒙皮阻力的影响。

\begin{sourceblock}
\includegraphics[page=161,pagebox=mediabox,trim={45.00bp 312.94bp 44.37bp 328.30bp},clip,width=0.92\textwidth,max height=0.38\textheight]{Ishii.pdf}
\end{sourceblock}
\begin{sourceblock}
\includegraphics[page=161,pagebox=mediabox,trim={45.00bp 249.27bp 44.37bp 379.93bp},clip,width=0.92\textwidth,max height=0.38\textheight]{Ishii.pdf}
\end{sourceblock}
with

\begin{sourceblock}
\includegraphics[page=161,pagebox=mediabox,trim={45.00bp 173.57bp 44.37bp 448.06bp},clip,width=0.92\textwidth,max height=0.38\textheight]{Ishii.pdf}
\end{sourceblock}
\begin{sourceblock}
\includegraphics[page=161,pagebox=mediabox,trim={45.00bp 112.23bp 44.37bp 516.97bp},clip,width=0.92\textwidth,max height=0.38\textheight]{Ishii.pdf}
\end{sourceblock}
如果我们假设 ki p 、 σ 、 H m M 、 i a 、 k Γ 和 ki ρ 已知，则三个
应为 21 H 、 l i v 和 1 i M 指定本构关系，以确定界面几何形状、运动和广义阻力。此外，我们

\section*{8.3 界面能量传递}\phantomsection\addcontentsline{toc}{section}{8.3 界面能量传递}
第 k 相的宏观界面总能量传递由 k E 表示，它仅在相能量方程平均后出现，而混合界面能量源项为 m E。这三项应满足界面能量传递条件，即界面处的平衡方程。由于式(5-28)和式(5-29)给出的k E 和m E 的关系是由局部瞬时变量表示的，因此它们不能以其原始形式用于宏观公式中。现在我们用宏观变量来转换这些关系，作为建立界面本构关系的第一步。

由于其实际重要性，我们从分析界面热能传递项 k Λ 开始，然后继续研究
k E。由式(5-39)、式(5-19)、(5-23)、(5-28)的定义可得

\begin{sourceblock}
\includegraphics[page=162,pagebox=mediabox,trim={45.00bp 260.58bp 44.37bp 305.96bp},clip,width=0.92\textwidth,max height=0.38\textheight]{Ishii.pdf}
\end{sourceblock}
我们类比式（5-31）定义界面处的虚拟内能，因此

\begin{sourceblock}
\includegraphics[page=162,pagebox=mediabox,trim={45.00bp 137.86bp 44.37bp 452.03bp},clip,width=0.92\textwidth,max height=0.38\textheight]{Ishii.pdf}
\end{sourceblock}
每单位界面面积的热量输入定义为

\begin{sourceblock}
\includegraphics[page=162,pagebox=mediabox,trim={45.00bp 60.09bp 44.37bp 560.41bp},clip,width=0.92\textwidth,max height=0.38\textheight]{Ishii.pdf}
\end{sourceblock}
\begin{sourceblock}
\includegraphics[page=163,pagebox=mediabox,trim={45.00bp 545.32bp 44.37bp 83.95bp},clip,width=0.92\textwidth,max height=0.38\textheight]{Ishii.pdf}
\end{sourceblock}
为了检查上式右侧的第二组，我们引入由下式定义的波动分量

\begin{sourceblock}
\includegraphics[page=163,pagebox=mediabox,trim={45.00bp 473.92bp 44.37bp 166.24bp},clip,width=0.92\textwidth,max height=0.38\textheight]{Ishii.pdf}
\end{sourceblock}
然后我们有

\begin{sourceblock}
\includegraphics[page=163,pagebox=mediabox,trim={45.00bp 343.95bp 44.37bp 221.72bp},clip,width=0.92\textwidth,max height=0.38\textheight]{Ishii.pdf}
\end{sourceblock}
既然我们有

\begin{sourceblock}
\includegraphics[page=163,pagebox=mediabox,trim={45.00bp 270.58bp 44.37bp 353.73bp},clip,width=0.92\textwidth,max height=0.38\textheight]{Ishii.pdf}
\end{sourceblock}
式(8-31)右边第一项变为

\begin{sourceblock}
\includegraphics[page=163,pagebox=mediabox,trim={45.00bp 194.56bp 44.37bp 427.03bp},clip,width=0.92\textwidth,max height=0.38\textheight]{Ishii.pdf}
\end{sourceblock}
第二项可以重新排列为以下形式

\begin{sourceblock}
\includegraphics[page=163,pagebox=mediabox,trim={45.00bp 61.24bp 44.37bp 506.51bp},clip,width=0.92\textwidth,max height=0.38\textheight]{Ishii.pdf}
\end{sourceblock}
\begin{sourceblock}
\includegraphics[page=164,pagebox=mediabox,trim={45.00bp 433.30bp 44.37bp 52.00bp},clip,width=0.92\textwidth,max height=0.38\textheight]{Ishii.pdf}
\end{sourceblock}
因此，它代表总阻力的波动分量。因此，我们将阻力 T ki W 引起的湍流功定义为

\begin{sourceblock}
\includegraphics[page=164,pagebox=mediabox,trim={45.00bp 348.22bp 44.37bp 279.59bp},clip,width=0.92\textwidth,max height=0.38\textheight]{Ishii.pdf}
\end{sourceblock}
将式（8-33）和式（8-34）用式（8-36）代入式（8-31）可得

\begin{sourceblock}
\includegraphics[page=164,pagebox=mediabox,trim={45.00bp 247.90bp 44.37bp 350.34bp},clip,width=0.92\textwidth,max height=0.38\textheight]{Ishii.pdf}
\end{sourceblock}
由式（8-29）和式（8-37）可知，宏观界面热能传递 k Λ 为

\begin{sourceblock}
\includegraphics[page=164,pagebox=mediabox,trim={45.00bp 133.90bp 44.37bp 464.34bp},clip,width=0.92\textwidth,max height=0.38\textheight]{Ishii.pdf}
\end{sourceblock}
现在我们类比式(8-27)和式(5-37)引入界面处第k相的虚焓，从而

\begin{sourceblock}
\includegraphics[page=164,pagebox=mediabox,trim={45.00bp 50.14bp 44.37bp 581.98bp},clip,width=0.92\textwidth,max height=0.38\textheight]{Ishii.pdf}
\end{sourceblock}
\begin{sourceblock}
\includegraphics[page=165,pagebox=mediabox,trim={45.00bp 524.62bp 44.37bp 80.21bp},clip,width=0.92\textwidth,max height=0.38\textheight]{Ishii.pdf}
\end{sourceblock}
从 k Λ 、k M 和 k Γ 的关系式可以直接得到 k E，因此从式（8-21）和式（8-40）可以得到以下结果

\begin{sourceblock}
\includegraphics[page=165,pagebox=mediabox,trim={45.00bp 411.34bp 44.37bp 186.35bp},clip,width=0.92\textwidth,max height=0.38\textheight]{Ishii.pdf}
\end{sourceblock}
k Λ 和 k E 的表达式给出了第 k 相界面热能通量以及以界面平均值表示的总能量。

现在我们继续分析混合能源项 m E。假设

\begin{sourceblock}
\includegraphics[page=165,pagebox=mediabox,trim={45.00bp 288.20bp 44.37bp 336.76bp},clip,width=0.92\textwidth,max height=0.38\textheight]{Ishii.pdf}
\end{sourceblock}
式(5-29)可近似为

\begin{sourceblock}
\includegraphics[page=165,pagebox=mediabox,trim={45.00bp 215.22bp 44.37bp 406.68bp},clip,width=0.92\textwidth,max height=0.38\textheight]{Ishii.pdf}
\end{sourceblock}
我们在这里回顾一下，界面速度的表面散度是表面积膨胀（Aris，1962）。因此，我们有

\begin{sourceblock}
\includegraphics[page=165,pagebox=mediabox,trim={45.00bp 126.91bp 44.37bp 495.77bp},clip,width=0.92\textwidth,max height=0.38\textheight]{Ishii.pdf}
\end{sourceblock}
因此，结合表面张力梯度较小的假设，我们可以将式（8-43）近似为

\begin{sourceblock}
\includegraphics[page=166,pagebox=mediabox,trim={45.00bp 566.50bp 44.37bp 55.39bp},clip,width=0.92\textwidth,max height=0.38\textheight]{Ishii.pdf}
\end{sourceblock}
\begin{sourceblock}
\includegraphics[page=166,pagebox=mediabox,trim={45.00bp 495.38bp 44.37bp 107.68bp},clip,width=0.92\textwidth,max height=0.38\textheight]{Ishii.pdf}
\end{sourceblock}
我们注意到，方程（8-45）右侧的第一项考虑了与面积变化相关的表面能变化的影响，而第二项和最后一项代表表面张力所做的平均功。最后一项表示平均曲率变化对混合能量源的影响。结合式（8-21）和（8-26）我们有

\begin{sourceblock}
\includegraphics[page=166,pagebox=mediabox,trim={45.00bp 321.33bp 44.37bp 264.34bp},clip,width=0.92\textwidth,max height=0.38\textheight]{Ishii.pdf}
\end{sourceblock}
作为界面能量传递的总结，我们有以下关系。

总能量传递条件

\begin{sourceblock}
\includegraphics[page=166,pagebox=mediabox,trim={45.00bp 155.94bp 44.37bp 401.56bp},clip,width=0.92\textwidth,max height=0.38\textheight]{Ishii.pdf}
\end{sourceblock}
热能传递条件

\begin{sourceblock}
\includegraphics[page=166,pagebox=mediabox,trim={45.00bp 63.82bp 44.37bp 541.01bp},clip,width=0.92\textwidth,max height=0.38\textheight]{Ishii.pdf}
\end{sourceblock}
\begin{sourceblock}
\includegraphics[page=167,pagebox=mediabox,trim={45.00bp 488.82bp 44.37bp 55.55bp},clip,width=0.92\textwidth,max height=0.38\textheight]{Ishii.pdf}
\end{sourceblock}
由于这些关系现在用本体流体和界面的平均值来表达，因此它们可以被认为具有宏观形式。本构方程可以通过将界面变量与整体流体平均值和其他特征参数（例如 i a ）相关联来获得。鉴于方程（8-5）、（8-21）和（8-48），我们认识到双流体模型和漂移通量（混合）模型的必要界面本构关系之间存在相当大的差异。对于前一种模型，需要通过本构方程指定 1 Г 、 1 M 、 1 E 、 m M 和 m E，而对于后一种模型，只需提供 1 Г，

\begin{sourceblock}
\includegraphics[page=167,pagebox=mediabox,trim={45.00bp 317.01bp 44.37bp 315.11bp},clip,width=0.92\textwidth,max height=0.38\textheight]{Ishii.pdf}
\end{sourceblock}
比二流体模型简单。在扩散或漂移通量模型中，我们提供每一相速度之间的关系，因此只需要一个动量方程。然而，在双流体模型中，我们指定动量交换项 k M，然后同时求解两个动量方程。我们还注意到，在不做假设的情况下，两相的 k Λ 之和不会简化为 m E 这样的简单形式，因此预计在漂移通量模型中使用热能方程时应特别注意。

\chapter{双流体模型}
双流体模型（Ishii，1975，Ishii 和 Mishima，1984）是通过分别考虑每个相来制定的。因此，该模型用两组守恒方程来表示，这些方程控制着每个相中的质量、动量和能量的平衡。然而，由于一相的平均场并不独立于另一相，因此这些平衡方程中出现了相互作用项。由 k Γ 、 k M 和 表示的项
k E 是从界面转移到第 k 相的质量、动量和能量。由于这些量也应遵守界面处的平衡定律，因此我们从局部跳跃条件导出了界面输运条件。因此，具有三个界面输运条件的六个微分场方程控制着宏观两相流系统。

在双流体模型公式中，各相的传递过程由各自的平衡方程表示，因此预计该模型可以比漂移通量（或混合）模型预测更详细的变化和相相互作用。然而，这意味着二流体模型不仅在涉及的场方程数量方面而且在必要的本构方程方面都要复杂得多。显然，这些本构方程应该准确地显示模型的有用性。对于交互项尤其如此
k Γ 、k M 和 k E，因为在场方程中没有这些界面交换，两相本质上是独立的。这些交互项决定了阶段之间的耦合程度，因此每个阶段的传递过程都受到这些项的很大影响。

双流体模型的真正重要性在于它可以考虑相之间的动态和非平衡相互作用。这是通过使用每个相的动量方程和两个独立的速度场以及公式中的两个能量方程来完成的。因此，预计双流体模型可用于
然而，如果两相强烈耦合（其中各相的响应是同时的，使得两相接近机械和热平衡，或者波传播牢固地互锁），则二流体模型会给系统带来实际应用中不必要的复杂性。此外，可以说二流体模型非常适合局部波传播和相关稳定性问题的研究。然而，如果关注系统中两相混合物的总响应而不是每一相的局部行为，则漂移通量模型更简单，并且在大多数情况下对于解决问题是有效的。对于一般的三维流动，二流体模型比混合模型更好，因为相对速度相关性很难以一般的三维形式发展。

接下来，我们研究二流体模型的一般公式以及封闭方程组所需的各种本构方程。但需要注意的是，通过使未知数和方程的数量相同来封闭微分方程组并不意味着解的存在，也不保证其唯一性。然而，这是正确设置代表要分析的物理系统的数学模型的必要条件。

还应该记住，两相流系统的数学模型在任何意义上都没有牢固建立，需要一些额外的研究来完成一般两相流的三维模型。为了理解我们面临的困难，我们回顾一下，即使在单相湍流中，湍流通量的一般本构方程也尚未建立。鉴于三维二流体模型制定的现有技术，首先一般性地讨论模型的必要形式。因此，本章应被视为制定本构方程的框架和指导。下图总结了二流体模型公式的建立。

\section*{9.1 二流体模型场方程}\phantomsection\addcontentsline{toc}{section}{9.1 二流体模型场方程}
二流体模型连续性方程 二流体模型的特征是两个独立的速度场，它们指定每个相的运动。速度场最自然的选择显然是质量加权平均相速度 l k v。

\begin{sourceblock}
\includegraphics[page=170,pagebox=mediabox,trim={45.00bp 292.81bp 45.56bp 54.37bp},clip,width=\textwidth,max height=0.72\textheight]{Ishii.pdf}
\end{sourceblock}
因此，模型中使用的连续性方程的合适形式应为方程（5-21）

\begin{sourceblock}
\includegraphics[page=170,pagebox=mediabox,trim={45.00bp 195.40bp 44.37bp 429.59bp},clip,width=0.92\textwidth,max height=0.38\textheight]{Ishii.pdf}
\end{sourceblock}
界面传质条件见式(5-22)

\begin{sourceblock}
\includegraphics[page=170,pagebox=mediabox,trim={45.00bp 124.29bp 44.37bp 504.88bp},clip,width=0.92\textwidth,max height=0.38\textheight]{Ishii.pdf}
\end{sourceblock}
因此，由k Γ 表示的项表示从每单位体积的界面处的相变产生第k相质量的速率。出现这种情况是因为局部连续性方程已及时积分以获得宏观场方程。此外，可以说式（9-1）和（9-2）是质量守恒定律的一般表述
根据式（7-11）的对流导数，连续性方程变为

\begin{sourceblock}
\includegraphics[page=171,pagebox=mediabox,trim={45.00bp 506.25bp 44.37bp 121.88bp},clip,width=0.92\textwidth,max height=0.38\textheight]{Ishii.pdf}
\end{sourceblock}
因此我们有

\begin{sourceblock}
\includegraphics[page=171,pagebox=mediabox,trim={45.00bp 433.42bp 44.37bp 189.61bp},clip,width=0.92\textwidth,max height=0.38\textheight]{Ishii.pdf}
\end{sourceblock}
对于稳态流，方程（9-1）的时间导数下降，因此我们有

\begin{sourceblock}
\includegraphics[page=171,pagebox=mediabox,trim={45.00bp 375.40bp 44.37bp 265.21bp},clip,width=0.92\textwidth,max height=0.38\textheight]{Ishii.pdf}
\end{sourceblock}
如果每个相最初都是不可压缩的，则平均密度 ρ k 是恒定的。因此我们有

\begin{sourceblock}
\includegraphics[page=171,pagebox=mediabox,trim={45.00bp 284.50bp 44.37bp 336.62bp},clip,width=0.92\textwidth,max height=0.38\textheight]{Ishii.pdf}
\end{sourceblock}
此外，如果没有相位变化，连续性方程简化为

\begin{sourceblock}
\includegraphics[page=171,pagebox=mediabox,trim={45.00bp 196.24bp 44.37bp 425.62bp},clip,width=0.92\textwidth,max height=0.38\textheight]{Ishii.pdf}
\end{sourceblock}
可用于无相变的低速两相流。在这些条件下，两相系统的运动学完全由相重新分布（即对流和扩散）控制。上述方程的形式类似于单相可压缩流的形式。

相位连续性方程的一般形式由式（9-1）以矢量表示形式给出。考虑到实际的重要性，我们将方程表达在直角坐标系和柱坐标系中。因此，在直角坐标 ( ) , , x y z 中我们有

\begin{sourceblock}
\includegraphics[page=172,pagebox=mediabox,trim={45.00bp 536.11bp 44.37bp 55.00bp},clip,width=0.92\textwidth,max height=0.38\textheight]{Ishii.pdf}
\end{sourceblock}
如果流动被限制在二维，那么它代表平面流动。在这种情况下，关于x的偏导数可以从方程(9-8)中去掉。我们还注意到，对于没有相变的稳定平面流，可以引入流函数。

柱坐标系中的连续性方程变为

\begin{sourceblock}
\includegraphics[page=172,pagebox=mediabox,trim={45.00bp 383.24bp 44.37bp 210.24bp},clip,width=0.92\textwidth,max height=0.38\textheight]{Ishii.pdf}
\end{sourceblock}
如果不依赖于 θ 方向，则称流动是轴对称的，因此对于此类流动，我们有

\begin{sourceblock}
\includegraphics[page=172,pagebox=mediabox,trim={45.00bp 300.39bp 44.37bp 324.57bp},clip,width=0.92\textwidth,max height=0.38\textheight]{Ishii.pdf}
\end{sourceblock}
还可以针对没有相变的稳定轴对称流引入流函数，从而可以从公式中消除连续性方程。

二流体模型动量方程在二流体模型公式中，动量守恒由两个具有界面动量传递条件的动量方程来表达。正如前面提到的，适当的场方程应该用质心或各相的质心速度 l k v 来表示，因此从方程（5-26）我们有两个动量方程：

\begin{sourceblock}
\includegraphics[page=172,pagebox=mediabox,trim={45.00bp 65.40bp 44.37bp 525.62bp},clip,width=0.92\textwidth,max height=0.38\textheight]{Ishii.pdf}
\end{sourceblock}
\begin{sourceblock}
\includegraphics[page=173,pagebox=mediabox,trim={45.00bp 543.27bp 44.37bp 85.93bp},clip,width=0.92\textwidth,max height=0.38\textheight]{Ishii.pdf}
\end{sourceblock}
with

\begin{sourceblock}
\includegraphics[page=173,pagebox=mediabox,trim={45.00bp 483.94bp 44.37bp 157.30bp},clip,width=0.92\textwidth,max height=0.38\textheight]{Ishii.pdf}
\end{sourceblock}
我们在这里注意到，每个相的动量方程都有一个界面源项 k M，它通过方程（9-12）耦合两个相的运动。

这里， 21 H 和 σ 是界面的平均曲率和表面张力，而 m M 给出的术语考虑了平均曲率变化的影响。

鉴于第8章1.2节，式(9-11)可改写为

\begin{sourceblock}
\includegraphics[page=173,pagebox=mediabox,trim={45.00bp 281.23bp 44.37bp 287.86bp},clip,width=0.92\textwidth,max height=0.38\textheight]{Ishii.pdf}
\end{sourceblock}
因此，通过使用方程（7-11）的对流导数，第k相运动方程变为

\begin{sourceblock}
\includegraphics[page=173,pagebox=mediabox,trim={45.00bp 173.27bp 44.37bp 430.90bp},clip,width=0.92\textwidth,max height=0.38\textheight]{Ishii.pdf}
\end{sourceblock}
或者将式(8-20)代入式(9-15)得

\begin{sourceblock}
\includegraphics[page=173,pagebox=mediabox,trim={45.00bp 50.14bp 44.37bp 525.89bp},clip,width=0.92\textwidth,max height=0.38\textheight]{Ishii.pdf}
\end{sourceblock}
直角坐标系中的运动方程 ( ) , , x y z 可以给出如下
x -component

\begin{sourceblock}
\includegraphics[page=174,pagebox=mediabox,trim={45.00bp 349.82bp 44.37bp 158.71bp},clip,width=0.92\textwidth,max height=0.38\textheight]{Ishii.pdf}
\end{sourceblock}
y -component

\begin{sourceblock}
\includegraphics[page=174,pagebox=mediabox,trim={45.00bp 159.63bp 44.37bp 348.20bp},clip,width=0.92\textwidth,max height=0.38\textheight]{Ishii.pdf}
\end{sourceblock}
z -component

\begin{sourceblock}
\includegraphics[page=174,pagebox=mediabox,trim={45.00bp 80.20bp 44.37bp 538.74bp},clip,width=0.92\textwidth,max height=0.38\textheight]{Ishii.pdf}
\end{sourceblock}
\begin{sourceblock}
\includegraphics[page=175,pagebox=mediabox,trim={45.00bp 494.09bp 44.37bp 55.67bp},clip,width=0.92\textwidth,max height=0.38\textheight]{Ishii.pdf}
\end{sourceblock}
我们在这里注意到，对于平面流，运动方程的 x 分量下降。此外，所有关于x的偏导数都应该从运动方程的y和z分量中消除，即方程（9-18）和（9-19）。

柱坐标系中的运动方程 ( ) , , r z θ 变为 r 分量

\begin{sourceblock}
\includegraphics[page=175,pagebox=mediabox,trim={45.00bp 191.07bp 44.37bp 281.84bp},clip,width=0.92\textwidth,max height=0.38\textheight]{Ishii.pdf}
\end{sourceblock}
θ -component

\begin{sourceblock}
\includegraphics[page=175,pagebox=mediabox,trim={45.00bp 73.88bp 44.37bp 505.73bp},clip,width=0.92\textwidth,max height=0.38\textheight]{Ishii.pdf}
\end{sourceblock}
\begin{sourceblock}
\includegraphics[page=176,pagebox=mediabox,trim={45.00bp 501.03bp 44.37bp 55.37bp},clip,width=0.92\textwidth,max height=0.38\textheight]{Ishii.pdf}
\end{sourceblock}
z -component

\begin{sourceblock}
\includegraphics[page=176,pagebox=mediabox,trim={45.00bp 275.58bp 44.37bp 195.25bp},clip,width=0.92\textwidth,max height=0.38\textheight]{Ishii.pdf}
\end{sourceblock}
对于轴对称流，具有偏导数的项
θ 从方程中消失。此外，如果流动不存在绕 z 轴的循环运动，则 θ 方向上的速度也为零，从而可以消除方程的整个 θ 分量。对于两相管流的许多实际问题，这是一个足够精确的模型。

二流体模型能量方程 能量守恒的最基本形式是通过考虑总能量的平衡来表达的。对于双流体模型公式，我们有两个具有界面能量传递条件的总能量方程。因此从式（5-32）我们有

\begin{sourceblock}
\includegraphics[page=177,pagebox=mediabox,trim={45.00bp 516.40bp 44.37bp 56.05bp},clip,width=0.92\textwidth,max height=0.38\textheight]{Ishii.pdf}
\end{sourceblock}
这里我们假设 k g 是常数，即 l k k = g g。界面输运条件（8-48）耦合了两相的能量传输过程，因此我们有

\begin{sourceblock}
\includegraphics[page=177,pagebox=mediabox,trim={45.00bp 417.51bp 44.37bp 211.69bp},clip,width=0.92\textwidth,max height=0.38\textheight]{Ishii.pdf}
\end{sourceblock}
where

\begin{sourceblock}
\includegraphics[page=177,pagebox=mediabox,trim={45.00bp 345.49bp 44.37bp 279.26bp},clip,width=0.92\textwidth,max height=0.38\textheight]{Ishii.pdf}
\end{sourceblock}
这些关系表明界面能量传递项的总和

每个相的 k E 与表面能的时间变化率和表面张力所做的功保持平衡。我们在这里注意到，H m E 给出的项考虑了平均曲率变化的影响。如果我们根据方程（8-48）给出的界面变量使用 k E 的详细表达式，则总能量方程可以重写为

\begin{sourceblock}
\includegraphics[page=177,pagebox=mediabox,trim={45.00bp 92.81bp 44.37bp 415.73bp},clip,width=0.92\textwidth,max height=0.38\textheight]{Ishii.pdf}
\end{sourceblock}
通过对对流导数进行变换，方程（7-13）可以写成

\begin{sourceblock}
\includegraphics[page=178,pagebox=mediabox,trim={45.00bp 498.42bp 44.37bp 56.03bp},clip,width=0.92\textwidth,max height=0.38\textheight]{Ishii.pdf}
\end{sourceblock}
方程(9-23)描述了从观察者在固定点处看到的能量传递，方程(9-27)表示跟随具有重心速度l k v 的流体的能量传递。

在许多实际的传热问题中，使用热能方程代替总能量方程是很方便的。对于具有附加热量的低速两相流尤其如此，其中机械项与高传热速率相比是微不足道的。因此，通过回顾方程（5-38）和方程（5-39），热能方程由下式给出

\begin{sourceblock}
\includegraphics[page=178,pagebox=mediabox,trim={45.00bp 285.59bp 44.37bp 303.62bp},clip,width=0.92\textwidth,max height=0.38\textheight]{Ishii.pdf}
\end{sourceblock}
将式(8-40)中kΛ的表达式代入上式，得

\begin{sourceblock}
\includegraphics[page=178,pagebox=mediabox,trim={45.00bp 115.90bp 44.37bp 423.40bp},clip,width=0.92\textwidth,max height=0.38\textheight]{Ishii.pdf}
\end{sourceblock}
该方程也可以根据对流导数进行变换为

\begin{sourceblock}
\includegraphics[page=179,pagebox=mediabox,trim={45.00bp 480.12bp 44.37bp 55.28bp},clip,width=0.92\textwidth,max height=0.38\textheight]{Ishii.pdf}
\end{sourceblock}
这是描述从随着第 k 相的质心移动的观察者处看到的热能交换的方程。

\begin{sourceblock}
\includegraphics[page=179,pagebox=mediabox,trim={45.00bp 381.85bp 44.37bp 232.37bp},clip,width=0.92\textwidth,max height=0.38\textheight]{Ishii.pdf}
\end{sourceblock}
\begin{sourceblock}
\includegraphics[page=179,pagebox=mediabox,trim={45.00bp 333.10bp 44.37bp 307.88bp},clip,width=0.92\textwidth,max height=0.38\textheight]{Ishii.pdf}
\end{sourceblock}
那么式(9-30)简化为

\begin{sourceblock}
\includegraphics[page=179,pagebox=mediabox,trim={45.00bp 204.84bp 44.37bp 363.80bp},clip,width=0.92\textwidth,max height=0.38\textheight]{Ishii.pdf}
\end{sourceblock}
现在我们将上述热能方程展开到两个具有实际意义的坐标系中。因此，在直角坐标系中

\begin{sourceblock}
\includegraphics[page=179,pagebox=mediabox,trim={45.00bp 150.23bp 44.37bp 495.53bp},clip,width=0.92\textwidth,max height=0.38\textheight]{Ishii.pdf}
\end{sourceblock}
\begin{sourceblock}
\includegraphics[page=179,pagebox=mediabox,trim={45.00bp 60.02bp 44.37bp 519.17bp},clip,width=0.92\textwidth,max height=0.38\textheight]{Ishii.pdf}
\end{sourceblock}
\begin{sourceblock}
\includegraphics[page=180,pagebox=mediabox,trim={45.00bp 341.52bp 44.37bp 55.62bp},clip,width=0.92\textwidth,max height=0.38\textheight]{Ishii.pdf}
\end{sourceblock}
对于平面流，关于 x 的偏导数下降并且速度的 x 分量为零。

在柱坐标 ( ) , , r z θ 中，热能方程变为

\begin{sourceblock}
\includegraphics[page=180,pagebox=mediabox,trim={45.00bp 58.82bp 44.37bp 395.48bp},clip,width=0.92\textwidth,max height=0.38\textheight]{Ishii.pdf}
\end{sourceblock}
\begin{sourceblock}
\includegraphics[page=181,pagebox=mediabox,trim={45.00bp 478.69bp 44.37bp 56.54bp},clip,width=0.92\textwidth,max height=0.38\textheight]{Ishii.pdf}
\end{sourceblock}
对于轴对称流，关于 θ 的偏导数从方程中下降。此外，如果流动不存在绕 z 轴的循环运动，则速度的 θ 分量为零。对于管道中的许多两相流，特别是垂直管道流来说，这是一个很好的近似。

可以看出，总能量方程和热能方程的完整形式都相当复杂，因此进行一些简化对于解决实际问题非常重要。我们研究以下几个特殊情况。如果传热和相变主导能量交换，那么我们可以忽略机械效应产生的项。在此条件下，式(9-30)可简化为

\begin{sourceblock}
\includegraphics[page=181,pagebox=mediabox,trim={45.00bp 276.15bp 44.37bp 347.36bp},clip,width=0.92\textwidth,max height=0.38\textheight]{Ishii.pdf}
\end{sourceblock}
除了可压缩波传播和/或高速流动条件下的问题外，上述方程足以满足许多两相流分析。

\begin{sourceblock}
\includegraphics[page=181,pagebox=mediabox,trim={45.00bp 85.72bp 44.37bp 436.93bp},clip,width=0.92\textwidth,max height=0.38\textheight]{Ishii.pdf}
\end{sourceblock}
如果我们使用柱坐标 ( ) , , r z θ 我们有

\begin{sourceblock}
\includegraphics[page=182,pagebox=mediabox,trim={45.00bp 489.80bp 44.37bp 55.12bp},clip,width=0.92\textwidth,max height=0.38\textheight]{Ishii.pdf}
\end{sourceblock}
此外，如果流动是轴对称的，并且轴向方向上的热传递可以忽略不计，则它会简化为以下形式

\begin{sourceblock}
\includegraphics[page=182,pagebox=mediabox,trim={45.00bp 360.17bp 44.37bp 219.33bp},clip,width=0.92\textwidth,max height=0.38\textheight]{Ishii.pdf}
\end{sourceblock}
它是方程（9-35）的简化形式，但重要的传热机制保留在上面的方程中。

\section*{9.2 二流体模型本构关系}\phantomsection\addcontentsline{toc}{section}{9.2 二流体模型本构关系}
\subsection*{9.2.1 熵不等式}\phantomsection\addcontentsline{toc}{subsection}{9.2.1 熵不等式}
本章开头讨论了构建二流体模型的一般方案。显然，宏观场方程（9-1）、（9-11）和（9-23）以及界面输运条件（9-2）、（9-12）和（9-24）不足以描述任何特定系统，因为变量的数量超过了可用方程的数量。需要指定特定材料组的材料和响应特性的附加信息。这些通常称为本构方程，如第 2 章中详细解释的那样。本节的目的是检查封闭方程组所需的本构方程。总是可以引入更详细的机制和变量来区分材料和转移机制的各种影响，然后使方程组复杂化。因此，我们将用最简单且合理通用的方程组来讨论本构关系最重要的方面，即决定论原理。为此我们考虑两个
类比第二章，我们继续研究宏观领域的熵不等式。因此，通过将式（5-8）和（5-10）的平均过程应用于不等式（2-23）和（2-85），我们得到

\begin{sourceblock}
\includegraphics[page=183,pagebox=mediabox,trim={45.00bp 421.81bp 44.37bp 145.30bp},clip,width=0.92\textwidth,max height=0.38\textheight]{Ishii.pdf}
\end{sourceblock}
and

\begin{sourceblock}
\includegraphics[page=183,pagebox=mediabox,trim={45.00bp 323.92bp 44.37bp 273.21bp},clip,width=0.92\textwidth,max height=0.38\textheight]{Ishii.pdf}
\end{sourceblock}
其中我们将内部身体加热 k q 设为零。我们记得在第 2 章中，为了获得界面处的简单边界条件，界面熵生成 a 被假设为零。我们在这里遵循完全相同的方法，因此我们有

\begin{sourceblock}
\includegraphics[page=183,pagebox=mediabox,trim={45.00bp 226.91bp 44.37bp 414.34bp},clip,width=0.92\textwidth,max height=0.38\textheight]{Ishii.pdf}
\end{sourceblock}
因此，我们得到

\begin{sourceblock}
\includegraphics[page=183,pagebox=mediabox,trim={45.00bp 95.23bp 44.37bp 473.14bp},clip,width=0.92\textwidth,max height=0.38\textheight]{Ishii.pdf}
\end{sourceblock}
然而，对于大多数实际问题，最后一个条件可以近似为

\begin{sourceblock}
\includegraphics[page=184,pagebox=mediabox,trim={45.00bp 527.44bp 44.37bp 55.48bp},clip,width=0.92\textwidth,max height=0.38\textheight]{Ishii.pdf}
\end{sourceblock}
即式(2-107)的宏观形式。

方程（9-43）的第一和第二条件可用于通过两个参数 i Τ 和 m ti v 来代替流体温度和流体切向速度，而最后一个条件在宏观公式中对于设置界面的能级仍然非常重要。由于在较低减压下，密度比较大，而在较高减压下，表面张力影响较小，式(9-44)可近似为

\begin{sourceblock}
\includegraphics[page=184,pagebox=mediabox,trim={45.00bp 372.37bp 44.37bp 264.83bp},clip,width=0.92\textwidth,max height=0.38\textheight]{Ishii.pdf}
\end{sourceblock}
因此，蒸气几乎总是非常接近界面处的饱和状态。方程（9-45）非常简单，甚至在两相流问题的局部即时公式中也得到广泛应用。从以上讨论可以看出，式（9-42）的结果可以用单个方程（9-45）来表示，因为式（9-43）的其他条件只需用界面温度和切向速度i Τ 和 m ti v 代替即可满足。现在我们研究式（9-40）给出的各相的熵不等式。如果界面温度的波动不重要，那么我们可以类推式（8-9），得到以下近似值

\begin{sourceblock}
\includegraphics[page=184,pagebox=mediabox,trim={45.00bp 194.14bp 44.37bp 447.14bp},clip,width=0.92\textwidth,max height=0.38\textheight]{Ishii.pdf}
\end{sourceblock}
则式(9-40)可以用第8章1.1节的界面宏观变量表示为

\begin{sourceblock}
\includegraphics[page=184,pagebox=mediabox,trim={45.00bp 53.27bp 44.37bp 515.39bp},clip,width=0.92\textwidth,max height=0.38\textheight]{Ishii.pdf}
\end{sourceblock}
虽然上述方程可以通过具有正粘度和导热系数的局部瞬时公式来满足，但它仍然对宏观本构方程施加了一些限制。换句话说，由局部瞬时公式得到的k Δ 的精确形式满足式(9-47)，但它并不能保证各种本构方程的等式左边总是为正。可以说，只有当一组受连续性、动量和能量方程约束的本构方程满足不等式（9-47）时，该公式才是一致的。在应用上述不等式之前，我们将讨论宏观模型最重要的特征之一，该特征仅在平均后出现。

\subsection*{9.2.2 状态方程}\phantomsection\addcontentsline{toc}{subsection}{9.2.2 状态方程}
可以说，即使原始的局部瞬时公式具有第2章给出的带有标准状态方程的简单线性本构方程，通过平均得到的宏观模型也可能不具有如此简单的本构方程。这是因为公式中出现了局部瞬时波动的统计效应。一般来说，这些统计效应不仅取决于宏观变量的当前状态，而且还取决于达到当前状态的过程。例如，具有相同能量 l k u 和密度 k ρ 值的流体粒子可以具有完全不同的平均温度或压力值。这些都表明宏观场具有记忆材料的特征(Truesdell 1969)。因此，本构方程一般是由过去过程的泛函给出的。这使得宏观本构方程的分析变得极其复杂和困难。显然，该公式将产生一组耦合积分微分方程。为了避免这些困难，我们必须做出一些假设，而牺牲模型的准确性。我们知道，所有材料都表现出记忆衰退的特性（Coleman and Noll，1960）。因此，记忆在公式中的作用的重要性取决于有效记忆的时间跨度与宏观过程的时间常数的比率。

现在让我们检查与方程（2-24）相对应的平均状态方程。我们有

\begin{sourceblock}
\includegraphics[page=185,pagebox=mediabox,trim={45.00bp 73.22bp 44.37bp 548.62bp},clip,width=0.92\textwidth,max height=0.38\textheight]{Ishii.pdf}
\end{sourceblock}
\begin{sourceblock}
\includegraphics[page=186,pagebox=mediabox,trim={45.00bp 560.83bp 44.37bp 55.84bp},clip,width=0.92\textwidth,max height=0.38\textheight]{Ishii.pdf}
\end{sourceblock}
通过对上式求平均值，我们得到

\begin{sourceblock}
\includegraphics[page=186,pagebox=mediabox,trim={45.00bp 325.27bp 44.37bp 134.79bp},clip,width=0.92\textwidth,max height=0.38\textheight]{Ishii.pdf}
\end{sourceblock}
这里我们使用了身份

\begin{sourceblock}
\includegraphics[page=186,pagebox=mediabox,trim={45.00bp 188.92bp 44.37bp 377.47bp},clip,width=0.92\textwidth,max height=0.38\textheight]{Ishii.pdf}
\end{sourceblock}
with

\begin{sourceblock}
\includegraphics[page=186,pagebox=mediabox,trim={45.00bp 127.12bp 44.37bp 513.04bp},clip,width=0.92\textwidth,max height=0.38\textheight]{Ishii.pdf}
\end{sourceblock}
方程（9-49）表明，一般来说，我们没有平均变量的简单状态方程。内能 l k u 、熵 l k s 和密度 k ρ 之间的关系受到界面输运和变量波动的统计效应的影响。

如果波动分量足够小于这些变量的宏观变化，则这些平均值之间会出现简单的静态关系，因此状态方程的线性展开是一个很好的近似。在这种情况下我们有

\begin{sourceblock}
\includegraphics[page=187,pagebox=mediabox,trim={45.00bp 514.64bp 44.37bp 123.31bp},clip,width=0.92\textwidth,max height=0.38\textheight]{Ishii.pdf}
\end{sourceblock}
with

\begin{sourceblock}
\includegraphics[page=187,pagebox=mediabox,trim={45.00bp 436.54bp 44.37bp 184.67bp},clip,width=0.92\textwidth,max height=0.38\textheight]{Ishii.pdf}
\end{sourceblock}
and

\begin{sourceblock}
\includegraphics[page=187,pagebox=mediabox,trim={45.00bp 352.99bp 44.37bp 263.58bp},clip,width=0.92\textwidth,max height=0.38\textheight]{Ishii.pdf}
\end{sourceblock}
这些关系对于两相流成立，每个相本身在 t Δ 时间间隔内处于接近平衡状态。此后我们假设宏观场中各相都服从静态状态方程Eq.(9-52)。这是一个相当重要且实用的假设，使我们能够与标准单相流公式并行地构建双流体模型及其本构方程。在上述条件下，类推式（2-24）、（2-25）和（2-26）有以下关系

\begin{sourceblock}
\includegraphics[page=187,pagebox=mediabox,trim={45.00bp 71.91bp 44.37bp 426.87bp},clip,width=0.92\textwidth,max height=0.38\textheight]{Ishii.pdf}
\end{sourceblock}
基本状态方程也可以用热量状态方程和热状态方程的组合来表示，因此

\begin{sourceblock}
\includegraphics[page=188,pagebox=mediabox,trim={45.00bp 540.18bp 44.37bp 96.97bp},clip,width=0.92\textwidth,max height=0.38\textheight]{Ishii.pdf}
\end{sourceblock}
\begin{sourceblock}
\includegraphics[page=188,pagebox=mediabox,trim={45.00bp 488.17bp 44.37bp 149.03bp},clip,width=0.92\textwidth,max height=0.38\textheight]{Ishii.pdf}
\end{sourceblock}
或者如果我们将焓作为变量，则变为

\begin{sourceblock}
\includegraphics[page=188,pagebox=mediabox,trim={45.00bp 423.18bp 44.37bp 213.15bp},clip,width=0.92\textwidth,max height=0.38\textheight]{Ishii.pdf}
\end{sourceblock}
\begin{sourceblock}
\includegraphics[page=188,pagebox=mediabox,trim={45.00bp 371.17bp 44.37bp 266.03bp},clip,width=0.92\textwidth,max height=0.38\textheight]{Ishii.pdf}
\end{sourceblock}
鉴于其巨大的实际重要性，我们现在研究几种热力学二阶导数。恒压比热

\begin{sourceblock}
\includegraphics[page=188,pagebox=mediabox,trim={45.00bp 255.10bp 44.37bp 331.85bp},clip,width=0.92\textwidth,max height=0.38\textheight]{Ishii.pdf}
\end{sourceblock}
and

\begin{sourceblock}
\includegraphics[page=188,pagebox=mediabox,trim={45.00bp 166.11bp 44.37bp 443.02bp},clip,width=0.92\textwidth,max height=0.38\textheight]{Ishii.pdf}
\end{sourceblock}
类似地，热膨胀系数 k β 和等温压缩系数
Tk κ 定义为

\begin{sourceblock}
\includegraphics[page=188,pagebox=mediabox,trim={45.00bp 64.17bp 44.37bp 544.96bp},clip,width=0.92\textwidth,max height=0.38\textheight]{Ishii.pdf}
\end{sourceblock}
\begin{sourceblock}
\includegraphics[page=189,pagebox=mediabox,trim={45.00bp 556.83bp 44.37bp 55.61bp},clip,width=0.92\textwidth,max height=0.38\textheight]{Ishii.pdf}
\end{sourceblock}
其中 Tk a 是等温声速。在这四个导数中，我们有以下恒等式

\begin{sourceblock}
\includegraphics[page=189,pagebox=mediabox,trim={45.00bp 465.13bp 44.37bp 157.33bp},clip,width=0.92\textwidth,max height=0.38\textheight]{Ishii.pdf}
\end{sourceblock}
已知，如果式（9-55）成立，则只有三个热力学二阶导数是独立的，其他的可以由这三个导数得到。我们来介绍一下比热容比

\begin{sourceblock}
\includegraphics[page=189,pagebox=mediabox,trim={45.00bp 368.86bp 44.37bp 257.43bp},clip,width=0.92\textwidth,max height=0.38\textheight]{Ishii.pdf}
\end{sourceblock}
\begin{sourceblock}
\includegraphics[page=189,pagebox=mediabox,trim={45.00bp 285.51bp 44.37bp 303.48bp},clip,width=0.92\textwidth,max height=0.38\textheight]{Ishii.pdf}
\end{sourceblock}
其中 Sk a 是等熵声速。然后我们有

\begin{sourceblock}
\includegraphics[page=189,pagebox=mediabox,trim={45.00bp 207.88bp 44.37bp 415.29bp},clip,width=0.92\textwidth,max height=0.38\textheight]{Ishii.pdf}
\end{sourceblock}
这表明等熵声速总是大于等温声速，因为从系统的稳定性来看，0 Tk κ ≥，因此1 k γ ≥。热力学二阶导数或热和热量状态方程的重要性与测量的可能性有关。例如，流体压力和温度相对容易测量，因此可以通过实验构造式（9-58）和（9-59）形式的状态方程。饱和条件 经典饱和条件由下式给出

\begin{sourceblock}
\includegraphics[page=190,pagebox=mediabox,trim={45.00bp 552.54bp 44.37bp 55.54bp},clip,width=0.92\textwidth,max height=0.38\textheight]{Ishii.pdf}
\end{sourceblock}
then

\begin{sourceblock}
\includegraphics[page=190,pagebox=mediabox,trim={45.00bp 491.84bp 44.37bp 146.11bp},clip,width=0.92\textwidth,max height=0.38\textheight]{Ishii.pdf}
\end{sourceblock}
因此我们有关系

\begin{sourceblock}
\includegraphics[page=190,pagebox=mediabox,trim={45.00bp 432.13bp 44.37bp 209.19bp},clip,width=0.92\textwidth,max height=0.38\textheight]{Ishii.pdf}
\end{sourceblock}
假设与式（2-99）相同。

\subsection*{9.2.3 决定论}\phantomsection\addcontentsline{toc}{subsection}{9.2.3 决定论}
在本分析中，我们假设存在静态状态方程Eq.(9-55)。从决定论的原理来看，我们应该能够从过去的历史来预测现在的状态。必要条件是方程组是闭的，或者未知数的个数与方程组的个数相同。我们看到场方程（9-1）、（9-11）和（9-23）、界面条件（9-2）、（9-12）和（9-24）以及状态方程（9-55）不满足这个条件。因此，有必要添加几个本构方程来表达平均分子扩散、湍流传递和界面交换的传递机制。

通过采用热状态方程和热量状态方程，方程（9-56）和（9-57），双流体模型公式中出现的变量是： 1. 质量守恒 l , , ,； k k k k α ρ Γ v

\begin{sourceblock}
\includegraphics[page=190,pagebox=mediabox,trim={45.00bp 142.66bp 44.37bp 462.48bp},clip,width=0.92\textwidth,max height=0.38\textheight]{Ishii.pdf}
\end{sourceblock}
4. 状态方程 l , ,。 k k i u Τ Τ 其中 k = 1 和 2。因此，变量总数为 33。对于正确设置的模型，我们也应该有相同数量的方程。这些可以分为以下几组。

Τ T 2 湍流能量传递 k
Τ q 2 6) 体力场 l k g 2 7) 界面输运质量 1 Γ 1 动量 1 M 1 能量 1 E 1 8) 界面源动量 m M 1 能量 m E 1 9) 状态方程热状态方程 2 热量状态方程 2 10) 湍流动能 l l k k e u − 2 11) 指定界面温度 i 的相变条件Τ1
这表明我们也有三十三个方程，因此公式是一致的。但需要注意的是，上面所示的本构方程是以最原始的形式表达的，因此这些方程很可能通过一些附加参数与相同数量的补充本构方程相互耦合。此外，如果要使用熵不等式，则应在公式中引入方程（9-55）。

\subsection*{9.2.4 平均分子扩散通量}\phantomsection\addcontentsline{toc}{subsection}{9.2.4 平均分子扩散通量}
粘性应力张量 k T 和 k q 的本构方程可以通过使用恒等式 (9-50) 来研究。为简单起见，我们假设流体是牛顿流体，并且

\begin{sourceblock}
\includegraphics[page=192,pagebox=mediabox,trim={45.00bp 321.37bp 44.37bp 288.86bp},clip,width=0.92\textwidth,max height=0.38\textheight]{Ishii.pdf}
\end{sourceblock}
\begin{sourceblock}
\includegraphics[page=192,pagebox=mediabox,trim={45.00bp 296.65bp 44.37bp 350.51bp},clip,width=0.92\textwidth,max height=0.38\textheight]{Ishii.pdf}
\end{sourceblock}
\begin{sourceblock}
\includegraphics[page=192,pagebox=mediabox,trim={45.00bp 243.98bp 44.37bp 373.53bp},clip,width=0.92\textwidth,max height=0.38\textheight]{Ishii.pdf}
\end{sourceblock}
其中 k ' v 是第 k 相速度相对于 l k v 的脉动分量。很容易看出，当界面流体速度与平均速度之间的差异较大时，应力张量的第二部分变得很重要。因此，它考虑了界面运动和质量传递对平均变形的影响。让我们定义界面额外变形张量

\begin{sourceblock}
\includegraphics[page=192,pagebox=mediabox,trim={45.00bp 72.04bp 44.37bp 518.40bp},clip,width=0.92\textwidth,max height=0.38\textheight]{Ishii.pdf}
\end{sourceblock}
\begin{sourceblock}
\includegraphics[page=193,pagebox=mediabox,trim={45.00bp 546.49bp 44.37bp 82.72bp},clip,width=0.92\textwidth,max height=0.38\textheight]{Ishii.pdf}
\end{sourceblock}
因此，我们有

\begin{sourceblock}
\includegraphics[page=193,pagebox=mediabox,trim={45.00bp 487.00bp 44.37bp 154.29bp},clip,width=0.92\textwidth,max height=0.38\textheight]{Ishii.pdf}
\end{sourceblock}
如果公式中包含额外变形张量的影响，则应给出指定每个相的 ki D 的本构方程。一般来说，由于影响 ki D 的机制多种多样，它被认为是相当复杂的，但是，在特殊条件下，它可以简化为简单的形式。例如，如果c相是分散流中的连续相，并且界面的运动非常规则，相变的影响很小，则方程（9-73）和方程（4-62）可以近似为

\begin{sourceblock}
\includegraphics[page=193,pagebox=mediabox,trim={45.00bp 294.52bp 44.37bp 292.04bp},clip,width=0.92\textwidth,max height=0.38\textheight]{Ishii.pdf}
\end{sourceblock}
对于更一般的情况，我们可以将方程（9-73）近似为

\begin{sourceblock}
\includegraphics[page=193,pagebox=mediabox,trim={45.00bp 220.61bp 44.37bp 405.36bp},clip,width=0.92\textwidth,max height=0.38\textheight]{Ishii.pdf}
\end{sourceblock}
其中 i c a 代表 c 相的迁移率。传导传热 遵循傅立叶热传导定律方程（2-41）的流体的平均热通量 k q 可由下式给出

\begin{sourceblock}
\includegraphics[page=193,pagebox=mediabox,trim={45.00bp 86.01bp 44.37bp 531.45bp},clip,width=0.92\textwidth,max height=0.38\textheight]{Ishii.pdf}
\end{sourceblock}
其中我们使用了恒等式 (9-50) 并假设

\begin{sourceblock}
\includegraphics[page=194,pagebox=mediabox,trim={45.00bp 580.33bp 44.37bp 59.06bp},clip,width=0.92\textwidth,max height=0.38\textheight]{Ishii.pdf}
\end{sourceblock}
此外，如果我们假设界面处的热平衡，

\begin{sourceblock}
\includegraphics[page=194,pagebox=mediabox,trim={45.00bp 530.73bp 44.37bp 112.77bp},clip,width=0.92\textwidth,max height=0.38\textheight]{Ishii.pdf}
\end{sourceblock}
\begin{sourceblock}
\includegraphics[page=194,pagebox=mediabox,trim={45.00bp 498.94bp 44.37bp 123.16bp},clip,width=0.92\textwidth,max height=0.38\textheight]{Ishii.pdf}
\end{sourceblock}
则式(9-78)简化为

\begin{sourceblock}
\includegraphics[page=194,pagebox=mediabox,trim={45.00bp 417.82bp 44.37bp 200.84bp},clip,width=0.92\textwidth,max height=0.38\textheight]{Ishii.pdf}
\end{sourceblock}
其中我们使用了方程（4-62）。值得注意的是，第二项代表浓度梯度引起的热通量，在某种程度上它类似于单相混合物中的杜福尔效应（Hirschfelder 等，1954）。

\subsection*{9.2.5 湍流}\phantomsection\addcontentsline{toc}{subsection}{9.2.5 湍流}
湍流应力张量 - 混合长度模型 即使在单相流中，编写湍流本构方程时遇到的困难也是相当大的。湍流分析的基本问题是制定平均场方程的闭合方案。有两种不同的方法被广泛用于研究湍流的传输机制。第一种方法基于湍流本构方程的唯象构造。最好的代表是普朗特的混合长度假说，他通过类比气体运动学理论提出了湍流模型。

第二种方法是使用更准确的动力学方程来描述湍流传输。这可以通过采用动量方程的较高矩来完成。这样，动力学方程的数量可以根据需要增加。然而，该集合不是封闭的，因为作为力矩方程中的附加通量出现的湍流相关项始终比其他项高一个数量级。因此，这些方程永远无法在数学上闭合。因此，有必要进行一些近似并仅使用有限数量的动力学方程。与基于更高层次的统计理论相比
对于许多工程问题，混合长度模型仍然是获得解决方案的主要手段，特别是针对壁引起的湍流。包含耦合的高阶矩方程几乎总是需要大量的计算机计算，而在许多情况下，甚至积分方法也足以满足工程要求。

即使在单相流中，湍流的统计理论也没有牢固建立，并且该方法通常涉及相当复杂的方程组。因此，除了第 12 章 1.4 节中的分散两相流之外，我们不会讨论它们在两相流系统中的应用。由于其简单性，我们现在研究二流体模型公式中湍流通量的唯象方法。

根据应力张量的标准分析（Truesdell 和 Toupin，1960；Aris，1962；Slattery，1972），我们假设局部湍流应力 k
如果我们知道该点的相速度和周围相的变形，则可以确定 Τ T。上述假设满足局部作用的本构原理。此外，如果我们使用材料框架无差异原理，我们得出应力张量仅取决于变形张量的结论

\begin{sourceblock}
\includegraphics[page=195,pagebox=mediabox,trim={45.00bp 300.88bp 44.37bp 340.34bp},clip,width=0.92\textwidth,max height=0.38\textheight]{Ishii.pdf}
\end{sourceblock}
其中体积变形张量 kb D 和界面额外变形张量分别由式（9-74）和（9-73）给出。由体积变形引起的湍流应力可称为剪切引起的湍流，而由界面额外变形引起的湍流应力可称为气泡引起的湍流。

实际上，正如 Lumley (1970) 所讨论的那样，即使在单相湍流中，这两个本构原则也可能无法在严格意义上得到满足。由于可供依赖的实验事实很少，因此我们考虑上述假设有效时的简单情况。因此，条件下允许的最一般形式是

\begin{sourceblock}
\includegraphics[page=195,pagebox=mediabox,trim={45.00bp 126.94bp 44.37bp 512.06bp},clip,width=0.92\textwidth,max height=0.38\textheight]{Ishii.pdf}
\end{sourceblock}
其中系数 0 k a 、 1 k a 和 2 k a 是由 tr k D 、 : k k D D 和 det k D 给出的变形张量 k D 的三个不变量的函数。这些是，即，其自身的迹、双点积和行列式
分别为 kD。因此，我们有

\begin{sourceblock}
\includegraphics[page=196,pagebox=mediabox,trim={45.00bp 578.26bp 44.37bp 58.57bp},clip,width=0.92\textwidth,max height=0.38\textheight]{Ishii.pdf}
\end{sourceblock}
除了三个不变量之外，系数的参数还有： 流体密度 k ρ；粘度 k μ ;空隙率 k α；距墙 A 的距离；界面面积浓度 i a；和平均曲率21 H。式(9-83)和式(9-84)给出的湍流应力张量的表达式仍然非常复杂。然而，如果我们使用类似于单相流中的混合长度假设，结果就会简化为简单的形式。

首先，我们假设方程（9-83）的应力张量仅取决于第二项，即牛顿假设。然后我们有

\begin{sourceblock}
\includegraphics[page=196,pagebox=mediabox,trim={45.00bp 415.12bp 44.37bp 225.85bp},clip,width=0.92\textwidth,max height=0.38\textheight]{Ishii.pdf}
\end{sourceblock}
其中 T k μ 是湍流粘度。此外，这里系数 1 k a 取为

\begin{sourceblock}
\includegraphics[page=196,pagebox=mediabox,trim={45.00bp 339.16bp 44.37bp 297.67bp},clip,width=0.92\textwidth,max height=0.38\textheight]{Ishii.pdf}
\end{sourceblock}
这里需要注意的是，由于混合长度模型是针对剪切引起的湍流，因此使用体积变形张量 kb D 代替总变形 k D。由于其重要性，我们将在第 12 章中对其进行更详细的讨论。因此，从量纲分析中，我们定义

\begin{sourceblock}
\includegraphics[page=196,pagebox=mediabox,trim={45.00bp 215.79bp 44.37bp 411.69bp},clip,width=0.92\textwidth,max height=0.38\textheight]{Ishii.pdf}
\end{sourceblock}
\begin{sourceblock}
\includegraphics[page=196,pagebox=mediabox,trim={45.00bp 134.87bp 44.37bp 457.42bp},clip,width=0.92\textwidth,max height=0.38\textheight]{Ishii.pdf}
\end{sourceblock}
最终的表达式就变成了

\begin{sourceblock}
\includegraphics[page=196,pagebox=mediabox,trim={45.00bp 74.44bp 44.37bp 565.83bp},clip,width=0.92\textwidth,max height=0.38\textheight]{Ishii.pdf}
\end{sourceblock}
为了使模型可视化，让我们考虑一个非常简单的两相管流。通过采用没有相变的完全开发的流程，我们有

\begin{sourceblock}
\includegraphics[page=197,pagebox=mediabox,trim={45.00bp 380.94bp 44.37bp 132.50bp},clip,width=0.92\textwidth,max height=0.38\textheight]{Ishii.pdf}
\end{sourceblock}
\begin{sourceblock}
\includegraphics[page=197,pagebox=mediabox,trim={45.00bp 139.95bp 44.37bp 288.83bp},clip,width=0.92\textwidth,max height=0.38\textheight]{Ishii.pdf}
\end{sourceblock}
如果我们排除非常靠近墙壁的区域，则可以从函数 * T k μ 的参数中删除第一个无量纲群（局部雷诺数）。因此，在这种情况下，湍流的本构方程 Eq.(9-88) 仅取决于表示流动中某一点的平均几何构型的静态参数。

湍流传热-混合长度模型 湍流能通量由式(5-46)定义。从方程中我们可以看出，它由三部分组成，即内能的湍流传递、动能的湍流传递以及湍流所做的功。对于许多实际的两相流系统，后两种效应的作用不如单相流情况下的第一种效应那么重要。因此，我们构建了一个湍流热流模型，主要考虑热能传输的影响，即方程（5-46）的第一项和最后一项给出了焓传输。与式（9-81）类似，我们假设

\begin{sourceblock}
\includegraphics[page=198,pagebox=mediabox,trim={45.00bp 382.26bp 44.37bp 197.66bp},clip,width=0.92\textwidth,max height=0.38\textheight]{Ishii.pdf}
\end{sourceblock}
其中湍流能量传输系数表示为

\begin{sourceblock}
\includegraphics[page=198,pagebox=mediabox,trim={45.00bp 320.97bp 44.37bp 315.81bp},clip,width=0.92\textwidth,max height=0.38\textheight]{Ishii.pdf}
\end{sourceblock}
从量纲分析，我们引入

\begin{sourceblock}
\includegraphics[page=198,pagebox=mediabox,trim={45.00bp 245.56bp 44.37bp 378.50bp},clip,width=0.92\textwidth,max height=0.38\textheight]{Ishii.pdf}
\end{sourceblock}
这里 pk c 和 A 分别是比热和距壁的距离或混合长度。那么无量纲参数*k
Τ Κ 是四个相似组的函数：

\begin{sourceblock}
\includegraphics[page=198,pagebox=mediabox,trim={45.00bp 131.39bp 44.37bp 481.89bp},clip,width=0.92\textwidth,max height=0.38\textheight]{Ishii.pdf}
\end{sourceblock}
因此，湍流热通量可由下式给出

\begin{sourceblock}
\includegraphics[page=198,pagebox=mediabox,trim={45.00bp 51.69bp 44.37bp 566.95bp},clip,width=0.92\textwidth,max height=0.38\textheight]{Ishii.pdf}
\end{sourceblock}
\subsection*{9.2.6 界面输运本构关系}\phantomsection\addcontentsline{toc}{subsection}{9.2.6 界面输运本构关系}
由熵不等式(9-47)、热能方程(9-28)和状态方程(9-55)可以看出，熵产生与质量k Γ 、广义阻力的界面输运有关

\begin{sourceblock}
\includegraphics[page=199,pagebox=mediabox,trim={45.00bp 344.86bp 44.37bp 205.46bp},clip,width=0.92\textwidth,max height=0.38\textheight]{Ishii.pdf}
\end{sourceblock}
在这里，我们的分析基于这些效应独立满足熵不等式的假设。不可逆热力学的标准理论（De Groot 和 Mazur，1962）给出了获得线性本构方程的简单方法。为此，首先我们应该将熵不等式中的项排列成通量和势的合适组合（通量根据势线性展开）。这里我们应该特别注意，因为每个阶段都有方程（9-97）中的两个不等式；传质项k Г 和广义阻力ik M 应满足跳跃条件（8-5）和（8-19）。由于在许多实际问题中，H m M 和 t k M 的数量级远小于阻力本身，因此我们可以近似为 2

\begin{sourceblock}
\includegraphics[page=199,pagebox=mediabox,trim={45.00bp 157.82bp 44.37bp 480.01bp},clip,width=0.92\textwidth,max height=0.38\textheight]{Ishii.pdf}
\end{sourceblock}
遵循不平等

\begin{sourceblock}
\includegraphics[page=199,pagebox=mediabox,trim={45.00bp 88.03bp 44.37bp 523.51bp},clip,width=0.92\textwidth,max height=0.38\textheight]{Ishii.pdf}
\end{sourceblock}
\begin{sourceblock}
\includegraphics[page=200,pagebox=mediabox,trim={45.00bp 517.29bp 44.37bp 55.15bp},clip,width=0.92\textwidth,max height=0.38\textheight]{Ishii.pdf}
\end{sourceblock}
此外，如果我们忽略质量传递引起的推力和界面处的法向应力，则从式（2-104）和（9-44）我们有

\begin{sourceblock}
\includegraphics[page=200,pagebox=mediabox,trim={45.00bp 417.77bp 44.37bp 193.65bp},clip,width=0.92\textwidth,max height=0.38\textheight]{Ishii.pdf}
\end{sourceblock}
式(8-11)给出的界面处总动量通量可以简化为

\begin{sourceblock}
\includegraphics[page=200,pagebox=mediabox,trim={45.00bp 343.18bp 44.37bp 297.92bp},clip,width=0.92\textwidth,max height=0.38\textheight]{Ishii.pdf}
\end{sourceblock}
因此，界面处的压力应与以下关系相关：

\begin{sourceblock}
\includegraphics[page=200,pagebox=mediabox,trim={45.00bp 284.20bp 44.37bp 357.03bp},clip,width=0.92\textwidth,max height=0.38\textheight]{Ishii.pdf}
\end{sourceblock}
方程（9-100）自动满足，方程（8-22）、（8-23）和（8-25）作为界面动量传递条件的法向分量。

在下文中，我们假设本体流体和界面处的相平均值之间的差异的影响对于密度和压力可以忽略不计，但对于温度则不然。因此我们取

\begin{sourceblock}
\includegraphics[page=200,pagebox=mediabox,trim={45.00bp 158.20bp 44.37bp 483.03bp},clip,width=0.92\textwidth,max height=0.38\textheight]{Ishii.pdf}
\end{sourceblock}
\begin{sourceblock}
\includegraphics[page=200,pagebox=mediabox,trim={45.00bp 114.15bp 44.37bp 527.85bp},clip,width=0.92\textwidth,max height=0.38\textheight]{Ishii.pdf}
\end{sourceblock}
在这些假设下，界面输运项的简单线性本构方程可以表示为以下形式

\begin{sourceblock}
\includegraphics[page=201,pagebox=mediabox,trim={45.00bp 578.47bp 44.37bp 58.73bp},clip,width=0.92\textwidth,max height=0.38\textheight]{Ishii.pdf}
\end{sourceblock}
\begin{sourceblock}
\includegraphics[page=201,pagebox=mediabox,trim={45.00bp 532.18bp 44.37bp 108.99bp},clip,width=0.92\textwidth,max height=0.38\textheight]{Ishii.pdf}
\end{sourceblock}
\begin{sourceblock}
\includegraphics[page=201,pagebox=mediabox,trim={45.00bp 479.34bp 44.37bp 153.75bp},clip,width=0.92\textwidth,max height=0.38\textheight]{Ishii.pdf}
\end{sourceblock}
其中传输系数 k b Γ 、 1 M b 和 E k b 被认为是正标量。

界面传质项 我们假设方程（9-104）中的传递系数 k b Γ 是以下参数的函数

\begin{sourceblock}
\includegraphics[page=201,pagebox=mediabox,trim={45.00bp 349.17bp 44.37bp 286.95bp},clip,width=0.92\textwidth,max height=0.38\textheight]{Ishii.pdf}
\end{sourceblock}
为了简化上式，我们首先引入一个无量纲参数

\begin{sourceblock}
\includegraphics[page=201,pagebox=mediabox,trim={45.00bp 249.92bp 44.37bp 363.44bp},clip,width=0.92\textwidth,max height=0.38\textheight]{Ishii.pdf}
\end{sourceblock}
雅各布数定义为

\begin{sourceblock}
\includegraphics[page=201,pagebox=mediabox,trim={45.00bp 166.41bp 44.37bp 450.51bp},clip,width=0.92\textwidth,max height=0.38\textheight]{Ishii.pdf}
\end{sourceblock}
\begin{sourceblock}
\includegraphics[page=201,pagebox=mediabox,trim={45.00bp 96.34bp 44.37bp 520.48bp},clip,width=0.92\textwidth,max height=0.38\textheight]{Ishii.pdf}
\end{sourceblock}
因此，界面传质项可以重写为

\begin{sourceblock}
\includegraphics[page=202,pagebox=mediabox,trim={45.00bp 494.04bp 44.37bp 55.01bp},clip,width=0.92\textwidth,max height=0.38\textheight]{Ishii.pdf}
\end{sourceblock}
其中无量纲函数 * k b Γ 可以用四个相似群表示为

\begin{sourceblock}
\includegraphics[page=202,pagebox=mediabox,trim={45.00bp 403.79bp 44.37bp 213.83bp},clip,width=0.92\textwidth,max height=0.38\textheight]{Ishii.pdf}
\end{sourceblock}
雅各布数由方程（9-109）定义，是可用能量的规模。众所周知，它是气泡生长分析中的一个重要参数。

现在让我们研究一些特殊情况，其中 1 Γ 的本构方程可以简化为简单形式。在许多实际工程问题中，我们可以假设气相处于饱和状态，因此可以取

\begin{sourceblock}
\includegraphics[page=202,pagebox=mediabox,trim={45.00bp 261.37bp 44.37bp 375.83bp},clip,width=0.92\textwidth,max height=0.38\textheight]{Ishii.pdf}
\end{sourceblock}
在此条件下，式(9-111)简化为

\begin{sourceblock}
\includegraphics[page=202,pagebox=mediabox,trim={45.00bp 173.64bp 44.37bp 441.32bp},clip,width=0.92\textwidth,max height=0.38\textheight]{Ishii.pdf}
\end{sourceblock}
\begin{sourceblock}
\includegraphics[page=202,pagebox=mediabox,trim={45.00bp 79.77bp 44.37bp 512.34bp},clip,width=0.92\textwidth,max height=0.38\textheight]{Ishii.pdf}
\end{sourceblock}
界面阻力 ik M 的一般表达式由方程（9-105）假设。现在我们进一步假设系数 1
M b 取决于以下参数。

\begin{sourceblock}
\includegraphics[page=203,pagebox=mediabox,trim={45.00bp 448.96bp 44.37bp 187.87bp},clip,width=0.92\textwidth,max height=0.38\textheight]{Ishii.pdf}
\end{sourceblock}
那么从量纲分析，我们可以将式(9-105)改写为

\begin{sourceblock}
\includegraphics[page=203,pagebox=mediabox,trim={45.00bp 388.35bp 44.37bp 252.22bp},clip,width=0.92\textwidth,max height=0.38\textheight]{Ishii.pdf}
\end{sourceblock}
where

\begin{sourceblock}
\includegraphics[page=203,pagebox=mediabox,trim={45.00bp 311.62bp 44.37bp 313.36bp},clip,width=0.92\textwidth,max height=0.38\textheight]{Ishii.pdf}
\end{sourceblock}
\begin{sourceblock}
\includegraphics[page=203,pagebox=mediabox,trim={45.00bp 230.62bp 44.37bp 364.90bp},clip,width=0.92\textwidth,max height=0.38\textheight]{Ishii.pdf}
\end{sourceblock}
这里我们定义了界面雷诺数

\begin{sourceblock}
\includegraphics[page=203,pagebox=mediabox,trim={45.00bp 153.77bp 44.37bp 468.70bp},clip,width=0.92\textwidth,max height=0.38\textheight]{Ishii.pdf}
\end{sourceblock}
和相变效应数

\begin{sourceblock}
\includegraphics[page=203,pagebox=mediabox,trim={45.00bp 79.67bp 44.37bp 543.82bp},clip,width=0.92\textwidth,max height=0.38\textheight]{Ishii.pdf}
\end{sourceblock}
对于分散流态，关于阻力的研究有很多。对于具有恒定直径的实心球形颗粒的稀悬浮液的流动，分析相对容易。然而，随着分散相的空隙率增加或壁效应变得重要，问题变得越来越复杂。显然，对于具有可变形界面、界面传质和湍流的流动，阻力相关性应该广泛依赖于实验数据。

关于分散流阻力定律的一些结果可以在 Brodkey (1967)、Soo (1967)、Wallis (1969)、Schlichting (1979)、Happel 和 Brenner (1965) 中找到。我们在第 9 章的 1.4 节中讨论了重要的特殊情况。更一般和完整的建模和讨论将在第 12 章中介绍。

界面热通量 界面传热本构方程由式（9-106）假设。首先我们引入一个无量纲传热系数1

\begin{sourceblock}
\includegraphics[page=204,pagebox=mediabox,trim={45.00bp 316.23bp 44.37bp 328.30bp},clip,width=0.92\textwidth,max height=0.38\textheight]{Ishii.pdf}
\end{sourceblock}
\begin{sourceblock}
\includegraphics[page=204,pagebox=mediabox,trim={45.00bp 263.73bp 44.37bp 331.90bp},clip,width=0.92\textwidth,max height=0.38\textheight]{Ishii.pdf}
\end{sourceblock}
然后我们有

\begin{sourceblock}
\includegraphics[page=204,pagebox=mediabox,trim={45.00bp 197.06bp 44.37bp 439.70bp},clip,width=0.92\textwidth,max height=0.38\textheight]{Ishii.pdf}
\end{sourceblock}
\begin{sourceblock}
\includegraphics[page=204,pagebox=mediabox,trim={45.00bp 109.55bp 44.37bp 479.80bp},clip,width=0.92\textwidth,max height=0.38\textheight]{Ishii.pdf}
\end{sourceblock}
其中普朗特数定义为

\begin{sourceblock}
\includegraphics[page=205,pagebox=mediabox,trim={45.00bp 562.90bp 44.37bp 57.07bp},clip,width=0.92\textwidth,max height=0.38\textheight]{Ishii.pdf}
\end{sourceblock}
\begin{sourceblock}
\includegraphics[page=205,pagebox=mediabox,trim={45.00bp 542.14bp 44.37bp 90.30bp},clip,width=0.92\textwidth,max height=0.38\textheight]{Ishii.pdf}
\end{sourceblock}
k N 和相变效应数i pch N 分别由式(9-120)和(9-121)给出。有趣的是，这里实际上无量纲参数 1
* E b 是界面努塞尔数。此外，如果我们给出 1 Γ 的本构方程，那么仅提供一个本构关系就足够了
界面传热，1 i a q ''。然而，也可以给出一个
对于 1 i a q '' 和 2 i a q '' 的本构关系，则等价于由于宏观能量跳跃条件已知 1 Γ。

界面剪切应力 由于气泡分散在连续相剪切层中，因此界面剪切应力近似为连续相剪切应力。因此，我们有

\begin{sourceblock}
\includegraphics[page=205,pagebox=mediabox,trim={45.00bp 334.12bp 44.37bp 307.28bp},clip,width=0.92\textwidth,max height=0.38\textheight]{Ishii.pdf}
\end{sourceblock}
界面动量源 在原始动量跳跃条件下，有两个不同的信息：法向跳跃和切向跳跃平衡。由于我们在平均公式中保留了这一特殊特性，因此除了方程（9-12）给出的界面动量传递条件之外，我们还获得了阻力平衡（8-19）。此外，通过忽略质量推力效应并使用方程（9-103）的假设，我们从方程（8-21）得到

\begin{sourceblock}
\includegraphics[page=205,pagebox=mediabox,trim={45.00bp 146.98bp 44.37bp 457.21bp},clip,width=0.92\textwidth,max height=0.38\textheight]{Ishii.pdf}
\end{sourceblock}
因此，根据式（8-19）、（9-12）和（9-13），我们有

\begin{sourceblock}
\includegraphics[page=205,pagebox=mediabox,trim={45.00bp 87.16bp 44.37bp 554.08bp},clip,width=0.92\textwidth,max height=0.38\textheight]{Ishii.pdf}
\end{sourceblock}
and

\begin{sourceblock}
\includegraphics[page=206,pagebox=mediabox,trim={45.00bp 568.48bp 44.37bp 58.51bp},clip,width=0.92\textwidth,max height=0.38\textheight]{Ishii.pdf}
\end{sourceblock}
这里我们有界面的热状态方程

\begin{sourceblock}
\includegraphics[page=206,pagebox=mediabox,trim={45.00bp 505.38bp 44.37bp 131.77bp},clip,width=0.92\textwidth,max height=0.38\textheight]{Ishii.pdf}
\end{sourceblock}
由于界面动量传递方程（9-129）的法向分量规定了两相之间的机械平衡条件，因此需要通过本构方程规定平均曲率21 H。一个简单的情况是假设界面几何形状完全不规则，因此我们可以取 21 0 H =。对于分散两相流，平均曲率是颗粒半径的倒数。如果流体颗粒尺寸均匀，则半径由 3 i a α 给出。因此，α与i a 的比值通常是一个重要的长度尺度。 6 i a α 被称为索特平均直径，因此平均曲率基本上与索特平均直径的倒数成比例。

然而，参数 i a 的重要性仍然存在于宏观公式中，因为它代表两相之间的可用接触面积。很明显，质量、动量和能量的界面传输受到单位体积表面积浓度 i a 的显著影响。一般来说，本构方程为
21 H 和 i a 极其复杂，因为它们是决定宏观场中局部几何构型的参数。显然，我们可以直接或间接提供这些信息。

直接信息意味着我们对流程结构有先验知识。那么在各种变量以及初始条件和边界条件下获得 21 H 和 i a 的关系并不困难。例如，对于气泡或液滴流，这可以很容易地完成，而无需气泡（或液滴）的相变、聚结或分裂。还可以通过 21 H 和 i a 的本构方程以各种参数（例如 k α 、 σ 、 l l 2 1 − v v 、 k μ 、 k ）给出有关流动结构的间接信息
然后我们可以求解整个方程组来找到局部几何配置。这很困难，因为几何构型具有长期记忆，并且在大多数情况下不遵守局部作用原理。这意味着初始条件以及壁面效应对流动几何形状非常重要。

由于一般情况下遇到的困难，让我们从一个简单的情况开始讨论上述本构方程。现在让

\begin{sourceblock}
\includegraphics[page=207,pagebox=mediabox,trim={45.00bp 526.96bp 44.37bp 109.87bp},clip,width=0.92\textwidth,max height=0.38\textheight]{Ishii.pdf}
\end{sourceblock}
体积 V 中相 2 的表面积由下式给出

\begin{sourceblock}
\includegraphics[page=207,pagebox=mediabox,trim={45.00bp 461.98bp 44.37bp 174.85bp},clip,width=0.92\textwidth,max height=0.38\textheight]{Ishii.pdf}
\end{sourceblock}
完全相同的论证可以在时域中进行，因此我们有

\begin{sourceblock}
\includegraphics[page=207,pagebox=mediabox,trim={45.00bp 383.97bp 44.37bp 252.81bp},clip,width=0.92\textwidth,max height=0.38\textheight]{Ishii.pdf}
\end{sourceblock}
and

\begin{sourceblock}
\includegraphics[page=207,pagebox=mediabox,trim={45.00bp 310.60bp 44.37bp 316.73bp},clip,width=0.92\textwidth,max height=0.38\textheight]{Ishii.pdf}
\end{sourceblock}
那么我们可以假设

\begin{sourceblock}
\includegraphics[page=207,pagebox=mediabox,trim={45.00bp 221.72bp 44.37bp 390.50bp},clip,width=0.92\textwidth,max height=0.38\textheight]{Ishii.pdf}
\end{sourceblock}
其中 i C 是考虑分散相形状和尺寸的因素。因此我们可以将 21 H 的本构方程写为

\begin{sourceblock}
\includegraphics[page=207,pagebox=mediabox,trim={45.00bp 136.63bp 44.37bp 490.94bp},clip,width=0.92\textwidth,max height=0.38\textheight]{Ishii.pdf}
\end{sourceblock}
这里，对于相当均匀的球形液滴或气泡，因子 i C 为 1，并且它不会变化太大，除非分散相具有相当细长的形状。方程（9-136）给出的关系是静态或几何关系，我们可以采用更一般的形式：

\begin{sourceblock}
\includegraphics[page=208,pagebox=mediabox,trim={45.00bp 520.34bp 44.37bp 59.20bp},clip,width=0.92\textwidth,max height=0.38\textheight]{Ishii.pdf}
\end{sourceblock}
我们将这种关系称为状态的几何方程。从方程（9-137）我们可以看出，平均曲率取决于空隙率、表面积浓度和空隙率梯度。

然而，我们仍然应该有一个更重要的 i a 本构方程。认为关于其他参数的 i a 信息实际上是局部即时公式的解的一部分。在双流体公式中包含 i a 的最通用方法是再引入一个界面面积浓度的传输方程：

\begin{sourceblock}
\includegraphics[page=208,pagebox=mediabox,trim={45.00bp 343.34bp 44.37bp 281.53bp},clip,width=0.92\textwidth,max height=0.38\textheight]{Ishii.pdf}
\end{sourceblock}
使用这种方法，源项考虑了气泡或液滴的膨胀或塌陷、聚结、分解和界面不稳定性。显然，应该提供 l i v 和 L φ 的本构方程。界面面积输运方程是描述相间表面积变化的基本方程。由于其重要性，第 10 章对其进行了详细讨论。在某些情况下，平衡方程 (9-138) 可以用更简单的代数本构关系代替，例如

\begin{sourceblock}
\includegraphics[page=208,pagebox=mediabox,trim={45.00bp 191.84bp 44.37bp 446.11bp},clip,width=0.92\textwidth,max height=0.38\textheight]{Ishii.pdf}
\end{sourceblock}
\begin{sourceblock}
\includegraphics[page=208,pagebox=mediabox,trim={45.00bp 126.15bp 44.37bp 492.47bp},clip,width=0.92\textwidth,max height=0.38\textheight]{Ishii.pdf}
\end{sourceblock}
其中相 2 是分散相。然而，在许多实际问题中，与 1 i M 或
2米。在这种情况下，我们可以将H m M 设置为零。

界面能量源 界面能量源m E 由式(8-45)给出。显然，包含界面热能项，即方程（8-45）给出的项，使公式显著复杂化。除了极少数情况外，对于涉及相变时的潜热的大能量交换，可以忽略该术语。因此，我们将方程（8-45）近似为

\begin{sourceblock}
\includegraphics[page=209,pagebox=mediabox,trim={45.00bp 466.29bp 44.37bp 158.53bp},clip,width=0.92\textwidth,max height=0.38\textheight]{Ishii.pdf}
\end{sourceblock}
不需要任何额外的本构方程。

为了完成式（8-48）给出的界面能量传递条件的模型，我们应该提供湍流动能、平均速度与平均界面相速度之差以及阻力T ki W 产生的界面湍流通量的本构方程。由于在具有相变的两相流中，这些项的数量级与热项相比相对较小，因此，我们可以假设

\begin{sourceblock}
\includegraphics[page=209,pagebox=mediabox,trim={45.00bp 297.55bp 44.37bp 325.87bp},clip,width=0.92\textwidth,max height=0.38\textheight]{Ishii.pdf}
\end{sourceblock}
\begin{sourceblock}
\includegraphics[page=209,pagebox=mediabox,trim={45.00bp 255.93bp 44.37bp 385.10bp},clip,width=0.92\textwidth,max height=0.38\textheight]{Ishii.pdf}
\end{sourceblock}
\begin{sourceblock}
\includegraphics[page=209,pagebox=mediabox,trim={45.00bp 211.96bp 44.37bp 427.03bp},clip,width=0.92\textwidth,max height=0.38\textheight]{Ishii.pdf}
\end{sourceblock}
与式（9-142）类似，我们取体相

\begin{sourceblock}
\includegraphics[page=209,pagebox=mediabox,trim={45.00bp 153.93bp 44.37bp 484.65bp},clip,width=0.92\textwidth,max height=0.38\textheight]{Ishii.pdf}
\end{sourceblock}
那么，式（8-48）可以简化为以下形式

\begin{sourceblock}
\includegraphics[page=209,pagebox=mediabox,trim={45.00bp 55.42bp 44.37bp 542.22bp},clip,width=0.92\textwidth,max height=0.38\textheight]{Ishii.pdf}
\end{sourceblock}
\begin{sourceblock}
\includegraphics[page=210,pagebox=mediabox,trim={45.00bp 530.19bp 44.37bp 99.02bp},clip,width=0.92\textwidth,max height=0.38\textheight]{Ishii.pdf}
\end{sourceblock}
那么，热能传递条件（8-49）可以近似为

\begin{sourceblock}
\includegraphics[page=210,pagebox=mediabox,trim={45.00bp 460.29bp 44.37bp 165.73bp},clip,width=0.92\textwidth,max height=0.38\textheight]{Ishii.pdf}
\end{sourceblock}
and

\begin{sourceblock}
\includegraphics[page=210,pagebox=mediabox,trim={45.00bp 351.06bp 44.37bp 235.99bp},clip,width=0.92\textwidth,max height=0.38\textheight]{Ishii.pdf}
\end{sourceblock}
其中我们使用方程（9-128）来消除表面张力项。将式(9-148)与式(9-149)联立可得

\begin{sourceblock}
\includegraphics[page=210,pagebox=mediabox,trim={45.00bp 226.17bp 44.37bp 359.47bp},clip,width=0.92\textwidth,max height=0.38\textheight]{Ishii.pdf}
\end{sourceblock}
上式右侧的第一组是表面张力的影响，第二组来自界面阻力功，第三项与界面剪切所做的功有关。因此，我们可以假设对于相对低速的流动

\begin{sourceblock}
\includegraphics[page=210,pagebox=mediabox,trim={45.00bp 114.27bp 44.37bp 514.94bp},clip,width=0.92\textwidth,max height=0.38\textheight]{Ishii.pdf}
\end{sourceblock}
这是局部即时公式的众所周知的关系。我们在这里展示了可以将这个重要且有用的公式应用于宏观两相流问题的条件。

\section*{9.3 二流体模型制定}\phantomsection\addcontentsline{toc}{section}{9.3 二流体模型制定}
双流体模型公式的最一般情况已在第 9 章 1.2.3 节中结合决定论原理进行了讨论。我们现在将结合前两节的结果来建立一个现实的公式。我们已经对界面变量和本构方程做出了一些假设，因此目前的分析并不是局部瞬时变量的微观场到宏观场的完整映射。相反，它应该被视为基于各种本构假设的近似理论。本节中给出的结果在一定程度上被简化为现实，但对于两相流系统分析中遇到的大多数工程问题来说它足够通用。

首先，我们列出为获得模型而做出的所有重要假设。

• 平均值平滑的基本假设 第 4 章 1.3 节 • 状态方程的存在性 Eq.(9-55) • 传输特性 k μ 和 k K 在时间平均区间内为常数 Eq.(9-71) Eq.(9-79) • 界面变量近似为

\begin{sourceblock}
\includegraphics[page=211,pagebox=mediabox,trim={45.00bp 244.11bp 44.37bp 361.23bp},clip,width=0.92\textwidth,max height=0.38\textheight]{Ishii.pdf}
\end{sourceblock}
ki k p p ≈ Eq.(9-103) • 忽略质量传递引起的界面法向应力和推力 Eq.(9-127) • 忽略湍流动能或能量传递 Eq.(9-145) • 忽略界面能量传递条件下的机械相互作用项 Eq.(9-148) • 均匀体积力 Eq.(5-50) 在这些条件下，我们有以下场方程 连续性方程（9-1）

\begin{sourceblock}
\includegraphics[page=211,pagebox=mediabox,trim={45.00bp 51.74bp 44.37bp 573.14bp},clip,width=0.92\textwidth,max height=0.38\textheight]{Ishii.pdf}
\end{sourceblock}
方程（9-15）的运动方程

\begin{sourceblock}
\includegraphics[page=212,pagebox=mediabox,trim={45.00bp 504.50bp 44.37bp 80.70bp},clip,width=0.92\textwidth,max height=0.38\textheight]{Ishii.pdf}
\end{sourceblock}
值得注意的是，上式中的最后一项被保留，尽管在大多数情况下它非常小，因为对于某些情况（例如水平流）它可能很重要。

由式(9-30)得到热能方程

\begin{sourceblock}
\includegraphics[page=212,pagebox=mediabox,trim={45.00bp 337.63bp 44.37bp 233.79bp},clip,width=0.92\textwidth,max height=0.38\textheight]{Ishii.pdf}
\end{sourceblock}
Τ α ⋅∇⋅ v T，因为它被认为主要贡献于在公式中被忽略的湍流动能变化。这两组三个平衡方程描述了宏观领域中质量、动量和能量守恒的物理定律。

这里我们忽略了湍流功项 l ( ) k k k
由它们自己的场方程控制的两个相通过下面给出的三个界面输运条件耦合。

由式(9-2)得到界面传质条件

\begin{sourceblock}
\includegraphics[page=212,pagebox=mediabox,trim={45.00bp 158.19bp 44.37bp 471.02bp},clip,width=0.92\textwidth,max height=0.38\textheight]{Ishii.pdf}
\end{sourceblock}
式（9-127）和（9-128）中的界面动量传递条件

\begin{sourceblock}
\includegraphics[page=212,pagebox=mediabox,trim={45.00bp 72.27bp 44.37bp 556.94bp},clip,width=0.92\textwidth,max height=0.38\textheight]{Ishii.pdf}
\end{sourceblock}
法线分量满足

\begin{sourceblock}
\includegraphics[page=213,pagebox=mediabox,trim={45.00bp 582.16bp 44.37bp 59.05bp},clip,width=0.92\textwidth,max height=0.38\textheight]{Ishii.pdf}
\end{sourceblock}
由式(9-151)可知界面热能传递条件

\begin{sourceblock}
\includegraphics[page=213,pagebox=mediabox,trim={45.00bp 505.47bp 44.37bp 123.74bp},clip,width=0.92\textwidth,max height=0.38\textheight]{Ishii.pdf}
\end{sourceblock}
那么，根据连续性公理我们有

\begin{sourceblock}
\includegraphics[page=213,pagebox=mediabox,trim={45.00bp 450.12bp 44.37bp 191.05bp},clip,width=0.92\textwidth,max height=0.38\textheight]{Ishii.pdf}
\end{sourceblock}
每个相的状态方程由方程（9-55）或方程（9-58）和（9-59）给出，因此我们有状态热量方程

\begin{sourceblock}
\includegraphics[page=213,pagebox=mediabox,trim={45.00bp 374.47bp 44.37bp 261.90bp},clip,width=0.92\textwidth,max height=0.38\textheight]{Ishii.pdf}
\end{sourceblock}
和热状态方程

\begin{sourceblock}
\includegraphics[page=213,pagebox=mediabox,trim={45.00bp 309.37bp 44.37bp 327.83bp},clip,width=0.92\textwidth,max height=0.38\textheight]{Ishii.pdf}
\end{sourceblock}
而表面的状态方程由方程（9-130）给出

\begin{sourceblock}
\includegraphics[page=213,pagebox=mediabox,trim={45.00bp 244.38bp 44.37bp 392.77bp},clip,width=0.92\textwidth,max height=0.38\textheight]{Ishii.pdf}
\end{sourceblock}
界面温度由相变条件 (9-44) 给出：

\begin{sourceblock}
\includegraphics[page=213,pagebox=mediabox,trim={45.00bp 166.89bp 44.37bp 454.74bp},clip,width=0.92\textwidth,max height=0.38\textheight]{Ishii.pdf}
\end{sourceblock}
where

\begin{sourceblock}
\includegraphics[page=213,pagebox=mediabox,trim={45.00bp 102.37bp 44.37bp 534.83bp},clip,width=0.92\textwidth,max height=0.38\textheight]{Ishii.pdf}
\end{sourceblock}
是经典的饱和条件。对于许多实际情况，我们可以将方程（9-163）近似为

\begin{sourceblock}
\includegraphics[page=214,pagebox=mediabox,trim={45.00bp 578.47bp 44.37bp 58.73bp},clip,width=0.92\textwidth,max height=0.38\textheight]{Ishii.pdf}
\end{sourceblock}
其中 g p 是气相压力。粘性应力 k T 的本构方程由下式给出

\begin{sourceblock}
\includegraphics[page=214,pagebox=mediabox,trim={45.00bp 503.20bp 44.37bp 138.06bp},clip,width=0.92\textwidth,max height=0.38\textheight]{Ishii.pdf}
\end{sourceblock}
这里体积和界面额外变形张量 kb D 和 ki D 由方程（9-74）和（9-77）给出。因此我们有

\begin{sourceblock}
\includegraphics[page=214,pagebox=mediabox,trim={45.00bp 418.69bp 44.37bp 210.51bp},clip,width=0.92\textwidth,max height=0.38\textheight]{Ishii.pdf}
\end{sourceblock}
and

\begin{sourceblock}
\includegraphics[page=214,pagebox=mediabox,trim={45.00bp 345.88bp 44.37bp 280.07bp},clip,width=0.92\textwidth,max height=0.38\textheight]{Ishii.pdf}
\end{sourceblock}
系数i k a 表示第k相的迁移率。由式(9-85)、(9-86)、(9-87)、(9-88)和(9-89)的混合长度假设得到了湍流应力T k T的本构方程，因此

\begin{sourceblock}
\includegraphics[page=214,pagebox=mediabox,trim={45.00bp 224.75bp 44.37bp 391.40bp},clip,width=0.92\textwidth,max height=0.38\textheight]{Ishii.pdf}
\end{sourceblock}
这里无量纲湍流粘度 * T k μ 是以下参数的函数

\begin{sourceblock}
\includegraphics[page=214,pagebox=mediabox,trim={45.00bp 128.81bp 44.37bp 488.83bp},clip,width=0.92\textwidth,max height=0.38\textheight]{Ishii.pdf}
\end{sourceblock}
我们在这里回顾一下，A 和 i a 分别是距壁的距离和表面积浓度。如果我们排除非常靠近墙壁的区域，则可以从 * T k μ 的参数中删除第一个参数。平均传导热通量 k q 的本构关系由式（9-81）给出，因此有

\begin{sourceblock}
\includegraphics[page=215,pagebox=mediabox,trim={45.00bp 562.89bp 44.37bp 55.75bp},clip,width=0.92\textwidth,max height=0.38\textheight]{Ishii.pdf}
\end{sourceblock}
有趣的是，第二项由于浓度梯度而代表了体相和界面之间温差的影响，即热非平衡。

\begin{sourceblock}
\includegraphics[page=215,pagebox=mediabox,trim={45.00bp 446.79bp 44.37bp 156.51bp},clip,width=0.92\textwidth,max height=0.38\textheight]{Ishii.pdf}
\end{sourceblock}
其中无量纲电导率 * T k Κ 是以下参数的函数

\begin{sourceblock}
\includegraphics[page=215,pagebox=mediabox,trim={45.00bp 355.79bp 44.37bp 261.86bp},clip,width=0.92\textwidth,max height=0.38\textheight]{Ishii.pdf}
\end{sourceblock}
我们在这里注意到，如果我们排除非常靠近墙壁的区域，则可以从函数 * T k Κ 的参数中删除第一个无量纲组。传质项 1 Γ 由式（9-111）和（9-112）给出，因此我们有

\begin{sourceblock}
\includegraphics[page=215,pagebox=mediabox,trim={45.00bp 183.58bp 44.37bp 384.80bp},clip,width=0.92\textwidth,max height=0.38\textheight]{Ishii.pdf}
\end{sourceblock}
\begin{sourceblock}
\includegraphics[page=215,pagebox=mediabox,trim={45.00bp 100.79bp 44.37bp 491.26bp},clip,width=0.92\textwidth,max height=0.38\textheight]{Ishii.pdf}
\end{sourceblock}
Jk N 表示的群是式（9-109）和（9-110）定义的雅各布数，它是相位变化最重要的参数。一个
鉴于界面动量传递条件，方程（9-156），我们应该提供两个本构方程来指定曳力 1 i M 和表面张力的影响。界面阻力由式（9-117）和（9-119）给出，因此我们有

\begin{sourceblock}
\includegraphics[page=216,pagebox=mediabox,trim={45.00bp 492.40bp 44.37bp 148.21bp},clip,width=0.92\textwidth,max height=0.38\textheight]{Ishii.pdf}
\end{sourceblock}
其中系数 1 * M b 预计是以下参数的函数

\begin{sourceblock}
\includegraphics[page=216,pagebox=mediabox,trim={45.00bp 397.61bp 44.37bp 220.03bp},clip,width=0.92\textwidth,max height=0.38\textheight]{Ishii.pdf}
\end{sourceblock}
\begin{sourceblock}
\includegraphics[page=216,pagebox=mediabox,trim={45.00bp 375.95bp 44.37bp 275.29bp},clip,width=0.92\textwidth,max height=0.38\textheight]{Ishii.pdf}
\end{sourceblock}
k N 和相变效应数i pch N 分别由式(9-120)和(9-121)定义。此外，几何状态方程Eq.(9-137)指定了界面的平均曲率

\begin{sourceblock}
\includegraphics[page=216,pagebox=mediabox,trim={45.00bp 280.60bp 44.37bp 346.93bp},clip,width=0.92\textwidth,max height=0.38\textheight]{Ishii.pdf}
\end{sourceblock}
式（9-158）给出的界面热能传递条件需要界面处传热的本构方程。因此，从方程（9-123）和（9-124）我们有

\begin{sourceblock}
\includegraphics[page=216,pagebox=mediabox,trim={45.00bp 189.98bp 44.37bp 446.78bp},clip,width=0.92\textwidth,max height=0.38\textheight]{Ishii.pdf}
\end{sourceblock}
where

\begin{sourceblock}
\includegraphics[page=216,pagebox=mediabox,trim={45.00bp 108.59bp 44.37bp 509.05bp},clip,width=0.92\textwidth,max height=0.38\textheight]{Ishii.pdf}
\end{sourceblock}
普朗特数的定义由式(9-125)给出。在许多实际问题中，可以假设分散相处于热平衡状态，则本构方程（9-179）简化为

\begin{sourceblock}
\includegraphics[page=217,pagebox=mediabox,trim={45.00bp 518.18bp 44.37bp 106.75bp},clip,width=0.92\textwidth,max height=0.38\textheight]{Ishii.pdf}
\end{sourceblock}
其中源项 L φ 由已经出现的各种参数表示。一般而言，预计本构方程为
L φ 很难获得，除非流动几何形状非常简单，即气泡流或液滴流。平均界面速度
l i v 大约可以由下式给出

\begin{sourceblock}
\includegraphics[page=217,pagebox=mediabox,trim={45.00bp 396.58bp 44.37bp 227.72bp},clip,width=0.92\textwidth,max height=0.38\textheight]{Ishii.pdf}
\end{sourceblock}
其中下标d表示分散相。

分散流的 H m M 本构方程可由方程 1 给出。 （9-140），因此我们有

\begin{sourceblock}
\includegraphics[page=217,pagebox=mediabox,trim={45.00bp 307.06bp 44.37bp 329.77bp},clip,width=0.92\textwidth,max height=0.38\textheight]{Ishii.pdf}
\end{sourceblock}
双流体配方中出现的基本变量是

\begin{sourceblock}
\includegraphics[page=217,pagebox=mediabox,trim={45.00bp 171.28bp 44.37bp 392.73bp},clip,width=0.92\textwidth,max height=0.38\textheight]{Ishii.pdf}
\end{sourceblock}
因此，未知数总数为三十六，我们有： 六个平衡方程方程（9-152）、（9-153）和（9-154）； • 三个界面条件方程（9-155）、（9-156）和（9-158）； • 界面机械条件方程(9-157)； • 界面化学条件Eq.(9-163)； （相变） • 饱和条件Eq.(9-164)；

• T k T Eq.(9-169) 的两个本构方程；

• ki T Eq.(9-126) 的两个本构方程； • k q 的两个本构方程 Eq.(9-171)； • T k q 的两个本构方程 Eq.(9-172)； • 相变本构律Eq.(9-174)； （或类似于方程（9-179）的2 i i a q ′′的本构方程） • 阻力本构关系方程（9-176）； • 几何状态方程Eq.(9-178)；

• 1 i i a q ′′ 的本构方程Eq.(9-179)； • 表面积平衡方程Eq.(9-181)； • H m M Eq.(9-183) 的本构方程；

• ki p 的本构方程（对于大多数情况）Eq.(9-103)。这表明我们有三十六个方程。因此未知方程和方程的总数是相同的。因此，我们的描述在数学意义上是一致和完整的，尽管这不能保证模型解的唯一性，甚至不能保证解的存在性。这些应该通过解决各种简单的情况来检查，然后正如我们在第 2 章中讨论的那样，应将结果与实验数据进行比较，以验证和改进模型。这非常重要
需要注意的是，1 Γ 、1 i q '' 和 2 i q '' 与代表宏观能量跳跃条件的界面热能传递条件相关。

因此，可以使用 2 i q ′′ 的本构关系来代替相变本构关系，式（1） （9-174）。由于热通量比相变更容易建模，因此这是一种更实用的方法。

\section*{9.4 各种特殊情况}\phantomsection\addcontentsline{toc}{section}{9.4 各种特殊情况}
缩放参数上面已经给出了二流体模型的一般公式。在下面的分析中，我们从场方程中获得一些重要的标度参数。在进行详细研究之前，我们回顾一下，无量纲群可以从守恒方程、边界条件和本构律中获得。两个不同系统的相似性只能通过包括所有这些组来讨论。这将是
在下面的分析中，下标o表示参考参数选择为常数。特征长度为 o L，而时间常数为 o τ。对于大多数问题，它被视为 o L 与速度尺度的比率，但是，对于振荡流，它可以是振荡周期。下面，我们定义无量纲参数，其数量级被认为是1。

\begin{sourceblock}
\includegraphics[page=219,pagebox=mediabox,trim={45.00bp 50.14bp 44.37bp 329.06bp},clip,width=0.92\textwidth,max height=0.38\textheight]{Ishii.pdf}
\end{sourceblock}
\begin{sourceblock}
\includegraphics[page=220,pagebox=mediabox,trim={45.00bp 548.21bp 44.37bp 76.72bp},clip,width=0.92\textwidth,max height=0.38\textheight]{Ishii.pdf}
\end{sourceblock}
将这些新参数代入场方程中，我们得到以下结果。

方程（9-152）的无量纲连续性方程

\begin{sourceblock}
\includegraphics[page=220,pagebox=mediabox,trim={45.00bp 431.14bp 44.37bp 191.56bp},clip,width=0.92\textwidth,max height=0.38\textheight]{Ishii.pdf}
\end{sourceblock}
方程（9-153）的无量纲运动方程

\begin{sourceblock}
\includegraphics[page=220,pagebox=mediabox,trim={45.00bp 250.71bp 44.37bp 265.68bp},clip,width=0.92\textwidth,max height=0.38\textheight]{Ishii.pdf}
\end{sourceblock}
式（9-154）中的无量纲热能方程

\begin{sourceblock}
\includegraphics[page=220,pagebox=mediabox,trim={45.00bp 50.14bp 44.37bp 446.77bp},clip,width=0.92\textwidth,max height=0.38\textheight]{Ishii.pdf}
\end{sourceblock}
\begin{sourceblock}
\includegraphics[page=221,pagebox=mediabox,trim={45.00bp 540.45bp 44.37bp 57.59bp},clip,width=0.92\textwidth,max height=0.38\textheight]{Ishii.pdf}
\end{sourceblock}
这里我们引入了几个缩放参数，定义为

\begin{sourceblock}
\includegraphics[page=221,pagebox=mediabox,trim={45.00bp 145.06bp 44.37bp 153.73bp},clip,width=0.92\textwidth,max height=0.38\textheight]{Ishii.pdf}
\end{sourceblock}
前两个参数，斯特劳哈尔数和相变数，是运动学组。欧拉数、雷诺数、弗劳德数和阻力数是动态群，因为它们是动量方程中出现的各种力的比率。类似地，佩克莱特、埃克特和界面加热数是衡量各种能量传递机制的能量组。从定义（9-188）和非的形式
在我们讨论通过考虑尺度参数的限制条件可以获得的各种特殊情况之前，让我们研究界面转移、相间机械状态和相变（或化学状态）条件的无量纲形式。由式(9-155)和式(9-184)可得到无量纲界面传质条件，从而

\begin{sourceblock}
\includegraphics[page=222,pagebox=mediabox,trim={45.00bp 426.27bp 44.37bp 202.94bp},clip,width=0.92\textwidth,max height=0.38\textheight]{Ishii.pdf}
\end{sourceblock}
从方程（9-156）、（9-157）和（9-184）中，界面动量传递条件由下式给出

\begin{sourceblock}
\includegraphics[page=222,pagebox=mediabox,trim={45.00bp 340.23bp 44.37bp 288.98bp},clip,width=0.92\textwidth,max height=0.38\textheight]{Ishii.pdf}
\end{sourceblock}
和机械状态的条件

\begin{sourceblock}
\includegraphics[page=222,pagebox=mediabox,trim={45.00bp 282.31bp 44.37bp 358.30bp},clip,width=0.92\textwidth,max height=0.38\textheight]{Ishii.pdf}
\end{sourceblock}
类似地，相变条件(9-163)变为

\begin{sourceblock}
\includegraphics[page=222,pagebox=mediabox,trim={45.00bp 201.46bp 44.37bp 413.76bp},clip,width=0.92\textwidth,max height=0.38\textheight]{Ishii.pdf}
\end{sourceblock}
此外，能量传递条件 (9-158) 变为

\begin{sourceblock}
\includegraphics[page=222,pagebox=mediabox,trim={45.00bp 119.50bp 44.37bp 499.98bp},clip,width=0.92\textwidth,max height=0.38\textheight]{Ishii.pdf}
\end{sourceblock}
在这些方程中，我们引入了以下缩放参数

\begin{sourceblock}
\includegraphics[page=222,pagebox=mediabox,trim={45.00bp 50.14bp 44.37bp 576.76bp},clip,width=0.92\textwidth,max height=0.38\textheight]{Ishii.pdf}
\end{sourceblock}
\begin{sourceblock}
\includegraphics[page=223,pagebox=mediabox,trim={45.00bp 571.78bp 44.37bp 52.54bp},clip,width=0.92\textwidth,max height=0.38\textheight]{Ishii.pdf}
\end{sourceblock}
\begin{sourceblock}
\includegraphics[page=223,pagebox=mediabox,trim={45.00bp 511.78bp 44.37bp 112.54bp},clip,width=0.92\textwidth,max height=0.38\textheight]{Ishii.pdf}
\end{sourceblock}
通过结合表面张力和欧拉数，我们得到韦伯数：

\begin{sourceblock}
\includegraphics[page=223,pagebox=mediabox,trim={45.00bp 423.16bp 44.37bp 199.54bp},clip,width=0.92\textwidth,max height=0.38\textheight]{Ishii.pdf}
\end{sourceblock}
这表明我们在双流体配方中获得了两个韦伯数，因此使用表面张力数更方便。

转换后的热函比将相热函变化缩放为潜热。如果压力远低于临界压力，这个数字通常很小。然而，通过研究方程（9-192）可以获得最重要的简化。如果表面张力数或密度比非常小，我们有

\begin{sourceblock}
\includegraphics[page=223,pagebox=mediabox,trim={45.00bp 238.51bp 44.37bp 351.58bp},clip,width=0.92\textwidth,max height=0.38\textheight]{Ishii.pdf}
\end{sourceblock}
则由式(9-157)可得

\begin{sourceblock}
\includegraphics[page=223,pagebox=mediabox,trim={45.00bp 177.37bp 44.37bp 459.83bp},clip,width=0.92\textwidth,max height=0.38\textheight]{Ishii.pdf}
\end{sourceblock}
最简单的情况发生，如果表面张力数很小，那么

\begin{sourceblock}
\includegraphics[page=223,pagebox=mediabox,trim={45.00bp 112.38bp 44.37bp 524.77bp},clip,width=0.92\textwidth,max height=0.38\textheight]{Ishii.pdf}
\end{sourceblock}
这表明两相处于机械平衡。

现在我们研究一些重要的特殊情况。

无相变的流动 如果流动无相变，我们可以设置

\begin{sourceblock}
\includegraphics[page=224,pagebox=mediabox,trim={45.00bp 542.14bp 44.37bp 97.16bp},clip,width=0.92\textwidth,max height=0.38\textheight]{Ishii.pdf}
\end{sourceblock}
然后场方程中由相变数加权的所有项都从公式中删除。在这种情况下，通常更方便地将热能方程转换为 pk c 和 k Τ。因此，根据热量状态方程 (9-58)，我们有

\begin{sourceblock}
\includegraphics[page=224,pagebox=mediabox,trim={45.00bp 364.53bp 44.37bp 194.14bp},clip,width=0.92\textwidth,max height=0.38\textheight]{Ishii.pdf}
\end{sourceblock}
\begin{sourceblock}
\includegraphics[page=224,pagebox=mediabox,trim={114.63bp 423.73bp 173.60bp 199.63bp},clip,width=\textwidth,max height=0.72\textheight]{Ishii.pdf}
\end{sourceblock}
\begin{sourceblock}
\includegraphics[page=224,pagebox=mediabox,trim={45.00bp 324.68bp 44.37bp 309.53bp},clip,width=0.92\textwidth,max height=0.38\textheight]{Ishii.pdf}
\end{sourceblock}
\begin{sourceblock}
\includegraphics[page=224,pagebox=mediabox,trim={45.00bp 279.15bp 44.37bp 345.76bp},clip,width=0.92\textwidth,max height=0.38\textheight]{Ishii.pdf}
\end{sourceblock}
\begin{sourceblock}
\includegraphics[page=224,pagebox=mediabox,trim={45.00bp 201.70bp 44.37bp 403.00bp},clip,width=0.92\textwidth,max height=0.38\textheight]{Ishii.pdf}
\end{sourceblock}
and

\begin{sourceblock}
\includegraphics[page=224,pagebox=mediabox,trim={45.00bp 92.11bp 44.37bp 495.65bp},clip,width=0.92\textwidth,max height=0.38\textheight]{Ishii.pdf}
\end{sourceblock}
这里我们将式（9-202）代入式（9-154）。

现在我们可以将普朗特数定义为

\begin{sourceblock}
\includegraphics[page=225,pagebox=mediabox,trim={45.00bp 570.88bp 44.37bp 52.32bp},clip,width=0.92\textwidth,max height=0.38\textheight]{Ishii.pdf}
\end{sourceblock}
佩克莱特数应修改为

\begin{sourceblock}
\includegraphics[page=225,pagebox=mediabox,trim={45.00bp 496.84bp 44.37bp 126.28bp},clip,width=0.92\textwidth,max height=0.38\textheight]{Ishii.pdf}
\end{sourceblock}
此外，我们注意到方程（9-205）右侧的第二项简化为理想气体或不可压缩流体的简单形式

\begin{sourceblock}
\includegraphics[page=225,pagebox=mediabox,trim={45.00bp 386.76bp 44.37bp 213.63bp},clip,width=0.92\textwidth,max height=0.38\textheight]{Ishii.pdf}
\end{sourceblock}
如果埃克特数非常小，或者热传递主导能量交换，则可以使用具有实际重要性的能量方程的简单形式。然后我们有

\begin{sourceblock}
\includegraphics[page=225,pagebox=mediabox,trim={45.00bp 287.62bp 44.37bp 338.51bp},clip,width=0.92\textwidth,max height=0.38\textheight]{Ishii.pdf}
\end{sourceblock}
其中压缩性效应和粘性耗散项已从方程（9-205）中删除。此外，如果两相不可压缩且具有与温度无关的传输特性，则能量方程可以与连续性和动量方程解耦。

无相变的等温流 在这种情况下，整个能量方程可以从公式中删除。我们有

\begin{sourceblock}
\includegraphics[page=225,pagebox=mediabox,trim={45.00bp 138.70bp 44.37bp 502.06bp},clip,width=0.92\textwidth,max height=0.38\textheight]{Ishii.pdf}
\end{sourceblock}
因此这种流动称为正压流动。此外，如果压力变化或等温压缩性很小，则可以认为流动是不可压缩的。然后我们有

\begin{sourceblock}
\includegraphics[page=225,pagebox=mediabox,trim={45.00bp 53.20bp 44.37bp 588.04bp},clip,width=0.92\textwidth,max height=0.38\textheight]{Ishii.pdf}
\end{sourceblock}
\begin{sourceblock}
\includegraphics[page=226,pagebox=mediabox,trim={45.00bp 485.68bp 44.37bp 106.73bp},clip,width=0.92\textwidth,max height=0.38\textheight]{Ishii.pdf}
\end{sourceblock}
此外，如果系统具有固定的界面几何形状，则由于几何本构关系，公式可以简化为简单的形式。方程(9-178)和(9-181)，以及阻力定律方程(9-176)，可以很容易地得到。下面给出的关于分散流态的一些结果可以应用于这种情况。

流体颗粒的分散流 在下面的分析中，我们分别使用下标 c 和 d 表示连续相和分散相。因此我们设置：

阶段1→阶段c；连续相；阶段2→阶段d；分散相。为简单起见，我们假设分散相具有球形几何形状，在任何点和时间都具有相当均匀的直径。然后，根据几何状态方程（9-136）或方程（9-137），我们有

\begin{sourceblock}
\includegraphics[page=226,pagebox=mediabox,trim={45.00bp 215.44bp 44.37bp 408.70bp},clip,width=0.92\textwidth,max height=0.38\textheight]{Ishii.pdf}
\end{sourceblock}
where

\begin{sourceblock}
\includegraphics[page=226,pagebox=mediabox,trim={45.00bp 114.78bp 44.37bp 482.33bp},clip,width=0.92\textwidth,max height=0.38\textheight]{Ishii.pdf}
\end{sourceblock}
因此d R 可以被认为是流体粒子的平均半径。体积平衡方程可写成以下形式

\begin{sourceblock}
\includegraphics[page=227,pagebox=mediabox,trim={45.00bp 565.84bp 44.37bp 55.73bp},clip,width=0.92\textwidth,max height=0.38\textheight]{Ishii.pdf}
\end{sourceblock}
where

\begin{sourceblock}
\includegraphics[page=227,pagebox=mediabox,trim={45.00bp 489.64bp 44.37bp 134.50bp},clip,width=0.92\textwidth,max height=0.38\textheight]{Ishii.pdf}
\end{sourceblock}
是每个流体颗粒可用的自由体积，其中 d n 是颗粒数密度。右侧代表由于聚结而产生的体积源和由于颗粒崩解而产生的体积源。

我们在没有源项或汇项的简单情况下演示了方程（9-138）的推导。因此，通过考虑平均属性的流体粒子，我们可以近似

\begin{sourceblock}
\includegraphics[page=227,pagebox=mediabox,trim={45.00bp 345.70bp 44.37bp 275.58bp},clip,width=0.92\textwidth,max height=0.38\textheight]{Ishii.pdf}
\end{sourceblock}
\begin{sourceblock}
\includegraphics[page=227,pagebox=mediabox,trim={45.00bp 256.83bp 44.37bp 339.45bp},clip,width=0.92\textwidth,max height=0.38\textheight]{Ishii.pdf}
\end{sourceblock}
我们在这里注意到，如果颗粒直径变化很大，那么系数 i C 就不是常数。在这种情况下，由于 iC 的变化，我们应该有一个附加项，因为流体颗粒的平均表面积和体积与根据平均直径计算的并不完全相同。

现在让我们研究流体粒子系统中的阻力本构方程。 Hadamard 将著名的斯托克斯定律扩展到无限纳维-斯托克斯流体中球形流体粒子的蠕动运动（Brodkey，1967；Soo，1967）。因此，作用在流体粒子上的总力由下式给出

\begin{sourceblock}
\includegraphics[page=227,pagebox=mediabox,trim={45.00bp 57.65bp 44.37bp 558.74bp},clip,width=0.92\textwidth,max height=0.38\textheight]{Ishii.pdf}
\end{sourceblock}
\begin{sourceblock}
\includegraphics[page=228,pagebox=mediabox,trim={45.00bp 532.64bp 44.37bp 58.85bp},clip,width=0.92\textwidth,max height=0.38\textheight]{Ishii.pdf}
\end{sourceblock}
粒子雷诺数为

\begin{sourceblock}
\includegraphics[page=228,pagebox=mediabox,trim={45.00bp 456.57bp 44.37bp 164.69bp},clip,width=0.92\textwidth,max height=0.38\textheight]{Ishii.pdf}
\end{sourceblock}
这里c v ∞ 和d v 是原状流速和粒子速度。由上可知，我们有

\begin{sourceblock}
\includegraphics[page=228,pagebox=mediabox,trim={45.00bp 362.48bp 44.37bp 253.94bp},clip,width=0.92\textwidth,max height=0.38\textheight]{Ishii.pdf}
\end{sourceblock}
Hadamard 给出的阻力定律在雷诺数约为 1 时效果良好。对于更高的雷诺数，我们有 Levich (1962) 和 Chao (1962) 给出的结果：

\begin{sourceblock}
\includegraphics[page=228,pagebox=mediabox,trim={45.00bp 265.72bp 44.37bp 360.77bp},clip,width=0.92\textwidth,max height=0.38\textheight]{Ishii.pdf}
\end{sourceblock}
\begin{sourceblock}
\includegraphics[page=228,pagebox=mediabox,trim={45.00bp 187.23bp 44.37bp 418.55bp},clip,width=0.92\textwidth,max height=0.38\textheight]{Ishii.pdf}
\end{sourceblock}
分别。我们还注意到 Clift 等人所做的审查工作。 （1978）在这些联系中。在更高的雷诺数下，
0 44。牛顿给出的 DC ∞ ≈ 可用于液滴。然而，对于气泡来说，界面变形会导致椭圆形或球冠气泡。

通过结合这些结果，我们可以将阻力系数设置为雷诺数 ( ) d Re 和两个粘度之比的函数，从而

\begin{sourceblock}
\includegraphics[page=229,pagebox=mediabox,trim={45.00bp 565.84bp 44.37bp 55.73bp},clip,width=0.92\textwidth,max height=0.38\textheight]{Ishii.pdf}
\end{sourceblock}
上述关系总结了无限介质中单个流体粒子的理想情况。

在一般情况下，我们假设界面阻力可以由具有方程（9-117）和（9-119）形式的本构关系给出。对于受式(9-213)、(9-214)和(9-218)限制的分散流，我们可以通过引入阻力系数来简化一般阻力本构关系
直流定义为

\begin{sourceblock}
\includegraphics[page=229,pagebox=mediabox,trim={45.00bp 401.71bp 44.37bp 213.30bp},clip,width=0.92\textwidth,max height=0.38\textheight]{Ishii.pdf}
\end{sourceblock}
鉴于式（9-221）和（9-226），我们重新定义合适的粒子雷诺数为

\begin{sourceblock}
\includegraphics[page=229,pagebox=mediabox,trim={45.00bp 309.88bp 44.37bp 311.26bp},clip,width=0.92\textwidth,max height=0.38\textheight]{Ishii.pdf}
\end{sourceblock}
如果我们使用粘度比作为无量纲群，则分散相的雷诺数是多余的。

因此，根据方程（9-119）和（9-227），我们假设阻力定律可以由下式给出：

\begin{sourceblock}
\includegraphics[page=229,pagebox=mediabox,trim={45.00bp 194.74bp 44.37bp 426.85bp},clip,width=0.92\textwidth,max height=0.38\textheight]{Ishii.pdf}
\end{sourceblock}
其中* f 是考虑其他粒子的影响和相位变化的校正因子。可以说，如果i pch N 很大，则由于流体粒子的快速膨胀或塌陷，方程(9-228)的线性校正就无法应用。第 12 章详细讨论了多粒子系统中的阻力。

\chapter{界面面积输运}
界面输运项与界面面积和局部传递机制密切相关，例如界面附近的湍流程度和驱动势。基本上，质量、动量和能量的界面传输与界面面积浓度和驱动力成正比。该面积浓度定义为混合物每单位体积的界面面积，表征了运动学效应；因此，它必然与两相流的结构有关。相间传输的驱动力表征了局部传输机制，必须单独对它们进行建模。

由于界面传输速率可以被视为界面通量和可用界面面积的乘积，因此界面面积浓度的建模至关重要。在两相流分析中，空隙率和界面面积浓度代表两个基本的一阶几何参数。因此，它们与两相流态密切相关。然而，两相流态的概念很难在局部点进行数学量化，因为它通常是在接近系统尺度的尺度上定义的。这可能表明直接通过输运方程对界面面积浓度的变化进行建模，即界面区域输运方程。与使用流态转变标准和界面面积浓度的态相关本构关系的传统方法相比，这是一种更好的方法。对于多相流的三维公式尤其如此。

本章给出了界面区域输运方程的详细推导和必要的本构关系，建立了双流体模型中界面面积浓度的动态闭合关系。考虑到小球形和大帽气泡传输机制的显著差异，推导了两组传输方程。

\section*{10.1 三维界面面积输运方程}\phantomsection\addcontentsline{toc}{section}{10.1 三维界面面积输运方程}
玻尔兹曼输运方程通过粒子分布函数的积分微分方程描述粒子输运。由于流体粒子的界面面积与粒子数密切相关，因此可以根据玻尔兹曼输运方程建立界面面积输运方程（Kocamustafaogullari和Ishii，1995；Ishii和Kim，2004）。

考虑连续介质中的流体粒子系统，其中由于粒子相互作用（例如聚结和分解）而存在流体粒子的源和汇。令 ( ) f V, , ,t x v 为每单位混合物和气泡体积的颗粒数密度分布函数。假设它是连续的，并指定在给定时间 t 、空间范围 δ x 中以粒子速度 v 移动的流体粒子的可能数密度，其体积中心位于 x，粒子体积介于 V 和 V V δ + 之间。假设t～t t δ + 时间间隔内颗粒速度变化较小，则单位混合物和气泡体积的颗粒数密度分布函数可简化为( ) f V, ,t x。对于给定颗粒尺寸的均匀颗粒速度的假设对于大多数两相流来说是实用的。然而，对于中子输运来说，中子的能量跨越多个数量级，是输运理论的本质。因此，速度依赖性不可忽视。然后，我们可以编写两相流系统

\begin{sourceblock}
\includegraphics[page=231,pagebox=mediabox,trim={45.00bp 195.98bp 44.37bp 367.59bp},clip,width=0.92\textwidth,max height=0.38\textheight]{Ishii.pdf}
\end{sourceblock}
其中δμ是μ空间中的体积元素。在方程的右侧，j S 和 ph S 是由于第 j 个粒子相互作用（例如分解或聚结以及相变引起的）每单位混合物体积的粒子源和汇速率。将泰勒级数中的方程（10-1）左侧第一项在 t δ 中展开并除以 t δμδ，方程： (10-1) 简化为

\begin{sourceblock}
\includegraphics[page=231,pagebox=mediabox,trim={45.00bp 61.90bp 44.37bp 562.96bp},clip,width=0.92\textwidth,max height=0.38\textheight]{Ishii.pdf}
\end{sourceblock}
\subsection*{10.1.1 数输运方程}\phantomsection\addcontentsline{toc}{subsection}{10.1.1 数输运方程}
在两相流应用中，式（10-2）给出的颗粒输运方程过于详细而无法在实践中使用，因此需要更宏观的公式。这可以通过对从最小 V 到最大 V 的所有尺寸粒子的体积积分方程（10-2）并应用莱布尼茨积分规则来完成。我们得到粒子数输运方程为

\begin{sourceblock}
\includegraphics[page=232,pagebox=mediabox,trim={45.00bp 375.21bp 44.37bp 249.72bp},clip,width=0.92\textwidth,max height=0.38\textheight]{Ishii.pdf}
\end{sourceblock}
其中，体积 min V 和 max V 的气泡的分布函数假设近似为零。这里，方程的左侧表示总粒子数密度及其对流的时间变化率。右侧的两项分别表示由于粒子相互作用（例如粒子解体或聚结）引起的数源和汇速率以及由于相变引起的数源速率。

在式(10-3)中，单位混合物体积中各种尺寸颗粒的总数以及数源和汇率分别定义为

\begin{sourceblock}
\includegraphics[page=232,pagebox=mediabox,trim={45.00bp 203.19bp 44.37bp 426.21bp},clip,width=0.92\textwidth,max height=0.38\textheight]{Ishii.pdf}
\end{sourceblock}
and

\begin{sourceblock}
\includegraphics[page=232,pagebox=mediabox,trim={45.00bp 135.28bp 44.37bp 494.23bp},clip,width=0.92\textwidth,max height=0.38\textheight]{Ishii.pdf}
\end{sourceblock}
此外，pm v 是按粒子数加权的平均局部粒子速度，定义为

\begin{sourceblock}
\includegraphics[page=233,pagebox=mediabox,trim={45.00bp 547.30bp 44.37bp 58.21bp},clip,width=0.92\textwidth,max height=0.38\textheight]{Ishii.pdf}
\end{sourceblock}
\subsection*{10.1.2 体积输运方程}\phantomsection\addcontentsline{toc}{subsection}{10.1.2 体积输运方程}
颗粒体积（或空隙率）输运方程可以通过将方程（10-2）乘以颗粒体积V并在所有尺寸颗粒的体积上积分来获得。然后，考虑到所关注的两相流由连续液体介质中的分散气泡组成，空隙率传输方程由下式给出

\begin{sourceblock}
\includegraphics[page=233,pagebox=mediabox,trim={45.00bp 354.60bp 44.37bp 228.25bp},clip,width=0.92\textwidth,max height=0.38\textheight]{Ishii.pdf}
\end{sourceblock}
其中 V\$ 表示体积 V 的时间导数。这里，分散相（或气相）的空隙率和平均体积中心速度分别定义为

\begin{sourceblock}
\includegraphics[page=233,pagebox=mediabox,trim={45.00bp 260.67bp 44.37bp 368.83bp},clip,width=0.92\textwidth,max height=0.38\textheight]{Ishii.pdf}
\end{sourceblock}
and

\begin{sourceblock}
\includegraphics[page=233,pagebox=mediabox,trim={45.00bp 165.64bp 44.37bp 439.84bp},clip,width=0.92\textwidth,max height=0.38\textheight]{Ishii.pdf}
\end{sourceblock}
在方程（10-7）中，左侧第三项归因于由于沿流场的压力变化而引起的颗粒体积的变化（通过膨胀或收缩）。为了更好地表示这个术语，我们假设

\begin{sourceblock}
\includegraphics[page=233,pagebox=mediabox,trim={45.00bp 65.29bp 44.37bp 557.24bp},clip,width=0.92\textwidth,max height=0.38\textheight]{Ishii.pdf}
\end{sourceblock}
\begin{sourceblock}
\includegraphics[page=234,pagebox=mediabox,trim={45.00bp 446.55bp 44.37bp 172.16bp},clip,width=0.92\textwidth,max height=0.38\textheight]{Ishii.pdf}
\end{sourceblock}
此外，由于蒸发过程的传质由下式给出

\begin{sourceblock}
\includegraphics[page=234,pagebox=mediabox,trim={45.00bp 368.56bp 44.37bp 253.60bp},clip,width=0.92\textwidth,max height=0.38\textheight]{Ishii.pdf}
\end{sourceblock}
其中 g Γ 是单位混合物体积质量的总变化率，
ph η 是每单位混合物体积中成核源产生的体积速率，定义为

\begin{sourceblock}
\includegraphics[page=234,pagebox=mediabox,trim={45.00bp 274.17bp 44.37bp 358.25bp},clip,width=0.92\textwidth,max height=0.38\textheight]{Ishii.pdf}
\end{sourceblock}
体积源可以写为

\begin{sourceblock}
\includegraphics[page=234,pagebox=mediabox,trim={45.00bp 155.60bp 44.37bp 423.70bp},clip,width=0.92\textwidth,max height=0.38\textheight]{Ishii.pdf}
\end{sourceblock}
因此，将式（10-14）与式（10-11）结合起来，代入式（10-7），得到空隙率输运方程的最终形式为

\begin{sourceblock}
\includegraphics[page=235,pagebox=mediabox,trim={45.00bp 523.73bp 44.37bp 55.63bp},clip,width=0.92\textwidth,max height=0.38\textheight]{Ishii.pdf}
\end{sourceblock}
其中方程左侧的前两项代表 g α 的变化和对流的时间速率，其余项分别代表由于体积变化、粒子相互作用和相变而导致的 g α 变化率。

通过重新排列方程（10-15），有趣的是我们有

\begin{sourceblock}
\includegraphics[page=235,pagebox=mediabox,trim={45.00bp 393.70bp 44.37bp 226.74bp},clip,width=0.92\textwidth,max height=0.38\textheight]{Ishii.pdf}
\end{sourceblock}
在这里，值得注意的是，质量守恒需要

\begin{sourceblock}
\includegraphics[page=235,pagebox=mediabox,trim={45.00bp 321.73bp 44.37bp 302.25bp},clip,width=0.92\textwidth,max height=0.38\textheight]{Ishii.pdf}
\end{sourceblock}
因此，从式（10-16）和（10-17）我们得到恒等式

\begin{sourceblock}
\includegraphics[page=235,pagebox=mediabox,trim={45.00bp 253.55bp 44.37bp 375.82bp},clip,width=0.92\textwidth,max height=0.38\textheight]{Ishii.pdf}
\end{sourceblock}
方程（10-18）同时满足体积和质量守恒。

\subsection*{10.1.3 界面面积输运方程}\phantomsection\addcontentsline{toc}{subsection}{10.1.3 界面面积输运方程}
界面面积浓度的传输方程可以通过先前公式中应用的类似方法获得。因此，将方程（10-2）乘以体积为 V 的颗粒的表面积，
( ) i A V，（与坐标系无关）并将其积分到所有粒子的体积上，我们得到

\begin{sourceblock}
\includegraphics[page=236,pagebox=mediabox,trim={45.00bp 525.81bp 44.37bp 55.83bp},clip,width=0.92\textwidth,max height=0.38\textheight]{Ishii.pdf}
\end{sourceblock}
其中体积在 min V 和 min V 之间的所有流体颗粒的平均值 i a
max V 和界面速度分别由下式给出

\begin{sourceblock}
\includegraphics[page=236,pagebox=mediabox,trim={45.00bp 444.40bp 44.37bp 185.02bp},clip,width=0.92\textwidth,max height=0.38\textheight]{Ishii.pdf}
\end{sourceblock}
and

\begin{sourceblock}
\includegraphics[page=236,pagebox=mediabox,trim={45.00bp 349.18bp 44.37bp 256.25bp},clip,width=0.92\textwidth,max height=0.38\textheight]{Ishii.pdf}
\end{sourceblock}
现在，考虑到式（10-19）左边的第三项，我们将表面积为 i A 、体积为 V 的流体颗粒的体积当量直径 e D 和表面当量直径 s D 定义为

\begin{sourceblock}
\includegraphics[page=236,pagebox=mediabox,trim={45.00bp 254.58bp 44.37bp 359.44bp},clip,width=0.92\textwidth,max height=0.38\textheight]{Ishii.pdf}
\end{sourceblock}
因此，将它们与方程（10-19）结合并回忆由方程（10-14）给出的体积源，可以得到界面面积输运方程：

\begin{sourceblock}
\includegraphics[page=236,pagebox=mediabox,trim={45.00bp 111.02bp 44.37bp 467.65bp},clip,width=0.92\textwidth,max height=0.38\textheight]{Ishii.pdf}
\end{sourceblock}
其中左侧第三项表示由于颗粒体积变化而引起的界面面积浓度的变化。在推导式(10-23)时，为了简化方程，假设比率( ) s e D D 为常数。虽然这个近似值可能不合适

\begin{sourceblock}
\includegraphics[page=237,pagebox=mediabox,trim={50.00bp 443.40bp 44.37bp 55.00bp},clip,width=\textwidth,max height=0.72\textheight]{Ishii.pdf}
\end{sourceblock}
Δ A i (Ishii 和 Kim, 2004)

\begin{center}\small 图10-1。流体颗粒聚结和分解过程的图示\end{center}
对于扭曲或块状气泡，它是球形和帽状气泡的良好近似。本质上，该直径比是形状因子，并且对于相似的颗粒形状，该因子可以被认为是常数。

为了封闭方程组，方程（10-23）的右侧代表界面面积浓度的源和汇速率，必须由本构关系指定。鉴于此，我们定义

\begin{sourceblock}
\includegraphics[page=237,pagebox=mediabox,trim={45.00bp 227.19bp 44.37bp 372.39bp},clip,width=0.92\textwidth,max height=0.38\textheight]{Ishii.pdf}
\end{sourceblock}
and

\begin{sourceblock}
\includegraphics[page=237,pagebox=mediabox,trim={45.00bp 126.29bp 44.37bp 473.38bp},clip,width=0.92\textwidth,max height=0.38\textheight]{Ishii.pdf}
\end{sourceblock}
此外，注意到 j φ 可以用流体粒子在一定的粒子相互作用过程后表面积的变化来表示，我们可以写成

\begin{sourceblock}
\includegraphics[page=238,pagebox=mediabox,trim={45.00bp 586.18bp 44.37bp 55.03bp},clip,width=0.92\textwidth,max height=0.38\textheight]{Ishii.pdf}
\end{sourceblock}
其中 j R 可以为每个相互作用机制进行机械建模，i A Δ 取决于给定的相互作用机制；例如分解或聚结过程。

为了指定 i A Δ，请考虑球形颗粒的聚结和破碎过程，如图 10-1 所示。这里，假设给定的过程是相同尺寸的粒子之间的二元过程。然后，由于粒子的总体积应该守恒，我们可以写成

\begin{sourceblock}
\includegraphics[page=238,pagebox=mediabox,trim={45.00bp 436.21bp 44.37bp 206.60bp},clip,width=0.92\textwidth,max height=0.38\textheight]{Ishii.pdf}
\end{sourceblock}
其中下标1和2分别表示体积较小和较大的颗粒。因此，假设相互作用过程是一个二元过程，对于近球形颗粒，经过一次相互作用过程后表面积的变化可以得到为

\begin{sourceblock}
\includegraphics[page=238,pagebox=mediabox,trim={45.00bp 342.03bp 44.37bp 299.06bp},clip,width=0.92\textwidth,max height=0.38\textheight]{Ishii.pdf}
\end{sourceblock}
and

\begin{sourceblock}
\includegraphics[page=238,pagebox=mediabox,trim={45.00bp 287.08bp 44.37bp 354.13bp},clip,width=0.92\textwidth,max height=0.38\textheight]{Ishii.pdf}
\end{sourceblock}
其中减号和加号分别用于表示一次相互作用过程后表面积的减少和增加。此外，回顾方程（10-4）给出的定义，粒子数密度 n 可以通过 i a 和 g α 指定：

\begin{sourceblock}
\includegraphics[page=238,pagebox=mediabox,trim={45.00bp 193.18bp 44.37bp 423.51bp},clip,width=0.92\textwidth,max height=0.38\textheight]{Ishii.pdf}
\end{sourceblock}
这样

\begin{sourceblock}
\includegraphics[page=238,pagebox=mediabox,trim={45.00bp 117.76bp 44.37bp 509.64bp},clip,width=0.92\textwidth,max height=0.38\textheight]{Ishii.pdf}
\end{sourceblock}
形状因子定义为

\begin{sourceblock}
\includegraphics[page=239,pagebox=mediabox,trim={45.00bp 563.74bp 44.37bp 57.76bp},clip,width=0.92\textwidth,max height=0.38\textheight]{Ishii.pdf}
\end{sourceblock}
其中气泡索特平均直径由下式给出

\begin{sourceblock}
\includegraphics[page=239,pagebox=mediabox,trim={45.00bp 488.80bp 44.37bp 134.32bp},clip,width=0.92\textwidth,max height=0.38\textheight]{Ishii.pdf}
\end{sourceblock}
因此，将这些与方程（10-26）结合起来，表面源汇率 j φ 可以由下式给出：

\begin{sourceblock}
\includegraphics[page=239,pagebox=mediabox,trim={45.00bp 397.70bp 44.37bp 223.80bp},clip,width=0.92\textwidth,max height=0.38\textheight]{Ishii.pdf}
\end{sourceblock}
\begin{sourceblock}
\includegraphics[page=239,pagebox=mediabox,trim={45.00bp 339.11bp 44.37bp 275.50bp},clip,width=0.92\textwidth,max height=0.38\textheight]{Ishii.pdf}
\end{sourceblock}
其中 bc D 是临界气泡尺寸。这应根据给定的成核过程来确定；即，本体沸腾或冷凝过程的临界空腔尺寸，以及壁成核的气泡离开尺寸。对于大多数两相流，壁成核是主要机制。

将上面给出的本构关系代入式(10-23)，得到界面面积输运方程为

\begin{sourceblock}
\includegraphics[page=239,pagebox=mediabox,trim={45.00bp 141.27bp 44.37bp 437.27bp},clip,width=0.92\textwidth,max height=0.38\textheight]{Ishii.pdf}
\end{sourceblock}
其中左侧代表界面面积浓度的变化率和对流率。右侧的每一项分别表示由于压力变化、各种颗粒相互作用和相变引起的颗粒体积变化而引起的界面面积浓度的变化率。由式（10-36）可知，
j R 应该根据给定的粒子独立建模

\section*{10.2 一组界面面积输运方程}\phantomsection\addcontentsline{toc}{section}{10.2 一组界面面积输运方程}
当所关注的流体颗粒的输运现象在给定的两相流系统中没有显著变化，并且颗粒相互作用后形状保持相似时，它们的特征输运现象是相似的，并且可以用一个输运方程来描述。然而，当不同形状和尺寸的流体颗粒同时存在时，它们的传输机制可能会显著不同。在这种情况下，可能需要采用多个传输方程来描述流体粒子传输。

鉴于此，我们首先考虑连续液体介质中分散气泡的两相流系统（即气泡流），其中所有存在的气泡都可以归类为一组。在这种流动条件下，假设气泡的形状是球形的，并且它们的传输现象受到类似的特征阻力。因此，考虑到单群传输中的球形形状，方程（10-32）中定义的ψ可以近似为

\begin{sourceblock}
\includegraphics[page=240,pagebox=mediabox,trim={45.00bp 270.21bp 44.37bp 357.75bp},clip,width=0.92\textwidth,max height=0.38\textheight]{Ishii.pdf}
\end{sourceblock}
因为气泡索特平均直径约等于体积当量直径。此外，注意到由于成核而导致的临界气泡尺寸比平均气泡索特平均直径小得多，我们可以假设

\begin{sourceblock}
\includegraphics[page=240,pagebox=mediabox,trim={45.00bp 157.54bp 44.37bp 464.06bp},clip,width=0.92\textwidth,max height=0.38\textheight]{Ishii.pdf}
\end{sourceblock}
\begin{sourceblock}
\includegraphics[page=240,pagebox=mediabox,trim={45.00bp 87.49bp 44.37bp 515.68bp},clip,width=0.92\textwidth,max height=0.38\textheight]{Ishii.pdf}
\end{sourceblock}
分散气泡的界面面积输运方程，或单组界面面积输运方程，由下式给出

\begin{sourceblock}
\includegraphics[page=241,pagebox=mediabox,trim={45.00bp 523.31bp 44.37bp 55.27bp},clip,width=0.92\textwidth,max height=0.38\textheight]{Ishii.pdf}
\end{sourceblock}
左侧表示界面面积浓度的总变化率，而右侧分别表示由于颗粒体积变化、各种颗粒相互作用和相变而导致的界面面积浓度变化率。值得注意的是，由于偏离尺寸小于索特平均直径，因此可以忽略 ph η 的影响。

\section*{10.3 两组界面面积输运方程}\phantomsection\addcontentsline{toc}{section}{10.3 两组界面面积输运方程}
在气液两相流系统中，根据给定的流态，存在各种气泡形状和尺寸。因此，为了建立描述各种两相流状态下气泡传输的界面面积输运方程，模型必须考虑不同类型气泡传输特性的差异。由于阻力，气泡形状和尺寸的这些变化导致其传输机制存在显著差异。此外，这种流动条件下的气泡相互作用机制与单组传输中的气泡相互作用机制可能有很大不同。

在大多数两相流条件下，气泡可分为五种类型：球形；扭曲的；帽;泰勒；和搅动湍流的泡沫。然而，考虑到它们的传输特性，它们可以分为两大类，即第一组包括球形气泡和扭曲气泡，而第二组包括帽形气泡、泰勒气泡和搅动湍流气泡。因此，在本分析中，给出了采用两个传输方程的方法来描述大范围的两相流条件下的气泡传输。即，第 1 组输运方程描述了球形和扭曲气泡的输运，第 2 组输运方程描述了帽型、泰勒型和搅动湍流气泡的输运。

在一组公式中，输运方程通过所有尺寸颗粒体积上的积分过程进行平均，因为假设颗粒的形状及其输运现象在给定的颗粒体积范围内相似。然而，在两组公式中，每个传输方程的积分极限应受气泡体积的限制，从而确定气泡组。鉴于此，我们将c V定义为临界气泡体积

\begin{sourceblock}
\includegraphics[page=242,pagebox=mediabox,trim={45.00bp 532.22bp 44.37bp 92.72bp},clip,width=0.92\textwidth,max height=0.38\textheight]{Ishii.pdf}
\end{sourceblock}
在此期间，气泡变成盖状，并且由于尾流区域较大，阻力效应开始偏离较小气泡的阻力效应。因此，第 1 组气泡存在于 min V 至 c V 的范围内，而第 2 组气泡存在于 c V 至 max V 的范围内。

\subsection*{10.3.1 二群粒子数输运方程}\phantomsection\addcontentsline{toc}{subsection}{10.3.1 二群粒子数输运方程}
通过在以 c V 为界的不同积分极限范围内对式（10-2）进行积分，可以容易地得到两组粒子数输运方程；即，对于第 1 组，从 min V 到 c V 以及从 c V 到
第 2 组的 max V。在两组公式中，与一组公式中一样，( ) f V, ,t x 描述每单位混合物和气泡体积的颗粒数密度分布函数。假设它是连续的，指定在给定时间 t 、空间范围 δ x 中以速度 v 移动的流体粒子的可能数密度，其体积中心位于 x，粒子体积在 V 和 V dV + 之间。然后，组 1 和组 2 的数传输方程由下式给出

\begin{sourceblock}
\includegraphics[page=242,pagebox=mediabox,trim={45.00bp 247.55bp 44.37bp 373.98bp},clip,width=0.92\textwidth,max height=0.38\textheight]{Ishii.pdf}
\end{sourceblock}
and

\begin{sourceblock}
\includegraphics[page=242,pagebox=mediabox,trim={45.00bp 170.56bp 44.37bp 450.94bp},clip,width=0.92\textwidth,max height=0.38\textheight]{Ishii.pdf}
\end{sourceblock}
其中方程中的下标 1 和 2 分别表示组 1 和组 2，1 pm v 和 2 pm v 是按每个气泡组的粒子数加权的平均局部粒子速度，因此它们的定义为

\begin{sourceblock}
\includegraphics[page=243,pagebox=mediabox,trim={45.00bp 490.00bp 44.37bp 58.69bp},clip,width=0.92\textwidth,max height=0.38\textheight]{Ishii.pdf}
\end{sourceblock}
在式(10-42)和式(10-43)中，式左边表示各气泡群的流体粒子数的时间变化率和对流率。右侧的每一项代表每个气泡组的颗粒体积变化、颗粒相互作用和相变通过组间转移引起的颗粒数量变化率。在这里，有趣的是，在两组公式中，有一些术语解释了由颗粒体积变化引起的组间转移，而这些项在一组公式中没有出现。这是因为当存在两组气泡时，由于气泡分布函数的变化，一组中颗粒体积的变化可能成为另一组中颗粒体积的变化。当两个方程相加得到总流体粒子数传递方程时，这些组间传递项消失。

\subsection*{10.3.2 两组空隙率输运方程}\phantomsection\addcontentsline{toc}{subsection}{10.3.2 两组空隙率输运方程}
以类似的方式可以得到两组空隙率输运方程。将方程（10-2）乘以粒子体积 V，并将其在每个组的指定限制上积分，我们得到

\begin{sourceblock}
\includegraphics[page=243,pagebox=mediabox,trim={45.00bp 143.01bp 44.37bp 442.01bp},clip,width=0.92\textwidth,max height=0.38\textheight]{Ishii.pdf}
\end{sourceblock}
and

\begin{sourceblock}
\includegraphics[page=244,pagebox=mediabox,trim={45.00bp 529.62bp 44.37bp 55.44bp},clip,width=0.92\textwidth,max height=0.38\textheight]{Ishii.pdf}
\end{sourceblock}
对于组 1 和组 2，方程左侧的第三项分别表示由于颗粒体积变化而导致的空隙率变化率。它们由下式给出：

\begin{sourceblock}
\includegraphics[page=244,pagebox=mediabox,trim={45.00bp 425.44bp 44.37bp 194.76bp},clip,width=0.92\textwidth,max height=0.38\textheight]{Ishii.pdf}
\end{sourceblock}
and

\begin{sourceblock}
\includegraphics[page=244,pagebox=mediabox,trim={45.00bp 348.52bp 44.37bp 271.50bp},clip,width=0.92\textwidth,max height=0.38\textheight]{Ishii.pdf}
\end{sourceblock}
分别适用于第 1 组和第 2 组，其中 c f 是具有临界体积 c V 或 ( ) c f V 的气泡的分布函数。这里，体积源 ( ) V V 可以用总传质速率来表示，如前面在一组公式中所示。然而，由于两组气泡的存在及其相互作用，必须考虑两组之间的传质速率。因此，通过将下标 ij 表示为从组 i 到组 j 的组间传输，每组的体积源由下式给出

\begin{sourceblock}
\includegraphics[page=244,pagebox=mediabox,trim={45.00bp 139.52bp 44.37bp 439.52bp},clip,width=0.92\textwidth,max height=0.38\textheight]{Ishii.pdf}
\end{sourceblock}
and

\begin{sourceblock}
\includegraphics[page=245,pagebox=mediabox,trim={45.00bp 523.52bp 44.37bp 55.52bp},clip,width=0.92\textwidth,max height=0.38\textheight]{Ishii.pdf}
\end{sourceblock}
分别为组 1 和组 2，其中 12 m Δ 表示从组 1 到组 2 的组间传质速率。组 12 m Δ 之间传质的本构关系将在后面讨论。因此，方程（10-49）和（10-50）需要以下恒等式

\begin{sourceblock}
\includegraphics[page=245,pagebox=mediabox,trim={45.00bp 397.30bp 44.37bp 211.18bp},clip,width=0.92\textwidth,max height=0.38\textheight]{Ishii.pdf}
\end{sourceblock}
and

\begin{sourceblock}
\includegraphics[page=245,pagebox=mediabox,trim={45.00bp 306.28bp 44.37bp 302.20bp},clip,width=0.92\textwidth,max height=0.38\textheight]{Ishii.pdf}
\end{sourceblock}
此外，通过将两个方程相加，我们得到气相的连续性方程为

\begin{sourceblock}
\includegraphics[page=245,pagebox=mediabox,trim={45.00bp 217.59bp 44.37bp 406.49bp},clip,width=0.92\textwidth,max height=0.38\textheight]{Ishii.pdf}
\end{sourceblock}
具有以下本构关系

\begin{sourceblock}
\includegraphics[page=245,pagebox=mediabox,trim={45.00bp 164.68bp 44.37bp 476.64bp},clip,width=0.92\textwidth,max height=0.38\textheight]{Ishii.pdf}
\end{sourceblock}
\begin{sourceblock}
\includegraphics[page=245,pagebox=mediabox,trim={45.00bp 120.70bp 44.37bp 520.54bp},clip,width=0.92\textwidth,max height=0.38\textheight]{Ishii.pdf}
\end{sourceblock}
and

\begin{sourceblock}
\includegraphics[page=245,pagebox=mediabox,trim={45.00bp 50.14bp 44.37bp 575.72bp},clip,width=0.92\textwidth,max height=0.38\textheight]{Ishii.pdf}
\end{sourceblock}
\begin{sourceblock}
\includegraphics[page=246,pagebox=mediabox,trim={45.00bp 402.65bp 44.37bp 146.15bp},clip,width=0.92\textwidth,max height=0.38\textheight]{Ishii.pdf}
\end{sourceblock}
and

\begin{sourceblock}
\includegraphics[page=246,pagebox=mediabox,trim={45.00bp 251.39bp 44.37bp 296.62bp},clip,width=0.92\textwidth,max height=0.38\textheight]{Ishii.pdf}
\end{sourceblock}
分别为第 1 组和第 2 组。这里，sc D 是表面积和体积为 ic A 和 c V 的群边界的临界气泡尺寸。此外，χ是解释组间转移贡献的系数，这将在下一节中详细讨论。

在方程(10-57)和(10-58)中，方程的左侧表示各组空隙率的时间变化率和对流率，右侧表示由于体积变化引起的空隙率的变化率。这包括组间转移和各种粒子相互作用。此外，由于方程的左侧对应于每个气泡组的连续性方程，因此需要以下恒等式；

\begin{sourceblock}
\includegraphics[page=247,pagebox=mediabox,trim={45.00bp 547.57bp 44.37bp 57.46bp},clip,width=0.92\textwidth,max height=0.38\textheight]{Ishii.pdf}
\end{sourceblock}
and

\begin{sourceblock}
\includegraphics[page=247,pagebox=mediabox,trim={45.00bp 454.57bp 44.37bp 150.46bp},clip,width=0.92\textwidth,max height=0.38\textheight]{Ishii.pdf}
\end{sourceblock}
分别为组 1 和组 2，这表明气泡体积守恒。这里，方程（10-59）和（10-60）中的第一项表示气泡群边界处的群间传递，第二项表示由于给定气泡群的各种流体粒子相互作用而产生的源和汇。

\subsection*{10.3.3 两组界面面积输运方程}\phantomsection\addcontentsline{toc}{subsection}{10.3.3 两组界面面积输运方程}
类似地，如在一组界面面积输运方程公式中一样，将方程（10-2）乘以体积为 V 的颗粒的表面积，该表面积与坐标系无关。然后，在定义每个气泡组的体积上积分后，我们得到

\begin{sourceblock}
\includegraphics[page=247,pagebox=mediabox,trim={45.00bp 200.26bp 44.37bp 384.55bp},clip,width=0.92\textwidth,max height=0.38\textheight]{Ishii.pdf}
\end{sourceblock}
and

\begin{sourceblock}
\includegraphics[page=247,pagebox=mediabox,trim={45.00bp 93.56bp 44.37bp 497.71bp},clip,width=0.92\textwidth,max height=0.38\textheight]{Ishii.pdf}
\end{sourceblock}
其中每个气泡组的平均界面速度定义为

\begin{sourceblock}
\includegraphics[page=248,pagebox=mediabox,trim={45.00bp 489.28bp 44.37bp 58.65bp},clip,width=0.92\textwidth,max height=0.38\textheight]{Ishii.pdf}
\end{sourceblock}
在式(10-61)和(10-62)中，方程左边的第三项表示由于颗粒体积变化引起的界面面积浓度的变化，使得

\begin{sourceblock}
\includegraphics[page=248,pagebox=mediabox,trim={45.00bp 385.25bp 44.37bp 234.78bp},clip,width=0.92\textwidth,max height=0.38\textheight]{Ishii.pdf}
\end{sourceblock}
and

\begin{sourceblock}
\includegraphics[page=248,pagebox=mediabox,trim={45.00bp 308.21bp 44.37bp 311.82bp},clip,width=0.92\textwidth,max height=0.38\textheight]{Ishii.pdf}
\end{sourceblock}
其中 ic c c A fV 归因于两组气泡之间气泡相互作用导致的组间转移。因此，当 0 c f → 时，组间相互作用没有贡献。然而，实际上，当存在两个气泡群时，cf 是有限的，并且该群间传递项作为每个气泡群的界面面积浓度的源或汇起着重要作用。

为了纳入这种组间转移的贡献，应指定粒子分布函数。然而，对不同尺寸气泡的两相流中颗粒分布函数的精确数学描述需要利用原始的玻尔兹曼输运方程和统计力学。为了我们的目的，我们需要开发一个简单的综合传输方程。因此，在本分析中，为简单起见，假设颗粒分布呈线性分布或均匀分布，如图10-2所示。其中，1 p V 是指定值 1 f 的第 1 组气泡中的峰值气泡体积，并定义为

\begin{sourceblock}
\includegraphics[page=248,pagebox=mediabox,trim={45.00bp 54.58bp 44.37bp 572.88bp},clip,width=0.92\textwidth,max height=0.38\textheight]{Ishii.pdf}
\end{sourceblock}
\begin{sourceblock}
\includegraphics[page=249,pagebox=mediabox,trim={50.00bp 401.16bp 44.37bp 55.00bp},clip,width=\textwidth,max height=0.72\textheight]{Ishii.pdf}
\par\vspace{0.45\baselineskip}
\small 图10-2。流体颗粒分布函数轮廓的线性近似（Ishii 和
\end{sourceblock}
Kim, 2004)

那么，第 1 组气泡的数密度可以表示为

\begin{sourceblock}
\includegraphics[page=249,pagebox=mediabox,trim={45.00bp 292.98bp 44.37bp 336.20bp},clip,width=0.92\textwidth,max height=0.38\textheight]{Ishii.pdf}
\end{sourceblock}
这产生

\begin{sourceblock}
\includegraphics[page=249,pagebox=mediabox,trim={45.00bp 214.77bp 44.37bp 403.08bp},clip,width=0.92\textwidth,max height=0.38\textheight]{Ishii.pdf}
\end{sourceblock}
现在，考虑如图 10-3(a) 至 10-3(c) 所示的三种限制情况，使得情况 I：常数 c f f = =，因此

\begin{sourceblock}
\includegraphics[page=249,pagebox=mediabox,trim={45.00bp 120.40bp 44.37bp 497.36bp},clip,width=0.92\textwidth,max height=0.38\textheight]{Ishii.pdf}
\end{sourceblock}
\begin{sourceblock}
\includegraphics[page=249,pagebox=mediabox,trim={45.00bp 90.03bp 44.37bp 532.79bp},clip,width=0.92\textwidth,max height=0.38\textheight]{Ishii.pdf}
\end{sourceblock}
\begin{sourceblock}
\includegraphics[page=249,pagebox=mediabox,trim={45.00bp 50.14bp 44.37bp 553.28bp},clip,width=0.92\textwidth,max height=0.38\textheight]{Ishii.pdf}
\end{sourceblock}
\begin{sourceblock}
\includegraphics[page=250,pagebox=mediabox,trim={50.00bp 77.25bp 44.37bp 55.00bp},clip,width=\textwidth,max height=0.72\textheight]{Ishii.pdf}
\par\vspace{0.45\baselineskip}
\small 图10-3。流体颗粒分布的限制条件 (a) f = 常数， (b) ψ → V min / V c，
\end{sourceblock}
(c) χ→1（石井和金，2004）

\begin{sourceblock}
\includegraphics[page=251,pagebox=mediabox,trim={45.00bp 574.53bp 44.37bp 66.68bp},clip,width=0.92\textwidth,max height=0.38\textheight]{Ishii.pdf}
\end{sourceblock}
\begin{sourceblock}
\includegraphics[page=251,pagebox=mediabox,trim={45.00bp 533.46bp 44.37bp 68.68bp},clip,width=0.92\textwidth,max height=0.38\textheight]{Ishii.pdf}
\end{sourceblock}
因此，假设 min c V V >>，并定义任意系数 χ 为

\begin{sourceblock}
\includegraphics[page=251,pagebox=mediabox,trim={45.00bp 465.88bp 44.37bp 175.33bp},clip,width=0.92\textwidth,max height=0.38\textheight]{Ishii.pdf}
\end{sourceblock}
对于每个极限情况，我们得到

\begin{sourceblock}
\includegraphics[page=251,pagebox=mediabox,trim={45.00bp 347.11bp 44.37bp 230.54bp},clip,width=0.92\textwidth,max height=0.38\textheight]{Ishii.pdf}
\end{sourceblock}
并由方程(10-68)

\begin{sourceblock}
\includegraphics[page=251,pagebox=mediabox,trim={45.00bp 270.28bp 44.37bp 347.58bp},clip,width=0.92\textwidth,max height=0.38\textheight]{Ishii.pdf}
\end{sourceblock}
因此，回想一下，当
0 c f →，常数 χ 的边界应为

\begin{sourceblock}
\includegraphics[page=251,pagebox=mediabox,trim={45.00bp 204.56bp 44.37bp 437.30bp},clip,width=0.92\textwidth,max height=0.38\textheight]{Ishii.pdf}
\end{sourceblock}
在提供 χ 的解析解时，我们有三个未知数和三个方程，这样

\begin{sourceblock}
\includegraphics[page=251,pagebox=mediabox,trim={45.00bp 86.16bp 44.37bp 508.75bp},clip,width=0.92\textwidth,max height=0.38\textheight]{Ishii.pdf}
\end{sourceblock}
\begin{sourceblock}
\includegraphics[page=252,pagebox=mediabox,trim={45.00bp 534.88bp 44.37bp 109.50bp},clip,width=0.92\textwidth,max height=0.38\textheight]{Ishii.pdf}
\end{sourceblock}
\begin{sourceblock}
\includegraphics[page=252,pagebox=mediabox,trim={45.00bp 506.43bp 44.37bp 118.31bp},clip,width=0.92\textwidth,max height=0.38\textheight]{Ishii.pdf}
\end{sourceblock}
\begin{sourceblock}
\includegraphics[page=252,pagebox=mediabox,trim={45.00bp 457.53bp 44.37bp 136.68bp},clip,width=0.92\textwidth,max height=0.38\textheight]{Ishii.pdf}
\end{sourceblock}
我们可以将方程（10-64）和（10-65）重写为

\begin{sourceblock}
\includegraphics[page=252,pagebox=mediabox,trim={45.00bp 371.79bp 44.37bp 241.19bp},clip,width=0.92\textwidth,max height=0.38\textheight]{Ishii.pdf}
\end{sourceblock}
and

\begin{sourceblock}
\includegraphics[page=252,pagebox=mediabox,trim={45.00bp 254.03bp 44.37bp 323.48bp},clip,width=0.92\textwidth,max height=0.38\textheight]{Ishii.pdf}
\end{sourceblock}
respectively.

因此，结合上述方程，两组界面面积输运方程由下式给出

\begin{sourceblock}
\includegraphics[page=252,pagebox=mediabox,trim={45.00bp 74.37bp 44.37bp 467.17bp},clip,width=0.92\textwidth,max height=0.38\textheight]{Ishii.pdf}
\end{sourceblock}
\begin{sourceblock}
\includegraphics[page=253,pagebox=mediabox,trim={45.00bp 491.98bp 44.37bp 55.93bp},clip,width=0.92\textwidth,max height=0.38\textheight]{Ishii.pdf}
\end{sourceblock}
分别为第 1 组和第 2 组。其中，方程的左侧表示每组 i a 的变化和对流的时间速率。右侧的每一项代表由于颗粒体积变化、组间转移、各种颗粒相互作用和每组相变而引起的 i a 变化率。总界面面积输运方程也可以通过将两个方程相加得到，使得

\begin{sourceblock}
\includegraphics[page=253,pagebox=mediabox,trim={45.00bp 296.38bp 44.37bp 284.25bp},clip,width=0.92\textwidth,max height=0.38\textheight]{Ishii.pdf}
\end{sourceblock}
其中下标k表示气泡群。

在此分析中，我们通过假设颗粒分布中的线性轮廓或均匀轮廓来演示由于两组气泡之间的气泡相互作用而导致的组间转移。然而，在两相流应用中，实际上可以假设采用均匀的颗粒分布轮廓。

\subsection*{10.3.4 本构关系}\phantomsection\addcontentsline{toc}{subsection}{10.3.4 本构关系}
本节总结了求解界面面积输运方程所需的本构关系，如下：

\begin{sourceblock}
\includegraphics[page=253,pagebox=mediabox,trim={45.00bp 80.07bp 44.37bp 532.49bp},clip,width=0.92\textwidth,max height=0.38\textheight]{Ishii.pdf}
\end{sourceblock}
\begin{sourceblock}
\includegraphics[page=254,pagebox=mediabox,trim={45.00bp 571.60bp 44.37bp 57.73bp},clip,width=0.92\textwidth,max height=0.38\textheight]{Ishii.pdf}
\end{sourceblock}
\begin{sourceblock}
\includegraphics[page=254,pagebox=mediabox,trim={45.00bp 507.82bp 44.37bp 114.80bp},clip,width=0.92\textwidth,max height=0.38\textheight]{Ishii.pdf}
\end{sourceblock}
\begin{sourceblock}
\includegraphics[page=254,pagebox=mediabox,trim={45.00bp 441.77bp 44.37bp 179.73bp},clip,width=0.92\textwidth,max height=0.38\textheight]{Ishii.pdf}
\end{sourceblock}
\begin{sourceblock}
\includegraphics[page=254,pagebox=mediabox,trim={45.00bp 379.84bp 44.37bp 243.31bp},clip,width=0.92\textwidth,max height=0.38\textheight]{Ishii.pdf}
\end{sourceblock}
\begin{sourceblock}
\includegraphics[page=254,pagebox=mediabox,trim={45.00bp 337.91bp 44.37bp 303.40bp},clip,width=0.92\textwidth,max height=0.38\textheight]{Ishii.pdf}
\end{sourceblock}
\begin{sourceblock}
\includegraphics[page=254,pagebox=mediabox,trim={45.00bp 293.93bp 44.37bp 347.31bp},clip,width=0.92\textwidth,max height=0.38\textheight]{Ishii.pdf}
\end{sourceblock}
and

\begin{sourceblock}
\includegraphics[page=254,pagebox=mediabox,trim={45.00bp 220.49bp 44.37bp 402.58bp},clip,width=0.92\textwidth,max height=0.38\textheight]{Ishii.pdf}
\end{sourceblock}
对于连续性方程，第 1 组和第 2 组气泡之间由于稳定状态下气泡相互作用（无相变效应）而产生的净传质速率可以通过两组气泡相互作用的建模获得：

\begin{sourceblock}
\includegraphics[page=254,pagebox=mediabox,trim={45.00bp 95.81bp 44.37bp 517.24bp},clip,width=0.92\textwidth,max height=0.38\textheight]{Ishii.pdf}
\end{sourceblock}
其中，2 , j η 是由于两组气泡之间的第 j 个相互作用（例如气泡合并和破裂）而从第 1 组气泡到第 2 组气泡的净体积转移。

在上面给出的本构关系中，方程（10-84）定义的数源和汇速率应通过对导致界面面积浓度变化的主要粒子相互作用进行机械建模来建立。考虑到气液两相流的范围广泛，导致颗粒聚结或崩解的主要气泡相互作用机制可概括如下：

• 随机碰撞（RC R）：通过湍流涡流驱动的随机碰撞进行合并； • 尾流夹带（WE R）：由于后一个粒子在前一个粒子的尾迹中加速，通过碰撞而合并； • 湍流冲击(TI R)：湍流涡流冲击时的分解； • 剪切(SO R)：围绕帽泡的底边剪切； • 表面不稳定性（SI R）：由于表面不稳定性导致大帽气泡破裂； • 上升速度（RV R ）：由于气泡上升速度差异而产生的碰撞； • 层流剪切 (LS R)：由于粘性流体中的层流剪切而导致的破裂；以及； • 速度梯度(VG R)：由于速度梯度而发生的碰撞。

\chapter{界面面积输运的本构建模}
二流体模型广泛应用于目前的两相流分析代码中，如核反应堆安全分析代码RELAP5、TRAC、CATHARE等。在传统模型中，界面面积浓度由经验相关性给出。这些相关性基于两相流态和态转变标准，它们不能动态地表示界面结构的变化。这种静态方法存在以下缺点。

1. 流态转变准则是稳态、完全发展的流的代数关系。它们没有完全反映界面结构变化的真实动态本质。因此，不能正确考虑进入和发展流动的影响，也不能正确考虑制度之间的逐渐过渡。 2.基于流态转变标准的方法是两步方法，需要流态转变标准和每种流态的界面面积相关性。转变标准和界面面积相关性的复合误差可能非常显著。 3. 转变标准和流态相关的界面相关性在某些特定操作条件和几何形状的有限参数范围内有效。其中大部分是从简单的空气水实验和唯象模型中获得的。通常没有正确考虑几何形状和流体属性的尺度效应。当应用于高低压汽水瞬变时，这些模型可能会导致显著的差异、人为不连续性和数值不稳定。

第10章引入了一种基于物理的方法，即界面区域传输方程，来动态获得界面面积浓度。在气泡流态中，可以假设气泡的形状接近球形且具有相似的尺寸，因此一组界面面积输运方程足以描述界面面积输运现象。然而，在更广义的气液两相流中，例如帽泡流、段塞流和搅动湍流流，存在不同尺寸和形状的气泡，例如球形气泡、扭曲气泡、帽泡、段塞泡或搅动湍流气泡。由于阻力和气泡相互作用机制的差异，气泡尺寸和形状的这些变化显著影响气泡传输现象。在开发适用于各种两相流的传输方程时，应考虑气泡形状和尺寸以及特征传输现象的差异。鉴于此，气泡被分为两组：球形/扭曲气泡作为第 1 组，帽/段塞/搅动湍流气泡作为第 2 组。在第 10 章中，提出了处理两组气泡的一般方法，并制定了两组界面面积输运方程。在将两组界面面积输运方程应用到二流体模型中时，需要对传统二流体模型进行一些修改。这主要是因为两组气泡的引入需要两个气体速度场，而传统的二流体模型仅通过动量方程提供一种气体速度。

本章介绍了修改后的双流体模型，该模型已准备好在两组界面面积输运方程的方法中实施。可以为两组气泡写出两个动量方程，尽管求解两个气体动量方程还不太实用。然而，对于完全三维流，这可能是必要的，而对于一维流，提出了一种简化的方法。在这种情况下，通过组合两个气体动量方程保留气相平均速度的动量方程。应指定与第 1 组和第 2 组气泡之间的速度差相关的附加术语。可以通过考虑压力梯度和一般阻力，基于第 1 组和第 2 组气泡的简化动量方程来估计该速度差。此外，在一维简化中，可以应用修改的漂移通量模型来求解速度差。除此之外，本章还演示了一组和两组界面区域输运方程中汇和源项的建模。

\section*{11.1 修正的二组界面面积输运方程的二流体模型}\phantomsection\addcontentsline{toc}{section}{11.1 修正的二组界面面积输运方程的二流体模型}
\subsection*{11.1.1 传统二流体模型}\phantomsection\addcontentsline{toc}{subsection}{11.1.1 传统二流体模型}
正如第 9 章所讨论的，三维二维流体模型是通过使用时间或统计平均获得的。该模型用两组守恒方程来表示，这些方程控制着每个相中的质量、动量和能量的平衡。然而，由于一相的平均场并不独立于另一相，因此相互作用项作为源项出现在场方程中。对于大多数实际应用，双流体模型可以简化为第 9 章中的以下形式（Ishii，1977；Ishii 和 Mishima，1984）。

气相连续性方程

\begin{sourceblock}
\includegraphics[page=258,pagebox=mediabox,trim={45.00bp 363.15bp 44.37bp 260.03bp},clip,width=0.92\textwidth,max height=0.38\textheight]{Ishii.pdf}
\end{sourceblock}
液相连续性方程

\begin{sourceblock}
\includegraphics[page=258,pagebox=mediabox,trim={45.00bp 287.10bp 44.37bp 331.90bp},clip,width=0.92\textwidth,max height=0.38\textheight]{Ishii.pdf}
\end{sourceblock}
气相动量方程

\begin{sourceblock}
\includegraphics[page=258,pagebox=mediabox,trim={45.00bp 169.61bp 44.37bp 410.49bp},clip,width=0.92\textwidth,max height=0.38\textheight]{Ishii.pdf}
\end{sourceblock}
液相动量方程

\begin{sourceblock}
\includegraphics[page=258,pagebox=mediabox,trim={45.00bp 50.14bp 44.37bp 529.36bp},clip,width=0.92\textwidth,max height=0.38\textheight]{Ishii.pdf}
\end{sourceblock}
气相热能方程

\begin{sourceblock}
\includegraphics[page=259,pagebox=mediabox,trim={45.00bp 508.42bp 44.37bp 83.16bp},clip,width=0.92\textwidth,max height=0.38\textheight]{Ishii.pdf}
\end{sourceblock}
液相热能方程

\begin{sourceblock}
\includegraphics[page=259,pagebox=mediabox,trim={45.00bp 363.05bp 44.37bp 187.99bp},clip,width=0.92\textwidth,max height=0.38\textheight]{Ishii.pdf}
\end{sourceblock}
其中，k Г 、ik M 、ki T 、ki q ′′ 和 k φ 分别是质量生成、广义界面阻力、界面剪切应力、界面热通量和耗散。为了简单起见，上式中省略了求平均的数学符号，用 k 表示 k T 、ki T 和 k q
μ T 、 ki T 和 C k q。式（11-1）至式（11-6）中，单位体积质量的产生、单位体积广义阻力和单位体积界面能量传递构成界面输运项。界面转移的跳跃条件为

\begin{sourceblock}
\includegraphics[page=259,pagebox=mediabox,trim={45.00bp 144.33bp 44.37bp 439.71bp},clip,width=0.92\textwidth,max height=0.38\textheight]{Ishii.pdf}
\end{sourceblock}
\subsection*{11.1.2 两组空隙率和界面面积输运方程}\phantomsection\addcontentsline{toc}{subsection}{11.1.2 两组空隙率和界面面积输运方程}
两组空隙率输运方程由下式给出

\begin{sourceblock}
\includegraphics[page=260,pagebox=mediabox,trim={45.00bp 565.82bp 44.37bp 57.37bp},clip,width=0.92\textwidth,max height=0.38\textheight]{Ishii.pdf}
\end{sourceblock}
其中，对于第 1 组和第 2 组，k 分别 =1 和 2。 gk Г 是由于相变而导致的第 k 组气泡的质量生成速率，12 m Δ 表示由于气泡相互作用和流体动力学效应而从第 1 组气泡到第 2 组气泡的净组间传质速率

\begin{sourceblock}
\includegraphics[page=260,pagebox=mediabox,trim={45.00bp 449.19bp 44.37bp 168.32bp},clip,width=0.92\textwidth,max height=0.38\textheight]{Ishii.pdf}
\end{sourceblock}
其中 2 、 j η 和 phk η 分别是从第 1 组气泡到第 2 组气泡的净组间空隙率传输以及由于相变而产生的气体体积的源项和汇项。 χ 是组间传递系数，1 * c D 是无量纲气泡直径，定义为

\begin{sourceblock}
\includegraphics[page=260,pagebox=mediabox,trim={45.00bp 340.78bp 44.37bp 286.69bp},clip,width=0.92\textwidth,max height=0.38\textheight]{Ishii.pdf}
\end{sourceblock}
其中 crit D 是第 1 组和第 2 组边界处气泡的体积当量直径。

两组界面面积输运方程由下式给出

\begin{sourceblock}
\includegraphics[page=260,pagebox=mediabox,trim={45.00bp 179.19bp 44.37bp 381.83bp},clip,width=0.92\textwidth,max height=0.38\textheight]{Ishii.pdf}
\end{sourceblock}
\begin{sourceblock}
\includegraphics[page=260,pagebox=mediabox,trim={45.00bp 50.14bp 44.37bp 506.73bp},clip,width=0.92\textwidth,max height=0.38\textheight]{Ishii.pdf}
\end{sourceblock}
\subsection*{11.1.3 修正的二流体模型}\phantomsection\addcontentsline{toc}{subsection}{11.1.3 修正的二流体模型}
接下来，将双流体模型修改为两组界面面积输运方程（Sun et al., 2003）。气相的多场模型给出了一般形式。一般来说，可以假设第 1 组和第 2 组气泡的压力和温度大致相同。但两组的速度并不相同，因此原则上需要引入两个连续性方程和两个动量方程。基于上述假设，第 1 组和第 2 组气泡的气相密度相同。这导致气相连续性方程为
第 1 组气泡的连续性方程

\begin{sourceblock}
\includegraphics[page=261,pagebox=mediabox,trim={45.00bp 332.52bp 44.37bp 290.64bp},clip,width=0.92\textwidth,max height=0.38\textheight]{Ishii.pdf}
\end{sourceblock}
第 2 组气泡的连续性方程

\begin{sourceblock}
\includegraphics[page=261,pagebox=mediabox,trim={45.00bp 258.60bp 44.37bp 364.56bp},clip,width=0.92\textwidth,max height=0.38\textheight]{Ishii.pdf}
\end{sourceblock}
这里，12 m Δ 是由于流体动力学机制引起的组间传质。此外，如果引入以下恒等式，

\begin{sourceblock}
\includegraphics[page=261,pagebox=mediabox,trim={45.00bp 98.81bp 44.37bp 449.63bp},clip,width=0.92\textwidth,max height=0.38\textheight]{Ishii.pdf}
\end{sourceblock}
则式（11-13）和式（11-14）求和即可恢复传统的连续性方程，即式（11-1）。液相的连续性方程与方程（11-2）相同

\begin{sourceblock}
\includegraphics[page=262,pagebox=mediabox,trim={45.00bp 582.71bp 44.37bp 58.11bp},clip,width=0.92\textwidth,max height=0.38\textheight]{Ishii.pdf}
\end{sourceblock}
由于两组气泡的引入，动量方程更加复杂。与气相的连续性方程不同，由于动量方程的复杂性（至少对于一维公式而言），第 1 组和第 2 组气泡不希望有两个动量方程。如果我们假设第 2 组气泡是除液相和第 1 组气泡之外的“第三相”，并忽略第 1 组和第 2 组气泡之间的直接动量相互作用，则第 1 组和第 2 组气泡的两个动量方程可以写为

\begin{sourceblock}
\includegraphics[page=262,pagebox=mediabox,trim={45.00bp 357.30bp 44.37bp 222.88bp},clip,width=0.92\textwidth,max height=0.38\textheight]{Ishii.pdf}
\end{sourceblock}
and

\begin{sourceblock}
\includegraphics[page=262,pagebox=mediabox,trim={45.00bp 236.21bp 44.37bp 343.90bp},clip,width=0.92\textwidth,max height=0.38\textheight]{Ishii.pdf}
\end{sourceblock}
然后，结合方程（11-17）和（11-18）得出

\begin{sourceblock}
\includegraphics[page=262,pagebox=mediabox,trim={45.00bp 50.14bp 44.37bp 464.74bp},clip,width=0.92\textwidth,max height=0.38\textheight]{Ishii.pdf}
\end{sourceblock}
\begin{sourceblock}
\includegraphics[page=263,pagebox=mediabox,trim={45.00bp 564.72bp 44.37bp 76.12bp},clip,width=0.92\textwidth,max height=0.38\textheight]{Ishii.pdf}
\end{sourceblock}
符合式(11-15)中的定义。有趣的是，由于不同气泡群中气泡速度的差异，式（11-19）右侧第一项是附加扩散项。

然而，式(11-19)过于复杂，不适合一般应用。如前所述，对于大多数实际应用，两组气泡的压力可以近似相同，使得

\begin{sourceblock}
\includegraphics[page=263,pagebox=mediabox,trim={45.00bp 429.88bp 44.37bp 211.37bp},clip,width=0.92\textwidth,max height=0.38\textheight]{Ishii.pdf}
\end{sourceblock}
此外，可以假设两组气泡的界面剪切非常相似，使得：

\begin{sourceblock}
\includegraphics[page=263,pagebox=mediabox,trim={45.00bp 357.88bp 44.37bp 283.36bp},clip,width=0.92\textwidth,max height=0.38\textheight]{Ishii.pdf}
\end{sourceblock}
我们还有以下定义来进一步简化方程（11-19）

\begin{sourceblock}
\includegraphics[page=263,pagebox=mediabox,trim={45.00bp 156.26bp 44.37bp 340.56bp},clip,width=0.92\textwidth,max height=0.38\textheight]{Ishii.pdf}
\end{sourceblock}
因此，式(11-19)可简化为

\begin{sourceblock}
\includegraphics[page=263,pagebox=mediabox,trim={45.00bp 55.23bp 44.37bp 537.70bp},clip,width=0.92\textwidth,max height=0.38\textheight]{Ishii.pdf}
\end{sourceblock}
\begin{sourceblock}
\includegraphics[page=264,pagebox=mediabox,trim={45.00bp 540.29bp 44.37bp 73.96bp},clip,width=0.92\textwidth,max height=0.38\textheight]{Ishii.pdf}
\end{sourceblock}
可以合理地假设主体流体中的平均应力和界面处的平均应力大致相同。因此，

\begin{sourceblock}
\includegraphics[page=264,pagebox=mediabox,trim={45.00bp 466.67bp 44.37bp 173.89bp},clip,width=0.92\textwidth,max height=0.38\textheight]{Ishii.pdf}
\end{sourceblock}
那么式(11-23)进一步简化为

\begin{sourceblock}
\includegraphics[page=264,pagebox=mediabox,trim={45.00bp 317.16bp 44.37bp 232.61bp},clip,width=0.92\textwidth,max height=0.38\textheight]{Ishii.pdf}
\end{sourceblock}
应针对第 1 组和第 2 组气泡单独建模广义界面阻力项 1 ig M 和 2 ig M。

此外，液相的动量方程与式（11-3）具有相同的形式。因此，

\begin{sourceblock}
\includegraphics[page=264,pagebox=mediabox,trim={45.00bp 155.27bp 44.37bp 420.76bp},clip,width=0.92\textwidth,max height=0.38\textheight]{Ishii.pdf}
\end{sourceblock}
with

\begin{sourceblock}
\includegraphics[page=264,pagebox=mediabox,trim={45.00bp 93.80bp 44.37bp 546.16bp},clip,width=0.92\textwidth,max height=0.38\textheight]{Ishii.pdf}
\end{sourceblock}
可以合理地假设液相和界面处的平均应力大致相同。因此，

\begin{sourceblock}
\includegraphics[page=265,pagebox=mediabox,trim={45.00bp 582.71bp 44.37bp 57.86bp},clip,width=0.92\textwidth,max height=0.38\textheight]{Ishii.pdf}
\end{sourceblock}
则式(11-26)进一步简化为

\begin{sourceblock}
\includegraphics[page=265,pagebox=mediabox,trim={45.00bp 460.33bp 44.37bp 116.69bp},clip,width=0.92\textwidth,max height=0.38\textheight]{Ishii.pdf}
\end{sourceblock}
在上述推导中，假设两组气泡的压力和温度基本相同。那么，与动量方程类似，气相的热能方程可以表示为

\begin{sourceblock}
\includegraphics[page=265,pagebox=mediabox,trim={45.00bp 311.63bp 44.37bp 279.94bp},clip,width=0.92\textwidth,max height=0.38\textheight]{Ishii.pdf}
\end{sourceblock}
应用以下定义的情况

\begin{sourceblock}
\includegraphics[page=265,pagebox=mediabox,trim={45.00bp 139.11bp 44.37bp 384.11bp},clip,width=0.92\textwidth,max height=0.38\textheight]{Ishii.pdf}
\end{sourceblock}
\begin{sourceblock}
\includegraphics[page=265,pagebox=mediabox,trim={45.00bp 67.77bp 44.37bp 533.66bp},clip,width=0.92\textwidth,max height=0.38\textheight]{Ishii.pdf}
\end{sourceblock}
\begin{sourceblock}
\includegraphics[page=266,pagebox=mediabox,trim={45.00bp 488.68bp 44.37bp 81.06bp},clip,width=0.92\textwidth,max height=0.38\textheight]{Ishii.pdf}
\end{sourceblock}
液相侧的界面传热为

\begin{sourceblock}
\includegraphics[page=266,pagebox=mediabox,trim={45.00bp 412.60bp 44.37bp 209.45bp},clip,width=0.92\textwidth,max height=0.38\textheight]{Ishii.pdf}
\end{sourceblock}
请注意，应满足以下界面能量条件

\begin{sourceblock}
\includegraphics[page=266,pagebox=mediabox,trim={45.00bp 353.51bp 44.37bp 286.24bp},clip,width=0.92\textwidth,max height=0.38\textheight]{Ishii.pdf}
\end{sourceblock}
在上面的推导中，引入了非常复杂的界面输运项。为了用两组界面面积输运方程求解修正的双流体模型，应为附加变量指定各种本构关系、界面输运项和边界条件。这些变量可以概括为

\begin{sourceblock}
\includegraphics[page=266,pagebox=mediabox,trim={45.00bp 199.60bp 44.37bp 397.14bp},clip,width=0.92\textwidth,max height=0.38\textheight]{Ishii.pdf}
\end{sourceblock}
\subsection*{11.1.4 两种气体速度场的建模}\phantomsection\addcontentsline{toc}{subsection}{11.1.4 两种气体速度场的建模}
对于强一维流动，引入两个气体动量方程可能会带来不必要的复杂化。在这种情况下，气体混合物动量方程和指定第 1 组和第 2 组气体速度之间相对速度的附加本构关系就足够了。重要的是要确保这样做不会过度指定未知数，因为未知数的数量应等于可用方程的数量。

气泡速度的差异可能与局部滑移有关：

\begin{sourceblock}
\includegraphics[page=267,pagebox=mediabox,trim={45.00bp 557.51bp 44.37bp 83.32bp},clip,width=0.92\textwidth,max height=0.38\textheight]{Ishii.pdf}
\end{sourceblock}
为了获得气相和液相之间的局部相对速度，可以在漂移通量模型公式中采用与 Ishii (1977) 类似的方法。

第 1 组气泡的动量方程，即式（11-17）可以通过代入连续性方程并考虑式（11-20）和（11-21）以及 g gi p p = 的假设，写成以下形式。

\begin{sourceblock}
\includegraphics[page=267,pagebox=mediabox,trim={45.00bp 395.16bp 44.37bp 184.89bp},clip,width=0.92\textwidth,max height=0.38\textheight]{Ishii.pdf}
\end{sourceblock}
类似地，我们获得第 2 组气泡和液相的动量方程为

\begin{sourceblock}
\includegraphics[page=267,pagebox=mediabox,trim={45.00bp 285.12bp 44.37bp 316.78bp},clip,width=0.92\textwidth,max height=0.38\textheight]{Ishii.pdf}
\end{sourceblock}
and

\begin{sourceblock}
\includegraphics[page=267,pagebox=mediabox,trim={45.00bp 188.16bp 44.37bp 413.74bp},clip,width=0.92\textwidth,max height=0.38\textheight]{Ishii.pdf}
\end{sourceblock}
为了获得局部相对速度相关性，我们考虑一种特殊条件，例如无相变且横向压力梯度可忽略不计的稳态条件。在没有相变效应的情况下，各相的界面速度和相速度可以认为是相等的，即

\begin{sourceblock}
\includegraphics[page=267,pagebox=mediabox,trim={45.00bp 79.66bp 44.37bp 561.60bp},clip,width=0.92\textwidth,max height=0.38\textheight]{Ishii.pdf}
\end{sourceblock}
各相的压力可近似为

\begin{sourceblock}
\includegraphics[page=268,pagebox=mediabox,trim={45.00bp 586.90bp 44.37bp 54.34bp},clip,width=0.92\textwidth,max height=0.38\textheight]{Ishii.pdf}
\end{sourceblock}
在这些近似下，对于近一维流动，上述动量方程可以表示为

\begin{sourceblock}
\includegraphics[page=268,pagebox=mediabox,trim={45.00bp 516.88bp 44.37bp 124.36bp},clip,width=0.92\textwidth,max height=0.38\textheight]{Ishii.pdf}
\end{sourceblock}
\begin{sourceblock}
\includegraphics[page=268,pagebox=mediabox,trim={45.00bp 472.90bp 44.37bp 168.34bp},clip,width=0.92\textwidth,max height=0.38\textheight]{Ishii.pdf}
\end{sourceblock}
and

\begin{sourceblock}
\includegraphics[page=268,pagebox=mediabox,trim={45.00bp 412.48bp 44.37bp 228.30bp},clip,width=0.92\textwidth,max height=0.38\textheight]{Ishii.pdf}
\end{sourceblock}
由界面力平衡，即式(11-27)，将上述三个方程相加可得

\begin{sourceblock}
\includegraphics[page=268,pagebox=mediabox,trim={45.00bp 340.53bp 44.37bp 286.81bp},clip,width=0.92\textwidth,max height=0.38\textheight]{Ishii.pdf}
\end{sourceblock}
其中 m τ M 是与混合物横向应力梯度相关的力，由下式给出

\begin{sourceblock}
\includegraphics[page=268,pagebox=mediabox,trim={45.00bp 246.55bp 44.37bp 370.67bp},clip,width=0.92\textwidth,max height=0.38\textheight]{Ishii.pdf}
\end{sourceblock}
with

\begin{sourceblock}
\includegraphics[page=268,pagebox=mediabox,trim={45.00bp 189.32bp 44.37bp 451.41bp},clip,width=0.92\textwidth,max height=0.38\textheight]{Ishii.pdf}
\end{sourceblock}
而 m ρ 定义为

\begin{sourceblock}
\includegraphics[page=268,pagebox=mediabox,trim={45.00bp 127.48bp 44.37bp 513.30bp},clip,width=0.92\textwidth,max height=0.38\textheight]{Ishii.pdf}
\end{sourceblock}
公式（11-45）还假设公式（11-28）对于大多数应用都有效。

此外，由式（11-45）可知，重力场可以用动量方程中未知的压力场代替，为

\begin{sourceblock}
\includegraphics[page=269,pagebox=mediabox,trim={45.00bp 569.80bp 44.37bp 57.66bp},clip,width=0.92\textwidth,max height=0.38\textheight]{Ishii.pdf}
\end{sourceblock}
这使得该方法可以在微重力条件下应用。在稳态下，广义界面阻力通过忽略虚拟质量力、巴塞特力和非阻力（例如升力）来近似：

\begin{sourceblock}
\includegraphics[page=269,pagebox=mediabox,trim={45.00bp 470.73bp 44.37bp 152.31bp},clip,width=0.92\textwidth,max height=0.38\textheight]{Ishii.pdf}
\end{sourceblock}
其中 1 d r 和 1 DC 分别为 Group-1 气泡的阻力半径和阻力系数，Group-1 气泡的相对速度定义为

\begin{sourceblock}
\includegraphics[page=269,pagebox=mediabox,trim={45.00bp 402.87bp 44.37bp 238.36bp},clip,width=0.92\textwidth,max height=0.38\textheight]{Ishii.pdf}
\end{sourceblock}
因此，在稳态下，利用式（11-49）和（11-50），我们可以将式（11-42）重写为

\begin{sourceblock}
\includegraphics[page=269,pagebox=mediabox,trim={45.00bp 277.85bp 44.37bp 307.48bp},clip,width=0.92\textwidth,max height=0.38\textheight]{Ishii.pdf}
\end{sourceblock}
或采用以下形式

\begin{sourceblock}
\includegraphics[page=269,pagebox=mediabox,trim={45.00bp 196.56bp 44.37bp 420.10bp},clip,width=0.92\textwidth,max height=0.38\textheight]{Ishii.pdf}
\end{sourceblock}
同样，对于第 2 组气泡，我们有以下公式

\begin{sourceblock}
\includegraphics[page=269,pagebox=mediabox,trim={45.00bp 117.53bp 44.37bp 499.03bp},clip,width=0.92\textwidth,max height=0.38\textheight]{Ishii.pdf}
\end{sourceblock}
其中 2 d r 和 2 DC 分别是第 2 组气泡的阻力半径和阻力系数，第 2 组气泡的相对速度定义为

\begin{sourceblock}
\includegraphics[page=270,pagebox=mediabox,trim={45.00bp 586.89bp 44.37bp 54.32bp},clip,width=0.92\textwidth,max height=0.38\textheight]{Ishii.pdf}
\end{sourceblock}
由式（11-53）和（11-54），结合式（11-36），我们可以求解两组气泡之间的局部滑移，即 1 2 g g − v v。对于一维流动，可以利用第 14 章中讨论的一维漂移通量模型来指定速度差。将一维漂移通量模型修改为两组界面面积输运方程：

\begin{sourceblock}
\includegraphics[page=270,pagebox=mediabox,trim={45.00bp 460.89bp 44.37bp 183.49bp},clip,width=0.92\textwidth,max height=0.38\textheight]{Ishii.pdf}
\end{sourceblock}
其中 gk v 、0 k C 和 gjk V 分别是 k 组气泡的空隙率加权平均气体速度、分布参数和空隙率加权平均漂移速度。然后，速度差由下式给出

\begin{sourceblock}
\includegraphics[page=270,pagebox=mediabox,trim={45.00bp 325.63bp 44.37bp 282.31bp},clip,width=0.92\textwidth,max height=0.38\textheight]{Ishii.pdf}
\end{sourceblock}
两组气泡的分布参数应从某些流动几何形状的实验数据中获得。此外，如果我们假设对于某些流动，两组气泡的分布参数本质上相似，则可以近似获得以下简化形式：

\begin{sourceblock}
\includegraphics[page=270,pagebox=mediabox,trim={45.00bp 215.85bp 44.37bp 425.38bp},clip,width=0.92\textwidth,max height=0.38\textheight]{Ishii.pdf}
\end{sourceblock}
\section*{11.2 一组界面面积输运方程中源项和汇项的建模}\phantomsection\addcontentsline{toc}{section}{11.2 一组界面面积输运方程中源项和汇项的建模}
为了模拟由颗粒聚结和破碎引起的界面区域输运方程中的积分源项和汇项，一般方法将气泡分为两组：球形/扭曲气泡组和帽/块状气泡组，产生两个涉及组内和组间相互作用的界面区域输运方程，如图11-1所示。如图11-2所示，这些相互作用的机制可以概括为五类：液体湍流驱动的随机碰撞导致的聚结；由于尾流而合并

\begin{sourceblock}
\includegraphics[page=271,pagebox=mediabox,trim={50.00bp 352.50bp 44.37bp 55.00bp},clip,width=\textwidth,max height=0.72\textheight]{Ishii.pdf}
\par\vspace{0.45\baselineskip}
\small 图 11-1。两组气泡可能相互作用的分类（Hibiki 和 Ishii，
\end{sourceblock}
\begin{sourceblock}
\includegraphics[page=271,pagebox=mediabox,trim={50.00bp 64.86bp 44.37bp 327.64bp},clip,width=\textwidth,max height=0.72\textheight]{Ishii.pdf}
\par\vspace{0.45\baselineskip}
\small 图 11-2。两组气泡相互作用的示意图（Ishii 等人，2002）
\end{sourceblock}
实际上，当两相气泡流的空隙率很小时，不存在帽泡或段塞泡。然后，两组界面面积输运方程被简化为一组，而不涉及两组之间的相互作用：

\begin{sourceblock}
\includegraphics[page=272,pagebox=mediabox,trim={45.00bp 362.69bp 44.37bp 223.60bp},clip,width=0.92\textwidth,max height=0.38\textheight]{Ishii.pdf}
\end{sourceblock}
本节简要解释了一组界面区域输运方程中源项和汇项的一些模型。

\subsection*{11.2.1 Wu 等人建模的源项和汇项。 (1998)}\phantomsection\addcontentsline{toc}{subsection}{11.2.1 Wu 等人建模的源项和汇项。 (1998)}
吴等人。 (1998)考虑了绝热气泡流中界面区域传输的三种机制，即由于液体湍流驱动的随机碰撞而导致的聚结、由于尾流夹带而导致的聚结以及由于湍流涡流的影响而导致的破裂。则式(11-59)进一步简化为

\begin{sourceblock}
\includegraphics[page=272,pagebox=mediabox,trim={45.00bp 134.69bp 44.37bp 451.60bp},clip,width=0.92\textwidth,max height=0.38\textheight]{Ishii.pdf}
\end{sourceblock}
A. 随机碰撞导致的气泡合并 为了模拟连续介质中湍流驱动的气泡合并速率，气泡随机碰撞速率至关重要。假设这些碰撞仅发生在相邻气泡之间，因为长程相互作用是由大气泡驱动的

\begin{sourceblock}
\includegraphics[page=273,pagebox=mediabox,trim={50.00bp 412.98bp 44.37bp 55.00bp},clip,width=\textwidth,max height=0.72\textheight]{Ishii.pdf}
\par\vspace{0.45\baselineskip}
\small 图 11-3。两个接近气泡的几何定义（Wu et al., 1998）
\end{sourceblock}
输送气泡群而不导致显著相对运动的涡流（Prince 和 Blanch，1990；Coulaloglou 和 Tavlarides，1976）。如图11-3所示，相邻两个大小相同的球形气泡之间，一次碰撞的时间间隔C t Δ 定义为

\begin{sourceblock}
\includegraphics[page=273,pagebox=mediabox,trim={45.00bp 274.90bp 44.37bp 325.36bp},clip,width=0.92\textwidth,max height=0.38\textheight]{Ishii.pdf}
\end{sourceblock}
其中，tu 为两个气泡的均方根接近速度，TL 为一次碰撞时两个气泡之间的平均行进距离。这近似为

\begin{sourceblock}
\includegraphics[page=273,pagebox=mediabox,trim={45.00bp 168.54bp 44.37bp 448.30bp},clip,width=0.92\textwidth,max height=0.38\textheight]{Ishii.pdf}
\end{sourceblock}
式中 E D 为含有 1 个气泡的混合体积的有效直径，b D 为气泡直径。由于一次碰撞的气泡行进长度从 E D 到 ( ) E b D D − 变化，因此在式(11-62)中引入因子 δ 来表征平均效应（而 δ ′ 是考虑 E D 和 E b D D − 之间比例符号的集体参数）

1 3 b g D α )。对于较小的空隙率，δ ' 的作用较小，因为 E D 远大于 b D。然而，重要的是行进长度是否与平均气泡尺寸相当。当空隙率

\begin{sourceblock}
\includegraphics[page=274,pagebox=mediabox,trim={45.00bp 568.49bp 44.37bp 55.78bp},clip,width=0.92\textwidth,max height=0.38\textheight]{Ishii.pdf}
\end{sourceblock}
\begin{sourceblock}
\includegraphics[page=274,pagebox=mediabox,trim={45.00bp 484.50bp 44.37bp 125.92bp},clip,width=0.92\textwidth,max height=0.38\textheight]{Ishii.pdf}
\end{sourceblock}
因此，两个彼此相向运动的气泡的碰撞频率 RC f 由下式给出

\begin{sourceblock}
\includegraphics[page=274,pagebox=mediabox,trim={45.00bp 393.16bp 44.37bp 223.68bp},clip,width=0.92\textwidth,max height=0.38\textheight]{Ishii.pdf}
\end{sourceblock}
然而，由于气泡并不总是向彼此移动，因此此处考虑气泡向相邻气泡移动的概率 C P 来修改碰撞率。通过假设六方密堆积结构，该概率由下式给出

\begin{sourceblock}
\includegraphics[page=274,pagebox=mediabox,trim={45.00bp 279.57bp 44.37bp 346.27bp},clip,width=0.92\textwidth,max height=0.38\textheight]{Ishii.pdf}
\end{sourceblock}
其中 g,crit α 是当中心气泡无法通过相邻气泡之间的自由空间时的临界空隙率。实际上，相邻的气泡不断运动，临界空隙率可能非常接近致密堆积极限。这导致

\begin{sourceblock}
\includegraphics[page=274,pagebox=mediabox,trim={45.00bp 160.27bp 44.37bp 458.67bp},clip,width=0.92\textwidth,max height=0.38\textheight]{Ishii.pdf}
\end{sourceblock}
随后，具有气泡数密度 b n 的混合物的碰撞频率由下式给出

\begin{sourceblock}
\includegraphics[page=274,pagebox=mediabox,trim={45.00bp 55.72bp 44.37bp 551.02bp},clip,width=0.92\textwidth,max height=0.38\textheight]{Ishii.pdf}
\end{sourceblock}
\begin{sourceblock}
\includegraphics[page=275,pagebox=mediabox,trim={45.00bp 550.95bp 44.37bp 55.85bp},clip,width=0.92\textwidth,max height=0.38\textheight]{Ishii.pdf}
\end{sourceblock}
上述碰撞率的函数依赖性与 Coulaloglou 和 Tavlarides (1976) 为液-液液滴流系统提出的模型一致，类似于理想气体中的粒子碰撞模型。不同之处在于，当前模型在括号中包含一个额外项，它涵盖了气泡的平均自由程与平均气泡尺寸相当的情况。然而，目前形式的模型仍然不完整，因为无论相邻气泡距离多远，只要存在有限的接近速度，就会发生碰撞。实际上，当平均距离很大时，不应计算碰撞，因为相邻气泡之间碰撞的相对运动范围受到与气泡尺寸相当的涡流尺寸的限制。为了考虑这种影响，建议对方程（11-67）采用以下修改因子

\begin{sourceblock}
\includegraphics[page=275,pagebox=mediabox,trim={45.00bp 315.83bp 44.37bp 300.77bp},clip,width=0.92\textwidth,max height=0.38\textheight]{Ishii.pdf}
\end{sourceblock}
其中 T C 和 t L 分别是一个可调参数，具体取决于流体的特性以及将相邻气泡驱动在一起的涡流的平均大小。假设这些涡流与平均气泡尺寸处于同一数量级，因为较小的涡流不会为气泡提供相当大的整体运动。然而，较大的涡流会输送气泡群，而不会引起气泡之间显著的相对运动。此后，气泡碰撞频率的最终形式由下式给出

\begin{sourceblock}
\includegraphics[page=275,pagebox=mediabox,trim={45.00bp 93.57bp 44.37bp 470.45bp},clip,width=0.92\textwidth,max height=0.38\textheight]{Ishii.pdf}
\end{sourceblock}
对于每次碰撞，可能不会发生合并，因此许多研究人员提出了碰撞效率（Oolman 和 Blanch，
模型中使用 C λ 来描述每次碰撞后合并现象的随机性。然而，恒定的聚结效率只是一个近似值，需要进一步努力对效率进行机械建模。方程（11-69）中的平均气泡脉动速度 tu 与长度尺度两点之间的均方根液体脉动速度差 b D 成正比，由下式给出
1 3 1 3 b D ε 其中ε 是连续介质每单位质量的能量耗散率（Rotta，1972）。因此，由于随机碰撞引起的气泡合并而导致的界面面积浓度的降低率，
RC φ，由下式给出

\begin{sourceblock}
\includegraphics[page=276,pagebox=mediabox,trim={45.00bp 206.37bp 44.37bp 316.43bp},clip,width=0.92\textwidth,max height=0.38\textheight]{Ishii.pdf}
\end{sourceblock}
其中 RC Γ 是一个可调节的参数，具体取决于流体的性质，经实验确定为 0.016。式(11-70)中的常数设定为g,max α =0.75且T C =3。

B. 由于尾流夹带导致的气泡合并 当气泡进入前导气泡的尾流区域时，它们将加速并可能与前一个气泡碰撞（Otake 等，1977；Bilicki 和 Kestin，1987；Stewart，1995）。对于液体介质中带有尾流区域的球形气泡，定义了有效尾流体积 W V，其中后续气泡可能与前导气泡发生碰撞

\begin{sourceblock}
\includegraphics[page=277,pagebox=mediabox,trim={45.00bp 532.51bp 44.37bp 71.75bp},clip,width=0.92\textwidth,max height=0.38\textheight]{Ishii.pdf}
\end{sourceblock}
假设尾流区气泡追上前一个气泡的平均时间间隔为 W t Δ，则单位混合体积的碰撞率 WE R 应满足

\begin{sourceblock}
\includegraphics[page=277,pagebox=mediabox,trim={45.00bp 396.50bp 44.37bp 168.82bp},clip,width=0.92\textwidth,max height=0.38\textheight]{Ishii.pdf}
\end{sourceblock}
其中 rW u 是前导气泡和尾流区域气泡之间的平均相对速度。 Schlichting (1979)给出了无量纲相对速度的解析表达式为

\begin{sourceblock}
\includegraphics[page=277,pagebox=mediabox,trim={45.00bp 294.64bp 44.37bp 327.71bp},clip,width=0.92\textwidth,max height=0.38\textheight]{Ishii.pdf}
\end{sourceblock}
其中 rW v 、 r v 、 DC 、 A 、 β 和 y 分别为： 前导气泡与尾流区域气泡之间的相对速度；前导气泡与液相之间的相对速度；阻力系数；气泡的正面面积；混合长度与尾流宽度之比；以及从前导气泡中心测量的沿流动方向的距离。然后通过在临界距离上对 rW v 进行积分来获得尾流区域中的平均相对速度 rW u：

\begin{sourceblock}
\includegraphics[page=277,pagebox=mediabox,trim={45.00bp 90.39bp 44.37bp 458.75bp},clip,width=0.92\textwidth,max height=0.38\textheight]{Ishii.pdf}
\end{sourceblock}
其中 ( ) W b F L D 是 W b L D 的函数，因为 β 通常被假定为常数 (Schlichting, 1979)。 ( ) W b F L D 的确切形式并不重要

\begin{sourceblock}
\includegraphics[page=278,pagebox=mediabox,trim={45.00bp 447.49bp 44.37bp 146.55bp},clip,width=0.92\textwidth,max height=0.38\textheight]{Ishii.pdf}
\end{sourceblock}
式中，WE 为可调参数，主要由有效尾流长度与气泡尺寸之比和聚结效率决定，实验确定为0.0076。

C. 湍流冲击导致的气泡破裂对于湍流涡流冲击导致的二元气泡破裂，其驱动力来自连续介质中湍流涡流的惯性力惯性F，而保持力为表面张力，
张力 F .以特征长度驱使子气泡分开
b D 在时间间隔 B t Δ 内，简单的动量平衡方法给出以下关系。

\begin{sourceblock}
\includegraphics[page=278,pagebox=mediabox,trim={45.00bp 241.60bp 44.37bp 381.52bp},clip,width=0.92\textwidth,max height=0.38\textheight]{Ishii.pdf}
\end{sourceblock}
这里，由于液体和气体的密度比较大，气泡的惯性由虚拟质量决定。重新排列方程（11-76）可得到以下平均气泡破裂频率

\begin{sourceblock}
\includegraphics[page=278,pagebox=mediabox,trim={45.00bp 140.55bp 44.37bp 482.56bp},clip,width=0.92\textwidth,max height=0.38\textheight]{Ishii.pdf}
\end{sourceblock}
速度 tu 被假定为长度为 b D 的两点之间的均方根速度差，这意味着只有尺寸等于气泡尺寸的涡流才能打破气泡。 Crit We 是一个集体常数，指定为临界韦伯数。由于湍流波动的共振激发，气泡破裂的临界 We 值在很大范围内变化（Sevik 和 Park，1973）。

在均匀湍流中，气泡与具有足够能量来破碎气泡的涡流碰撞的概率，即破碎效率 B λ 大约为（Coulaloglou 和 Tavlarrides，1976）

\begin{sourceblock}
\includegraphics[page=279,pagebox=mediabox,trim={45.00bp 501.14bp 44.37bp 119.46bp},clip,width=0.92\textwidth,max height=0.38\textheight]{Ishii.pdf}
\end{sourceblock}
其中 2 t,crit u 是临界均方脉动速度

\begin{sourceblock}
\includegraphics[page=279,pagebox=mediabox,trim={45.00bp 331.29bp 44.37bp 197.55bp},clip,width=0.92\textwidth,max height=0.38\textheight]{Ishii.pdf}
\end{sourceblock}
实验确定可调参数TI Γ 和crit We 分别为0.17 和6.0。该表达式与之前的模型（Prince 和 Blanch，1990）不同，因为当韦伯数小于临界 We 时，破裂率为零。这一独特的功能允许气泡合并和破裂过程的解耦。在低液体流速和小的空隙率下，湍流脉动很小，因此不会计算破碎。这允许对聚结项中的可调节参数进行微调，而与气泡破裂无关。

D. 一维一组模型 界面面积输运方程的最简单形式是通过对方程（11-60）应用横截面积平均而获得的一维公式

\begin{sourceblock}
\includegraphics[page=279,pagebox=mediabox,trim={45.00bp 50.14bp 44.37bp 533.51bp},clip,width=0.92\textwidth,max height=0.38\textheight]{Ishii.pdf}
\end{sourceblock}
\begin{sourceblock}
\includegraphics[page=280,pagebox=mediabox,trim={45.00bp 513.16bp 44.37bp 108.58bp},clip,width=0.92\textwidth,max height=0.38\textheight]{Ishii.pdf}
\end{sourceblock}
如果忽略气泡的内部循环，这与按空隙率加权的传统面积平均气体速度相同。面积平均源项和汇项的精确数学表达式将涉及许多协方差，这可能使一维问题进一步复杂化。然而，由于这些局部项最初是从混合物的有限体积单元获得的，如果将流路的水力直径视为有限元的长度尺度，则面积平均源项和汇项对平均参数的函数依赖性应大致相同。因此，式（11-70）、（11-75）和（11-79）在横截面积内平均参数的情况下，仍然适用于式（11-80）中的面积平均源汇项。

在式（11-70）、（11-75）和（11-79）中，应规定单位混合料质量的能量耗散率。在完整的二流体模型中，ε 来自其自身的本构关系，例如两相 k - ε 模型（Lopez de Bertodano 等，1994）。然而，对于一维分析，该项可以通过简单的代数方程近似为

\begin{sourceblock}
\includegraphics[page=280,pagebox=mediabox,trim={45.00bp 231.58bp 44.37bp 395.90bp},clip,width=0.92\textwidth,max height=0.38\textheight]{Ishii.pdf}
\end{sourceblock}
其中 m v 、H D 和 TW f 是平均混合物速度、流路的水力直径和两相摩擦系数。

\subsection*{11.2.2 Hibiki 和 Ishii (2000a) 建模的源项和汇项}\phantomsection\addcontentsline{toc}{subsection}{11.2.2 Hibiki 和 Ishii (2000a) 建模的源项和汇项}
Hibiki 和 Ishii (2000a; 2002c) 讨论了尾流夹带对界面区域传输的贡献。尾流夹带将在气泡到段塞流的转变、段塞流和搅动湍流中发挥重要作用。对于小直径管中的气泡流或对于非常低的流量条件也可能很重要，因为气泡的横向波动很小。然而，对于相对高流量的条件，即使在尾流区域捕获的气泡也很容易由于液体湍流而离开尾流区域，导致尾流夹带对界面面积的贡献较小
A.由于随机碰撞而导致的气泡合并气泡合并被认为是由于液相中的湍流引起的气泡随机碰撞而发生的。为了估计气泡与气泡碰撞频率，假设气泡的运动行为类似于理想气体分子（Coulaloglou 和 Tavlarides，1977）。根据气体动力学理论 (Loeb, 1927)，气泡随机碰撞频率 RC f 可以通过假设相同的气泡速度 C u 来表示，作为可用于碰撞的表面的函数，

\begin{sourceblock}
\includegraphics[page=281,pagebox=mediabox,trim={45.00bp 375.58bp 44.37bp 234.63bp},clip,width=0.92\textwidth,max height=0.38\textheight]{Ishii.pdf}
\end{sourceblock}
考虑到气泡的排除体积，表面积和体积由下式给出

\begin{sourceblock}
\includegraphics[page=281,pagebox=mediabox,trim={45.00bp 288.58bp 44.37bp 334.60bp},clip,width=0.92\textwidth,max height=0.38\textheight]{Ishii.pdf}
\end{sourceblock}
\begin{sourceblock}
\includegraphics[page=281,pagebox=mediabox,trim={45.00bp 209.34bp 44.37bp 395.80bp},clip,width=0.92\textwidth,max height=0.38\textheight]{Ishii.pdf}
\end{sourceblock}
其中b N 、b D 、V 、b n 和g α 分别表示气泡数、气泡直径、控制体积、气泡数密度和空隙率。将变量 C β ( ≤ 1) 引入到排除体积中，以便考虑高空隙率区域的排除体积的重叠。尽管它可能是空隙率的函数，但为了简单起见，将其视为常数。由该假设引起的失真将通过稍后介绍的气泡聚结率的最终方程中的调整参数来调整。

各向同性湍流惯性子范围内两点 b D 之间的平均脉动速度差由下式给出（Hinze，1959）

\begin{sourceblock}
\includegraphics[page=282,pagebox=mediabox,trim={45.00bp 583.18bp 44.37bp 57.86bp},clip,width=0.92\textwidth,max height=0.38\textheight]{Ishii.pdf}
\end{sourceblock}
其中 ε 表示能量耗散。考虑到气泡之间的相对运动，平均气泡速度由下式给出

\begin{sourceblock}
\includegraphics[page=282,pagebox=mediabox,trim={45.00bp 510.16bp 44.37bp 130.88bp},clip,width=0.92\textwidth,max height=0.38\textheight]{Ishii.pdf}
\end{sourceblock}
其中 C γ 是常数。当空隙率接近最大空隙率 C,max α 时，碰撞频率将增加到无穷大。由于 74.1\% 的体积实际上被根据面心立方晶格密排的相同球体占据，因此 C,max α 可假定为 0.741。最后，我们得到

\begin{sourceblock}
\includegraphics[page=282,pagebox=mediabox,trim={45.00bp 365.35bp 44.37bp 255.26bp},clip,width=0.92\textwidth,max height=0.38\textheight]{Ishii.pdf}
\end{sourceblock}
其中 C γ ′ 是可调整值。为了获得气泡聚结率，需要确定聚结效率。 Coulaloglou 和 Tavlarides (1977) 给出了聚结效率 C λ 的表达式，它是气泡聚结所需时间 C t 和两个气泡接触时间 C τ 的函数。

\begin{sourceblock}
\includegraphics[page=282,pagebox=mediabox,trim={45.00bp 222.77bp 44.37bp 398.83bp},clip,width=0.92\textwidth,max height=0.38\textheight]{Ishii.pdf}
\end{sourceblock}
Oolman 和 Blanch (1986a; 1986b) 给出了气泡合并所需的时间，用于使相同尺寸的气泡之间的液膜变薄：

\begin{sourceblock}
\includegraphics[page=282,pagebox=mediabox,trim={45.00bp 117.76bp 44.37bp 505.36bp},clip,width=0.92\textwidth,max height=0.38\textheight]{Ishii.pdf}
\end{sourceblock}
其中 f ρ 、 σ 、 init δ 和 crit δ 分别是液体密度、界面张力、初始膜厚度和发生破裂的临界膜厚度。 Levich (1962) 从尺寸考虑推导出湍流中的接触时间。

\begin{sourceblock}
\includegraphics[page=283,pagebox=mediabox,trim={45.00bp 565.43bp 44.37bp 58.17bp},clip,width=0.92\textwidth,max height=0.38\textheight]{Ishii.pdf}
\end{sourceblock}
其中 br 是气泡半径。最后，我们得到

\begin{sourceblock}
\includegraphics[page=283,pagebox=mediabox,trim={45.00bp 487.39bp 44.37bp 129.46bp},clip,width=0.92\textwidth,max height=0.38\textheight]{Ishii.pdf}
\end{sourceblock}
Kirkpatrick 和 Locket (1974) 估计空气-水系统中薄膜的初始厚度为 1 × 10 -4 m，而最终薄膜厚度通常取为 1 × 10 -8 m (Kim 和 Lee, 1987)。因此，空气-水系统的实验系数 C K 确定为 1.29。界面面积浓度的下降率 RC φ 表示为

\begin{sourceblock}
\includegraphics[page=283,pagebox=mediabox,trim={45.00bp 301.00bp 44.37bp 277.08bp},clip,width=0.92\textwidth,max height=0.38\textheight]{Ishii.pdf}
\end{sourceblock}
可调节变量 RC 一定是排除体积的重叠、气泡变形和气泡速度分布的函数。然而，为简单起见，可将可调变量假设为常数，并通过实验确定气泡流的值为 0.0314。

B. 湍流冲击导致气泡破裂 气泡破裂被认为是由于湍流涡流与气泡的碰撞而发生的。为了估计气泡-涡流碰撞频率，假设涡流和气泡的运动行为类似于理想气体分子（Coulaloglou 和 Tavlarides，1977）。此外，气泡涡流碰撞率的建模还做出了以下假设（Prince 和 Blanch，1990）： (i) 湍流是各向同性的； (ii) 感兴趣的涡流大小 e D 位于惯性子范围内； (iii) 尺寸从 e b c D 到 b D 的涡流可以打碎 b D 尺寸的气泡，因为较大的涡流倾向于输送气泡而不是打碎它，而较小的涡流没有足够的能量来打碎气泡。阿兹贝尔和

\begin{sourceblock}
\includegraphics[page=284,pagebox=mediabox,trim={45.00bp 529.67bp 44.37bp 73.81bp},clip,width=0.92\textwidth,max height=0.38\textheight]{Ishii.pdf}
\end{sourceblock}
其中 ( ) e e N k 表示每体积流体的波数 e k (= 2 e D ) 的涡流数量。这里，每体积两相混合物的波数涡流数 ( ) e e n k 由下式给出

\begin{sourceblock}
\includegraphics[page=284,pagebox=mediabox,trim={45.00bp 444.47bp 44.37bp 170.98bp},clip,width=0.92\textwidth,max height=0.38\textheight]{Ishii.pdf}
\end{sourceblock}
根据气体动力学理论 (Loeb, 1927)，气泡涡流随机碰撞频率 TI f 可以通过假设相同的气泡速度 B u 来表示，作为可用于碰撞的表面 B S 和可用于碰撞的体积 B U 的函数

\begin{sourceblock}
\includegraphics[page=284,pagebox=mediabox,trim={45.00bp 331.60bp 44.37bp 271.85bp},clip,width=0.92\textwidth,max height=0.38\textheight]{Ishii.pdf}
\end{sourceblock}
考虑到气泡和涡流的排除体积，表面积和体积由下式给出

\begin{sourceblock}
\includegraphics[page=284,pagebox=mediabox,trim={45.00bp 187.30bp 44.37bp 380.53bp},clip,width=0.92\textwidth,max height=0.38\textheight]{Ishii.pdf}
\end{sourceblock}
\begin{sourceblock}
\includegraphics[page=284,pagebox=mediabox,trim={45.00bp 60.00bp 44.37bp 497.15bp},clip,width=0.92\textwidth,max height=0.38\textheight]{Ishii.pdf}
\end{sourceblock}
\begin{sourceblock}
\includegraphics[page=285,pagebox=mediabox,trim={45.00bp 582.71bp 44.37bp 58.12bp},clip,width=0.92\textwidth,max height=0.38\textheight]{Ishii.pdf}
\end{sourceblock}
其中 ( ) S e F c 和 ( ) V e F c 是由 e b D D 定义的 e c 的函数。将变量 B β ( ≤ 1) 引入到排除体积中，以考虑高空隙率区域的排除体积的重叠。尽管它可能是空隙率的函数，但为了简单起见，将其视为常数。由该假设引起的失真将通过稍后介绍的气泡破裂率的最终方程中的调整参数来调整。

\begin{sourceblock}
\includegraphics[page=285,pagebox=mediabox,trim={45.00bp 412.75bp 44.37bp 198.65bp},clip,width=0.92\textwidth,max height=0.38\textheight]{Ishii.pdf}
\end{sourceblock}
这里，考虑气泡和涡流之间的相对运动，平均相对速度 B u 可以表示为

\begin{sourceblock}
\includegraphics[page=285,pagebox=mediabox,trim={45.00bp 339.16bp 44.37bp 284.37bp},clip,width=0.92\textwidth,max height=0.38\textheight]{Ishii.pdf}
\end{sourceblock}
\begin{sourceblock}
\includegraphics[page=285,pagebox=mediabox,trim={45.00bp 259.26bp 44.37bp 337.55bp},clip,width=0.92\textwidth,max height=0.38\textheight]{Ishii.pdf}
\end{sourceblock}
其中 ( ) B e c γ ′ 和 B,max α 是分别取决于 e c 和最大允许空隙率的可调变量。如果假定气泡大小几乎相同的涡流能够破碎气泡，则方程（11-101）中的最大允许空隙率 B,max α 可近似取与 C,max α 相同的值，即 0.741。因此，气泡-涡流随机碰撞频率的函数形式（方程（11-101））看起来与气泡-气泡随机碰撞频率（方程（11-88））相似。

为了获得气泡破碎率，需要确定破碎效率B λ。破碎效率以单个涡流的平均能量 e E 和气泡破碎所需的平均能量 B E 表示（Prince 和 Blanch，1990；Coulaloglou 和 Tavlarides，1977；Tsouris 和 Tavlarides，1994）

\begin{sourceblock}
\includegraphics[page=286,pagebox=mediabox,trim={45.00bp 565.84bp 44.37bp 55.73bp},clip,width=0.92\textwidth,max height=0.38\textheight]{Ishii.pdf}
\end{sourceblock}
对于二元破裂，即气泡破裂为两个气泡，所需能量 B E 可简单计算为破裂为小子气泡和大子气泡所需能量的平均值，如下所示

\begin{sourceblock}
\includegraphics[page=286,pagebox=mediabox,trim={45.00bp 481.18bp 44.37bp 160.07bp},clip,width=0.92\textwidth,max height=0.38\textheight]{Ishii.pdf}
\end{sourceblock}
\begin{sourceblock}
\includegraphics[page=286,pagebox=mediabox,trim={99.51bp 407.95bp 50.18bp 203.83bp},clip,width=\textwidth,max height=0.72\textheight]{Ishii.pdf}
\end{sourceblock}
1 3 2 b,max b,min b D D D = = 为 2 0 260。 b D πσ。因此，破裂能 B E 近似为 2 0 230。 b D πσ 通过对两种极端情况的分裂能求平均值。这里应该注意的是，通过对方程（11-103）进行平均得到的 B E (= 2 0 230 . b D πσ ) 与假设二元破碎成两个相等尺寸的气泡的 B E (= 2 0 260 . b D πσ ) 之间的相对差异约为 13\%。因此，对小子气泡和大子气泡尺寸的假设可能不会显著影响 B E 的估计。作用于气泡破裂的单个涡流的平均能量可简单地计算为

\begin{sourceblock}
\includegraphics[page=286,pagebox=mediabox,trim={45.00bp 243.46bp 44.37bp 362.39bp},clip,width=0.92\textwidth,max height=0.38\textheight]{Ishii.pdf}
\end{sourceblock}
其中 e 是单个涡流的能量，由下式给出

\begin{sourceblock}
\includegraphics[page=286,pagebox=mediabox,trim={45.00bp 171.66bp 44.37bp 457.54bp},clip,width=0.92\textwidth,max height=0.38\textheight]{Ishii.pdf}
\end{sourceblock}
其中，em 是每个涡流的质量。从式（11-94）、（11-95）、（11-99）、（11-104）和（11-105）中，作用于气泡破裂的单涡流的平均能量由下式给出：

\begin{sourceblock}
\includegraphics[page=287,pagebox=mediabox,trim={45.00bp 506.68bp 44.37bp 58.37bp},clip,width=0.92\textwidth,max height=0.38\textheight]{Ishii.pdf}
\end{sourceblock}
Prince 和 Blanch (1990) 设置了不会导致气泡破裂的最小涡流尺寸，即涡流小于气泡尺寸的 20\%，即 3 e c =0.2。因此，单个涡流的平均能量可以表示为

\begin{sourceblock}
\includegraphics[page=287,pagebox=mediabox,trim={45.00bp 422.14bp 44.37bp 218.78bp},clip,width=0.92\textwidth,max height=0.38\textheight]{Ishii.pdf}
\end{sourceblock}
破碎效率的最终形式由下式给出

\begin{sourceblock}
\includegraphics[page=287,pagebox=mediabox,trim={45.00bp 337.29bp 44.37bp 266.50bp},clip,width=0.92\textwidth,max height=0.38\textheight]{Ishii.pdf}
\end{sourceblock}
其中B K 是常数，为1.59 (=0.230/0.145)。界面面积浓度的增加率 TI φ 可表示为

\begin{sourceblock}
\includegraphics[page=287,pagebox=mediabox,trim={45.00bp 193.01bp 44.37bp 385.07bp},clip,width=0.92\textwidth,max height=0.38\textheight]{Ishii.pdf}
\end{sourceblock}
其中 TI Γ 是一个可调变量。可调变量 Г B 肯定是排除体积的重叠、气泡变形、气泡速度分布以及涡流尺寸与气泡尺寸之比的函数。然而，为简单起见，可将可调变量假设为常数，并通过实验确定气泡流的值为 0.0209。

这里应该注意的是，对于一维分析，每单位质量的能量耗散率可以简单地从机械能方程（Bello，1968）中获得：

\begin{sourceblock}
\includegraphics[page=288,pagebox=mediabox,trim={45.00bp 568.78bp 44.37bp 55.36bp},clip,width=0.92\textwidth,max height=0.38\textheight]{Ishii.pdf}
\end{sourceblock}
其中 j 、 m ρ 、 P 和 z 分别表示混合物体积通量、混合物密度、压力和距测试段入口的轴向位置。

\subsection*{11.2.3 Hibiki 等人建模的源项和汇项。 (2001b)}\phantomsection\addcontentsline{toc}{subsection}{11.2.3 Hibiki 等人建模的源项和汇项。 (2001b)}
响等人。 (2001b) 讨论了在低液体速度下相对较小直径的管中界面区域传输的主要机制，其中气泡破裂可以忽略不计。这里，相对小直径的管被定义为具有相对高的气泡尺寸与管直径比的管。在这种直径相对较小的管中，由于壁的存在，径向气泡运动​​将受到限制，导致微不足道的气泡随机碰撞，而气泡沿流动方向排列，导致显著的尾流夹带。因此，Hibiki 等人。 (2001b) 考虑到气泡聚结机制对管直径的依赖性，由于气泡聚结而开发了汇项。

A. 由于随机碰撞而导致的气泡合并 使用与 Hibiki 和 Ishii (2000a) 中相同的模型。

\begin{sourceblock}
\includegraphics[page=288,pagebox=mediabox,trim={45.00bp 254.40bp 44.37bp 362.45bp},clip,width=0.92\textwidth,max height=0.38\textheight]{Ishii.pdf}
\end{sourceblock}
B. 由于尾流夹带而导致的气泡合并 该模型是通过修改 Wu 等人提出的模型而开发的。 （1998）。

\begin{sourceblock}
\includegraphics[page=288,pagebox=mediabox,trim={45.00bp 147.42bp 44.37bp 469.43bp},clip,width=0.92\textwidth,max height=0.38\textheight]{Ishii.pdf}
\end{sourceblock}
其中 WE Γ 和 C K 分别是空气-水系统的可调节参数，实验确定为 0.082，实验常数确定为 1.29。

C. 管尺寸对界面面积输运机制的影响上述简单的考虑表明，主要机制是

\begin{sourceblock}
\includegraphics[page=289,pagebox=mediabox,trim={45.00bp 318.25bp 44.37bp 301.72bp},clip,width=0.92\textwidth,max height=0.38\textheight]{Ishii.pdf}
\end{sourceblock}
函数 ( ) b f D D 可以根据实验数据近似为

\begin{sourceblock}
\includegraphics[page=289,pagebox=mediabox,trim={45.00bp 234.27bp 44.37bp 391.54bp},clip,width=0.92\textwidth,max height=0.38\textheight]{Ishii.pdf}
\end{sourceblock}
考虑管尺寸效应的界面面积输运方程有望用于预测中小管中气泡流的界面面积输运。

\section*{11.3 两组界面面积输运方程中源汇项的建模}\phantomsection\addcontentsline{toc}{section}{11.3 两组界面面积输运方程中源汇项的建模}
不同流态下的界面结构发生巨大变化。对于盖泡流、段塞流和搅动湍流，气泡根据其几何和物理特征分为两组。球形和扭曲气泡被归类为第 1 组，帽状气泡、团状气泡和搅动湍流气泡被归类为第 2 组。这两类气泡受到不同的聚结/分解机制的影响。因此，一个

\subsection*{11.3.1 Hibiki 和 Ishii (2000b) 建模的源项和汇项}\phantomsection\addcontentsline{toc}{subsection}{11.3.1 Hibiki 和 Ishii (2000b) 建模的源项和汇项}
Hibiki 和 Ishii (2000b) 开发了中等直径管中绝热气泡到段塞过渡流的两组界面面积输运方程，并使用在 50.8 mm 直径管中获取的垂直空气-水流数据对其进行了评估（Hibiki 等人，2001a）。接下来，简要解释界面区域传输机制的分类以及建模的源和汇项。

A. 界面区域传输机制的分类 第 1 组和第 2 组气泡之间的边界可以通过以下方式确定（Ishii 和 Zuber，1979）

\begin{sourceblock}
\includegraphics[page=290,pagebox=mediabox,trim={45.00bp 350.24bp 44.37bp 274.72bp},clip,width=0.92\textwidth,max height=0.38\textheight]{Ishii.pdf}
\end{sourceblock}
其中 crit D 是第 1 组和第 2 组气泡之间边界处的气泡体积当量直径。式（11-115）给出了大气压力下空气-水系统的大约10mm的值。

为了模拟由气泡合并和破裂引起的两组界面面积输运方程中的积分源项和汇项，可以根据所属气泡组将气泡相互作用的可能组合分为八类（见图11-1）： （1）气泡（组1）合并成气泡（组1）； (2) 气泡（第1组）破碎成气泡（第1组）； (3) 气泡（第1组和第2组）合并为气泡（第2组）； (4) 气泡（第 2 组）破裂为气泡（第 1 组和第 2 组）； (5)气泡(第1组)合并成气泡(第2组)； (6) 气泡（第2组）破碎成气泡（第1组）； (7)气泡(第2组)合并成气泡(第2组)； (8)气泡（第2组）破碎成气泡（第2组）。如表 11-1 所示，Hibiki 和 Ishii (2000b) 考虑了三种主要的气泡相互作用：(1) 湍流驱动的随机碰撞导致的合并； (2)尾流夹带引起的聚结； (3) 湍流涡流冲击导致的破碎。他们假设，在气泡到段塞的过渡流中，由于帽泡或段塞气泡的剪切而导致的气泡聚结可能是微不足道的。他们还假设，由于表面不稳定性而导致的气泡破裂在适度的情况下可以忽略不计。

符号 机制 交互参数

\begin{sourceblock}
\includegraphics[page=291,pagebox=mediabox,trim={45.00bp 433.30bp 44.37bp 92.87bp},clip,width=0.92\textwidth,max height=0.38\textheight]{Ishii.pdf}
\end{sourceblock}
由于表面不稳定，管尺寸小于气泡破裂的极限。

Hibiki 和 Ishii (2000b) 开发了二组模型，其中包含由于湍流冲击和尾流夹带而产生的必要的组间耦合项，以及由于组 2 中的尾流夹带和湍流冲击而产生的源和汇项（参见表 11-1）。这里，排除了一些其他机制，例如由于盖泡之间随机碰撞而导致的聚结，因为该模型是针对中等直径管中的气泡到段塞过渡流而开发的。在这种情况下，盖泡和段状气泡将围绕管中心上升，导致盖泡之间随机碰撞的作用较小。该模型还省略了另外两个机制。由于帽泡（第 2 组）完全破裂为小气泡（第 1 组）以及小气泡（第 1 组）合并为帽泡（第 2 组），因此它们是互换术语。由于在数据库的实验条件下（Hibiki et al., 2001a），盖泡与小泡的直径之比约为10至20，因此不太可能发生这些气泡的交换。最终，表 11-1 中列出的六个项被视为两组界面面积输运方程中的源项和汇项，应用于中等直径管中的气泡到段塞过渡流。

B. 简化的两组界面面积输运方程 这里，假设等温流动条件。另外，忽略系数χ。这是因为在组边界处只有很少部分的气泡可以仅仅由于膨胀而转移到另一组。因此，两组界面面积输运方程简化为

\begin{sourceblock}
\includegraphics[page=292,pagebox=mediabox,trim={45.00bp 565.83bp 44.37bp 57.35bp},clip,width=0.92\textwidth,max height=0.38\textheight]{Ishii.pdf}
\end{sourceblock}
\begin{sourceblock}
\includegraphics[page=292,pagebox=mediabox,trim={45.00bp 485.26bp 44.37bp 116.14bp},clip,width=0.92\textwidth,max height=0.38\textheight]{Ishii.pdf}
\end{sourceblock}
\begin{sourceblock}
\includegraphics[page=292,pagebox=mediabox,trim={45.00bp 403.17bp 44.37bp 198.22bp},clip,width=0.92\textwidth,max height=0.38\textheight]{Ishii.pdf}
\end{sourceblock}
C. 建模汇和源项的总结 建模汇和源项总结如下。在一维公式中，所有两相参数，例如 g α 、 i a 和 Sm D，都是面积平均值。为了简单起见，在下面的公式中省略了代表面积平均值的符号。

由于随机碰撞而导致气泡合并

\begin{sourceblock}
\includegraphics[page=292,pagebox=mediabox,trim={45.00bp 232.39bp 44.37bp 384.46bp},clip,width=0.92\textwidth,max height=0.38\textheight]{Ishii.pdf}
\end{sourceblock}
由于尾流夹带而导致气泡合并

\begin{sourceblock}
\includegraphics[page=292,pagebox=mediabox,trim={45.00bp 102.79bp 44.37bp 465.05bp},clip,width=0.92\textwidth,max height=0.38\textheight]{Ishii.pdf}
\end{sourceblock}
\begin{sourceblock}
\includegraphics[page=293,pagebox=mediabox,trim={45.00bp 556.72bp 44.37bp 55.61bp},clip,width=0.92\textwidth,max height=0.38\textheight]{Ishii.pdf}
\end{sourceblock}
\begin{sourceblock}
\includegraphics[page=293,pagebox=mediabox,trim={45.00bp 437.50bp 44.37bp 127.46bp},clip,width=0.92\textwidth,max height=0.38\textheight]{Ishii.pdf}
\end{sourceblock}
气泡因湍流冲击而破裂

\begin{sourceblock}
\includegraphics[page=293,pagebox=mediabox,trim={45.00bp 359.28bp 44.37bp 257.59bp},clip,width=0.92\textwidth,max height=0.38\textheight]{Ishii.pdf}
\end{sourceblock}
\begin{sourceblock}
\includegraphics[page=293,pagebox=mediabox,trim={45.00bp 241.38bp 44.37bp 327.85bp},clip,width=0.92\textwidth,max height=0.38\textheight]{Ishii.pdf}
\end{sourceblock}
\begin{sourceblock}
\includegraphics[page=293,pagebox=mediabox,trim={45.00bp 173.59bp 44.37bp 443.26bp},clip,width=0.92\textwidth,max height=0.38\textheight]{Ishii.pdf}
\end{sourceblock}
\begin{sourceblock}
\includegraphics[page=293,pagebox=mediabox,trim={45.00bp 50.50bp 44.37bp 513.47bp},clip,width=0.92\textwidth,max height=0.38\textheight]{Ishii.pdf}
\end{sourceblock}
\subsection*{11.3.2 Fu 和 Ishii (2002a) 建模的源项和汇项}\phantomsection\addcontentsline{toc}{subsection}{11.3.2 Fu 和 Ishii (2002a) 建模的源项和汇项}
Fu 和 Ishii (2002a) 开发了中等直径管中气泡流、段塞流和搅动湍流的两组界面面积输运方程，并使用在 50.8 mm 直径管中获取的垂直空气-水流数据对其进行了评估。接下来，简要解释界面区域传输机制的分类以及建模的源和汇项。

A. 界面区域传输机制的分类 Fu 和 Ishii (2002a) 采用了五种主要的气泡相互作用： (1) 湍流驱动的随机碰撞导致的聚结； (2)尾流夹带引起的聚结； (3)湍流涡流冲击造成的破碎； （四）因剪断造成的断裂； (5)由于气泡表面流动不稳定性而导致大帽气泡破裂。鉴于将所有源项和汇项纳入界面面积输运方程的复杂性和实验验证的难度，他们根据相互作用项的性质和数量级进行了分析，以简化相互作用项。实验证实，大多数群间相互作用是由尾流夹带和第 1 组气泡与第 2 组气泡之间的剪断引起的（Fu 和 Ishii，2002a）。此外，第 2 组气泡之间的尾流夹带主要控制第 2 组气泡数量，这显著影响流结构和组间相互作用的强度。段状气泡尾流区域的高湍流强度和活跃的涡流-气泡相互作用显著增强了由于表面不稳定性而导致的第 2 组气泡破裂。

由于相似的性质，Group-1 和 Group-2 气泡之间的随机碰撞可以包含在 Group-1 到 Group-2 气泡的尾流夹带中。此外，由于第 1 组气泡数密度较低且尾流区域外的湍流强度较低，第 1 组气泡在第 2 组气泡头部碰撞的贡献可能较小。同时，在中等直径流（2.5 cm ≤ D ≤ 10 cm）中，第 2 组气泡之间的随机碰撞也可以忽略不计，因为主要相互作用通常位于前导气泡的尾流区域内，并且合并机制可以视为第 2 组气泡之间的尾流夹带。同样，导致从第 2 组气泡生成第 1 组气泡的湍流分解可以被视为剪切效应的一部分，并且可能不会
符号 机制 交互参数

\begin{sourceblock}
\includegraphics[page=295,pagebox=mediabox,trim={45.00bp 479.26bp 44.37bp 92.93bp},clip,width=0.92\textwidth,max height=0.38\textheight]{Ishii.pdf}
\end{sourceblock}
\begin{sourceblock}
\includegraphics[page=295,pagebox=mediabox,trim={45.00bp 410.26bp 44.37bp 190.12bp},clip,width=0.92\textwidth,max height=0.38\textheight]{Ishii.pdf}
\end{sourceblock}
( ) 2 TI φ 湍流冲击 (2)→(2)+(2)

( ) 2 12 , SO φ 剪切 (2)→(2)+(1)

需要单独建模。此外，由于表面不稳定性导致的第2族气泡的崩解被认为非常小，并且可以与湍流冲击引起的第2族气泡的崩解相结合。最终，表 11-2 中列出的九项被视为两组界面面积输运方程中的汇项和源项，应用于中等直径管中的气泡流、段塞流和搅动湍流。

B. 简化的两组界面面积输运方程 这里，假设等温流动条件。另外，忽略系数χ。这是因为在组边界处只有很少部分的气泡可以仅仅由于膨胀而转移到另一组。因此，两组界面面积输运方程简化为

\begin{sourceblock}
\includegraphics[page=295,pagebox=mediabox,trim={45.00bp 95.40bp 44.37bp 507.75bp},clip,width=0.92\textwidth,max height=0.38\textheight]{Ishii.pdf}
\end{sourceblock}
\begin{sourceblock}
\includegraphics[page=296,pagebox=mediabox,trim={45.00bp 545.97bp 44.37bp 55.66bp},clip,width=0.92\textwidth,max height=0.38\textheight]{Ishii.pdf}
\end{sourceblock}
\begin{sourceblock}
\includegraphics[page=296,pagebox=mediabox,trim={45.00bp 460.71bp 44.37bp 139.66bp},clip,width=0.92\textwidth,max height=0.38\textheight]{Ishii.pdf}
\end{sourceblock}
C. 假设的 Group-2 气泡形状和气泡数密度分布
第 2 组气泡由盖泡和块状气泡组成。当直径比 , b cross D D 超过一定限度时，气泡形状会受到壁效应的影响，其中，b cross D 是气泡横截面直径，D 是管直径。此外，还有两个重要参数被认为对于确定气泡形状至关重要。它们是粘度

\begin{sourceblock}
\includegraphics[page=296,pagebox=mediabox,trim={72.09bp 285.97bp 190.58bp 315.97bp},clip,width=\textwidth,max height=0.72\textheight]{Ishii.pdf}
\end{sourceblock}
1 2 /
, b cross D D ≤ 0.6，与在无限介质中一样，壁对帽泡形状造成的变形很小。此时，帽泡的形状可以近似为一段球体，尾流角接近50°。当直径比b cross D D 超过约0.6的值时，管直径成为控制气泡锋面形状的控制长度，然后该气泡被称为段塞气泡。截面半径a、气泡高度h、尾流角wθ等几何参数的定义如图11-4所示。研究表明（Clift 等人，1978），段塞可以被认为由两部分组成，一个是形状与段塞长度无关的圆形鼻部区域，另一个是被环形液体膜包围的近圆柱形部分。还验证了当f N μ ≤ 0.032且* D >10时，粘度和表面张力可以忽略不计，并且可以很好地应用势流理论中的气泡形状。中等直径管道中的空气-水流量（对于50.8毫米直径管道，f N μ =2.36 × 10 -3 且* D ≅ 19）满足上述要求。因此，可以根据伯努利方程的应用来预测气泡形状（Mishima 和 Ishii，1984）。

\begin{center}\small 图11-5给出了简化的气泡数密度分布。假设所有气泡群在相应的气泡体积范围内具有平坦的数密度分布。分布值\end{center}
\begin{sourceblock}
\includegraphics[page=297,pagebox=mediabox,trim={50.00bp 461.28bp 44.37bp 55.00bp},clip,width=\textwidth,max height=0.72\textheight]{Ishii.pdf}
\par\vspace{0.45\baselineskip}
\small 图11-4。盖帽气泡和泰勒气泡的几何参数的定义（Fu 和
\end{sourceblock}
\begin{sourceblock}
\includegraphics[page=297,pagebox=mediabox,trim={50.00bp 224.82bp 44.37bp 218.86bp},clip,width=\textwidth,max height=0.72\textheight]{Ishii.pdf}
\par\vspace{0.45\baselineskip}
\small 图11-5。简化的气泡数密度分布图（Fu 和 Ishii，
\end{sourceblock}
2002a)

对于第 1 组气泡、帽状气泡和段状气泡，函数分别表示为 1 f 、 c f 和 s f。

D. 建模源项和汇项的总结 建模源项和汇项总结如下。在一维公式中，所有两相参数（例如 g α 、 i a 和 Sm D ）都是面积平均值。为了简单起见，在下面的公式中省略了代表面积平均值的符号。

由于随机碰撞而导致气泡合并

\begin{sourceblock}
\includegraphics[page=298,pagebox=mediabox,trim={45.00bp 554.79bp 44.37bp 83.30bp},clip,width=0.92\textwidth,max height=0.38\textheight]{Ishii.pdf}
\end{sourceblock}
\begin{sourceblock}
\includegraphics[page=298,pagebox=mediabox,trim={45.00bp 505.78bp 44.37bp 132.51bp},clip,width=0.92\textwidth,max height=0.38\textheight]{Ishii.pdf}
\end{sourceblock}
\begin{sourceblock}
\includegraphics[page=298,pagebox=mediabox,trim={45.00bp 456.76bp 44.37bp 181.44bp},clip,width=0.92\textwidth,max height=0.38\textheight]{Ishii.pdf}
\end{sourceblock}
\begin{sourceblock}
\includegraphics[page=298,pagebox=mediabox,trim={45.00bp 407.54bp 44.37bp 230.51bp},clip,width=0.92\textwidth,max height=0.38\textheight]{Ishii.pdf}
\end{sourceblock}
where

\begin{sourceblock}
\includegraphics[page=298,pagebox=mediabox,trim={45.00bp 277.22bp 44.37bp 290.47bp},clip,width=0.92\textwidth,max height=0.38\textheight]{Ishii.pdf}
\end{sourceblock}
\begin{sourceblock}
\includegraphics[page=298,pagebox=mediabox,trim={45.00bp 218.08bp 44.37bp 407.72bp},clip,width=0.92\textwidth,max height=0.38\textheight]{Ishii.pdf}
\end{sourceblock}
湍流脉动速度（或均方根速度） 2 t u 由各向同性湍流强度 2 , t isot u 和尾流湍流强度 2 , twake u 组成，如下

\begin{sourceblock}
\includegraphics[page=298,pagebox=mediabox,trim={45.00bp 133.53bp 44.37bp 491.17bp},clip,width=0.92\textwidth,max height=0.38\textheight]{Ishii.pdf}
\end{sourceblock}
\begin{sourceblock}
\includegraphics[page=298,pagebox=mediabox,trim={45.00bp 85.46bp 44.37bp 555.48bp},clip,width=0.92\textwidth,max height=0.38\textheight]{Ishii.pdf}
\end{sourceblock}
\begin{sourceblock}
\includegraphics[page=299,pagebox=mediabox,trim={45.00bp 563.74bp 44.37bp 57.76bp},clip,width=0.92\textwidth,max height=0.38\textheight]{Ishii.pdf}
\end{sourceblock}
where

\begin{sourceblock}
\includegraphics[page=299,pagebox=mediabox,trim={45.00bp 487.78bp 44.37bp 136.55bp},clip,width=0.92\textwidth,max height=0.38\textheight]{Ishii.pdf}
\end{sourceblock}
\begin{sourceblock}
\includegraphics[page=299,pagebox=mediabox,trim={45.00bp 424.72bp 44.37bp 201.50bp},clip,width=0.92\textwidth,max height=0.38\textheight]{Ishii.pdf}
\end{sourceblock}
\begin{sourceblock}
\includegraphics[page=299,pagebox=mediabox,trim={45.00bp 356.41bp 44.37bp 260.44bp},clip,width=0.92\textwidth,max height=0.38\textheight]{Ishii.pdf}
\end{sourceblock}
\begin{sourceblock}
\includegraphics[page=299,pagebox=mediabox,trim={45.00bp 262.33bp 44.37bp 319.52bp},clip,width=0.92\textwidth,max height=0.38\textheight]{Ishii.pdf}
\end{sourceblock}
\begin{sourceblock}
\includegraphics[page=299,pagebox=mediabox,trim={45.00bp 184.69bp 44.37bp 419.79bp},clip,width=0.92\textwidth,max height=0.38\textheight]{Ishii.pdf}
\end{sourceblock}
\begin{sourceblock}
\includegraphics[page=299,pagebox=mediabox,trim={45.00bp 138.03bp 44.37bp 500.09bp},clip,width=0.92\textwidth,max height=0.38\textheight]{Ishii.pdf}
\end{sourceblock}
\begin{sourceblock}
\includegraphics[page=299,pagebox=mediabox,trim={45.00bp 88.95bp 44.37bp 549.17bp},clip,width=0.92\textwidth,max height=0.38\textheight]{Ishii.pdf}
\end{sourceblock}
\begin{sourceblock}
\includegraphics[page=300,pagebox=mediabox,trim={45.00bp 580.36bp 44.37bp 58.27bp},clip,width=0.92\textwidth,max height=0.38\textheight]{Ishii.pdf}
\end{sourceblock}
Otherwise,

\begin{sourceblock}
\includegraphics[page=300,pagebox=mediabox,trim={45.00bp 517.18bp 44.37bp 121.08bp},clip,width=0.92\textwidth,max height=0.38\textheight]{Ishii.pdf}
\end{sourceblock}
\begin{sourceblock}
\includegraphics[page=300,pagebox=mediabox,trim={45.00bp 484.52bp 44.37bp 155.64bp},clip,width=0.92\textwidth,max height=0.38\textheight]{Ishii.pdf}
\end{sourceblock}
\begin{sourceblock}
\includegraphics[page=300,pagebox=mediabox,trim={45.00bp 428.97bp 44.37bp 209.09bp},clip,width=0.92\textwidth,max height=0.38\textheight]{Ishii.pdf}
\end{sourceblock}
\begin{sourceblock}
\includegraphics[page=300,pagebox=mediabox,trim={45.00bp 379.97bp 44.37bp 258.23bp},clip,width=0.92\textwidth,max height=0.38\textheight]{Ishii.pdf}
\end{sourceblock}
\begin{sourceblock}
\includegraphics[page=300,pagebox=mediabox,trim={45.00bp 330.94bp 44.37bp 307.25bp},clip,width=0.92\textwidth,max height=0.38\textheight]{Ishii.pdf}
\end{sourceblock}
\begin{sourceblock}
\includegraphics[page=300,pagebox=mediabox,trim={45.00bp 272.56bp 44.37bp 350.45bp},clip,width=0.92\textwidth,max height=0.38\textheight]{Ishii.pdf}
\end{sourceblock}
\begin{sourceblock}
\includegraphics[page=300,pagebox=mediabox,trim={45.00bp 213.58bp 44.37bp 409.43bp},clip,width=0.92\textwidth,max height=0.38\textheight]{Ishii.pdf}
\end{sourceblock}
\begin{sourceblock}
\includegraphics[page=300,pagebox=mediabox,trim={45.00bp 89.85bp 44.37bp 471.51bp},clip,width=0.92\textwidth,max height=0.38\textheight]{Ishii.pdf}
\end{sourceblock}
\begin{sourceblock}
\includegraphics[page=301,pagebox=mediabox,trim={45.00bp 579.99bp 44.37bp 58.07bp},clip,width=0.92\textwidth,max height=0.38\textheight]{Ishii.pdf}
\end{sourceblock}
\begin{sourceblock}
\includegraphics[page=301,pagebox=mediabox,trim={45.00bp 521.56bp 44.37bp 101.45bp},clip,width=0.92\textwidth,max height=0.38\textheight]{Ishii.pdf}
\end{sourceblock}
where

\begin{sourceblock}
\includegraphics[page=301,pagebox=mediabox,trim={45.00bp 461.92bp 44.37bp 178.93bp},clip,width=0.92\textwidth,max height=0.38\textheight]{Ishii.pdf}
\end{sourceblock}
\begin{sourceblock}
\includegraphics[page=301,pagebox=mediabox,trim={45.00bp 400.66bp 44.37bp 223.75bp},clip,width=0.92\textwidth,max height=0.38\textheight]{Ishii.pdf}
\end{sourceblock}
\begin{sourceblock}
\includegraphics[page=301,pagebox=mediabox,trim={45.00bp 322.28bp 44.37bp 287.58bp},clip,width=0.92\textwidth,max height=0.38\textheight]{Ishii.pdf}
\end{sourceblock}
\begin{sourceblock}
\includegraphics[page=301,pagebox=mediabox,trim={45.00bp 279.03bp 44.37bp 361.88bp},clip,width=0.92\textwidth,max height=0.38\textheight]{Ishii.pdf}
\end{sourceblock}
\begin{sourceblock}
\includegraphics[page=301,pagebox=mediabox,trim={45.00bp 232.05bp 44.37bp 407.12bp},clip,width=0.92\textwidth,max height=0.38\textheight]{Ishii.pdf}
\end{sourceblock}
\begin{sourceblock}
\includegraphics[page=301,pagebox=mediabox,trim={45.00bp 186.31bp 44.37bp 456.51bp},clip,width=0.92\textwidth,max height=0.38\textheight]{Ishii.pdf}
\end{sourceblock}
\begin{sourceblock}
\includegraphics[page=301,pagebox=mediabox,trim={45.00bp 141.34bp 44.37bp 499.58bp},clip,width=0.92\textwidth,max height=0.38\textheight]{Ishii.pdf}
\end{sourceblock}
\begin{sourceblock}
\includegraphics[page=301,pagebox=mediabox,trim={45.00bp 107.72bp 44.37bp 532.44bp},clip,width=0.92\textwidth,max height=0.38\textheight]{Ishii.pdf}
\end{sourceblock}
等式(11-142)至(11-145)和(11-147)。在大多数情况下，段塞气泡的最大横截面空隙率 2 m α 可指定为 0.81，l C 是确定为 0.1 的可调参数。

气泡因湍流冲击而破裂

\begin{sourceblock}
\includegraphics[page=302,pagebox=mediabox,trim={45.00bp 487.51bp 44.37bp 80.97bp},clip,width=0.92\textwidth,max height=0.38\textheight]{Ishii.pdf}
\end{sourceblock}
其中 We 是由下式定义的韦伯数

\begin{sourceblock}
\includegraphics[page=302,pagebox=mediabox,trim={45.00bp 415.20bp 44.37bp 208.49bp},clip,width=0.92\textwidth,max height=0.38\textheight]{Ishii.pdf}
\end{sourceblock}
\begin{sourceblock}
\includegraphics[page=302,pagebox=mediabox,trim={45.00bp 352.25bp 44.37bp 266.67bp},clip,width=0.92\textwidth,max height=0.38\textheight]{Ishii.pdf}
\end{sourceblock}
where

\begin{sourceblock}
\includegraphics[page=302,pagebox=mediabox,trim={45.00bp 221.84bp 44.37bp 345.45bp},clip,width=0.92\textwidth,max height=0.38\textheight]{Ishii.pdf}
\end{sourceblock}
* s V 可由下式确定

\begin{sourceblock}
\includegraphics[page=302,pagebox=mediabox,trim={45.00bp 152.56bp 44.37bp 474.77bp},clip,width=0.92\textwidth,max height=0.38\textheight]{Ishii.pdf}
\end{sourceblock}
实验观察到，对于中小直径管道，与其他三种机制相比，2 ( ) TI φ 项非常小。因此，在应用中可以忽略简化方程（Fu and Ishii，2002b）。

由于剪切而导致气泡破裂

\begin{sourceblock}
\includegraphics[page=303,pagebox=mediabox,trim={45.00bp 558.99bp 44.37bp 82.07bp},clip,width=0.92\textwidth,max height=0.38\textheight]{Ishii.pdf}
\end{sourceblock}
\begin{sourceblock}
\includegraphics[page=303,pagebox=mediabox,trim={45.00bp 512.97bp 44.37bp 128.06bp},clip,width=0.92\textwidth,max height=0.38\textheight]{Ishii.pdf}
\end{sourceblock}
\begin{sourceblock}
\includegraphics[page=303,pagebox=mediabox,trim={45.00bp 466.90bp 44.37bp 174.14bp},clip,width=0.92\textwidth,max height=0.38\textheight]{Ishii.pdf}
\end{sourceblock}
where

\begin{sourceblock}
\includegraphics[page=303,pagebox=mediabox,trim={45.00bp 374.70bp 44.37bp 233.00bp},clip,width=0.92\textwidth,max height=0.38\textheight]{Ishii.pdf}
\end{sourceblock}
\begin{sourceblock}
\includegraphics[page=303,pagebox=mediabox,trim={45.00bp 328.18bp 44.37bp 310.45bp},clip,width=0.92\textwidth,max height=0.38\textheight]{Ishii.pdf}
\end{sourceblock}
\begin{sourceblock}
\includegraphics[page=303,pagebox=mediabox,trim={45.00bp 265.20bp 44.37bp 358.43bp},clip,width=0.92\textwidth,max height=0.38\textheight]{Ishii.pdf}
\end{sourceblock}
\begin{sourceblock}
\includegraphics[page=303,pagebox=mediabox,trim={45.00bp 222.03bp 44.37bp 419.00bp},clip,width=0.92\textwidth,max height=0.38\textheight]{Ishii.pdf}
\end{sourceblock}
\begin{sourceblock}
\includegraphics[page=303,pagebox=mediabox,trim={45.00bp 176.02bp 44.37bp 465.02bp},clip,width=0.92\textwidth,max height=0.38\textheight]{Ishii.pdf}
\end{sourceblock}
其中 g ν 是气相的运动粘度。源项和汇项中的系数值列于表11-2。

\subsection*{11.3.3 Sun 等人建模的源项和汇项。 (2004a)}\phantomsection\addcontentsline{toc}{subsection}{11.3.3 Sun 等人建模的源项和汇项。 (2004a)}
孙等人。 (2004a) 开发了受限通道中气泡流、帽湍流和搅动湍流的两组界面面积输运方程，并使用获取的垂直空气-水流数据对其进行了评估
符号 机制 交互参数

\begin{sourceblock}
\includegraphics[page=304,pagebox=mediabox,trim={45.00bp 433.30bp 44.37bp 92.93bp},clip,width=0.92\textwidth,max height=0.38\textheight]{Ishii.pdf}
\end{sourceblock}
\begin{sourceblock}
\includegraphics[page=304,pagebox=mediabox,trim={45.00bp 364.30bp 44.37bp 236.11bp},clip,width=0.92\textwidth,max height=0.38\textheight]{Ishii.pdf}
\end{sourceblock}
( ) 2 11 , TI φ 湍流冲击 (2)→(1)+(1)

( ) 2 12 , TI φ 湍流冲击 (2)→(1)+(2)

( ) 2 TI φ 湍流冲击 (2)→(2)+(2)

( ) 2 12 , SO φ 剪切 (2)→(2)+(1)

( ) 2 SI φ 表面不稳定性 (2)→(2)+(2)

在宽度 W 为 200 mm、间隙 G 为 10 mm 的矩形通道中。由于测试段宽度较大，在测试段中没有观察到稳定的段塞流状态（Sun et al., 2004a）。接下来，简要解释界面区域传输机制的分类以及建模的源和汇项。

A. 界面区域传输机制的分类 Sun 等人。 （2004a）采用了五种主要的气泡相互作用：（1）由于湍流驱动的随机碰撞而合并； (2)尾流夹带引起的聚结； (3)湍流涡流冲击造成的破碎； （四）因剪断造成的断裂； (5)由于气泡表面流动不稳定而导致大帽气泡破裂。表 11-3 中列出的 14 项被视为源项和汇项
B. 简化的两组界面面积输运方程 这里，假设等温流动条件。因此，两组界面面积输运方程简化为

\begin{sourceblock}
\includegraphics[page=305,pagebox=mediabox,trim={45.00bp 436.48bp 44.37bp 148.19bp},clip,width=0.92\textwidth,max height=0.38\textheight]{Ishii.pdf}
\end{sourceblock}
\begin{sourceblock}
\includegraphics[page=305,pagebox=mediabox,trim={45.00bp 324.82bp 44.37bp 245.07bp},clip,width=0.92\textwidth,max height=0.38\textheight]{Ishii.pdf}
\end{sourceblock}
\begin{sourceblock}
\includegraphics[page=305,pagebox=mediabox,trim={45.00bp 202.72bp 44.37bp 360.93bp},clip,width=0.92\textwidth,max height=0.38\textheight]{Ishii.pdf}
\end{sourceblock}
C. 假设的第 2 组气泡形状和气泡数密度分布
为了对源项和汇项进行分析建模，应简化气泡形状。对于第 1 组气泡，可以假定为球形。然而，对于第 2 组气泡，应考虑感兴趣的测试部分的独特几何形状。由于在大气压下的绝热空气-水系统中，由式(11-115)定义的第1组和第2组气泡之间的边界约为10mm，因此假设帽泡夹在两个平行的平坦壁之间，使得帽泡的厚度为G，如图11-6所示。这里，R 和 2a

\begin{sourceblock}
\includegraphics[page=306,pagebox=mediabox,trim={50.00bp 479.76bp 44.37bp 55.00bp},clip,width=\textwidth,max height=0.72\textheight]{Ishii.pdf}
\par\vspace{0.45\baselineskip}
\small 图 11-6。帽泡几何参数的定义（Sun 等人，2004a）
\end{sourceblock}
\begin{sourceblock}
\includegraphics[page=306,pagebox=mediabox,trim={50.00bp 288.36bp 44.37bp 200.38bp},clip,width=\textwidth,max height=0.72\textheight]{Ishii.pdf}
\par\vspace{0.45\baselineskip}
\small 图 11-7。简化的气泡数密度分布图（Sun 等人，
\end{sourceblock}
2004a)

是帽泡的曲率半径和底部宽度，w θ 是尾流角。考虑到 Clift 等人给出的尾流角相关性，假定 50° 尾流角是感兴趣的流动条件下所有第 2 组气泡的合理近似值。 （1978）。此外，鉴于受限流的特征以及基础物理，
选择2a作为确定群边界和最大稳定气泡尺寸的特征长度。

\begin{center}\small 图11-7给出了简化的气泡数密度分布。假设所有气泡群在相应的气泡体积范围内具有平坦的数密度分布。第 1 组气泡和第 2 组气泡的分布函数值分别表示为 1 f 和 2 f。图中，min V 为系统中最小气泡的体积，c,max V 为最大稳定气泡的体积，对应于\end{center}
\begin{sourceblock}
\includegraphics[page=307,pagebox=mediabox,trim={45.00bp 570.43bp 44.37bp 54.47bp},clip,width=0.92\textwidth,max height=0.38\textheight]{Ishii.pdf}
\end{sourceblock}
然而，应该注意的是，系统中可能不存在这种尺寸的气泡，并且 c,max D 仅提供可能的最大气泡尺寸的上限，超过该上限，则假定气泡会立即破裂。 1,max V 和2,max V 分别是在给定流动条件下，假设均匀气泡数密度分布，组1 和组2 的最大气泡体积。此外，c V 是第 1 组和第 2 组气泡边界处的临界气泡体积，对应于临界 D，超过该临界值，气泡会变成盖形并被归类为第 2 组气泡。对于窄通道，Group-1 和 Group-2 气泡之间的边界可以通过以下方式确定：

\begin{sourceblock}
\includegraphics[page=307,pagebox=mediabox,trim={45.00bp 376.24bp 44.37bp 246.11bp},clip,width=0.92\textwidth,max height=0.38\textheight]{Ishii.pdf}
\end{sourceblock}
其中 crit D 是第 1 组和第 2 组气泡之间边界处的气泡体积当量直径。

D. 模拟源项和汇项的总结 由于随机碰撞而导致的气泡合并

\begin{sourceblock}
\includegraphics[page=307,pagebox=mediabox,trim={45.00bp 179.74bp 44.37bp 389.42bp},clip,width=0.92\textwidth,max height=0.38\textheight]{Ishii.pdf}
\end{sourceblock}
\begin{sourceblock}
\includegraphics[page=307,pagebox=mediabox,trim={45.00bp 61.30bp 44.37bp 504.80bp},clip,width=0.92\textwidth,max height=0.38\textheight]{Ishii.pdf}
\end{sourceblock}
\begin{sourceblock}
\includegraphics[page=308,pagebox=mediabox,trim={45.00bp 547.95bp 44.37bp 93.92bp},clip,width=0.92\textwidth,max height=0.38\textheight]{Ishii.pdf}
\end{sourceblock}
对于Sun等人的应用。 （2004a）。

\begin{sourceblock}
\includegraphics[page=308,pagebox=mediabox,trim={45.00bp 357.97bp 44.37bp 153.35bp},clip,width=0.92\textwidth,max height=0.38\textheight]{Ishii.pdf}
\end{sourceblock}
\begin{sourceblock}
\includegraphics[page=308,pagebox=mediabox,trim={45.00bp 239.69bp 44.37bp 326.00bp},clip,width=0.92\textwidth,max height=0.38\textheight]{Ishii.pdf}
\end{sourceblock}
\begin{sourceblock}
\includegraphics[page=308,pagebox=mediabox,trim={45.00bp 146.13bp 44.37bp 442.19bp},clip,width=0.92\textwidth,max height=0.38\textheight]{Ishii.pdf}
\end{sourceblock}
where

\begin{sourceblock}
\includegraphics[page=308,pagebox=mediabox,trim={45.00bp 73.30bp 44.37bp 550.85bp},clip,width=0.92\textwidth,max height=0.38\textheight]{Ishii.pdf}
\end{sourceblock}
\begin{sourceblock}
\includegraphics[page=309,pagebox=mediabox,trim={45.00bp 496.96bp 44.37bp 58.57bp},clip,width=0.92\textwidth,max height=0.38\textheight]{Ishii.pdf}
\end{sourceblock}
\begin{sourceblock}
\includegraphics[page=309,pagebox=mediabox,trim={45.00bp 376.27bp 44.37bp 189.79bp},clip,width=0.92\textwidth,max height=0.38\textheight]{Ishii.pdf}
\end{sourceblock}
由于尾流夹带而导致气泡合并

\begin{sourceblock}
\includegraphics[page=309,pagebox=mediabox,trim={45.00bp 307.12bp 44.37bp 333.74bp},clip,width=0.92\textwidth,max height=0.38\textheight]{Ishii.pdf}
\end{sourceblock}
\begin{sourceblock}
\includegraphics[page=309,pagebox=mediabox,trim={45.00bp 250.07bp 44.37bp 373.35bp},clip,width=0.92\textwidth,max height=0.38\textheight]{Ishii.pdf}
\end{sourceblock}
\begin{sourceblock}
\includegraphics[page=309,pagebox=mediabox,trim={45.00bp 125.09bp 44.37bp 433.89bp},clip,width=0.92\textwidth,max height=0.38\textheight]{Ishii.pdf}
\end{sourceblock}
\begin{sourceblock}
\includegraphics[page=309,pagebox=mediabox,trim={45.00bp 60.29bp 44.37bp 558.63bp},clip,width=0.92\textwidth,max height=0.38\textheight]{Ishii.pdf}
\end{sourceblock}
\begin{sourceblock}
\includegraphics[page=310,pagebox=mediabox,trim={45.00bp 566.74bp 44.37bp 56.27bp},clip,width=0.92\textwidth,max height=0.38\textheight]{Ishii.pdf}
\end{sourceblock}
\begin{sourceblock}
\includegraphics[page=310,pagebox=mediabox,trim={45.00bp 508.20bp 44.37bp 118.34bp},clip,width=0.92\textwidth,max height=0.38\textheight]{Ishii.pdf}
\end{sourceblock}
\begin{sourceblock}
\includegraphics[page=310,pagebox=mediabox,trim={45.00bp 450.28bp 44.37bp 172.73bp},clip,width=0.92\textwidth,max height=0.38\textheight]{Ishii.pdf}
\end{sourceblock}
气泡因湍流冲击而破裂

\begin{sourceblock}
\includegraphics[page=310,pagebox=mediabox,trim={45.00bp 294.82bp 44.37bp 260.77bp},clip,width=0.92\textwidth,max height=0.38\textheight]{Ishii.pdf}
\end{sourceblock}
\begin{sourceblock}
\includegraphics[page=310,pagebox=mediabox,trim={45.00bp 189.59bp 44.37bp 389.75bp},clip,width=0.92\textwidth,max height=0.38\textheight]{Ishii.pdf}
\end{sourceblock}
\begin{sourceblock}
\includegraphics[page=310,pagebox=mediabox,trim={45.00bp 93.32bp 44.37bp 497.25bp},clip,width=0.92\textwidth,max height=0.38\textheight]{Ishii.pdf}
\end{sourceblock}
\begin{sourceblock}
\includegraphics[page=311,pagebox=mediabox,trim={45.00bp 524.52bp 44.37bp 58.22bp},clip,width=0.92\textwidth,max height=0.38\textheight]{Ishii.pdf}
\end{sourceblock}
由于剪切而导致气泡破裂

\begin{sourceblock}
\includegraphics[page=311,pagebox=mediabox,trim={45.00bp 424.15bp 44.37bp 186.28bp},clip,width=0.92\textwidth,max height=0.38\textheight]{Ishii.pdf}
\end{sourceblock}
其中 rb v 和 , crit SO We 分别是大气泡相对于帽泡底部附近液膜的相对速度和临界韦伯数，2 ,max We 定义为

\begin{sourceblock}
\includegraphics[page=311,pagebox=mediabox,trim={45.00bp 325.88bp 44.37bp 275.26bp},clip,width=0.92\textwidth,max height=0.38\textheight]{Ishii.pdf}
\end{sourceblock}
在圆管中的向上流动中，当大帽或块状气泡上升时，液相被推开并作为气泡侧界面和壁之间的液膜向下流动。然而，在Sun等人考虑的流道中。 (2004a)，当大帽气泡上升时，气泡侧界面和壁之间的液膜可能几乎保持静止，因为在流道宽度方向上有更多的自由空间可用于液相。当帽泡速度高且发生剪切时，情况可能更是如此。鉴于此，盖泡相对于气泡底部周围液膜的相对速度 rb v 可以通过第 2 组气泡在主流方向上的速度来估计。

\begin{sourceblock}
\includegraphics[page=311,pagebox=mediabox,trim={45.00bp 66.86bp 44.37bp 500.52bp},clip,width=0.92\textwidth,max height=0.38\textheight]{Ishii.pdf}
\end{sourceblock}
\begin{sourceblock}
\includegraphics[page=312,pagebox=mediabox,trim={45.00bp 514.34bp 44.37bp 57.33bp},clip,width=0.92\textwidth,max height=0.38\textheight]{Ishii.pdf}
\end{sourceblock}
由于表面不稳定而导致气泡破裂

\begin{sourceblock}
\includegraphics[page=312,pagebox=mediabox,trim={45.00bp 394.35bp 44.37bp 193.77bp},clip,width=0.92\textwidth,max height=0.38\textheight]{Ishii.pdf}
\end{sourceblock}
源项和汇项中的系数值列于表11-3。

组边界处的组间传递系数 组边界处的组间传递系数通过实验确定为

\begin{sourceblock}
\includegraphics[page=312,pagebox=mediabox,trim={45.00bp 251.57bp 44.37bp 369.93bp},clip,width=0.92\textwidth,max height=0.38\textheight]{Ishii.pdf}
\end{sourceblock}
这种相关性是基于有限的实验数据库获得的（Sun et al., 2004a）。然而，在一般的帽湍流和搅动湍流中，1 g α 的值通常在 0.05 到 0.40 之间。因此，该相关性可能适用于大多数这些流动条件，因为开发该相关性的数据库涵盖了类似的范围

\begin{sourceblock}
\includegraphics[page=312,pagebox=mediabox,trim={45.00bp 156.58bp 44.37bp 484.65bp},clip,width=0.92\textwidth,max height=0.38\textheight]{Ishii.pdf}
\end{sourceblock}
\section*{11.4 界面区域输运方程中相变项的建模}\phantomsection\addcontentsline{toc}{section}{11.4 界面区域输运方程中相变项的建模}
要在沸腾流中使用界面面积输运方程，应考虑由于相变而产生的源项和汇项。相变项由壁成核项和体成核项组成或

\begin{sourceblock}
\includegraphics[page=313,pagebox=mediabox,trim={45.00bp 531.57bp 44.37bp 68.69bp},clip,width=0.92\textwidth,max height=0.38\textheight]{Ishii.pdf}
\end{sourceblock}
其中 H ψ 和 C A 分别是沸腾通道的加热周长和横截面积。从式（11-207）中可以清楚地看出，活性成核位点密度n N 、气泡离开频率f 和气泡离开直径d D 模型对于计算壁成核源项是必不可少的。这些项的建模目前是沸腾流研究中的一个具有挑战性的课题。下面简要描述了最近对这些术语的建模工作。

\subsection*{11.4.1 Kocamustafaogullari 和 Ishii (1983) 以及 Hibiki 和 Ishii (2003b) 模拟的活性成核位点密度}\phantomsection\addcontentsline{toc}{subsection}{11.4.1 Kocamustafaogullari 和 Ishii (1983) 以及 Hibiki 和 Ishii (2003b) 模拟的活性成核位点密度}
A. Kocamustafaogullari 和 Ishii 的模型 (1983) Kocamustafaogullari 和 Ishii (1983) 通过参数研究关联了现有的活性成核位点密度数据。他们假设池沸腾 np N 中的活性成核位点密度受到表面条件和流体热物理性质的影响，并根据水数据建立了以下相关性。

\begin{sourceblock}
\includegraphics[page=313,pagebox=mediabox,trim={45.00bp 251.19bp 44.37bp 389.04bp},clip,width=0.92\textwidth,max height=0.38\textheight]{Ishii.pdf}
\end{sourceblock}
where

\begin{sourceblock}
\includegraphics[page=313,pagebox=mediabox,trim={45.00bp 191.71bp 44.37bp 449.33bp},clip,width=0.92\textwidth,max height=0.38\textheight]{Ishii.pdf}
\end{sourceblock}
\begin{sourceblock}
\includegraphics[page=313,pagebox=mediabox,trim={45.00bp 141.51bp 44.37bp 496.86bp},clip,width=0.92\textwidth,max height=0.38\textheight]{Ishii.pdf}
\end{sourceblock}
\begin{sourceblock}
\includegraphics[page=313,pagebox=mediabox,trim={45.00bp 79.33bp 44.37bp 544.40bp},clip,width=0.92\textwidth,max height=0.38\textheight]{Ishii.pdf}
\end{sourceblock}
\begin{sourceblock}
\includegraphics[page=314,pagebox=mediabox,trim={45.00bp 530.59bp 44.37bp 96.35bp},clip,width=0.92\textwidth,max height=0.38\textheight]{Ishii.pdf}
\end{sourceblock}
Kocamustafaogullari 和 Ishii (1983) 给出了气泡离开直径
d D 通过将 Fritz 提出的气泡离开直径 dF D 的相关性修改为

\begin{sourceblock}
\includegraphics[page=314,pagebox=mediabox,trim={45.00bp 443.66bp 44.37bp 197.40bp},clip,width=0.92\textwidth,max height=0.38\textheight]{Ishii.pdf}
\end{sourceblock}
where

\begin{sourceblock}
\includegraphics[page=314,pagebox=mediabox,trim={45.00bp 369.40bp 44.37bp 255.06bp},clip,width=0.92\textwidth,max height=0.38\textheight]{Ishii.pdf}
\end{sourceblock}
\begin{sourceblock}
\includegraphics[page=314,pagebox=mediabox,trim={45.00bp 261.56bp 44.37bp 317.42bp},clip,width=0.92\textwidth,max height=0.38\textheight]{Ishii.pdf}
\end{sourceblock}
其中 f p 、fg i 和 R 分别是液体压力、潜热和基于分子量的气体常数。例如，水蒸气的R值为462 J/(kg·K)。对于 g f ρ ρ 和 ( ) 1 fg sat g sat i T RTT Δ，式(11-215)简化为

\begin{sourceblock}
\includegraphics[page=314,pagebox=mediabox,trim={45.00bp 152.48bp 44.37bp 448.27bp},clip,width=0.92\textwidth,max height=0.38\textheight]{Ishii.pdf}
\end{sourceblock}
这里应该注意的是，对于沸腾的池，近似值
( ) sat w w sat T T T T Δ Δ ≈ ≡ −，其中 w T 为壁温，用于式(11-215)。

考虑有效液体过热度之间的差异

e T Δ 和壁过热度以及池沸腾和对流沸腾之间的机制相似性，Kocamustafaogullari 和 Ishii 假设池沸腾形成了活性成核位点密度相关性，

\begin{sourceblock}
\includegraphics[page=315,pagebox=mediabox,trim={45.00bp 534.74bp 44.37bp 106.33bp},clip,width=0.92\textwidth,max height=0.38\textheight]{Ishii.pdf}
\end{sourceblock}
其中 S 是抑制因子（Chen，1966）。有效临界空腔尺寸* ce R 可通过式（11-215）和（11-217）计算。因此，对流系统中的无量纲活性成核位点密度 nc N 表示为

\begin{sourceblock}
\includegraphics[page=315,pagebox=mediabox,trim={45.00bp 436.64bp 44.37bp 201.24bp},clip,width=0.92\textwidth,max height=0.38\textheight]{Ishii.pdf}
\end{sourceblock}
B. Hibiki 和 Ishii 的模型 (2003b) 现有实验数据表明，壁面过热度的功率
随着壁过热度的增加，w T Δ（或无量纲最小空腔尺寸的幂 * c R ）的范围从 2 到 6（或 -2 到 -6）（Hibiki 和 Ishii，2003b）。为了解释功率对壁过热度的依赖性，Hibiki 和 Ishii (2003b) 使用表面上实际存在的空腔的尺寸和锥角分布对活性成核位点密度 n N 进行机械建模：

\begin{sourceblock}
\includegraphics[page=315,pagebox=mediabox,trim={45.00bp 263.18bp 44.37bp 366.93bp},clip,width=0.92\textwidth,max height=0.38\textheight]{Ishii.pdf}
\end{sourceblock}
其中 n N 、 β 、 max R 和 r 分别是平均空腔密度、锥角的一半、最大空腔尺寸和空腔半径。腔体半径和锥角一半的腔体数量分布函数近似为

\begin{sourceblock}
\includegraphics[page=315,pagebox=mediabox,trim={45.00bp 151.79bp 44.37bp 473.57bp},clip,width=0.92\textwidth,max height=0.38\textheight]{Ishii.pdf}
\end{sourceblock}
and

\begin{sourceblock}
\includegraphics[page=315,pagebox=mediabox,trim={45.00bp 80.25bp 44.37bp 544.19bp},clip,width=0.92\textwidth,max height=0.38\textheight]{Ishii.pdf}
\end{sourceblock}
\begin{sourceblock}
\includegraphics[page=316,pagebox=mediabox,trim={45.00bp 508.91bp 44.37bp 107.08bp},clip,width=0.92\textwidth,max height=0.38\textheight]{Ishii.pdf}
\end{sourceblock}
如果我们使用最大型腔尺寸，则该表达式是通用的。然而，这导致活性成核位点密度的高度复杂的相关性。因此，需要获得一个满足整体物理现象的简单表达式。

\begin{sourceblock}
\includegraphics[page=316,pagebox=mediabox,trim={45.00bp 351.35bp 44.37bp 238.41bp},clip,width=0.92\textwidth,max height=0.38\textheight]{Ishii.pdf}
\end{sourceblock}
其次，考虑极大 sat T Δ 的情况，即 0 m c R =，导致 ( ) exp c R λ = ∞。因此，对于极大的 sat T Δ，活性成核位点密度变得极大，这在物理上是合理的。因此，从这两个限制条件可以得出，n N 的一般表达式可以近似由式（11-223）给出。在计算活性成核位点密度时，任何温度下的接触角值可能不容易获得。因此，我们将在最终确定方程（11-223）时使用室温下的接触角R θ 值。这种近似产生的误差将通过在方程（11-223）中引入修正因子（ ） f ρ + 来补偿：

\begin{sourceblock}
\includegraphics[page=316,pagebox=mediabox,trim={45.00bp 155.32bp 44.37bp 435.94bp},clip,width=0.92\textwidth,max height=0.38\textheight]{Ishii.pdf}
\end{sourceblock}
\begin{sourceblock}
\includegraphics[page=316,pagebox=mediabox,trim={45.00bp 115.09bp 44.37bp 487.64bp},clip,width=0.92\textwidth,max height=0.38\textheight]{Ishii.pdf}
\end{sourceblock}
\begin{sourceblock}
\includegraphics[page=316,pagebox=mediabox,trim={45.00bp 66.18bp 44.37bp 554.23bp},clip,width=0.92\textwidth,max height=0.38\textheight]{Ishii.pdf}
\end{sourceblock}
\begin{sourceblock}
\includegraphics[page=317,pagebox=mediabox,trim={50.00bp 375.84bp 44.37bp 55.00bp},clip,width=\textwidth,max height=0.72\textheight]{Ishii.pdf}
\par\vspace{0.45\baselineskip}
\small 图 11-8。活性成核位点密度模型与Basu等人的数据的比较。
\end{sourceblock}
（响和石井，2003b）

where

\begin{sourceblock}
\includegraphics[page=317,pagebox=mediabox,trim={45.00bp 260.17bp 44.37bp 359.57bp},clip,width=0.92\textwidth,max height=0.38\textheight]{Ishii.pdf}
\end{sourceblock}
这里需要说明的是，式（11-224）中的气体过热度近似为壁面过热度。巴苏等人。 (2002) 通过实验证明，对于给定的液体-表面组合，质量速度和液体过冷度对活性成核位点密度没有影响。图11-8显示了活性成核位点模型与Basu等人的数据的比较。 （2002）。图中，实线和虚线分别表示 Hibiki 和 Ishii 模型（2003b）以及 Kocamustafaogullari 和 Ishii 模型（1983）的预测。可以看出，Hibiki 和 Ishii 的模型在高壁过热度区域与 Kocamustafaogullari 和 Ishii 的模型一致，并且正确地表示了低壁过热度区域中活性成核位点密度对壁过热度的依赖性。

\begin{sourceblock}
\includegraphics[page=318,pagebox=mediabox,trim={50.00bp 370.98bp 44.37bp 55.00bp},clip,width=\textwidth,max height=0.72\textheight]{Ishii.pdf}
\par\vspace{0.45\baselineskip}
\small 图11-9。成核位置处蒸汽泡的力平衡（Situ 等人，2008）
\end{sourceblock}
\subsection*{11.4.2 Situ 等人建模的气泡偏离尺寸。 (2008)}\phantomsection\addcontentsline{toc}{subsection}{11.4.2 Situ 等人建模的气泡偏离尺寸。 (2008)}
目前尚未建立准确、适用范围广的气泡脱离尺寸机理模型。司徒等人。 (2008) 通过考虑作用在气泡上的力的平衡，对气泡偏离尺寸进行了初步建模。图 11-9 显示了作用在气泡成核位置处的力的平衡（Situ 等人，2005 年；2008 年）。这里
sx F 、dux F 、sl F 、sy F 、duy F 、p F 、g F 和 qs F 分别是 x 方向的表面张力、x 方向的非定常曳力（生长力）、剪切升力、y 方向的表面张力、y 方向的非定常曳力、压力、重力和准定常力。当气泡违反沿流动方向的力平衡时，气泡就会离开：

\begin{sourceblock}
\includegraphics[page=318,pagebox=mediabox,trim={45.00bp 113.56bp 44.37bp 520.66bp},clip,width=0.92\textwidth,max height=0.38\textheight]{Ishii.pdf}
\end{sourceblock}
在气泡离开的瞬间，由于气泡在壁上的接触面积变为零，因此表面张力可以忽略不计。流动方向上的生长力由下式给出（Klausner et al., 1993）

\begin{sourceblock}
\includegraphics[page=319,pagebox=mediabox,trim={45.00bp 570.73bp 44.37bp 55.33bp},clip,width=0.92\textwidth,max height=0.38\textheight]{Ishii.pdf}
\end{sourceblock}
其中 br 、 br 、 br 和 i θ 分别是气泡半径、气泡半径对时间的导数、气泡半径对时间的二阶导数和倾斜角 (= 18 π )。采用修正的 Zuber 气泡生长方程（Zuber，1961），可以推导出生长力为

\begin{sourceblock}
\includegraphics[page=319,pagebox=mediabox,trim={45.00bp 444.51bp 44.37bp 179.25bp},clip,width=0.92\textwidth,max height=0.38\textheight]{Ishii.pdf}
\end{sourceblock}
其中 b 和 f α 分别是常数和热扩散率。有效雅各布数定义为

\begin{sourceblock}
\includegraphics[page=319,pagebox=mediabox,trim={45.00bp 358.78bp 44.37bp 264.27bp},clip,width=0.92\textwidth,max height=0.38\textheight]{Ishii.pdf}
\end{sourceblock}
其中 pf c 、S 和 w T 分别为恒压比热、考虑液体流动和壁温影响的抑制因子。压力、重力和准稳态力 (Klausner et al., 1993) 给出为

\begin{sourceblock}
\includegraphics[page=319,pagebox=mediabox,trim={45.00bp 235.75bp 44.37bp 393.46bp},clip,width=0.92\textwidth,max height=0.38\textheight]{Ishii.pdf}
\end{sourceblock}
and

\begin{sourceblock}
\includegraphics[page=319,pagebox=mediabox,trim={45.00bp 152.44bp 44.37bp 462.26bp},clip,width=0.92\textwidth,max height=0.38\textheight]{Ishii.pdf}
\end{sourceblock}
其中 r v 、f ν 、Re b N 和 n 是相间的相对速度、液体的运动粘度、气泡雷诺数和常数 (=0.65)。

\subsection*{11.4.3 Euh 等人建模的气泡离开频率。 (2010)}\phantomsection\addcontentsline{toc}{subsection}{11.4.3 Euh 等人建模的气泡离开频率。 (2010)}
目前还没有建立一个准确、适用范围广的气泡脱离频率的机理模型。尤等人。 （2010）通过无量纲分析并使用可用数据库开发了气泡离开频率的初步相关性

\begin{sourceblock}
\includegraphics[page=320,pagebox=mediabox,trim={45.00bp 493.90bp 44.37bp 148.88bp},clip,width=0.92\textwidth,max height=0.38\textheight]{Ishii.pdf}
\end{sourceblock}
无量纲气泡离开频率 fd N 和代表核态沸腾传热的无量纲热通量 qNB N 定义为

\begin{sourceblock}
\includegraphics[page=320,pagebox=mediabox,trim={45.00bp 392.68bp 44.37bp 233.18bp},clip,width=0.92\textwidth,max height=0.38\textheight]{Ishii.pdf}
\end{sourceblock}
and

\begin{sourceblock}
\includegraphics[page=320,pagebox=mediabox,trim={45.00bp 317.50bp 44.37bp 304.47bp},clip,width=0.92\textwidth,max height=0.38\textheight]{Ishii.pdf}
\end{sourceblock}
其中 t α 和 qNB q '' 是热扩散率和核态沸腾热通量。

\subsection*{11.4.4 由 Park 等人建模的冷凝引起的汇项。 (2007)}\phantomsection\addcontentsline{toc}{subsection}{11.4.4 由 Park 等人建模的冷凝引起的汇项。 (2007)}
帕克等人。 (2007) 通过假设由于壁面沸腾引起的气泡成核和由于冷凝引起的气泡破裂是对称现象，对由于本体冷凝引起的汇项进行了建模。基于这一假设，过冷条件下的冷凝区域可以分为两个区域：传热控制区域和惯性控制区域。在传热控制区域，冷凝努塞尔数方法是合适的，而在惯性控制区域，气泡会在短时间内因周围液体的惯性而破裂。本体冷凝汇项 BC φ 的模型解释如下。

A. 气泡破裂的简化机制 由于沸腾而产生气泡和由于冷凝而导致气泡消失是对称现象。图 11-10 显示了一个典型的

\begin{sourceblock}
\includegraphics[page=321,pagebox=mediabox,trim={50.00bp 377.76bp 44.37bp 55.00bp},clip,width=\textwidth,max height=0.72\textheight]{Ishii.pdf}
\par\vspace{0.45\baselineskip}
\small 图 11-10。沸腾和冷凝过程中气泡尺寸的变化（Park 等人，
\end{sourceblock}
2007)

沸腾和冷凝过程中气泡尺寸的变化（Zuber，1961）。气泡从加热表面突然产生（A→B），并且初始气泡生长非常快（B→C），但随着尺寸增加，生长速度减慢（C→D）。当气泡达到最大直径 (D) 并离开过热区域时，它开始在环境过冷液体中塌陷 (D → E)。在低过冷度下，气泡破裂主要受传热控制，但随着气泡尺寸的减小（E→F）和液体过冷度的增加，气泡不再能维持其界面边界，并在短时间内（F→G）因惯性而破裂。根据之前对过冷液体中气泡运动行为的研究，假设气泡的生命周期中有四个区域。区域I是环境过热条件下的气泡产生区域，区域II是过热条件下的气泡生长区域，区域III是过冷条件下的气泡收缩区域，区域IV是过冷条件下的气泡破裂区域。 I区和IV区的居住时间很短。图11-10中，B D 为核沸腾过程中惯性控制气泡生长过程中B点气泡直径，F D 为气泡凝结过程中气泡受惯性控制时F点气泡直径，G D 为G点临界溃灭气泡直径。

如图11-10所示，过冷工况下的冷凝区域可分为两个区域：传热控制区域（区域III）和惯性控制区域（区域IV）。在传热控制区域，努塞尔数方法是合适的。然而，在惯性控制区域，气泡会在短时间内因周围液体的惯性而破裂。

当饱和条件下的蒸汽泡被环境过冷液体包围时，通过界面发生冷凝。最初，冷凝现象是由通过界面的传热控制的。随着冷凝过程继续通过蒸汽泡和环境液体之间的界面，气泡直径减小，气泡内的蒸汽温度也减小。当气泡尺寸达到边界气泡尺寸 (F) 时，尺寸不会因通过界面的传热而逐渐减小，而是会因机械惯性而非常迅速地崩溃。气泡的任意截止直径可用于识别气泡的突然破裂和气泡尺寸的逐渐减小。在惯性控制区域的末端，临界破裂气泡尺寸可以根据力平衡中的液体温度来计算。当气泡尺寸减小到临界破裂气泡直径（G）后，它就会消失。为了在冷凝过程中保持气泡形状，气泡与周围液体之间需要有一定的压力差。如果压力差相当小，气泡就无法再存在并消失。临界破裂气泡尺寸 G D 可以根据蒸汽泡和环境液体之间的压力平衡计算得出（Park 等人，2007）

\begin{sourceblock}
\includegraphics[page=322,pagebox=mediabox,trim={45.00bp 192.09bp 44.37bp 404.73bp},clip,width=0.92\textwidth,max height=0.38\textheight]{Ishii.pdf}
\end{sourceblock}
B.区域识别和边界气泡直径传热控制区域和惯性控制区域之间的过渡标准可以根据先前的理论和实验观察确定（Rayleigh，1917；Plesset，1954；Zwick 和 Plesset，1955）。对于接近崩溃结束的快速崩溃的气泡来说，球形是不稳定的（Plesset，1954）。根据对液体中塌陷蒸汽泡的分析（Rayleigh，1917；Zwick 和 Plesset，1955），气泡尺寸在其寿命结束时迅速减小。假设气泡尺寸开始迅速减小后，气泡处于惯性控制区域。

通过初始气泡直径* b D 无量纲化的气泡直径用于确定将惯性控制区域与传热控制区域分开的标准。根据瑞利解 (1917)，当无量纲气泡直径在 0.4 到 0.6 之间时，气泡尺寸会发生快速变化。因此，假设区域边界处的无量纲气泡直径为 0.4。 Rayleigh 解（1917）的气泡破裂时间比 Zwick 和 Plesset（1955）的解要短一些。这是因为后者考虑了传热效应。

C.传热控制区域中的气泡破裂时间气泡破裂时间是从通过气泡界面的能量平衡得出的。塌陷气泡界面的能量平衡方程由下式给出

\begin{sourceblock}
\includegraphics[page=323,pagebox=mediabox,trim={45.00bp 391.96bp 44.37bp 249.27bp},clip,width=0.92\textwidth,max height=0.38\textheight]{Ishii.pdf}
\end{sourceblock}
其中 cm 、 c h 和 b A 分别是冷凝传质速率、冷凝传热系数和气泡界面面积。式(11-237)左边和右边的项分别表示为式(11-238)和(11-239)。

\begin{sourceblock}
\includegraphics[page=323,pagebox=mediabox,trim={45.00bp 281.31bp 44.37bp 344.73bp},clip,width=0.92\textwidth,max height=0.38\textheight]{Ishii.pdf}
\end{sourceblock}
and

\begin{sourceblock}
\includegraphics[page=323,pagebox=mediabox,trim={45.00bp 226.96bp 44.37bp 414.27bp},clip,width=0.92\textwidth,max height=0.38\textheight]{Ishii.pdf}
\end{sourceblock}
其中 b V 、 f k 、 ( ) sub sat f T T T Δ ≡ − 和 Nuc N 分别是气泡的体积、热导率、液体过冷度和冷凝努塞尔数。

将式（11-238）和（11-239）代入式（11-237）得

\begin{sourceblock}
\includegraphics[page=323,pagebox=mediabox,trim={45.00bp 112.54bp 44.37bp 510.51bp},clip,width=0.92\textwidth,max height=0.38\textheight]{Ishii.pdf}
\end{sourceblock}
式(11-240)从初始气泡直径0 D或图11-10中D点气泡直径到边界气泡直径的积分
F D 给出传热控制区域 c t 中的气泡停留时间：

\begin{sourceblock}
\includegraphics[page=324,pagebox=mediabox,trim={45.00bp 567.76bp 44.37bp 55.29bp},clip,width=0.92\textwidth,max height=0.38\textheight]{Ishii.pdf}
\end{sourceblock}
D. 冷凝汇项这里我们考虑两个冷凝汇项，这是由于气泡与周围液体之间的界面热传递 HC φ for b F D D > 以及通过气泡-液体界面 IC φ 的机械力平衡
b F D D < .前一项表示气泡数密度恒定且气泡尺寸变化时界面面积浓度的变化，后一项则说明由于整个气泡破裂导致的数密度变化导致的界面面积浓度变化。惯性控制区域 ( ) c b F p D D < 中尺寸小于边界气泡尺寸的气泡比例是计算冷凝汇项的重要参数。

\begin{sourceblock}
\includegraphics[page=324,pagebox=mediabox,trim={45.00bp 310.72bp 44.37bp 287.51bp},clip,width=0.92\textwidth,max height=0.38\textheight]{Ishii.pdf}
\end{sourceblock}
无量纲时间 * t 与无量纲气泡直径 * b D 之间的关系表示为（Rayleigh，1917）

\begin{sourceblock}
\includegraphics[page=324,pagebox=mediabox,trim={45.00bp 204.92bp 44.37bp 406.04bp},clip,width=0.92\textwidth,max height=0.38\textheight]{Ishii.pdf}
\end{sourceblock}
where

\begin{sourceblock}
\includegraphics[page=324,pagebox=mediabox,trim={45.00bp 130.78bp 44.37bp 496.70bp},clip,width=0.92\textwidth,max height=0.38\textheight]{Ishii.pdf}
\end{sourceblock}
\begin{sourceblock}
\includegraphics[page=324,pagebox=mediabox,trim={45.00bp 101.74bp 44.37bp 515.63bp},clip,width=0.92\textwidth,max height=0.38\textheight]{Ishii.pdf}
\end{sourceblock}
\begin{sourceblock}
\includegraphics[page=324,pagebox=mediabox,trim={45.00bp 53.74bp 44.37bp 541.50bp},clip,width=0.92\textwidth,max height=0.38\textheight]{Ishii.pdf}
\end{sourceblock}
\begin{sourceblock}
\includegraphics[page=325,pagebox=mediabox,trim={45.00bp 522.78bp 44.37bp 94.07bp},clip,width=0.92\textwidth,max height=0.38\textheight]{Ishii.pdf}
\end{sourceblock}
where

\begin{sourceblock}
\includegraphics[page=325,pagebox=mediabox,trim={45.00bp 450.16bp 44.37bp 177.32bp},clip,width=0.92\textwidth,max height=0.38\textheight]{Ishii.pdf}
\end{sourceblock}
\begin{sourceblock}
\includegraphics[page=325,pagebox=mediabox,trim={45.00bp 365.86bp 44.37bp 235.06bp},clip,width=0.92\textwidth,max height=0.38\textheight]{Ishii.pdf}
\end{sourceblock}
其中 b n 和 t α 分别是气泡数密度和热扩散率。雅各布数定义为

\begin{sourceblock}
\includegraphics[page=325,pagebox=mediabox,trim={45.00bp 280.12bp 44.37bp 342.92bp},clip,width=0.92\textwidth,max height=0.38\textheight]{Ishii.pdf}
\end{sourceblock}
其中 pf c 是液体的比热。由于式（11-246）计算的惯性控制区的停留时间与式（11-241）计算的传热控制区的停留时间相比非常小，因此将传热控制区的停留时间视为总停留时间。因此，由于惯性控制区域中的凝结而产生的凝结汇项表示为

\begin{sourceblock}
\includegraphics[page=325,pagebox=mediabox,trim={45.00bp 127.18bp 44.37bp 497.22bp},clip,width=0.92\textwidth,max height=0.38\textheight]{Ishii.pdf}
\end{sourceblock}
总冷凝汇项 CO φ 由两个区域中的冷凝汇项之和给出：

\begin{sourceblock}
\includegraphics[page=326,pagebox=mediabox,trim={45.00bp 562.90bp 44.37bp 55.74bp},clip,width=0.92\textwidth,max height=0.38\textheight]{Ishii.pdf}
\end{sourceblock}
凝结汇项预测的关键是在式(11-251)中实现精确的凝结努塞尔数模型。 Park 等人对现有努塞尔数模型进行了回顾。 （2007）。

\chapter{界面输运的流体动力学本构关系}
在分析界面力和相之间的相对运动时，首先考虑每个相的动量方程。假设本体流体和界面处的平均压力和应力大致相同，k 相动量方程为

\begin{sourceblock}
\includegraphics[page=327,pagebox=mediabox,trim={45.00bp 273.34bp 44.37bp 328.35bp},clip,width=0.92\textwidth,max height=0.38\textheight]{Ishii.pdf}
\end{sourceblock}
其中 k T 、 T k T 、 ki T 和 ik M 是平均粘性应力张量、平均湍流应力张量、界面剪切应力和广义界面阻力。混合物动量守恒要求

\begin{sourceblock}
\includegraphics[page=327,pagebox=mediabox,trim={45.00bp 166.96bp 44.37bp 462.32bp},clip,width=0.92\textwidth,max height=0.38\textheight]{Ishii.pdf}
\end{sourceblock}
这是平均动量跳跃条件的修改形式。使用二流体模型分析两相流需要平均湍流应力张量和广义界面阻力的本构方程。

在宏观两相流分析（例如一维两相流分析）中，除了壁面剪切贡献之外，平均湍流应力项可以忽略不计，而在微观气泡流分析中，湍流模型（例如混合长度模型和 k - ε 模型）具有
在二流体动量方程中，通过本构关系建模的最重要的项是广义阻力 id M，它指定了界面表面力。模拟该力的最简单方法是将各种已知界面力的线性组合公式化为

\begin{sourceblock}
\includegraphics[page=328,pagebox=mediabox,trim={45.00bp 436.60bp 44.37bp 168.64bp},clip,width=0.92\textwidth,max height=0.38\textheight]{Ishii.pdf}
\end{sourceblock}
其中 d B 、DF 、V F 、B F 、L F 、W F 和 T F 分别是典型颗粒的体积、标准曳力、虚拟质量力、巴塞特力、升力、壁升力和湍流色散力。

方程中各项的意义如下。左侧的项是作用于分散相的广义界面阻力的组合。右侧第一项是稳态条件下的蒙皮阻力和形状阻力。第二项是当相对速度变化时加速周围相的表观质量所需的力。第三项称为巴塞特力，是加速度对粘性阻力和边界层发展的影响。第四项是由于流体旋转而垂直于相对速度的升力。第五项是由于壁附近粒子周围的速度分布变化而产生的壁升力。最后一项是由于浓度梯度而产生的湍流色散力。在宏观两相流分析中，例如一维两相流分析，不考虑除标准曳力和虚拟质量力之外的力，而在三维流的微观分析中，还考虑附加力，例如升力和湍流分散力。

本章按照 Ishii 和 Zuber (1979)、Ishii 和 Chawla (1979)、Ishii 和 Mishima (1984) 详细讨论了界面输运的本构方程和多相流的界面流体力学。在下面的讨论中，为了简单起见，除了第 12 章的 1.4 节之外，省略了表示时间平均的符号。

\section*{12.1 多粒子系统中的瞬态力}\phantomsection\addcontentsline{toc}{section}{12.1 多粒子系统中的瞬态力}
这两个瞬态项的形式尚未确定。由于它们在瞬态条件下和数值稳定性问题中的重要性，需要在该领域进行进一步的研究。

巴塞特力由下式给出

\begin{sourceblock}
\includegraphics[page=329,pagebox=mediabox,trim={45.00bp 477.64bp 44.37bp 145.12bp},clip,width=0.92\textwidth,max height=0.38\textheight]{Ishii.pdf}
\end{sourceblock}
其中 m μ 是混合物粘度。该术语表示由于颗粒瞬态加速期间边界层或粘性流的发展而产生的额外阻力。导数 k D Dt 是相对于速度 k v 的对流导数。下标c和d分别代表连续相和分散相。混合物粘度的详细表达式在第12章的1.2节中给出。由于巴塞特力的时间积分形式复杂，在实际的两相流分析中不考虑巴塞特力。 Clift 等人给出了该术语对更高雷诺数的一些评估。 （1978）。

在这里，我们将针对 ( ) d d c D D ϋ − v v 随时间的弱函数情况推导出 Basset 力的近似形式。那么式(12-4)可以改写为

\begin{sourceblock}
\includegraphics[page=329,pagebox=mediabox,trim={45.00bp 224.66bp 44.37bp 365.21bp},clip,width=0.92\textwidth,max height=0.38\textheight]{Ishii.pdf}
\end{sourceblock}
其中 τ 是巴塞特力的时间尺度，可以根据动量穿透深度估计为

\begin{sourceblock}
\includegraphics[page=329,pagebox=mediabox,trim={45.00bp 157.48bp 44.37bp 483.75bp},clip,width=0.92\textwidth,max height=0.38\textheight]{Ishii.pdf}
\end{sourceblock}
将式(12-6)代入式(12-5)得

\begin{sourceblock}
\includegraphics[page=329,pagebox=mediabox,trim={45.00bp 86.13bp 44.37bp 538.11bp},clip,width=0.92\textwidth,max height=0.38\textheight]{Ishii.pdf}
\end{sourceblock}
巴塞特力 B d F 与界面巴塞特力的关系为

\begin{sourceblock}
\includegraphics[page=330,pagebox=mediabox,trim={45.00bp 567.87bp 44.37bp 58.03bp},clip,width=0.92\textwidth,max height=0.38\textheight]{Ishii.pdf}
\end{sourceblock}
\begin{sourceblock}
\includegraphics[page=330,pagebox=mediabox,trim={45.00bp 495.69bp 44.37bp 107.37bp},clip,width=0.92\textwidth,max height=0.38\textheight]{Ishii.pdf}
\end{sourceblock}
\begin{sourceblock}
\includegraphics[page=330,pagebox=mediabox,trim={45.00bp 407.80bp 44.37bp 193.34bp},clip,width=0.92\textwidth,max height=0.38\textheight]{Ishii.pdf}
\end{sourceblock}
莱希等人。 (1978)研究了虚质量项本构方程的必要条件。根据本构方程无差异的要求，他们确定虚拟质量力V d F 应满足

\begin{sourceblock}
\includegraphics[page=330,pagebox=mediabox,trim={45.00bp 295.40bp 44.37bp 306.38bp},clip,width=0.92\textwidth,max height=0.38\textheight]{Ishii.pdf}
\end{sourceblock}
鉴于 Zuber (1964a) 对浓度影响的研究和上述框架无差异条件，这里提出了 V d F 的新形式。由于颗粒相对于流体的加速度，会产生加速阻力。这应该与引起的质量 * c d ρ B 和与框架无关的相对加速度矢量成正比。因此，

\begin{sourceblock}
\includegraphics[page=330,pagebox=mediabox,trim={45.00bp 172.47bp 44.37bp 451.19bp},clip,width=0.92\textwidth,max height=0.38\textheight]{Ishii.pdf}
\end{sourceblock}
无限介质中单个粒子的诱导质量 * c d ρ B 的值可以由势论获得。因此，对于球形颗粒，V d F 在 0 d α → 处的极限值为

\begin{sourceblock}
\includegraphics[page=330,pagebox=mediabox,trim={45.00bp 75.64bp 44.37bp 549.46bp},clip,width=0.92\textwidth,max height=0.38\textheight]{Ishii.pdf}
\end{sourceblock}
从这个极限可以看出

\begin{sourceblock}
\includegraphics[page=331,pagebox=mediabox,trim={45.00bp 571.68bp 44.37bp 57.53bp},clip,width=0.92\textwidth,max height=0.38\textheight]{Ishii.pdf}
\end{sourceblock}
and

\begin{sourceblock}
\includegraphics[page=331,pagebox=mediabox,trim={45.00bp 511.95bp 44.37bp 124.11bp},clip,width=0.92\textwidth,max height=0.38\textheight]{Ishii.pdf}
\end{sourceblock}
如果式(12-12)中的λ为常数，则λ的值应为2。

浓度对 * d B 的影响可以通过 Zuber (1964a) 使用的方法考虑。因此，根据在外球体内移动的球体的感生质量的解，* d B 可以近似为：

\begin{sourceblock}
\includegraphics[page=331,pagebox=mediabox,trim={45.00bp 387.76bp 44.37bp 236.55bp},clip,width=0.92\textwidth,max height=0.38\textheight]{Ishii.pdf}
\end{sourceblock}
其中 d α 是分散相的体积分数。在假设 λ = 常数的情况下，由式（12-12）和式（12-16）得到虚拟质量力的本构方程为（Ishii 和 Mishima，1984）

\begin{sourceblock}
\includegraphics[page=331,pagebox=mediabox,trim={45.00bp 287.80bp 44.37bp 336.35bp},clip,width=0.92\textwidth,max height=0.38\textheight]{Ishii.pdf}
\end{sourceblock}
上式表明，每个颗粒的虚拟质量力 V d F 随着颗粒浓度的增加而显著增加。该关系意味着浓度对动态耦合的影响可以通过 ( ) ( ) 1 2 1 d d α α + − 因子来缩放。 Mokeyev（1977）采用电流体动力学模拟方法，通过电场势确定速度势，并获得经验函数
0 5 2 1 *。。 d d d B B α = +。式（12-17）的理论结果与该相关式相比较。可以使用伯努利方程通过简单的势流分析得出段塞流中虚拟质量力的相关性。首先考虑直径为D的管子中长度为b L 、直径为b D 的圆柱形气泡，见图12-1(a)。那么段塞气泡部分的空隙率由下式给出

\begin{sourceblock}
\includegraphics[page=331,pagebox=mediabox,trim={45.00bp 59.62bp 44.37bp 569.67bp},clip,width=0.92\textwidth,max height=0.38\textheight]{Ishii.pdf}
\end{sourceblock}
\begin{sourceblock}
\includegraphics[page=332,pagebox=mediabox,trim={114.15bp 472.09bp 114.50bp 68.35bp},clip,width=\textwidth,max height=0.72\textheight]{Ishii.pdf}
\end{sourceblock}
(a) 圆柱形气泡

\begin{sourceblock}
\includegraphics[page=332,pagebox=mediabox,trim={94.71bp 306.43bp 94.94bp 233.71bp},clip,width=\textwidth,max height=0.72\textheight]{Ishii.pdf}
\end{sourceblock}
(b) 球边圆柱形气泡图 12-1。用于虚拟质量力分析的弹流模型（Ishii 和 Mishima，1984）

\begin{sourceblock}
\includegraphics[page=332,pagebox=mediabox,trim={45.00bp 177.46bp 44.37bp 424.55bp},clip,width=0.92\textwidth,max height=0.38\textheight]{Ishii.pdf}
\end{sourceblock}
其中 L 是节距。现在让连续相相对于气泡加速。由于沿着薄膜部分的加速度，这将产生作用在气泡上的压力。通过简单的一维分析，可以发现该力为

\begin{sourceblock}
\includegraphics[page=332,pagebox=mediabox,trim={45.00bp 66.69bp 44.37bp 557.53bp},clip,width=0.92\textwidth,max height=0.38\textheight]{Ishii.pdf}
\end{sourceblock}
2 4 d b b B D L π =，因此单位体积的虚拟质量力变为

\begin{sourceblock}
\includegraphics[page=333,pagebox=mediabox,trim={45.00bp 529.61bp 44.37bp 94.77bp},clip,width=0.92\textwidth,max height=0.38\textheight]{Ishii.pdf}
\end{sourceblock}
这里，第二种形式是通过将段塞气泡部分中的空隙率近似为0 8 来获得的。 b α ≅。考虑的第二种情况是一串球边圆柱形气泡，见图12-1(b)。伯努利方程的应用

\begin{sourceblock}
\includegraphics[page=333,pagebox=mediabox,trim={45.00bp 417.20bp 44.37bp 207.67bp},clip,width=0.92\textwidth,max height=0.38\textheight]{Ishii.pdf}
\end{sourceblock}
其中 Φ 和 Ω 分别是该几何体在相对加速度下的速度势和势函数

\begin{sourceblock}
\includegraphics[page=333,pagebox=mediabox,trim={45.00bp 329.20bp 44.37bp 292.71bp},clip,width=0.92\textwidth,max height=0.38\textheight]{Ishii.pdf}
\end{sourceblock}
其中对近似值 0 8 进行了简化。 b α ≅（石井和三岛，1981）。对于一串球形气泡的极限情况，
b b L D =，上式简化为

\begin{sourceblock}
\includegraphics[page=333,pagebox=mediabox,trim={45.00bp 230.68bp 44.37bp 393.75bp},clip,width=0.92\textwidth,max height=0.38\textheight]{Ishii.pdf}
\end{sourceblock}
如果 b b L D >>，则 b L L 可以近似为 d b α α。因此，对于长段状气泡，方程（12-23）本质上收敛于方程（12-22）给出的简单解。式（12-23）给出的段塞流的虚拟质量力以没有大对流加速度的情况下的相对加速度来表示。然而，如果对流加速度不能忽略，则式(12-12)形式的特殊对流导数可能更合适。因此对于一般情况，

\begin{sourceblock}
\includegraphics[page=333,pagebox=mediabox,trim={45.00bp 78.37bp 44.37bp 544.59bp},clip,width=0.92\textwidth,max height=0.38\textheight]{Ishii.pdf}
\end{sourceblock}
该公式也可应用于搅动湍流。

现在，可以通过引入诱导质量系数来检查分散流方程（12-17）和段塞流方程（12-23）的解
M C 定义为

\begin{sourceblock}
\includegraphics[page=334,pagebox=mediabox,trim={45.00bp 516.58bp 44.37bp 107.66bp},clip,width=0.92\textwidth,max height=0.38\textheight]{Ishii.pdf}
\end{sourceblock}
where

\begin{sourceblock}
\includegraphics[page=334,pagebox=mediabox,trim={45.00bp 393.27bp 44.37bp 180.79bp},clip,width=0.92\textwidth,max height=0.38\textheight]{Ishii.pdf}
\end{sourceblock}
M C 与 d α 的关系图如图 12-2 所示。由于两相之间的耦合更强，虚拟质量力随着分散相的空隙率的增加而增加。上述两个解的交集出现在空隙率介于 0.66 和 0.75 之间的位置。对于较低的空隙率，气泡流的虚拟质量力小于段塞流的虚拟质量力。这意味着，在气泡流结构中，气相对加速的阻力比段塞流结构中的阻力小，如果
0 66。 d α < .这也可能表明加速的段塞流在 0 66 时有分解成气泡流的趋势。 d α < .对于 0 66。 d α > 即使在瞬态条件下，段塞流也应该相当稳定。

由于流动几何形状的相似性，湍流的虚拟质量力可以通过方程（12-24）给出的段塞流的解来近似。在液体分散流中，虚拟质量力变得比蒸气分散流中的虚拟质量力小得多。这种下降是由方程（12-17）中使用的连续相密度的变化引起的。通过将 c ρ 从 f ρ 更改为 g ρ，液滴流的虚拟质量力变得微不足道。这也表明，环形和环形分散流中的虚拟质量力应显著减小。

在当前的计算流体动力学（CFD）代码中，巴塞特力经常被忽略，虚拟质量系数被设置为较高的值以稳定计算。分散流中巴塞特力与虚拟质量力的比值可由式（1）获得。 (12-9) 和 (12-10) 为

\begin{sourceblock}
\includegraphics[page=335,pagebox=mediabox,trim={50.00bp 362.52bp 44.37bp 55.00bp},clip,width=\textwidth,max height=0.72\textheight]{Ishii.pdf}
\par\vspace{0.45\baselineskip}
\small 图12-2。分散流和段塞流流态的虚拟质量系数（Ishii 和 Mishima，
\end{sourceblock}
\begin{sourceblock}
\includegraphics[page=335,pagebox=mediabox,trim={50.00bp 250.45bp 44.37bp 317.62bp},clip,width=\textwidth,max height=0.72\textheight]{Ishii.pdf}
\par\vspace{0.45\baselineskip}
\small 图 12-3 表明巴塞特力与虚拟质量力一起可以稳定数值模拟并且很重要。
\end{sourceblock}
\section*{12.2 多粒子系统中的阻力}\phantomsection\addcontentsline{toc}{section}{12.2 多粒子系统中的阻力}
在稳态条件下作用在颗粒上的标准阻力可以根据基于相对速度的阻力系数 DC 给出：

\begin{sourceblock}
\includegraphics[page=335,pagebox=mediabox,trim={45.00bp 97.33bp 44.37bp 531.88bp},clip,width=0.92\textwidth,max height=0.38\textheight]{Ishii.pdf}
\end{sourceblock}
其中 d A 是典型粒子的投影面积，r v 是由 r dc = − v v v 给出的相对速度。那么 D d F 与界面阻力的关系为

\begin{sourceblock}
\includegraphics[page=336,pagebox=mediabox,trim={50.00bp 346.62bp 44.37bp 55.00bp},clip,width=\textwidth,max height=0.72\textheight]{Ishii.pdf}
\par\vspace{0.45\baselineskip}
\small 图12-3。巴塞特力与虚拟质量力的比较
\end{sourceblock}
\begin{sourceblock}
\includegraphics[page=336,pagebox=mediabox,trim={45.00bp 258.57bp 44.37bp 367.36bp},clip,width=0.92\textwidth,max height=0.38\textheight]{Ishii.pdf}
\end{sourceblock}
\begin{sourceblock}
\includegraphics[page=336,pagebox=mediabox,trim={45.00bp 182.50bp 44.37bp 416.58bp},clip,width=0.92\textwidth,max height=0.38\textheight]{Ishii.pdf}
\end{sourceblock}
下面从单粒子体系开始详细说明分散两相流中阻力系数DC的本构关系。

\subsection*{12.2.1 单颗粒阻力系数}\phantomsection\addcontentsline{toc}{subsection}{12.2.1 单颗粒阻力系数}
单个固体颗粒、液滴或气泡在无限介质中的运动过去已被广泛研究，例如参见 Peebles 和 Garber (1953)、Harmathy (1960) 和 Wallis (1974)。接下来我们

\begin{sourceblock}
\includegraphics[page=337,pagebox=mediabox,trim={45.00bp 519.47bp 44.37bp 94.67bp},clip,width=0.92\textwidth,max height=0.38\textheight]{Ishii.pdf}
\end{sourceblock}
其中 D F 是阻力，d r 是粒子的半径。为了根据相对速度计算曳力 D F，我们应该确定独立于方程（12-32）的 D C ∞ 本构关系。对于单粒子阻力相关性，两个相似性参数很重要。它们是颗粒雷诺数和粘度数

\begin{sourceblock}
\includegraphics[page=337,pagebox=mediabox,trim={45.00bp 389.56bp 44.37bp 233.85bp},clip,width=0.92\textwidth,max height=0.38\textheight]{Ishii.pdf}
\end{sourceblock}
and

\begin{sourceblock}
\includegraphics[page=337,pagebox=mediabox,trim={45.00bp 299.10bp 44.37bp 304.78bp},clip,width=0.92\textwidth,max height=0.38\textheight]{Ishii.pdf}
\end{sourceblock}
对单粒子阻力的广泛研究表明，在大多数情况下，阻力系数是雷诺数的函数（见图12-4）。然而，确切的函数形式取决于粒子是固体粒子、液滴还是气泡。简而言之，对于固体球形颗粒系统，我们有粘性状态，其中雷诺数依赖性为
D C ∞ 是明显的，并且牛顿状态，其中 D C ∞ 独立于
Re N ∞。在粘性状态下的清洁流体球的情况下，与固体颗粒的相关性预测的值相比，D C ∞ 可以减少高达 33 \%。这是通过流体颗粒内的内部循环来解释的。然而，少量的杂质就足以消除这种减阻作用。因此，对于大多数实际应用，流体颗粒系统中的阻力定律可以近似于达到一定颗粒尺寸的固体颗粒系统中的阻力定律。超过这一点，粒子形状的扭曲和不规则运动都会变得明显。在这种扭曲粒子状态下，D C ∞ 不依赖于粘度，而是随粒子半径线性增加。由于流体动力学不稳定，D C ∞ 存在上限，并且

\begin{sourceblock}
\includegraphics[page=338,pagebox=mediabox,trim={50.00bp 343.50bp 44.37bp 55.00bp},clip,width=\textwidth,max height=0.72\textheight]{Ishii.pdf}
\par\vspace{0.45\baselineskip}
\small 图12-4。单粒子阻力系数（Ishii 和 Chawla，1979）
\end{sourceblock}
颗粒达到帽泡条件或最大液滴尺寸。这些阻力系数的状态如图 12-4 所示。

对于粘性状态，函数 DC 由经验相关性给出：

\begin{sourceblock}
\includegraphics[page=338,pagebox=mediabox,trim={45.00bp 204.96bp 44.37bp 421.90bp},clip,width=0.92\textwidth,max height=0.38\textheight]{Ishii.pdf}
\end{sourceblock}
当雷诺数较小时 ( 1 Re N ∞ < )，上述相关性本质上简化为众所周知的斯托克斯阻力定律 24 Re D C N ∞ =。粘性状态的相关性表明阻力系数对雷诺数的依赖性随着雷诺数值的增加而减小。

在固体颗粒中，阻力系数基本恒定在大约

\begin{sourceblock}
\includegraphics[page=338,pagebox=mediabox,trim={45.00bp 71.26bp 44.37bp 569.37bp},clip,width=0.92\textwidth,max height=0.38\textheight]{Ishii.pdf}
\end{sourceblock}
对于液滴或气泡等流体颗粒，我们具有以颗粒形状扭曲和不规则运动为特征的流态。在这种扭曲的颗粒状态下，实验数据表明终端速度与颗粒尺寸无关（见图12-5）。在图 12-5 中，无量纲终端速度 * r v ∞ 和折减半径 * d r 定义为

\begin{sourceblock}
\includegraphics[page=339,pagebox=mediabox,trim={45.00bp 461.79bp 44.37bp 161.29bp},clip,width=0.92\textwidth,max height=0.38\textheight]{Ishii.pdf}
\end{sourceblock}
可以看出，阻力系数 D C ∞ 不依赖于粘度，但它应该与颗粒的半径成正比（Harmathy，1960）。从物理上讲，这表明阻力是由颗粒的扭曲和偏转运动控制的，并且颗粒形状的变化朝向有效横截面的增加。因此，D C ∞ 可以通过粒子的平均半径而不是雷诺数来缩放（Harmathy，1960）。然后我们有

\begin{sourceblock}
\includegraphics[page=339,pagebox=mediabox,trim={45.00bp 321.21bp 44.37bp 303.77bp},clip,width=0.92\textwidth,max height=0.38\textheight]{Ishii.pdf}
\end{sourceblock}
这里，使用基于终端速度的流体颗粒尺寸。

因此，粘性流和扭曲颗粒流之间的流态转变可以用粘数给出，如图12-4和12-5所示。然而，由于在这种情况下，终端速度可以与属性唯一相关，因此方程（12-37）可以根据终端速度或雷诺数概括为

\begin{sourceblock}
\includegraphics[page=339,pagebox=mediabox,trim={45.00bp 184.45bp 44.37bp 444.76bp},clip,width=0.92\textwidth,max height=0.38\textheight]{Ishii.pdf}
\end{sourceblock}
随着气泡尺寸进一步增大，气泡变成球形，阻力系数达到恒定值

\begin{sourceblock}
\includegraphics[page=339,pagebox=mediabox,trim={45.00bp 102.49bp 44.37bp 526.72bp},clip,width=0.92\textwidth,max height=0.38\textheight]{Ishii.pdf}
\end{sourceblock}
从扭曲气泡状态到球帽气泡状态的转变发生在

\begin{sourceblock}
\includegraphics[page=340,pagebox=mediabox,trim={50.00bp 337.38bp 44.37bp 55.00bp},clip,width=\textwidth,max height=0.72\textheight]{Ishii.pdf}
\par\vspace{0.45\baselineskip}
\small 图12-5。单粒子系统的终端速度（Ishii 和 Chawla，1979）
\end{sourceblock}
\begin{sourceblock}
\includegraphics[page=340,pagebox=mediabox,trim={45.00bp 251.90bp 44.37bp 373.07bp},clip,width=0.92\textwidth,max height=0.38\textheight]{Ishii.pdf}
\end{sourceblock}
对于液滴，根据式（12-37），阻力系数可能进一步增加。然而，最终液滴变得不稳定并分解成小液滴。这个极限可以由著名的韦伯数准则给出。通过引入韦伯数 2 2 g r d We v r ρ σ ≡，其中 r v 是相对速度，我们可以给出近似为 12 We 的稳定性准则。由于式(12-37)对应的终端速度为 ( )

\begin{sourceblock}
\includegraphics[page=340,pagebox=mediabox,trim={112.29bp 140.41bp 128.84bp 471.07bp},clip,width=\textwidth,max height=0.72\textheight]{Ishii.pdf}
\end{sourceblock}
\begin{sourceblock}
\includegraphics[page=340,pagebox=mediabox,trim={45.00bp 88.94bp 44.37bp 536.03bp},clip,width=0.92\textwidth,max height=0.38\textheight]{Ishii.pdf}
\end{sourceblock}
\begin{sourceblock}
\includegraphics[page=341,pagebox=mediabox,trim={50.00bp 317.94bp 44.37bp 55.00bp},clip,width=\textwidth,max height=0.72\textheight]{Ishii.pdf}
\par\vspace{0.45\baselineskip}
\small 图12-6。由于不稳定而导致大气泡破裂（500 帧/秒）
\end{sourceblock}
对应于液滴的最大阻力系数 4 DC =。如果液滴界面的稳定性受泰勒不稳定性控制，则特征液滴半径由下式给出

\begin{sourceblock}
\includegraphics[page=341,pagebox=mediabox,trim={45.00bp 193.46bp 44.37bp 431.51bp},clip,width=0.92\textwidth,max height=0.38\textheight]{Ishii.pdf}
\end{sourceblock}
这可能是液滴尺寸的更实际的上限。还值得注意的是，在高度湍流（Hinze，1959）或压力冲击条件下（Dinh 等，2003；Theofanous 等，2004），稳定极限韦伯数可能远小于 12。

帽泡保持一定的规则形状，伴流角约为 50o，但也存在最大稳定帽泡直径（Grace 等，1978；Clift 等，1978；Kocamstafaogullari 等，1984；Miller，1993）。这种不稳定性如图12-6所示。科卡穆斯塔福古拉里等人。 (1984) 使用了基于沿帽泡表面的 Kelvin-Helmholtz 不稳定性的稳定性分析。通过比较表面波停留时间

\begin{sourceblock}
\includegraphics[page=342,pagebox=mediabox,trim={45.00bp 519.20bp 44.37bp 105.77bp},clip,width=0.92\textwidth,max height=0.38\textheight]{Ishii.pdf}
\end{sourceblock}
这个结果很重要，因为它定义了可能出现段塞流的较小管道和由于泰勒气泡的不稳定性而无法形成段塞流的较大管道之间的边界。对于管道直径
( ) 2 40 ,max cap D r g σ Δρ < = 可形成稳定的段塞流。然而对于
2，最大帽D r >> 气泡流之后是帽湍流，其中在较高气体通量下存在多个相互作用的帽气泡。

利用上述阻力系数，通过平衡压力、重力和阻力，我们可以得到无限介质中的终端速度。各种颗粒和流态的结果总结在图 12-5 中。

\subsection*{12.2.2 分散两相流阻力系数}\phantomsection\addcontentsline{toc}{subsection}{12.2.2 分散两相流阻力系数}
A. 颗粒和流态的影响 上一节已经表明，单颗粒系统的阻力相关性不仅取决于流态，还取决于颗粒的性质。即固体颗粒、液滴或气泡。因此，对于多粒子系统，这些差异也有望在确定阻力相关性方面发挥核心作用。在本研究中，通过考虑以下流态，多粒子阻力相关性与单粒子系统并行开发

\begin{sourceblock}
\includegraphics[page=342,pagebox=mediabox,trim={45.00bp 176.34bp 44.37bp 437.74bp},clip,width=0.92\textwidth,max height=0.38\textheight]{Ishii.pdf}
\end{sourceblock}
\begin{sourceblock}
\includegraphics[page=342,pagebox=mediabox,trim={45.00bp 50.14bp 44.37bp 506.21bp},clip,width=0.92\textwidth,max height=0.38\textheight]{Ishii.pdf}
\end{sourceblock}
B. 粘性状态（未扭曲粒子状态） 该状态的特点是粘性对粒子运动的强烈影响。对于流体-颗粒系统，只有当颗粒形状没有因界面不稳定性或湍流流体运动而扭曲时，才会出现这种情况。为了建立多粒子阻力相关性，引入了几种相似假设。首先，假设粘性区域的阻力系数可以作为粒子雷诺数的函数给出。因此，

\begin{sourceblock}
\includegraphics[page=343,pagebox=mediabox,trim={45.00bp 364.00bp 44.37bp 279.53bp},clip,width=0.92\textwidth,max height=0.38\textheight]{Ishii.pdf}
\end{sourceblock}
其中雷诺数根据混合物粘度 m μ 定义为

\begin{sourceblock}
\includegraphics[page=343,pagebox=mediabox,trim={45.00bp 274.61bp 44.37bp 348.82bp},clip,width=0.92\textwidth,max height=0.38\textheight]{Ishii.pdf}
\end{sourceblock}
阻力系数Eq.(12-29)的介绍以及Eq.(12-29)的使用(12-44)基于这样的假设：两相混合物中颗粒运动的阻力可以通过考虑代表性颗粒与周围流体之间的相对运动引起的剪切的局部阻力来评估。其他颗粒对阻力的影响源于颗粒对流场变形的阻力。由于颗粒比流体更坚硬，不易变形，因此颗粒将施加一个作用在流体上的力系统。由于额外的应力，原始颗粒的运动阻力增加，这在它看来是由于粘度的增加而引起的。

因此，在分析悬浮颗粒的运动时，应使用混合物粘度（Burgers，1941；Zuber，1964a）。预计混合物粘度是浓度、流体粘度和颗粒粘度的函数。分散相的粘度考虑了界面的流动性，是电阻的量度
在本分析中，我们基于最大堆积 dm α，将流体颗粒的混合物粘度的线性相关性（Taylor，1932）扩展到固体颗粒的幂关系（Roscoe，1952）。流体-粒子系统的泰勒粘度模型由下式给出

\begin{sourceblock}
\includegraphics[page=344,pagebox=mediabox,trim={45.00bp 440.56bp 44.37bp 183.76bp},clip,width=0.92\textwidth,max height=0.38\textheight]{Ishii.pdf}
\end{sourceblock}
仅适用于 1 d α <<。固体颗粒系统的简单幂律粘度模型由下式给出

\begin{sourceblock}
\includegraphics[page=344,pagebox=mediabox,trim={45.00bp 349.54bp 44.37bp 270.02bp},clip,width=0.92\textwidth,max height=0.38\textheight]{Ishii.pdf}
\end{sourceblock}
这表明混合物的粘度在接近最大堆积时迅速增加。还要注意，方程（12-47）在小的情况下线性展开
d α 为1 2 5。 m c d μ μ α = +，与式(12-46)类似。固体颗粒系统的最大堆积 dm α 范围为 0.5 至 0.74。然而，0 62。 dm α = 足以满足大多数实际情况。对于气泡流，由于气泡的变形，理论 dm α 可能要高得多。在没有湍流运动和粒子聚结的情况下，流体-粒子系统中的空隙率可高达 0.95。通过将 dm α 视为单位，我们可以在分析中包含这些泡沫或致密堆积状态。因此，对于流体-粒子系统，我们取 1 dm α =。结合以上两个表达式，可以得出所有浓度下的固体颗粒系统和流体-颗粒系统的以下模型

\begin{sourceblock}
\includegraphics[page=344,pagebox=mediabox,trim={50.00bp 115.63bp 44.37bp 486.51bp},clip,width=\textwidth,max height=0.72\textheight]{Ishii.pdf}
\par\vspace{0.45\baselineskip}
\small 图 12-7 将此混合物粘度模型与固体颗粒系统的各种现有模型进行了比较（Eilers，1941；Roscoe，1952；Brinkman，1952；Frankel 和 Acrivos，1967；Landel 等，1965；Thomas，1965）。请注意，对于固体粒子系统，d μ 接近 Infini。因此，如果我们取方程（12-48）的极限，则粘度比项变为单位并且
\end{sourceblock}
\begin{sourceblock}
\includegraphics[page=345,pagebox=mediabox,trim={50.00bp 360.66bp 44.37bp 55.00bp},clip,width=\textwidth,max height=0.72\textheight]{Ishii.pdf}
\par\vspace{0.45\baselineskip}
\small 图 12-7。当前混合物粘度模型与现有固体模型的比较
\end{sourceblock}
粒子系统（Ishii 和 Chawla，1979）

相关性简化为幂律，方程（12-47）。通过在关联式中考虑分散相粘度的影响，该模型比传统关联式具有优势，因为它不仅限于颗粒流，还可以应用于液滴和气泡流。

使用最大填充的推荐值，我们可以通过简单的幂律来近似混合物粘度：

\begin{sourceblock}
\includegraphics[page=345,pagebox=mediabox,trim={45.00bp 199.90bp 44.37bp 439.34bp},clip,width=0.92\textwidth,max height=0.38\textheight]{Ishii.pdf}
\end{sourceblock}
where

\begin{sourceblock}
\includegraphics[page=345,pagebox=mediabox,trim={45.00bp 89.53bp 44.37bp 500.90bp},clip,width=0.92\textwidth,max height=0.38\textheight]{Ishii.pdf}
\end{sourceblock}
固体颗粒系统的表达式仅适用于 d α 的中等值。这些关系如图12-8所示。

分析中引入的第二个相似性假设是，在粘性状态下，单粒子系统和多粒子系统之间存在完全相似性。因此，多粒子阻力系数

\begin{sourceblock}
\includegraphics[page=346,pagebox=mediabox,trim={45.00bp 492.58bp 44.37bp 109.96bp},clip,width=0.92\textwidth,max height=0.38\textheight]{Ishii.pdf}
\end{sourceblock}
式(12-51)给出的关系如图12-9所示。这种相关性表明阻力系数随着体积浓度 d α 的增加而增加。通过比较单颗粒和多颗粒系统，图 12-10 清楚地显示了固体颗粒系统的这一趋势。

由 ( ) ( ) Re Re D D C N C N ∞ = 给出的相似性判据以及基于混合物粘度的雷诺数首次被引入到 Stokes 体系中的固体颗粒系统中 (Hawksley, 1951; Zuber, 1964a)。请注意，由于使用了广义阻力定律以及混合物粘度模型对流体颗粒系统的适用性，本模型不限于固体颗粒系统或斯托克斯体系。

C. 牛顿体系 在牛顿体系中，涡旋系统在粒子后面形成，它的离开产生了一个大的尾流区域。阻力主要由流动分离产生的涡流决定。因此，对于单粒子系统，阻力近似与惯性力成正比，并且阻力系数可以认为是常数。

\begin{sourceblock}
\includegraphics[page=346,pagebox=mediabox,trim={45.00bp 187.90bp 44.37bp 426.72bp},clip,width=0.92\textwidth,max height=0.38\textheight]{Ishii.pdf}
\end{sourceblock}
函数 ( ) d E α 可以通过考虑无限介质中终端速度的特殊情况来获得。根据重力、压力和阻力之间的力平衡，我们有

\begin{sourceblock}
\includegraphics[page=346,pagebox=mediabox,trim={45.00bp 88.60bp 44.37bp 538.73bp},clip,width=0.92\textwidth,max height=0.38\textheight]{Ishii.pdf}
\end{sourceblock}
对于单粒子系统，这可以简化为

\begin{sourceblock}
\includegraphics[page=347,pagebox=mediabox,trim={50.00bp 360.66bp 44.37bp 55.00bp},clip,width=\textwidth,max height=0.72\textheight]{Ishii.pdf}
\par\vspace{0.45\baselineskip}
\small 图12-8。各种系统的混合物粘度（Ishii 和 Chawla，1979）
\end{sourceblock}
\begin{sourceblock}
\includegraphics[page=347,pagebox=mediabox,trim={50.00bp 63.72bp 44.37bp 319.48bp},clip,width=\textwidth,max height=0.72\textheight]{Ishii.pdf}
\par\vspace{0.45\baselineskip}
\small 图12-9。粘性状态下的阻力系数（Ishii 和 Chawla，1979）
\end{sourceblock}
\begin{sourceblock}
\includegraphics[page=348,pagebox=mediabox,trim={50.00bp 358.92bp 44.37bp 55.00bp},clip,width=\textwidth,max height=0.72\textheight]{Ishii.pdf}
\par\vspace{0.45\baselineskip}
\small 图 12-10。粘性状态下浓度对阻力系数的影响（Ishii 和 Chawla，
\end{sourceblock}
1979)

\begin{sourceblock}
\includegraphics[page=348,pagebox=mediabox,trim={45.00bp 261.36bp 44.37bp 365.51bp},clip,width=0.92\textwidth,max height=0.38\textheight]{Ishii.pdf}
\end{sourceblock}
通过将多粒子系统与具有相同粒子尺寸的单粒子系统进行比较，我们有

\begin{sourceblock}
\includegraphics[page=348,pagebox=mediabox,trim={45.00bp 170.40bp 44.37bp 452.53bp},clip,width=0.92\textwidth,max height=0.38\textheight]{Ishii.pdf}
\end{sourceblock}
由于雷诺数可以是速度的函数，因此方程（12-55）是终端速度 r v 的隐式方程。如果我们考虑方程（12-35）和（12-51）给出的粘性阻力定律，则方程（12-55）变为

\begin{sourceblock}
\includegraphics[page=348,pagebox=mediabox,trim={45.00bp 68.40bp 44.37bp 555.51bp},clip,width=0.92\textwidth,max height=0.38\textheight]{Ishii.pdf}
\end{sourceblock}
\begin{sourceblock}
\includegraphics[page=349,pagebox=mediabox,trim={45.00bp 546.11bp 44.37bp 55.67bp},clip,width=0.92\textwidth,max height=0.38\textheight]{Ishii.pdf}
\end{sourceblock}
\begin{sourceblock}
\includegraphics[page=349,pagebox=mediabox,trim={45.00bp 517.06bp 44.37bp 96.65bp},clip,width=0.92\textwidth,max height=0.38\textheight]{Ishii.pdf}
\end{sourceblock}
\begin{sourceblock}
\includegraphics[page=349,pagebox=mediabox,trim={45.00bp 468.05bp 44.37bp 125.57bp},clip,width=0.92\textwidth,max height=0.38\textheight]{Ishii.pdf}
\end{sourceblock}
\begin{sourceblock}
\includegraphics[page=349,pagebox=mediabox,trim={45.00bp 370.34bp 44.37bp 220.85bp},clip,width=0.92\textwidth,max height=0.38\textheight]{Ishii.pdf}
\end{sourceblock}
然而，由于终端速度 r v ∞ 通过式（12-54）与雷诺数唯一相关，因此 Re N ∞ 可以用粒子的半径代替。因此，

\begin{sourceblock}
\includegraphics[page=349,pagebox=mediabox,trim={45.00bp 258.39bp 44.37bp 353.17bp},clip,width=0.92\textwidth,max height=0.38\textheight]{Ishii.pdf}
\end{sourceblock}
where

\begin{sourceblock}
\includegraphics[page=349,pagebox=mediabox,trim={45.00bp 161.02bp 44.37bp 436.42bp},clip,width=0.92\textwidth,max height=0.38\textheight]{Ishii.pdf}
\end{sourceblock}
对于单粒子系统，从粘性状态到牛顿状态的转变发生在 34 65 * 处。 d r =（或 990 Re N ∞ ）。在该粒径下，方程（12-52）减少为

\begin{sourceblock}
\includegraphics[page=349,pagebox=mediabox,trim={45.00bp 55.35bp 44.37bp 558.71bp},clip,width=0.92\textwidth,max height=0.38\textheight]{Ishii.pdf}
\end{sourceblock}
\begin{sourceblock}
\includegraphics[page=350,pagebox=mediabox,trim={45.00bp 483.48bp 44.37bp 121.74bp},clip,width=0.92\textwidth,max height=0.38\textheight]{Ishii.pdf}
\end{sourceblock}
其中 0 45 34 65 *。在。 D d C r ∞ = =。鉴于式(12-48)有0 62。 dm α =，我们得到

\begin{sourceblock}
\includegraphics[page=350,pagebox=mediabox,trim={45.00bp 348.21bp 44.37bp 229.54bp},clip,width=0.92\textwidth,max height=0.38\textheight]{Ishii.pdf}
\end{sourceblock}
阻力系数随着体积浓度的增加而增加

d α 由式(12-64)绘制如图12-11所示。这意味着，由于相之间的耦合更强，平衡相对速度通常随着浓度的增加而降低。

D. 扭曲流体粒子体系 在扭曲流体粒子体系中，单粒子阻力系数仅取决于粒子半径和流体性质，而不取决于
速度或粘度，即 ( ) 4 3 D d C r g Δρ σ ∞ =，如所讨论的
哈马西（1960）着。因此，对于固定直径的粒子，D C ∞ 变为常数。在考虑具有相同半径的多粒子系统的阻力系数时，必须考虑其他粒子的存在对流场的限制。因此，在这种情况下，D C 预计与 D C ∞ 不同。由于湍流涡旋对尾流区域的强烈贡献，粒子会因其他粒子而受到增加的阻力，其方式与固体粒子系统的牛顿体系基本相似，其中
DC ∞ 在尾流主导的流动条件下也是恒定的。因此，我们假设，无论这些体系中 D C ∞ 的差异如何，扭曲流体粒子体系中阻力增加的影响都可以通过类似于牛顿体系中的表达式来预测。

\begin{sourceblock}
\includegraphics[page=351,pagebox=mediabox,trim={50.00bp 362.58bp 44.37bp 55.00bp},clip,width=\textwidth,max height=0.72\textheight]{Ishii.pdf}
\par\vspace{0.45\baselineskip}
\small 图 12-11。牛顿体系的阻力系数（Ishii 和 Chawla，1979）
\end{sourceblock}
在此假设下，方程（12-63）可以与方程（12-63）给出的 D C ∞ 的适当表达式一起使用。 （12-49）。因此，

\begin{sourceblock}
\includegraphics[page=351,pagebox=mediabox,trim={45.00bp 228.27bp 44.37bp 376.85bp},clip,width=0.92\textwidth,max height=0.38\textheight]{Ishii.pdf}
\end{sourceblock}
鉴于式（12-49）给出的近似，上述相关性简化为

\begin{sourceblock}
\includegraphics[page=351,pagebox=mediabox,trim={45.00bp 131.85bp 44.37bp 484.49bp},clip,width=0.92\textwidth,max height=0.38\textheight]{Ishii.pdf}
\end{sourceblock}
\begin{sourceblock}
\includegraphics[page=351,pagebox=mediabox,trim={45.00bp 50.14bp 44.37bp 543.93bp},clip,width=0.92\textwidth,max height=0.38\textheight]{Ishii.pdf}
\end{sourceblock}
\begin{sourceblock}
\includegraphics[page=352,pagebox=mediabox,trim={50.00bp 353.46bp 44.37bp 55.00bp},clip,width=\textwidth,max height=0.72\textheight]{Ishii.pdf}
\par\vspace{0.45\baselineskip}
\small 图 12-12。扭曲粒子状态下的阻力系数（Ishii 和 Chawla，1979）
\end{sourceblock}
对于液滴-液体流 ( 1 75 . n = )

\begin{sourceblock}
\includegraphics[page=352,pagebox=mediabox,trim={45.00bp 244.38bp 44.37bp 373.21bp},clip,width=0.92\textwidth,max height=0.38\textheight]{Ishii.pdf}
\end{sourceblock}
对于液滴-气体流 ( 2 5 . n = )

\begin{sourceblock}
\includegraphics[page=352,pagebox=mediabox,trim={45.00bp 161.46bp 44.37bp 456.13bp},clip,width=0.92\textwidth,max height=0.38\textheight]{Ishii.pdf}
\end{sourceblock}
上述三种关系如图12-12所示。相关性的形式表明，相间的动量耦合随着粒子浓度的增加而增加，就像牛顿体系的情况一样。

E. 搅动湍流流态随着流体粒子半径的进一步增加，由于形成大尾流，尾流和气泡边界层可能重叠
在搅拌-湍流流态中，由于上面讨论的流体动力学条件，典型的颗粒相对于平均体积通量j而不是连续相的平均速度移动。因此，阻力系数和阻力相似定律定义中的参考速度应该是漂移速度而不是相对速度。因此，阻力应由下式给出

\begin{sourceblock}
\includegraphics[page=353,pagebox=mediabox,trim={45.00bp 351.48bp 44.37bp 277.73bp},clip,width=0.92\textwidth,max height=0.38\textheight]{Ishii.pdf}
\end{sourceblock}
这里的漂移速度 dj V 是分散相相对于混合物体积中心速度的相对速度。它还可以通过以下方式与相之间的真实相对速度相关：

\begin{sourceblock}
\includegraphics[page=353,pagebox=mediabox,trim={45.00bp 268.90bp 44.37bp 372.37bp},clip,width=0.92\textwidth,max height=0.38\textheight]{Ishii.pdf}
\end{sourceblock}
其中总通量 j（体积中心速度）由下式给出

\begin{sourceblock}
\includegraphics[page=353,pagebox=mediabox,trim={45.00bp 209.98bp 44.37bp 431.29bp},clip,width=0.92\textwidth,max height=0.38\textheight]{Ishii.pdf}
\end{sourceblock}
在搅拌湍流流态中，一些颗粒应该已经达到与帽泡转变或液滴分裂相对应的变形极限。通过使用漂移速度作为以下形式的参考速度，可以将此限制作为韦伯数准则（Wallis，1969）的扩展给出

\begin{sourceblock}
\includegraphics[page=353,pagebox=mediabox,trim={45.00bp 73.51bp 44.37bp 540.94bp},clip,width=0.92\textwidth,max height=0.38\textheight]{Ishii.pdf}
\end{sourceblock}
\begin{sourceblock}
\includegraphics[page=354,pagebox=mediabox,trim={45.00bp 465.49bp 44.37bp 159.82bp},clip,width=0.92\textwidth,max height=0.38\textheight]{Ishii.pdf}
\end{sourceblock}
其中阻力表达式的参考 d r 由 ( ) ( ) 给出
2 2 1 4 或 6 d c r d r v σ ρ α [ ] = −，因为这些粒子控制相对运动。上述方程意味着基于相间真实相对速度的表观阻力系数应由下式给出

\begin{sourceblock}
\includegraphics[page=354,pagebox=mediabox,trim={45.00bp 357.48bp 44.37bp 271.73bp},clip,width=0.92\textwidth,max height=0.38\textheight]{Ishii.pdf}
\end{sourceblock}
由式（12-75）的形式可知，阻力系数随着颗粒体积浓度的增加而减小，如图12-13所示。因此，在搅动湍流流态中，d α 对 DC 的影响与其他流态中相反。这种特殊的趋势可以通过较大颗粒尾迹后面夹带其他颗粒的效应来解释。这种夹带促进了分散相的沟道而不增加阻力。随着体积浓度的增加，颗粒之间的相互作用朝着减小阻力的方向增加。

F. 段塞流流态 受限通道中分散两相流的限制情况之一是段塞流。当气泡的体积变得非常大时，气泡的形状会显著变形以适应通道的几何形状。气泡的直径变得接近管道的直径，并且薄液膜将气泡与管壁隔开。气泡呈细长的子弹形状，带有帽形鼻子。可以通过使用单个气泡鼻部周围的势流分析来研究这些气泡在相对无粘性流体中的运动。因此，Dumitrescu (1943) 分析得出上升速度为

\begin{sourceblock}
\includegraphics[page=354,pagebox=mediabox,trim={45.00bp 52.50bp 44.37bp 588.28bp},clip,width=0.92\textwidth,max height=0.38\textheight]{Ishii.pdf}
\end{sourceblock}
\begin{sourceblock}
\includegraphics[page=355,pagebox=mediabox,trim={50.00bp 362.58bp 44.37bp 55.00bp},clip,width=\textwidth,max height=0.72\textheight]{Ishii.pdf}
\par\vspace{0.45\baselineskip}
\small 图 12-13。搅拌湍流的阻力系数（Ishii 和 Chawla，1979）
\end{sourceblock}
其中 D 是水力直径。这一结果也与Dumitrescu (1943)和White and Beardmore (1962)的实验数据非常吻合。

在具有气泡链的流动系统中，应考虑浓度和速度分布的影响。一般来说，由于速度分布，核心速度高于横截面积平均速度。因此，基于平均速度的相对速度大于核心中的局部相对速度。这种效应被称为分布参数效应，Bankoff (1960)、Zuber 和 Findlay (1965) 以及 Ishii (1977) 等人进行了广泛的研究。当使用平均速度时，尼克林等人的结果。 （1962）和尼尔（1963）表明

\begin{sourceblock}
\includegraphics[page=355,pagebox=mediabox,trim={45.00bp 133.84bp 44.37bp 507.40bp},clip,width=0.92\textwidth,max height=0.38\textheight]{Ishii.pdf}
\end{sourceblock}
这里左边是气泡的漂移速度，即
dj d V v j = −。上式可以改写为

\begin{sourceblock}
\includegraphics[page=355,pagebox=mediabox,trim={45.00bp 57.88bp 44.37bp 583.36bp},clip,width=0.92\textwidth,max height=0.38\textheight]{Ishii.pdf}
\end{sourceblock}
\begin{sourceblock}
\includegraphics[page=356,pagebox=mediabox,trim={50.00bp 362.46bp 44.37bp 55.00bp},clip,width=\textwidth,max height=0.72\textheight]{Ishii.pdf}
\par\vspace{0.45\baselineskip}
\small 图 12-14。段塞流的阻力系数（Ishii 和 Chawla，1979）
\end{sourceblock}
其中 1 2 . core j j == 和 core j 可以被认为是磁芯中的局部总通量。在这种情况下，d core v j − 是核心中的局部漂移速度。由式(12-71)给出的关系式可知，局部相对速度r v 应满足

\begin{sourceblock}
\includegraphics[page=356,pagebox=mediabox,trim={45.00bp 237.76bp 44.37bp 403.48bp},clip,width=0.92\textwidth,max height=0.38\textheight]{Ishii.pdf}
\end{sourceblock}
这与 Dumitrescu 在 0 d α → 处的结果一致。通过将我们的讨论限制在局部阻力系数，我们可以将上述半经验结果重新转化为段塞流状态下阻力系数的相关性。鉴于式(12-79)和(12-53)，我们得到

\begin{sourceblock}
\includegraphics[page=356,pagebox=mediabox,trim={45.00bp 124.11bp 44.37bp 504.95bp},clip,width=0.92\textwidth,max height=0.38\textheight]{Ishii.pdf}
\end{sourceblock}
对于大多数实际应用，2 d r D 可以近似为 0.9。然后，

\begin{sourceblock}
\includegraphics[page=356,pagebox=mediabox,trim={45.00bp 66.76bp 44.37bp 574.29bp},clip,width=0.92\textwidth,max height=0.38\textheight]{Ishii.pdf}
\end{sourceblock}
\clearpage
\begin{center}\small 表 12-1。多粒子系统中的局部阻力系数（Ishii 和 Chawla，1979）\end{center}
\begin{sourceblock}
\includegraphics[page=357,pagebox=mediabox,trim={47.95bp 87.07bp 49.08bp 196.75bp},clip,width=\textwidth,max height=0.72\textheight]{Ishii.pdf}
\end{sourceblock}
\section*{12.3 其他力量}\phantomsection\addcontentsline{toc}{section}{12.3 其他力量}
在多维两相流分析中，除了标准阻力、巴塞特力和虚拟质量力外，还考虑升力和湍流分散力等一些力。从式（12-3）可以看出，这些力通常线性地添加到标准阻力、巴塞特力和虚拟质量力中。在实际的两相流中，气泡后面的尾流可能会完全改变液体的湍流结构，因此升力可能与湍流引起的力紧密耦合。由于升力和其他侧向力很小并且可能相互紧密耦合，因此很难通过实验来识别每个力。因此，是否可以将这些力线性地添加到标准阻力、巴塞特力和虚拟质量力中可能存在争议。目前的发展现状是，一些力的本构方程是基于推测而提出的，并且没有明确给出本构方程的适用流量范围。因此，与标准阻力不同，这种侧向力的本构方程尚未得到很好的发展。然而，这些力在预测三维气泡分布中发挥着重要作用。下面对多维两相流分析中常用的升力和湍流弥散力的一些本构方程进行简要解释。 Akiyama 和 Aritomi (2002) 也对升力建模进行了回顾。

\subsection*{12.3.1 升力}\phantomsection\addcontentsline{toc}{subsection}{12.3.1 升力}
首先考虑单个球形颗粒相对于均匀简单剪切穿过非常粘稠的液体，如图 12-15 所示。粒子受到升力 LS V d
− ∞ F，垂直于流动方向（Saffman，1965）

\begin{sourceblock}
\includegraphics[page=358,pagebox=mediabox,trim={45.00bp 164.53bp 44.37bp 459.70bp},clip,width=0.92\textwidth,max height=0.38\textheight]{Ishii.pdf}
\end{sourceblock}
对于正相对速度和速度梯度（ 0 r v ∞ > 和 0 f dv dx > ），升力将粒子推向负 x 方向。

接下来考虑将单个球形颗粒置于无粘性流体的弱剪切流中。粒子受到升力 LS I d
− ∞ F，如（Auton，1987）

\begin{sourceblock}
\includegraphics[page=359,pagebox=mediabox,trim={50.00bp 386.46bp 44.37bp 55.00bp},clip,width=\textwidth,max height=0.72\textheight]{Ishii.pdf}
\par\vspace{0.45\baselineskip}
\small 图 12-15。剪切流中粒子的示意图
\end{sourceblock}
\begin{sourceblock}
\includegraphics[page=359,pagebox=mediabox,trim={45.00bp 302.28bp 44.37bp 326.93bp},clip,width=0.92\textwidth,max height=0.38\textheight]{Ishii.pdf}
\end{sourceblock}
Legendre 和 Magnaudet (1998) 通过求解完整的纳维-斯托克斯方程，以数值方式研究了在粘性线性剪切流中稳定移动的球形气泡周围的三维流动。他们假设气泡表面是干净的，因此外流遵循零剪切应力条件，并且不会引起气泡的任何旋转。他们通过数值模拟，建立了有限雷诺数下粘性线性剪切流中球形气泡升力的经验相关性：

\begin{sourceblock}
\includegraphics[page=359,pagebox=mediabox,trim={45.00bp 155.28bp 44.37bp 473.93bp},clip,width=0.92\textwidth,max height=0.38\textheight]{Ishii.pdf}
\end{sourceblock}
球形气泡的升力系数 L C ∞ 由下式给出

\begin{sourceblock}
\includegraphics[page=359,pagebox=mediabox,trim={45.00bp 93.36bp 44.37bp 546.27bp},clip,width=0.92\textwidth,max height=0.38\textheight]{Ishii.pdf}
\end{sourceblock}
where

\begin{sourceblock}
\includegraphics[page=360,pagebox=mediabox,trim={45.00bp 496.65bp 44.37bp 55.43bp},clip,width=0.92\textwidth,max height=0.38\textheight]{Ishii.pdf}
\end{sourceblock}
\begin{sourceblock}
\includegraphics[page=360,pagebox=mediabox,trim={45.00bp 435.35bp 44.37bp 190.41bp},clip,width=0.92\textwidth,max height=0.38\textheight]{Ishii.pdf}
\end{sourceblock}
单粒子系统的雷诺数 Re N ∞ 由式（12-33）定义。式(12-85)在0 1 500 Re 范围内有效。 N ∞ ≤ ≤ 且 0 01 0 25。。 S G ≤ ≤。在 20 世纪 80 年代和 90 年代，进行了大量的实验来确定确定气泡横向迁移特征的重要参数。实验表明，相对较小和较大的气泡倾向于分别向通道壁和中心迁移（Zun，1988；Liu，1993；Hibiki 和 Ishii，1999）。对泊肃叶流中单个气泡的数值模拟（Tomiyama 等，1993；1995）表明，气泡向管道中心的迁移与变形气泡后面的倾斜尾流密切相关。因此，已经表明，气泡尺寸以及气泡尾流和气泡周围剪切场之间的复杂相互作用在气泡横向迁移中起着重要作用（Serizawa 和 Kataoka，1988；1994）。富山等人。 (2002) 测量了甘油水溶液简单剪切流中单个气泡的气泡轨迹，以评估作用于单个气泡上的横向升力。基于实验结果，他们假设倾斜尾流引起的升力与剪切引起的升力具有相同的函数形式，并提出了升力系数的经验相关性。

Hibiki 和 Ishii (2007) 基于 Legendre 和 Magnaudet (1998) 的剪切升力模型提出了升力系数的相关性：

\begin{sourceblock}
\includegraphics[page=360,pagebox=mediabox,trim={45.00bp 105.72bp 44.37bp 533.90bp},clip,width=0.92\textwidth,max height=0.38\textheight]{Ishii.pdf}
\end{sourceblock}
其中 xi ∞ 是考虑气泡变形对升力影响的修正因子。该系数是根据 Tomiyama 等人的数据确定的。 (2002) 在 10 5 5 的条件下拍摄。 log M N − ≤ ≤

\begin{sourceblock}
\includegraphics[page=361,pagebox=mediabox,trim={50.00bp 399.06bp 44.37bp 55.00bp},clip,width=\textwidth,max height=0.72\textheight]{Ishii.pdf}
\par\vspace{0.45\baselineskip}
\small 图 12-16。升力系数对气泡尺寸的依赖性
\end{sourceblock}
2 8 . - , 1 39 5 74。。 Eo N ≤ ≤ 且 1 0 8 3。 f dv dx s − ≤ ≤，其中 M N 和 Eo N 分别是 Morton 数和 Eötvös 数，定义为

\begin{sourceblock}
\includegraphics[page=361,pagebox=mediabox,trim={45.00bp 281.92bp 44.37bp 340.24bp},clip,width=0.92\textwidth,max height=0.38\textheight]{Ishii.pdf}
\end{sourceblock}
\begin{sourceblock}
\includegraphics[page=361,pagebox=mediabox,trim={45.00bp 217.65bp 44.37bp 405.23bp},clip,width=0.92\textwidth,max height=0.38\textheight]{Ishii.pdf}
\end{sourceblock}
如图12-16所示，基于Tomiyama等人的数据提出了以下修正因子的相关性。 （2002）作为

\begin{sourceblock}
\includegraphics[page=361,pagebox=mediabox,trim={45.00bp 147.50bp 44.37bp 492.90bp},clip,width=0.92\textwidth,max height=0.38\textheight]{Ishii.pdf}
\end{sourceblock}
其中 * d D 是气泡直径与扭曲气泡极限处的气泡直径之比，如下所示

\begin{sourceblock}
\includegraphics[page=361,pagebox=mediabox,trim={45.00bp 57.33bp 44.37bp 566.97bp},clip,width=0.92\textwidth,max height=0.38\textheight]{Ishii.pdf}
\end{sourceblock}
从式（12-91）可以看出，升力方向在 * 处反转，d crit D ∞ = 0.522，有趣的是大约是扭曲气泡极限处气泡直径的一半。这里应该注意的是，在弱线性剪切流中沿短轴移动的椭球升力系数的解析解表明该系数具有恒定的正符号（Naciri，1992）。因此，简单的气泡变形可能不足以引起升力的负号。据推测，剪切流中的气泡方向、尾流结构的改变和气泡形状可能会影响升力方向。方程（12-91）还表明，在测试的实验条件下，修正因子与无量纲剪切速率无关。

基于阻力模型和升力模型的相似性，Hibiki 和 Ishii (2007) 假设修正因子取常数值 1，极限 ψ ∞ = −，在 ( ) 0 643 * *，。 d d 极限 D D ∞ ∞ ≥ =。我们可以通过将式（12-88）中的 Re N ∞ 和 s G ∞ 替换为
Re N 和 s G，其中多粒子系统的雷诺数 Re N 是由方程 (12-45) 定义的，无量纲速度梯度 s G 是由下式定义的

\begin{sourceblock}
\includegraphics[page=362,pagebox=mediabox,trim={45.00bp 306.64bp 44.37bp 319.66bp},clip,width=0.92\textwidth,max height=0.38\textheight]{Ishii.pdf}
\end{sourceblock}
因此，多粒子系统中的净横向升力 L d M 近似为

\begin{sourceblock}
\includegraphics[page=362,pagebox=mediabox,trim={45.00bp 85.20bp 44.37bp 401.51bp},clip,width=0.92\textwidth,max height=0.38\textheight]{Ishii.pdf}
\end{sourceblock}
由式（12-94）计算的升力与空气-水气泡流数据（Wang et al., 1987）得到了合理的一致性，这意味着

\begin{sourceblock}
\includegraphics[page=363,pagebox=mediabox,trim={50.00bp 400.62bp 44.37bp 55.00bp},clip,width=\textwidth,max height=0.72\textheight]{Ishii.pdf}
\par\vspace{0.45\baselineskip}
\small 图 12-17。单颗粒系统中升力系数对气泡直径的依赖性
\end{sourceblock}
升力模型方程（12-94）有望用于预测多粒子气泡流中的净横向升力。未来的研究应进一步努力检验方程（12-94）对多粒子系统的适用性。如上所述，人们对升力的了解仍然很少，需要进一步的实验和数值研究来了解升力（Sridhar 和 Katz，1995；Ervin 和 Tryggvason，1997；Loth 等，1997）。图12-17显示了20℃、0.1MPa下空气-水系统升力系数与气泡直径的关系。实线和虚线分别是 Hibiki-Ishii 模型 (2007) 计算出的可变形气泡在 0 1 处的升力系数。 S G ∞ = 以及 Legendre 和 Magnaudet 的球形粒子模型 (1998)。两个模型之间的差异表明气泡变形对升力系数的影响。

\subsection*{12.3.2 壁提升（壁润滑）力}\phantomsection\addcontentsline{toc}{subsection}{12.3.2 壁提升（壁润滑）力}
引入了壁升力 W d M，并将其解释为由于壁附近粒子周围的速度分布的变化（Antal 等人，1991）。该力用于预测观察到的并流层流向上和向下流动的空隙轮廓。该壁升力类似于润滑力，作用在壁附近的气泡上，以防止气泡接触壁。

考虑在壁附近移动的气泡周围液体的排出。墙壁的防滑条件会减缓排水速度

\begin{sourceblock}
\includegraphics[page=364,pagebox=mediabox,trim={45.00bp 428.69bp 44.37bp 188.11bp},clip,width=0.92\textwidth,max height=0.38\textheight]{Ishii.pdf}
\end{sourceblock}
其中 ( ) ( ) g f W g f W [ ] = − − ⋅ −     v v v n v v n \& , 1 0 104 0 06。。 W r C v = − − 且
2 0 147。 WC = .这里，bw d 和W n 分别是气泡与壁之间的距离和壁表面上的单位向外法向量。方程(12-95)表明壁升力的方向在( ) 0 0 707 0 408 处反转。。 d r y r v = +，且不考虑近壁面气泡变形的影响。已经进行了一些实验和分析工作来改善壁升力（Tomiyama等人，1995；Hosokawa和Tomiyama，2003；2006；Frank等人，2008），但尚未获得对壁升力的全面理解。

\subsection*{12.3.3 湍流色散力}\phantomsection\addcontentsline{toc}{subsection}{12.3.3 湍流色散力}
从物理上讲，湍流扩散是作用在气泡上的力的波动分量的结果。在最简单的情况下，某一点的湍流扩散（或分散）力是轨迹与该点相交的所有气泡上阻力波动分量的整体平均值（Lopez de Bertodano 等，2006）。 Lopez de Bertodano (1998) 根据 Reeks (1991; 1992) 获得的动力学方程推导了湍流色散力 T d M：

\begin{sourceblock}
\includegraphics[page=364,pagebox=mediabox,trim={45.00bp 111.58bp 44.37bp 515.69bp},clip,width=0.92\textwidth,max height=0.38\textheight]{Ishii.pdf}
\end{sourceblock}
斯托克斯数定义为

\begin{sourceblock}
\includegraphics[page=365,pagebox=mediabox,trim={45.00bp 571.78bp 44.37bp 52.54bp},clip,width=0.92\textwidth,max height=0.38\textheight]{Ishii.pdf}
\end{sourceblock}
其中 b τ 和 e τ 分别是气泡和涡流的弛豫时间。涡流的弛豫时间表示为涡流周转时间 et τ 和涡流交叉时间 ec τ 的组合。方程 (12-96) 适用于颗粒或气泡小于剪切引起涡流的剪切引起湍流。

对于 et ec τ τ，方程（12-96）与基于分子色散力（Lahey et al., 1993）类比推导的湍流色散力一致：

\begin{sourceblock}
\includegraphics[page=365,pagebox=mediabox,trim={45.00bp 423.10bp 44.37bp 218.15bp},clip,width=0.92\textwidth,max height=0.38\textheight]{Ishii.pdf}
\end{sourceblock}
其中 T C 和 f k 分别是系数和液相的总湍流动能。方程 (12-98) 适用于带有小气泡的剪切引起的湍流。

对于 e b τ τ 或 ec et τ τ，方程（12-96）可以简化为（Lopez de Bertodano et al., 2006）

\begin{sourceblock}
\includegraphics[page=365,pagebox=mediabox,trim={45.00bp 297.76bp 44.37bp 329.57bp},clip,width=0.92\textwidth,max height=0.38\textheight]{Ishii.pdf}
\end{sourceblock}
其中 BI t ν 是涡流扩散率，可以使用 Sato 等人的模型计算。 （1981）。方程 (12-99) 适用于由大于剪切引起的涡流的盖泡产生的湍流运动。洛佩斯·德·贝托达诺等人。 (2006)应用方程(12-99)来限制气泡流。导出湍流色散力的另一种方法是基于相间动量力波动部分的时间平均。伯恩斯等人。 (2004) 根据相间阻力的 Favre 平均值导出了湍流色散力：

\begin{sourceblock}
\includegraphics[page=365,pagebox=mediabox,trim={45.00bp 128.47bp 44.37bp 490.48bp},clip,width=0.92\textwidth,max height=0.38\textheight]{Ishii.pdf}
\end{sourceblock}
其中 t ν 和 Sct N 分别是湍流运动粘度和湍流施密特数 (=0.9)。相间阻力的动量传递系数定义为

\begin{sourceblock}
\includegraphics[page=366,pagebox=mediabox,trim={45.00bp 569.80bp 44.37bp 53.32bp},clip,width=0.92\textwidth,max height=0.38\textheight]{Ishii.pdf}
\end{sourceblock}
这里需要注意的是，式(12-100)与式(12-99)类似。尚未获得对湍流色散力的全面理解。

\section*{12.4 多粒子系统中的湍流}\phantomsection\addcontentsline{toc}{section}{12.4 多粒子系统中的湍流}
在单相流分析中，湍流结构已被广泛研究，并提出了几种湍流模型。湍流模型通常分为零方程模型、一方程模型、二方程模型和应力方程模型。纳维-斯托克斯方程的大涡模拟和直接数值模拟也是可能的。然而，在两相流分析中，由于两相流湍流的复杂性，对湍流建模的研究有限。下面，简要解释气泡流态中的一些初步湍流模型。在模型中，由于气相和液相之间的密度差较大，通常忽略气相的湍流动能。为了强调时间平均，对符号应用了上划线。 Akiyama 和 Aritomi (2002) 也对这个主题进行了综述。

A. 零方程模型 零方程湍流模型通常表示为没有微分方程来确定雷诺应力的模型。接下来，佐藤等人提出的模型。 （1981）将作为零方程湍流模型的例子进行简要解释。

考虑二维完全发展的气泡流，例如垂直管道中或两个平行平坦壁之间的流。 y 轴和 z 轴分别垂直且平行于主流方向。该模型假设 (i) 只有液相有助于动量传递，(ii) 液相中有两种独立于气泡搅拌和依赖于气泡搅拌的湍流。然后，y 和 z 方向上的速度 , f y v 和 , f z v 表示为

\begin{sourceblock}
\includegraphics[page=366,pagebox=mediabox,trim={45.00bp 108.15bp 44.37bp 510.70bp},clip,width=0.92\textwidth,max height=0.38\textheight]{Ishii.pdf}
\end{sourceblock}
\begin{sourceblock}
\includegraphics[page=366,pagebox=mediabox,trim={45.00bp 63.15bp 44.37bp 577.05bp},clip,width=0.92\textwidth,max height=0.38\textheight]{Ishii.pdf}
\end{sourceblock}
f τ，可表示为

\begin{sourceblock}
\includegraphics[page=367,pagebox=mediabox,trim={45.00bp 517.95bp 44.37bp 122.25bp},clip,width=0.92\textwidth,max height=0.38\textheight]{Ishii.pdf}
\end{sourceblock}
其中 y 是从通道壁测量的法向距离。这里，涡流扩散率 ε ′ 和 ε ′′ 定义为

\begin{sourceblock}
\includegraphics[page=367,pagebox=mediabox,trim={45.00bp 430.46bp 44.37bp 193.53bp},clip,width=0.92\textwidth,max height=0.38\textheight]{Ishii.pdf}
\end{sourceblock}
\begin{sourceblock}
\includegraphics[page=367,pagebox=mediabox,trim={45.00bp 357.50bp 44.37bp 241.75bp},clip,width=0.92\textwidth,max height=0.38\textheight]{Ishii.pdf}
\end{sourceblock}
因此，只要给出 ε ' 和 ε '' 即可计算湍流应力分布。

对于充分发展的湍流气泡流，独立于气泡搅动的涡流扩散率 ε ′ 可以根据普朗特混合长度理论和靠近光滑壁区域的阻尼因子凭经验确定：

\begin{sourceblock}
\includegraphics[page=367,pagebox=mediabox,trim={45.00bp 171.41bp 44.37bp 406.57bp},clip,width=0.92\textwidth,max height=0.38\textheight]{Ishii.pdf}
\end{sourceblock}
这里，* f f y yv ν + ϋ 和 * f f R Rv ν + ϋ 其中 * f v 和 R 分别是由 W f τ ρ 和管道或一半管道的半径定义的摩擦速度

\begin{sourceblock}
\includegraphics[page=367,pagebox=mediabox,trim={114.03bp 126.79bp 205.88bp 498.31bp},clip,width=\textwidth,max height=0.72\textheight]{Ishii.pdf}
\end{sourceblock}
具有平行平坦壁的通道的宽度。取决于气泡搅动的涡流扩散率 ε '' 可以根据自由湍流的虚拟运动粘度（例如固体后面的尾流）和靠近光滑壁的区域中的阻尼因子凭经验确定，如下所示

\begin{sourceblock}
\includegraphics[page=368,pagebox=mediabox,trim={45.00bp 562.31bp 44.37bp 57.19bp},clip,width=0.92\textwidth,max height=0.38\textheight]{Ishii.pdf}
\end{sourceblock}
其中B d 和r v ∞ 分别是由式(12-109)给出的气泡的平均直径和静止液体中的终端速度。

\begin{sourceblock}
\includegraphics[page=368,pagebox=mediabox,trim={45.00bp 438.26bp 44.37bp 145.35bp},clip,width=0.92\textwidth,max height=0.38\textheight]{Ishii.pdf}
\end{sourceblock}
where m
B d 是气泡的横截面平均直径。湍流应力是根据上述方程和边界条件计算得出的。

2.2 一方程模型 一方程湍流模型通常表示为只有一个湍流动能守恒微分方程和混合长度及其他湍流源项的本构方程的模型。下面，以 Kataoka 和 Serizawa (1995) 提出的模型作为一方程湍流模型的例子进行简要说明。

考虑圆管中稳定、充分发展的绝热气泡流。 y 轴和 z 轴分别垂直且平行于主流方向。湍流应力，T

\begin{sourceblock}
\includegraphics[page=368,pagebox=mediabox,trim={45.00bp 199.40bp 44.37bp 403.53bp},clip,width=0.92\textwidth,max height=0.38\textheight]{Ishii.pdf}
\end{sourceblock}
\begin{sourceblock}
\includegraphics[page=368,pagebox=mediabox,trim={45.00bp 141.16bp 44.37bp 490.23bp},clip,width=0.92\textwidth,max height=0.38\textheight]{Ishii.pdf}
\end{sourceblock}
其中 f ′ v 是液体速度脉动矢量。因此，只要给定 TP l 和 t v，就可以计算湍流应力分布。假设两相流的混合长度由剪切引起的湍流的混合长度的线性叠加表示

\begin{sourceblock}
\includegraphics[page=369,pagebox=mediabox,trim={45.00bp 547.83bp 44.37bp 96.43bp},clip,width=0.92\textwidth,max height=0.38\textheight]{Ishii.pdf}
\end{sourceblock}
单相液体流中剪切引起的湍流的混合长度由下式给出

\begin{sourceblock}
\includegraphics[page=369,pagebox=mediabox,trim={45.00bp 452.31bp 44.37bp 161.69bp},clip,width=0.92\textwidth,max height=0.38\textheight]{Ishii.pdf}
\end{sourceblock}
基于机械模型给出了气泡引起的湍流的混合长度，其中相同体积的液体伴随着气泡通过控制面的湍流运动而交换：

\begin{sourceblock}
\includegraphics[page=369,pagebox=mediabox,trim={45.00bp 250.50bp 44.37bp 268.52bp},clip,width=0.92\textwidth,max height=0.38\textheight]{Ishii.pdf}
\end{sourceblock}
其中 B d 是气泡直径。湍流速度 t v 由液相湍流动能方程计算得出：

\begin{sourceblock}
\includegraphics[page=369,pagebox=mediabox,trim={45.00bp 63.78bp 44.37bp 467.99bp},clip,width=0.92\textwidth,max height=0.38\textheight]{Ishii.pdf}
\end{sourceblock}
\begin{sourceblock}
\includegraphics[page=370,pagebox=mediabox,trim={45.00bp 563.69bp 44.37bp 57.65bp},clip,width=0.92\textwidth,max height=0.38\textheight]{Ishii.pdf}
\end{sourceblock}
其中 k 是液相的湍流动能，由下式给出
2 ′ ′ = ⋅ v v f f k。 ( ) 1 0 56。 β = , ( ) 2 0 38。 β = , 1 0 18。 γ = , ( ) 1 0 075。 K =，并且 ( ) 2 1 0。 K = 是系数。在该模型中，湍流速度由式（12-111）确定，假设等边湍流为

\begin{sourceblock}
\includegraphics[page=370,pagebox=mediabox,trim={45.00bp 450.45bp 44.37bp 178.61bp},clip,width=0.92\textwidth,max height=0.38\textheight]{Ishii.pdf}
\end{sourceblock}
方程中各项的意义如下。右侧的第一项、第二项和第三项分别表示湍流扩散、由于剪切而产生的湍流以及湍流​​耗散。第四项和第五项分别表示由于小尺度界面引起的湍流吸收和由于气泡相对运动而产生的湍流。最后一项表示非常靠近壁的数值误差的补偿。湍流应力是根据上述方程和边界条件计算得出的。

C. 二方程模型 二方程湍流模型通常表示为用两个微分方程来确定雷诺应力的模型。接下来，由 Lopez de Bertodano 等人提出的 k - ε 模型。 （1994）将作为二方程湍流模型的例子进行简要解释。

考虑稳定、充分发展的绝热稀气泡流。假定湍流应力张量 T f T 由剪切引起的 (SI) 湍流应力张量 SI f T 和气泡引起的 (BI) 湍流张量 BI f T 的线性叠加表示为

\begin{sourceblock}
\includegraphics[page=370,pagebox=mediabox,trim={45.00bp 172.18bp 44.37bp 452.47bp},clip,width=0.92\textwidth,max height=0.38\textheight]{Ishii.pdf}
\end{sourceblock}
剪切引起的湍流应力由下式计算

\begin{sourceblock}
\includegraphics[page=370,pagebox=mediabox,trim={45.00bp 100.69bp 44.37bp 528.53bp},clip,width=0.92\textwidth,max height=0.38\textheight]{Ishii.pdf}
\end{sourceblock}
其中 t ν 、 A 和 SI k 分别是湍流运动粘度、湍流各向异性张量和剪切引起的湍流引起的湍流动能。对于各向同性湍流，= A I。

气泡引起的湍流计算如下

\begin{sourceblock}
\includegraphics[page=371,pagebox=mediabox,trim={45.00bp 545.53bp 44.37bp 80.52bp},clip,width=0.92\textwidth,max height=0.38\textheight]{Ishii.pdf}
\end{sourceblock}
\begin{sourceblock}
\includegraphics[page=371,pagebox=mediabox,trim={45.00bp 454.11bp 44.37bp 138.65bp},clip,width=0.92\textwidth,max height=0.38\textheight]{Ishii.pdf}
\end{sourceblock}
其中 vm C 是虚拟体积系数，围绕球体的势流值为 1/2。对于低空隙率和高空隙率情况，建议分别使用 2.0 和 1.2 值。由气泡引起的湍流引起的湍流动能 BI k 可由下式获得

\begin{sourceblock}
\includegraphics[page=371,pagebox=mediabox,trim={45.00bp 344.53bp 44.37bp 260.77bp},clip,width=0.92\textwidth,max height=0.38\textheight]{Ishii.pdf}
\end{sourceblock}
然后，将式（12-119）和（12-118）代入式（12-117），得

\begin{sourceblock}
\includegraphics[page=371,pagebox=mediabox,trim={45.00bp 209.55bp 44.37bp 354.23bp},clip,width=0.92\textwidth,max height=0.38\textheight]{Ishii.pdf}
\end{sourceblock}
因此，只要给出 t ν 和 SI k，就可以计算剪应力分布。

\begin{sourceblock}
\includegraphics[page=371,pagebox=mediabox,trim={45.00bp 98.98bp 44.37bp 515.64bp},clip,width=0.92\textwidth,max height=0.38\textheight]{Ishii.pdf}
\end{sourceblock}
由剪切引起的湍流和气泡引起的湍流引起的湍流运动粘度分别如式(12-123)和(12-124)所示。

\begin{sourceblock}
\includegraphics[page=372,pagebox=mediabox,trim={45.00bp 562.53bp 44.37bp 64.06bp},clip,width=0.92\textwidth,max height=0.38\textheight]{Ishii.pdf}
\end{sourceblock}
其中 SI ε 是剪切引起的湍流的耗散。

\begin{sourceblock}
\includegraphics[page=372,pagebox=mediabox,trim={45.00bp 494.65bp 44.37bp 134.68bp},clip,width=0.92\textwidth,max height=0.38\textheight]{Ishii.pdf}
\end{sourceblock}
剪切引起的湍流引起的湍流动能和耗散可以通过式（12-125）给出的剪切引起的湍流动能传递方程和式（12-127）给出的剪切引起的湍流耗散率传递方程计算。

\begin{sourceblock}
\includegraphics[page=372,pagebox=mediabox,trim={45.00bp 380.75bp 44.37bp 240.83bp},clip,width=0.92\textwidth,max height=0.38\textheight]{Ishii.pdf}
\end{sourceblock}
其中 ( ) 1 0。 k σ = 是一个常数，SI P 是剪切引起的湍流的产生，由下式给出

\begin{sourceblock}
\includegraphics[page=372,pagebox=mediabox,trim={45.00bp 303.51bp 44.37bp 333.83bp},clip,width=0.92\textwidth,max height=0.38\textheight]{Ishii.pdf}
\end{sourceblock}
\begin{sourceblock}
\includegraphics[page=372,pagebox=mediabox,trim={45.00bp 238.21bp 44.37bp 382.73bp},clip,width=0.92\textwidth,max height=0.38\textheight]{Ishii.pdf}
\end{sourceblock}
其中 ( ) 1 3。 ε σ = , ( ) 1 1 44。 C ε = , 和 ( ) 2 1 92。 C ε = 是常数。湍流应力是根据上述方程和边界条件计算得出的。

\chapter{漂移通量模型}
漂移通量模型的基本概念是将混合物视为一个整体，而不是单独考虑两个相。显然，漂移通量模型公式比二流体模型更简单，但是它需要一些严格的本构假设，导致两相流的一些重要特征丢失。然而，正是漂移通量模型的这种简单性使其在许多工程应用中非常有用。正如两相流系统动力学分析的情况一样，工程问题中所需的信息通常是总混合物的响应，而不是每个组成相的响应（Tong，1965）。此外，如果这些混合物响应已知，则可以轻松地对每个相的局部行为进行详细分析。

漂移通量模型的另一个重要方面涉及系统的缩放，该系统可直接应用于两相流实验和工程系统的规划和设计。通过使用漂移通量模型公式和混合物特性可以有效地研究两个不同系统的相似性。与双流体模型相比，漂移通量模型最重要的方面是减少了公式中所需的场和本构方程的总数。漂移通量模型用四个场方程表示：混合连续性；势头;能量方程；和气体连续性方程。

因此，可以看出，漂移通量模型遵循用于分析气体混合物或可混溶液体的动力学的标准方法。人们普遍认为，漂移通量模型适用于两个组分的动力学紧密耦合的混合物。这表明相同的论点可以用于宏观两相流。漂移通量模型在许多实际工程系统中的有用性来自于这样一个事实：即使是局部弱耦合的两相混合物也可以被考虑，因为相对较大的轴向
在漂移通量模型公式中，我们只有四个场方程，因此，从原来的六个场方程中消除了一个能量方程和一个动量方程。然后，相对运动和能量差应该用附加的本构方程来表达。换句话说，动态的相互作用关系被本构规律所取代。此外，重要的是根据混合物质心制定模型，以保留广泛变量的加性特征，如第 4 章所述。

在本章中，我们开发了混合模型的一般公式（Ishii，1975），然后讨论在实际应用中重要的各种特殊情况（Ishii，1977）。由于我们在第 9 章中对双流体模型的场和本构方程进行了详细分析，因此只要有帮助，我们就会回顾并使用这些结果来建立漂移通量模型公式。下图总结了漂移通量模型公式的建立。在这里，我们看到了两个阶段之间的运动学、机械和热关系的特殊重要性。显然，从公式中消除两个动量方程之一需要相之间的运动学关系，因此，相对速度应由本构关系给出。类似地，通过仅使用混合能量方程来平衡两相流中的能量，应该给出相之间的热关系。

\section*{13.1 漂移通量模型场方程}\phantomsection\addcontentsline{toc}{section}{13.1 漂移通量模型场方程}
基于质心和漂移通量速度的公式 漂移通量模型的四个基本场方程的最一般形式已在第 5 章的 1.3 节中给出。在本节中，我们首先通过对双流体模型本构方程的一些分析，将这些方程转化为更实际的形式。然后我们讨论一些对实际应用很重要的适当简化。在这里，我们根据混合物连续性、动量和热能方程以及其中一相的连续性方程来制定模型。这些方程可以简化为以下形式

\begin{sourceblock}
\includegraphics[page=375,pagebox=mediabox,trim={66.45bp 320.59bp 67.52bp 79.69bp},clip,width=\textwidth,max height=0.72\textheight]{Ishii.pdf}
\end{sourceblock}
由式(5-40)得到的混合连续性方程

\begin{sourceblock}
\includegraphics[page=375,pagebox=mediabox,trim={45.00bp 233.95bp 44.37bp 387.88bp},clip,width=0.92\textwidth,max height=0.38\textheight]{Ishii.pdf}
\end{sourceblock}
式（5-41）中第2相的连续性方程

\begin{sourceblock}
\includegraphics[page=375,pagebox=mediabox,trim={45.00bp 154.06bp 44.37bp 467.81bp},clip,width=0.92\textwidth,max height=0.38\textheight]{Ishii.pdf}
\end{sourceblock}
\begin{sourceblock}
\includegraphics[page=375,pagebox=mediabox,trim={45.00bp 129.44bp 44.37bp 517.72bp},clip,width=0.92\textwidth,max height=0.38\textheight]{Ishii.pdf}
\end{sourceblock}
\begin{sourceblock}
\includegraphics[page=375,pagebox=mediabox,trim={45.00bp 50.14bp 44.37bp 540.22bp},clip,width=0.92\textwidth,max height=0.38\textheight]{Ishii.pdf}
\end{sourceblock}
\begin{sourceblock}
\includegraphics[page=376,pagebox=mediabox,trim={45.00bp 556.96bp 44.37bp 84.28bp},clip,width=0.92\textwidth,max height=0.38\textheight]{Ishii.pdf}
\end{sourceblock}
式(13-3)中，体力场被取为常数。混合物动量方程右侧的最后一项表示表面张力对混合物的影响。

正如我们在第 5 章 1.3 节中讨论的那样，获得适当的混合物热能方程存在相当大的困难。那里已经表明，我们有两种不同的方法来获得该方程。将各相的热能方程相加，得到式（5-53）。因此，根据式（9-154）和（9-150），我们有

\begin{sourceblock}
\includegraphics[page=376,pagebox=mediabox,trim={45.00bp 267.80bp 44.37bp 247.90bp},clip,width=0.92\textwidth,max height=0.38\textheight]{Ishii.pdf}
\end{sourceblock}
根据混合物性质和界面动量传递条件的定义，我们有

\begin{sourceblock}
\includegraphics[page=376,pagebox=mediabox,trim={45.00bp 179.38bp 44.37bp 442.63bp},clip,width=0.92\textwidth,max height=0.38\textheight]{Ishii.pdf}
\end{sourceblock}
为了简单起见，我们定义三种不同的效果如下

\begin{sourceblock}
\includegraphics[page=376,pagebox=mediabox,trim={45.00bp 106.23bp 44.37bp 522.97bp},clip,width=0.92\textwidth,max height=0.38\textheight]{Ishii.pdf}
\end{sourceblock}
\begin{sourceblock}
\includegraphics[page=377,pagebox=mediabox,trim={45.00bp 567.46bp 44.37bp 54.40bp},clip,width=0.92\textwidth,max height=0.38\textheight]{Ishii.pdf}
\end{sourceblock}
and

\begin{sourceblock}
\includegraphics[page=377,pagebox=mediabox,trim={45.00bp 487.83bp 44.37bp 133.75bp},clip,width=0.92\textwidth,max height=0.38\textheight]{Ishii.pdf}
\end{sourceblock}
那么热能方程（13-5）简化为

\begin{sourceblock}
\includegraphics[page=377,pagebox=mediabox,trim={45.00bp 374.80bp 44.37bp 210.54bp},clip,width=0.92\textwidth,max height=0.38\textheight]{Ishii.pdf}
\end{sourceblock}
当上述方程的最后三项不能被忽略时，混合热能方程的漂移通量模型面临相当大的困难。这些术语源自粘性耗散、表面张力作用和界面机械能传递。如果项 Φ i
式（13-9）给出的m对热能交换有显著贡献。显然，在这种情况下，漂移通量模型需要相对速度以及界面机械能传递的本构方程。

将式（9-141）和（9-127）代入式（5-55），然后使用式（9-145）给出的近似，可以得到热能方程的另一种形式。这两种不同形式的热能方程并没有给出相同的漂移通量模型公式，正如第 5 章 1.3 节中所讨论的那样。然而，我们使用方程（13-10）给出的关系，该关系是通过将每个相的热能方程相加而获得的，因为它比基于方程（5-55）的形式具有相对简单的形式。

上述四个场方程(13-1)、(13-2)、(13-3)和(13-10)陈述了控制宏观混合场的平衡定律。它们是通过将时间平均应用于具有界面不连续性的两相流系统而获得的。我们注意到混合物连续性、动量和能量方程与单相流的方程有些相似。实际上，混合连续性方程的形式与没有内部不连续性的连续体的形式完全相同。这是通过使用正确定义的混合物属性来完成的。混合物动量
( ) Τ + T T。在混合热能方程中，我们有单相流方程中没有出现的三个附加项。方程（13-10）右侧的第二项是由于各相相对于混合物质心的相对运动而传输能量而产生的能量扩散。回顾方程（13-8）和（13-9），由 m 给出的项
σ Φ 和 i m Φ 分别表示表面张力效应和界面机械能传递的贡献。在正常情况下，这两项和能量耗散项几乎总是可以忽略不计。

方程（13-3）和（13-10）给出的混合动量和热能方程描述了静止观察者的动量和能量交换。因此，具有混合物质心速度的对流通量和相对于混合物重心定义的附加扩散通量出现在方程中。将这两个方程转化为式(7-14)的对流导数如下

\begin{sourceblock}
\includegraphics[page=378,pagebox=mediabox,trim={45.00bp 232.08bp 44.37bp 358.61bp},clip,width=0.92\textwidth,max height=0.38\textheight]{Ishii.pdf}
\end{sourceblock}
and

\begin{sourceblock}
\includegraphics[page=378,pagebox=mediabox,trim={45.00bp 119.84bp 44.37bp 468.63bp},clip,width=0.92\textwidth,max height=0.38\textheight]{Ishii.pdf}
\end{sourceblock}
上述两个方程描述了从以速度 m v 移动的观察者看来动量和能量的传递。由于其特殊的形式，式(13-11)被称为运动方程。

几个坐标系中的场方程鉴于实际重要性，我们将这四个场方程表达在两个不同的坐标系中。由于导数可以直接从标准向量微积分中得出（Aris，1962；McConnell，1957），因此我们仅在下面列出结果。

在直角坐标 ( ) , , x y z 中，我们有混合物质量守恒

\begin{sourceblock}
\includegraphics[page=379,pagebox=mediabox,trim={45.00bp 463.21bp 44.37bp 158.57bp},clip,width=0.92\textwidth,max height=0.38\textheight]{Ishii.pdf}
\end{sourceblock}
相 2 的质量漂移通量

\begin{sourceblock}
\includegraphics[page=379,pagebox=mediabox,trim={45.00bp 315.88bp 44.37bp 233.34bp},clip,width=0.92\textwidth,max height=0.38\textheight]{Ishii.pdf}
\end{sourceblock}
用于保持混合物动量守恒

x -component

\begin{sourceblock}
\includegraphics[page=379,pagebox=mediabox,trim={45.00bp 151.63bp 44.37bp 404.80bp},clip,width=0.92\textwidth,max height=0.38\textheight]{Ishii.pdf}
\end{sourceblock}
y -component

\begin{sourceblock}
\includegraphics[page=379,pagebox=mediabox,trim={45.00bp 79.23bp 44.37bp 545.74bp},clip,width=0.92\textwidth,max height=0.38\textheight]{Ishii.pdf}
\end{sourceblock}
\begin{sourceblock}
\includegraphics[page=380,pagebox=mediabox,trim={45.00bp 498.61bp 44.37bp 55.60bp},clip,width=0.92\textwidth,max height=0.38\textheight]{Ishii.pdf}
\end{sourceblock}
z -component

\begin{sourceblock}
\includegraphics[page=380,pagebox=mediabox,trim={45.00bp 326.83bp 44.37bp 195.76bp},clip,width=0.92\textwidth,max height=0.38\textheight]{Ishii.pdf}
\end{sourceblock}
用于混合物热能平衡

\begin{sourceblock}
\includegraphics[page=380,pagebox=mediabox,trim={45.00bp 50.14bp 44.37bp 366.95bp},clip,width=0.92\textwidth,max height=0.38\textheight]{Ishii.pdf}
\end{sourceblock}
类似地，在柱坐标 ( ) , , r z θ 中，我们有混合物质量守恒

\begin{sourceblock}
\includegraphics[page=381,pagebox=mediabox,trim={45.00bp 461.09bp 44.37bp 132.35bp},clip,width=0.92\textwidth,max height=0.38\textheight]{Ishii.pdf}
\end{sourceblock}
相 2 的质量漂移通量

\begin{sourceblock}
\includegraphics[page=381,pagebox=mediabox,trim={45.00bp 325.72bp 44.37bp 233.51bp},clip,width=0.92\textwidth,max height=0.38\textheight]{Ishii.pdf}
\end{sourceblock}
用于保持混合物动量守恒

r -component

\begin{sourceblock}
\includegraphics[page=381,pagebox=mediabox,trim={45.00bp 130.97bp 44.37bp 396.23bp},clip,width=0.92\textwidth,max height=0.38\textheight]{Ishii.pdf}
\end{sourceblock}
θ -component

\begin{sourceblock}
\includegraphics[page=381,pagebox=mediabox,trim={45.00bp 61.23bp 44.37bp 563.74bp},clip,width=0.92\textwidth,max height=0.38\textheight]{Ishii.pdf}
\end{sourceblock}
\begin{sourceblock}
\includegraphics[page=382,pagebox=mediabox,trim={45.00bp 503.85bp 44.37bp 54.82bp},clip,width=0.92\textwidth,max height=0.38\textheight]{Ishii.pdf}
\end{sourceblock}
z -component

\begin{sourceblock}
\includegraphics[page=382,pagebox=mediabox,trim={45.00bp 337.20bp 44.37bp 190.01bp},clip,width=0.92\textwidth,max height=0.38\textheight]{Ishii.pdf}
\end{sourceblock}
用于混合物热能平衡

\begin{sourceblock}
\includegraphics[page=382,pagebox=mediabox,trim={45.00bp 64.26bp 44.37bp 370.27bp},clip,width=0.92\textwidth,max height=0.38\textheight]{Ishii.pdf}
\end{sourceblock}
\section*{13.2 漂移-通量（或混合）模型本构关系}\phantomsection\addcontentsline{toc}{section}{13.2 漂移-通量（或混合）模型本构关系}
可见，基于四个场方程的漂移-通量模型是二流体模型的近似理论。为了完成漂移通量模型，有必要提供混合物的几个本构关系。我们可以考虑两种不同的方法来实现我们的目的。第一种方法是从混合场方程和混合熵不等式开始对本构方程进行分析，然后将各种本构公理直接应用于混合物并且独立于二流体模型。第二种方法是通过双流体模型公式化简得到必要的本构方程。

起初，遵循前一种方法似乎更合乎逻辑，因为它是混合物模型的自足且独立的公式。但事实上，它也面临着巨大的困难，不容忽视。主要问题是由于两相通常不处于热平衡，因此不可能引入混合温度。这表明我们不能期望在宏观混合物性质方面存在简单的状态方程。

可见，热非平衡条件和界面结构是相变的控制因素。此外，两相之间的运动学和机械状态很大程度上受到界面性质和结构的影响。为了将这些重要的影响带入漂移通量模型公式中，使用双流体模型的简化比前一种方法更简单、更现实。因此，在本节中，我们将与第 9 章的 1.2 节并行地对混合本构方程进行分析。

确定性原理 式(13-1)、(13-2)、(13-3)和(13-10)给出了漂移-通量模型场方程，但不足以完全描述系统。从决定论原理来看，有必要提供额外的本构方程来指定某些宏观两相混合物的响应特征。为了保持
考虑到上述考虑，我们在漂移通量模型公式中出现了以下变量。

1. 状态方程： l , , , , , ; k k k k i p i ρ Τ Τ σ 2. 质量守恒： , ; m m ρ v 3. 动量守恒： , , , , , ; m km m k p Τ α V M T T 4. 能量守恒： , , , , , ; i m m m m i Τ μ σ Φ Φ Φ q q 5. 漂移通量方程： 2； Г 其中​​ k = 1 和 2。因此变量总数为二十七个。对于正确制定的​​漂移通量模型，我们也应该有相同数量的方程。这些可以分为以下几组。

方程 方程数 1) 场方程 混合质量 Eq.(13-1) 1 混合动量 Eq.(13-3) 1 混合能量 Eq.(13-10) 1 漂移通量 Eq.(13-2) 1 2) 连续性公理 1 2 1 α α = − Eq.(4-13) 1 3) 相热状态方程的状态方程Eq.(9-56) 2 热态方程 Eq.(9-57) 2 m ρ 的定义 Eq.(4-66) 1 m i 的定义 Eq.(4-74) 1 m p 的定义 Eq.(4-72) 1 4) 界面状态方程 ( ) i σ σ Τ = Eq.(9-130) 1 5) 漂移通量恒等式速度

\begin{sourceblock}
\includegraphics[page=384,pagebox=mediabox,trim={45.00bp 78.81bp 44.37bp 556.32bp},clip,width=0.92\textwidth,max height=0.38\textheight]{Ishii.pdf}
\end{sourceblock}
\begin{sourceblock}
\includegraphics[page=385,pagebox=mediabox,trim={45.00bp 478.27bp 44.37bp 146.22bp},clip,width=0.92\textwidth,max height=0.38\textheight]{Ishii.pdf}
\end{sourceblock}
10) 力学本构方程
粘性应力 T 1 湍流应力 T T 1 混合动量源 m M 1 11) 能量本构方程 传导热通量 q 1 湍流热通量 T q 1 耗散项 m
μ Φ 1 表面张力效应 m
σ Φ 1 机械能效应 i m Φ 1 12) 相变质量产生的本构方程 2 Γ 1 这表明我们也有二十七个方程，因此目前的公式是一致的。我们在这里注意到，这些场和本构方程是决定论原理所必需的，但是它们并不能确保解的存在。对于我们的系统来说，证明问题是正确设置的，即公正设置（Truesdell and Toupin，1960）是非常困难的，因为它涉及到存在性、唯一性以及适当的初始条件和边界条件。通常，只能针对非常简单的问题类别进行检查。现在我们对上述本构方程进行详细讨论。

状态方程和混合性质 连续性公理要求界面不会在有限的时间间隔内停留在一点（参见第 5 章的 1.4 节），因此我们有

\begin{sourceblock}
\includegraphics[page=385,pagebox=mediabox,trim={45.00bp 75.96bp 44.37bp 565.35bp},clip,width=0.92\textwidth,max height=0.38\textheight]{Ishii.pdf}
\end{sourceblock}
混合物密度 m ρ 定义为

\begin{sourceblock}
\includegraphics[page=386,pagebox=mediabox,trim={45.00bp 584.14bp 44.37bp 57.09bp},clip,width=0.92\textwidth,max height=0.38\textheight]{Ishii.pdf}
\end{sourceblock}
每个相的热状态方程由下式给出

\begin{sourceblock}
\includegraphics[page=386,pagebox=mediabox,trim={45.00bp 521.47bp 44.37bp 115.73bp},clip,width=0.92\textwidth,max height=0.38\textheight]{Ishii.pdf}
\end{sourceblock}
\begin{sourceblock}
\includegraphics[page=386,pagebox=mediabox,trim={45.00bp 469.38bp 44.37bp 167.77bp},clip,width=0.92\textwidth,max height=0.38\textheight]{Ishii.pdf}
\end{sourceblock}
混合物压力与相压力的关系为

\begin{sourceblock}
\includegraphics[page=386,pagebox=mediabox,trim={45.00bp 410.14bp 44.37bp 231.09bp},clip,width=0.92\textwidth,max height=0.38\textheight]{Ishii.pdf}
\end{sourceblock}
由式(13-22)、式(13-23)和式(13-24)给出的混合物密度可以被认为是混合物热状态方程，其受到两相之间的热、机械和化学关系的约束。

混合物的热函定义为

\begin{sourceblock}
\includegraphics[page=386,pagebox=mediabox,trim={45.00bp 283.78bp 44.37bp 338.36bp},clip,width=0.92\textwidth,max height=0.38\textheight]{Ishii.pdf}
\end{sourceblock}
每个相的热量状态方程

\begin{sourceblock}
\includegraphics[page=386,pagebox=mediabox,trim={45.00bp 220.38bp 44.37bp 415.95bp},clip,width=0.92\textwidth,max height=0.38\textheight]{Ishii.pdf}
\end{sourceblock}
\begin{sourceblock}
\includegraphics[page=386,pagebox=mediabox,trim={45.00bp 168.37bp 44.37bp 468.00bp},clip,width=0.92\textwidth,max height=0.38\textheight]{Ishii.pdf}
\end{sourceblock}
将式(13-27)和式(13-28)代入式(13-26)，得到混合热量状态方程，该方程显示了m i 对温度、压力和局部空隙率的依赖性。然而，与混合物热状态方程的情况一样，它受到两相之间的温差和压差本构方程以及相变条件方程（9-163）的约束。

宏观场中界面的热状态方程可以近似为

\begin{sourceblock}
\includegraphics[page=387,pagebox=mediabox,trim={45.00bp 540.18bp 44.37bp 96.97bp},clip,width=0.92\textwidth,max height=0.38\textheight]{Ishii.pdf}
\end{sourceblock}
它可以被认为是无表面质量界面的基本状态方程。

运动学本构方程正如上一节所解释的，我们应该为相的相对运动提供一个本构方程。由于在漂移通量模型公式中我们消除了一个动量方程，因此运动学本构方程作为相对运动方程。然而，可以预期，通过用运动学关系代替动量方程，两相之间的动态相互作用将会消失。

每个相的扩散速度通过恒等式相关

\begin{sourceblock}
\includegraphics[page=387,pagebox=mediabox,trim={45.00bp 350.92bp 44.37bp 290.31bp},clip,width=0.92\textwidth,max height=0.38\textheight]{Ishii.pdf}
\end{sourceblock}
因此，我们应该通过运动学本构方程仅提供一种扩散速度。然而，由于扩散速度 km V 可以通过第 4 章 1.6 节中的定义与相之间的相对速度或漂移速度相关：

\begin{sourceblock}
\includegraphics[page=387,pagebox=mediabox,trim={45.00bp 201.70bp 44.37bp 386.37bp},clip,width=0.92\textwidth,max height=0.38\textheight]{Ishii.pdf}
\end{sourceblock}
本构方程可以用上述任何速度给出。

两相之间的相对速度取决于作用在界面上的拖曳力以及界面几何形状。因此，可以预期，只要混合物的界面结构发生变化，相对速度就会发生变化。第 9 章 1.4 节表明，在分散两相流中，阻力相关性应以漂移速度 l ( ) d − j v 和基于该速度的雷诺数来表示，方程（9-223）和（9-224）。这表明相间相对运动的运动学本构方程最好根据分散相的漂移速度来研究，正如 Zuber (1964b)、Zuber 等人提出的那样。 (1964) 以及 Zuber 和 Staub (1966)。

鉴于上述参考文献中获得的结果，可以说漂移速度是无限介质中单个颗粒的终端速度 ∞ v 和连续相的空隙率的函数。为了考虑浓度梯度引起的漂移效应，我们提出了以下形式的线性本构律

\begin{sourceblock}
\includegraphics[page=388,pagebox=mediabox,trim={45.00bp 491.49bp 44.37bp 132.48bp},clip,width=0.92\textwidth,max height=0.38\textheight]{Ishii.pdf}
\end{sourceblock}
其中 d D α 是基于空隙率 d α 的漂移系数。右侧第一项考虑了重力和力的影响，这通常是漂移速度的主要部分。 Zuber 等人对气泡流态中的这一术语进行了详细分析。 (1964) 和 Ishii (1977) 还证明了漂移速度的本构方程可以从各种流态的双流体模型中推导出来。在没有壁并且在没有相变的稳态条件下，无限介质中的多粒子系统本质上减少为重力和阻力主导的一维流，因为平均空隙和速度分布变得平坦。求解每个相的动量方程即可得出相对速度定律。因此，我们在第 12 章中使用阻力相关性的结果。

对于粘性状态，漂移速度可由下式给出

\begin{sourceblock}
\includegraphics[page=388,pagebox=mediabox,trim={45.00bp 211.96bp 44.37bp 363.26bp},clip,width=0.92\textwidth,max height=0.38\textheight]{Ishii.pdf}
\end{sourceblock}
where

\begin{sourceblock}
\includegraphics[page=388,pagebox=mediabox,trim={45.00bp 140.62bp 44.37bp 483.69bp},clip,width=0.92\textwidth,max height=0.38\textheight]{Ishii.pdf}
\end{sourceblock}
and

\begin{sourceblock}
\includegraphics[page=388,pagebox=mediabox,trim={45.00bp 73.42bp 44.37bp 559.76bp},clip,width=0.92\textwidth,max height=0.38\textheight]{Ishii.pdf}
\end{sourceblock}
\begin{sourceblock}
\includegraphics[page=389,pagebox=mediabox,trim={45.00bp 563.74bp 44.37bp 57.74bp},clip,width=0.92\textwidth,max height=0.38\textheight]{Ishii.pdf}
\end{sourceblock}
\begin{sourceblock}
\includegraphics[page=389,pagebox=mediabox,trim={45.00bp 447.60bp 44.37bp 109.40bp},clip,width=0.92\textwidth,max height=0.38\textheight]{Ishii.pdf}
\end{sourceblock}
对于扭曲的流体粒子状态，漂移速度可以由下式给出

\begin{sourceblock}
\includegraphics[page=389,pagebox=mediabox,trim={45.00bp 330.95bp 44.37bp 251.91bp},clip,width=0.92\textwidth,max height=0.38\textheight]{Ishii.pdf}
\end{sourceblock}
8 3 0 11 1。 N μ ψ ψ ≥ + 其中 N μ 是由方程（12-34）给出的粘度数。

上述标准适用于（ ）

对于搅拌-湍流流态，漂移速度可由下式给出

\begin{sourceblock}
\includegraphics[page=389,pagebox=mediabox,trim={45.00bp 184.06bp 44.37bp 390.18bp},clip,width=0.92\textwidth,max height=0.38\textheight]{Ishii.pdf}
\end{sourceblock}
在 dj V 的精确表达式中，比例常数 2 适用于气泡流，1.57 适用于液滴流。然而，考虑到预测阻力系数的不确定性，这种差异以及空隙率的影响可以忽略不计。

对于段塞流流态，漂移速度可由下式给出

\begin{sourceblock}
\includegraphics[page=389,pagebox=mediabox,trim={45.00bp 53.74bp 44.37bp 567.72bp},clip,width=0.92\textwidth,max height=0.38\textheight]{Ishii.pdf}
\end{sourceblock}
\begin{sourceblock}
\includegraphics[page=390,pagebox=mediabox,trim={50.00bp 386.10bp 44.37bp 55.00bp},clip,width=\textwidth,max height=0.72\textheight]{Ishii.pdf}
\par\vspace{0.45\baselineskip}
\small 图13-1。高雷诺固体颗粒系统实验数据比较
\end{sourceblock}
编号（石井，1977）

其中 D 是管的直径。

\begin{center}\small 图 13-1 将当前的分析结果 Eq.(13-37) 与固体颗粒流系统的经验相关性进行了比较（Richardson 和 Zaki，1954）。在连续相相对较高的空隙率下的一致性非常好。在 d α 值非常高时，方程（13-37）预测的漂移速度比 Richardson-Zaki 相关性低得多。然而，Richardson和Zaki的原始实验数据也表明了这一趋势，这是由式（13-37）预测的。图 13-2 显示了气泡和液滴-液体流态中的相对速度。 Lackme (1973) 的这些数据清楚地表明了气泡流和液滴流之间相对速度的浓度依赖性的差异。模型已正确预测了这些特征。图13-3和图13-4分别对间歇式和逆流气泡流以及液-液分散系统中的理论预测和实验数据进行了进一步比较。理论预测与数据非常吻合。图 13-5 比较了搅拌湍流的预测和 Yoshida 和 Akita (1965) 的实验数据。数据取自具有 7.7 至 60 cm 各种柱直径的空气-水亚硫酸钠溶液系统。从图中可以看出，该理论低估了较小柱径实验的气体通量。然而，对于较大的柱直径，预测和数据之间的一致性变得越来越令人满意。这个\end{center}
\begin{sourceblock}
\includegraphics[page=391,pagebox=mediabox,trim={50.00bp 372.78bp 44.37bp 55.00bp},clip,width=\textwidth,max height=0.72\textheight]{Ishii.pdf}
\par\vspace{0.45\baselineskip}
\small 图13-2。扭曲中气泡系统与液滴分散系统的区别
\end{sourceblock}
\begin{sourceblock}
\includegraphics[page=391,pagebox=mediabox,trim={50.00bp 106.56bp 44.37bp 307.36bp},clip,width=\textwidth,max height=0.72\textheight]{Ishii.pdf}
\par\vspace{0.45\baselineskip}
\small 图13-3。气体体积通量与扭曲气泡状态数据的比较
\end{sourceblock}
流动系统（Ishii 和 Chawla，1979）

\begin{sourceblock}
\includegraphics[page=392,pagebox=mediabox,trim={50.00bp 400.68bp 44.37bp 55.00bp},clip,width=\textwidth,max height=0.72\textheight]{Ishii.pdf}
\par\vspace{0.45\baselineskip}
\small 图13-4。预测的扭曲粒子区域漂移速度与数据的比较
\end{sourceblock}
\begin{sourceblock}
\includegraphics[page=392,pagebox=mediabox,trim={50.00bp 132.66bp 44.37bp 279.46bp},clip,width=\textwidth,max height=0.72\textheight]{Ishii.pdf}
\par\vspace{0.45\baselineskip}
\small 图13-5。基于搅动流态漂移的预测气体体积通量比较
\end{sourceblock}
数据速度（Ishii 和 Chawla，1979）

两相之间的热力学状态正如我们在上一节中讨论的那样，有必要指定两相之间的机械、热和化学状态。简化的法向动量跃变条件给出了两相之间的力学关系，因此我们有

\begin{sourceblock}
\includegraphics[page=393,pagebox=mediabox,trim={45.00bp 388.03bp 44.37bp 253.26bp},clip,width=0.92\textwidth,max height=0.38\textheight]{Ishii.pdf}
\end{sourceblock}
正如我们从上式中看到的，只有当平均曲率21 H 较大时，即对于具有小流体颗粒直径的气泡或液滴流，压力差才可能是重要的。因此，在许多可以应用漂移通量模型的实际工程问题中，两相之间的压力差可以忽略不计。然后我们有

\begin{sourceblock}
\includegraphics[page=393,pagebox=mediabox,trim={45.00bp 274.99bp 44.37bp 366.30bp},clip,width=0.92\textwidth,max height=0.38\textheight]{Ishii.pdf}
\end{sourceblock}
相间的化学状态决定了相变的条件，由式(9-163)给出：

\begin{sourceblock}
\includegraphics[page=393,pagebox=mediabox,trim={45.00bp 186.69bp 44.37bp 434.94bp},clip,width=0.92\textwidth,max height=0.38\textheight]{Ishii.pdf}
\end{sourceblock}
然而，这可以近似为

\begin{sourceblock}
\includegraphics[page=393,pagebox=mediabox,trim={45.00bp 122.17bp 44.37bp 515.03bp},clip,width=0.92\textwidth,max height=0.38\textheight]{Ishii.pdf}
\end{sourceblock}
其中 g p 表示气相压力。如果可以应用机械平衡条件（13-42），那么我们有

\begin{sourceblock}
\includegraphics[page=394,pagebox=mediabox,trim={50.00bp 421.50bp 44.37bp 55.00bp},clip,width=\textwidth,max height=0.72\textheight]{Ishii.pdf}
\par\vspace{0.45\baselineskip}
\small 图 13-6。轴向温度分布（Ishii，1975）
\end{sourceblock}
\begin{sourceblock}
\includegraphics[page=394,pagebox=mediabox,trim={45.00bp 343.98bp 44.37bp 293.17bp},clip,width=0.92\textwidth,max height=0.38\textheight]{Ishii.pdf}
\end{sourceblock}
两相之间热状态的本构方程是难以获得的本构方程之一，因为它指定了热不平衡的程度。首先我们注意到它可以用 1 i Τ Τ − 或 2 i Τ Τ − 的关系代替。如果两个关系都给出，那么质量传递 2 Γ 的本构方程就变得多余了。认识到这一点非常重要：在许多实际问题中，某一相与界面近似处于热平衡状态，因此

\begin{sourceblock}
\includegraphics[page=394,pagebox=mediabox,trim={45.00bp 204.76bp 44.37bp 436.44bp},clip,width=0.92\textwidth,max height=0.38\textheight]{Ishii.pdf}
\end{sourceblock}
例如，在沸腾系统中，我们可以假设

\begin{sourceblock}
\includegraphics[page=394,pagebox=mediabox,trim={45.00bp 117.23bp 44.37bp 499.29bp},clip,width=0.92\textwidth,max height=0.38\textheight]{Ishii.pdf}
\end{sourceblock}
如图13-6所示。

力学本构方程：T 和 T T
各相的平均粘性应力已在第 9 章 1.2 节中获得，因此我们有

\begin{sourceblock}
\includegraphics[page=395,pagebox=mediabox,trim={45.00bp 530.19bp 44.37bp 98.98bp},clip,width=0.92\textwidth,max height=0.38\textheight]{Ishii.pdf}
\end{sourceblock}
其中右侧第一项具有基于变形张量的粘性应力的标准形式。第二项考虑了界面的影响，界面额外变形张量ki D 由式(9-73)定义。通过代入漂移速度的定义，我们得到

\begin{sourceblock}
\includegraphics[page=395,pagebox=mediabox,trim={45.00bp 367.89bp 44.37bp 219.24bp},clip,width=0.92\textwidth,max height=0.38\textheight]{Ishii.pdf}
\end{sourceblock}
这表明如果相对速度和界面的影响

\begin{sourceblock}
\includegraphics[page=395,pagebox=mediabox,trim={45.00bp 310.83bp 44.37bp 321.29bp},clip,width=0.92\textwidth,max height=0.38\textheight]{Ishii.pdf}
\end{sourceblock}
对于分散两相流，如果传质效应不显著，则界面额外变形应力张量可以采用简单的形式，如式（9-76）。在这种情况下，方程（13-48）可以简化为以下形式

\begin{sourceblock}
\includegraphics[page=395,pagebox=mediabox,trim={45.00bp 189.06bp 44.37bp 411.40bp},clip,width=0.92\textwidth,max height=0.38\textheight]{Ishii.pdf}
\end{sourceblock}
其中 j 和 dj V 分别是混合物的体积通量和分散相的漂移速度。此外，我们注意到 j 可以通过方程（4-93）与 m v 和 dj V 相关。上述结果很重要，因为它表明在分散流中，混合变形张量应该基于体积中心的速度而不是质量中心的速度。混合物粘度由下式给出

\begin{sourceblock}
\includegraphics[page=395,pagebox=mediabox,trim={45.00bp 68.92bp 44.37bp 572.31bp},clip,width=0.92\textwidth,max height=0.38\textheight]{Ishii.pdf}
\end{sourceblock}
现在让我们考虑混合动量方程 (13-3) 中出现的混合湍流应力 T T。从定义中不难看出
T T 是每个相 T k k α T 的湍流应力之和。在第9章1.2节中，我们将混合长度理论应用到双流体模型公式中，从而得到了T k T 的本构方程。然而，方程（9-89）给出的这些表达式并不适合漂移通量模型公式，因为它们是根据每个相的变量而不是混合物的变量编写的。

获得混合物湍流通量 T T 可以考虑两种方法，即利用混合物速度和漂移速度的定义从各相的湍流通量推导 T T，或者类比单相流根据混合物性质建立混合长度模型。如果两相强烈耦合或者涡流的尺寸与分散相的特征尺寸相比很大，则后一种方法是合理的。然后我们有

\begin{sourceblock}
\includegraphics[page=396,pagebox=mediabox,trim={45.00bp 325.00bp 44.37bp 316.06bp},clip,width=0.92\textwidth,max height=0.38\textheight]{Ishii.pdf}
\end{sourceblock}
where

\begin{sourceblock}
\includegraphics[page=396,pagebox=mediabox,trim={45.00bp 249.61bp 44.37bp 377.04bp},clip,width=0.92\textwidth,max height=0.38\textheight]{Ishii.pdf}
\end{sourceblock}
无量纲系数(*T m μ)对应于混合长度常数。

对于分散流，我们可以获得 T T 的不同表达式。如果分散相的尺寸与湍流涡流的尺寸相当，则混合物应力的主要部分由连续相的尺寸给出。因此，我们有

\begin{sourceblock}
\includegraphics[page=396,pagebox=mediabox,trim={45.00bp 128.98bp 44.37bp 512.08bp},clip,width=0.92\textwidth,max height=0.38\textheight]{Ishii.pdf}
\end{sourceblock}
其中我们使用了方程（9-89）。连续相c D 的总变形张量由式(9-82)、(9-74)和(9-76)给出。因此，我们有

\begin{sourceblock}
\includegraphics[page=397,pagebox=mediabox,trim={45.00bp 565.84bp 44.37bp 55.77bp},clip,width=0.92\textwidth,max height=0.38\textheight]{Ishii.pdf}
\end{sourceblock}
这里我们使用了定义（4-91）和恒等式（4-95）。对于许多实际情况，分散相的漂移速度可以仅视为 d α 的函数，正如我们结合运动学本构方程所讨论的那样。如果我们假设漂移速度恒定，则可以进一步简化，然后，通过结合方程（13-50）和（13-54），总混合物应力可以近似为

\begin{sourceblock}
\includegraphics[page=397,pagebox=mediabox,trim={45.00bp 361.19bp 44.37bp 200.08bp},clip,width=0.92\textwidth,max height=0.38\textheight]{Ishii.pdf}
\end{sourceblock}
对于完全开发的管流，我们有

\begin{sourceblock}
\includegraphics[page=397,pagebox=mediabox,trim={45.00bp 256.51bp 44.37bp 337.60bp},clip,width=0.92\textwidth,max height=0.38\textheight]{Ishii.pdf}
\end{sourceblock}
分散流的混合动量源 m M 可由下式给出

\begin{sourceblock}
\includegraphics[page=397,pagebox=mediabox,trim={45.00bp 182.26bp 44.37bp 454.57bp},clip,width=0.92\textwidth,max height=0.38\textheight]{Ishii.pdf}
\end{sourceblock}
对于过渡流程，我们可以采取

\begin{sourceblock}
\includegraphics[page=397,pagebox=mediabox,trim={45.00bp 123.00bp 44.37bp 516.12bp},clip,width=0.92\textwidth,max height=0.38\textheight]{Ishii.pdf}
\end{sourceblock}
能量本构方程 各相的平均传导热通量已在第 9 章 1.2 节中得到，因此我们有

\begin{sourceblock}
\includegraphics[page=398,pagebox=mediabox,trim={45.00bp 562.90bp 44.37bp 55.77bp},clip,width=0.92\textwidth,max height=0.38\textheight]{Ishii.pdf}
\end{sourceblock}
其中右侧第一项具有温度梯度引起的传导传热的标准意义。第二项考虑了接口的影响。通过重新排列各项，上式可简化为

\begin{sourceblock}
\includegraphics[page=398,pagebox=mediabox,trim={45.00bp 451.41bp 44.37bp 171.48bp},clip,width=0.92\textwidth,max height=0.38\textheight]{Ishii.pdf}
\end{sourceblock}
平均传导热通量的这种形式表明这个概念

\begin{sourceblock}
\includegraphics[page=398,pagebox=mediabox,trim={45.00bp 375.03bp 44.37bp 251.20bp},clip,width=0.92\textwidth,max height=0.38\textheight]{Ishii.pdf}
\end{sourceblock}
代表热非平衡效应。

对于许多实际系统，压降对热力学性质的影响可以忽略不计，正如我们已经讨论过的，分散相的温度可以用界面温度来近似。因此通过采取

\begin{sourceblock}
\includegraphics[page=398,pagebox=mediabox,trim={45.00bp 253.96bp 44.37bp 362.50bp},clip,width=0.92\textwidth,max height=0.38\textheight]{Ishii.pdf}
\end{sourceblock}
热通量的本构方程变为

\begin{sourceblock}
\includegraphics[page=398,pagebox=mediabox,trim={45.00bp 184.57bp 44.37bp 443.65bp},clip,width=0.92\textwidth,max height=0.38\textheight]{Ishii.pdf}
\end{sourceblock}
此外，湍流热通量可以与湍流应力张量并行开发。然而，这里应该特别小心，因为混合物温度没有明确定义。根据各相湍流能量传递本构方程（9-92），我们有

\begin{sourceblock}
\includegraphics[page=398,pagebox=mediabox,trim={45.00bp 55.29bp 44.37bp 562.23bp},clip,width=0.92\textwidth,max height=0.38\textheight]{Ishii.pdf}
\end{sourceblock}
\begin{sourceblock}
\includegraphics[page=399,pagebox=mediabox,trim={45.00bp 521.41bp 44.37bp 106.89bp},clip,width=0.92\textwidth,max height=0.38\textheight]{Ishii.pdf}
\end{sourceblock}
根据式(9-94)可得

\begin{sourceblock}
\includegraphics[page=399,pagebox=mediabox,trim={45.00bp 449.56bp 44.37bp 174.49bp},clip,width=0.92\textwidth,max height=0.38\textheight]{Ishii.pdf}
\end{sourceblock}
因此我们得到

\begin{sourceblock}
\includegraphics[page=399,pagebox=mediabox,trim={45.00bp 376.45bp 44.37bp 251.85bp},clip,width=0.92\textwidth,max height=0.38\textheight]{Ishii.pdf}
\end{sourceblock}
这里无量纲系数*T c K 对应于热混合长度常数。我们预计它取决于电导率、表面积浓度、平均曲率和空隙率，如方程（9-95）所示。

\begin{sourceblock}
\includegraphics[page=399,pagebox=mediabox,trim={45.00bp 310.53bp 44.37bp 330.33bp},clip,width=0.92\textwidth,max height=0.38\textheight]{Ishii.pdf}
\end{sourceblock}
σ Φ 和 i m Φ，我们仅注意到，如果必须在分析中包含由于粘性耗散、表面张力和机械能相互作用而产生的这些影响，则应通过三个本构方程来指定它们。从方程（13-7）、（13-8）和（13-9）这些术语的定义可以明显看出，这种本构关系预计会相当复杂。这意味着如果不能忽略这些影响，使用热能方程的大部分优点就会减弱。

相变本构方程 界面传质本构方程由式（9-111）给出。此外，我们注意到，对于漂移通量模型，有必要提供有关两相之间热状态的信息。其简化形式可以由方程（13-47）给出，这对于大多数实际问题都是有用的。因此，对于分散流态，我们有

\begin{sourceblock}
\includegraphics[page=399,pagebox=mediabox,trim={45.00bp 73.48bp 44.37bp 543.50bp},clip,width=0.92\textwidth,max height=0.38\textheight]{Ishii.pdf}
\end{sourceblock}
\begin{sourceblock}
\includegraphics[page=400,pagebox=mediabox,trim={45.00bp 523.62bp 44.37bp 94.05bp},clip,width=0.92\textwidth,max height=0.38\textheight]{Ishii.pdf}
\end{sourceblock}
几何本构方程 如果平均曲率 21 H 和表面积浓度的影响必须明确地包含在公式中，我们应该给出两个附加的几何本构方程。一般来说，这些由方程（9-137）和（9-138）给出，但是对于分散两相流，它们可以简化为方程（9-213）和（9-215）。

\section*{13.3 漂移-通量（或混合）模型公式}\phantomsection\addcontentsline{toc}{section}{13.3 漂移-通量（或混合）模型公式}
漂移通量模型公式的场和本构方程的一般情况已在第 13 章的 1.1 和 1.2 节中讨论。我们已经注意到公式中混合物质心速度以及漂移速度的重要性。将我们当前的模型称为漂移通量模型可能是合适的，以强调漂移速度 kj V 考虑了两相之间相对运动的影响。

\subsection*{13.3.1 漂移-通量模型}\phantomsection\addcontentsline{toc}{subsection}{13.3.1 漂移-通量模型}
鉴于式(4-91)和式(4-89)分别给出了ki V和km V的定义，漂移通量模型的场方程可以给出如下：

由式(13-1)得到的混合连续性方程

\begin{sourceblock}
\includegraphics[page=400,pagebox=mediabox,trim={45.00bp 146.19bp 44.37bp 475.66bp},clip,width=0.92\textwidth,max height=0.38\textheight]{Ishii.pdf}
\end{sourceblock}
方程（13-2）中第2相的连续性方程

\begin{sourceblock}
\includegraphics[page=400,pagebox=mediabox,trim={45.00bp 67.93bp 44.37bp 548.30bp},clip,width=0.92\textwidth,max height=0.38\textheight]{Ishii.pdf}
\end{sourceblock}
由式(13-3)得到的混合动量方程

\begin{sourceblock}
\includegraphics[page=401,pagebox=mediabox,trim={45.00bp 504.28bp 44.37bp 80.25bp},clip,width=0.92\textwidth,max height=0.38\textheight]{Ishii.pdf}
\end{sourceblock}
由式(13-10)得到的混合热能方程

\begin{sourceblock}
\includegraphics[page=401,pagebox=mediabox,trim={45.00bp 341.43bp 44.37bp 206.38bp},clip,width=0.92\textwidth,max height=0.38\textheight]{Ishii.pdf}
\end{sourceblock}
在这里，我们根据混合物特性、空隙率 2 α 和漂移速度 2 j V 制定了模型。该模型对于分散两相流最有效，因为对于这种情况，本构方程可以简化为实际形式，如第 13 章 1.2 节中讨论的那样。

需要再次强调的是，漂移-通量模型对于两相流系统分析是有用的。如果两相的运动强耦合，则尤其如此。由于其简单性，该模型可用于进行现实的相似性分析以及解决许多重要的工程问题。

\subsection*{13.3.2 缩放参数}\phantomsection\addcontentsline{toc}{subsection}{13.3.2 缩放参数}
让我们用下标 o 来表示参考参数。特征长度尺度为o L，时间尺度取为o L 与速度尺度的比值。然后我们定义数量级为1的无量纲参数如下：

\begin{sourceblock}
\includegraphics[page=401,pagebox=mediabox,trim={45.00bp 70.93bp 44.37bp 548.74bp},clip,width=0.92\textwidth,max height=0.38\textheight]{Ishii.pdf}
\end{sourceblock}
\begin{sourceblock}
\includegraphics[page=402,pagebox=mediabox,trim={45.00bp 413.61bp 44.37bp 57.04bp},clip,width=0.92\textwidth,max height=0.38\textheight]{Ishii.pdf}
\end{sourceblock}
将这些新参数代入场方程中，我们得到以下结果：

无量纲混合连续性方程

\begin{sourceblock}
\includegraphics[page=402,pagebox=mediabox,trim={45.00bp 299.42bp 44.37bp 325.54bp},clip,width=0.92\textwidth,max height=0.38\textheight]{Ishii.pdf}
\end{sourceblock}
第 2 相的无量纲连续性方程

\begin{sourceblock}
\includegraphics[page=402,pagebox=mediabox,trim={45.00bp 223.74bp 44.37bp 397.42bp},clip,width=0.92\textwidth,max height=0.38\textheight]{Ishii.pdf}
\end{sourceblock}
混合物的无量纲动量方程

\begin{sourceblock}
\includegraphics[page=402,pagebox=mediabox,trim={45.00bp 73.14bp 44.37bp 476.38bp},clip,width=0.92\textwidth,max height=0.38\textheight]{Ishii.pdf}
\end{sourceblock}
混合物的无量纲热能方程

\begin{sourceblock}
\includegraphics[page=403,pagebox=mediabox,trim={45.00bp 446.75bp 44.37bp 55.42bp},clip,width=0.92\textwidth,max height=0.38\textheight]{Ishii.pdf}
\end{sourceblock}
这里我们定义了

\begin{sourceblock}
\includegraphics[page=403,pagebox=mediabox,trim={45.00bp 121.12bp 44.37bp 248.26bp},clip,width=0.92\textwidth,max height=0.38\textheight]{Ishii.pdf}
\end{sourceblock}
我们注意到，这八组是基于漂移通量模型公式的混合物的缩放参数。这些组与 Ishii (1971) 中的一维两相流模型获得的组类似。然而，雷诺数和佩克莱特数是一个例外，因为
应明确区分缩放参数和相似性组，因为它们不是同义词。相似群可以从场方程和边界以及指定所有本构方程的初始条件获得。因此，除非明确定义系统和问题，否则无法详细讨论相似性标准。缩放参数使我们能够做出各种假设和近似，因为它们给出了场方程中出现的各种项的数量级。

式(13-79)的前三组、pch N 、DN 、N ρ 为运动学群。如果相位变化数pch N 远大于漂移数D N，则系统由相位变化控制。然而，如果 D pch N N >>，则系统由相位的重新分配控制。 Re N 、Fr N 和 N σ 表示的数是动态群，因为它们缩放了混合动量方程中出现的各种力，而佩克莱特数和埃克特数是能量群。我们注意到埃克特数 Ec N 的特别重要性。显然，除了高速流动外，Ec N 都很小，因此 m 给出的项

\begin{sourceblock}
\includegraphics[page=404,pagebox=mediabox,trim={45.00bp 340.53bp 44.37bp 292.86bp},clip,width=0.92\textwidth,max height=0.38\textheight]{Ishii.pdf}
\end{sourceblock}
从式（13-76）、（13-77）和（13-78）中，我们设置质量、动量和能量的漂移传输不由相同的缩放参数加权。由于相间相对运动而产生的相漂移和焓传输具有大致相同的量级 D N，因此这两项应在相同的条件下处理。然而，动量漂移由 2 D N N ρ 加权。因此，根据该缩放参数的大小，可以独立于其他两个效应而忽略漂移应力张量。

我们注意到，在缩放参数中，我们引入了由 m μ 和 m K 给出的混合物粘度和电导率。只要涉及缩放参数，这些参数就不需要精确定义，因为只有它们的数量级才是重要的。我们可以选择k μ 和k K 中较大的粘度和电导率，或者定义为2

\begin{sourceblock}
\includegraphics[page=404,pagebox=mediabox,trim={45.00bp 117.80bp 44.37bp 499.34bp},clip,width=0.92\textwidth,max height=0.38\textheight]{Ishii.pdf}
\end{sourceblock}
系统边界处的示例。

然而，如果涉及相似组，则应使用应力张量和热通量的精确本构方程。因此，应根据本构方程的形式和变量，选择正确的参考参数。基于

\subsection*{13.3.3 均质流模型}\phantomsection\addcontentsline{toc}{subsection}{13.3.3 均质流模型}
如果漂移数 D N 远小于相变数
pch N，则系统是反应（相变）控制的，并且鉴于第 2 相的连续性方程，质量的漂移或扩散可以忽略不计。此外，如果漂移数 D N 远小于 1，则场方程中的所有漂移项和界面机械能传递效应 i m Φ 都可以忽略不计。然而，我们可能不会从公式中删除第 2 相的连续性方程，因为它考虑了 Zuber 和 Dougherty (1967) 结合一维模型讨论的热非平衡效应。因此，对于一般均匀流模型，我们有以下四个场方程：

\begin{sourceblock}
\includegraphics[page=405,pagebox=mediabox,trim={45.00bp 217.61bp 44.37bp 262.66bp},clip,width=0.92\textwidth,max height=0.38\textheight]{Ishii.pdf}
\end{sourceblock}
在这种情况下，分散两相流的力学本构方程可以简化为简单形式。因此，根据式（13-50），我们有

\begin{sourceblock}
\includegraphics[page=405,pagebox=mediabox,trim={45.00bp 126.67bp 44.37bp 506.14bp},clip,width=0.92\textwidth,max height=0.38\textheight]{Ishii.pdf}
\end{sourceblock}
湍流应力由式（13-52）给出：

\begin{sourceblock}
\includegraphics[page=405,pagebox=mediabox,trim={45.00bp 70.90bp 44.37bp 570.07bp},clip,width=0.92\textwidth,max height=0.38\textheight]{Ishii.pdf}
\end{sourceblock}
此外，热通量的本构方程由下式给出

\begin{sourceblock}
\includegraphics[page=406,pagebox=mediabox,trim={45.00bp 547.45bp 44.37bp 80.85bp},clip,width=0.92\textwidth,max height=0.38\textheight]{Ishii.pdf}
\end{sourceblock}
and

\begin{sourceblock}
\includegraphics[page=406,pagebox=mediabox,trim={45.00bp 480.37bp 44.37bp 147.87bp},clip,width=0.92\textwidth,max height=0.38\textheight]{Ishii.pdf}
\end{sourceblock}
其中，假设分散相与界面处于热平衡，即式（13-62）。如果埃克特数 Ec N 和表面数 N σ 较小，则混合动量方程中的毛细管力 m M，压缩性效应 m Dp Dt，耗散项 m
能量方程中的σ Φ 可以去掉。

μ Φ，以及表面张力效应 m

\subsection*{13.3.4 密度传播模型}\phantomsection\addcontentsline{toc}{subsection}{13.3.4 密度传播模型}
对于许多具有实际意义的系统，特别是对于没有大声相互作用的高减压系统，可以合理地假设每个相本质上是不可压缩的。然后我们有

\begin{sourceblock}
\includegraphics[page=406,pagebox=mediabox,trim={45.00bp 281.85bp 44.37bp 359.28bp},clip,width=0.92\textwidth,max height=0.38\textheight]{Ishii.pdf}
\end{sourceblock}
此外，我们考虑第 2 相的漂移速度近似仅是 2 α 的函数的情况。我们记得，对于许多实际流动状态来说，这是一个有效的假设，因为它已在第 7 章的 1.3 节和第 13 章的 1.2 节中讨论过。因此，我们有

\begin{sourceblock}
\includegraphics[page=406,pagebox=mediabox,trim={45.00bp 179.56bp 44.37bp 459.39bp},clip,width=0.92\textwidth,max height=0.38\textheight]{Ishii.pdf}
\end{sourceblock}
为了简单起见，我们还假设表面数 N σ 和埃克特数 Ec N 都很小，使得

\begin{sourceblock}
\includegraphics[page=406,pagebox=mediabox,trim={45.00bp 110.60bp 44.37bp 532.34bp},clip,width=0.92\textwidth,max height=0.38\textheight]{Ishii.pdf}
\end{sourceblock}
则由式（7-31）、（7-38）、（13-72）、（13-73）得到如下场方程：

\begin{sourceblock}
\includegraphics[page=407,pagebox=mediabox,trim={45.00bp 402.90bp 44.37bp 55.66bp},clip,width=0.92\textwidth,max height=0.38\textheight]{Ishii.pdf}
\end{sourceblock}
其中 K C 和 j 由下式给出

\begin{sourceblock}
\includegraphics[page=407,pagebox=mediabox,trim={45.00bp 328.87bp 44.37bp 295.54bp},clip,width=0.92\textwidth,max height=0.38\textheight]{Ishii.pdf}
\end{sourceblock}
and

\begin{sourceblock}
\includegraphics[page=407,pagebox=mediabox,trim={45.00bp 248.80bp 44.37bp 373.37bp},clip,width=0.92\textwidth,max height=0.38\textheight]{Ishii.pdf}
\end{sourceblock}
目前的密度波传播模型最适合分散（或混合）两相流状态。在这种情况下，相2应视为分散相。那么应力的本构方程可由式（13-50）和（13-54）给出，而热通量由式（13-63）和（13-67）给出。

\chapter{一维漂移通量模型}
两相流总是涉及一相相对于另一相的一些相对运动；因此，两相流问题应该用两个速度场来表示。一般的瞬态两相流问题可以通过使用二流体模型或漂移通量模型来表达，具体取决于相之间的动态耦合程度。在双流体模型中，每个相都被单独考虑；因此，该模型是根据两组守恒方程来表述的，这些方程控制着每个相的质量、动量和能量的平衡。然而，在公式中引入两个动量方程（如双流体模型的情况）会由于数学复杂性和指定两相之间的界面相互作用项的不确定性而带来相当大的困难（Delhaye，1968；Vernier 和 Delhaye，1968；Bouré 和 Réocruex，1972；Ishii，1975）。由于相动量方程中界面相互作用项选择不当而导致的数值不稳定是很常见的。因此，在制定二流体模型时需要仔细研究界面本构方程。例如，有人建议（Réocruex，1974）交互作用项应包括某些条件下的一阶时间和空间导数。

通过根据漂移通量模型制定两相问题，可以显著减少与双流体模型相关的这些困难（Zuber，1967）。在该模型中，整个混合物的运动由混合物动量方程表示，相间的相对运动由运动学本构方程考虑。因此，漂移通量模型的基本概念是将混合物视为一个整体而不是两个分离的相。基于混合物平衡方程的漂移通量模型的公式比基于每个相单独的平衡方程的双流体模型更简单。与漂移通量模型相关的最重要的假设是两相的动力学可以通过混合动量方程来表达，其中
在漂移通量模型中，速度场用混合物质心速度和气相漂移速度（即相对于混合物体积中心的蒸气速度）来表示。热非平衡的影响通过相变本构方程容纳在漂移通量模型中，该方程指定了每单位体积的传质速率。由于界面处的质量和动量传递速率取决于界面的结构，因此漂移速度和蒸汽生成的这些本构方程是流态的函数（Zuber 和 Dougherty，1967；Ishii 等，1975）。

与更严格的双流体公式相比，漂移通量模型是一个近似公式。然而，由于其简单性和适用于实际感兴趣的各种两相流问题，漂移通量模型具有相当重要的意义。特别是，通过对横截面积上的局部漂移通量公式进行平均而获得的一维漂移通量模型对于涉及流体流动和传热的复杂工程问题非常有用，因为场方程可以简化为准一维形式。通过面积平均，通道内垂直于主流方向的变量变化信息基本上丢失。因此，壁与流体之间动量和能量的传递应通过经验相关性或简化模型来表达。在本章中，我们开发了漂移通量模型的一般一维公式，并讨论了实际应用中重要的各种特殊情况。为了简单起见，本章的公式中省略了时间平均的数学符号。 Ishii (1977)对该模型进行了广泛的回顾。

\section*{14.1 三维漂移通量模型面积平均}\phantomsection\addcontentsline{toc}{section}{14.1 三维漂移通量模型面积平均}
采用时间平均法或统计平均法得到了漂移-通量模型的三维形式。结果可总结如下：

由式(13-70)得到的混合连续性方程

\begin{sourceblock}
\includegraphics[page=409,pagebox=mediabox,trim={45.00bp 74.47bp 44.37bp 550.78bp},clip,width=0.92\textwidth,max height=0.38\textheight]{Ishii.pdf}
\end{sourceblock}
由式（13-71）得出的分散相连续性方程

\begin{sourceblock}
\includegraphics[page=410,pagebox=mediabox,trim={45.00bp 540.64bp 44.37bp 80.96bp},clip,width=0.92\textwidth,max height=0.38\textheight]{Ishii.pdf}
\end{sourceblock}
由式(13-72)得到的混合动量方程

\begin{sourceblock}
\includegraphics[page=410,pagebox=mediabox,trim={45.00bp 431.71bp 44.37bp 157.84bp},clip,width=0.92\textwidth,max height=0.38\textheight]{Ishii.pdf}
\end{sourceblock}
由式(13-73)得到混合物的焓-能方程

\begin{sourceblock}
\includegraphics[page=410,pagebox=mediabox,trim={45.00bp 281.60bp 44.37bp 266.02bp},clip,width=0.92\textwidth,max height=0.38\textheight]{Ishii.pdf}
\end{sourceblock}
获得一维漂移通量模型的合理方法是在横截面上积分三维漂移通量模型，然后引入适当的平均值。横截面积上的简单面积平均值 A 定义为

\begin{sourceblock}
\includegraphics[page=410,pagebox=mediabox,trim={45.00bp 174.34bp 44.37bp 457.34bp},clip,width=0.92\textwidth,max height=0.38\textheight]{Ishii.pdf}
\end{sourceblock}
空隙率加权平均值由下式给出

\begin{sourceblock}
\includegraphics[page=410,pagebox=mediabox,trim={45.00bp 100.54bp 44.37bp 522.88bp},clip,width=0.92\textwidth,max height=0.38\textheight]{Ishii.pdf}
\end{sourceblock}
在后续分析中，认为任意横截面积内各相密度 d ρ 和 c ρ 是均匀的，因此 k k ρ ρ =。对于大多数实际的两相流问题，该假设是有效的，因为

\begin{sourceblock}
\includegraphics[page=411,pagebox=mediabox,trim={45.00bp 519.91bp 44.37bp 121.29bp},clip,width=0.92\textwidth,max height=0.38\textheight]{Ishii.pdf}
\end{sourceblock}
k 相加权平均速度的轴向分量为

\begin{sourceblock}
\includegraphics[page=411,pagebox=mediabox,trim={45.00bp 443.51bp 44.37bp 179.92bp},clip,width=0.92\textwidth,max height=0.38\textheight]{Ishii.pdf}
\end{sourceblock}
其中速度的标量表达式对应于矢量的轴向分量。然后混合速度定义为

\begin{sourceblock}
\includegraphics[page=411,pagebox=mediabox,trim={45.00bp 352.55bp 44.37bp 269.92bp},clip,width=0.92\textwidth,max height=0.38\textheight]{Ishii.pdf}
\end{sourceblock}
体积通量由下式给出

\begin{sourceblock}
\includegraphics[page=411,pagebox=mediabox,trim={45.00bp 293.81bp 44.37bp 347.28bp},clip,width=0.92\textwidth,max height=0.38\textheight]{Ishii.pdf}
\end{sourceblock}
平均混合物焓也应按密度加权；因此，

\begin{sourceblock}
\includegraphics[page=411,pagebox=mediabox,trim={45.00bp 215.56bp 44.37bp 406.90bp},clip,width=0.92\textwidth,max height=0.38\textheight]{Ishii.pdf}
\end{sourceblock}
适当的平均漂移速度定义为

\begin{sourceblock}
\includegraphics[page=411,pagebox=mediabox,trim={45.00bp 154.86bp 44.37bp 485.31bp},clip,width=0.92\textwidth,max height=0.38\textheight]{Ishii.pdf}
\end{sourceblock}
如果测量各相的体积流量 k Q 和平均空隙率 d α，则可以通过实验确定漂移速度。这是因为式(14-12)可以转化为

\begin{sourceblock}
\includegraphics[page=411,pagebox=mediabox,trim={45.00bp 52.54bp 44.37bp 574.03bp},clip,width=0.92\textwidth,max height=0.38\textheight]{Ishii.pdf}
\end{sourceblock}
\begin{sourceblock}
\includegraphics[page=412,pagebox=mediabox,trim={45.00bp 464.73bp 44.37bp 117.87bp},clip,width=0.92\textwidth,max height=0.38\textheight]{Ishii.pdf}
\end{sourceblock}
and

\begin{sourceblock}
\includegraphics[page=412,pagebox=mediabox,trim={45.00bp 391.53bp 44.37bp 231.75bp},clip,width=0.92\textwidth,max height=0.38\textheight]{Ishii.pdf}
\end{sourceblock}
在漂移通量公式中，通过给定 dj V 的本构关系解决 d α 和 m v 的问题。因此，问题解决后，可以利用式(14-14)恢复各相速度的解。

通过对方程（14-1）-（14-4）进行面积平均并使用各种平均值，我们得到
混合连续性方程

\begin{sourceblock}
\includegraphics[page=412,pagebox=mediabox,trim={45.00bp 238.52bp 44.37bp 385.83bp},clip,width=0.92\textwidth,max height=0.38\textheight]{Ishii.pdf}
\end{sourceblock}
分散相的连续性方程

\begin{sourceblock}
\includegraphics[page=412,pagebox=mediabox,trim={45.00bp 162.55bp 44.37bp 457.03bp},clip,width=0.92\textwidth,max height=0.38\textheight]{Ishii.pdf}
\end{sourceblock}
混合动量方程

\begin{sourceblock}
\includegraphics[page=412,pagebox=mediabox,trim={45.00bp 56.16bp 44.37bp 536.62bp},clip,width=0.92\textwidth,max height=0.38\textheight]{Ishii.pdf}
\end{sourceblock}
\begin{sourceblock}
\includegraphics[page=413,pagebox=mediabox,trim={45.00bp 526.53bp 44.37bp 55.48bp},clip,width=0.92\textwidth,max height=0.38\textheight]{Ishii.pdf}
\end{sourceblock}
混合物的焓能方程

\begin{sourceblock}
\includegraphics[page=413,pagebox=mediabox,trim={45.00bp 335.38bp 44.37bp 171.76bp},clip,width=0.92\textwidth,max height=0.38\textheight]{Ishii.pdf}
\end{sourceblock}
式中，T zz zz τ τ + 为应力张量轴向法向分量，dc h Δ 为相间的焓差；因此，
dc d c h h h Δ = −。协方差项表示乘积平均值与两个变量平均值乘积之间的差异，使得 ( ) ( ) k k k k k k k k k COV v v v α ρ ψ α ρ ψ ≡ −。如果 k ψ 或 k v 的轮廓是平坦的，则协方差项减小到零。式(14-18)中( ) 2 mm m m f v v D ρ 表示的项为两相摩擦压降。我们在这里注意到，与相间相对运动相关的质量、动量和能量扩散的影响明确地出现在漂移通量公式中，因为场方程左侧的对流项是用混合速度表示的。本配方中的这些扩散效应以分散相的漂移速度 dj V 来表达。这可以用函数形式表示为

\begin{sourceblock}
\includegraphics[page=413,pagebox=mediabox,trim={45.00bp 111.01bp 44.37bp 530.16bp},clip,width=0.92\textwidth,max height=0.38\textheight]{Ishii.pdf}
\end{sourceblock}
为了考虑跨界面的质量传递，还应该给出 d Γ 的本构方程。在函数形式中，该相变本构方程可以写为

\begin{sourceblock}
\includegraphics[page=414,pagebox=mediabox,trim={45.00bp 565.90bp 44.37bp 55.74bp},clip,width=0.92\textwidth,max height=0.38\textheight]{Ishii.pdf}
\end{sourceblock}
如果给出了其中一相的漂移速度的适当本构关系，则上述公式可以扩展到非分散两相流，例如环形流。

\section*{14.2 一维漂移速度}\phantomsection\addcontentsline{toc}{section}{14.2 一维漂移速度}
\subsection*{14.2.1 分散两相流}\phantomsection\addcontentsline{toc}{subsection}{14.2.1 分散两相流}
为了获得一维漂移通量模型的运动学本构方程，我们必须对通道横截面上的局部漂移速度进行平均。第 13 章 1.2 节阐述了受限通道中局部漂移速度 dj V 的本构关系。现在我们将其与方程（14-12）定义的平均漂移速度 dj V 联系起来。由式(14-6)和式(14-12)可得

\begin{sourceblock}
\includegraphics[page=414,pagebox=mediabox,trim={45.00bp 315.76bp 44.37bp 306.45bp},clip,width=0.92\textwidth,max height=0.38\textheight]{Ishii.pdf}
\end{sourceblock}
where

\begin{sourceblock}
\includegraphics[page=414,pagebox=mediabox,trim={45.00bp 235.72bp 44.37bp 386.37bp},clip,width=0.92\textwidth,max height=0.38\textheight]{Ishii.pdf}
\end{sourceblock}
and

\begin{sourceblock}
\includegraphics[page=414,pagebox=mediabox,trim={45.00bp 159.70bp 44.37bp 463.71bp},clip,width=0.92\textwidth,max height=0.38\textheight]{Ishii.pdf}
\end{sourceblock}
方程（14-22）右侧第二项是浓度分布和体积通量分布之间的协方差；因此也可以表示为 ( ) d d COV j α α。因子 0 C 已被多位作者（Zuber 和 Findlay，1965；Nicklin 等人，1962；Neal，1963）用于气泡流或段塞流，被称为分布参数。 Bankoff (1960) 的早期工作中也使用了该参数的逆。从物理上讲，这种效应是由于分散相局部存在的事实而产生的。

0 C 的值可以根据空隙率 d α 和总体积通量 j 的假设分布确定（Zuber 和 Findlay，1965），或根据实验数据（Zuber 等人，1967）。通过假设管道中 j 和 d α 的幂律分布，我们有

\begin{sourceblock}
\includegraphics[page=415,pagebox=mediabox,trim={45.00bp 405.84bp 44.37bp 146.58bp},clip,width=0.92\textwidth,max height=0.38\textheight]{Ishii.pdf}
\end{sourceblock}
其中 0 j 、0 d α 、dW α 、r 和 W R 分别是中心处 j 和 α 的值、壁部空隙率、径向距离和管道半径。将这些分布代入式（14-24）给出的 0 C 的定义，我们得到

\begin{sourceblock}
\includegraphics[page=415,pagebox=mediabox,trim={45.00bp 291.55bp 44.37bp 328.06bp},clip,width=0.92\textwidth,max height=0.38\textheight]{Ishii.pdf}
\end{sourceblock}
基于上述假设轮廓的分布参数在参考文献中进一步讨论（Zuber 等人，1967）。

此时式(14-22)可改写为

\begin{sourceblock}
\includegraphics[page=415,pagebox=mediabox,trim={45.00bp 189.52bp 44.37bp 437.04bp},clip,width=0.92\textwidth,max height=0.38\textheight]{Ishii.pdf}
\end{sourceblock}
其中 d v 和 j 是实验中容易获得的参数，特别是在绝热条件下。因此，该方程给出了平均速度 d v 与平均体积通量 j 的关系图。如果整个通道的浓度分布是均匀的，则分布参数的值等于 1。此外，如果局部漂移 dj V 的影响小到可以忽略不计，则流动基本上变得均匀。在这种情况下，平均速度和通量之间的关系简化为以 45° 角通过原点的直线。实验数据与该均匀流线的偏差表明
该图的一个重要特征是，对于具有充分发展的空隙和速度剖面的两相流状态，数据点聚集在一条直线周围（见图14-1-14-3）。当局部漂移速度恒定或小到可以忽略不计时，这种趋势尤其明显。因此，对于给定的流态，分布参数 0 C 的值可以从这些线的斜率获得，而该线与平均速度轴的截距可以解释为加权平均局部漂移速度 dj V。 Zuber 等人的广泛研究。 (1967) 表明 0 C 取决于压力、通道几何形状，或许还取决于流速。 Hancox 和 Nicoll (1972) 还指出了过冷沸腾和形成的空隙分布对分布参数的重要影响。这里，基于 Ishii (1977) 的研究，提出了气泡流状态中分布参数的简单相关性。首先，通过考虑完全发展的气泡流，我们假设 0 C 取决于密度比 g f ρ ρ 和基于液体特性的雷诺数 f GD μ，其中 G 、D 和 f μ 分别是液体的总质量流量、水力直径和粘度。因此，

\begin{sourceblock}
\includegraphics[page=416,pagebox=mediabox,trim={50.00bp 75.30bp 44.37bp 309.51bp},clip,width=\textwidth,max height=0.72\textheight]{Ishii.pdf}
\par\vspace{0.45\baselineskip}
\small 图14-1。完全开发的空气-水流数据（Ishii，1977）
\end{sourceblock}
\begin{sourceblock}
\includegraphics[page=417,pagebox=mediabox,trim={50.00bp 376.32bp 44.37bp 55.00bp},clip,width=\textwidth,max height=0.72\textheight]{Ishii.pdf}
\par\vspace{0.45\baselineskip}
\small 图14-2。汽水并流上流、并流下流实验数据
\end{sourceblock}
\begin{sourceblock}
\includegraphics[page=417,pagebox=mediabox,trim={50.00bp 111.00bp 44.37bp 303.82bp},clip,width=\textwidth,max height=0.72\textheight]{Ishii.pdf}
\par\vspace{0.45\baselineskip}
\small 图14-3。加热型并流上流和并流下流实验数据
\end{sourceblock}
Santowax-R 系统（Ishii，1977）

单相湍流分布和最大速度与平均速度之比给出了在 0 d α → 和 0 g f ρ ρ → 时理论极限值为 0 C，因为在这种情况下所有气泡都应集中在中心区域。因此，根据 Nukuradse (1932) 的圆管实验数据，给出了最大速度与平均速度之比，我们有

\begin{sourceblock}
\includegraphics[page=418,pagebox=mediabox,trim={45.00bp 473.22bp 44.37bp 145.72bp},clip,width=0.92\textwidth,max height=0.38\textheight]{Ishii.pdf}
\end{sourceblock}
为 0 d α → 和 0 g f ρ ρ →。此外，当密度比接近1时，分布参数0 C 也应变为1。因此，

\begin{sourceblock}
\includegraphics[page=418,pagebox=mediabox,trim={45.00bp 404.05bp 44.37bp 237.30bp},clip,width=0.92\textwidth,max height=0.38\textheight]{Ishii.pdf}
\end{sourceblock}
为 1 g f ρ ρ →。基于这些限制和各种实验数据，全面
发展的流动，分布参数可以近似地给出

\begin{sourceblock}
\includegraphics[page=418,pagebox=mediabox,trim={45.00bp 326.52bp 44.37bp 314.33bp},clip,width=0.92\textwidth,max height=0.38\textheight]{Ishii.pdf}
\end{sourceblock}
其中密度组缩放横向空隙分布中每个相的惯性效应。从物理上讲，方程（14-31）模拟了较轻相迁移到高速区域的趋势，从而导致中心区域出现更高的空隙浓度（Bankoff，1960）。对于层流，C ∞ 为 2，但由于速度梯度较大，0 C 在低空隙率下对 d α 非常敏感。在较宽的雷诺数范围内，f GD μ，方程（14-29）可以近似为 1 2 C。 ∞ ≅ 对于圆管中的流动。此外，对于矩形通道，实验数据显示该值约为1.35。因此，对于充分发展的湍流气泡流，

\begin{sourceblock}
\includegraphics[page=418,pagebox=mediabox,trim={45.00bp 120.93bp 44.37bp 492.25bp},clip,width=0.92\textwidth,max height=0.38\textheight]{Ishii.pdf}
\par\vspace{0.45\baselineskip}
\small 图14-4和图14-5将上述相关性与各种实验数据进行了比较。图中的每个点代表 5 到 150 个数据点。例如，图14-1所示的Smissaert(1963)的原始实验数据在图14-4中以单图呈现。每个
\end{sourceblock}
\begin{sourceblock}
\includegraphics[page=419,pagebox=mediabox,trim={50.00bp 399.06bp 44.37bp 55.00bp},clip,width=\textwidth,max height=0.72\textheight]{Ishii.pdf}
\par\vspace{0.45\baselineskip}
\small 图14-4。圆管中完全发展的流动的分布参数（Ishii，1977）
\end{sourceblock}
\begin{center}\small 图 14-1 中的点可以通过使用平均局部漂移速度 dj V 的现有相关性来获得 0 C 的相应值。但鉴于分散相平均速度与总通量之间存在较强的线性关系，图14-4和图14-5中均采用线性拟合得到的0℃平均值。\end{center}
在速度-通量平面（参见图14-2和14-3）中，可以容易地识别三种运行模式。在第一象限，流动基本上是并流向上；因此，液相和气相都向上流动。在第二象限，气相向上移动；然而，混合物存在净向下流动。 Consequently, the flow is countercurrent. The cocurrent downflow operation should appear in the third quadrant of the velocity-flux plane, as shown in Figs.14-2 and 14-3.这些数据表明，方程（14-27）描述的基本特征对于具有相同分布参数值 0 C 的并流向上和向下流动均有效。这一事实证明了将漂移速度与平均局部漂移速度 dj V 和 0 C 相关联的有用性。在加热的两相系统中，可能会发生空隙轮廓从凹形到凸形的变化。凹形空隙率分布是由壁成核和气泡向通道中心的延迟横向迁移引起的。在这些条件下，大多数气泡最初位于成核壁附近，尽管即使在绝热流动中，小气泡也倾向于在低空隙率的壁区域附近积聚。 The concave profile is particularly pronounced in the subcooled boiling regime,

\begin{sourceblock}
\includegraphics[page=420,pagebox=mediabox,trim={50.00bp 399.06bp 44.37bp 55.00bp},clip,width=\textwidth,max height=0.72\textheight]{Ishii.pdf}
\par\vspace{0.45\baselineskip}
\small 图14-5。矩形通道中完全发展的流动的分布参数（Ishii，
\end{sourceblock}
1977)

因为这里只有壁边界层被加热到饱和温度以上，而核心液体被过冷。这种温度分布将引起核心区域中迁移气泡的破裂，并由此产生从壁到过冷液体的潜热传输。通过多孔管壁将气体注入流动液体中也可以获得类似的凹形轮廓（Rose 和 Griffith，1965）。

在空隙仍然集中在靠近壁的区域中，蒸汽的平均速度可能小于液体的平均速度，因为大部分液体以高核心速度移动。然而，随着越来越多的蒸汽沿着通道产生，空隙率轮廓从凹面变为凸面并变得完全发展。对于由于成核或气体注入而在壁处产生空隙的流动，分布参数 0 C 在两相流区域的开始处应具有接近于零的值。这从式(14-24)中0 C 的定义也可以看出。因此，我们有

\begin{sourceblock}
\includegraphics[page=420,pagebox=mediabox,trim={45.00bp 103.90bp 44.37bp 519.52bp},clip,width=0.92\textwidth,max height=0.38\textheight]{Ishii.pdf}
\end{sourceblock}
随着横截面平均空隙率的增加，局部空隙率的峰值从近壁区域向中心区域移动。随着空隙轮廓的发展，这将导致 0 C 值增加。

\begin{sourceblock}
\includegraphics[page=421,pagebox=mediabox,trim={50.00bp 401.40bp 44.37bp 55.00bp},clip,width=\textwidth,max height=0.72\textheight]{Ishii.pdf}
\par\vspace{0.45\baselineskip}
\small 图 14-6。由于沸腾而产生流动的分布参数（矩形管道的数据已修改1.2/1.35倍，以获得相应的数据
\end{sourceblock}
圆管。）（石井，1977）

鉴于上述基本特征和各种实验数据（Zuber et al., 1967; Maurer, 1956; Pierre, 1965; Marchaterre, 1956），提出了以下简单的相关性（Ishii, 1977）

\begin{sourceblock}
\includegraphics[page=421,pagebox=mediabox,trim={45.00bp 257.46bp 44.37bp 371.56bp},clip,width=0.92\textwidth,max height=0.38\textheight]{Ishii.pdf}
\end{sourceblock}
该表达式表明在 0 0 25 d 给出的区域中形成的空隙轮廓的重要性。 α < < ;超出该区域，0℃的值迅速接近完全发展的流动（见图14-6）。因此，对于 0 g Γ >，我们得到

\begin{sourceblock}
\includegraphics[page=421,pagebox=mediabox,trim={45.00bp 100.91bp 44.37bp 479.30bp},clip,width=0.92\textwidth,max height=0.38\textheight]{Ishii.pdf}
\end{sourceblock}
对于湍流状态下的大多数液滴或颗粒流，由于湍流混合和壁附近的颗粒滑移，体积通量分布非常平坦，这增加了体积通量。浓度

\begin{sourceblock}
\includegraphics[page=422,pagebox=mediabox,trim={45.00bp 504.94bp 44.37bp 136.30bp},clip,width=0.92\textwidth,max height=0.38\textheight]{Ishii.pdf}
\end{sourceblock}
在这种情况下，局部滑移就变得很重要。

根据局部本构方程计算dj V，即积分变换，式(14-23)；因此，它将需要有关空洞轮廓的附加信息（Ishii，1976）。由于该轮廓通常不为人所知，因此我们进行以下简化近似。假设用平均值代替局部空隙率 d α 和应力梯度的无量纲差异，则由局部滑移引起的平均漂移速度 dj V 可以通过与参考文献（Ishii，1976）中给出的局部本构关系相同的表达式来预测。这些近似值适用于具有相对平坦的空隙率轮廓的流动；此外，从一维模型的整体简单性来看，它们可以被认为是可以接受的。

对于充分发展的垂直流，流体中和分散相中的应力分布应该相似；因此，剪切梯度对平均局部漂移速度的影响可以忽略不计。在这些条件下，我们得到以下结果。

未扭曲粒子状态

\begin{sourceblock}
\includegraphics[page=422,pagebox=mediabox,trim={45.00bp 111.21bp 44.37bp 415.90bp},clip,width=0.92\textwidth,max height=0.38\textheight]{Ishii.pdf}
\end{sourceblock}
为 34 65 * d r。 ≥ .未扭曲粒子状态的极限情况是斯托克斯状态，其中平均漂移速度降低为

\begin{sourceblock}
\includegraphics[page=423,pagebox=mediabox,trim={45.00bp 568.85bp 44.37bp 55.51bp},clip,width=0.92\textwidth,max height=0.38\textheight]{Ishii.pdf}
\end{sourceblock}
\begin{sourceblock}
\includegraphics[page=423,pagebox=mediabox,trim={45.00bp 539.66bp 44.37bp 74.43bp},clip,width=0.92\textwidth,max height=0.38\textheight]{Ishii.pdf}
\end{sourceblock}
\begin{sourceblock}
\includegraphics[page=423,pagebox=mediabox,trim={45.00bp 489.71bp 44.37bp 104.52bp},clip,width=0.92\textwidth,max height=0.38\textheight]{Ishii.pdf}
\end{sourceblock}
这里 n 的值取决于粘度：

\begin{sourceblock}
\includegraphics[page=423,pagebox=mediabox,trim={45.00bp 392.73bp 44.37bp 207.65bp},clip,width=0.92\textwidth,max height=0.38\textheight]{Ishii.pdf}
\end{sourceblock}
搅动湍流流态

\begin{sourceblock}
\includegraphics[page=423,pagebox=mediabox,trim={45.00bp 314.74bp 44.37bp 306.76bp},clip,width=0.92\textwidth,max height=0.38\textheight]{Ishii.pdf}
\end{sourceblock}
这里平均混合物粘度（Ishii，1976）由下式给出

\begin{sourceblock}
\includegraphics[page=423,pagebox=mediabox,trim={45.00bp 234.76bp 44.37bp 384.81bp},clip,width=0.92\textwidth,max height=0.38\textheight]{Ishii.pdf}
\end{sourceblock}
最大包装值，0 62 dm。 α =，建议用于固体颗粒-流体系统，尽管它的范围可以从 0.5 到 0.74。然而，对于气泡流，dm α 的理论值可能要高得多。如果我们考虑气泡流中空隙率的标准范围，则 dm α 可以近似为 1 dm α =。因此，对于气泡流，混合物粘度变为

\begin{sourceblock}
\includegraphics[page=423,pagebox=mediabox,trim={45.00bp 93.70bp 44.37bp 529.72bp},clip,width=0.92\textwidth,max height=0.38\textheight]{Ishii.pdf}
\end{sourceblock}
然而，对于低颗粒浓度的颗粒流，即

\footnotetext[1]{d α << , m μ 可以近似为}
\begin{sourceblock}
\includegraphics[page=424,pagebox=mediabox,trim={45.00bp 566.80bp 44.37bp 56.62bp},clip,width=0.92\textwidth,max height=0.38\textheight]{Ishii.pdf}
\end{sourceblock}
在分散相完全悬浮的水平流中，保持颗粒悬浮的横向混合可以显著影响各相的应力梯度；因此，应力梯度效应不可忽视。然而，鉴于目前的技术水平，假设 0 dj V ≅ 可以用作一阶近似，
特别是在高通量流中。正如本章末尾所解释的，实际的局部漂移速度还取决于摩擦引起的压力梯度，因此，严格意义上，即使在水平流中，实际的局部漂移速度也不为零。

对于高通量流动，与协方差项 ( ) 0 1 C j − 相比，局部漂移 dj V 对平均漂移速度的影响较小。因此，通过忽略前者，我们有

\begin{sourceblock}
\includegraphics[page=424,pagebox=mediabox,trim={45.00bp 355.72bp 44.37bp 267.58bp},clip,width=0.92\textwidth,max height=0.38\textheight]{Ishii.pdf}
\end{sourceblock}
对于气泡流，上述方程对适用的空隙率范围施加了一个条件；因此，我们应该有 ( ) ( ) 0 1 m d d c C ρ α ρ ρ > − −。这里，通过计算总体积通量与终端速度之比，可以获得高通量流和低通量流之间边界的简单标准。如果该比值大于10，则可以认为该流是高通量流。

受限通道中分散两相流的另一个极限情况是段塞流。当气泡的体积非常大时，气泡的形状会显著变形以适应通道的几何形状。气泡的直径接近管道的直径，薄膜将气泡与管壁隔开。气泡呈子弹状，带有帽状鼻子。这些气泡在相对无粘性流体中的运动可以通过使用球体周围的势流分析来研究（Dumitrescu，1943），并且结果与实验数据一致。因此，

\begin{sourceblock}
\includegraphics[page=424,pagebox=mediabox,trim={45.00bp 94.69bp 44.37bp 526.84bp},clip,width=0.92\textwidth,max height=0.38\textheight]{Ishii.pdf}
\end{sourceblock}
这最初是由 Nicklin 等人提出的。 （1962）和尼尔（1963）。

\subsection*{14.2.2 环形两相流}\phantomsection\addcontentsline{toc}{subsection}{14.2.2 环形两相流}
在环形两相流中，相之间的相对运动由界面几何形状、体力场和界面动量传递控制。通过考虑结构化两相流的宏观效应，建立了环形两相流中蒸汽漂移速度的本构方程（Ishii 等，1976）。假设稳态绝热两相环形流具有恒定的单相特性，我们对每个相有以下一维动量方程。

\begin{sourceblock}
\includegraphics[page=425,pagebox=mediabox,trim={45.00bp 426.22bp 44.37bp 197.56bp},clip,width=0.92\textwidth,max height=0.38\textheight]{Ishii.pdf}
\end{sourceblock}
and

\begin{sourceblock}
\includegraphics[page=425,pagebox=mediabox,trim={45.00bp 347.83bp 44.37bp 272.50bp},clip,width=0.92\textwidth,max height=0.38\textheight]{Ishii.pdf}
\end{sourceblock}
其中 i τ 、 wf τ 、 i P 和 wf P 分别是界面剪切、壁剪切、界面润湿周长和壁润湿周长。水力直径和湿周长比由 4 wf D A P ≡ 和

\begin{sourceblock}
\includegraphics[page=425,pagebox=mediabox,trim={73.59bp 249.79bp 78.50bp 353.05bp},clip,width=\textwidth,max height=0.72\textheight]{Ishii.pdf}
\end{sourceblock}
重力修正项 2 3 wf wf f f f z f v v g τ ρ Δρ δ = −，其中
wf f 可以通过标准摩擦系数关系式给出： 16 wf f f Re = 对于层流膜流和 0 25 0 0791。 wf f f。 Re − = 对于湍流。这里液膜雷诺数由 f f f f Re j D ρ μ = 给出。同样，
界面剪切可表示为 2 i i gr r r f v v τ ρ = 其中
界面摩擦系数由 ( ) 0 005 1 75 1 i g f 给出。 α [ ] = + −     对于粗糙的波浪薄膜（Wallis，1969）。

根据定义，蒸气漂移速度与 r v 有关，即 ( ) 1 gj gr V v α = −。因此，通过消除动量方程中的压力梯度，我们得到层状薄膜

\begin{sourceblock}
\includegraphics[page=426,pagebox=mediabox,trim={45.00bp 546.61bp 44.37bp 57.58bp},clip,width=0.92\textwidth,max height=0.38\textheight]{Ishii.pdf}
\end{sourceblock}
对于一部动荡的电影

\begin{sourceblock}
\includegraphics[page=426,pagebox=mediabox,trim={45.00bp 403.17bp 44.37bp 148.15bp},clip,width=0.92\textwidth,max height=0.38\textheight]{Ishii.pdf}
\end{sourceblock}
这里，当绝对符号内的项变为负数时，取负根。式(14-49)和式(14-50)表示的漂移速度便于分析稳态绝热或热平衡流动，因为在这些情况下，可以容易地获得f j 的值。

在一般的漂移通量模型公式中，gj V 应以混合物速度 m v 而不是 f j 来表示，因为 m v 是公式中使用的速度。根据定义，我们有

\begin{sourceblock}
\includegraphics[page=426,pagebox=mediabox,trim={45.00bp 237.82bp 44.37bp 384.28bp},clip,width=0.92\textwidth,max height=0.38\textheight]{Ishii.pdf}
\end{sourceblock}
将式（14-51）代入式（14-49），可得层状薄膜

\begin{sourceblock}
\includegraphics[page=426,pagebox=mediabox,trim={45.00bp 82.08bp 44.37bp 460.96bp},clip,width=0.92\textwidth,max height=0.38\textheight]{Ishii.pdf}
\end{sourceblock}
这对于由下式给出的层流范围有效

\begin{sourceblock}
\includegraphics[page=427,pagebox=mediabox,trim={45.00bp 518.74bp 44.37bp 58.51bp},clip,width=0.92\textwidth,max height=0.38\textheight]{Ishii.pdf}
\end{sourceblock}
\begin{sourceblock}
\includegraphics[page=427,pagebox=mediabox,trim={93.45bp 471.31bp 46.04bp 154.33bp},clip,width=\textwidth,max height=0.72\textheight]{Ishii.pdf}
\end{sourceblock}
( ) ( ) 1 g m m g g v α ρ α ρ ≤ −，流动是并流向上的，而对于
gj V 大于上述限值时，液体向下流动。解决方案
对于湍流薄膜流的情况则稍微复杂一些。为了方便起见，我们介绍一下以下参数。

\begin{sourceblock}
\includegraphics[page=427,pagebox=mediabox,trim={45.00bp 253.88bp 44.37bp 270.92bp},clip,width=0.92\textwidth,max height=0.38\textheight]{Ishii.pdf}
\end{sourceblock}
那么，对于向上的液体流动，我们有

\begin{sourceblock}
\includegraphics[page=427,pagebox=mediabox,trim={45.00bp 132.52bp 44.37bp 441.43bp},clip,width=0.92\textwidth,max height=0.38\textheight]{Ishii.pdf}
\end{sourceblock}
\begin{sourceblock}
\includegraphics[page=427,pagebox=mediabox,trim={139.83bp 82.09bp 172.40bp 535.21bp},clip,width=\textwidth,max height=0.72\textheight]{Ishii.pdf}
\end{sourceblock}
适用于 2 m v cb a ≥ 的条件。然而，在转型过程中
由 2 m c v cb a − ≤ < 给出的状态，其中液膜流量为
向下，薄膜上的界面剪切力向上，蒸气漂移速度变为

\begin{sourceblock}
\includegraphics[page=428,pagebox=mediabox,trim={45.00bp 562.46bp 44.37bp 57.39bp},clip,width=0.92\textwidth,max height=0.38\textheight]{Ishii.pdf}
\end{sourceblock}
\begin{sourceblock}
\includegraphics[page=428,pagebox=mediabox,trim={45.00bp 482.71bp 44.37bp 109.38bp},clip,width=0.92\textwidth,max height=0.38\textheight]{Ishii.pdf}
\end{sourceblock}
适用于并流向下流动。

仅当满足以下湍流准则时，才能应用上述解决方案。

\begin{sourceblock}
\includegraphics[page=428,pagebox=mediabox,trim={45.00bp 305.22bp 44.37bp 238.24bp},clip,width=0.92\textwidth,max height=0.38\textheight]{Ishii.pdf}
\end{sourceblock}
对于一部湍流电影来说，这些结果并不具有非常简单的形式。然而，如果混合物速度的绝对值较大，使得流动基本上是并流且重力效应较小，则湍流解可以用简单的形式来近似

\begin{sourceblock}
\includegraphics[page=428,pagebox=mediabox,trim={45.00bp 154.74bp 44.37bp 431.49bp},clip,width=0.92\textwidth,max height=0.38\textheight]{Ishii.pdf}
\end{sourceblock}
在平均液体速度远小于蒸汽速度的简化假设下，对漂移速度方程（14-59）进行变换即可得到滑移率 g f v v。然后我们有

\begin{sourceblock}
\includegraphics[page=429,pagebox=mediabox,trim={45.00bp 548.64bp 44.37bp 57.64bp},clip,width=0.92\textwidth,max height=0.38\textheight]{Ishii.pdf}
\end{sourceblock}
对于管道中的环形流动，其中 g xi α =。上述滑移率的表达式与 Fauske (1962) 得到的相似，即
g f f g v v ρ ρ =，与空隙率无关。

对于 0 8 1 g 的范围，方程（14-60）中考虑空隙率的因子大致在 0.24 到 1 之间变化。 α < <。因此，对于湍流薄膜，Fauske 相关性应在高空隙率下给出相当准确的结果。

环形流的漂移速度相关性已用混合物速度表示，因为 m v 是一般漂移通量模型公式中的基本变量。然而，根据总体积通量 j 解析 gj V 的表达式也是有趣且重要的，因为 j 是用于关联分散两相流状态中的 gj V 的变量。

通过考虑湍流薄膜流状态并使用定义

( ) 1 f g g gj j j V α α = − −，我们可以求解式(14-50)得到平均漂移
速度 gj V。结果并不具有简单的形式；然而，对于大多数实际情况，它可以通过 j 的线性函数来近似。

\begin{sourceblock}
\includegraphics[page=429,pagebox=mediabox,trim={45.00bp 213.17bp 44.37bp 374.38bp},clip,width=0.92\textwidth,max height=0.38\textheight]{Ishii.pdf}
\end{sourceblock}
\begin{sourceblock}
\includegraphics[page=429,pagebox=mediabox,trim={45.00bp 50.14bp 44.37bp 468.77bp},clip,width=0.92\textwidth,max height=0.38\textheight]{Ishii.pdf}
\end{sourceblock}
\begin{sourceblock}
\includegraphics[page=429,pagebox=mediabox,trim={74.91bp 143.23bp 229.58bp 470.89bp},clip,width=\textwidth,max height=0.72\textheight]{Ishii.pdf}
\end{sourceblock}
\begin{sourceblock}
\includegraphics[page=430,pagebox=mediabox,trim={45.00bp 526.23bp 44.37bp 95.56bp},clip,width=0.92\textwidth,max height=0.38\textheight]{Ishii.pdf}
\end{sourceblock}
这表明环形流中的表观 0 C 应接近一致。

\subsection*{14.2.3 环形雾流}\phantomsection\addcontentsline{toc}{subsection}{14.2.3 环形雾流}
随着环形流中气体速度的增加，液体从薄膜夹带到气芯流中。根据为夹带开始而制定的标准（Ishii 和 Grolmes，1975），粗糙湍流膜流的临界气体速度可由下式给出：

\begin{sourceblock}
\includegraphics[page=430,pagebox=mediabox,trim={45.00bp 296.50bp 44.37bp 267.57bp},clip,width=0.92\textwidth,max height=0.38\textheight]{Ishii.pdf}
\end{sourceblock}
远大于环形雾流状态下的液体通量。那么，对于水或钠等弱粘性流体，上述关系式可以替换为

\begin{sourceblock}
\includegraphics[page=430,pagebox=mediabox,trim={45.00bp 207.25bp 44.37bp 411.70bp},clip,width=0.92\textwidth,max height=0.38\textheight]{Ishii.pdf}
\end{sourceblock}
如果满足不等式（14-65），则应考虑液滴夹带进入气芯流；否则可应用环流相关式（14-62）。

通过结合先前的分散流和纯环形流的结果，可以很容易地得出环形雾流中 gj V 的相关性。在任何横截面处，气芯中夹带的液体占总液体面积的面积分数用 d E 表示，并且横截面积平均空隙率用 g α 表示。然后薄膜面积分数由下式给出

\begin{sourceblock}
\includegraphics[page=431,pagebox=mediabox,trim={45.00bp 550.03bp 44.37bp 57.81bp},clip,width=0.92\textwidth,max height=0.38\textheight]{Ishii.pdf}
\end{sourceblock}
单独气芯中的平均液滴分数由下式给出

\begin{sourceblock}
\includegraphics[page=431,pagebox=mediabox,trim={45.00bp 435.43bp 44.37bp 151.11bp},clip,width=0.92\textwidth,max height=0.38\textheight]{Ishii.pdf}
\end{sourceblock}
因此，应在环流相关式（式（14-62））中使用核心α，以获得核心与薄膜之间的相对运动，而
液滴α应该用于分散流关联中以获得液滴和气体核心流之间的滑移。

通过分别用 gc v 、 fc v 和 ff v 表示气芯速度、液滴速度和薄膜速度，总体积通量由下式给出

\begin{sourceblock}
\includegraphics[page=431,pagebox=mediabox,trim={45.00bp 303.91bp 44.37bp 313.50bp},clip,width=0.92\textwidth,max height=0.38\textheight]{Ishii.pdf}
\end{sourceblock}
此外，通过用核心 j 表示基于核心面积的核心中的总体积通量，我们可以从环形相关性方程（14-62）中得到：

\begin{sourceblock}
\includegraphics[page=431,pagebox=mediabox,trim={45.00bp 213.46bp 44.37bp 389.46bp},clip,width=0.92\textwidth,max height=0.38\textheight]{Ishii.pdf}
\end{sourceblock}
\begin{sourceblock}
\includegraphics[page=431,pagebox=mediabox,trim={74.85bp 119.23bp 188.78bp 494.89bp},clip,width=\textwidth,max height=0.72\textheight]{Ishii.pdf}
\end{sourceblock}
从分散流相关性可以看出，对于扭曲液滴或搅拌液滴流态，漂移速度可以近似由下式给出

\begin{sourceblock}
\includegraphics[page=432,pagebox=mediabox,trim={45.00bp 560.12bp 44.37bp 57.46bp},clip,width=0.92\textwidth,max height=0.38\textheight]{Ishii.pdf}
\end{sourceblock}
这里我们使用了基于 ( ) 1 1 g α − << 的近似。然而，根据核心气体速度，应该使用更小的颗粒的分散流漂移速度相关性。当液滴由液膜夹带产生时，对于斯托克斯体系之外的未扭曲粒子体系，建议采用以下近似形式（Ishii，1976）。

\begin{sourceblock}
\includegraphics[page=432,pagebox=mediabox,trim={45.00bp 410.50bp 44.37bp 205.48bp},clip,width=0.92\textwidth,max height=0.38\textheight]{Ishii.pdf}
\end{sourceblock}
其中粒子半径可以根据剪切波峰处的韦伯数准则来近似。因此，

\begin{sourceblock}
\includegraphics[page=432,pagebox=mediabox,trim={45.00bp 325.72bp 44.37bp 298.51bp},clip,width=0.92\textwidth,max height=0.38\textheight]{Ishii.pdf}
\end{sourceblock}
仅当总体积通量足够高以引起波峰破碎时，上述关系才适用。因此，当

\begin{sourceblock}
\includegraphics[page=432,pagebox=mediabox,trim={45.00bp 209.20bp 44.37bp 399.31bp},clip,width=0.92\textwidth,max height=0.38\textheight]{Ishii.pdf}
\end{sourceblock}
综合以上结果，我们得到

\begin{sourceblock}
\includegraphics[page=432,pagebox=mediabox,trim={45.00bp 72.93bp 44.37bp 489.11bp},clip,width=0.92\textwidth,max height=0.38\textheight]{Ishii.pdf}
\end{sourceblock}
\begin{sourceblock}
\includegraphics[page=433,pagebox=mediabox,trim={45.00bp 488.81bp 44.37bp 55.88bp},clip,width=0.92\textwidth,max height=0.38\textheight]{Ishii.pdf}
\end{sourceblock}
其中后一个表达式适用于式（14-73）给出的条件。如果颗粒半径很小，则对相间相对运动的本质贡献来自式（14-74）第一项，核心流可被认为是均匀分散流。在这种情况下，方程（14-74）简化为

\begin{sourceblock}
\includegraphics[page=433,pagebox=mediabox,trim={45.00bp 302.93bp 44.37bp 262.19bp},clip,width=0.92\textwidth,max height=0.38\textheight]{Ishii.pdf}
\end{sourceblock}
该表达式显示，就夹带液体分数而言，漂移速度呈线性下降，这可以在各种实验数据中观察到（Alia 等人，1965 年；Cravarolo 等人，1964 年）。

\section*{14.3 对流通量的协方差}\phantomsection\addcontentsline{toc}{section}{14.3 对流通量的协方差}
在一维漂移通量模型中，动量和能量对流通量被分为三项：混合对流通量；漂移对流通量；和协方差项，如式（14-18）和（14-19）所示。换句话说，混合物的对流通量 ψ 可以写为

\begin{sourceblock}
\includegraphics[page=433,pagebox=mediabox,trim={45.00bp 55.11bp 44.37bp 527.25bp},clip,width=0.92\textwidth,max height=0.38\textheight]{Ishii.pdf}
\end{sourceblock}
因此，对于动量通量，我们有 k ψ k v = 和 dc Δψ ( ) 1 dj g V α = −。对于焓通量，我们有 k k h ψ = 和 dc Δψ
d c h h = −，如果各相处于热平衡状态，则相当于潜热。

为了闭合控制方程组，我们必须指定这些协方差项的关系。这可以通过引入动量和能量通量的分布参数来完成。如果我们将通量的分布参数定义为

\begin{sourceblock}
\includegraphics[page=434,pagebox=mediabox,trim={45.00bp 425.56bp 44.37bp 197.93bp},clip,width=0.92\textwidth,max height=0.38\textheight]{Ishii.pdf}
\end{sourceblock}
协方差项变为

\begin{sourceblock}
\includegraphics[page=434,pagebox=mediabox,trim={45.00bp 338.41bp 44.37bp 278.13bp},clip,width=0.92\textwidth,max height=0.38\textheight]{Ishii.pdf}
\end{sourceblock}
对于动量通量，分布参数定义为

\begin{sourceblock}
\includegraphics[page=434,pagebox=mediabox,trim={45.00bp 257.74bp 44.37bp 362.73bp},clip,width=0.92\textwidth,max height=0.38\textheight]{Ishii.pdf}
\end{sourceblock}
从物理上讲，vk C 表示空隙和动量通量分布对 k 相横截面积平均动量通量的影响。 vk C 的定量研究可以通过考虑圆形管道中的对称流动并引入幂律表达式与第 14 章 1.2.1 节中 0 C 的分析并行进行。因此我们假设

\begin{sourceblock}
\includegraphics[page=434,pagebox=mediabox,trim={45.00bp 127.54bp 44.37bp 494.04bp},clip,width=0.92\textwidth,max height=0.38\textheight]{Ishii.pdf}
\end{sourceblock}
and

\begin{sourceblock}
\includegraphics[page=434,pagebox=mediabox,trim={45.00bp 50.56bp 44.37bp 571.01bp},clip,width=0.92\textwidth,max height=0.38\textheight]{Ishii.pdf}
\end{sourceblock}
为简单起见，假设空隙和速度分布相似；即，n m = .该假设广泛应用于质量传递问题。如果考虑到蒸汽通量以及由此产生的空隙浓度极大地影响速度分布，那么对于完全发展的两相流来说，这也可能是合理的。在此假设下，可以证明

\begin{sourceblock}
\includegraphics[page=435,pagebox=mediabox,trim={45.00bp 421.02bp 44.37bp 171.69bp},clip,width=0.92\textwidth,max height=0.38\textheight]{Ishii.pdf}
\end{sourceblock}
\begin{sourceblock}
\includegraphics[page=435,pagebox=mediabox,trim={45.00bp 378.53bp 44.37bp 220.33bp},clip,width=0.92\textwidth,max height=0.38\textheight]{Ishii.pdf}
\end{sourceblock}
\begin{sourceblock}
\includegraphics[page=435,pagebox=mediabox,trim={45.00bp 337.23bp 44.37bp 264.71bp},clip,width=0.92\textwidth,max height=0.38\textheight]{Ishii.pdf}
\end{sourceblock}
然而，根据式（14-26），体积通量分布参数 0 C 变为

\begin{sourceblock}
\includegraphics[page=435,pagebox=mediabox,trim={45.00bp 254.24bp 44.37bp 374.73bp},clip,width=0.92\textwidth,max height=0.38\textheight]{Ishii.pdf}
\end{sourceblock}
因此，在 n 的标准范围内，参数 vg C 可以近似给出：

\begin{sourceblock}
\includegraphics[page=435,pagebox=mediabox,trim={45.00bp 184.01bp 44.37bp 457.30bp},clip,width=0.92\textwidth,max height=0.38\textheight]{Ishii.pdf}
\end{sourceblock}
对于蒸气分散流状态下的液相，1 fw α ≅ 且

\begin{sourceblock}
\includegraphics[page=435,pagebox=mediabox,trim={45.00bp 98.99bp 44.37bp 518.84bp},clip,width=0.92\textwidth,max height=0.38\textheight]{Ishii.pdf}
\end{sourceblock}
对于环形流，动量协方差项也可以通过使用蒸汽和液体流的标准速度分布来计算。这样我们就得到

\begin{sourceblock}
\includegraphics[page=436,pagebox=mediabox,trim={45.00bp 557.07bp 44.37bp 56.02bp},clip,width=0.92\textwidth,max height=0.38\textheight]{Ishii.pdf}
\end{sourceblock}
上述各个相的结果现在可用于研究混合协方差项。通过将混合动量分布参数定义为

\begin{sourceblock}
\includegraphics[page=436,pagebox=mediabox,trim={45.00bp 457.73bp 44.37bp 165.69bp},clip,width=0.92\textwidth,max height=0.38\textheight]{Ishii.pdf}
\end{sourceblock}
协方差项变为

\begin{sourceblock}
\includegraphics[page=436,pagebox=mediabox,trim={45.00bp 304.19bp 44.37bp 243.68bp},clip,width=0.92\textwidth,max height=0.38\textheight]{Ishii.pdf}
\end{sourceblock}
综合以上分析，( ) vd vc C C − 的数量级与( ) 1 vm C − 的数量级相同或更小；因此，对于几乎所有情况，式（14-89）右边的最后一项都可以忽略。该项可能仅在近临界状态且 m dj v V ≅ 时才重要。然而，一般来说，当密度比接近 1 时，dj V 变得微不足道。因此，在上述条件下，对流项本身变得相对较小。因此，即使在这种情况下，该术语也可能被删除。因此我们有

\begin{sourceblock}
\includegraphics[page=436,pagebox=mediabox,trim={45.00bp 95.47bp 44.37bp 487.10bp},clip,width=0.92\textwidth,max height=0.38\textheight]{Ishii.pdf}
\end{sourceblock}
\begin{sourceblock}
\includegraphics[page=437,pagebox=mediabox,trim={45.00bp 584.21bp 44.37bp 57.10bp},clip,width=0.92\textwidth,max height=0.38\textheight]{Ishii.pdf}
\end{sourceblock}
然而，在近临界状态下，vm C 还取决于空隙率和密度比。此外，在完全发展的流动或发展的流动中，在非常低的空隙率下，vm C 的值应减小到方程（14-87）给出的单相流的值。可以通过方程（14-34）中 0 C 相关性中使用的类似空隙率校正项来考虑空隙轮廓发展为幂律给出的效果的影响。回想一下湍流 1 0 vm C。 = 在 0 d α →，我们得到圆管

\begin{sourceblock}
\includegraphics[page=437,pagebox=mediabox,trim={45.00bp 426.62bp 44.37bp 180.45bp},clip,width=0.92\textwidth,max height=0.38\textheight]{Ishii.pdf}
\end{sourceblock}
它既可以用于完全开发的流程，也可以用于开发流程。

对于湍流环流流态，根据方程（14-87）和（14-88），我们有 1 02 vm C。 ≅.出于所有实际目的，这可以进一步近似为

\begin{sourceblock}
\includegraphics[page=437,pagebox=mediabox,trim={45.00bp 317.09bp 44.37bp 324.13bp},clip,width=0.92\textwidth,max height=0.38\textheight]{Ishii.pdf}
\end{sourceblock}
实际上，从方程（14-92）给出的值到方程（14-93）给出的值的过渡是通过搅动-环形（或段塞-环形）流状态逐渐进行的，其中搅动和环形流的特征交替。如果首选 vm C 的单一相关性，则无论流动状态转换如何，则可以通过以下给出的简单修改将方程（14-92）安全地外推到更高空隙率状态：

\begin{sourceblock}
\includegraphics[page=437,pagebox=mediabox,trim={45.00bp 183.62bp 44.37bp 447.85bp},clip,width=0.92\textwidth,max height=0.38\textheight]{Ishii.pdf}
\end{sourceblock}
通过假设空隙、速度和焓分布，可以对焓协方差项进行类似的分析。一般来说，

\begin{sourceblock}
\includegraphics[page=437,pagebox=mediabox,trim={45.00bp 58.84bp 44.37bp 533.64bp},clip,width=0.92\textwidth,max height=0.38\textheight]{Ishii.pdf}
\end{sourceblock}
\begin{sourceblock}
\includegraphics[page=438,pagebox=mediabox,trim={45.00bp 515.44bp 44.37bp 80.83bp},clip,width=0.92\textwidth,max height=0.38\textheight]{Ishii.pdf}
\end{sourceblock}
\begin{sourceblock}
\includegraphics[page=438,pagebox=mediabox,trim={45.00bp 488.50bp 44.37bp 155.88bp},clip,width=0.92\textwidth,max height=0.38\textheight]{Ishii.pdf}
\end{sourceblock}
fs h 是蒸气和液体的饱和焓。由于在这种情况下，焓曲线完全是每个相的，所以分布参数变得统一；即 1 hg hf hm C C C = = =。同样明显的是，如果其中一个相处于饱和状态，则该相的 hk C 变为 1。

在单相区域中，分布参数可以根据速度和焓的假设分布来计算。使用标准

\begin{sourceblock}
\includegraphics[page=438,pagebox=mediabox,trim={80.97bp 354.79bp 69.74bp 260.47bp},clip,width=\textwidth,max height=0.72\textheight]{Ishii.pdf}
\end{sourceblock}
可以表明，对于正在发展的流和完全发展的流来说，协方差项都小到可以忽略不计。

从上述两种限制情况，我们可以得出结论，只有在高度非平衡流动中，焓协变项才可能变得重要。即使在这种情况下，与相变相关的能量也比与横向温度分布变化相关的能量大得多。因此，除了高度瞬态的情况外，可以忽略焓协方差。因此，

\begin{sourceblock}
\includegraphics[page=438,pagebox=mediabox,trim={45.00bp 205.23bp 44.37bp 419.74bp},clip,width=0.92\textwidth,max height=0.38\textheight]{Ishii.pdf}
\end{sourceblock}
\section*{14.4 各种流动条件下的一维漂移-通量相关性}\phantomsection\addcontentsline{toc}{section}{14.4 各种流动条件下的一维漂移-通量相关性}
本节总结了具有实际意义的各种流动条件下的一维漂移通量模型的本构方程。

\subsection*{14.4.1 向上气泡流的本构方程}\phantomsection\addcontentsline{toc}{subsection}{14.4.1 向上气泡流的本构方程}
考虑气泡横向迁移特性，对圆管内向上绝热气泡流分布参数的本构方程（14-32）进行了改进（Hibiki and Ishii，2002b；2003c）

\begin{sourceblock}
\includegraphics[page=439,pagebox=mediabox,trim={45.00bp 419.50bp 44.37bp 148.11bp},clip,width=0.92\textwidth,max height=0.38\textheight]{Ishii.pdf}
\end{sourceblock}
其中 f Re 由 f f j D ν 定义。式(14-98)中的气泡直径Sm D 可以通过气泡直径相关性来预测(Hibiki and Ishii, 2002a)。从式（14-98）可以看出，随着液体雷诺数的增加，式（14-98）预测的分布参数渐近逼近式（14-32）。气液气泡流中漂移速度的本构方程由下式给出

\begin{sourceblock}
\includegraphics[page=439,pagebox=mediabox,trim={45.00bp 271.24bp 44.37bp 347.57bp},clip,width=0.92\textwidth,max height=0.38\textheight]{Ishii.pdf}
\end{sourceblock}
\subsection*{14.4.2 向上绝热环空和内热环空本构方程}\phantomsection\addcontentsline{toc}{subsection}{14.4.2 向上绝热环空和内热环空本构方程}
方程（14-98）和（14-99）的适用性已通过在大气压下垂直同心环空中采集的向上空气-水湍流泡状流数据得到证实（Hibiki et al., 2003b）。进一步的研究表明，C ∞ =1.1 可以更好地预测环空（Ozar 等，2008）。考虑通道几何差异，从式（14-35）推导出内热环空沸腾泡流分布参数的本构方程（Hibiki et al., 2003a；Julia et al., 2009）

\begin{sourceblock}
\includegraphics[page=439,pagebox=mediabox,trim={45.00bp 50.14bp 44.37bp 569.62bp},clip,width=0.92\textwidth,max height=0.38\textheight]{Ishii.pdf}
\end{sourceblock}
\subsection*{14.4.3 向下两相流本构方程}\phantomsection\addcontentsline{toc}{subsection}{14.4.3 向下两相流本构方程}
所有流态的向下两相流分布参数的本构方程由下式给出（Goda et al., 2003）

\begin{sourceblock}
\includegraphics[page=440,pagebox=mediabox,trim={45.00bp 375.61bp 44.37bp 161.50bp},clip,width=0.92\textwidth,max height=0.38\textheight]{Ishii.pdf}
\end{sourceblock}
其中 gj j j V == *。所有流态下向下两相流的漂移速度的本构方程近似为

\begin{sourceblock}
\includegraphics[page=440,pagebox=mediabox,trim={45.00bp 277.84bp 44.37bp 340.97bp},clip,width=0.92\textwidth,max height=0.38\textheight]{Ishii.pdf}
\end{sourceblock}
这些分布参数和漂移速度的本构方程是由一维数据建立的，并且尚未通过详细的局部流动数据单独验证。因此，它们不应该单独使用。

\subsection*{14.4.4 鼓泡或沸腾池系统的本构方程}\phantomsection\addcontentsline{toc}{subsection}{14.4.4 鼓泡或沸腾池系统的本构方程}
在鼓泡或池沸腾系统中，容器直径与长度的比率通常比强制对流系统大。值得注意的是，在低流量的大容器中可能会形成再循环流动模式。这可能会显著影响横向速度和空隙率分布。鼓泡或沸腾池系统的漂移速度 ( 0 f j = ) 的本构方程由下式给出（Kataoka 和 Ishii，1987）

\begin{sourceblock}
\includegraphics[page=441,pagebox=mediabox,trim={45.00bp 588.44bp 44.37bp 53.56bp},clip,width=0.92\textwidth,max height=0.38\textheight]{Ishii.pdf}
\end{sourceblock}
\begin{sourceblock}
\includegraphics[page=441,pagebox=mediabox,trim={45.00bp 492.24bp 44.37bp 55.75bp},clip,width=0.92\textwidth,max height=0.38\textheight]{Ishii.pdf}
\end{sourceblock}
\begin{sourceblock}
\includegraphics[page=441,pagebox=mediabox,trim={45.00bp 464.31bp 44.37bp 149.49bp},clip,width=0.92\textwidth,max height=0.38\textheight]{Ishii.pdf}
\end{sourceblock}
\begin{sourceblock}
\includegraphics[page=441,pagebox=mediabox,trim={45.00bp 410.22bp 44.37bp 183.84bp},clip,width=0.92\textwidth,max height=0.38\textheight]{Ishii.pdf}
\end{sourceblock}
\begin{sourceblock}
\includegraphics[page=441,pagebox=mediabox,trim={45.00bp 379.50bp 44.37bp 231.49bp},clip,width=0.92\textwidth,max height=0.38\textheight]{Ishii.pdf}
\end{sourceblock}
鼓泡或池沸腾系统中分布参数的本构方程根据通道几何形状给出，如方程(14-32)。

\subsection*{14.4.5 大直径管道系统的本构方程}\phantomsection\addcontentsline{toc}{subsection}{14.4.5 大直径管道系统的本构方程}
在大直径通道 ( ( ) 40 H D g σ Δρ ≥ ) 中，由于界面不稳定性，段状气泡无法持续存在，并且它们会破裂以覆盖气泡。再循环流动模式可以在低流速下在大直径通道中形成。测试部分入口处的流态和发展流中的流态转变也可能对液体再循环模式产生影响。液体再循环、入口流态和流态转变可能显著影响横向速度和空隙率分布。大直径管道系统中向上气泡流的漂移速度的本构方程近似为（Hibiki 和 Ishii，2003a）

\begin{sourceblock}
\includegraphics[page=441,pagebox=mediabox,trim={45.00bp 133.46bp 44.37bp 492.21bp},clip,width=0.92\textwidth,max height=0.38\textheight]{Ishii.pdf}
\end{sourceblock}
\begin{sourceblock}
\includegraphics[page=441,pagebox=mediabox,trim={76.23bp 81.07bp 130.52bp 538.69bp},clip,width=\textwidth,max height=0.72\textheight]{Ishii.pdf}
\end{sourceblock}
1 4 2 g g f j j g σ Δρ ρ + ≡ .大直径管道系统中向上气泡流分布参数的本构方程为
入口流态为均匀分布的气泡流的情况

\begin{sourceblock}
\includegraphics[page=442,pagebox=mediabox,trim={45.00bp 444.95bp 44.37bp 81.05bp},clip,width=0.92\textwidth,max height=0.38\textheight]{Ishii.pdf}
\end{sourceblock}
\begin{sourceblock}
\includegraphics[page=442,pagebox=mediabox,trim={45.00bp 412.20bp 44.37bp 198.32bp},clip,width=0.92\textwidth,max height=0.38\textheight]{Ishii.pdf}
\end{sourceblock}
入口流型为盖泡流或段塞流的情况

\begin{sourceblock}
\includegraphics[page=442,pagebox=mediabox,trim={45.00bp 257.18bp 44.37bp 278.44bp},clip,width=0.92\textwidth,max height=0.38\textheight]{Ishii.pdf}
\end{sourceblock}
这些分布参数和漂移速度的本构方程是由一维数据建立的，并且尚未通过详细的局部流动数据单独验证。因此，它们不应该单独使用。在段塞流、搅动流和环流流态中，分布参数效应比局部滑移效应占主导地位，即
0 gj V C j + + <<。因此，第 14 章 1.2 节中详述的本构方程可适用于此类流态。

\subsection*{14.4.6 减重条件下的本构方程}\phantomsection\addcontentsline{toc}{subsection}{14.4.6 减重条件下的本构方程}
为了将漂移通量模型的适用范围扩展到重力减小的条件下，第 14 章 1.2 节中详述的漂移速度本构方程已重新表述，除了重力压力损失之外还考虑了摩擦压力损失（Hibiki 等，2006）

气泡流状态（扭曲流体粒子状态）

\begin{sourceblock}
\includegraphics[page=443,pagebox=mediabox,trim={63.75bp 543.37bp 110.48bp 81.49bp},clip,width=\textwidth,max height=0.72\textheight]{Ishii.pdf}
\end{sourceblock}
\begin{sourceblock}
\includegraphics[page=443,pagebox=mediabox,trim={45.00bp 492.71bp 44.37bp 82.72bp},clip,width=0.92\textwidth,max height=0.38\textheight]{Ishii.pdf}
\end{sourceblock}
\begin{sourceblock}
\includegraphics[page=443,pagebox=mediabox,trim={45.00bp 404.89bp 44.37bp 196.84bp},clip,width=0.92\textwidth,max height=0.38\textheight]{Ishii.pdf}
\end{sourceblock}
\begin{sourceblock}
\includegraphics[page=443,pagebox=mediabox,trim={45.00bp 315.18bp 44.37bp 267.66bp},clip,width=0.92\textwidth,max height=0.38\textheight]{Ishii.pdf}
\end{sourceblock}
方程 (14-108) 对于 ( ) \{ \} ( ) \{ \} 成立

\begin{sourceblock}
\includegraphics[page=443,pagebox=mediabox,trim={45.00bp 237.46bp 44.37bp 369.65bp},clip,width=0.92\textwidth,max height=0.38\textheight]{Ishii.pdf}
\end{sourceblock}
\begin{sourceblock}
\includegraphics[page=443,pagebox=mediabox,trim={45.00bp 203.85bp 44.37bp 399.93bp},clip,width=0.92\textwidth,max height=0.38\textheight]{Ishii.pdf}
\end{sourceblock}
段塞流状态

\begin{sourceblock}
\includegraphics[page=443,pagebox=mediabox,trim={45.00bp 71.92bp 44.37bp 489.64bp},clip,width=0.92\textwidth,max height=0.38\textheight]{Ishii.pdf}
\end{sourceblock}
搅动流状态

\begin{sourceblock}
\includegraphics[page=444,pagebox=mediabox,trim={63.75bp 543.25bp 199.28bp 82.27bp},clip,width=\textwidth,max height=0.72\textheight]{Ishii.pdf}
\end{sourceblock}
\begin{sourceblock}
\includegraphics[page=444,pagebox=mediabox,trim={45.00bp 483.77bp 44.37bp 83.02bp},clip,width=0.92\textwidth,max height=0.38\textheight]{Ishii.pdf}
\end{sourceblock}
环形流态 在分离流中，两相之间的局部相对速度无法定义（Hibiki 和 Ishii，2003c）。如果气核中夹带小液滴或液膜中夹带小气泡，则由于气体和液体速度较大，局部相对速度可能近似为零，导致0 gj V ≈。这种近似在环形流中是可以接受的，其中从薄膜到气芯流的液体夹带小得可以忽略不计。

通过考虑重力对空隙分布的影响，对分布参数的本构方程进行了改进（Hibiki et al., 2006）：

气泡流动状态

\begin{sourceblock}
\includegraphics[page=444,pagebox=mediabox,trim={45.00bp 65.56bp 44.37bp 367.45bp},clip,width=0.92\textwidth,max height=0.38\textheight]{Ishii.pdf}
\end{sourceblock}
段塞流状态

\begin{sourceblock}
\includegraphics[page=445,pagebox=mediabox,trim={45.00bp 523.98bp 44.37bp 99.56bp},clip,width=0.92\textwidth,max height=0.38\textheight]{Ishii.pdf}
\end{sourceblock}
搅动流状态

\begin{sourceblock}
\includegraphics[page=445,pagebox=mediabox,trim={45.00bp 448.44bp 44.37bp 175.10bp},clip,width=0.92\textwidth,max height=0.38\textheight]{Ishii.pdf}
\end{sourceblock}
环形流态

\begin{sourceblock}
\includegraphics[page=445,pagebox=mediabox,trim={45.00bp 243.88bp 44.37bp 252.75bp},clip,width=0.92\textwidth,max height=0.38\textheight]{Ishii.pdf}
\end{sourceblock}
\subsection*{14.4.7 棒束几何本构方程}\phantomsection\addcontentsline{toc}{subsection}{14.4.7 棒束几何本构方程}
通过假设空隙率和混合物体积通量的幂律分布为，分析推导了方阵棒束几何形状中向上绝热两相流的分布参数的本构方程（Julia 等人，2009）

\begin{sourceblock}
\includegraphics[page=445,pagebox=mediabox,trim={45.00bp 91.13bp 44.37bp 531.93bp},clip,width=0.92\textwidth,max height=0.38\textheight]{Ishii.pdf}
\end{sourceblock}
The detailed analysis indicates that the distribution parameter depends on the ratio of rod diameter, 0 D , to pitch distance, 0 P , as

\begin{sourceblock}
\includegraphics[page=446,pagebox=mediabox,trim={45.00bp 510.81bp 44.37bp 55.49bp},clip,width=0.92\textwidth,max height=0.38\textheight]{Ishii.pdf}
\end{sourceblock}
公式(14-119)适用于0.2 0 0 D P ≤ ≤ 0.8。借助气泡层厚度模型（Hibiki 等人，2003a），上述方程已扩展到方形阵列棒束几何形状中的沸腾两相流：

\begin{sourceblock}
\includegraphics[page=446,pagebox=mediabox,trim={45.00bp 303.09bp 44.37bp 222.94bp},clip,width=0.92\textwidth,max height=0.38\textheight]{Ishii.pdf}
\end{sourceblock}
众所周知，当气泡直径与通道尺寸之比增大时，气泡上升速度会减小（Clift 等人，1978）。通过考虑气泡尺寸因子 sf B，对方程 (14- 99) 进行了修改，如下所示

\begin{sourceblock}
\includegraphics[page=446,pagebox=mediabox,trim={45.00bp 198.25bp 44.37bp 402.62bp},clip,width=0.92\textwidth,max height=0.38\textheight]{Ishii.pdf}
\end{sourceblock}
考虑棒壁对气泡上升速度影响的气泡尺寸因子由下式给出

\begin{sourceblock}
\includegraphics[page=446,pagebox=mediabox,trim={45.00bp 65.81bp 44.37bp 512.98bp},clip,width=0.92\textwidth,max height=0.38\textheight]{Ishii.pdf}
\end{sourceblock}
\begin{sourceblock}
\includegraphics[page=447,pagebox=mediabox,trim={45.00bp 558.95bp 44.37bp 58.83bp},clip,width=0.92\textwidth,max height=0.38\textheight]{Ishii.pdf}
\end{sourceblock}
\subsection*{14.4.8 池杆束几何本构方程}\phantomsection\addcontentsline{toc}{subsection}{14.4.8 池杆束几何本构方程}
池棒束几何形状中的两相流很复杂，因为再循环流型可能在低流速下形成。根据流动条件，两相流特征长度可以从子通道水力直径变为通道箱长度。池棒束几何形状的漂移速度本构方程由式（14-105）近似。池杆束几何分布参数的本构方程由下式给出（Chen et al., 2010）

\begin{sourceblock}
\includegraphics[page=447,pagebox=mediabox,trim={45.00bp 296.61bp 44.37bp 255.89bp},clip,width=0.92\textwidth,max height=0.38\textheight]{Ishii.pdf}
\end{sourceblock}
这些分布参数和漂移速度的本构方程是根据一维数据得出的，并且尚未通过详细的局部流动数据单独验证。因此，它们不应该单独使用。

\chapter{一维双流体模型}
双流体模型是两相系统热流体动力学最详细、最准确的宏观表述。在双流体模型中，场方程由每个相的质量、动量和能量方程组成的六个守恒方程表示。由于这些场方程是从局部瞬时平衡方程的适当平均获得的，因此相相互作用项出现在每个平均平衡方程中。这些术语表示通过相之间的界面输运的质量、动量和能量。界面输运项的存在是二流体模型公式最重要的特征之一。这些项决定了相变速率以及相间机械和热不平衡的程度，因此它们是应该准确建模的基本闭合关系。然而，由于测量和建模方面存在相当大的困难，界面输运项的可靠且准确的闭合关系尚未完全开发。尽管二流体模型存在这些缺点，但是没有任何替代品可以准确地模拟两相弱耦合的两相现象。这些例子有：

• Sudden mixing of two phases; • Transient flooding and flow reversal; • Transient countercurrent flow; • Two-phase flow with sudden acceleration. A three-dimensional, two-fluid model has been obtained by using temporal or statistical averaging.考虑到实际工程问题，通过对横截面积上的局部二流体公式进行平均而获得的一维二流体模型对于涉及流体流动和传热的复杂工程问题是有用的。这是由于

\section*{15.1 三维二流体模型面积平均}\phantomsection\addcontentsline{toc}{section}{15.1 三维二流体模型面积平均}
通过时间或统计平均方法获得了二流体模型的三维形式。对于大多数实际应用，Ishii（1975）开发的模型可以简化为以下形式：

连续性方程

\begin{sourceblock}
\includegraphics[page=449,pagebox=mediabox,trim={45.00bp 320.61bp 44.37bp 304.67bp},clip,width=0.92\textwidth,max height=0.38\textheight]{Ishii.pdf}
\end{sourceblock}
动量方程

\begin{sourceblock}
\includegraphics[page=449,pagebox=mediabox,trim={45.00bp 197.50bp 44.37bp 375.68bp},clip,width=0.92\textwidth,max height=0.38\textheight]{Ishii.pdf}
\end{sourceblock}
焓能方程

\begin{sourceblock}
\includegraphics[page=449,pagebox=mediabox,trim={45.00bp 90.92bp 44.37bp 501.38bp},clip,width=0.92\textwidth,max height=0.38\textheight]{Ishii.pdf}
\end{sourceblock}
这里 k Г 、 ik M 、 ki T 、 ki q ′′ 和 k Φ 是质量生成、广义界面阻力、界面剪切应力、界面热通量和耗散，

\begin{sourceblock}
\includegraphics[page=450,pagebox=mediabox,trim={45.00bp 507.33bp 44.37bp 121.89bp},clip,width=0.92\textwidth,max height=0.38\textheight]{Ishii.pdf}
\end{sourceblock}
从上面的场方程可以看出，方程的右侧出现了几个界面输运项。由于这些界面输运项也应遵守界面处的平衡定律，因此可以从局部跳跃条件的平均值获得界面输运条件（Ishii，1975）。它们由下式给出：

\begin{sourceblock}
\includegraphics[page=450,pagebox=mediabox,trim={45.00bp 309.25bp 44.37bp 240.82bp},clip,width=0.92\textwidth,max height=0.38\textheight]{Ishii.pdf}
\end{sourceblock}
因此，对于界面输运项来说，ik M 、 i fi a q '' 和 i gi a q '' 的本构方程是必要的。焓界面输运条件表明，如果可以忽略机械能传递项，则指定两相界面处的热通量相​​当于 k Γ 的本构关系（Ishii，1975）。这方面大大简化了界面转移项本构关系的发展。

获得一维模型的合理方法是在横截面上积分三维模型，然后引入适当的平均值。横截面积A上的简单面积平均值由式(14-5)定义，空隙率加权平均值由式(14-6)给出。在后续的分析中，认为任意横截面积内各相的密度是均匀的，使得 k k ρ ρ =。对于大多数实际的两相流问题，该假设是有效的，因为通道内的横向压力梯度相对较小。 k 相加权平均速度的轴向分量为

\begin{sourceblock}
\includegraphics[page=451,pagebox=mediabox,trim={45.00bp 566.74bp 44.37bp 56.68bp},clip,width=0.92\textwidth,max height=0.38\textheight]{Ishii.pdf}
\end{sourceblock}
其中速度的标量表达式对应于矢量的轴向分量。通过对式（15-1）至（15-3）进行面积平均，并进行一些适用于大多数实际问题的简化，可以得到以下场方程
连续性方程

\begin{sourceblock}
\includegraphics[page=451,pagebox=mediabox,trim={45.00bp 426.72bp 44.37bp 197.62bp},clip,width=0.92\textwidth,max height=0.38\textheight]{Ishii.pdf}
\end{sourceblock}
动量方程

\begin{sourceblock}
\includegraphics[page=451,pagebox=mediabox,trim={45.00bp 261.61bp 44.37bp 268.44bp},clip,width=0.92\textwidth,max height=0.38\textheight]{Ishii.pdf}
\end{sourceblock}
其中 kw α 和 kw τ 分别是壁面应力和壁面剪切应力下的平均空隙率。压力差和空隙梯度项对于水平分层流很重要。除此情况外，该术语可以被忽略。 d k M 是总界面剪切力，由下式给出

\begin{sourceblock}
\includegraphics[page=451,pagebox=mediabox,trim={45.00bp 162.64bp 44.37bp 459.48bp},clip,width=0.92\textwidth,max height=0.38\textheight]{Ishii.pdf}
\end{sourceblock}
右侧第一项是广义粒子阻力，这对于分散流很重要。第二项是界面剪切和空隙梯度的影响。该术语对于分离流尤其重要。在对流项中，k 相动量的分布参数 vk C 是由于变量乘积的平均值与平均变量乘积之间的差异而出现的。

焓能方程

\begin{sourceblock}
\includegraphics[page=452,pagebox=mediabox,trim={45.00bp 483.17bp 44.37bp 80.10bp},clip,width=0.92\textwidth,max height=0.38\textheight]{Ishii.pdf}
\end{sourceblock}
其中 h xi 和 kw q '' 分别是受热周界和壁面热通量。

hk C 是 k 相焓的分布参数。从界面处的宏观跳跃条件来看，界面转移项之间存在以下关系：

\begin{sourceblock}
\includegraphics[page=452,pagebox=mediabox,trim={45.00bp 298.38bp 44.37bp 251.80bp},clip,width=0.92\textwidth,max height=0.38\textheight]{Ishii.pdf}
\end{sourceblock}
\section*{15.2 一维本构关系的特殊考虑}\phantomsection\addcontentsline{toc}{section}{15.2 一维本构关系的特殊考虑}
\subsection*{15.2.1 场方程中的协方差效应}\phantomsection\addcontentsline{toc}{subsection}{15.2.1 场方程中的协方差效应}
在一维模型中，非常仔细地检查各种变量的横向分布及其对平衡和本构方程的影响是至关重要的。如果操作不当，所得的两相流配方可能会不一致。不正确的建模或忽视分布效应不仅可能导致模型非常不准确，而且还会导致各种数值不稳定。分布效应可以分为两组。第一个是协方差效应，它直接影响场方程中对流项的形式。第二个效应出现在各种局部本构关系的平均中。下面分别讨论这两种效应。

对流项的协方差定义为

\begin{sourceblock}
\includegraphics[page=453,pagebox=mediabox,trim={45.00bp 581.10bp 44.37bp 59.08bp},clip,width=0.92\textwidth,max height=0.38\textheight]{Ishii.pdf}
\end{sourceblock}
为了闭合控制方程组，我们必须指定这些协方差项的关系。这可以通过引入动量和能量通量的分布参数来完成。如果我们将通量的分布参数定义为

\begin{sourceblock}
\includegraphics[page=453,pagebox=mediabox,trim={45.00bp 464.87bp 44.37bp 158.65bp},clip,width=0.92\textwidth,max height=0.38\textheight]{Ishii.pdf}
\end{sourceblock}
协方差项变为

\begin{sourceblock}
\includegraphics[page=453,pagebox=mediabox,trim={45.00bp 377.59bp 44.37bp 238.96bp},clip,width=0.92\textwidth,max height=0.38\textheight]{Ishii.pdf}
\end{sourceblock}
对于动量通量，分布参数定义为

\begin{sourceblock}
\includegraphics[page=453,pagebox=mediabox,trim={45.00bp 296.92bp 44.37bp 323.56bp},clip,width=0.92\textwidth,max height=0.38\textheight]{Ishii.pdf}
\end{sourceblock}
从物理上讲，vk C 表示空隙和动量通量分布对 k 相横截面积平均动量通量的影响。 vk C 的定量研究可以通过考虑圆形管道中的对称流动并引入幂律表达式与漂移通量模型分析并行来进行（Zuber 和 Findlay，1965；Ishii，1977）。以下是 Ishii (1977) 得出的总结。

因此，对于气泡流、段塞流和搅动湍流，假设

\begin{sourceblock}
\includegraphics[page=453,pagebox=mediabox,trim={45.00bp 140.75bp 44.37bp 480.86bp},clip,width=0.92\textwidth,max height=0.38\textheight]{Ishii.pdf}
\end{sourceblock}
and

\begin{sourceblock}
\includegraphics[page=453,pagebox=mediabox,trim={45.00bp 63.76bp 44.37bp 557.81bp},clip,width=0.92\textwidth,max height=0.38\textheight]{Ishii.pdf}
\end{sourceblock}
\begin{sourceblock}
\includegraphics[page=454,pagebox=mediabox,trim={45.00bp 434.02bp 44.37bp 158.66bp},clip,width=0.92\textwidth,max height=0.38\textheight]{Ishii.pdf}
\end{sourceblock}
其中 0 k k kw Δα α α = −。漂移通量模型的体积通量分布参数 0 C 由下式给出

\begin{sourceblock}
\includegraphics[page=454,pagebox=mediabox,trim={45.00bp 348.12bp 44.37bp 280.50bp},clip,width=0.92\textwidth,max height=0.38\textheight]{Ishii.pdf}
\end{sourceblock}
其中 0 C 可以通过以下经验相关性给出（Ishii，1977）

\begin{sourceblock}
\includegraphics[page=454,pagebox=mediabox,trim={45.00bp 290.98bp 44.37bp 350.24bp},clip,width=0.92\textwidth,max height=0.38\textheight]{Ishii.pdf}
\end{sourceblock}
为了在圆管中充分发展流动。对于矩形通道中的过冷沸腾或流动，分别参见式（14-35）或式（14-32）。因此，在 n 的标准范围内，参数 vg C 可以近似给出：

\begin{sourceblock}
\includegraphics[page=454,pagebox=mediabox,trim={45.00bp 190.00bp 44.37bp 451.29bp},clip,width=0.92\textwidth,max height=0.38\textheight]{Ishii.pdf}
\end{sourceblock}
对于蒸气分散流状态下的液相，1 fw α 和
0 1 f α < .则由式（15-18）可知，对于标准范围

\begin{sourceblock}
\includegraphics[page=454,pagebox=mediabox,trim={45.00bp 104.98bp 44.37bp 512.85bp},clip,width=0.92\textwidth,max height=0.38\textheight]{Ishii.pdf}
\end{sourceblock}
对于环形流，动量协方差项也可以通过使用蒸汽和液体流的标准速度分布来计算。这样我们就得到

\begin{sourceblock}
\includegraphics[page=455,pagebox=mediabox,trim={45.00bp 558.72bp 44.37bp 55.70bp},clip,width=0.92\textwidth,max height=0.38\textheight]{Ishii.pdf}
\end{sourceblock}
类似地，焓通量的分布参数可以定义为

\begin{sourceblock}
\includegraphics[page=455,pagebox=mediabox,trim={45.00bp 472.78bp 44.37bp 150.65bp},clip,width=0.92\textwidth,max height=0.38\textheight]{Ishii.pdf}
\end{sourceblock}
对于热平衡流， g gs h h = 和 f fs h h =，其中 gs h 和 fs H 是蒸气和液体的饱和焓。在这种情况下，由于每个相的焓曲线完全平坦，因此分布参数变得一致；即 1 hg hf C C ==。同样明显的是，如果其中一个相处于饱和状态，则该相的 hk C 变为 1。在单相区域中，分布参数可以根据速度和焓的假设分布来计算。使用标准

\begin{sourceblock}
\includegraphics[page=455,pagebox=mediabox,trim={80.97bp 321.01bp 69.56bp 294.25bp},clip,width=\textwidth,max height=0.72\textheight]{Ishii.pdf}
\end{sourceblock}
可以表明，对于正常条件下的发展流和完全发展流，协方差项都小得可以忽略不计。然后

\begin{sourceblock}
\includegraphics[page=455,pagebox=mediabox,trim={45.00bp 267.03bp 44.37bp 374.72bp},clip,width=0.92\textwidth,max height=0.38\textheight]{Ishii.pdf}
\end{sourceblock}
因此，除了高度瞬态的情况外，焓协方差可以忽略不计。

\subsection*{15.2.2 相分布对本构关系的影响}\phantomsection\addcontentsline{toc}{subsection}{15.2.2 相分布对本构关系的影响}
传统双流体模型的最大缺点在于对式（15-9）定义的界面剪切d k M 本构方程进行建模。当二流体模型应用于非分离流时尤其如此。问题是双重的：

1. 平均阻力 ik z M 的建模； 2. 界面剪切效应的建模 k i z α τ −∇ ⋅。这些将在下面分别讨论。

对于分散两相流，平均界面阻力项可近似由下式给出

\begin{sourceblock}
\includegraphics[page=456,pagebox=mediabox,trim={45.00bp 569.80bp 44.37bp 57.53bp},clip,width=0.92\textwidth,max height=0.38\textheight]{Ishii.pdf}
\end{sourceblock}
这里，仅考虑 ik M 的稳态阻力部分，因为它是最重要的项。上述近似形式是基于实验观察到整个流道中局部相对速度 r v 相对均匀（Serizawa et al.，1975；Hibiki and Ishii，1999；Hibiki et al.，2001a）以及大多数两相流中局部相对速度远小于相速度的事实。

然而，重要的一点是平均阻力应该与平均局部相对速度 r v 相关，由下式给出：

\begin{sourceblock}
\includegraphics[page=456,pagebox=mediabox,trim={45.00bp 408.52bp 44.37bp 223.16bp},clip,width=0.92\textwidth,max height=0.38\textheight]{Ishii.pdf}
\end{sourceblock}
而不是由下式给出的区域平均平均速度之间的差异

\begin{sourceblock}
\includegraphics[page=456,pagebox=mediabox,trim={45.00bp 339.10bp 44.37bp 302.16bp},clip,width=0.92\textwidth,max height=0.38\textheight]{Ishii.pdf}
\end{sourceblock}
一般来说，

\begin{sourceblock}
\includegraphics[page=456,pagebox=mediabox,trim={45.00bp 281.14bp 44.37bp 360.05bp},clip,width=0.92\textwidth,max height=0.38\textheight]{Ishii.pdf}
\end{sourceblock}
这两个相对速度之间的差异可能非常大。原因是在一维公式中，两相之间的滑移 r v 是由两种完全不同的效应引起的；即局部相对运动以及相位和速度分布的积分效应。这两种效应的存在已经众所周知（Zuber 和 Findlay，1965；Ishii，1977；Bankoff，1960）。第一个效应是局部点处两相之间的真实相对运动，不需要任何进一步的解释。分布的第二个效应是由于面积平均而产生的。例如，如果分散相更集中在高速核心区域，则分散相的平均速度应远高于集中在低速壁区域附近的连续相的平均速度。即使当两个相以相同的速度局部移动时也是如此。

根据漂移通量模型公式，可以看出 r v 的近似表达式由下式给出

\begin{sourceblock}
\includegraphics[page=457,pagebox=mediabox,trim={45.00bp 564.40bp 44.37bp 57.34bp},clip,width=0.92\textwidth,max height=0.38\textheight]{Ishii.pdf}
\end{sourceblock}
适用于气泡流、段塞流和搅动湍流。因此，从流态准则（Ishii and Mishima，1981；Ishii，1977）来看，它适用于以下条件。

\begin{sourceblock}
\includegraphics[page=457,pagebox=mediabox,trim={45.00bp 423.87bp 44.37bp 160.91bp},clip,width=0.92\textwidth,max height=0.38\textheight]{Ishii.pdf}
\end{sourceblock}
当管径相对较小时，该标准有效。有关更一般的条件，请参阅 Ishii 和 Mishima (1981)。简单情况下 0 C 的本构方程由方程（15-20）给出。方程（15-26）和方程（15-30）给出的曳力表达式补偿了由于相位和速度分布而产生的滑移。传统的双流体模型从未考虑到 r v 和 r v 之间的这种差异。在大多数两相流系统中，由于相分布而产生的滑移远大于相间的局部滑移。因此，忽略上述效应将导致气泡、段塞和搅动湍流流态中的空隙率和速度的预测出现较大误差。因此，在这些流态中，即使是双流体模型的稳态预测也不如漂移通量模型的稳态预测。这是传统双流体模型最显著的缺点之一，应该在未来的所有分析中予以纠正。

\subsection*{15.2.3 界面剪切项}\phantomsection\addcontentsline{toc}{subsection}{15.2.3 界面剪切项}
\begin{sourceblock}
\includegraphics[page=457,pagebox=mediabox,trim={154.05bp 135.19bp 129.38bp 477.61bp},clip,width=\textwidth,max height=0.72\textheight]{Ishii.pdf}
\end{sourceblock}
由 d k M 表示的总界面剪切力有两个来源：即广义阻力 ik z M 以及界面剪切和空隙梯度的贡献 k i z α τ −∇ ⋅ 如式（15-9）所示。在分离流中，第二项占主导地位。例如，对于管内的环形流，可以证明：

\begin{sourceblock}
\includegraphics[page=458,pagebox=mediabox,trim={45.00bp 528.31bp 44.37bp 55.38bp},clip,width=0.92\textwidth,max height=0.38\textheight]{Ishii.pdf}
\end{sourceblock}
其中 i 是气芯的润湿周长。在这种情况下，gi τ 的本构关系可以用标准界面摩擦系数给出：

\begin{sourceblock}
\includegraphics[page=458,pagebox=mediabox,trim={45.00bp 440.14bp 44.37bp 196.67bp},clip,width=0.92\textwidth,max height=0.38\textheight]{Ishii.pdf}
\end{sourceblock}
其中 r g f v v v = −。界面摩擦系数 i f 有许多相关性。 Wallis 相关性由下式给出

\begin{sourceblock}
\includegraphics[page=458,pagebox=mediabox,trim={45.00bp 350.97bp 44.37bp 280.10bp},clip,width=0.92\textwidth,max height=0.38\textheight]{Ishii.pdf}
\end{sourceblock}
适用于薄膜粗糙的情况。

对于环形流，传统的双流体模型已正确考虑了界面剪切项。然而，该术语在气泡流、段塞流和搅动湍流中的影响通常被忽略。包含该术语对于相间界面动量耦合的正确建模非常重要。为了获得该界面剪切项的本构关系，需要进行一些假设，因为它需要有关空隙和剪切应力分布的信息。为此，假设以下幂律分布

\begin{sourceblock}
\includegraphics[page=458,pagebox=mediabox,trim={45.00bp 159.77bp 44.37bp 461.84bp},clip,width=0.92\textwidth,max height=0.38\textheight]{Ishii.pdf}
\end{sourceblock}
由此以及式(15-16)的空隙剖面可以看出

\begin{sourceblock}
\includegraphics[page=458,pagebox=mediabox,trim={45.00bp 88.47bp 44.37bp 537.52bp},clip,width=0.92\textwidth,max height=0.38\textheight]{Ishii.pdf}
\end{sourceblock}
where g α is the void fraction of dispersed phase. By introducing the distribution parameter C τ given by

\begin{sourceblock}
\includegraphics[page=459,pagebox=mediabox,trim={45.00bp 571.47bp 44.37bp 57.52bp},clip,width=0.92\textwidth,max height=0.38\textheight]{Ishii.pdf}
\end{sourceblock}
分散两相流的界面剪切项变为

\begin{sourceblock}
\includegraphics[page=459,pagebox=mediabox,trim={45.00bp 501.52bp 44.37bp 124.54bp},clip,width=0.92\textwidth,max height=0.38\textheight]{Ishii.pdf}
\end{sourceblock}
其中 C τ 预计非常接近 1。在水平通道中，即使在稳态条件下，该项也会导致相之间的滑移。由于宏观动量跳跃条件，包含该术语不会改变两相混合物的整体动量平衡。然而，它表明相之间的动量相互作用通过界面剪切和空隙梯度分布受到壁剪切应力的影响。

\chapter{考虑控制体内结构材料的双流体模型}
几位作者已经很好地描述了微分平衡方程的局部瞬时公式（Ishii，1975；Ishii 和 Hibiki，2006；Drew 和 Passman，1998）。这些方程主要由于界面的存在而遇到各种困难。应用平均方法获得了宏观两相流场方程和平均值的本构方程。其中包括局部体积平均值（Sha et al., 1989）、时间平均（Ishii, 1975；Ishii and Hibiki, 2006）和整体平均值（Hill, 1998）。一些作者讨论了复合平均数（Delhaye et al., 1981, Lahey et al., 1988, Sha et al., 1989）。如图16-1所示，平均公式根据关注的长度尺度进行分类，即核热工水力系统分析代码中经常使用的一维公式、考虑流道中结构材料存在的多孔介质方法和计算流体动力学代码中使用的三维公式。

流道中结构材料的存在使得水力计算变得更加困难。几位作者（Slattery，1972；Todreas 和 Kazimi，1990）将局部体积平均方法用于单相流来解释结构材料。 Hughes (1976) 和 Sha 等人将相同的方法应用于两相流。 （1984、1989）。他们利用局部体积平均技术与平均微分守恒方程并行推导一些有用的参数，例如孔隙率和相分数。界面面积和结构材料接触面积对于完成推导都很重要，但这些变量仍然具有波动成分。因此，需要时间平均来获得最终的微分方程（Kolev，2002）。的

\begin{sourceblock}
\includegraphics[page=461,pagebox=mediabox,trim={50.00bp 180.54bp 44.37bp 55.00bp},clip,width=\textwidth,max height=0.72\textheight]{Ishii.pdf}
\par\vspace{0.45\baselineskip}
\small 图16-1。空间尺度与计算方法的关系(Lee et al., 2009) Sha等人采用了复合平均，即将时间平均应用于局部体积平均方程。 （1989）。由于时间和空间平均算子之间的交换律，两种方法应该给出相同的结果（Delhaye et al., 1981）。 Sha 和 Chao (2007) 为多相流守恒方程开发了一种新颖的多孔介质公式。
\end{sourceblock}
鉴于多孔介质方法在热工水力系统分析代码开发中的重要性，我们制定了时间-体积

\begin{center}\small 表16-1。性质、通量和来源（Lee et al., 2009）\end{center}
\begin{sourceblock}
\includegraphics[page=462,pagebox=mediabox,trim={45.00bp 557.73bp 44.37bp 83.38bp},clip,width=0.92\textwidth,max height=0.38\textheight]{Ishii.pdf}
\end{sourceblock}
\begin{sourceblock}
\includegraphics[page=462,pagebox=mediabox,trim={45.00bp 476.38bp 41.37bp 98.74bp},clip,width=\textwidth,max height=0.72\textheight]{Ishii.pdf}
\end{sourceblock}
本章中考虑控制体积中的结构材料的平均双流体模型（Lee et al., 2009）。首先，我们简要讨论时间平均二流体模型和推导时间-体积平均二流体模型的重要数学定理。其次，我们对流道上的时间平均双流体模型进行平均，以获得时间-体积平均双流体模型。最后，我们讨论了该公式中使用的本构方程的一些问题。

\section*{16.1 时间平均二流体模型}\phantomsection\addcontentsline{toc}{section}{16.1 时间平均二流体模型}
几位作者已经很好地描述了微分平衡方程的局部瞬时公式（Ishii，1975；Ishii 和 Hibiki，2006；Drew 和 Passman，1998）：

\begin{sourceblock}
\includegraphics[page=462,pagebox=mediabox,trim={45.00bp 243.93bp 44.37bp 381.04bp},clip,width=0.92\textwidth,max height=0.38\textheight]{Ishii.pdf}
\end{sourceblock}
其中k ρ 、k ψ 、k v 、k J 和k φ 分别是k 相的密度、广延特性、速度、通量和来源。表 16-1 显示了连续性、动量和能量守恒方程的场变量。平衡方程由界面平衡方程、各种边界条件和状态方程补充。

时间平均方法也已得到很好的建立（Ishii，1975；Ishii 和 Hibiki，2006）。 k 相任意性质 k ψ 的平衡方程可以表示为

\begin{sourceblock}
\includegraphics[page=462,pagebox=mediabox,trim={45.00bp 82.89bp 44.37bp 542.07bp},clip,width=0.92\textwidth,max height=0.38\textheight]{Ishii.pdf}
\end{sourceblock}
其中量上的条表示时间平均操作，k I 是 k 相平衡方程中的界面源项。

通过定义相密度函数 k M 我们可以定义各种权重变量。

\begin{sourceblock}
\includegraphics[page=463,pagebox=mediabox,trim={45.00bp 517.15bp 44.37bp 93.77bp},clip,width=0.92\textwidth,max height=0.38\textheight]{Ishii.pdf}
\end{sourceblock}
然后我们有

\begin{sourceblock}
\includegraphics[page=463,pagebox=mediabox,trim={45.00bp 460.54bp 44.37bp 180.70bp},clip,width=0.92\textwidth,max height=0.38\textheight]{Ishii.pdf}
\end{sourceblock}
场变量 k F 的相位平均值也定义为

\begin{sourceblock}
\includegraphics[page=463,pagebox=mediabox,trim={45.00bp 386.20bp 44.37bp 241.28bp},clip,width=0.92\textwidth,max height=0.38\textheight]{Ishii.pdf}
\end{sourceblock}
如果k F 由相密度k ρ 和性质k ψ 的乘积表示，则相密度加权平均变量可以定义如下。

\begin{sourceblock}
\includegraphics[page=463,pagebox=mediabox,trim={45.00bp 298.18bp 44.37bp 326.13bp},clip,width=0.92\textwidth,max height=0.38\textheight]{Ishii.pdf}
\end{sourceblock}
也可以写成

\begin{sourceblock}
\includegraphics[page=463,pagebox=mediabox,trim={45.00bp 221.20bp 44.37bp 403.13bp},clip,width=0.92\textwidth,max height=0.38\textheight]{Ishii.pdf}
\end{sourceblock}
那么时间平均平衡方程可以用加权平均变量表示为

\begin{sourceblock}
\includegraphics[page=463,pagebox=mediabox,trim={45.00bp 108.15bp 44.37bp 491.39bp},clip,width=0.92\textwidth,max height=0.38\textheight]{Ishii.pdf}
\end{sourceblock}
在这里，我们将时间波动项定义如下。

\begin{sourceblock}
\includegraphics[page=463,pagebox=mediabox,trim={45.00bp 50.14bp 44.37bp 590.11bp},clip,width=0.92\textwidth,max height=0.38\textheight]{Ishii.pdf}
\end{sourceblock}
\begin{sourceblock}
\includegraphics[page=464,pagebox=mediabox,trim={50.00bp 387.06bp 44.37bp 55.00bp},clip,width=\textwidth,max height=0.72\textheight]{Ishii.pdf}
\par\vspace{0.45\baselineskip}
\small 图16-2。控制体积，包括结构材料（Lee et al., 2009）然后
\end{sourceblock}
\begin{sourceblock}
\includegraphics[page=464,pagebox=mediabox,trim={45.00bp 310.91bp 44.37bp 328.83bp},clip,width=0.92\textwidth,max height=0.38\textheight]{Ishii.pdf}
\end{sourceblock}
因此我们有

\begin{sourceblock}
\includegraphics[page=464,pagebox=mediabox,trim={45.00bp 217.46bp 44.37bp 388.38bp},clip,width=0.92\textwidth,max height=0.38\textheight]{Ishii.pdf}
\end{sourceblock}
\section*{16.2 本地音量平均操作}\phantomsection\addcontentsline{toc}{section}{16.2 本地音量平均操作}
\subsection*{16.2.1 参数和平均量的定义}\phantomsection\addcontentsline{toc}{subsection}{16.2.1 参数和平均量的定义}
体积平均过程通常应用于由于相界面而导致的不连续函数。这里，局部体积平均过程应用于通过时间平均过程平滑的场变量。因此，相界面并不重要，而仅应考虑固体结构和流体之间的界面。图 16-2 显示了局部体积元素。时间-体积平均平衡方程可

\begin{sourceblock}
\includegraphics[page=465,pagebox=mediabox,trim={45.00bp 521.66bp 44.37bp 122.87bp},clip,width=0.92\textwidth,max height=0.38\textheight]{Ishii.pdf}
\par\vspace{0.45\baselineskip}
\small 图16-2中F A 和S A 分别是控制体积边界上的自由表面和结构材料表面。因此，控制体积 T A 的总边界表面由下式给出
\end{sourceblock}
\begin{sourceblock}
\includegraphics[page=465,pagebox=mediabox,trim={45.00bp 440.66bp 44.37bp 203.87bp},clip,width=0.92\textwidth,max height=0.38\textheight]{Ishii.pdf}
\end{sourceblock}
体积孔隙率和表面孔隙率通常分别如式(16-14)和(16-15)所示定义。

\begin{sourceblock}
\includegraphics[page=465,pagebox=mediabox,trim={45.00bp 356.38bp 44.37bp 271.31bp},clip,width=0.92\textwidth,max height=0.38\textheight]{Ishii.pdf}
\end{sourceblock}
\begin{sourceblock}
\includegraphics[page=465,pagebox=mediabox,trim={45.00bp 298.48bp 44.37bp 329.21bp},clip,width=0.92\textwidth,max height=0.38\textheight]{Ishii.pdf}
\end{sourceblock}
\begin{sourceblock}
\includegraphics[page=465,pagebox=mediabox,trim={45.00bp 229.42bp 44.37bp 382.53bp},clip,width=0.92\textwidth,max height=0.38\textheight]{Ishii.pdf}
\end{sourceblock}
如果数量k R 是常数，则式(16-16)表明相位平均数量不等于该常数值。因此，我们通常更喜欢使用由 (Whitaker, 1985) 定义的本征相位平均值

\begin{sourceblock}
\includegraphics[page=465,pagebox=mediabox,trim={45.00bp 121.42bp 44.37bp 507.64bp},clip,width=0.92\textwidth,max height=0.38\textheight]{Ishii.pdf}
\end{sourceblock}
由于 k R 不存在于固体结构内部，因此我们有

\begin{sourceblock}
\includegraphics[page=465,pagebox=mediabox,trim={45.00bp 57.95bp 44.37bp 576.29bp},clip,width=0.92\textwidth,max height=0.38\textheight]{Ishii.pdf}
\end{sourceblock}
这两个平均数量的关系为

\begin{sourceblock}
\includegraphics[page=466,pagebox=mediabox,trim={45.00bp 542.06bp 44.37bp 93.69bp},clip,width=0.92\textwidth,max height=0.38\textheight]{Ishii.pdf}
\end{sourceblock}
局部表面 A 上标量、向量或张量 k R 的平均值定义为

\begin{sourceblock}
\includegraphics[page=466,pagebox=mediabox,trim={45.00bp 461.32bp 44.37bp 167.74bp},clip,width=0.92\textwidth,max height=0.38\textheight]{Ishii.pdf}
\end{sourceblock}
\begin{sourceblock}
\includegraphics[page=466,pagebox=mediabox,trim={45.00bp 392.32bp 44.37bp 219.63bp},clip,width=0.92\textwidth,max height=0.38\textheight]{Ishii.pdf}
\end{sourceblock}
由于 k R 不存在于固体结构内部，因此我们有

\begin{sourceblock}
\includegraphics[page=466,pagebox=mediabox,trim={45.00bp 328.85bp 44.37bp 305.39bp},clip,width=0.92\textwidth,max height=0.38\textheight]{Ishii.pdf}
\end{sourceblock}
这两个平均数量的关系为

\begin{sourceblock}
\includegraphics[page=466,pagebox=mediabox,trim={45.00bp 267.07bp 44.37bp 368.66bp},clip,width=0.92\textwidth,max height=0.38\textheight]{Ishii.pdf}
\end{sourceblock}
\subsection*{16.2.2 一些重要定理}\phantomsection\addcontentsline{toc}{subsection}{16.2.2 一些重要定理}
散度定理表示为

\begin{sourceblock}
\includegraphics[page=466,pagebox=mediabox,trim={45.00bp 165.82bp 44.37bp 468.37bp},clip,width=0.92\textwidth,max height=0.38\textheight]{Ishii.pdf}
\end{sourceblock}
其中 F n 是垂直于表面 F A 的单位向量，I n 是垂直于控制体积 I A 边界内的结构材料表面的单位向量。固有局部体积平均值的散度由下式表示（Whitaker，1985）

\begin{sourceblock}
\includegraphics[page=466,pagebox=mediabox,trim={45.00bp 60.88bp 44.37bp 572.96bp},clip,width=0.92\textwidth,max height=0.38\textheight]{Ishii.pdf}
\end{sourceblock}
\begin{sourceblock}
\includegraphics[page=467,pagebox=mediabox,trim={45.00bp 549.82bp 44.37bp 84.02bp},clip,width=0.92\textwidth,max height=0.38\textheight]{Ishii.pdf}
\end{sourceblock}
\section*{16.3 时间-体积平均二流体模型公式}\phantomsection\addcontentsline{toc}{section}{16.3 时间-体积平均二流体模型公式}
\subsection*{16.3.1 仅具有体积孔隙率的配方}\phantomsection\addcontentsline{toc}{subsection}{16.3.1 仅具有体积孔隙率的配方}
将式(16-11)对F V 积分，除以体积元素T V，可得时间-体积平均平衡方程为

\begin{sourceblock}
\includegraphics[page=467,pagebox=mediabox,trim={45.00bp 338.92bp 44.37bp 253.60bp},clip,width=0.92\textwidth,max height=0.38\textheight]{Ishii.pdf}
\end{sourceblock}
对于固定体积（即静止且不可变形的体积），时间导数和体积积分的阶数可以互换。因此，式(16-27)左边第一项可表示为

\begin{sourceblock}
\includegraphics[page=467,pagebox=mediabox,trim={45.00bp 209.28bp 44.37bp 382.42bp},clip,width=0.92\textwidth,max height=0.38\textheight]{Ishii.pdf}
\end{sourceblock}
将三维莱布尼茨规则应用于第二项，得出

\begin{sourceblock}
\includegraphics[page=467,pagebox=mediabox,trim={45.00bp 72.42bp 44.37bp 485.86bp},clip,width=0.92\textwidth,max height=0.38\textheight]{Ishii.pdf}
\end{sourceblock}
如果我们忽略结构表面处液相相对于气相的滑移速度，则结构表面处的相速度可近似为零（I k A = v 处为 0）。因此

\begin{sourceblock}
\includegraphics[page=468,pagebox=mediabox,trim={45.00bp 507.33bp 44.37bp 97.90bp},clip,width=0.92\textwidth,max height=0.38\textheight]{Ishii.pdf}
\end{sourceblock}
方程（16-27）左边的第二项近似由下式给出

\begin{sourceblock}
\includegraphics[page=468,pagebox=mediabox,trim={45.00bp 425.31bp 44.37bp 203.75bp},clip,width=0.92\textwidth,max height=0.38\textheight]{Ishii.pdf}
\end{sourceblock}
类似地，式(16-27)右边第一项可表示为

\begin{sourceblock}
\includegraphics[page=468,pagebox=mediabox,trim={45.00bp 213.02bp 44.37bp 290.81bp},clip,width=0.92\textwidth,max height=0.38\textheight]{Ishii.pdf}
\end{sourceblock}
式(16-27)右边的第二项和第三项可以分别表示为式(16-33)和(16-34)。

\begin{sourceblock}
\includegraphics[page=468,pagebox=mediabox,trim={45.00bp 131.32bp 44.37bp 499.58bp},clip,width=0.92\textwidth,max height=0.38\textheight]{Ishii.pdf}
\end{sourceblock}
\begin{sourceblock}
\includegraphics[page=468,pagebox=mediabox,trim={45.00bp 75.34bp 44.37bp 553.73bp},clip,width=0.92\textwidth,max height=0.38\textheight]{Ishii.pdf}
\end{sourceblock}
那么式(16-27)可以改写为

\begin{sourceblock}
\includegraphics[page=469,pagebox=mediabox,trim={45.00bp 477.18bp 44.37bp 80.93bp},clip,width=0.92\textwidth,max height=0.38\textheight]{Ishii.pdf}
\end{sourceblock}
在后续分析中，我们假设流道上的 k ρ 为常数。方程（16-35）可以进一步简化为

\begin{sourceblock}
\includegraphics[page=469,pagebox=mediabox,trim={45.00bp 324.87bp 44.37bp 231.57bp},clip,width=0.92\textwidth,max height=0.38\textheight]{Ishii.pdf}
\end{sourceblock}
其中湍流定义为

\begin{sourceblock}
\includegraphics[page=469,pagebox=mediabox,trim={45.00bp 264.45bp 44.37bp 375.69bp},clip,width=0.92\textwidth,max height=0.38\textheight]{Ishii.pdf}
\end{sourceblock}
在式(16-35)中，仅利用体积孔隙率。

这里我们将一些平均数量定义为

\begin{sourceblock}
\includegraphics[page=469,pagebox=mediabox,trim={45.00bp 164.79bp 44.37bp 447.56bp},clip,width=0.92\textwidth,max height=0.38\textheight]{Ishii.pdf}
\end{sourceblock}
\begin{sourceblock}
\includegraphics[page=469,pagebox=mediabox,trim={45.00bp 89.59bp 44.37bp 520.58bp},clip,width=0.92\textwidth,max height=0.38\textheight]{Ishii.pdf}
\end{sourceblock}
\begin{sourceblock}
\includegraphics[page=470,pagebox=mediabox,trim={45.00bp 540.98bp 44.37bp 70.29bp},clip,width=0.92\textwidth,max height=0.38\textheight]{Ishii.pdf}
\end{sourceblock}
\begin{sourceblock}
\includegraphics[page=470,pagebox=mediabox,trim={45.00bp 472.94bp 44.37bp 143.36bp},clip,width=0.92\textwidth,max height=0.38\textheight]{Ishii.pdf}
\end{sourceblock}
\begin{sourceblock}
\includegraphics[page=470,pagebox=mediabox,trim={45.00bp 419.43bp 44.37bp 211.12bp},clip,width=0.92\textwidth,max height=0.38\textheight]{Ishii.pdf}
\end{sourceblock}
假设固体结构和流体混合物之间的边界处的湍流可以忽略不计，这可能是合理的。最后可以得到各相的时间-体积平均平衡方程

\begin{sourceblock}
\includegraphics[page=470,pagebox=mediabox,trim={45.00bp 219.94bp 44.37bp 301.85bp},clip,width=0.92\textwidth,max height=0.38\textheight]{Ishii.pdf}
\end{sourceblock}
对于质量守恒方程，我们将参数设置为

\begin{sourceblock}
\includegraphics[page=470,pagebox=mediabox,trim={45.00bp 166.24bp 44.37bp 475.11bp},clip,width=0.92\textwidth,max height=0.38\textheight]{Ishii.pdf}
\end{sourceblock}
其中 k Γ 是 k 相的质量生成。然后我们得到质量守恒方程为

\begin{sourceblock}
\includegraphics[page=470,pagebox=mediabox,trim={45.00bp 77.70bp 44.37bp 547.51bp},clip,width=0.92\textwidth,max height=0.38\textheight]{Ishii.pdf}
\end{sourceblock}
对于动量守恒方程，我们将参数设置为

\begin{sourceblock}
\includegraphics[page=471,pagebox=mediabox,trim={45.00bp 586.24bp 44.37bp 55.11bp},clip,width=0.92\textwidth,max height=0.38\textheight]{Ishii.pdf}
\end{sourceblock}
其中 k T 、k p 、I 、k T 、kg 和 k M 分别是 k 相的应力张量、压力、单位张量、粘性应力张量、重力加速度和动量源。然后我们得到动量守恒方程为

\begin{sourceblock}
\includegraphics[page=471,pagebox=mediabox,trim={45.00bp 318.85bp 44.37bp 153.49bp},clip,width=0.92\textwidth,max height=0.38\textheight]{Ishii.pdf}
\end{sourceblock}
其中湍流 ( ) T T k k equation− J T 定义为

\begin{sourceblock}
\includegraphics[page=471,pagebox=mediabox,trim={45.00bp 259.12bp 44.37bp 381.04bp},clip,width=0.92\textwidth,max height=0.38\textheight]{Ishii.pdf}
\end{sourceblock}
对于能量守恒方程，我们将参数设置为

\begin{sourceblock}
\includegraphics[page=471,pagebox=mediabox,trim={45.00bp 168.40bp 44.37bp 438.84bp},clip,width=0.92\textwidth,max height=0.38\textheight]{Ishii.pdf}
\end{sourceblock}
其中 k u 、k q 、k q \$ 和 k E 分别是 k 相的内能、热通量、物体加热和通过界面的总能量增益。身体加热项 k k q ρ \$ 经常被忽略。那么我们得到能量守恒方程为

\begin{sourceblock}
\includegraphics[page=472,pagebox=mediabox,trim={45.00bp 387.86bp 44.37bp 68.50bp},clip,width=0.92\textwidth,max height=0.38\textheight]{Ishii.pdf}
\end{sourceblock}
式中，表观内能 k e 和湍流热通量 T k q 分别定义为式（16-51）和式（16-52）。

\begin{sourceblock}
\includegraphics[page=472,pagebox=mediabox,trim={45.00bp 296.89bp 44.37bp 320.58bp},clip,width=0.92\textwidth,max height=0.38\textheight]{Ishii.pdf}
\end{sourceblock}
\begin{sourceblock}
\includegraphics[page=472,pagebox=mediabox,trim={45.00bp 225.74bp 44.37bp 394.18bp},clip,width=0.92\textwidth,max height=0.38\textheight]{Ishii.pdf}
\end{sourceblock}
除了时间-体积平均双流体模型中出现的孔隙率项、协方差项和固体表面源项外，这些守恒方程与局部时间平均双流体模型类似。这里还需要说明的是，时间-体积平均二流体模型中时间脉动项的定义与时间平均二流体模型中的定义一致。因此，现有的本构方程可能适用于时间-体积平均二流体模型。

\subsection*{16.3.2 具有体积和表面孔隙率的配方}\phantomsection\addcontentsline{toc}{subsection}{16.3.2 具有体积和表面孔隙率的配方}
在上一节中，我们推导了仅具有体积孔隙率的传统多孔介质公式。在这种方法中分布式
利用式（16-25），式（16-36）可改写为

\begin{sourceblock}
\includegraphics[page=473,pagebox=mediabox,trim={45.00bp 353.04bp 44.37bp 210.56bp},clip,width=0.92\textwidth,max height=0.38\textheight]{Ishii.pdf}
\end{sourceblock}
这里我们引入以下特殊运算符

\begin{sourceblock}
\includegraphics[page=473,pagebox=mediabox,trim={45.00bp 284.32bp 44.37bp 344.73bp},clip,width=0.92\textwidth,max height=0.38\textheight]{Ishii.pdf}
\end{sourceblock}
然后我们可以将方程（16-53）改写为

\begin{sourceblock}
\includegraphics[page=473,pagebox=mediabox,trim={45.00bp 122.35bp 44.37bp 411.18bp},clip,width=0.92\textwidth,max height=0.38\textheight]{Ishii.pdf}
\end{sourceblock}
为了解释左边的第二项，我们考虑如图 16-3 所示的笛卡尔坐标系的特殊情况，并且特殊运算符在笛卡尔坐标系中由方程（16-56）表示（Todreas 和 Kazimi，1990）。

\begin{sourceblock}
\includegraphics[page=474,pagebox=mediabox,trim={50.00bp 375.00bp 44.37bp 55.00bp},clip,width=\textwidth,max height=0.72\textheight]{Ishii.pdf}
\par\vspace{0.45\baselineskip}
\small 图16-3。控制体积，包括笛卡尔坐标系中的结构材料
\end{sourceblock}
（李等人，2009）

\begin{sourceblock}
\includegraphics[page=474,pagebox=mediabox,trim={45.00bp 186.80bp 44.37bp 335.00bp},clip,width=0.92\textwidth,max height=0.38\textheight]{Ishii.pdf}
\end{sourceblock}
详细解释见附录。这里我们定义一些平均数量和协方差为

\begin{sourceblock}
\includegraphics[page=474,pagebox=mediabox,trim={45.00bp 88.26bp 44.37bp 523.81bp},clip,width=0.92\textwidth,max height=0.38\textheight]{Ishii.pdf}
\end{sourceblock}
\begin{sourceblock}
\includegraphics[page=475,pagebox=mediabox,trim={45.00bp 553.82bp 44.37bp 57.26bp},clip,width=0.92\textwidth,max height=0.38\textheight]{Ishii.pdf}
\end{sourceblock}
\begin{sourceblock}
\includegraphics[page=475,pagebox=mediabox,trim={45.00bp 476.71bp 44.37bp 130.33bp},clip,width=0.92\textwidth,max height=0.38\textheight]{Ishii.pdf}
\end{sourceblock}
\begin{sourceblock}
\includegraphics[page=475,pagebox=mediabox,trim={45.00bp 423.53bp 44.37bp 207.07bp},clip,width=0.92\textwidth,max height=0.38\textheight]{Ishii.pdf}
\end{sourceblock}
k ψ C 式(16-55)的第二项在笛卡尔坐标中表示为
引入方向协方差项 A

\begin{sourceblock}
\includegraphics[page=475,pagebox=mediabox,trim={45.00bp 50.14bp 44.37bp 285.38bp},clip,width=0.92\textwidth,max height=0.38\textheight]{Ishii.pdf}
\end{sourceblock}
where

\begin{sourceblock}
\includegraphics[page=476,pagebox=mediabox,trim={45.00bp 512.13bp 44.37bp 95.59bp},clip,width=0.92\textwidth,max height=0.38\textheight]{Ishii.pdf}
\end{sourceblock}
假设固体结构和流体混合物之间的边界处的湍流通量可以忽略不计，这可能是合理的。最后可以得到各相的时间-体积平均平衡方程

\begin{sourceblock}
\includegraphics[page=476,pagebox=mediabox,trim={45.00bp 310.25bp 44.37bp 210.56bp},clip,width=0.92\textwidth,max height=0.38\textheight]{Ishii.pdf}
\end{sourceblock}
对于质量守恒方程，我们将参数设置为方程（16-44）。然后我们得到质量守恒方程为

\begin{sourceblock}
\includegraphics[page=476,pagebox=mediabox,trim={45.00bp 226.52bp 44.37bp 398.73bp},clip,width=0.92\textwidth,max height=0.38\textheight]{Ishii.pdf}
\end{sourceblock}
对于动量守恒方程，我们将参数设置为Eq. （16-46）。然后我们得到动量守恒方程为

\begin{sourceblock}
\includegraphics[page=476,pagebox=mediabox,trim={45.00bp 80.19bp 44.37bp 481.49bp},clip,width=0.92\textwidth,max height=0.38\textheight]{Ishii.pdf}
\end{sourceblock}
\begin{sourceblock}
\includegraphics[page=477,pagebox=mediabox,trim={45.00bp 548.71bp 44.37bp 59.83bp},clip,width=0.92\textwidth,max height=0.38\textheight]{Ishii.pdf}
\end{sourceblock}
对于能量守恒方程，我们将参数设置为式（16-49）。那么我们得到能量守恒方程为

\begin{sourceblock}
\includegraphics[page=477,pagebox=mediabox,trim={45.00bp 291.39bp 44.37bp 159.87bp},clip,width=0.92\textwidth,max height=0.38\textheight]{Ishii.pdf}
\end{sourceblock}
这些方程与局部时间平均二流体模型相似，除了孔隙率项、协方差项和固体表面源项出现在时间-体积平均二流体模型中。这里还需要说明的是，时间-体积平均二流体模型中时间脉动项的定义与时间平均二流体模型中的定义一致。因此，现有的本构方程可能适用于时间-体积平均二流体模型。

\section*{16.4 时间-体积平均本构关系的特别考虑}\phantomsection\addcontentsline{toc}{section}{16.4 时间-体积平均本构关系的特别考虑}
\subsection*{16.4.1 场方程中的协方差效应}\phantomsection\addcontentsline{toc}{subsection}{16.4.1 场方程中的协方差效应}
在体积平均模型中，非常仔细地检查各种变量的空间分布及其对平衡和本构方程的影响至关重要。如果操作不当，所得的两相流配方可能会不一致。不正确的建模或忽视分布效应不仅可能导致模型严重不准确，而且可能导致各种数值不稳定。分布效应可以分为两组。第一个是协方差效应，它直接影响场方程中对流项的形式。第二个效应出现在各种局部本构关系的平均中。下面分别讨论这两种效应（Ishii 和 Mishima，1984）。

对流项的协方差定义为

\begin{sourceblock}
\includegraphics[page=478,pagebox=mediabox,trim={45.00bp 429.17bp 44.37bp 199.15bp},clip,width=0.92\textwidth,max height=0.38\textheight]{Ishii.pdf}
\end{sourceblock}
为了闭合控制方程组，我们必须指定这些协方差项的关系。如上一节所示，这可以通过引入动量和能量通量的分布参数来完成。然后根据协方差的定义，方程（16-62），协方差项变为

\begin{sourceblock}
\includegraphics[page=478,pagebox=mediabox,trim={45.00bp 279.85bp 44.37bp 321.61bp},clip,width=0.92\textwidth,max height=0.38\textheight]{Ishii.pdf}
\end{sourceblock}
对于动量通量，分布参数定义为

\begin{sourceblock}
\includegraphics[page=478,pagebox=mediabox,trim={45.00bp 189.43bp 44.37bp 420.81bp},clip,width=0.92\textwidth,max height=0.38\textheight]{Ishii.pdf}
\end{sourceblock}
物理上vk C 表示空隙和动量通量分布对k 相控制面平均动量通量的影响。 vk C 的定量研究可以通过假设空隙和速度剖面来进行（Ishii 和 Mishima，1984）。同理还可以得到能量方程的分布参数。

\subsection*{16.4.2 相分布对本构关系的影响}\phantomsection\addcontentsline{toc}{subsection}{16.4.2 相分布对本构关系的影响}
k相界面动量增益k M 表示为

\begin{sourceblock}
\includegraphics[page=479,pagebox=mediabox,trim={45.00bp 571.80bp 44.37bp 57.31bp},clip,width=0.92\textwidth,max height=0.38\textheight]{Ishii.pdf}
\end{sourceblock}
Γ M 、 kip 、 ik M 和 ki T 分别是由于质量生成、界面压力、总广义阻力和界面剪切应力而产生的界面动量增益。在时间-体积平均二流体模型中，界面剪切本构方程的建模是必不可少的。当双流体模型应用于除分离流之外的其他流时尤其如此（Ishii 和 Mishima，1984）。问题是双重的：

where k

\begin{sourceblock}
\includegraphics[page=479,pagebox=mediabox,trim={45.00bp 415.38bp 44.37bp 194.44bp},clip,width=0.92\textwidth,max height=0.38\textheight]{Ishii.pdf}
\end{sourceblock}
这些将在下面单独讨论。

对于分散两相流，时间平均阻力表示为

\begin{sourceblock}
\includegraphics[page=479,pagebox=mediabox,trim={45.00bp 346.78bp 44.37bp 280.55bp},clip,width=0.92\textwidth,max height=0.38\textheight]{Ishii.pdf}
\end{sourceblock}
其中 DC 、 br 、 f ρ 和 r v 分别是阻力系数、气泡半径、液体密度和相对速度。固有时间体积平均阻力表示为

\begin{sourceblock}
\includegraphics[page=479,pagebox=mediabox,trim={45.00bp 244.85bp 44.37bp 376.74bp},clip,width=0.92\textwidth,max height=0.38\textheight]{Ishii.pdf}
\end{sourceblock}
在一维时间平均二流体模型中（Ishii 和 Mishima，1984；Ishii 和 Hibiki，2006），基于现有实验数据假设分散两相流具有均匀相对速度（Hibiki 和 Ishii，1999；Hibiki 等，2001a）。如果这适用于这种情况，则通常占主导地位的平均阻力的 z 分量可表示为

\begin{sourceblock}
\includegraphics[page=479,pagebox=mediabox,trim={45.00bp 113.88bp 44.37bp 507.53bp},clip,width=0.92\textwidth,max height=0.38\textheight]{Ishii.pdf}
\end{sourceblock}
where

\begin{sourceblock}
\includegraphics[page=479,pagebox=mediabox,trim={45.00bp 50.14bp 44.37bp 584.09bp},clip,width=0.92\textwidth,max height=0.38\textheight]{Ishii.pdf}
\end{sourceblock}
rz v 相的面积平均速度之差由下式给出

\begin{sourceblock}
\includegraphics[page=480,pagebox=mediabox,trim={45.00bp 504.18bp 44.37bp 108.49bp},clip,width=0.92\textwidth,max height=0.38\textheight]{Ishii.pdf}
\end{sourceblock}
一般来说，

\begin{sourceblock}
\includegraphics[page=480,pagebox=mediabox,trim={45.00bp 444.16bp 44.37bp 191.66bp},clip,width=0.92\textwidth,max height=0.38\textheight]{Ishii.pdf}
\end{sourceblock}
下面，通过多孔介质方法推导漂移通量模型。漂移速度定义为

\begin{sourceblock}
\includegraphics[page=480,pagebox=mediabox,trim={45.00bp 374.98bp 44.37bp 264.08bp},clip,width=0.92\textwidth,max height=0.38\textheight]{Ishii.pdf}
\end{sourceblock}
然后，空隙率加权平均漂移速度由下式给出

\begin{sourceblock}
\includegraphics[page=480,pagebox=mediabox,trim={45.00bp 221.14bp 44.37bp 324.92bp},clip,width=0.92\textwidth,max height=0.38\textheight]{Ishii.pdf}
\end{sourceblock}
\begin{sourceblock}
\includegraphics[page=480,pagebox=mediabox,trim={45.00bp 187.54bp 44.37bp 453.64bp},clip,width=0.92\textwidth,max height=0.38\textheight]{Ishii.pdf}
\end{sourceblock}
\begin{sourceblock}
\includegraphics[page=480,pagebox=mediabox,trim={45.00bp 124.13bp 44.37bp 488.53bp},clip,width=0.92\textwidth,max height=0.38\textheight]{Ishii.pdf}
\end{sourceblock}
where

\begin{sourceblock}
\includegraphics[page=481,pagebox=mediabox,trim={45.00bp 553.85bp 44.37bp 57.24bp},clip,width=0.92\textwidth,max height=0.38\textheight]{Ishii.pdf}
\end{sourceblock}
应在每个计算单元中获得该分布参数。

\begin{sourceblock}
\includegraphics[page=481,pagebox=mediabox,trim={45.00bp 483.80bp 44.37bp 151.95bp},clip,width=0.92\textwidth,max height=0.38\textheight]{Ishii.pdf}
\end{sourceblock}
在时间-体积双流体模型中应考虑V rz v。在大多数两相流系统中，由于相分布而产生的滑移远大于相间的局部滑移。因此，忽略上述效应会导致气泡、段塞和搅动湍流流态中的空隙率和速度预测出现较大误差。

\subsection*{16.4.3 界面剪切项}\phantomsection\addcontentsline{toc}{subsection}{16.4.3 界面剪切项}
\begin{sourceblock}
\includegraphics[page=481,pagebox=mediabox,trim={45.00bp 346.24bp 44.37bp 286.43bp},clip,width=0.92\textwidth,max height=0.38\textheight]{Ishii.pdf}
\end{sourceblock}
k ki α −∇ ⋅ T | ¢ \textasciitilde{} £ ¡ ⁄ .在分离流中，第二项占主导地位。

例如，对于杆束几何形状中的环形流，z 分量为

\begin{sourceblock}
\includegraphics[page=481,pagebox=mediabox,trim={45.00bp 308.26bp 44.37bp 342.93bp},clip,width=0.92\textwidth,max height=0.38\textheight]{Ishii.pdf}
\end{sourceblock}
k ki α −∇ ⋅ T | ¢ \textasciitilde{} £ ¡ ⁄ 占主导地位，可以证明

\begin{sourceblock}
\includegraphics[page=481,pagebox=mediabox,trim={45.00bp 206.62bp 44.37bp 383.82bp},clip,width=0.92\textwidth,max height=0.38\textheight]{Ishii.pdf}
\end{sourceblock}
当控制体积内的实体结构的构型固定时，

\begin{sourceblock}
\includegraphics[page=481,pagebox=mediabox,trim={45.00bp 172.24bp 44.37bp 478.95bp},clip,width=0.92\textwidth,max height=0.38\textheight]{Ishii.pdf}
\end{sourceblock}
\begin{sourceblock}
\includegraphics[page=481,pagebox=mediabox,trim={45.00bp 147.22bp 44.37bp 495.47bp},clip,width=0.92\textwidth,max height=0.38\textheight]{Ishii.pdf}
\end{sourceblock}
k ki α −∇ ⋅ T | ¢ \textasciitilde{} £ ¡ ⁄ 在气泡流、段状流和搅动湍流中
纳入

对于相间界面动量耦合的正确建模非常重要。为了获得该界面剪切项的本构关系，需要进行一些假设，因为它需要有关空隙和剪切应力分布的信息。假设空隙率和界面剪切应力呈幂律分布，我们可以得出近似值

\begin{sourceblock}
\includegraphics[page=481,pagebox=mediabox,trim={45.00bp 57.22bp 44.37bp 593.98bp},clip,width=0.92\textwidth,max height=0.38\textheight]{Ishii.pdf}
\end{sourceblock}
form of control volume (Ishii and Mishima, 1984).

\subsection*{16.4.4 表面平均量和体积平均量之间的关系}\phantomsection\addcontentsline{toc}{subsection}{16.4.4 表面平均量和体积平均量之间的关系}
为了使用时间-体积平均双流体模型，我们需要确定表面和体积平均量之间的关系，

\begin{sourceblock}
\includegraphics[page=482,pagebox=mediabox,trim={45.00bp 462.64bp 44.37bp 148.16bp},clip,width=0.92\textwidth,max height=0.38\textheight]{Ishii.pdf}
\end{sourceblock}
关系如下： (1) 与特征配置长度相比，控制量的尺寸较大导致

\begin{sourceblock}
\includegraphics[page=482,pagebox=mediabox,trim={45.00bp 374.05bp 44.37bp 261.68bp},clip,width=0.92\textwidth,max height=0.38\textheight]{Ishii.pdf}
\end{sourceblock}
这个假设对于分散流来说是合理的。

(2) 如果控制体的大小与字段的特征配置长度相当，则结果如下

\begin{sourceblock}
\includegraphics[page=482,pagebox=mediabox,trim={45.00bp 262.39bp 44.37bp 373.34bp},clip,width=0.92\textwidth,max height=0.38\textheight]{Ishii.pdf}
\end{sourceblock}
不幸的是，Kolev (2002) 没有给出如何识别表面平均量和体积平均量之间关系的明确策略。这里我们有三个选项来识别它。

(i) 表面平均量由两个相邻单元之间的体积平均量的平均值确定。 (ii) 表面平均数量被确定为与逆风单元中的体积平均数量相同。 (iii)引入加权函数来计算表面上的量。表面平均量是通过加权函数估计的表面量来计算的。由于表面平均量和体积平均量之间的关系可能会影响数值稳定性，因此我们需要仔细考虑这种关系。

\section*{16.5 附录}\phantomsection\addcontentsline{toc}{section}{16.5 附录}
我们考虑笛卡尔坐标系的特殊情况。局部表面上 k R 的平均值定义为

\begin{sourceblock}
\includegraphics[page=483,pagebox=mediabox,trim={45.00bp 500.20bp 44.37bp 125.16bp},clip,width=0.92\textwidth,max height=0.38\textheight]{Ishii.pdf}
\end{sourceblock}
\begin{sourceblock}
\includegraphics[page=483,pagebox=mediabox,trim={45.00bp 436.96bp 44.37bp 184.27bp},clip,width=0.92\textwidth,max height=0.38\textheight]{Ishii.pdf}
\end{sourceblock}
\begin{sourceblock}
\includegraphics[page=483,pagebox=mediabox,trim={45.00bp 313.02bp 44.37bp 223.54bp},clip,width=0.92\textwidth,max height=0.38\textheight]{Ishii.pdf}
\end{sourceblock}
\begin{sourceblock}
\includegraphics[page=483,pagebox=mediabox,trim={45.00bp 241.18bp 44.37bp 369.03bp},clip,width=0.92\textwidth,max height=0.38\textheight]{Ishii.pdf}
\end{sourceblock}
定向表面孔隙率定义为

\begin{sourceblock}
\includegraphics[page=483,pagebox=mediabox,trim={45.00bp 167.20bp 44.37bp 457.12bp},clip,width=0.92\textwidth,max height=0.38\textheight]{Ishii.pdf}
\end{sourceblock}
然后我们有

\begin{sourceblock}
\includegraphics[page=483,pagebox=mediabox,trim={45.00bp 104.65bp 44.37bp 531.06bp},clip,width=0.92\textwidth,max height=0.38\textheight]{Ishii.pdf}
\end{sourceblock}
控制体积边界自由表面上的 k R 表面积分为

\begin{sourceblock}
\includegraphics[page=484,pagebox=mediabox,trim={45.00bp 452.50bp 44.37bp 68.60bp},clip,width=0.92\textwidth,max height=0.38\textheight]{Ishii.pdf}
\end{sourceblock}
where

\begin{sourceblock}
\includegraphics[page=484,pagebox=mediabox,trim={45.00bp 326.95bp 44.37bp 247.19bp},clip,width=0.92\textwidth,max height=0.38\textheight]{Ishii.pdf}
\end{sourceblock}
\begin{sourceblock}
\includegraphics[page=484,pagebox=mediabox,trim={45.00bp 137.78bp 44.37bp 361.59bp},clip,width=0.92\textwidth,max height=0.38\textheight]{Ishii.pdf}
\end{sourceblock}
\chapter{过冷沸腾流中的一维界面面积输运方程}
界面区域传输方程的一维形式可以通过对界面区域传输方程的三维形式应用横截面积平均来获得。然而，面积平均汇和源项的精确数学表达式将涉及许多协方差，这可能会使一维问题进一步复杂化。由于这些局部项最初是从混合物的有限体积元中获得的（Wu et al., 1998；Hibiki and Ishii, 2000a），如果将流路的水力直径视为有限元的长度尺度，则可以假设面积平均源项和汇项对平均参数的函数依赖性大致相同（Wu et al., 1998）。因此，假设具有在横截面积内平均的参数的三维汇和源项仍然适用于界面面积输运方程的一维形式中的面积平均汇和源项。该假设对于流道上相对均匀的局部流动参数（如绝热垂直流中的那些）是有效的。在此假设下，通过对气泡合并和破裂导致的界面面积浓度的汇项和源项进行建模，成功地建立了一维绝热流的界面区域输运方程。

这种方法可能不适用于过冷沸腾流。由于过冷沸腾流中的相分布模式在流道上不均匀，因此应考虑相分布模式的不均匀性来重新表述界面面积输运方程的一维形式。过冷沸腾流可以表征为两个独特的流动区域，即（1）沸腾两相（气泡
在本章中，我们使用气泡层区域和液体单相区域，利用两区域概念开发了过冷沸腾流中一维界面面积输运方程的近似公式，以避免一维公式中的许多协方差（Hibiki et al., 2003a）。然后我们演示了使用一维漂移通量模型中的分布参数来预测气泡层厚度的方法。

\section*{17.1 过冷沸腾流动界面面积输运方程的建立}\phantomsection\addcontentsline{toc}{section}{17.1 过冷沸腾流动界面面积输运方程的建立}
在第10章中，由玻尔兹曼输运方程得到了界面面积输运方程的三维形式：

\begin{sourceblock}
\includegraphics[page=486,pagebox=mediabox,trim={45.00bp 260.07bp 44.37bp 318.47bp},clip,width=0.92\textwidth,max height=0.38\textheight]{Ishii.pdf}
\end{sourceblock}
其中左侧代表界面面积浓度的变化率和对流率。右侧的每一项分别表示由于压力变化、各种颗粒相互作用和相变引起的颗粒体积变化导致的界面面积浓度的变化率。临界气泡尺寸bc D 应根据给定的成核过程确定；即，本体沸腾或冷凝过程的临界空腔尺寸，以及壁成核的气泡离开尺寸。对于大多数两相流，壁成核是主要机制。

一组界面面积输运方程的最简单形式是通过对方程（17-1）应用横截面积平均而获得的一维公式。然而，面积平均源项和汇项的精确数学表达式涉及许多协方差，这使得一维问题进一步复杂化。对于绝热垂直流，相分布模式可以认为是相对均匀的，因此可以忽略协方差(Wu et al., 1998)。因此，具有在横截面积内平均的参数的三维汇和源项仍然适用于界面面积输运方程的一维形式中的面积平均汇和源项。然而，对于过冷沸腾流，相分布模式可能不被假设为均匀，导致一维界面面积输运方程中存在许多协方差。为了避免协方差，引入以下简单模型来制定过冷沸腾流的一维界面面积输运方程。

对于过冷沸腾流，气泡主要存在于受热壁附近，而远离受热壁几乎不存在气泡。因此，流路可以分为两个区域，即（i）沸腾两相（气泡层）区域，其中空隙率分布可以假设是均匀的，以及（ii）液体单相区域，其中空隙率可以假设为零。因此，在气泡层区域上平均的一组界面面积输运方程可得为

\begin{sourceblock}
\includegraphics[page=487,pagebox=mediabox,trim={45.00bp 240.57bp 44.37bp 303.16bp},clip,width=0.92\textwidth,max height=0.38\textheight]{Ishii.pdf}
\end{sourceblock}
其中 B 是气泡层区域的平均数量。横截面积平均量可以作为气泡层区域上的平均量与B C A A 的乘积给出，其中B A 和C A 分别是气泡层区域的面积和横截面积。

这里需要注意的是，具有在气泡层区域内平均的参数的三维源和汇项仍然可以适用于式（17-2）中在气泡层区域内平均的源和汇项，理由如下：由于三维界面面积输运方程中的这些局部项最初是从混合物的有限体积单元获得的，所以在气泡层区域上平均的源和汇项对在气泡层区域上平均的参数的函数依赖关系应近似为

\begin{sourceblock}
\includegraphics[page=488,pagebox=mediabox,trim={50.00bp 146.52bp 44.37bp 55.00bp},clip,width=\textwidth,max height=0.72\textheight]{Ishii.pdf}
\par\vspace{0.45\baselineskip}
\small 图 17-1。通过气泡层区域中的均匀空隙轮廓建模的过冷沸腾流
\end{sourceblock}
（响等人，2003a）

如果将气泡层厚度视为有限元的长度尺度，则同样如此（Wu et al., 1998）。

\section*{17.2 气泡层厚度模型的建立}\phantomsection\addcontentsline{toc}{section}{17.2 气泡层厚度模型的建立}
如上一节所述，预计在过冷沸腾流中会出现加热壁附近的空峰。在过冷沸腾流中，可能无法假设流道上相对均匀的相分布，导致一维界面面积输运方程中存在许多协方差。为了避免协方差，引入图17-1所示的气泡层模型来制定一维界面面积输运方程。这里以内部加热的环形空间为例。在此模型中，流路被分为两个区域，即 (i) 沸腾两相（气泡层）区域，其中空隙率分布被假定为均匀，以及 (ii) 液体单相区域，其中空隙率被假定为零。图17-1中，gα，x，
0 R 、WP α 、WP x 和 R 分别是局部空隙率、从加热棒表面中心测量的径向坐标、加热棒的半径、假设方形空隙峰处的空隙率、气泡层厚度和外圆管的半径。接下来，将导出内部加热的环形空间中的气泡层厚度。

流道中混合物体积通量 j 的分布近似为

\begin{sourceblock}
\includegraphics[page=489,pagebox=mediabox,trim={45.00bp 261.59bp 44.37bp 353.90bp},clip,width=0.92\textwidth,max height=0.38\textheight]{Ishii.pdf}
\end{sourceblock}
其中n和r分别是从加热棒表面测量的指数和径向坐标，并且是横截面平均量。如图17-1所示，出于气泡层模型的目的，假设空隙率的分布为靠近壁面（气泡层区域）的方峰，并近似为

\begin{sourceblock}
\includegraphics[page=489,pagebox=mediabox,trim={45.00bp 136.72bp 44.37bp 484.04bp},clip,width=0.92\textwidth,max height=0.38\textheight]{Ishii.pdf}
\end{sourceblock}
分布参数 0 C 可以从混合物体积通量和空隙率分布中获得：

\begin{sourceblock}
\includegraphics[page=489,pagebox=mediabox,trim={45.00bp 51.74bp 44.37bp 573.01bp},clip,width=0.92\textwidth,max height=0.38\textheight]{Ishii.pdf}
\end{sourceblock}
由式(17-3)-(17-5)，解析得到过冷沸腾流的分布参数为

\begin{sourceblock}
\includegraphics[page=490,pagebox=mediabox,trim={45.00bp 210.32bp 44.37bp 93.29bp},clip,width=0.92\textwidth,max height=0.38\textheight]{Ishii.pdf}
\end{sourceblock}
面积平均空隙率可由下式获得

\begin{sourceblock}
\includegraphics[page=490,pagebox=mediabox,trim={45.00bp 138.92bp 44.37bp 485.71bp},clip,width=0.92\textwidth,max height=0.38\textheight]{Ishii.pdf}
\end{sourceblock}
由式(17-6)可以看出，我们可以估算出气泡层的厚度
WP x 前提是给出了分布参数 0 C 和指数 n。指数假设为 7。如第 14 章所述，内热环空中沸腾流分布参数的本构方程推导为（Hibiki 等，2003a；Julia 等，2009）

\begin{sourceblock}
\includegraphics[page=491,pagebox=mediabox,trim={45.00bp 548.43bp 44.37bp 68.48bp},clip,width=0.92\textwidth,max height=0.38\textheight]{Ishii.pdf}
\end{sourceblock}
然后由气泡层厚度WP x 可以得到气泡层面积B A，并且可以得到不带协方差项的一维界面面积输运方程。

第 14 章给出了各种流动几何形状的分布参数的本构方程。

\chapter*{参考文献}\addcontentsline{toc}{chapter}{参考文献}
\markboth{参考文献}{参考文献}
{\small\raggedright\noindent Adorni, N., Peterlongo, G., Ravetta, R. and Tacconi, F. A., 1964, Large Scale Experiments on\par}
{\small\raggedright\noindent Heat Transfer and Hydrodynamics with Steam-water Mixtures, CISE Report R-91, Italy. Akiyama, M. and Aritomi, M., 2002, Advanced Numerical Analysis of Two-phase Flow\par}
{\small\raggedright\noindent Dynamics –Multi-dimensional Flow Analysis- , Corona Publishing Co. Ltd., Tokyo, Japan Alia, P., Cravarolo, L., Hassid, A. and Pedrocchi, E., 1965, Liquid Volume Fraction in\par}
{\small\raggedright\noindent Adiabatic Two-phase Vertical Upflow-round Conduit, CISE Report -105, Italy. Antal, S. P., Lahey Jr, R. T. and Flaherty, J. E., 1991, Analysis of Phase Distribution in Fully\par}
{\small\raggedright\noindent Developed Laminar Bubbly Two-phase Flow, Int. J. Multiphase Flow 17: 635-652. Aris, R., 1962, Vectors, Tensors and the Basic Equations of Fluid Mechanics , Prentice-Hall,\par}
{\small\raggedright\noindent Englewood Cliffs, N.J. Arnold, G. S., Drew, D. A. and Lahey Jr., R. T., 1989, Derivation of Constitutive Equations\par}
{\small\raggedright\noindent for Interfacial Force and Reynolds Stress for a Suspension of Spheres Using Ensemble Cell Averaging, Chem. Eng. Comm. 86: 43-54. Auton, T. R., 1987, The Lift Force on a Spherical Body in a Rotational Flow, J. Fluid Mech.\par}
{\small\raggedright\noindent 183: 199-218. Azbel, D., 1981, Two-phase Flows in Chemical Engineering , Cambridge University Press,\par}
{\small\raggedright\noindent Cambridge, UK. Azbel, D. and Athanasios, I. L., 1983, A Mechanism of Liquid Entrainment. Handbook of\par}
{\small\raggedright\noindent Fluids in Motion , Ann Arbor Sci. Pub., Ann Arbor, MI. Baker, J. L. L., 1965, Flow-regime Transitions at Elevated Pressure in Vertical Two-phase\par}
{\small\raggedright\noindent Flow, Argonne National Lab. Report , ANL-7093. Bankoff, S. G., 1960, A Variable Density Single-fluid Model for Two-phase Flow with\par}
{\small\raggedright\noindent Particular Reference to Steam Water Flow, J. Heat Transfer 82: 265-272. Basu, N., Warrier, G. R. and Dhir V. K., 2002, Onset of Nucleate Boiling and Active\par}
{\small\raggedright\noindent Nucleation Site Density during Subcooled Flow Boiling, J. Heat Transfer 124: 717-728. Bello, J. K., 1968, Turbulent Flow in Channel with Parallel Walls, Moskva, Mir, in Russian. Bergonzoli, F. and Halfen, F. J., 1964, Heat Transfer and Void Formation during Forced\par}
{\small\raggedright\noindent Circulation Boiling of Organic Coolants, NAA-SR-8906, Atomics International. Bilicki, A. and Kestin, J., 1987, Transition Criteria for Two-phase Flow Patterns in Vertical\par}
{\small\raggedright\noindent Upward Flow, Int. J. Multiphase Flow 13: 283-294. Bird, R. B., Stewart, W. E. and Lightfoot, E. N., 1960, Transport Phenomena , John Wiley\par}
{\small\raggedright\noindent and Sons, Inc., New York.\par}
{\small\raggedright\noindent Bornhorst, W. J. and Hatsopoulos, G. N., 1967, Analysis of a Liquid Vapor Phase Change by\par}
{\small\raggedright\noindent the Methods of Irreversible Thermodynamics, J. Applied Mech. 89: 847-853. Boure, J. and Réocreux, M., 1972, General Equations of Two-phase Flows: Application to\par}
{\small\raggedright\noindent Critical Flows and to Non Steady Flows, 4th All Union Heat and Mass Transfer Conference , Minsk. Boure, J., 1973, Dynamique des Écoulements Diphasiques: Propagation des Petites\par}
{\small\raggedright\noindent Perturbations, CEA-R-4456 . Bowen, R. M., 1967, Toward a Thermodynamics and Mechanics of Mixtures, Arch. Rational\par}
{\small\raggedright\noindent Mech. Anal. 24: 370-403. Bridge, A. G., Lapidus, L and Elgin, J. C., 1964, The Mechanics of Vertical Gas-Liquid\par}
{\small\raggedright\noindent Fluidized System I: Counter-current Flow, AIChE J. 10: 819-827. Brinkman, H., 1952, The Viscosity of Concentrated Suspensions and Solutions, J. Chem.\par}
{\small\raggedright\noindent Phys. 20: 571. Brodkey, R. S., 1967, The Phenomena of Fluid Motion , Addison-Wesley. Brodkey, R. S., 1971, Transport Phenomena at the Liquid-Vapor Interface of Mercury Using\par}
{\small\raggedright\noindent a Radioactive Tracer, International Symposium on Two-phase Systems , Haifa. Burgers, J. M., 1941, On the Influence of the Concentration of a Suspension upon the\par}
{\small\raggedright\noindent Sedimentation Velocity (in Particular for a Suspension of Spherical Particles) Proc. K. Med. Akad. Wet. 44: 1045-1051 (1941); 45: 9-16 (1942) Burns, A. D., Frank, T., Hamill, I. and Shi, J. M., 2004, Proc. 5 th Int. Conf. Multiphase Flow ,\par}
{\small\raggedright\noindent Yokohama, Paper No. 392. Buevich, I., 1969, A Hydrodynamic Model of Disperse Systems, J. Applied Math. Mech. 33:\par}
{\small\raggedright\noindent 466-479. Buyevich, Y., 1972, Statistical Hydrodynamics of Disperse Systems. Part.I. Physical\par}
{\small\raggedright\noindent Background and General Equations, J. Fluid Mech. 49: 489-507. Carrier, G. F., 1958, Shock Waves in Dusty Gas, J. Fluid Mech. 4: 376-382. Callen, H. B., 1960, Thermodynamics , Wiley. Chao, B., 1962, Motion of Spherical Gas Bubbles in a Viscous Liquid at Large Reynolds\par}
{\small\raggedright\noindent Numbers, Phys. Fluids 5: 69-79. Chen, J. C., 1966, Correlation of Boiling Heat Transfer to Saturated Fluids in Convective\par}
{\small\raggedright\noindent Flow, Ind. Eng. Chem. Proc. Des. Devel. 5: 322-329. Chen, S. W., Liu, Y., Hibiki, T. and Ishii, M., 2010, Drift-flux Model in Pool Bundle System,\par}
{\small\raggedright\noindent Technical Report , School of Nuclear Engineering, Purdue University, IN, USA, PU/NE- 10-14. Clift, R., Grace, J. R. and Weber, M. E., 1978, Bubbles, Drops, and Particles , Academic\par}
{\small\raggedright\noindent Press, New York. Coleman, B. and Noll, W., 1960, An Approximate Theorem for Functionals with\par}
{\small\raggedright\noindent Applications in Continuum Mechanics, Arch. Rational Mech. Anal. 6: 355-370. Coleman, B. D., 1964, Thermodynamics of Materials with Memory, Arch. Rational Mech.\par}
{\small\raggedright\noindent Anal. 17: 1-46. Collier, J., 1972, Convective Boiling and Condensation , McGraw Hill, London. Coulaloglou, C. A. and Tavlarides, L. L., 1976, Drop Size Distributions and Coalescence\par}
{\small\raggedright\noindent Frequencies of Liquid-liquid Dispersion in Flow Vessels, AIChE J. 22: 289-297. Coulaloglou, C. A. and Tavlarides, L. L., 1977, Description of Interaction Processes in\par}
{\small\raggedright\noindent Agitated Liquid-liquid Dispersions, Chem. Eng. Sci. 32: 1289-1297. Cravarolo, L., Giorgini, A., Hassid, A. and Pedrocchi, E., 1964, A Device for the\par}
{\small\raggedright\noindent Measurement of Shear Stress on the Wall of a Conduit; Its Application in the Mean Density Determination in Two-phase Flow; Shear Stress Data in Two-phase Adiabatic Vertical Flow, CISE Report -82, Italy.\par}
{\small\raggedright\noindent Culick, F., 1964, Boltzman Equation Applied to a Problem of Two-phase Flow, Phys. Fluid\par}
{\small\raggedright\noindent 7: 1898-1904. De Groot, S. B. and Mazur, P., 1962, Non-equilibrium Thermodynamics , North Holland. De Jarlais, G., Ishii and M., Linehan, J., 1986, Hydrodynamic Stability of Inverted Annular\par}
{\small\raggedright\noindent Flow in an Adiabatic Simulation, J. Heat Transfer 108 : 84-91. Delhaye, J. M., 1968, Equations of Fondamentales des Écoulements Diphasiques, Part 1 and 2,\par}
{\small\raggedright\noindent CEA-R-3429 . Delhaye, J. M., 1969, General Equations of Two-phase Systems and their Application to Air-\par}
{\small\raggedright\noindent water Bubble Flow and to Steam-water Flashing Flow, ASME Paper 69-HT-63, 11th Heat Transfer Conference , Minneapolis. Delhaye, J. M., 1970, Contrubution à L’étude des Écoulements Diphasiques Eau-air et Eau-\par}
{\small\raggedright\noindent vapeur, Ph.D. Thesis , University of Grenoble. Delhaye, J. M., 1974, Jump Conditions and Entropy Sources in Two-phase Systems. Local\par}
{\small\raggedright\noindent Instant Formulation, Int. J. Multiphase Flow 1: 395-409. Delhaye, J., Giot, M. and Riethmuller, M. L., 1981, Thermohydraulics of Two-phase Systems\par}
{\small\raggedright\noindent for Industrial Design and Nuclear Engineering , Hemisphere Publishing Corp., New York. Dinh, T. N., Li, G. J. and Theofanous, T. G., 2003, An Investigation of Droplet Breakup in\par}
{\small\raggedright\noindent High Mach, Low Weber Number Regime, Proc. 41 st Aerospace Sci. Mtg and Exh. , Paper AIAA 2003-317, Reno, Nevada. Diunin, A. K., 1963, On the Mechanics of Snow Storms, Siberian Branch, Akademii Nauk\par}
{\small\raggedright\noindent SSSR , Novosibirsk. Drew, D. A., 1971, Averaged Field Equations for Two-phase Media, Studies Appl. Math. 1:\par}
{\small\raggedright\noindent 133-166. Drew, A. and Passman, S., 1998, Theory of Multi-component Fluids , Springer Verlag Inc.,\par}
{\small\raggedright\noindent New York. Dumitrescu, D. T., 1943, Stomung an einer Luftblase in Senkrechten Rohr, Z. Angew, Math.\par}
{\small\raggedright\noindent Mech. 23: 139-149. Eilers, H., 1941, The Viscosity of the Emulsion of Highly Viscous Substances as Function of\par}
{\small\raggedright\noindent Cencentration, Kolloid Z 97: 313-321. Ervin, E. A. and Tryggvason, G., 1997, The Rise of Bubbles in a Vertical Shear Flow, J.\par}
{\small\raggedright\noindent Fluids Eng. 119: 443-449. Euh, D. J., Ozar, B., Hibiki, T., Ishii, M. and Song, S. H., 2010, Characteristics of Bubble\par}
{\small\raggedright\noindent Departure Frequency in a Low-Pressure Subcooled Boiling Flow, J. Nucl. Sci. Technol. 47: 608-617. Fauske, H., 1962, Critical Two-phase, Steam Water Flow, Proc. Heat Transfer and Fluid\par}
{\small\raggedright\noindent Mechanics Institute , Stanford Univ. Press, pp.78. Fick, A., 1855, Über Diffusion, Ann. Der Phys. 94: 59-86. Frank, T., Zwart, P. J., Krepper, E., Prasser, H. M. and Lucas, D., 2008, Validation of CFD\par}
{\small\raggedright\noindent Models for Monoand Polydisperse Air-water Two-phase Flows in Pipes, Nucl. Eng. Des. 238: 647-659. Frankl, F. I., 1953, On the Theory of Motion of Sediment Suspensions, Soviet Physics\par}
{\small\raggedright\noindent Doklady, Academii Nauk SSSR , 92: 247-250. Frankel, N. A. and Acrivos, A., 1967, On the Viscosity of a Concentrated Suspension of Solid\par}
{\small\raggedright\noindent Spheres, Chem. Eng. Sci. 22: 847-853. Friedlander, S. K., 1977, Smoke, Dust and Haze , Wiley, New York. Fu, X. Y. and Ishii, M., 2002a, Two-group Interfacial Area Transport in Vertical Air-water\par}
{\small\raggedright\noindent Flow I. Mechanistic Model, Nucl. Eng. Des. 219: 143-168. Fu, X. Y. and Ishii, M., 2002b, Two-group Interfacial Area Transport in Vertical Air-water\par}
{\small\raggedright\noindent Flow II. Model Evaluation, Nucl. Eng. Des. 219: 169-190.\par}
{\small\raggedright\noindent Goda, H., Hibiki, T. Kim, S., Ishii, M. and Uhle, J., 2003, Drift-flux Model for Downward\par}
{\small\raggedright\noindent Two-phase Flow, Int. J. Heat Mass Transfer 46: 4835-4844. Gibbs, J. W., 1948, Collected Work of J. W. Gibbs , Yale University Press, Vol.1, New York. Goldstein, S., 1938, Modern Developments in Fluid Dynamics , Oxford University Press,\par}
{\small\raggedright\noindent London. Govier G. W. and Aziz, K, 1972, The Flow of Complex Mixtures in Pipes , Van Nostrant-\par}
{\small\raggedright\noindent Reinhold Co., New York. Grace, J. R., Wairegi, T. and Brophy, J., 1978, Break-up of Drops and Bubbles in Stagnant\par}
{\small\raggedright\noindent Media, Can. J. Chem. Eng. 556: 3-8. Hadamard, J., 1911, Mouvement Permanent Lent d’une Sphere Liquide Visqueuse dans un\par}
{\small\raggedright\noindent Liquid Visqueux, C. R. Acad. Sci. Paris Sér A-B 152: 1735-1739. Hancox, W. T., and Nicoll, W. B., 1972, Prediction of Time-dependent Diabatic Two-phase\par}
{\small\raggedright\noindent Water Flows, Prog. Heat Mass Transfer 6: 119-135. Happel, J. and Brfnner, H., 1965, Low Reynolds Number Hydrodynamics , Prentice-Hall. Harmathy, T. Z., 1960, Velocity of Large Drops and Bubbles in Media of Infinite and\par}
{\small\raggedright\noindent Restricted Extent, AIChE J. 6: 281-288. Hawksley, P. G. W., 1951, The Effect of Concentration on the Settling of Suspensions and\par}
{\small\raggedright\noindent Flow through Porous Media, Some Aspect of Fluid Flow , pp.114, Edward Arnold, London. Hayes, W. D., 1970, Kinematic Wave Theory, Proc. Royal Soc. London Ser. A Math. Phys.\par}
{\small\raggedright\noindent Sci. 320: 209-226. Helmhotz, H., 1868, Uber Discontinuirliche Flüssigkeitsbewegungen, Monatsber. Dtsch.\par}
{\small\raggedright\noindent Akad. Wiss. Berlin pp.215-228. Hewitt, G. and Hall Taylor, N. S., 1970, Annular Two-phase Flow , Pergamon Press, Oxford. Hibiki, T. and Ishii, M., 1999, Experimental Study on Interfacial Area Transport in Bubbly\par}
{\small\raggedright\noindent Two-phase Flows, Int. J. Heat Mass Transfer 42: 3019-3035. Hibiki, T. and Ishii, M., 2000a, One-group Interfacial Area Transport of Bubbly Flows in\par}
{\small\raggedright\noindent Vertical Round Tubes, Int. J. Heat Mass Transfer 43: 2711-2726. Hibiki, T. and Ishii, M., 2000b, Two-group Interfacial Area Transport Equations at Bubbly-\par}
{\small\raggedright\noindent to-slug Flow Transition, Nucl. Eng. Des. 202: 39-76. Hibiki, T., Ishii, M. and Xiao, Z., 2001a, Axial Interfacial Area Transport of Vertical Bubbly\par}
{\small\raggedright\noindent Flows, Int. J. Heat Mass Transfer 44: 1869-1888. Hibiki, T., Takamasa, T. and Ishii, M., 2001b, Interfacial Area Transport of Bubbly Flow in a\par}
{\small\raggedright\noindent Small Diameter Pipe, J. Nucl. Sci. Technol. 38: 614-620. Hibiki, T. and Ishii, M., 2002a, Interfacial Area Concentration of Bubbly Flow Systems,\par}
{\small\raggedright\noindent Chem. Eng. Sci. 57: 3967-3977. Hibiki, T. and Ishii, M., 2002b, Distribution Parameter and Drift Velocity of Drift-flux Model\par}
{\small\raggedright\noindent in Bubbly Flow, Int. J. Heat Mass Transfer 45: 707-721. Hibiki, T. and Ishii, M., 2002c, Development of One-group Interfacial Area Transport\par}
{\small\raggedright\noindent Equation in Bubbly Flow Systems, Int. J. Heat Mass Transfer 45: 2351-2372. Hibiki, T. and Ishii, M., 2003a, One-dimensional Drift-flux Model for Two-phase Flow in a\par}
{\small\raggedright\noindent Large Diameter Pipe, Int. J. Heat Mass Transfer 46: 1773-1790. Hibiki, T. and Ishii, M., 2003b, Active Nucleation Site Density in Boiling Systems, Int. J.\par}
{\small\raggedright\noindent Heat Mass Transfer 46: 2587-2601. Hibiki, T. and Ishii, M., 2003c, One-dimensional Drift-flux Model and Constitutive Equations\par}
{\small\raggedright\noindent for Relative Motion between Phases in Various Two-phase Flow Regimes, Int. J. Heat Mass Transfer 46: 4935-4948; Erratum, 48: 1222-1223 (2005). Hibiki, T. and Ishii, M., 2007, Lift Force in Bubbly Flow Systems, Chem. Eng. Sci. 62: 6457-\par}
{\small\raggedright\noindent 6474 (2007). Hibiki, T., Situ, R., Mi, Y. and Ishii, M., 2003a, Modeling of Bubble-layer Thickness for\par}
{\small\raggedright\noindent Formulation of One-dimensional Interfacial Area Transport Equation in Subcooled\par}
{\small\raggedright\noindent Boiling Two-phase Flow, Int. J. Heat Mass Transfer 46: 1409-1423; Erratum, 46: 3549- 3550 (2003). Hibiki, T., Situ, R., Mi, Y. and Ishii, M., 2003b, Local Flow Measurements of Vertical\par}
{\small\raggedright\noindent Upward Bubbly Flow in an Annulus, Int. J. Heat Mass Transfer 46: 1479-1496. Hibiki, T., Takamasa, T., Ishii, M. and Gabriel, K. S., 2006, One-dimensional Drift-flux\par}
{\small\raggedright\noindent Model at Microgravity Conditions, AIAA J. 44: 1635-1642. Hirschfelder, J. V., Curtiss, C. F. and Bird, R. B., 1954, Molecular Theory of Gases and\par}
{\small\raggedright\noindent Liquids , John Wiley and Sons, Inc., New York. Hinze, J. W., 1959, Turbulence , McGraw-Hill, New York. Hill, D., 1998, The Computer Simulation of Dispersed Two-phase Flows, Ph. D. Thesis ,\par}
{\small\raggedright\noindent Imperial College of Science Technology and Medicine, London. Hosokawa, S. and Tomiyama, A., 2003, Lateral Force Acting on a Deformed Single Bubble\par}
{\small\raggedright\noindent due to the Presence of Wall, Trans. JSME 69: 2214-2220. Hosokawa, S. and Tomiyama, A., 2009, Multi-fluid Simulation of Turbulent Bubbly Pipe\par}
{\small\raggedright\noindent Flows, Chem. Eng. Sci. 64: 5308-5318. Hughes, T. A., 1958, Steam-water Mixture Density Studies in a Natural Circulation High\par}
{\small\raggedright\noindent Pressure System, Babcock and Wilcox, G. Report No. 5435. Hughes, D., 1976, The Field Balance Equations for Two-phase Flows in Porous Media, Proc.\par}
{\small\raggedright\noindent Two-phase Flow and Heat Transfer Symp.-Workshop , vol.1, FL, USA, Two-phase Transport and Reactor Safety. Ishii, M. and Zuber, N., 1970, Thermally Induced Flow Instabilities in Two-phase Mixtures,\par}
{\small\raggedright\noindent Proc. 4th International Heat Transfer Conference , Paris. Ishii, M., 1971, Thermally Induced Flow Instabilities in Two-phase Mixture in Thermal\par}
{\small\raggedright\noindent Equilibrium, Ph.D. Thesis , Georgia Institute of Technology. Ishii, M., 1975, Thermo-fluid Dynamic Theory of Two-phase Flow , Collection de la Direction\par}
{\small\raggedright\noindent des Etudes et Researches d’Electricite de France, Eyrolles, Paris, France, 22. Ishii, M. and Grolmes, M. A., 1975, Inception Criteria for Droplet Entrainment in Two-phase\par}
{\small\raggedright\noindent Concurrent Film Flow, AIChE J. 21: 308-318. Ishii, M., Jones, O. C. and Zuber, N., 1975, Thermal Non-equilibrium Effects in Drift Flux\par}
{\small\raggedright\noindent Model of Two-phase Flow, Trans. Am. Nucl. Soc. 22: 263-264. Ishii, M., 1976, One-dimensional Drift-flux Modeling: One-dimensional Drift Velocity of\par}
{\small\raggedright\noindent Dispersed Flow in Confined Channel, Argonne National Lab. Report , ANL-76-49. Ishii, M., Chawla, T. C. and Zuber, N., 1976, Constitutive Equation for Vapor Drift Velocity\par}
{\small\raggedright\noindent in Two-phase Annular Flow, AIChE J. 22: 283-289. Ishii, M. and Chawla, T. C., 1979, Local Drag Laws in Dispersed Two-phase Flow, Argonne\par}
{\small\raggedright\noindent National Lab. Report , ANL-79-105. Ishii, M., 1977, One-dimensional Drift-flux Model and Constitutive Equations for Relative\par}
{\small\raggedright\noindent Motion between Phases in Various Two-phase Flow Regimes, Argonne National Lab. Report , ANL-77-47. Ishii, M. and Zuber, N., 1979, Drag Coefficient and Relative Velocity in Bubbly, Droplet or\par}
{\small\raggedright\noindent Particulate Flows, AIChE J. 25: 843-855. Ishii, M. and Mishima, K., 1981, Study of Two-fluid Model and Interfacial Area, Argonne\par}
{\small\raggedright\noindent National Lab Report ANL-80-111. Ishii, M. and Mishima, K., 1984, Two-fluid Model and Hydrodynamic Constitutive Relations,\par}
{\small\raggedright\noindent Nucl. Eng. Des. 82: 107-126. Ishii, M. and De Jarlais, G., 1987, Flow Visualization Study of Inverted Annular Flow of Post\par}
{\small\raggedright\noindent Dryout Heat Transfer Region, Nucl. Eng. Des. 99: 187-199. Ishii, M., Kim, S. and Uhle, J., 2002, Interfacial Area Transport Equation: Model\par}
{\small\raggedright\noindent Development and Benchmark Experiments, Int. J. Heat Mass Transfer 45: 3111-3123.\par}
{\small\raggedright\noindent Ishii, M. and Kim, S., 2004, Development of One-group and Two-group Interfacial Area\par}
{\small\raggedright\noindent Transport Equation, Nucl. Sci. Eng. 146: 257-273. Ishii, M. and Hibiki, T., 2006, Thermo-fluid Dynamics of Two-phase Flow , Springer, New\par}
{\small\raggedright\noindent York, USA. Julia, J. E., Hibiki, T., Ishii, M., Yun, B. J. and Park, G. C., 2009, Drift-flux Model in a Sub-\par}
{\small\raggedright\noindent channel of Rod Bundle Geometry, Int. J. Heat Mass Transfer 52: 3032-3041. Kalinin, A. V., 1970, Derivation of Fluid Mechanics Equations for a Two-phase Medium with\par}
{\small\raggedright\noindent Phase Changes, Heat Transfer Soviet Res. 2: 83-96. Kataoka, I. and Ishii, M. 1987, Drift Flux Model for Large Diameter Pipe and New\par}
{\small\raggedright\noindent Correlation for Pool Void Fraction, Int. J. Heat Mass Transfer 30: 1927-1939. Kataoka, I. and Serizawa, A., 1995, Modeling and Prediction of Turbulence in Bubbly Two-\par}
{\small\raggedright\noindent phase Flow, Proc. 2 nd Int. Conf. Multiphase Flow ’95 – Kyoto , pp. MO2-11-MO2-16. Kelly, F. D., 1964, A Reacting Continuum, Int. J. Engng. Sci. 2: 129-153. Kelvin, W., 1871, Hydrokinetic Solutions and Observations, London, Edinburgh and Dublin\par}
{\small\raggedright\noindent Philosophical Magazine and Journal of Science Ser.4 42: 362-377. Kim, W. K. and Lee, K. L., 1987, Coalescence Behavior of two Bubbles in Stagnant Liquids,\par}
{\small\raggedright\noindent J. Chem. Eng. Jpn. 20: 449-453. Kirkpatrick, R. D. and Lockett, M. J., 1974, The Influence of Approach Velocity on Bubble\par}
{\small\raggedright\noindent Coalescence, Chem. Eng. Sci. 29: 2363-2373. Klausner, J. F., Mei, R., Bernhard, D. M. and Zeng, L. Z., 1993, Vapor Bubble Departure in\par}
{\small\raggedright\noindent Forced Convection Boiling, Int. J. Heat Mass Transfer 36: 651-662. Kocamustafaogullari, G., 1971, Thermo-fluid Dynamics of Separated Two-phase Flow, Ph.D.\par}
{\small\raggedright\noindent Thesis , Georgia Institute of Technology. Kocamustafaogullari, G. and Ishii, M., 1983, Interfacial Area and Nucleation Site Density in\par}
{\small\raggedright\noindent Boiling Systems, Int. J. Heat Mass Transfer 26: 1377-1387. Kocamustafaogullari, G., Chen, I. Y. and Ishii, M., 1984, Unified Theory for Predicting\par}
{\small\raggedright\noindent Maximum Fluid Particle Size for Drops and Bubbles, Argonne National Lab. Report , ANL-84-67. Kocamustafaogullari, G. and Ishii, M., 1995, Foundation of the Interfacial Area Transport\par}
{\small\raggedright\noindent Equation and its Closure Relations, Int. J. Heat Mass Transfer 38: 481-493. Kolev, N., 2002, Multiphase Flow Dynamics 1 Fundamentals, 2: Mechanical and Thermal\par}
{\small\raggedright\noindent Interactions, Springer-Verlag. Kordyban, E., 1977, Some Characteristics of High Waves in Closed Channels Approaching\par}
{\small\raggedright\noindent Kelvin-Helmholtz Instability, J. Fluids Eng. 99: 389-346. Kotchine, N. E., 1926, Sur la Théorie des Ondes De-choc dans un Fluide, Bend. Circ. Mat.\par}
{\small\raggedright\noindent Palermo 50: 305-344. Kutateladze, S. S., 1952, Heat Transfer in Condensation and Boiling, Moscow, AEC-TR-3770 ,\par}
{\small\raggedright\noindent USAEC Technical Information Service. Kynch, G. J., 1952, A Theory of Sedimentation, Trans. Faraday Soc. 48: 166-176. Lackme, C., 1973, Two Regimes of a Spray Column in Counter-current Flow, AIChE Symp.\par}
{\small\raggedright\noindent Heat Transfer R. D. 70: 59-63. Lahey Jr., R. T., Cheng, L. Y., Drew, D. A. and Flaherty, J. E., 1978, The Effect of Virtual\par}
{\small\raggedright\noindent Mass on the Numerical Stability of Accelerating Two-phase Flows, AIChE 71st Annual Meeting , Miami Beach, Florida. Lahey, Jr., R. T. and Drew, D. A., 1988, The Three-dimensional Time and Volume Averaged\par}
{\small\raggedright\noindent Conservation Equations of Two-phase Flow, Advances in Nucl. Sci. Technol. 20: Plenum. Lahey Jr, R. T., Lopez de Bertodano, M. and Jones Jr., O. C., 1993, Phase Distribution in\par}
{\small\raggedright\noindent Complex Geometry Conduits, Nucl. Eng. Des. 141: 117-201. Lamb, H., 1945, Hydrodynamics , Dover, New York. Landau, L. D., 1941, Theory of Super Fluidity of Helium II, Physical Review 60: 166-176.\par}
{\small\raggedright\noindent Landel, R. F., Moser, B. G. and Bauman, A. J., 1965, Rheology of Concentrated Suspensions:\par}
{\small\raggedright\noindent Effects of a Surfactant, Proc. 4th Int. Congress on Rheology , Brown University, Part 2, pp.663. Lee, S. Y., Hibiki, T. and Ishii, M., 2009, Formulation of Time and Volume Averaged Two-\par}
{\small\raggedright\noindent fluid Model Considering Structural Materials in a Control Volume, Nucl. Eng. Des. 239: 127-139. Legendre, D. and Magnaudet, J., 1998, The Lift Force on a Spherical Bubble in a Viscous\par}
{\small\raggedright\noindent Linear Shear Flow, J. Fluid Mech. 368 : 81-126. Letan, R. and Kehat, E., 1967, Mechanics of a Spray Column, AIChE J. 13: 443-449. Levich, V. G., 1962, Physicochemical Hydrodynamics , Prentice-Hall. Levy, S., 1960, Steam Slip-theoretical Prediction from Momentum Model, J. Heat Transfer\par}
{\small\raggedright\noindent 82: 113-124. Lighthill, M. J. and Whitham, G. B., 1955, On the Kinematic Waves I. Flood Movement in\par}
{\small\raggedright\noindent Long Rivers, Proc. Royal Soc. London 229 : 281-316. Liu, T. J., 1993, Bubble Size and Entrance Length Effects on Void Development in a Vertical\par}
{\small\raggedright\noindent Channel, Int. J. Multiphase Flow 19: 99-113. Loeb, L. B., 1927, The Kinetic Theory of Gases , Dover, New York. Lopez de Bertodano, M., Lahey Jr., R. T. and Jones, O. C., 1994, Development of a k - ε\par}
{\small\raggedright\noindent Model for Bubbly Two-phase Flow, J. Fluids Eng. 116: 128-134. Lopez de Bertodano, 1998, Two Fluid Model for Two-phase Turbulent Jets, Nucl. Eng. Des.\par}
{\small\raggedright\noindent 179: 65-74. Lopez de Bertodano, M., Sun, X., Ishii, M. and Ulke, A., 2006, Phase Distribution in the Cap\par}
{\small\raggedright\noindent Bubble Regime in a Duct, J. Fluids Eng. 128: 811-818. Loth, E., Taeibi-Rahni, M. and Tryggvason, G., 1997, Deformable Bubbles in a Free Shear\par}
{\small\raggedright\noindent Layer, Int. J. Multiphase Flow 23: 977-1001. Lumley, J., 1970, Toward a Turbulent Constitutive Relation, J. Fluid Mech. 41: 413-434. Marchaterre, J. F., 1956, The Effect of Pressure on Boiling Density in Multiple Rectangular\par}
{\small\raggedright\noindent Channels, Argonne National Lab. Report , ANL-5522. Martinelli, R. C. and Nelson, D. B., 1948, Prediction of Pressure Drop during Forced\par}
{\small\raggedright\noindent Circulation Boiling of Water, Trans. ASME 70: 695-702. Maurer, G. W., 1956, A Method of Predicting Steady State Boiling Vapour Fractions in\par}
{\small\raggedright\noindent Reactor Coolant Channels, WAPD-BT-19. Maxwell, J., 1867, On the Dynamical Theory of Gases, Phil. Trans. Roy. Soc. London 157:\par}
{\small\raggedright\noindent 49-88. McConnell, A. S., 1957, Application of Tensor Analysis , Dover. Meyer, J. E., 1960, Conservation Laws in One-dimensional Hydrodynamics, Bettis Technical\par}
{\small\raggedright\noindent Review 61: WAPD-BT-20. Miles, J. W., 1957, On the Generation of Surface Waves by Shear Flows, I-IV, J. Fluid Mech.\par}
{\small\raggedright\noindent 3: 185-204 (1957); 6: 568-582 (1959); 6: 583-598 (1959); 13: 433-448 (1962). Mishima, K. and Ishii, M., 1980, Theoretical Prediction of Onset of Slug Flow, J. Fluid Eng.\par}
{\small\raggedright\noindent 102: 441-445. Mishima, K. and Ishii, M., 1984, Flow Regime Transition Criteria for Upward Two-phase\par}
{\small\raggedright\noindent Flow in Vertical Tubes, Int. J. Heat Mass Transfer 27: 723-737. Miller, J. W., 1993, An Experimental Analysis of Large Spherical Cap Bubbles Rising in an\par}
{\small\raggedright\noindent Extended Liquid, M. S. Thesis , Purdue University. Mokeyev, Yu. G., 1977, Effect of Particle Concentration on their Drag and Induced Mass,\par}
{\small\raggedright\noindent Fluid Mech.-Soviet Res. 6: 161-168. Muller, I., 1968, A Thermodynamics Theory of Mixtures of Fluids, Arch. Rational Mech.\par}
{\small\raggedright\noindent Anal. 28: 1-39.\par}
{\small\raggedright\noindent Murray, S. O., 1954, On the Mathematics of Fluidization I, Fundamental Equations and Wave\par}
{\small\raggedright\noindent Propagation, J. Fluid Mech. 21: 465-493. Naciri, M. A., 1992, Contribution à l’étude des forces exercées par un liquide sur une bulle de\par}
{\small\raggedright\noindent gaz: portance, masse ajouté. Ph D Thesis, L’ecole Centrale de Lyon, France. Neal, L. G., 1963, An Analysis of Slip in Gas-liquid Flow Applicable to the Bubble and Slug\par}
{\small\raggedright\noindent Flow Regimes, Kjeller Research Establishment Report, Norway KR-2. Nicklin, D. J., Wilkes, J. O. and Davidson, J. F., 1962, Two-phase Flow in Vertical Tubes,\par}
{\small\raggedright\noindent Trans. Inst. Chem. Eng. 40: 61-68. Nikuradse, J., 1932, Gesetzmäsigkeit der Turbulenten Strömung in Glatten Rohre, Forsch.\par}
{\small\raggedright\noindent Arb. Ing.-Wes. pp.356. Oolman, T. O. and Blanch, H. W., 1986a, Bubble Coalescence in Air-Sparged Bioreactors,\par}
{\small\raggedright\noindent Biotech. Bioeng. 28: 578-584. Oolman, T. O. and Blanch, H. W., 1986b, Bubble Coalescence in Stagnant Liquids, Chem.\par}
{\small\raggedright\noindent Eng. Commun. 43: 237-261. Otake, T., Tone, S., Nakao, K., and Mitsuhashi, Y., 1977, Coalescence and Breakup of\par}
{\small\raggedright\noindent Bubbles in Liquids, Chem. Eng. Sci. 32: 377-383. Ozar, B., Jeong, J. J., Dixit, A., Julia, J. E., Hibiki, T. and Ishii, M., 2008, Flow Structure of\par}
{\small\raggedright\noindent Gas-Liquid Two-phase Flow in an Annulus, Chem. Eng. Sci. 63: 3998-4011. Pai, S. I., 1962, Magnetogasdynamics and Plasma Dynamics , Wien Springer-Verlag. Pai, S. I., 1971, Fundamental Equations of a Mixture of Gas and Small Spherical Solid\par}
{\small\raggedright\noindent Particles from Simple Kinetic Theory, Int. Sym. on Two-phase Systems , Paper 6-6, Haifa, Israel. Pai, S. I., 1972, A New Classification of Two-phase Flows, J. Mach. Phys. Sci. 6: 137-161. Panton, R., 1968, Flow Properties for the Continuum View-point of a Non-Equilibrium Gas-\par}
{\small\raggedright\noindent Particle Mixture, J. Fluid Mech. 31: 273-303. Park, H. S., Lee, T. H., Hibiki, T., Baek W. P. and Ishii, M., 2007, Modeling of the\par}
{\small\raggedright\noindent Condensation Sink Term in an Interfacial Area Transport Equation, Int. J. Heat Mass Transfer 50: 5041-5053. Peebles, F. N. and Garber, H. J., 1953, Studies on the Motion of Gas Bubbles in Liquid, Chem.\par}
{\small\raggedright\noindent Eng. Prog. 49: 88-97. Petrick, M., 1962, A Study of Vapor Carryunder and Associated Problems, Argonne National\par}
{\small\raggedright\noindent Lab. Report , ANL-65-81. Phillips, M. C. and Riddiford, A. C., 1972, Dynamic Contact Angles 2. Velocity and\par}
{\small\raggedright\noindent Relaxation Effects for Various Liquids, J. Colloid Interface Sci. 41 : 77-85. Pierre, C. C. St., 1965, Frequency-response Analysis of Steam Voids to Sinusoidal Power\par}
{\small\raggedright\noindent Modulation in a Thin-walled Boiling Water Coolant Channel, Argonne National Lab. Report , ANL-7041. Plesset, M. S., 1954, On the Stability of Fluid Flows with Spherical Symmetry, J. Appl. Phys.\par}
{\small\raggedright\noindent 25: 96-98. Prigogine, I. and Mazur, P., 1951, On Two-phase Hydrodynamic Formulations and the\par}
{\small\raggedright\noindent Problem of Liquid Helium II, Physica 17: 661-679. Prince, M. J. and Blanch, H. W., 1990, Bubble Coalescence and Break-up in Air-Sparged\par}
{\small\raggedright\noindent Bubble Columns, AIChE J. 36: 1485-1497. Prosperetti, A., 1999, Some Considerations on the Modeling of Disperse Multiphase Flows by\par}
{\small\raggedright\noindent Averaged Equations, JSME Intl. J. , Ser. B, 42: 573-585. Rayleigh, J. W. S., 1917, On the Pressure Developed in a Liquid during the Collapse of a\par}
{\small\raggedright\noindent Spherical Cavity, Philos. Mag. 34: 94-98. Reeks, M. W., 1991, On a Kinetic Equation for the Transport of Particles in Turbulent Flows,\par}
{\small\raggedright\noindent Phys. Fluids A3: 446-456.\par}
{\small\raggedright\noindent Reeks, M. W., 1992, On the Continuum Equations for Dispersed Particles in Nonuniform\par}
{\small\raggedright\noindent Flows, Phys. Fluids A4: 1290-1303. Réocreux, M., Barriere, G. and Vernay, B., 1973, Etude Expérimentale des Débits Critiques\par}
{\small\raggedright\noindent en Écoulement Diphasique Eau-vapeur à Faible Titre sur un Canal à Divergent de 7 Degrés, Cen. G, Rapport TT 115. Réocreux, M., 1974, Contribution à l’étude des Débits Critiques en Éulement Diphasique eau\par}
{\small\raggedright\noindent Vapeur, Ph. D. Thesis , University of Grenoble. Richardson, J. F. and Zaki, W. N., 1954, Sedimentation and Fluidization: Part 1, Trans. Inst.\par}
{\small\raggedright\noindent Chem. Eng. 32: 35-53. Roscoe, R., 1952, The Viscosity of Suspensions of Rigid Spheres, Br. J. Appl. Phys. 3: 267-\par}
{\small\raggedright\noindent 269. Rose, S. C. and Griffith, P., 1965, Flow Properties of Bubbly Mixtures, ASME Paper 65-HT-\par}
{\small\raggedright\noindent 8. Rotta, J. C., 1972, Turbulence Stromungen , B. G. Teubner, Stuttgart, Germany. Rouhani, S. Z. and Becker, K. M., 1963, Measurement of Void Fractions for Flow of Boiling\par}
{\small\raggedright\noindent Heavy Water in a Vertical Round Duct, AE-106, Aktiebolaget Atomenergi, Sweden. Rybczynski, W., 1911, Über die Fortschreitende Bewegung einer Flüssigen Kugel in einem\par}
{\small\raggedright\noindent Zähen Medium, Bull. Int. Acad. Sci. Cracov. 1911A: 40-46. Saffman, P. G., 1965, The Lift on a Small Sphere in a Slow Shear Flow, J. Fluid Mech. 22:\par}
{\small\raggedright\noindent 385-400. Sato, Y., Sadatomi, M. and Sekoguchi, K., 1981, Momentum and Heat Transfer in Two-phase\par}
{\small\raggedright\noindent Bubble Flow - 1. Theory, Int. J. Multiphase Flow 7: 167-177. Schlichting, H., 1979, Boundary Layer Theory , McGraw-Hill Book Co. Schwartz, K., 1954, Investigation of Distribution of Density, Water and Steam Velocity and\par}
{\small\raggedright\noindent of the Pressure Drop in Vertical Horizontal Tubes, VDI Forschungsh. 20, Series B, 445. Schwartz, A. M. and Tejada, S. B., 1972, Studies of Dynamic Contact Angles on Solids, J.\par}
{\small\raggedright\noindent Colloid Interface Sci. 38: 359-375. Scriven, L. E., 1960, Dynamics of Fluid Interface, Equation of Motion for Newtonian Surface\par}
{\small\raggedright\noindent Fluids, Chem. Eng. Sci. 2: 98-108. Serizawa, A., Kataoka, I. and Michiyoshi, I., 1975, Turbulence Structure of Air-water Bubbly\par}
{\small\raggedright\noindent Flow, I, II and III, Int. J. Multiphase Flow 2: 221-259. Serizawa, A. and Kataoka, I., 1988, Phase Distribution in Two-phase Flow, Transient\par}
{\small\raggedright\noindent Phenomena in Multiphase Flow , Hemisphere, Washington DC, pp.179-224. Serizawa, A. and Kataoka, I., 1994, Dispersed Flow-I., Multiphase Science and Technology ,\par}
{\small\raggedright\noindent Vol. 8, Begell House Inc., New York, pp.125-194. Serrin, J., 1959, Handbuch der Physik , Vol.8/I, Springer-Verlag. Sevik, M. and Park, S. H., 1973, The Splitting of Drops and Bubbles by Turbulent Fluid Flow,\par}
{\small\raggedright\noindent J. Fluids Eng. 95: 53-60. Sha, W., Chao, B. T. and Soo, S. L., 1984, Porous-media Formulation for Multiphase Flow\par}
{\small\raggedright\noindent with Heat Transfer, Nucl. Eng. Des. 82: 93-106. Sha, W., Chao, B. T. and Soo, S. L., 1989, Timeand Volume-averaged Conservation\par}
{\small\raggedright\noindent Equations for Multi-phase Flow, Proc. 2nd Int. Symp. Multi-phase Flow and Heat Transfer , Sin Chiu. Sha, W. and Chao, B. T., 2007, Novel Porous-media Formulation for Multiphase Flow\par}
{\small\raggedright\noindent Conservation Equations, Nucl. Eng. Des. 237: 918-942. Situ, R., Hibiki, T., Ishii, M. and Mori, M., 2005, Bubble Lift-off Size in Forced Convective\par}
{\small\raggedright\noindent Subcooled Boiling Flow, Int. J. Heat Mass Transfer 48: 5536-5548. Situ, R., Ishii, M., Hibiki, T., Tu, J. Y., Yeoh, G. H. and Mori, M., 2008, Bubble Departure\par}
{\small\raggedright\noindent Frequency in Forced Convective Subcooled Boiling Flow, Int. J. Heat Mass Transfer 51: 6268-6282.\par}
{\small\raggedright\noindent Slattery, J. C., 1964, Surface – I. Momentum and Moment-of-momentum Balance for Moving\par}
{\small\raggedright\noindent Surfaces, Chem. Eng. Sci. 19: 379-385. Slattery, J. C., 1972, Momentum, Energy and Mass Transfer in Continua , McGraw-Hill Book\par}
{\small\raggedright\noindent Co. Smissaert, G. E., 1963, Two-component Two-phase Flow Parameters for Low Circulation\par}
{\small\raggedright\noindent Rages, Argonne National Lab. Report , ANL-67-55. Soo, S. L., 1967, Fluid Dynamics of Multiphase Systems , Ginn Blaisdell. Sridhar, G. and Katz, J., 1995, Drag and Lift Forces on Microscopic Bubbles Entrained by a\par}
{\small\raggedright\noindent Vortex, Phys. Fluids 7: 389-399. Standart, G., 1964, The Mass, Momentum and Energy Equations for Heterogeneous Flow\par}
{\small\raggedright\noindent Systems, Chem. Eng. Sci. 19: 227-236. Standart, G., 1968, The Second Law of Thermodynamics for Heterogeneous Flow Systems\par}
{\small\raggedright\noindent III. Effect of Conditions of Mechanical Equilibrium and Electroneutrally on Simulataneous Heat and Mass Transfer and Prigogine Theorem, Chem. Eng. Sci. 23: 279- 285. Stefan, J., 1871, Uber das Gleichgewicht und die Bewegung, Insbesondere die Diffusion von\par}
{\small\raggedright\noindent Gasmengen, Sitzgsber, Akad. Wiss. Wien 63: 63-124. Stewart, C. W., 1995, Bubble Interaction in Low-viscosity Liquids, Int. J. Multiphase Flow\par}
{\small\raggedright\noindent 21: 1037-1046. Stokes, G. G., 1851, On the Effect of Internal Friction of Fluids on the Motion of Pendulums,\par}
{\small\raggedright\noindent Trans. Cambr. Phil. Soc. 9, Part II: 8-106 or Coll. Papers III: 55. St. Pierre, C. C., 1965, Frequency-response Analysis of Steam Voids to Sinusoidal Power\par}
{\small\raggedright\noindent Modulation in a Thin-walled Boiling Water Coolant Channel, Argonne National Lab. Report , ANL-7041. Sun, X., Ishii, M. and Kelly, J. M., 2003, Modified Two-fluid Model for the Two-group\par}
{\small\raggedright\noindent Interfacial Area Transport Equation, Annals Nucl. Energy 30: 1601-1622. Sun, X., Kim, S., Ishii, M. and Beus, S. G., 2004a, Modeling of Bubble Coalescence and\par}
{\small\raggedright\noindent Disintegration in Confined Upward Two-phase Flow, Nucl. Eng. Des. 230: 3-26. Sun, X., Kim, S., Ishii, M. and Beus, S. G., 2004b, Model Evaluation of Two-group\par}
{\small\raggedright\noindent Interfacial Area Transport Equation for Confined Upward Flow, Nucl. Eng. Des. 230: 27- 47. Taylor, G. I., 1932, The Viscosity of a Fluid Containing Small Drops of Another Fluid, Proc.\par}
{\small\raggedright\noindent R. Soc. A138: 41-48. Taylor, G. I., 1934, The Formation of Emulsions in Definable Fields of Flow, Proc. Royal\par}
{\small\raggedright\noindent Soc. London A146: 501-523. Teletov, S. G., 1945, Fluid Dynamic Equations for Two-phase Fluids, Soviet Physics Doklady,\par}
{\small\raggedright\noindent Akademii Nauk SSSR 50: 99-102. Teletov, S. G., 1957, On the Problem of Fluid Dynamics of Two-phase Mixtures, I.\par}
{\small\raggedright\noindent Hydrodynamic and Energy Equations, Bull. the Moscow University 2: 15. Theofanous, T. G., Li, G. J. and Dinh, T. N., 2004, Aerobreakup in Rarefied Supersonic Gas\par}
{\small\raggedright\noindent Flows, J. Fluid Eng. 126: 516-527. Thomas, D. G., 1965, Transport Characteristics of Suspension: VIII A Note on Viscosity of\par}
{\small\raggedright\noindent Newtonian Suspensions of Uniform Spherical Particles, J. Colloid Sci. 20: 267-277. Thome, R. J., 1964, Effect of a Transverse Magnetic Field and Vertical Two-phase Flow\par}
{\small\raggedright\noindent through a Rectangular Channel, Argonne National Lab. Report , ANL-6854. Tomiyama, A., Zun, I., Sou, A. and Sakaguchi, T., 1993, Numerical Analysis of Bubble\par}
{\small\raggedright\noindent Motion with the VOF method, Nucl. Eng. Des. 141: 69-82. Tomiyama, A., Sou, A., Zun, I., Kanami, N. and Sakaguchi, T., 1995, Effects of Eötvös\par}
{\small\raggedright\noindent Number and Dimensionless Liquid Volumetric Flux on Lateral Motion of a Bubble in a Laminar Duct Flow, Advances in Multiphase Flow , Elsevier, pp.3-15.\par}
{\small\raggedright\noindent Tomiyama, A., Tamai, H., Zun, I. and Hosokawa, S., 2002, Transverse Migration of Single\par}
{\small\raggedright\noindent Bubbles in Simple Shear Flows, Chem. Eng. Sci. 57: 1849-1858. Todreas, N. E. and Kazimi, M. S., 1990, Nuclear Systems II, Elements of Thermal Hydraulic\par}
{\small\raggedright\noindent Design , Hemisphere Publishing Co., USA. Tong, L. S., 1965, Boiling Heat Transfer and Two-phase Flow , John Wiley and Sons, Inc.,\par}
{\small\raggedright\noindent New York. Truesdell, C. and Toupin, R., 1960, The Classical Field Theories, Handbuch der Physic ,\par}
{\small\raggedright\noindent Vol.3/I, Springer-Verlag. Truesdell, C., 1969, Rational Thermodynamics , McGraw-Hill Book Co. Tsouris, C. and Tavlarides, L. L., 1994, Breakage and Coalescence Models for Drops in\par}
{\small\raggedright\noindent Turbulent Dispersions, AIChE J. 40: 395-406. Tsuchiya, K., Miyahara, T. and Fan, L. S., 1989, Visualization of Bubble-wake Interactions\par}
{\small\raggedright\noindent for a Stream of Bubbles in a Two-dimensional Liquid-solid Fluidized Bed, Int. J. Multiphase Flow 15: 35-49. Vernier, P. and Delhaye, J. M., 1968, General Two-phase Flow Equations Applied to the\par}
{\small\raggedright\noindent Thermohydrodynamics of Boiling Nuclear Reactors, Energie Primaire 4: No.1. Von Karman, 1950, Unpublished Lectures (1950-1951) at Sorbonne and Published by\par}
{\small\raggedright\noindent Nachbar et al. in Quart. Appl. Math. 7: 43 (1959). Wallis, G. B., Steen, D. A., Brenner, S. N. and Turner T. M., 1964, Joint U. S. – Euratom\par}
{\small\raggedright\noindent Research and Development Program, Quarterly Progress Report , January, Dartmouth College. Wallis, G. B., 1969, One-dimensional Two-phase Flow , McGraw-Hill Book Co. Wallis, G. B., 1974, The Terminal Speed of Single Drops or Bubbles in an Infinite Medium,\par}
{\small\raggedright\noindent Int. J. Multiphase Flow 1: 491-511. Wang, S. K., Lee, S. J., Jones Jr., O. C. and Lahey Jr. R. T., 1987, 3-D Turbulence Structure\par}
{\small\raggedright\noindent and Phase Distribution Measurements in Bubbly Two-phase Flows, Int. J. Heat Mass Transfer 13: 327-343. Weatherburn, C. E., 1927, Differential Geometry of Three Dimensions , Cambridge University\par}
{\small\raggedright\noindent Press. Werther, J., 1974, Influence of the Bed Diameter on the Hydrodynamics of Gas Fluidized\par}
{\small\raggedright\noindent Beds, AIChE Symp. Ser. No. 141 70: 53. Whitaker, S., 1968, Introduction to Fluid Mechanics , Prentice-Hall, Inc. Whitaker, S., 1985, A Simple Geometrical Derivation of the Spatial Averaging Theorem,\par}
{\small\raggedright\noindent Chem. Eng. Edu. , pp.18-21, 50-52. White, E. T. and Beardmore, R. H., 1962, The Velocity of Rise of Single Cylindrical Air\par}
{\small\raggedright\noindent Bubbles through Liquid Contained in Vertical Tubes, Chem. Eng. Sci. 17: 351-361. Wu, Q. and Ishii, M., 1996, Interfacial Wave Instability of Co-current Two-phase Flow in\par}
{\small\raggedright\noindent Horizontal Channel, Int. J. Heat Mass Transfer 39: 2067-2075. Wu, Q., Kim, S., Ishii, M. and Beus, S. G., 1998, One-group Interfacial Area Transport in\par}
{\small\raggedright\noindent Vertical Bubbly Flow, Int. J. Heat Mass Transfer 41: 1103-1112. Wundt, H., 1967, Basic Relationships in n-components Diabatic Flow, EUR 3459e . Yoshida, F. and Akita, K., 1965, Performance of Gas Bubble Columns: Volumetric Liquid-\par}
{\small\raggedright\noindent phase Mass Transfer Coefficient and Gas Holdup, AIChE J. 11: 9-13. Zhang, D. Z., 1993, Ensemble Phase Averaged Equations for Multiphase Flows, Ph D. Thesis ,\par}
{\small\raggedright\noindent Johns Hopkins University. Zhang, D. Z., Prosperetti, A., 1994a, Averaged Equations for Inviscid Disperse Two-phase\par}
{\small\raggedright\noindent Flow, J. Fluid Mech. 267: 185-219. Zhang, D. Z., Prosperetti, A., 1994b, Ensemble Phase-averaged Equations for Bubbly Flows,\par}
{\small\raggedright\noindent Phys. Fluids 6: 2956-2970.\par}
{\small\raggedright\noindent Zuber, N., 1961, The Dynamics of Vapor Bubbles in Nonuniform Temperature Fields, Int. J.\par}
{\small\raggedright\noindent Heat Mass Transfer 2: 83-98. Zuber, N., 1964a, On the Dispersed Flow in the Laminar Flow Regime, Chem. Eng. Sci. 19:\par}
{\small\raggedright\noindent 897-917. Zuber, N., 1964b, On the Problem of Hydrodynamic Diffusion in Two-phase Flow Media,\par}
{\small\raggedright\noindent Proc. 2nd All Union Conference on Heat and Mass Transfer , Minsk, USSR, 3: 351. Zuber, N., Staub, F. W. and Bijwaard, G., 1964, Steady State and Transient Void Fraction in\par}
{\small\raggedright\noindent Two-phase Flow Systems, Vol.1, GEAP 5417 . Zuber, N. and Findlay, J. A., 1965, Average Volumetric Concentration in Two-phase Flow\par}
{\small\raggedright\noindent Systems, J. Heat Transfer 87: 453-468. Zuber, N. and Staub, F. W., 1966, Propagation and the Wave Form of the Volumetric\par}
{\small\raggedright\noindent Concentration in Boiling Water Forced Convection Systems under Oscillatory Conditions, Int. J. Heat Mass Transfer 9: 871-895. Zuber, N., 1967, Flow Excursions and Oscillations in Boiling, Two-phase Flow Systems with\par}
{\small\raggedright\noindent Heat Addition, Proc. Symp. Two-phase Flow Dynamics , 1: 1071. Zuber, N. and Dougherty, D. E., 1967, Liquid Metals Challenge to the Traditional Methods in\par}
{\small\raggedright\noindent Two-phase Flow Investigations, Proc. EURATOM Symposium on Two-phase Flow Dynamics , pp.1085. Zuber, N., Staub, F. W., Bijwaard, G. and Kroeger, P. G., 1967, Steady State and Transient\par}
{\small\raggedright\noindent Void Fraction in Two-phase Flow Systems, General Electric Co. Report GEAP-5417, vol.1 Zuber, N., 1971, Personal Communication at Georgia Institute of Technology. Zun, I., 1988, Transition from Wall Void Peaking to Core Void Peaking in Turbulent Bubbly\par}
{\small\raggedright\noindent Flow, Transient Phenomena in Multiphase Flow , Hemisphere, Washington DC, pp.225- 245. Zwick, S. A. and Plesset M. S., 1955, On the Dynamics of Small Vapor Bubbles in Liquids, J.\par}
{\small\raggedright\noindent Math. Phys. 33: 308-330.\par}
\chapter*{符号表}\addcontentsline{toc}{chapter}{符号表}
\markboth{符号表}{符号表}
下表给出两相流平均方程与界面输运模型中最常用符号的中文释义。符号、相指标和平均算子
保持原书约定；其后的完整原书符号表继续保留，以便核对全部模型系数。

\begingroup\small
\renewcommand{\arraystretch}{1.18}
\begin{longtable}{>{\centering\arraybackslash}p{0.23\textwidth}p{0.67\textwidth}}
\toprule 符号 & 中文释义 \\ \midrule
\endfirsthead
\toprule 符号 & 中文释义 \\ \midrule
\endhead
$\alpha_k$ & 第 $k$ 相的体积分数 \\
$\rho_k,\ \boldsymbol{v}_k$ & 第 $k$ 相的密度和速度 \\
$p,\ T$ & 压力和温度 \\
$a_i$ & 界面面积浓度 \\
$\Gamma_k$ & 第 $k$ 相的界面质量传递率 \\
$\boldsymbol{M}_k,\ E_k$ & 第 $k$ 相的界面动量和能量传递项 \\
$D_{Sm},\ D_b$ & Sauter 平均直径和气泡直径 \\
$C_D,\ C_L,\ C_M$ & 阻力、升力和虚质量系数 \\
$C_0$ & 漂移通量模型的分布参数 \\
$V_{gj}$ & 气相相对于混合物体积中心的漂移速度 \\
$j,\ j_g,\ j_f$ & 总体积通量、气相和液相表观速度 \\
$\sigma$ & 表面张力系数 \\
$\mu_k,\ \nu_k$ & 第 $k$ 相的动力黏度和运动黏度 \\
$Re,\ We,\ Eo,\ Fr$ & 雷诺数、韦伯数、厄特沃什数和弗劳德数 \\
$\langle\,\cdot\,\rangle$ & 平均算子；具体平均方式由上下文确定 \\
$\delta_I$ & 界面 delta 函数 \\
\bottomrule
\end{longtable}
\endgroup
\clearpage
\begin{sourceblock}
\includegraphics[page=504,pagebox=mediabox,trim={45.00bp 54.14bp 41.37bp 54.00bp},clip,width=\textwidth,max height=0.80\textheight]{Ishii.pdf}
\end{sourceblock}
\clearpage
\begin{sourceblock}
\includegraphics[page=505,pagebox=mediabox,trim={45.00bp 54.14bp 41.37bp 54.00bp},clip,width=\textwidth,max height=0.80\textheight]{Ishii.pdf}
\end{sourceblock}
\clearpage
\begin{sourceblock}
\includegraphics[page=506,pagebox=mediabox,trim={45.00bp 54.14bp 41.37bp 54.00bp},clip,width=\textwidth,max height=0.80\textheight]{Ishii.pdf}
\end{sourceblock}
\clearpage
\begin{sourceblock}
\includegraphics[page=507,pagebox=mediabox,trim={45.00bp 54.14bp 41.37bp 54.00bp},clip,width=\textwidth,max height=0.80\textheight]{Ishii.pdf}
\end{sourceblock}
\clearpage
\begin{sourceblock}
\includegraphics[page=508,pagebox=mediabox,trim={45.00bp 54.14bp 41.37bp 54.00bp},clip,width=\textwidth,max height=0.80\textheight]{Ishii.pdf}
\end{sourceblock}
\clearpage
\begin{sourceblock}
\includegraphics[page=509,pagebox=mediabox,trim={45.00bp 54.14bp 41.37bp 54.00bp},clip,width=\textwidth,max height=0.80\textheight]{Ishii.pdf}
\end{sourceblock}
\clearpage
\begin{sourceblock}
\includegraphics[page=510,pagebox=mediabox,trim={45.00bp 54.14bp 41.37bp 54.00bp},clip,width=\textwidth,max height=0.80\textheight]{Ishii.pdf}
\end{sourceblock}
\clearpage
\begin{sourceblock}
\includegraphics[page=511,pagebox=mediabox,trim={45.00bp 54.14bp 41.37bp 54.00bp},clip,width=\textwidth,max height=0.80\textheight]{Ishii.pdf}
\end{sourceblock}
\clearpage
\begin{sourceblock}
\includegraphics[page=512,pagebox=mediabox,trim={45.00bp 54.14bp 41.37bp 54.00bp},clip,width=\textwidth,max height=0.80\textheight]{Ishii.pdf}
\end{sourceblock}
\clearpage
\begin{sourceblock}
\includegraphics[page=513,pagebox=mediabox,trim={45.00bp 54.14bp 41.37bp 54.00bp},clip,width=\textwidth,max height=0.80\textheight]{Ishii.pdf}
\end{sourceblock}
\clearpage
\begin{sourceblock}
\includegraphics[page=514,pagebox=mediabox,trim={45.00bp 54.14bp 41.37bp 54.00bp},clip,width=\textwidth,max height=0.80\textheight]{Ishii.pdf}
\end{sourceblock}
\clearpage
\begin{sourceblock}
\includegraphics[page=515,pagebox=mediabox,trim={45.00bp 54.14bp 41.37bp 54.00bp},clip,width=\textwidth,max height=0.80\textheight]{Ishii.pdf}
\end{sourceblock}
\clearpage
\begin{sourceblock}
\includegraphics[page=516,pagebox=mediabox,trim={45.00bp 54.14bp 41.37bp 54.00bp},clip,width=\textwidth,max height=0.80\textheight]{Ishii.pdf}
\end{sourceblock}
\clearpage
\begin{sourceblock}
\includegraphics[page=517,pagebox=mediabox,trim={45.00bp 54.14bp 41.37bp 54.00bp},clip,width=\textwidth,max height=0.80\textheight]{Ishii.pdf}
\end{sourceblock}
\clearpage
\begin{sourceblock}
\includegraphics[page=518,pagebox=mediabox,trim={45.00bp 54.14bp 41.37bp 54.00bp},clip,width=\textwidth,max height=0.80\textheight]{Ishii.pdf}
\end{sourceblock}
\clearpage
\begin{sourceblock}
\includegraphics[page=519,pagebox=mediabox,trim={45.00bp 54.14bp 41.37bp 54.00bp},clip,width=\textwidth,max height=0.80\textheight]{Ishii.pdf}
\end{sourceblock}
\clearpage
\begin{sourceblock}
\includegraphics[page=520,pagebox=mediabox,trim={45.00bp 54.14bp 41.37bp 54.00bp},clip,width=\textwidth,max height=0.80\textheight]{Ishii.pdf}
\end{sourceblock}
\clearpage
\begin{sourceblock}
\includegraphics[page=521,pagebox=mediabox,trim={45.00bp 54.14bp 41.37bp 54.00bp},clip,width=\textwidth,max height=0.80\textheight]{Ishii.pdf}
\end{sourceblock}

\end{document}
